\documentclass[11pt,a4paper,openany]{ctexbook}
\usepackage[a4paper,margin=28mm,headheight=15pt,footskip=12mm]{geometry}
\usepackage{setspace,microtype}
\setstretch{1.12}
\setlength{\parindent}{2em}
\setlength{\parskip}{0.25em}

\usepackage{amsmath,amssymb,amsthm,mathtools,bm,mathrsfs}
\numberwithin{equation}{chapter}
\allowdisplaybreaks[2]
\usepackage{enumitem,array,booktabs,xcolor}
\usepackage{tikz,tikz-cd,pgfplots}
\usetikzlibrary{arrows.meta,calc,matrix,positioning,graphs}
\pgfplotsset{compat=1.18}
\usepackage{hyperref}
\usepackage[nameinlink,noabbrev]{cleveref}
\usepackage{fancyhdr}
\hypersetup{colorlinks=true,linkcolor=blue!60!black,citecolor=teal!60!black,urlcolor=blue!60!black}
\pdfstringdefDisableCommands{\def\C{C}\def\Q{Q}\def\Z{Z}\def\F{F}\def\varphi{phi}\def\Gamma{Gamma}\def\ell{l}\def\Sha{Sha}\def\GL{GL}\def\Cp{Cp}\def\mu{mu}}

\newcommand{\Q}{\mathbb{Q}}
\newcommand{\R}{\mathbb{R}}
\newcommand{\C}{\mathbb{C}}
\newcommand{\Z}{\mathbb{Z}}
\newcommand{\N}{\mathbb{N}}
\newcommand{\F}{\mathbb{F}}
\newcommand{\A}{\mathbb{A}}
\newcommand{\OK}{\mathcal{O}_K}
\newcommand{\OF}{\mathcal{O}_F}
\newcommand{\OL}{\mathcal{O}_L}
\newcommand{\Gal}{\operatorname{Gal}}
\newcommand{\GL}{\operatorname{GL}}
\newcommand{\Spec}{\operatorname{Spec}}
\newcommand{\Tr}{\operatorname{Tr}}
\newcommand{\Nm}{\operatorname{N}}
\newcommand{\Cl}{\operatorname{Cl}}
\newcommand{\Pic}{\operatorname{Pic}}
\newcommand{\Br}{\operatorname{Br}}
\newcommand{\Hom}{\operatorname{Hom}}
\newcommand{\coker}{\operatorname{coker}}
\newcommand{\Ker}{\operatorname{Ker}}
\newcommand{\Cok}{\operatorname{Cok}}
\newcommand{\im}{\operatorname{im}}
\newcommand{\ord}{\operatorname{ord}}
\newcommand{\Frob}{\operatorname{Frob}}
\newcommand{\Fil}{\operatorname{Fil}}
\newcommand{\vol}{\operatorname{vol}}
\newcommand{\Sh}{\operatorname{Sh}}
\newcommand{\Sha}{\mathsf{Sha}}

\theoremstyle{plain}
\newtheorem{theorem}{定理}[chapter]
\newtheorem{proposition}[theorem]{命题}
\newtheorem{lemma}[theorem]{引理}
\newtheorem{corollary}[theorem]{推论}
\theoremstyle{definition}
\newtheorem{definition}[theorem]{定义}
\newtheorem{construction}[theorem]{构造}
\newtheorem{example}[theorem]{例}
\newtheorem{exercise}[theorem]{习题}
\theoremstyle{remark}
\newtheorem{remark}[theorem]{注记}
\newtheorem{notation}[theorem]{记号}
\renewcommand{\proofname}{证明}

\pagestyle{fancy}
\fancyhf{}
\fancyhead[LE,RO]{\thepage}
\fancyhead[LO]{\nouppercase{\leftmark}}
\fancyhead[RE]{\nouppercase{\rightmark}}
\renewcommand{\headrulewidth}{0.4pt}
\renewcommand{\footrulewidth}{0pt}

\begin{document}
\frontmatter
\tableofcontents

\mainmatter
\setcounter{chapter}{-1}
\chapter{预备知识}

本章的作用是给全书建立共同语言。这一章先提醒我们：这一章有些内容也许已经在抽象代数中见过，但在代数数论中它们不是装饰性的预备知识，而是后面每一章都会用到的工具。局部化把一个全局环放到某个素理想附近观察；代数扩张和态射扩张解释“加入方程的根”以后域同态如何变化；Galois 扩张把域扩张的对称性组织成群；迹和范把扩张域中的元素压回基域；有限域、过滤、无穷扩张和特征标则为局部域、类域论与表示论准备语言。

因此阅读第零章时，重点不是背下术语，而是看懂每个工具将来要解决什么问题：一个数域的整数环为什么要局部化，为什么极小多项式控制一个代数数，为什么自同构群能反映中间域，以及为什么有限层信息需要用逆极限整理。

\section*{记号}
\phantomsection
\addcontentsline{toc}{section}{记号}

记号先统一，是为了避免后面每章重复解释。有限集合 $S$ 的元素个数记作 $|S|$；差集 $A\setminus B$ 表示所有属于 $A$ 而不属于 $B$ 的元素。整数环、有理数域、实数域、复数域分别记作 $\Z,\Q,\R,\C$，非负整数集合记作 $\N$，正实数乘法群记作 $\R_{+}^{\times}$。

若 $F$ 是域，则 $F^\times$ 是 $F$ 的非零元乘法群；若 $A$ 是环，则 $A^\times$ 是 $A$ 的单位群，也就是存在乘法逆元的元素全体。有限域记作 $\F_q$，其中 $q$ 是它的元素个数；$m$ 次单位根组成的群记作 $\mu_m$。若 $V$ 是 $F$ 上向量空间，则 $\dim_F V$ 表示维数。若 $E/F$ 是域扩张，则
\[
[E:F]=\dim_F E
\]
称为扩张次数。若 $H$ 是群 $G$ 的子群，则 $[G:H]$ 表示陪集集合 $G/H$ 的元素个数；若 $H$ 是正规子群，记作 $H\triangleleft G$。

映射常写成 $\varphi:x\mapsto y$，意思是 $\varphi(x)=y$。集合 $X$ 上的恒等映射记作 $\operatorname{id}_X$。逻辑符号 $\forall,\exists,\Rightarrow$ 分别表示“对所有”“存在”“推出”。除非特别说明，本章的环都指含单位元的交换环。

\subsection*{读这些记号时要注意什么}

代数数论里同一个对象经常同时有“元素”“理想”“域扩张”“群作用”几种面貌。例如素数 $p$ 是 $\Z$ 的元素，也是素理想 $(p)$ 的生成元；把 $\Z$ 在 $(p)$ 处局部化得到 $\Z_{(p)}$，把 $\Z$ 按 $p$ 进距离完成得到 $\Z_p$，再模掉 $(p)$ 得到剩余域 $\F_p$。这些记号是在同一件事的不同层次之间来回移动的坐标。

\section{局部化}

从 $\Z$ 到 $\Q$ 的过程是局部化的原型：为了让非零整数都能作分母，我们把每个 $n\ne0$ 都变成可逆元。一般环中不一定要倒转所有非零元。局部化的思想是指定一批允许作分母的元素，只强迫这些元素变成单位，而不改变其他地方的信息。代数数论中最常用的情形，是在一个素理想附近观察环：不属于这个素理想的元素都被视为“局部上不消失”，于是可以倒过来。

\begin{definition}
设 $R$ 是含 $1$ 的交换环。子集 $S\subset R$ 称为乘性子集，如果 $1\in S$，并且 $s,t\in S$ 时 $st\in S$。本节默认 $0\notin S$。在 $R\times S$ 上定义关系
\[
(a,r)\sim(b,s)\quad\Longleftrightarrow\quad \text{存在 }t\in S\text{ 使 }t(as-br)=0.
\]
等价类记作 $a/r$。所有等价类组成的集合记作 $S^{-1}R$，称为 $R$ 关于 $S$ 的分式环或局部化。
\end{definition}

定义里额外出现的 $t$ 是为了处理有零因子的环。若 $R$ 是整环，则 $t(as-br)=0$ 与 $as=br$ 等价，分式相等就退回熟悉的交叉相乘。若 $R$ 有零因子，不能直接消去 $t$，因此必须把“乘上某个允许分母以后相等”写进等价关系。

\begin{lemma}
上面的关系 $\sim$ 是等价关系。
\end{lemma}

\begin{proof}
自反性来自 $1(ar-ra)=0$。若 $(a,r)\sim(b,s)$，则有 $t(as-br)=0$。把等式乘以 $-1$，得到 $t(br-as)=0$，所以 $(b,s)\sim(a,r)$，这给出对称性。

传递性是唯一需要计算的部分。设 $(a,r)\sim(b,s)$ 且 $(b,s)\sim(c,u)$。取 $t_1,t_2\in S$ 使
\[
t_1(as-br)=0,\qquad t_2(bu-cs)=0.
\]
把第一式乘以 $t_2u$，把第二式乘以 $t_1r$，得到
\[
t_1t_2u(as-br)=0,\qquad t_1t_2r(bu-cs)=0.
\]
两式相加，中间的 $bru$ 项抵消：
\[
t_1t_2s(au-cr)=0.
\]
因为 $t_1t_2s\in S$，所以 $(a,r)\sim(c,u)$。因此 $\sim$ 是等价关系。
\end{proof}

\begin{proposition}
集合 $S^{-1}R$ 在运算
\[
\frac a r+\frac b s=\frac{as+br}{rs},\qquad
\frac a r\cdot\frac b s=\frac{ab}{rs}
\]
下成为交换环。映射
\[
\iota:R\longrightarrow S^{-1}R,\qquad a\longmapsto a/1
\]
是环同态，并且每个 $\iota(s)$ 都是 $S^{-1}R$ 的单位。
\end{proposition}

\begin{proof}
先验证运算不依赖代表元。设 $a/r=a'/r'$ 与 $b/s=b'/s'$。按定义，存在 $u,v\in S$ 使
\[
u(ar'-a'r)=0,\qquad v(bs'-b's)=0.
\]
对加法，两种代表给出的分子差为
\[
(as+br)r's'-(a's'+b'r')rs
=ss'(ar'-a'r)+rr'(bs'-b's).
\]
乘以 $uv$ 后，右边第一项含有 $u(ar'-a'r)$，第二项含有 $v(bs'-b's)$，所以整体为 $0$。这说明
\[
\frac{as+br}{rs}=\frac{a's'+b'r'}{r's'}.
\]
对乘法，两种代表给出的分子差为
\[
ab r's'-a'b'rs=bs'(ar'-a'r)+a'r(bs'-b's).
\]
同样乘以 $uv$ 后为 $0$，故 $ab/(rs)=a'b'/(r's')$。

运算良定以后，交换律、结合律和分配律都从 $R$ 中对应的律逐项传到等价类上。例如分配律可直接化为
\[
\frac a r\left(\frac b s+\frac c u\right)
=\frac{a(bu+cs)}{rsu}
=\frac{ab}{rs}+\frac{ac}{ru}.
\]
元素 $1/1$ 是乘法单位，$0/1$ 是加法单位，$a/r$ 的加法逆元为 $(-a)/r$。于是 $S^{-1}R$ 是交换环。对 $s\in S$，$(s/1)(1/s)=1/1$，所以 $s/1$ 是单位。
\end{proof}

\begin{proposition}
局部化满足如下普遍性质。若 $\varphi:R\to A$ 是环同态，并且每个 $\varphi(s)$ 在 $A$ 中可逆，则存在唯一环同态 $\widetilde\varphi:S^{-1}R\to A$，使
\[
\widetilde\varphi(a/s)=\varphi(a)\varphi(s)^{-1}.
\]
并且 $\widetilde\varphi\circ\iota=\varphi$。
\end{proposition}

\begin{proof}
先证明公式良定。若 $a/s=b/r$，则存在 $t\in S$ 使 $t(ar-bs)=0$。施加 $\varphi$ 得
\[
\varphi(t)\bigl(\varphi(a)\varphi(r)-\varphi(b)\varphi(s)\bigr)=0.
\]
由于 $\varphi(t)$ 可逆，可消去它，得到
\[
\varphi(a)\varphi(r)=\varphi(b)\varphi(s).
\]
再右乘可逆元 $\varphi(s)^{-1}\varphi(r)^{-1}$，得到
\[
\varphi(a)\varphi(s)^{-1}=\varphi(b)\varphi(r)^{-1}.
\]
因此 $\widetilde\varphi$ 不依赖分式代表。

它保持加法和乘法由同分母计算直接得到。例如
\[
\widetilde\varphi(a/s+b/r)
=\varphi(ar+bs)\varphi(sr)^{-1}
=\varphi(a)\varphi(s)^{-1}+\varphi(b)\varphi(r)^{-1}.
\]
唯一性来自每个 $a/s$ 都可写成 $(a/1)(s/1)^{-1}$；任何延拓 $\varphi$ 且使 $s/1$ 可逆的同态，都必须把它送到 $\varphi(a)\varphi(s)^{-1}$。
\end{proof}

\begin{definition}
具有唯一极大理想的环称为局部环。若 $\mathfrak p$ 是 $R$ 的素理想，取 $S=R\setminus\mathfrak p$，则
\[
R_{\mathfrak p}=S^{-1}R
\]
称为 $R$ 在 $\mathfrak p$ 处的局部化。
\end{definition}

\begin{proposition}
若 $\mathfrak p$ 是 $R$ 的素理想，则 $R_{\mathfrak p}$ 是局部环，其唯一极大理想为
\[
\mathfrak pR_{\mathfrak p}=\{a/s:a\in\mathfrak p,\ s\notin\mathfrak p\}.
\]
\end{proposition}

\begin{proof}
先说明 $\mathfrak pR_{\mathfrak p}$ 是理想。若 $a/s,b/r$ 的分子都在 $\mathfrak p$，则
\[
\frac a s+\frac b r=\frac{ar+bs}{sr}.
\]
因为 $a,b\in\mathfrak p$ 且 $\mathfrak p$ 是理想，所以 $ar+bs\in\mathfrak p$。再乘任意 $c/u$ 得
\[
\frac c u\cdot\frac a s=\frac{ca}{us},
\]
分子仍在 $\mathfrak p$。

若 $a/s\notin\mathfrak pR_{\mathfrak p}$，则 $a\notin\mathfrak p$。于是 $a\in S$，所以 $a/s$ 的逆元是 $s/a$。因此所有不在 $\mathfrak pR_{\mathfrak p}$ 中的元素都是单位。任何真理想都不能含单位，故每个真理想都包含在 $\mathfrak pR_{\mathfrak p}$ 中。这说明 $\mathfrak pR_{\mathfrak p}$ 是唯一极大理想。
\end{proof}

\begin{example}
若 $R$ 是整环，取 $S=R\setminus\{0\}$，则 $S^{-1}R$ 是 $R$ 的分式域，记作 $\operatorname{Frac}(R)$。这是 $\Z$ 生成 $\Q$ 的一般形式。
\end{example}

\begin{example}
若 $p$ 是素数，则
\[
\Z_{(p)}=\left\{a/b\in\Q:p\nmid b\right\}.
\]
它与 $\Z_p$ 不同：$\Z_{(p)}$ 是代数式的局部化，只允许不含 $p$ 的整数作分母；$\Z_p$ 是按 $p$ 进距离补上极限后的完成。也要区分 $\Z[1/p]$：它把 $p$ 变成单位，形如 $a/p^n$；而 $\Z_{(p)}$ 把所有不含 $p$ 的整数变成单位，却保留 $p$ 为非单位。
\end{example}

第一章研究 Dedekind 环时，理想分解可以在每个素理想处局部观察；第三章和第十章讨论完备域、赋值环时，也会反复使用“先局部化，再完成”的思路。局部化让全局问题分解成素点附近的问题，是代数数论的基本降维手段。

\section{代数扩张}

代数扩张讨论的是“把方程的根加入一个域”。如果 $\alpha$ 满足某个 $F$ 系数非零多项式，那么 $\alpha$ 在 $F$ 上不是自由变量，而是受一个代数方程约束的元素。这个约束由极小多项式记录。后面说数域、分裂域、Galois 群、迹和范，都以本节语言为基础。

\begin{definition}
设 $F$ 是域。以 $F[X]$ 表示系数在 $F$ 中的一元多项式环，以 $F(X)$ 表示有理函数域。若 $F$ 是域 $E$ 的子域，则称 $E/F$ 为域扩张。此时 $E$ 是 $F$-向量空间，维数
\[
[E:F]=\dim_F E
\]
称为扩张次数。若 $[E:F]<\infty$，称 $E/F$ 为有限扩张。
\end{definition}

\begin{definition}
设 $E/F$ 是域扩张，$\alpha\in E$。若存在非零多项式 $f(X)\in F[X]$ 使 $f(\alpha)=0$，则称 $\alpha$ 是 $F$ 上的代数元素。若 $E$ 中每个元素都在 $F$ 上代数，则称 $E/F$ 为代数扩张。
\end{definition}

\begin{proposition}
设 $\alpha\in E$ 是 $F$ 上的代数元素。令
\[
I(\alpha)=\{f(X)\in F[X]:f(\alpha)=0\}.
\]
则存在唯一首一不可约多项式 $m_{\alpha,F}(X)\in F[X]$ 生成 $I(\alpha)$。它称为 $\alpha$ 在 $F$ 上的极小多项式。
\end{proposition}

\begin{proof}
集合 $I(\alpha)$ 是 $F[X]$ 的理想。它非零，因为 $\alpha$ 代数。由于 $F[X]$ 是主理想整环，存在非零多项式 $m$ 使 $I(\alpha)=(m)$。把 $m$ 乘以其首项系数的倒数，可取 $m$ 为首一多项式。首一生成元唯一：若 $m,m'$ 都首一且生成同一主理想，则 $m=um'$，其中 $u\in F^\times$，比较首项系数得 $u=1$。

再证不可约。若 $m=gh$ 且 $g,h$ 次数都正，则
\[
0=m(\alpha)=g(\alpha)h(\alpha).
\]
域中没有零因子，因此 $g(\alpha)=0$ 或 $h(\alpha)=0$。这说明 $g\in I(\alpha)$ 或 $h\in I(\alpha)$，从而 $m$ 整除 $g$ 或 $h$。但 $g,h$ 的次数都小于 $m$，矛盾。因此 $m$ 不可约。
\end{proof}

\begin{proposition}
若 $E/F$ 是有限扩张，则 $E/F$ 是代数扩张。
\end{proposition}

\begin{proof}
取任意 $\alpha\in E$，设 $[E:F]=n$。向量空间 $E$ 中 $n+1$ 个元素
\[
1,\alpha,\alpha^2,\dots,\alpha^n
\]
必线性相关。因此存在不全为零的 $a_0,\dots,a_n\in F$，使
\[
a_0+a_1\alpha+\cdots+a_n\alpha^n=0.
\]
令 $f(X)=a_0+a_1X+\cdots+a_nX^n$。这是 $F[X]$ 中非零多项式，且 $f(\alpha)=0$。所以 $\alpha$ 在 $F$ 上代数。由于 $\alpha$ 任意，扩张 $E/F$ 为代数扩张。
\end{proof}

\begin{construction}
设 $p(X)\in F[X]$ 不可约。理想 $(p(X))$ 是 $F[X]$ 的极大理想，所以商环
\[
F[X]/(p(X))
\]
是域。记 $\xi=X+(p(X))$。在这个商域中有 $p(\xi)=0$。因此我们构造出一个包含 $F$ 的域，并在其中加入了 $p$ 的一个根。
\end{construction}

\begin{proposition}
若 $p(X)\in F[X]$ 不可约，$\alpha$ 是某个扩张域中满足 $p(\alpha)=0$ 的元素，则
\[
F(\alpha)\simeq F[X]/(p(X)),
\]
同构由 $X\mapsto\alpha$ 给出。若 $\deg p=n$，则 $1,\alpha,\ldots,\alpha^{n-1}$ 是 $F(\alpha)$ 的一组 $F$-基。
\end{proposition}

\begin{proof}
定义代入同态
\[
\operatorname{ev}_{\alpha}:F[X]\to F(\alpha),\qquad f(X)\mapsto f(\alpha).
\]
它的像是 $F[\alpha]$。因为 $p(\alpha)=0$，有 $(p)\subset\ker(\operatorname{ev}_{\alpha})$。多项式 $p$ 不可约，所以 $(p)$ 是极大理想；而核是包含 $(p)$ 的真理想，因为 $1$ 不在核中。因此 $\ker(\operatorname{ev}_{\alpha})=(p)$。同构定理给出
\[
F[X]/(p)\simeq F[\alpha].
\]
左边是域，所以 $F[\alpha]$ 已经是域，因而等于 $F(\alpha)$。

任意多项式除以 $p$ 后可写成商乘 $p$ 加次数小于 $n$ 的余式。代入 $\alpha$，商乘 $p$ 的部分为零，所以 $F(\alpha)$ 中每个元素都是 $1,\alpha,\dots,\alpha^{n-1}$ 的线性组合。若有非平凡线性关系，则得到次数小于 $n$ 的非零多项式使 $\alpha$ 为根，与 $p$ 是极小多项式矛盾。因此这些元素线性无关。
\end{proof}

\begin{definition}
若 $E,K$ 是域，映射 $\varphi:E\to K$ 称为域同态，若它保持加法和乘法。单射域同态称为单态射或嵌入，双射域同态称为同构，域到自身的同构称为自同构。若 $F$ 同时嵌入在 $E$ 和 $K$ 中，并且 $\varphi(a)=a$ 对所有 $a\in F$ 成立，则 $\varphi$ 称为 $F$-态射。
\end{definition}

\begin{lemma}
设 $E/F$ 为代数扩张，$\sigma:E\to E$ 是固定 $F$ 的单态射。则 $\sigma$ 是 $E$ 的自同构。
\end{lemma}

\begin{proof}
只需证明 $\sigma$ 满射。任取 $y\in E$，令 $m_y(X)$ 为 $y$ 在 $F$ 上的极小多项式。考虑 $m_y$ 在 $E$ 中的根集合
\[
R_y=\{x\in E:m_y(x)=0\}.
\]
该集合有限，因为域上的非零多项式根数不超过次数。若 $x\in R_y$，由于 $\sigma$ 固定 $F$，有
\[
m_y(\sigma(x))=\sigma(m_y(x))=0.
\]
所以 $\sigma$ 把 $R_y$ 单射地映入自身。有限集合上的单射是满射，于是存在 $x\in R_y$ 使 $\sigma(x)=y$。任意 $y$ 都在 $\sigma$ 的像中，因此 $\sigma$ 满射。
\end{proof}

\begin{theorem}
设 $\varphi_1,\dots,\varphi_n:E\to K$ 是两两不同的域单态射。把每个 $\varphi_i$ 看作从 $E$ 到 $K$ 的函数，则它们在 $K$ 上线性无关：若 $y_i\in K$ 且
\[
\sum_{i=1}^n y_i\varphi_i(x)=0\qquad(\forall x\in E),
\]
则 $y_1=\cdots=y_n=0$。
\end{theorem}

\begin{proof}
反设存在非平凡线性关系，取其中项数最少的一条：
\[
y_1\varphi_1+\cdots+y_r\varphi_r=0,\qquad y_i\ne0.
\]
这里等号表示对每个 $w\in E$ 都有 $\sum_i y_i\varphi_i(w)=0$。若 $r=1$，代入 $w=1$ 得 $y_1=0$，矛盾，所以 $r>1$。

由于 $\varphi_1\ne\varphi_2$，取 $x\in E$ 使 $\varphi_1(x)\ne\varphi_2(x)$。把原关系应用到 $xw$：
\[
\sum_{i=1}^r y_i\varphi_i(x)\varphi_i(w)=0.
\]
再把原关系乘以常数 $\varphi_1(x)$：
\[
\sum_{i=1}^r y_i\varphi_1(x)\varphi_i(w)=0.
\]
两式相减，得到
\[
\sum_{i=2}^r y_i(\varphi_i(x)-\varphi_1(x))\varphi_i(w)=0.
\]
其中 $i=2$ 的系数非零，所以这是一条项数更少的非平凡线性关系，矛盾。
\end{proof}

\begin{definition}
若域 $K$ 中每个非常数多项式都有根，等价地说每个非常数多项式都能分解为一次因子的乘积，则称 $K$ 为代数闭域。
\end{definition}

\begin{theorem}
设 $F$ 是域。
\begin{enumerate}[label=(\arabic*)]
\item 存在域扩张 $E/F$，使 $E$ 是代数闭域。
\item 存在代数扩张 $F^{\mathrm{alg}}/F$，使 $F^{\mathrm{alg}}$ 是代数闭域。
\end{enumerate}
\end{theorem}

\begin{proof}
先构造一个代数闭扩张。对 $F[X]$ 中每个非常数多项式 $f$ 引入一个新变量 $T_f$，并考虑多项式环 $F[T_f]$。取由所有 $f(T_f)$ 生成的理想。这个理想是真理想：任取有限多个多项式，只需逐次通过不可约因子的商域加入根，就能找到一个域使这些有限关系同时成立，因此有限生成部分不会生成 $1$。由 Zorn 引理或紧性论证，整个理想包含在一个极大理想中。商掉该极大理想后得到一个扩张域，其中 $F$ 上每个非常数多项式都有一个根。重复这个过程并取递增并，可得到代数闭扩张 $E/F$。

在这个代数闭扩张 $E$ 中，令
\[
F^{\mathrm{alg}}=\{x\in E:x\text{ 在 }F\text{ 上代数}\}.
\]
代数元素对加、减、乘和非零除法封闭：若 $x,y$ 都在 $F$ 上代数，则 $F(x,y)/F$ 是有限扩张，故 $x\pm y,xy$ 和 $x/y$ 都在一个有限扩张中，从而在 $F$ 上代数。因此 $F^{\mathrm{alg}}$ 是子域，且 $F^{\mathrm{alg}}/F$ 是代数扩张。

还要证明 $F^{\mathrm{alg}}$ 代数闭。设 $g(X)\in F^{\mathrm{alg}}[X]$ 非常数。因 $E$ 代数闭，$g$ 在 $E$ 中有根 $\beta$。多项式 $g$ 的有限多个系数生成有限代数扩张 $K/F$。元素 $\beta$ 是 $K$ 上代数的，所以 $K(\beta)/K$ 有限；于是 $K(\beta)/F$ 有限，$\beta$ 在 $F$ 上代数。故 $\beta\in F^{\mathrm{alg}}$。这说明 $F^{\mathrm{alg}}$ 中每个非常数多项式都有根。
\end{proof}

有理数域 $\Q$ 的代数闭包记作 $\Q^{\mathrm{alg}}$。它的元素称为代数数，$\Q^{\mathrm{alg}}$ 的有限子扩张称为数域。代数数论研究的基本对象就是这些数域及其中的整数环、理想、局部化、Galois 群和 $L$ 函数。

\section{态射扩张}

本节研究一个基本问题：若已经有一个基域嵌入 $\sigma:F\to L$，并且我们把某些代数元加入 $F$，那么 $\sigma$ 能否延拓到更大的域？延拓以后有多少种选择？这个问题是后面 Galois 理论的计数核心。一个扩张的“共轭根”越完整，延拓越多；若出现重根或缺根，延拓数就会下降。

\subsection{态射扩张一}

先看最简单的情形：只加入一个代数元 $\alpha$。直觉是，延拓 $\tau:F(\alpha)\to L$ 一旦固定了 $\tau|_F=\sigma$，就只剩一个自由选择：$\alpha$ 应当送到哪里。因为 $\alpha$ 满足极小多项式，所以 $\tau(\alpha)$ 必须满足对应变换后的多项式。

\begin{lemma}
设 $F$ 是域，$L$ 是代数闭域，$\sigma:F\to L$ 是单态射。考虑单代数扩张 $F(\alpha)/F$。则：
\begin{enumerate}[label=(\arabic*)]
\item $\sigma$ 可延拓为单态射 $\tau:F(\alpha)\to L$。
\item 设
\[
S(\sigma)=\{\tau:F(\alpha)\to L:\tau\text{ 是单态射且 }\tau|_F=\sigma\}.
\]
若 $m_{\alpha,F}(X)$ 是 $\alpha$ 在 $F$ 上的极小多项式，则 $|S(\sigma)|$ 等于多项式 $\sigma(m_{\alpha,F})(X)$ 在 $L$ 中互不相同的根的个数。
\end{enumerate}
\end{lemma}

\begin{proof}
把 $m=m_{\alpha,F}$ 写成 $m(X)=a_0+a_1X+\cdots+a_nX^n$，并定义
\[
\sigma(m)(X)=\sigma(a_0)+\sigma(a_1)X+\cdots+\sigma(a_n)X^n\in L[X].
\]
因为 $L$ 代数闭，$\sigma(m)$ 在 $L$ 中至少有一个根，记为 $\beta$。定义从多项式环到 $L$ 的同态
\[
\Phi_\beta:F[X]\longrightarrow L,\qquad f(X)\longmapsto \sigma(f)(\beta).
\]
这里 $\sigma(f)$ 表示对 $f$ 的每个系数施加 $\sigma$。由于 $\sigma(m)(\beta)=0$，理想 $(m)$ 落入 $\Phi_\beta$ 的核中，于是 $\Phi_\beta$ 诱导同态
\[
F[X]/(m)\longrightarrow L.
\]
而 $F(\alpha)\simeq F[X]/(m)$，其中 $X$ 的像对应 $\alpha$。因此得到延拓 $\tau_\beta:F(\alpha)\to L$，并且 $\tau_\beta(\alpha)=\beta$。

反过来，若 $\tau:F(\alpha)\to L$ 延拓 $\sigma$，则
\[
0=\tau(m(\alpha))=\sigma(m)(\tau(\alpha)),
\]
所以 $\tau(\alpha)$ 必是 $\sigma(m)$ 的根。并且 $F(\alpha)$ 由 $F$ 和 $\alpha$ 生成，因此 $\tau$ 完全由 $\tau|_F=\sigma$ 与 $\tau(\alpha)$ 决定。于是不同的延拓与 $\sigma(m)$ 的不同根一一对应。
\end{proof}

这个引理要读出两层意思。第一，延拓存在性靠的是目标域 $L$ 代数闭，因为我们必须保证变换后的极小多项式有根。第二，延拓数量不一定等于 $\deg m$，因为多项式可能有重根；这个缺口正是后面“可分性”的来源。

\begin{proposition}
设 $E/F$ 为代数扩张，$L$ 为代数闭域，$\sigma:F\to L$ 为单态射。则：
\begin{enumerate}[label=(\arabic*)]
\item $\sigma$ 可延拓为单态射 $\tau:E\to L$。
\item 若 $E$ 为代数闭域，$L/\sigma(F)$ 为代数扩张，且 $\tau:E\to L$ 是 $\sigma$ 的一个延拓，则 $\tau$ 是同构。
\end{enumerate}
\end{proposition}

\begin{proof}
先证明延拓存在。把所有部分延拓组成集合：一个元素是二元组 $(K,\tau_K)$，其中 $F\subset K\subset E$，$\tau_K:K\to L$ 是延拓 $\sigma$ 的单态射。按包含关系排序：$(K,\tau_K)\le (K',\tau_{K'})$ 若 $K\subset K'$ 且 $\tau_{K'}|_K=\tau_K$。

任意全序链的并仍给出一个部分延拓：定义域取所有 $K$ 的并，映射在每个 $K$ 上按 $\tau_K$ 给出；链条件保证定义一致。因此由 Zorn 引理，有极大部分延拓 $(K_0,\tau_0)$。若 $K_0\ne E$，取 $\alpha\in E\setminus K_0$。因为 $E/F$ 代数，$\alpha$ 在 $K_0$ 上仍代数。由态射扩张引理，$\tau_0$ 可延拓到 $K_0(\alpha)$，这与极大性矛盾。故 $K_0=E$。

再证明第二点。由于 $\tau$ 是单态射，$\tau(E)$ 是 $L$ 的子域，并且 $\tau(E)$ 代数闭：若 $g(X)\in \tau(E)[X]$，把系数用 $\tau$ 拉回到 $E$，得到 $E[X]$ 中多项式；因为 $E$ 代数闭，它有根，施加 $\tau$ 后得到 $g$ 在 $\tau(E)$ 中的根。另一方面，$L/\sigma(F)$ 是代数扩张，而 $\tau(E)$ 含有 $\sigma(F)$ 且代数闭。任取 $y\in L$，它在 $\sigma(F)$ 上代数，因此也在 $\tau(E)$ 上代数；它的极小多项式在代数闭域 $\tau(E)$ 中分裂，故 $y\in\tau(E)$。于是 $\tau(E)=L$，$\tau$ 是同构。
\end{proof}

\begin{corollary}
若 $E/F$ 与 $E'/F$ 都是代数扩张，且 $E,E'$ 都是代数闭域，则 $E$ 与 $E'$ 在 $F$ 上同构。
\end{corollary}

\begin{proof}
把恒等嵌入 $F\to E'$ 延拓为 $F$-单态射 $\tau:E\to E'$。由于 $E$ 代数闭且 $E'/F$ 代数，上一命题说明 $\tau$ 是同构。
\end{proof}

因此得到：任一域 $F$ 的代数闭包 $F^{\mathrm{alg}}$ 在 $F$-同构意义下唯一。它不是唯一地坐在某个大域中，而是“除同构外唯一”。在这里，数域就是 $\Q^{\mathrm{alg}}$ 的有限子扩张，代数数就是 $\Q^{\mathrm{alg}}$ 的元素。

\paragraph{分裂域}
设 $f(X)\in F[X]$ 是次数 $n\ge1$ 的多项式。域扩张 $K/F$ 称为 $f$ 在 $F$ 上的分裂域，若：
\begin{enumerate}[label=(\arabic*)]
\item $f$ 在 $K$ 中完全分裂，即 $f(X)=a(X-\alpha_1)\cdots(X-\alpha_n)$，其中 $a\in F$、$\alpha_i\in K$；
\item $K$ 由这些根生成，即 $K=F(\alpha_1,\dots,\alpha_n)$。
\end{enumerate}
第一条说根已经全部放进来，第二条说没有额外放入不需要的元素。分裂域的元素都可写成关于根的有理表达式
\[
\frac{p(\alpha_1,\dots,\alpha_n)}{q(\alpha_1,\dots,\alpha_n)},\qquad q(\alpha_1,\dots,\alpha_n)\ne0.
\]

\begin{proposition}
设 $E$ 与 $K$ 都是 $f(X)\in F[X]$ 在 $F$ 上的分裂域。则：
\begin{enumerate}[label=(\arabic*)]
\item $E$ 与 $K$ 在 $F$ 上同构。
\item 若 $K^{\mathrm{alg}}$ 是 $K$ 的代数闭包，$\tau:E\to K^{\mathrm{alg}}$ 是 $F$-单态射，则 $\tau(E)=K$。
\end{enumerate}
\end{proposition}

\begin{proof}
由上面的延拓命题，$F$ 的恒等嵌入可延拓为 $F$-单态射
\[
\tau:E\longrightarrow K^{\mathrm{alg}}.
\]
设 $E=F(\alpha_1,\dots,\alpha_n)$，其中 $\alpha_i$ 是 $f$ 的全部根。因为 $\tau$ 固定 $F$，有
\[
f(\tau(\alpha_i))=\tau(f(\alpha_i))=0.
\]
所以每个 $\tau(\alpha_i)$ 仍是 $f$ 的根。由于 $K$ 已经包含 $f$ 的全部根，$\tau(\alpha_i)\in K$，从而 $\tau(E)\subset K$。另一方面，$\tau(E)$ 也包含 $f$ 的全部根的一个完整集合，并且由这些根生成，所以它本身是一个分裂域。若把同样论证用于反向延拓 $K\to E^{\mathrm{alg}}$，得到两个分裂域互相嵌入并由同一组根生成；因此 $E$ 与 $K$ 在 $F$ 上同构。

对第二点，$\tau(E)$ 是 $K^{\mathrm{alg}}$ 中由 $f$ 的全部根生成的 $F$-子域，而 $K$ 也是由同一批根生成的 $F$-子域。故 $\tau(E)=K$。
\end{proof}

同样地，一组多项式 $\{f_i(X):i\in I\}$ 的分裂域，就是在代数闭包中由每个 $f_i$ 的全部根生成的最小子域。上面的证明逐个多项式照搬，说明这样的分裂域也在 $F$-同构意义下唯一，并且任何 $F$-嵌入都把它送到同一个分裂域。

\subsection{态射扩张二}

本小节把“延拓是否落回自身”整理成正规扩张。正规性关心的不是根是否重合，而是根是否全部在域里。

\paragraph{正规扩张的定义}
取定 $F$ 的代数闭包 $F^{\mathrm{alg}}$，设 $F\subset E\subset F^{\mathrm{alg}}$ 为代数扩张。若每个 $F$-单态射 $\sigma:E\to F^{\mathrm{alg}}$ 都满足 $\sigma(E)=E$，则称 $E/F$ 为正规扩张。

直观地说，正规扩张不允许“一个根在 $E$ 中，而它的共轭根跑到 $E$ 外面”。例如 $\Q(\sqrt2)/\Q$ 是正规扩张，因为 $X^2-2$ 的两个根 $\pm\sqrt2$ 都在其中；而 $\Q(\sqrt[3]{2})/\Q$ 不是正规扩张，因为 $X^3-2$ 的另外两个复根不在这个实域里。

\begin{proposition}
设 $F\subset E\subset F^{\mathrm{alg}}$。以下三个条件等价：
\begin{enumerate}[label=(\arabic*)]
\item $E/F$ 是正规扩张。
\item $E$ 是一组 $F[X]$ 中多项式的分裂域。
\item 若 $f(X)\in F[X]$ 是不可约多项式，且 $f$ 在 $E$ 中有一个根，则 $f$ 在 $E$ 中完全分裂。
\end{enumerate}
\end{proposition}

\begin{proof}
先证 $(2)\Rightarrow(1)$。设 $E$ 是多项式组 $\{f_i\}$ 的分裂域，$\sigma:E\to F^{\mathrm{alg}}$ 是 $F$-单态射。由于 $\sigma$ 固定 $F$，它把每个 $f_i$ 的根送到同一个 $f_i$ 的根。因此 $\sigma(E)$ 仍由这些多项式的全部根生成，故 $\sigma(E)=E$。

再证 $(1)\Rightarrow(3)$。设 $f\in F[X]$ 不可约，$\alpha\in E$ 是它的一个根。若 $\beta\in F^{\mathrm{alg}}$ 是 $f$ 的任意根，则由态射扩张引理，存在 $F$-同构 $F(\alpha)\to F(\beta)$ 把 $\alpha$ 送到 $\beta$。再由延拓命题，这个同构可延拓为 $F$-单态射 $E\to F^{\mathrm{alg}}$。正规性给出该单态射的像仍为 $E$，所以 $\beta\in E$。于是 $f$ 的所有根都在 $E$ 中。

最后证 $(3)\Rightarrow(2)$。对每个 $\alpha\in E$，取它在 $F$ 上的极小多项式 $m_{\alpha,F}$。由条件 (3)，每个 $m_{\alpha,F}$ 都在 $E$ 中完全分裂。因为集合 $\{\alpha:\alpha\in E\}$ 本身已经包含 $E$ 的每一个元素，所以由它们生成的 $F$-子域就是 $E$；而每个 $\alpha$ 又是对应多项式 $m_{\alpha,F}$ 的一个根，条件 (3) 还保证该多项式的其余根也都在 $E$ 中。于是 $E$ 正是这些极小多项式的全部根生成的域。因此 $E$ 是多项式组 $\{m_{\alpha,F}:\alpha\in E\}$ 的分裂域。
\end{proof}

\subsection{态射扩张三}

现在计数嵌入。若 $E/F$ 是代数扩张，$L$ 是代数闭域，$\sigma:F\to L$ 是单态射，记
\[
S_E(\sigma)=\{\tau:E\to L:\tau\text{ 是单态射且 }\tau|_F=\sigma\}.
\]
若换另一个代数闭目标 $L'$ 和另一个基域嵌入 $\sigma':F\to L'$，这个集合的基数不变。原因是两个代数闭包在基域上同构，延拓集合可通过该同构互相转移。因此把这个共同基数记作
\[
\langle E:F\rangle_s,
\]
称为 $E/F$ 的嵌入数，有限扩张时它将等于可分次数。

\begin{proposition}
设 $E\supset K\supset F$ 为域扩张。则：
\begin{enumerate}[label=(\arabic*)]
\item $\langle E:F\rangle_s=\langle E:K\rangle_s\langle K:F\rangle_s$。
\item 若 $[E:F]<\infty$，则 $\langle E:F\rangle_s\le [E:F]$。
\item 若 $[E:F]<\infty$，则 $\langle E:F\rangle_s=[E:F]$ 当且仅当 $\langle E:K\rangle_s=[E:K]$ 且 $\langle K:F\rangle_s=[K:F]$。
\end{enumerate}
\end{proposition}

\begin{proof}
先看乘法公式。给定一个 $F$-嵌入 $\tau:E\to L$，限制到 $K$ 得到一个 $F$-嵌入 $\tau|_K:K\to L$；在固定 $\tau|_K$ 后，$\tau$ 的剩余选择就是把 $E$ 作为 $K$-扩张继续嵌入。因此延拓可分两步数：先数 $K/F$ 的嵌入，再对每一个这样的嵌入数 $E/K$ 的延拓，得到 (1)。

对 (2)，若 $E/F$ 有限，取一条塔
\[
F=E_0\subset E_1\subset\cdots\subset E_r=E,\qquad E_i=E_{i-1}(\alpha_i).
\]
单扩张情形由态射扩张引理给出：嵌入数不超过极小多项式次数，也就是不超过 $[E_i:E_{i-1}]$。逐层相乘得
\[
\langle E:F\rangle_s\le \prod_i [E_i:E_{i-1}]=[E:F].
\]
对 (3)，由 (1) 与次数乘法 $[E:F]=[E:K][K:F]$ 比较即可。两个非负整数乘积达到最大值，当且仅当每一层都达到自己的最大值。
\end{proof}

\begin{proposition}
设 $E/F$ 为域扩张，$\alpha\in E$ 是 $F$ 上的代数元。则以下条件等价：
\begin{enumerate}[label=(\arabic*)]
\item $\alpha$ 在 $F$ 上的极小多项式没有重根。
\item $\langle F(\alpha):F\rangle_s=[F(\alpha):F]$。
\end{enumerate}
满足这些条件的元素称为 $F$ 上的可分元素。
\end{proposition}

\begin{proof}
设 $m_{\alpha,F}$ 的次数为 $n$。由单扩张理论，$[F(\alpha):F]=n$。由态射扩张引理，$\langle F(\alpha):F\rangle_s$ 等于 $m_{\alpha,F}$ 在代数闭包中互不相同根的个数。这个数等于 $n$，正好等价于 $m_{\alpha,F}$ 没有重根。
\end{proof}

\begin{theorem}
设 $E/F$ 为有限扩张。以下条件等价：
\begin{enumerate}[label=(\arabic*)]
\item $\langle E:F\rangle_s=[E:F]$。
\item $E$ 中任意元素都是 $F$ 上的可分元素。
\end{enumerate}
满足这些条件的有限扩张称为可分扩张。
\end{theorem}

\begin{proof}
若 (1) 成立，任取 $\alpha\in E$。由上一命题的塔公式应用于 $E\supset F(\alpha)\supset F$，有
\[
\langle E:F\rangle_s=\langle E:F(\alpha)\rangle_s\langle F(\alpha):F\rangle_s,
\]
同时
\[
[E:F]=[E:F(\alpha)][F(\alpha):F].
\]
左边达到右边最大值，迫使 $\langle F(\alpha):F\rangle_s=[F(\alpha):F]$，故 $\alpha$ 可分。

反过来，若 $E$ 中每个元素都可分，由本原元素定理可取 $\theta\in E$ 使 $E=F(\theta)$；或者更一般地取有限生成塔并逐步使用“每个生成元在前一层上仍可分”。每一步单扩张的嵌入数都等于次数，于是逐层相乘得 $\langle E:F\rangle_s=[E:F]$。
\end{proof}

对一般代数扩张 $L/F$，若它的每个有限生成子扩张都可分，则称 $L/F$ 可分。代数闭包 $F^{\mathrm{alg}}$ 中所有 $F$ 上可分元素组成的子域记作 $F^{\mathrm{sep}}$，称为 $F$ 的可分闭包。数域的特征为 $0$，所以本讲义讨论的许多数域扩张天然可分；但第十章以后进入特征 $p$ 的剩余域和完美化结构时，不可分现象会重新出现。

\paragraph{纯不可分扩张与可分次数}
若代数扩张 $P/F$ 满足：每个在 $F$ 上可分的元素都已经属于 $F$，则称 $P/F$ 为纯不可分扩张。设 $K/F$ 为代数扩张，令
\[
S=\{x\in K:x\text{ 在 }F\text{ 上可分}\}.
\]
则 $S$ 是 $K$ 的子域，$S/F$ 可分，且 $K/S$ 纯不可分。若 $[K:F]<\infty$，称 $[S:F]$ 为可分次数，记作 $[K:F]_s$；称 $[K:S]$ 为不可分次数，记作 $[K:F]_i$。

\begin{proposition}
设 $E/F$ 为有限扩张。则：
\begin{enumerate}[label=(\arabic*)]
\item $[E:F]_s=\langle E:F\rangle_s$。
\item 若 $E\supset K\supset F$，则
\[
[E:F]_s=[E:K]_s[K:F]_s,\qquad
[E:F]_i=[E:K]_i[K:F]_i.
\]
\end{enumerate}
\end{proposition}

\begin{proof}
令 $S$ 为 $E/F$ 的最大可分子域。因为 $E/S$ 纯不可分，纯不可分层不会增加嵌入数：一个嵌入在 $S$ 上确定后，纯不可分元素的像被方程 $X^{p^r}-a$ 唯一决定。因此 $\langle E:F\rangle_s=\langle S:F\rangle_s$。而 $S/F$ 可分且有限，所以 $\langle S:F\rangle_s=[S:F]$。这就是 (1)。

对 (2)，可分次数的乘法来自嵌入数的塔公式。不可分次数再由
\[
[E:F]=[E:F]_s[E:F]_i
\]
和普通次数乘法比较得到。
\end{proof}

\begin{proposition}
设 $\alpha$ 是特征 $p\ne0$ 的域 $F$ 上的代数元。
\begin{enumerate}[label=(\arabic*)]
\item 若 $\alpha$ 是 $F$ 上的不可分元素，则 $\alpha$ 的极小多项式形如 $X^{p^r}-a$，其中 $a\in F$。
\item 若 $\alpha$ 是 $X^{p^r}-a$ 的根，$a\in F$，并且 $\alpha\notin F$ 时相应极小多项式非一次，则 $\alpha$ 是 $F$ 上的不可分元素。
\end{enumerate}
\end{proposition}

\begin{proof}
在特征 $p$ 中，多项式有重根当且仅当它与导数有非平凡公因子。不可分不可约多项式的导数必须为零，否则不可约多项式与其导数互素，就不会有重根。导数为零意味着多项式只含 $X^p$ 的幂，即可写成 $g(X^p)$。若仍不可分，继续写成 $h(X^{p^2})$，直到得到 $X^{p^r}-a$ 型的纯不可分极小关系，也就是 $\alpha^{p^r}=a\in F$。

反过来，若 $\alpha^{p^r}=a\in F$，则 $\alpha$ 的共轭根在代数闭包中只有一个，因为
\[
X^{p^r}-a=(X-\alpha)^{p^r}
\]
在特征 $p$ 中成立。只要扩张非平凡，极小多项式有重根，因此 $\alpha$ 不可分。
\end{proof}

\begin{proposition}
设 $E/F$ 为有限可分扩张。则存在 $\theta\in E$ 使 $E=F(\theta)$。这样的 $\theta$ 称为扩张 $E/F$ 的本原元素。
\end{proposition}

\begin{proof}
只说明二生成情形，有限多个生成元可归纳处理。设 $E=F(\alpha,\beta)$。若 $F$ 是有限域，则 $E^\times$ 是有限循环群，取其生成元即可生成整个有限域扩张。若 $F$ 是无限域，考虑 $\theta_c=\alpha+c\beta$。对不同的 $F$-嵌入 $\sigma,\tau:E\to F^{\mathrm{alg}}$，坏的 $c$ 至多满足
\[
\sigma(\alpha)+c\sigma(\beta)=\tau(\alpha)+c\tau(\beta),
\]
这是关于 $c$ 的至多一个线性条件。嵌入只有有限多个，所以坏的 $c$ 只有有限多个。取避开它们的 $c\in F$，则不同嵌入在 $\theta_c$ 上取值不同，因此 $\theta_c$ 的共轭数达到 $[E:F]$，从而 $[F(\theta_c):F]=[E:F]$，故 $E=F(\theta_c)$。
\end{proof}

\begin{proposition}
设 $E/F$ 为有限可分扩张。则：
\begin{enumerate}[label=(\arabic*)]
\item 存在正规扩张 $K/F$ 使 $K\supset E$。
\item 可取这样的 $K$，使从 $E$ 到 $K$ 的 $F$-单态射共有 $[E:F]$ 个。
\end{enumerate}
\end{proposition}

\begin{proof}
由本原元素定理，取 $E=F(\theta)$。令 $m_{\theta,F}$ 为 $\theta$ 的极小多项式。取 $K$ 为 $m_{\theta,F}$ 的分裂域。则 $K/F$ 正规，且 $K$ 包含 $\theta$，所以包含 $E$。由于 $E/F$ 可分，$m_{\theta,F}$ 有 $[E:F]$ 个不同根；每个根给出一个 $F$-嵌入 $E=F(\theta)\to K$，反之每个嵌入都由 $\theta$ 的像决定，故嵌入数为 $[E:F]$。
\end{proof}

\paragraph{完全域的定义}
若域 $F$ 的所有代数扩张都是可分扩张，则称 $F$ 为完全域。

\begin{proposition}
\begin{enumerate}[label=(\arabic*)]
\item 特征为 $0$ 的域是完全域。
\item 有限域是完全域。
\end{enumerate}
\end{proposition}

\begin{proof}
若 $\operatorname{char}F=0$，任意不可约多项式 $f$ 的导数 $f'$ 不为零，且 $\deg f'<\deg f$。不可约性迫使 $\gcd(f,f')=1$，所以 $f$ 无重根。

若 $F$ 是有限域，设 $\operatorname{char}F=p$。Frobenius 映射 $x\mapsto x^p$ 是单射；有限集合上的单射必为满射。因此每个元素都有 $p$ 次根，$F=F^p$。在特征 $p$ 中，一个域完全等价于 $F=F^p$，所以有限域完全。
\end{proof}

\subsection{线性无交}

线性无交是“两个扩张没有额外线性关系”的精确说法。设 $E/F$ 是一个共同上方域，$E_1,E_2$ 是其中的 $F$-子代数。由乘法给出自然映射
\[
E_1\otimes_F E_2\longrightarrow E_1E_2,\qquad x\otimes y\longmapsto xy.
\]
若这个映射是同构，则称 $E_1,E_2$ 在 $F$ 上线性无交。这里 $E_1E_2$ 表示由 $E_1$ 与 $E_2$ 共同生成的子代数。

若 $E_1,E_2$ 是子域，线性无交可以用基来检验：$S\subset E_1$ 在 $F$ 上线性无关，当且仅当它在合成域 $E_1E_2$ 中仍然在 $E_2$ 上线性无关；等价地，$T\subset E_2$ 在 $F$ 上线性无关，当且仅当它在 $E_1E_2$ 中仍然在 $E_1$ 上线性无关。特别地，若 $[E_1:F]$ 与 $[E_2:F]$ 有限，线性无交通常体现为
\[
[E_1E_2:F]=[E_1:F][E_2:F].
\]

\begin{theorem}
设域 $F$ 的特征为 $p\ne0$，并在代数闭包中记
\[
F^{p^{-1}}=\{x\in F^{\mathrm{alg}}:x^p\in F\}.
\]
则代数扩张 $E/F$ 可分，当且仅当 $E$ 与 $F^{p^{-1}}$ 在 $F$ 上线性无交。
\end{theorem}

\begin{proof}
先说明为什么这一定理是在检测不可分性。域 $F^{p^{-1}}$ 是把 $F$ 中元素的 $p$ 次根一次性加入后得到的域；它只携带纯不可分方向的信息。若 $E/F$ 可分，则可分扩张与纯不可分扩张不会制造新的线性关系。具体地，任取 $E$ 中有限个在 $F$ 上线性无关的元素，它们生成一个有限可分子扩张 $E_0/F$。可分扩张的极小多项式在纯不可分基变换后仍保持不可约，因此这些元素在 $F^{p^{-1}}$ 上仍线性无关。于是 $E$ 与 $F^{p^{-1}}$ 线性无交。

反过来，若 $E/F$ 不可分，则存在 $\alpha\in E$ 在 $F$ 上不可分。按上面的不可分元素命题，某个 $p^r$ 次幂满足 $\alpha^{p^r}\in F$。取最小的 $r$，则 $\alpha^{p^{r-1}}\notin F$ 但其 $p$ 次幂在 $F$ 中，所以 $\alpha^{p^{r-1}}\in F^{p^{-1}}$。这个元素同时属于 $E$ 与 $F^{p^{-1}}$，却不属于 $F$。于是交中出现额外元素，张量映射不可能保持线性无交。故线性无交推出 $E/F$ 可分。
\end{proof}

这一定理的意义是：可分性可以被理解为“与纯不可分方向独立”。后面在处理完美域、剩余域和 $p$ 进周期环时，会反复检查一个基变换是否引入新的张量关系；判断方法仍是把原来在 $F$ 上线性无关的元素放到更大的基域上，逐项检查它们是否仍线性无关。

\section{Galois 扩张}
上一节已经说明：一个域扩张有两类基本问题。第一类是“根是否全部留在同一个域里”，这由正规性控制；第二类是“根是否彼此分开”，这由可分性控制。Galois 理论把这两类条件合在一起，并用自同构群来记录扩张的对称性。对初学者来说，本节最重要的转换是：
\[
\text{中间域}\quad \longleftrightarrow\quad \text{自同构群的子群}.
\]
这个转换不是形式装饰。后面研究素理想分解、Frobenius 元、类域论和 $L$ 函数时，许多数论问题都会先被翻译成 Galois 群中的子群、商群、共轭类和特征标问题。

\paragraph{固定域与固定群}
设 $K$ 是域。由 $K$ 的全部自同构组成的群记为 $\operatorname{Aut}(K)$，群运算是映射复合。若 $G$ 是 $\operatorname{Aut}(K)$ 的子群，定义
\[
I(G)=\{a\in K:\varphi(a)=a,\ \forall \varphi\in G\}.
\]
这个 $I(G)$ 称为 $G$ 的固定域。它确实是 $K$ 的子域：若 $a,b$ 被每个 $\varphi\in G$ 固定，则 $\varphi(a+b)=a+b$、$\varphi(ab)=ab$；若 $a\ne0$，则 $\varphi(a^{-1})=\varphi(a)^{-1}=a^{-1}$。因此加法、乘法和取逆都不离开 $I(G)$。

反过来，若 $E$ 是 $K$ 的子域，定义
\[
A(E)=\{\varphi\in\operatorname{Aut}(K):\varphi(a)=a,\ \forall a\in E\}.
\]
这个 $A(E)$ 称为固定 $E$ 的自同构群。它是 $\operatorname{Aut}(K)$ 的子群：恒等映射固定 $E$；两个固定 $E$ 的自同构复合后仍固定 $E$；若 $\varphi$ 固定 $E$，则对 $a\in E$ 有 $\varphi(a)=a$，再作用 $\varphi^{-1}$ 得 $\varphi^{-1}(a)=a$。

\begin{proposition}
设 $K$ 是域。映射 $I,A$ 有以下性质：
\begin{enumerate}[label=(\arabic*)]
\item 若 $G_1\supset G_2$ 是 $\operatorname{Aut}(K)$ 的子群，则 $I(G_1)\subset I(G_2)$。
\item 若 $F_1\supset F_2$ 是 $K$ 的子域，则 $A(F_1)\subset A(F_2)$。
\item 若 $F$ 是 $K$ 的子域，则 $I(A(F))\supset F$，并且 $A(I(A(F)))=A(F)$。
\item 若 $G$ 是 $\operatorname{Aut}(K)$ 的子群，则 $A(I(G))\supset G$，并且 $I(A(I(G)))=I(G)$。
\end{enumerate}
\end{proposition}

\begin{proof}
第 (1) 点说明“固定更多自同构的人更少”。若 $a\in I(G_1)$，则 $a$ 被 $G_1$ 中每个元素固定；因为 $G_2\subset G_1$，它当然也被 $G_2$ 中每个元素固定，所以 $a\in I(G_2)$。

第 (2) 点说明“要求固定更大的子域，自同构选择更少”。若 $\varphi\in A(F_1)$，则 $\varphi$ 固定 $F_1$ 的每个元素；由于 $F_2\subset F_1$，它也固定 $F_2$，所以 $\varphi\in A(F_2)$。

第 (3) 点先看包含关系。每个 $\varphi\in A(F)$ 都固定 $F$ 的每个元素，因此 $F$ 中每个元素都属于 $I(A(F))$，即 $I(A(F))\supset F$。由第 (2) 点，把 $F\subset I(A(F))$ 代入，得到
\[
A(I(A(F)))\subset A(F).
\]
反向包含也成立：若 $\varphi\in A(F)$，而 $x\in I(A(F))$，则 $x$ 按定义被 $A(F)$ 中所有自同构固定，特别被这个 $\varphi$ 固定。因此 $\varphi$ 固定 $I(A(F))$，即 $\varphi\in A(I(A(F)))$。两边合起来得到等号。

第 (4) 点完全对偶。若 $\varphi\in G$，则 $\varphi$ 固定所有被 $G$ 固定的元素，所以 $\varphi\in A(I(G))$，即 $A(I(G))\supset G$。由第 (1) 点，$G\subset A(I(G))$ 推出
\[
I(A(I(G)))\subset I(G).
\]
另一方面，$I(G)$ 的每个元素被 $A(I(G))$ 中所有自同构固定，所以 $I(G)\subset I(A(I(G)))$。因此等号成立。
\end{proof}

现在把上面的抽象记号放到一个扩张里。设 $E/F$ 是域扩张，记
\[
\operatorname{Aut}_F(E)=\{\varphi\in\operatorname{Aut}(E):\varphi(a)=a,\ \forall a\in F\}.
\]
它的元素称为固定 $F$ 的域自同构。记
\[
\mathcal F(E/F)=\{K:F\subset K\subset E\},\qquad
\mathcal G(E/F)=\{H:H\le \operatorname{Aut}_F(E)\}.
\]
于是有两个反向映射
\[
\mathcal F(E/F)\longrightarrow\mathcal G(E/F),\quad
K\longmapsto\operatorname{Aut}_K(E),
\]
和
\[
\mathcal G(E/F)\longrightarrow\mathcal F(E/F),\quad
H\longmapsto E^H,
\qquad
E^H=\{x\in E:\sigma(x)=x,\ \forall\sigma\in H\}.
\]
“反向”是因为中间域越大，要固定它的自同构越少；子群越大，被它固定的元素越少。

\begin{definition}
称可分正规代数扩张 $E/F$ 为 Galois 扩张。以
\[
\Gal(E/F)=\operatorname{Aut}_F(E)
\]
记其 Galois 群。
\end{definition}

代数数论的一个中心任务就是研究数域扩张的 Galois 群。直观地说，数域扩张把新的代数数加入原来的域；Galois 群记录这些新数的共轭替换方式。比如在 $\Q(\sqrt d)/\Q$ 中，非平凡自同构把 $\sqrt d$ 送到 $-\sqrt d$，这正是在交换方程 $X^2-d$ 的两个根。

\begin{theorem}
设 $E/F$ 是 Galois 扩张，记 $G=\Gal(E/F)$。则：
\begin{enumerate}[label=(\arabic*)]
\item $E^G=F$。
\item 若 $K\in\mathcal F(E/F)$，则 $E/K$ 是 Galois 扩张，并且 $E^{\Gal(E/K)}=K$。
\item 映射
\[
\mathcal F(E/F)\longrightarrow\mathcal G(E/F),\qquad
K\longmapsto\operatorname{Aut}_K(E)
\]
是单射。
\end{enumerate}
\end{theorem}

\begin{proof}
先证 (1)。包含 $F\subset E^G$ 来自定义：$\Gal(E/F)$ 的每个元素都固定 $F$。反过来取 $a\in E$ 且 $a\notin F$。因为 $E/F$ 可分，$a$ 在 $F$ 上的极小多项式没有重根；因为次数大于 $1$，它在代数闭包中至少有一个不同于 $a$ 的共轭根 $b$。由上一节的态射扩张引理，存在 $F$-同构 $F(a)\to F(b)$ 把 $a$ 送到 $b$。再把这个同构延拓到 $E$；正规性保证延拓后的像仍在 $E$ 中，所以得到某个 $\sigma\in\Gal(E/F)$ 满足 $\sigma(a)=b\ne a$。因此 $a$ 不属于 $E^G$。这说明 $E^G$ 中没有 $F$ 以外的元素，故 $E^G=F$。

再证 (2)。若 $F\subset K\subset E$，则 $E/K$ 仍是代数扩张。可分性向上保持：$E$ 中每个元素在 $F$ 上可分，其在 $K$ 上的极小多项式整除原来的多项式，仍无重根。正规性也向上保持：设 $f(X)\in K[X]$ 不可约且在 $E$ 中有根 $\alpha$。任取 $f$ 的另一个根 $\beta$，态射扩张把 $K(\alpha)\to K(\beta)$ 延拓到 $E$ 的 $K$-嵌入；由于 $E/F$ 正规，特别对这种固定 $K$ 的嵌入仍有像在 $E$ 中，因此 $\beta\in E$。于是 $f$ 在 $E$ 中完全分裂，$E/K$ 正规。故 $E/K$ 是 Galois 扩张。把 (1) 应用于扩张 $E/K$，得到 $E^{\Gal(E/K)}=K$。

最后证 (3)。若两个中间域 $K_1,K_2$ 有同一个像，即
\[
\operatorname{Aut}_{K_1}(E)=\operatorname{Aut}_{K_2}(E),
\]
对两边取固定域并用 (2)，得到
\[
K_1=E^{\Gal(E/K_1)}=E^{\Gal(E/K_2)}=K_2.
\]
所以该映射是单射。
\end{proof}

\begin{lemma}
设 $E/F$ 是可分代数扩张。若存在整数 $n$ 使得对每个 $\alpha\in E$ 都有 $[F(\alpha):F]\le n$，则 $[E:F]\le n$。
\end{lemma}

\begin{proof}
先看任意有限生成子扩张 $L=F(\alpha_1,\dots,\alpha_r)\subset E$。因为 $E/F$ 可分，$L/F$ 是有限可分扩张。本原元素定理给出某个 $\theta\in L$，使得 $L=F(\theta)$。由假设，
\[
[L:F]=[F(\theta):F]\le n.
\]
现在若 $[E:F]>n$，则可在 $E$ 中取到 $n+1$ 个 $F$-线性无关元素。它们生成某个有限子扩张 $L/F$，于是 $[L:F]\ge n+1$，这与刚才得到的 $[L:F]\le n$ 矛盾。因此 $[E:F]\le n$。
\end{proof}

\begin{proposition}
设 $K$ 是域，$G$ 是 $\operatorname{Aut}(K)$ 的有限子群，$n=|G|$。以 $F=K^G$ 记固定域。则：
\begin{enumerate}[label=(\arabic*)]
\item $K/F$ 是有限 Galois 扩张。
\item $\Gal(K/F)=G$。
\item $[K:F]=n$。
\end{enumerate}
\end{proposition}

\begin{proof}
先证明 $K/F$ 有限且次数至多 $n$。任取 $\alpha\in K$，考虑它在 $G$ 作用下的轨道
\[
G\alpha=\{\sigma(\alpha):\sigma\in G\}.
\]
多项式
\[
P_\alpha(X)=\prod_{\beta\in G\alpha}(X-\beta)
\]
的系数是这些轨道元素的对称多项式。对任意 $\tau\in G$，$\tau$ 只是重新排列集合 $G\alpha$，所以固定 $P_\alpha$ 的每个系数；这些系数属于 $F$。因此 $\alpha$ 是 $F$ 上次数不超过 $|G|=n$ 的多项式的根，从而 $[F(\alpha):F]\le n$。又因为 $P_\alpha$ 在 $K$ 中分裂且没有重复列入轨道元素，$\alpha$ 在 $F$ 上可分。由上一个引理，$[K:F]\le n$。

另一方面，$G$ 中的 $n$ 个自同构都是固定 $F$ 的 $F$-嵌入 $K\to K$。有限扩张的 $F$-嵌入数不超过扩张次数，所以
\[
n\le [K:F].
\]
于是 $[K:F]=n$。此时 $\operatorname{Aut}_F(K)$ 至少包含 $G$，而任意有限扩张的自同构数至多为次数，所以
\[
|G|=n=[K:F]\ge |\operatorname{Aut}_F(K)|\ge |G|.
\]
故 $\Gal(K/F)=\operatorname{Aut}_F(K)=G$。

最后说明 $K/F$ 是 Galois 扩张。我们已经知道 $K/F$ 有限，并且固定 $F$ 的自同构数达到最大值 $[K:F]$。这等价于每个 $F$-嵌入都落回 $K$ 且嵌入数没有因重根而减少，因此 $K/F$ 既正规又可分。
\end{proof}

\begin{proposition}
设 $E/F$ 是有限 Galois 扩张。则映射
\[
\mathcal F(E/F)\longrightarrow\mathcal G(E/F),\qquad
K\longmapsto\operatorname{Aut}_K(E)
\]
是满射。
\end{proposition}

\begin{proof}
任取子群 $H\le \Gal(E/F)$，令 $K=E^H$。若 $\sigma\in H$，则 $\sigma$ 按定义固定 $E^H=K$ 中的每个元素，所以 $\sigma\in \operatorname{Aut}_K(E)$，即 $H\subset \operatorname{Aut}_K(E)$。为了证明反向包含，应用上一个命题到域 $E$ 和有限自同构群 $H$：它给出
\[
[E:K]=|H|,\qquad \Gal(E/K)=H.
\]
而 $\operatorname{Aut}_K(E)=\Gal(E/K)$，所以 $K$ 的像正是 $H$。这证明了满射性。
\end{proof}

于是当 $E/F$ 是有限 Galois 扩张时，有双射
\[
\mathcal F(E/F)\longleftrightarrow \mathcal G(E/F),
\]
它由
\[
K\longmapsto \Gal(E/K),\qquad H\longmapsto E^H
\]
给出。这个双射称为 $E/F$ 的 Galois 对应。要特别记住它是反向的：域越大，对应子群越小。

若在这个对应下 $K\leftrightarrow H$，取 $\varphi\in G=\Gal(E/F)$，则
\[
\varphi K\leftrightarrow \varphi H\varphi^{-1}
\]
或等价地，按左、右作用约定也可写成 $\varphi K\leftrightarrow \varphi^{-1}H\varphi$。核心含义是不变的：共轭子群对应共轭中间域。因此，“中间域对所有 Galois 自同构稳定”会翻译成“对应子群是正规子群”。

\begin{theorem}
设 $E/F$ 是 Galois 扩张，取中间域 $E\supset K\supset F$。则：
\begin{enumerate}[label=(\arabic*)]
\item $K/F$ 是正规扩张，当且仅当 $\Gal(E/K)$ 是 $\Gal(E/F)$ 的正规子群。
\item 若 $K/F$ 是正规扩张，则限制同态
\[
\Gal(E/F)\longrightarrow \Gal(K/F),\qquad
\sigma\longmapsto \sigma|_K
\]
是满同态，核为 $\Gal(E/K)$。因此
\[
\Gal(K/F)\simeq \Gal(E/F)/\Gal(E/K).
\]
\end{enumerate}
\end{theorem}

\begin{proof}
设 $G=\Gal(E/F)$，$H=\Gal(E/K)$。先证明 (1)。若 $K/F$ 正规，则任取 $\sigma\in G$，因为 $\sigma$ 固定 $F$，所以 $\sigma(K)$ 是 $K$ 在 $E$ 内的一个 $F$-共轭像；正规性说明这个共轭像仍是 $K$。因此 $\sigma H\sigma^{-1}$ 固定 $\sigma(K)=K$，也就是 $\sigma H\sigma^{-1}=H$。所以 $H$ 在 $G$ 中正规。

反过来，若 $H$ 正规，则对任意 $\sigma\in G$，子群 $\sigma H\sigma^{-1}$ 的固定域是 $\sigma(K)$。正规性给出 $\sigma H\sigma^{-1}=H$，所以固定域相同，即 $\sigma(K)=K$。因此每个固定 $F$ 的自同构都把 $K$ 保持在自身。由于 $E/F$ 是正规可分扩张，$K/F$ 也可分；并且任意 $F$-嵌入 $K\to E$ 都来自某个 $E$ 的 $F$-自同构的限制，从而像仍是 $K$。故 $K/F$ 正规。

再证 (2)。若 $\sigma_1,\sigma_2\in\Gal(E/F)$，则 $(\sigma_1\sigma_2)|_K=\sigma_1|_K\circ\sigma_2|_K$，所以限制映射保持乘法，是群同态。它的核由那些在 $K$ 上恒等的 $\sigma\in G$ 组成，正是 $\Gal(E/K)$。剩下说明满射。取 $\tau\in\Gal(K/F)$。因为 $E/F$ 是正规代数扩张，态射扩张定理可把 $\tau:K\to K\subset E$ 延拓为某个 $F$-嵌入 $\sigma:E\to E$；正规性保证像仍是 $E$，所以 $\sigma\in\Gal(E/F)$，且 $\sigma|_K=\tau$。因此限制映射满。由群同态基本定理得到商群同构。
\end{proof}

\paragraph{合成域与生成子群}
若 $F\subset E_j\subset K$，把 $K$ 中包含所有 $E_j$ 的最小子域称为 $E_1,\dots,E_r$ 的合成域，记为 $E_1\cdots E_r$。它由所有 $E_j$ 的元素通过有限次加、减、乘、除生成。

同样地，若 $H_1,\dots,H_r$ 是群 $G$ 的子群，记 $H_1\cdots H_r$ 为 $G$ 中包含所有 $H_j$ 的最小子群，也就是由这些子群共同生成的子群。不要把这里的乘积误解为集合逐项相乘后自动闭合；真正需要的是“最小的子群”。

\begin{proposition}
设 $K/F$ 是有限 Galois 扩张，$G=\Gal(K/F)$。设 $E_j$ 是 $K$ 中包含 $F$ 的子域，$H_j$ 是 $G$ 的子群，并且 $E_j\leftrightarrow H_j$ 是 Galois 对应。则：
\[
\bigcap_{j=1}^r E_j\ \longleftrightarrow\ H_1\cdots H_r,
\qquad
E_1\cdots E_r\ \longleftrightarrow\ \bigcap_{j=1}^r H_j.
\]
\end{proposition}

\begin{proof}
先看交。一个元素 $x\in K$ 属于 $\bigcap_j E_j$，当且仅当它同时被每个 $H_j$ 固定；这又等价于它被由所有 $H_j$ 生成的子群 $H_1\cdots H_r$ 固定。因此
\[
\bigcap_j E_j=K^{H_1\cdots H_r}.
\]
这就是第一组对应。

再看合成域。自同构固定合成域 $E_1\cdots E_r$，当且仅当它固定每个 $E_j$；这等价于它同时属于每个 $H_j=\Gal(K/E_j)$。因此
\[
\Gal(K/E_1\cdots E_r)=\bigcap_j H_j.
\]
由 Galois 对应得到第二组对应。
\end{proof}

\begin{proposition}
设 $K/F$ 是有限 Galois 扩张，且 $G=\Gal(K/F)$ 是 Abel 群。按有限 Abel 群基本结构定理写
\[
G\simeq Z_1\times\cdots\times Z_s,
\]
其中 $Z_j$ 是 $p_j^{\,n_j}$ 阶循环群。令
\[
G_i=Z_1\times\cdots\times Z_{i-1}\times Z_{i+1}\times\cdots\times Z_s.
\]
按 Galois 对应取
\[
E_i=K^{Z_i},\qquad N_i=K^{G_i}.
\]
则：
\begin{enumerate}[label=(\arabic*)]
\item 在 Galois 对应下有
\[
N_1\cdots N_s\leftrightarrow G_1\cap\cdots\cap G_s=1,
\qquad
E_i\cap N_i\leftrightarrow Z_iG_i=G.
\]
\item $K=N_1\cdots N_s$。
\item $E_i=N_1\cdots N_{i-1}N_{i+1}\cdots N_s$。
\item $E_i\cap N_i=F$。
\end{enumerate}
\end{proposition}

\begin{proof}
因为 $G$ 是直积，$\bigcap_i G_i=1$，而 $Z_iG_i=G$。由上一命题，合成域 $N_1\cdots N_s$ 对应子群 $\bigcap_iG_i=1$。固定子群为 $1$ 的中间域是整个 $K$，所以 $K=N_1\cdots N_s$。

再看 $E_i\cap N_i$。交域对应生成子群，而 $E_i$ 对应 $Z_i$、$N_i$ 对应 $G_i$，所以 $E_i\cap N_i$ 对应由 $Z_i$ 与 $G_i$ 生成的子群，即整个 $G$。固定整个 $G$ 的域是 $F$，所以 $E_i\cap N_i=F$。

最后计算 $E_i$。合成域 $N_1\cdots N_{i-1}N_{i+1}\cdots N_s$ 对应子群
\[
\bigcap_{j\ne i}G_j.
\]
在直积分解中，一个元素属于所有 $G_j$（$j\ne i$），等价于它除第 $i$ 个坐标外的各坐标都为单位，也就是属于 $Z_i$。因此这个合成域对应 $Z_i$，即等于 $K^{Z_i}=E_i$。
\end{proof}

\begin{proposition}
设 $L$ 是域 $F$ 的代数闭包，$K\subset L$ 是 $F$ 的有限 Galois 扩张。在 $L$ 内取 $F$ 的任意扩张 $\widetilde F$，设
\[
\widetilde K=K\widetilde F,\qquad E=K\cap \widetilde F.
\]
则：
\begin{enumerate}[label=(\arabic*)]
\item $\widetilde K/\widetilde F$ 是 Galois 扩张。
\item 有自然同构
\[
\Gal(\widetilde K/\widetilde F)\simeq \Gal(K/E).
\]
\end{enumerate}
\end{proposition}

\begin{proof}
限制给出映射
\[
\rho:\Gal(\widetilde K/\widetilde F)\longrightarrow \Gal(K/E),\qquad
\sigma\longmapsto \sigma|_K.
\]
它是单射，因为 $\widetilde K=K\widetilde F$ 由 $K$ 与 $\widetilde F$ 共同生成，而 $\sigma$ 在 $\widetilde F$ 上恒等；若它在 $K$ 上也恒等，就在整个 $\widetilde K$ 上恒等。

说明满射。因为 $K/F$ 是有限 Galois 扩张，$K/E$ 也是有限 Galois 扩张。又 $E=K\cap\widetilde F$，有限 Galois 扩张与一个只在基域相交的代数扩张线性无交，所以 $K$ 与 $\widetilde F$ 在 $E$ 上线性无交。于是任意 $\tau\in\Gal(K/E)$ 与 $\widetilde F$ 上的恒等映射相容，可以唯一延拓为 $\widetilde K=K\widetilde F$ 的 $\widetilde F$-自同构。这个延拓限制到 $K$ 正是 $\tau$，故 $\rho$ 满射。

由于 $\Gal(K/E)$ 的每个元素都延拓到 $\widetilde K/\widetilde F$，$\widetilde K$ 中来自 $K$ 的共轭根仍全部留在 $\widetilde K$；可分性也由 $K/F$ 的可分性保持。因此 $\widetilde K/\widetilde F$ 是有限 Galois 扩张。
\end{proof}

\paragraph{无关 Galois 扩张}
设有域 $F\subset E_i\subset K$。若对每个 $i$ 都有
\[
E_i\cap E_1\cdots E_{i-1}E_{i+1}\cdots E_s=F,
\]
则称扩张 $E_i/F$ 彼此无关。这个条件的含义是：第 $i$ 个扩张没有被其余扩张的合成域提前包含任何非平凡部分。它是“线性无交”的 Galois 对应版本。

\begin{proposition}
设域扩张 $K/F$ 内有无关 Galois 扩张 $E_i/F$。取合成域
\[
E=E_1\cdots E_s.
\]
则 $E/F$ 为 Galois 扩张，并且自然限制映射给出同构
\[
\Gal(E/F)\simeq \Gal(E_1/F)\times\cdots\times\Gal(E_s/F).
\]
等价地，把 $\Gal(E_i/F)$ 识别为固定其余因子的子群后，$\Gal(E/F)$ 分解为这些互相独立的 Galois 对称性的直积。
\end{proposition}

\begin{proof}
先说明 $E/F$ 是 Galois 扩张。每个 $E_i/F$ 都正规可分。可分扩张的合成仍可分；正规扩张的合成仍正规，因为每个 $E_i$ 中出现的共轭根已经全部在 $E_i$ 内，合在一起后仍在 $E$ 内。因此 $E/F$ 是 Galois 扩张。

考虑限制映射
\[
\rho:\Gal(E/F)\longrightarrow \prod_{i=1}^s\Gal(E_i/F),\qquad
\sigma\longmapsto(\sigma|_{E_i})_i.
\]
它是单射，因为 $E$ 由所有 $E_i$ 生成；若 $\sigma$ 在每个 $E_i$ 上恒等，则在合成域 $E$ 上恒等。

为了证明满射，给定一组 $\tau_i\in\Gal(E_i/F)$。无关条件保证 $E_i$ 与其余因子的合成域在 $F$ 上没有非平凡交，因此这些自同构在公共部分上的取值相容。于是可以先在张量积意义下定义
\[
x_1\cdots x_s\longmapsto \tau_1(x_1)\cdots\tau_s(x_s),
\qquad x_i\in E_i,
\]
再借由线性无交把它下降为合成域 $E$ 上的自同构。这个自同构固定 $F$，且限制到每个 $E_i$ 正是 $\tau_i$。所以 $\rho$ 满射。
\end{proof}

最后记录一个以后常用的构造。设 $K/F$ 是有限扩张。则存在有限 Galois 扩张 $L/F$ 使得 $L\supset K$。构造如下：取 $K$ 的代数闭包 $K^{\mathrm{alg}}$，设
\[
\sigma_1,\dots,\sigma_r:K\longrightarrow K^{\mathrm{alg}}
\]
是所有不同的固定 $F$ 的域单态射。令
\[
L=\sigma_1(K)\cdots\sigma_r(K).
\]
这个 $L$ 包含 $K$ 的所有 $F$-共轭像，所以它是正规扩张；若 $K/F$ 可分，则这些嵌入已经给出全部可分共轭，$L/F$ 是有限 Galois 扩张。数域的特征为 $0$，有限扩张自动可分，所以在代数数论中我们经常把一个有限数域扩张放进这样的有限 Galois 扩张里，再用 Galois 群统一处理所有共轭信息。

\section{迹和范}
迹和范把扩张域中的元素压回基域。它们不是分析里的长度或大小，而是线性代数中的“迹”和“行列式”。设 $E/F$ 是 $n$ 维有限域扩张，取 $x\in E$。乘以 $x$ 给出一个 $F$-线性映射
\[
\lambda_x:E\longrightarrow E,\qquad y\longmapsto xy.
\]
选定 $F$-向量空间 $E$ 的一组基，就得到矩阵 $[\lambda_x]$。改变基只会把矩阵变成相似矩阵，所以特征多项式
\[
f_x(X)=\det(XI-[\lambda_x])
\]
与基的选择无关。写成
\[
f_x(X)=X^n-tX^{n-1}+\cdots+(-1)^n\nu,
\]
其中 $t,\nu\in F$。称 $f_x(X)$ 为 $x$ 在扩张 $E/F$ 中的特征多项式；称 $t$ 为 $x$ 的迹，记为
\[
\Tr_{E/F}(x)=t;
\]
称 $\nu$ 为 $x$ 的范，记为
\[
\Nm_{E/F}(x)=\nu.
\]
等价地，$\Tr_{E/F}(x)$ 是线性算子 $\lambda_x$ 的矩阵迹，$\Nm_{E/F}(x)$ 是 $\lambda_x$ 的行列式。这里的范属于 $F$，通常不是非负实数。

\begin{proposition}
设 $E/F$ 是 $n$ 维域扩张。
\begin{enumerate}[label=(\arabic*)]
\item $\Tr_{E/F}:E\to F$ 是 $F$-线性函数。也就是说，对 $x,y\in E$ 与 $a\in F$ 有
\[
\Tr_{E/F}(x+y)=\Tr_{E/F}(x)+\Tr_{E/F}(y),\qquad
\Tr_{E/F}(ax)=a\Tr_{E/F}(x).
\]
\item 以 $E^\times$ 记 $E$ 的非零元素乘法群。映射 $\Nm_{E/F}:E^\times\to F^\times$ 是群同态，即
\[
\Nm_{E/F}(xy)=\Nm_{E/F}(x)\Nm_{E/F}(y).
\]
\item 若 $1_E$ 是 $E$ 的单位元，则 $\Tr_{E/F}(1_E)=n1_F$。并且对 $x\in E$、$a\in F$ 有
\[
\Nm_{E/F}(ax)=a^n\Nm_{E/F}(x).
\]
\end{enumerate}
\end{proposition}

\begin{proof}
选定 $E$ 的一组 $F$-基。因为 $\lambda_{x+y}=\lambda_x+\lambda_y$，矩阵也相加；矩阵迹对加法是线性的，所以
\[
\Tr_{E/F}(x+y)=\operatorname{Tr}(\lambda_x+\lambda_y)
=\operatorname{Tr}(\lambda_x)+\operatorname{Tr}(\lambda_y).
\]
又因为 $\lambda_{ax}=a\lambda_x$，矩阵迹满足 $\operatorname{Tr}(aM)=a\operatorname{Tr}(M)$，得到 $\Tr_{E/F}(ax)=a\Tr_{E/F}(x)$。

对范，若 $x,y\ne0$，则 $\lambda_{xy}=\lambda_x\circ\lambda_y$。行列式对矩阵乘法相乘，所以
\[
\Nm_{E/F}(xy)=\det(\lambda_x\lambda_y)=\det(\lambda_x)\det(\lambda_y)
=\Nm_{E/F}(x)\Nm_{E/F}(y).
\]
由于 $x\ne0$ 时 $\lambda_x$ 是可逆线性映射，行列式非零，故范确实落在 $F^\times$。

最后，$\lambda_{1_E}$ 是 $n$ 维向量空间上的恒等映射，所以迹为 $n1_F$。同时 $\lambda_{ax}=a\lambda_x$，$n$ 维矩阵乘以标量 $a$ 时行列式乘上 $a^n$，故 $\Nm_{E/F}(ax)=a^n\Nm_{E/F}(x)$。
\end{proof}

\begin{proposition}
设有有限域扩张 $F\subset E\subset K$。对 $x\in K$ 有
\begin{enumerate}[label=(\arabic*)]
\item $\Tr_{K/F}(x)=\Tr_{E/F}(\Tr_{K/E}(x))$。
\item $\Nm_{K/F}(x)=\Nm_{E/F}(\Nm_{K/E}(x))$。
\end{enumerate}
\end{proposition}

\begin{proof}
把 $K$ 先看成 $E$-向量空间，再把 $E$ 看成 $F$-向量空间。取 $E/F$ 的一组基 $e_1,\dots,e_r$，取 $K/E$ 的一组基 $u_1,\dots,u_s$。于是 $e_i u_j$ 构成 $K/F$ 的一组基。

乘法映射 $\lambda_x:K\to K$ 先是 $E$-线性的。设它在 $K/E$ 基 $u_1,\dots,u_s$ 下的矩阵为 $B=(b_{ij})$，其中 $b_{ij}\in E$。把这个矩阵再按 $E/F$ 展开，就得到一个 $rs\times rs$ 的 $F$-矩阵。这个展开矩阵的迹等于
\[
\Tr_{E/F}(b_{11}+\cdots+b_{ss})
=\Tr_{E/F}(\Tr_{K/E}(x)),
\]
而左边正是 $\Tr_{K/F}(x)$。这给出迹的塔式公式。

行列式的计算使用同一个基展开原则，但对象从迹换成行列式。先在 $K/E$ 的基下计算 $\lambda_x$ 的 $E$-行列式 $\det_E B$，再把这个 $E$ 中的元素通过 $E/F$ 的范压到 $F$ 中。按块矩阵展开，$F$-行列式正是
\[
\Nm_{E/F}(\det_E B)=\Nm_{E/F}(\Nm_{K/E}(x)).
\]
这就是 $\Nm_{K/F}(x)=\Nm_{E/F}(\Nm_{K/E}(x))$。
\end{proof}

\paragraph{特征多项式与极小多项式}
设 $x\in E$，令 $h_x(X)$ 为 $x$ 在 $F$ 上的极小多项式：
\[
h_x(X)=X^m+a_{m-1}X^{m-1}+\cdots+a_0.
\]
则 $1,x,\dots,x^{m-1}$ 是 $F(x)/F$ 的一组基。若 $y_1,\dots,y_r$ 是 $E/F(x)$ 的一组基，那么
\[
y_1,xy_1,\dots,x^{m-1}y_1;\ \dots;\ y_r,xy_r,\dots,x^{m-1}y_r
\]
是 $E/F$ 的一组基。相对于这组基，$\lambda_x$ 的矩阵是由 $r=[E:F(x)]$ 个相同友矩阵组成的对角块矩阵。友矩阵的特征多项式正是 $h_x(X)$，所以
\[
f_x(X)=h_x(X)^{[E:F(x)]}.
\]
这条公式解释了一个容易混淆的点：$x$ 的极小多项式只看单生成子域 $F(x)$，而 $x$ 在 $E/F$ 中的特征多项式还要把 $E$ 相对 $F(x)$ 的维数重复进去。

现在设 $E/F$ 是可分扩张。由本原元素定理，取 $\theta\in E$ 使 $E=F(\theta)$。设 $K$ 是 $\theta$ 在 $F$ 上的分裂域，于是 $K\supset E$。令 $G=\Gal(K/F)$，令 $H=\Gal(K/E)$。取代表元
\[
\sigma_1H,\dots,\sigma_nH,\qquad n=[E:F],
\]
使它们遍历 $G/H$。于是限制映射
\[
\sigma_j|_E:E\longrightarrow K,\qquad 1\le j\le n,
\]
正好给出 $E$ 到正规闭包 $K$ 中的全部 $F$-态射。以下仍把 $\sigma_j|_E$ 简记为 $\sigma_j$。

\begin{proposition}
设 $E/F$ 是可分扩张，$x\in E$。则
\begin{enumerate}[label=(\arabic*)]
\item $\Tr_{E/F}(x)=\sigma_1(x)+\cdots+\sigma_n(x)$。
\item $\Nm_{E/F}(x)=\sigma_1(x)\sigma_2(x)\cdots\sigma_n(x)$。
\end{enumerate}
\end{proposition}

\begin{proof}
先处理 $E=F(x)$ 的情形。此时 $x$ 的极小多项式在 $K$ 中分解为
\[
h_x(X)=\prod_{j=1}^n (X-\sigma_j(x)).
\]
另一方面，$f_x(X)=h_x(X)$。比较 $X^{n-1}$ 的系数得到迹等于共轭和，比较常数项得到范等于共轭积。

一般情形中，令 $m=[F(x):F]$，$r=[E:F(x)]$。按上面的特征多项式公式，
\[
f_x(X)=h_x(X)^r.
\]
$x$ 在 $F$ 上的每个共轭在 $E$ 的所有嵌入中恰好重复 $r$ 次，因为一个 $F(x)$ 的嵌入可延拓为 $E$ 的 $r$ 个 $F$-嵌入。于是 $f_x$ 的全部根按重数列出正是
\[
\sigma_1(x),\dots,\sigma_n(x).
\]
特征多项式的次高项系数给出根和，常数项给出根积，因此得到两条公式。
\end{proof}

\paragraph{迹配对}
设 $V$ 是域 $F$ 上的向量空间。映射 $b:V\times V\to F$ 称为双线性型，若对任意 $x,y,z\in V$ 和 $a\in F$ 有
\[
b(ax+y,z)=ab(x,z)+b(y,z),\qquad
b(z,ax+y)=ab(z,x)+b(z,y).
\]
若左右两个零化子都为零：
\[
\{x\in V:b(x,y)=0,\ \forall y\in V\}=\{0\}
\]
和
\[
\{x\in V:b(y,x)=0,\ \forall y\in V\}=\{0\},
\]
则称 $b$ 为非退化双线性型。非退化的意思是：没有非零向量会和所有向量配对为零。

\begin{proposition}
设 $E/F$ 是有限域扩张。对 $x,y\in E$ 定义
\[
(x,y)=\Tr_{E/F}(xy).
\]
则：
\begin{enumerate}[label=(\arabic*)]
\item $(\cdot,\cdot):E\times E\to F$ 是对称双线性型。
\item 这个双线性型非退化，当且仅当 $E/F$ 是可分扩张。
\item 若 $E/F$ 可分，且 $u_1,\dots,u_n$ 是 $E/F$ 的一组基，则存在一组基 $v_1,\dots,v_n$ 使得
\[
\Tr_{E/F}(u_i v_j)=\delta_{ij}.
\]
\item 若 $E/F$ 可分，$\sigma_1,\dots,\sigma_n$ 是 $E$ 到正规闭包中的全部 $F$-态射，记 $A=(\sigma_j(u_i))_{i,j}$，则
\[
\det(A)^2=\det\bigl(\Tr_{E/F}(u_i u_j)\bigr).
\]
\end{enumerate}
\end{proposition}

\begin{proof}
第 (1) 点来自迹的线性性与域乘法交换性。对第一个变量，
\[
(ax+y,z)=\Tr_{E/F}((ax+y)z)
=a\Tr_{E/F}(xz)+\Tr_{E/F}(yz).
\]
第二个变量同样成立。又 $xy=yx$，所以 $(x,y)=(y,x)$。

证明 (2) 的可分方向。设 $E/F$ 可分，并取一组基 $u_1,\dots,u_n$。令 $A=(\sigma_j(u_i))$。若 $\det A=0$，则 $A$ 的行向量线性相关，存在不全为零的 $c_i\in F$ 使得
\[
\sum_i c_i\sigma_j(u_i)=0,\qquad j=1,\dots,n.
\]
这等价于元素 $u=\sum_i c_i u_i$ 在每个 $F$-嵌入 $\sigma_j$ 下都变成 $0$。嵌入是单射，所以 $u=0$，与 $u_i$ 是基矛盾。因此 $\det A\ne0$。由第 (4) 中即将验证的矩阵恒等式，迹配对的 Gram 矩阵行列式非零，所以配对非退化。

若 $E/F$ 不可分，则存在非零元素的可分共轭数量少于扩张次数。等价地，嵌入数小于 $[E:F]$，迹函数无法给出足够多的线性函数来分离 $E$ 的所有方向；具体地，迹配对的 Gram 行列式为零，因而存在非零 $x$ 使 $\Tr_{E/F}(xy)=0$ 对所有 $y\in E$ 成立。故配对退化。

第 (3) 点是线性代数。非退化配对给出同构
\[
E\longrightarrow \Hom_F(E,F),\qquad
v\longmapsto (u\mapsto \Tr_{E/F}(uv)).
\]
令 $u_1^\vee,\dots,u_n^\vee$ 为基 $u_1,\dots,u_n$ 的对偶基。由上述同构，每个 $u_j^\vee$ 唯一对应某个 $v_j\in E$，于是
\[
\Tr_{E/F}(u_i v_j)=u_j^\vee(u_i)=\delta_{ij}.
\]
这些 $v_j$ 线性无关且个数为 $n$，所以也是一组基。

最后证 (4)。由可分情形的迹公式，
\[
\Tr_{E/F}(u_i u_j)=\sum_{k=1}^n \sigma_k(u_i)\sigma_k(u_j).
\]
这正是矩阵乘积 $AA^{\mathsf T}$ 的第 $(i,j)$ 项。因此迹配对的 Gram 矩阵为 $AA^{\mathsf T}$。取行列式得到
\[
\det\bigl(\Tr_{E/F}(u_i u_j)\bigr)=\det(A)\det(A^{\mathsf T})=\det(A)^2.
\]
\end{proof}

\section{有限域}
有限域是代数数论的“剩余层”。数域整数环中的素理想 $\mathfrak p$ 模掉以后，得到的剩余域 $\mathcal O_F/\mathfrak p$ 是有限域；局部域的剩余域也通常是有限域。后面定义 Frobenius 自同构、Artin 符号、局部因子时，都会反复用到本节的基本事实。

先区分两个概念。一个环 $D$ 若满足 $D\setminus\{0\}$ 在乘法下是群，则称 $D$ 为可除环。可除环不预设乘法交换；域是乘法交换的可除环。本节第一件事是：有限情形中，不交换不会发生。

\begin{proposition}
有限可除环是域。
\end{proposition}

\begin{proof}
这是 Wedderburn 小定理。给出证明思路和关键环节，因为它告诉我们为什么“有限除法结构”自动交换。

设 $D$ 是有限可除环，$Z$ 是其中心。则 $Z$ 是有限域，设 $|Z|=q$。若 $D$ 不交换，则存在严格大于 $Z$ 的最大交换子域 $L\subset D$。所有这样的最大交换子域都含 $Z$，并且它们在 $D^\times$ 的共轭作用下互相移动。对任意 $a\in D^\times$，共轭 $x\mapsto axa^{-1}$ 保持加法和乘法，因此把交换子域送到交换子域。

证明的核心是类方程计数。若两个最大交换子域的交超过 $Z$，取交中的非中心元素 $u$。所有与 $u$ 交换的元素构成 $D$ 的一个真可除子环；它比 $D$ 小，按对 $|D|$ 的归纳假设是域，于是包含这两个最大交换子域生成的交换结构，和最大性矛盾。因此不同最大交换子域只能在 $Z$ 中相交。

现在取一个最大交换子域 $L$，设 $|L|=q^m$，$m>1$。因为有限域的乘法群是循环群，$L^\times$ 是循环群。共轭作用给出从 $D^\times$ 中 $L^\times$ 的正规化子到 $\operatorname{Aut}_Z(L)$ 的同态，而 $\operatorname{Aut}_Z(L)$ 由 Frobenius 生成，阶为 $m$。类方程把 $D^\times$ 分解为中心部分和非中心元素的共轭类；每个非中心元素的中心化子是某个最大交换子域的乘法群。于是非中心元素总数被 $q^m-q$ 的倍数控制，同时又必须与 $|D^\times|=q^N-1$ 的整除关系相容。把这些整除关系写出后得到 $q^m-1$ 同时整除一个严格较小的数，矛盾。

矛盾说明 $D$ 不可能非交换。因此 $D=Z$，乘法交换，$D$ 是域。
\end{proof}

上面证明中用到的有限域乘法群循环性，会在下面同时得到。以 $F^\times$ 记域 $F$ 的乘法群 $F\setminus\{0\}$。

\begin{proposition}
设 $K$ 是特征 $p>0$ 的代数闭域。对任意 $f\ge1$，$K$ 包含唯一一个有 $p^f$ 个元素的有限域 $F$。具体地，
\[
F=\{x\in K:x^{p^f}=x\}.
\]
并且 $F^\times$ 是方程 $X^{p^f-1}=1$ 在 $K$ 中的根集，$F^\times$ 是 $p^f-1$ 个元素的循环群。
\end{proposition}

\begin{proof}
令 $q=p^f$，考虑多项式 $X^q-X$。它的导数是 $qX^{q-1}-1=-1$，因为在特征 $p$ 中 $q=p^f=0$。所以它没有重根，在代数闭域 $K$ 中有 $q$ 个不同根。记这些根的集合为 $F$。

证明 $F$ 是域。若 $a,b\in F$，则 $a^q=a$、$b^q=b$。Frobenius 的 $f$ 次迭代满足
\[
(a+b)^q=a^q+b^q=a+b,\qquad
(ab)^q=a^qb^q=ab.
\]
若 $a\ne0$，由 $a^q=a$ 得 $a^{q-1}=1$，于是 $(a^{-1})^q=a^{-q}=a^{-1}$。所以 $F$ 对加法、乘法和非零元取逆封闭，确实是含 $q$ 个元素的域。

唯一性也来自同一个方程。若 $F'\subset K$ 是任意含 $q$ 个元素的有限域，则 $F'^\times$ 的阶是 $q-1$，所以每个非零 $x\in F'$ 满足 $x^{q-1}=1$，零元满足 $0^q=0$。因此 $F'$ 的每个元素都是 $X^q-X$ 的根。两边都有 $q$ 个元素，故 $F'=F$。

最后证明 $F^\times$ 循环。先说明一个一般事实：域的有限乘法子群都是循环群。设 $G$ 是域的有限乘法子群，令 $m$ 为 $G$ 的指数，即所有元素阶的最小公倍数。则 $G$ 中每个元素都是 $X^m-1$ 的根，所以 $|G|\le m$；但按定义 $m$ 又整除有限群的阶 $|G|$。因此 $m=|G|$。有限 Abel 群若指数等于阶，必为循环群。应用到 $G=F^\times$，得 $F^\times$ 是阶 $q-1$ 的循环群。它的元素正是 $X^{q-1}=1$ 的全部根。
\end{proof}

\begin{corollary}
除同构外，$q=p^f$ 个元素的有限域唯一。
\end{corollary}

\begin{proof}
任意 $q$ 元有限域都含有素域 $\F_p$，并可嵌入其代数闭包。按上一命题，它在代数闭包中必是 $X^q-X$ 的根集。因此两个 $q$ 元有限域在代数闭包中同构到同一个根集，故彼此同构。
\end{proof}

以后把 $q$ 个元素的有限域记为 $\F_q$。

\begin{corollary}
设 $f\ge1$、$f'\ge1$，令 $q=p^f$、$q'=p^{f'}$。则：
\begin{enumerate}[label=(\arabic*)]
\item $\F_{q'}\supset \F_q$ 当且仅当 $f\mid f'$。
\item 若 $\F_{q'}\supset\F_q$，则 $\F_{q'}/\F_q$ 是次数 $f'/f$ 的循环扩张，并且 $\Gal(\F_{q'}/\F_q)$ 由
\[
x\longmapsto x^q
\]
生成。
\item 设 $f'=fd$。对 $n\ge1$，集合
\[
\{x\in\F_{q'}:x^{q^n}=x\}
\]
是有 $q^r$ 个元素的子域，其中 $r=(d,n)$。
\end{enumerate}
\end{corollary}

\begin{proof}
先证 (1)。若 $\F_{q'}\supset \F_q$，则 $\F_{q'}$ 是 $\F_q$ 上的有限维向量空间，所以 $|\F_{q'}|=|\F_q|^d=q^d$，即 $p^{f'}=p^{fd}$，故 $f'=fd$。反过来若 $f\mid f'$，写 $f'=fd$，则
\[
q'-1=p^{fd}-1
\]
被 $q-1=p^f-1$ 整除，所以 $X^q-X$ 的根包含在 $X^{q'}-X$ 的根中；因此 $\F_q\subset\F_{q'}$。

第 (2) 点中，扩张次数由元素个数给出：
\[
[\F_{q'}:\F_q]=\log_q(q')=f'/f.
\]
映射 $\varphi:x\mapsto x^q$ 固定 $\F_q$，因为 $\F_q$ 中元素满足 $x^q=x$。它在 $\F_{q'}$ 上是自同构；其 $d=f'/f$ 次迭代是 $x\mapsto x^{q^d}=x^{q'}=x$。若 $0<n<d$，固定 $\varphi^n$ 的元素满足 $x^{q^n}=x$，至多 $q^n<q^d=q'$ 个，所以 $\varphi^n$ 不是恒等。故 $\varphi$ 的阶为 $d$。Galois 群的阶等于扩张次数 $d$，所以它由 $\varphi$ 生成，是循环群。

第 (3) 点实际上是在问 $\varphi^n$ 的固定域大小。循环群 $\langle\varphi\rangle$ 的阶是 $d$，元素 $\varphi^n$ 的固定域对应于它生成的子群。该子群阶为 $d/(d,n)$，所以固定域在 $\F_q$ 上的次数是 $(d,n)=r$。因此固定域有 $q^r$ 个元素。
\end{proof}

Frobenius 是有限域理论的核心。后面讨论素理想在 Galois 扩张中的 Frobenius 元时，本质上就是把剩余域上的 $x\mapsto x^q$ 提升回数域扩张的 Galois 群。

\section{过滤}
过滤把一个线性对象分成由浅到深的层次。代数数论里最先遇到的模型是局部域整数环的单位群过滤
\[
\mathcal O_K^\times\supset 1+\mathfrak p\supset 1+\mathfrak p^2\supset\cdots,
\]
以后在分歧群、完备化和 $p$ 进 Hodge 理论中还会反复出现。这里先只讨论有限维向量空间，因为这一层语言已经足够支撑后面用到的“逐层看对象”的方法。

\begin{definition}[递减过滤、穷举性、分离性]
设 $V$ 是域 $F$ 上的向量空间。$V$ 的一个递减过滤是指一族子空间 $\{\Fil^i(V)\}_{i\in\Z}$，满足对每个 $i\in\Z$ 都有 $\Fil^{i+1}(V)\subset \Fil^i(V)$。若 $\bigcup_i\Fil^i(V)=V$，称过滤是穷举的；若 $\bigcap_i\Fil^i(V)=0$，称过滤是分离的。关联分级空间定义为
\[
\operatorname{gr}(V)=\bigoplus_{i\in\Z}\operatorname{gr}^i(V),
\qquad
\operatorname{gr}^i(V)=\Fil^i(V)/\Fil^{i+1}(V).
\]
\end{definition}

这一定义的重点不是把 $V$ 写成直和。一般而言，$V$ 不等于各层商空间的直和；关联分级只记录“从第 $i+1$ 层下降到第 $i$ 层时新出现的部分”。初学时可以把 $\Fil^i(V)$ 想成“精度至少为 $i$ 的元素”，把 $\operatorname{gr}^i(V)$ 想成“精度刚好在第 $i$ 层贡献的信息”。

\begin{example}[一维过滤]
令 $V=Fe$。给定整数 $a$，定义
\[
\Fil^i(V)=
\begin{cases}
V,&i\le a,\\
0,&i>a.
\end{cases}
\]
这是递减、穷举且分离的过滤。它的关联分级只有一层非零：$\operatorname{gr}^a(V)=V$，而 $i\ne a$ 时 $\operatorname{gr}^i(V)=0$。这个整数 $a$ 就是该一维过滤空间的跃点。
\end{example}

\begin{proposition}[子空间与商空间的诱导过滤]
设 $V$ 带递减过滤，$W\subset V$ 是子空间。定义
\[
\Fil^i(W)=W\cap\Fil^i(V),\qquad
\Fil^i(V/W)=(\Fil^i(V)+W)/W.
\]
则这两族子空间分别给出 $W$ 和 $V/W$ 的递减过滤。如果 $V$ 的过滤穷举，则二者穷举；如果 $V$ 的过滤分离，则 $W$ 的过滤分离。
\end{proposition}

\begin{proof}
先看 $W$。因为 $\Fil^{i+1}(V)\subset\Fil^i(V)$，所以 $W\cap\Fil^{i+1}(V)\subset W\cap\Fil^i(V)$，这给出递减性。若 $w\in W$，且 $V$ 的过滤穷举，则存在 $i$ 使 $w\in\Fil^i(V)$，于是 $w\in\Fil^i(W)$。若 $w$ 落在所有 $\Fil^i(W)$ 中，则它落在所有 $\Fil^i(V)$ 中；当 $V$ 的过滤分离时，$w=0$。

再看商空间。若 $\bar v\in(\Fil^{i+1}(V)+W)/W$，取代表元 $v\in\Fil^{i+1}(V)+W$。由于 $\Fil^{i+1}(V)\subset\Fil^i(V)$，同一个代表元也属于 $\Fil^i(V)+W$，于是商过滤递减。若 $V$ 的过滤穷举，任取 $\bar v\in V/W$，选代表元 $v\in V$，存在 $i$ 使 $v\in\Fil^i(V)$，从而 $\bar v\in\Fil^i(V/W)$。
\end{proof}

商过滤的分离性需要额外条件，例如 $W$ 在由过滤决定的拓扑中闭。有限维且过滤只有有限多个跃点时这一点自动成立；本书此处只需要上述自然构造本身。

\begin{definition}[过滤同态、核、余核]
设 $V'$ 与 $V$ 是带递减过滤的有限维 $F$-向量空间。一个过滤同态是线性映射 $T:V'\to V$，满足对每个 $i\in\Z$ 都有 $T(\Fil^i(V'))\subset \Fil^i(V)$。它的核和余核带自然过滤：
\[
\Fil^i(\Ker T)=\Ker T\cap\Fil^i(V'),
\qquad
\Fil^i(\operatorname{Cok}T)=\bigl(\Fil^i(V)+T(V')\bigr)/T(V').
\]
\end{definition}

\begin{proof}[公式检验]
核公式是子空间过滤的特例。余核公式是商空间过滤的特例，因为
\[
\operatorname{Cok}T=V/T(V').
\]
需要注意的是，$T(V')$ 未必继承为每层的满像；余核过滤只要求先在目标空间取 $\Fil^i(V)$，再模去整个像 $T(V')$。
\end{proof}

\begin{example}[底层同构不一定是过滤同构]
取同一个一维向量空间 $Fe$。在源空间 $V'$ 上设
\[
\Fil^0(V')=V',\qquad \Fil^1(V')=0,
\]
在目标空间 $V$ 上设
\[
\Fil^0(V)=V,\qquad \Fil^1(V)=V,\qquad \Fil^2(V)=0.
\]
恒等映射 $T:V'\to V$ 满足 $T(\Fil^0(V'))\subset\Fil^0(V)$ 且 $T(\Fil^1(V'))\subset\Fil^1(V)$，所以它是过滤同态；但它的逆映射不保持过滤，因为目标中的 $\Fil^1(V)=V$ 会被送到源空间的 $V'$，而源空间的 $\Fil^1(V')$ 是 $0$。因此底层线性代数的同构，在过滤范畴中可能不是同构。
\end{example}

\begin{definition}[严格同态]
带递减过滤的有限维 $F$-向量空间组成范畴 $\operatorname{Fil}_F$。对过滤同态 $T:V'\to V$，按范畴定义
\[
\operatorname{Img}(T)=\Ker(V\to\operatorname{Cok}T),
\qquad
\operatorname{Coim}(T)=\operatorname{Cok}(\Ker T\to V').
\]
若自然态射 $\operatorname{Coim}(T)\to\operatorname{Img}(T)$ 是过滤同构，则称 $T$ 是严格同态。
\end{definition}

\begin{proposition}[严格性的逐层判据]
过滤同态 $T:V'\to V$ 是严格的，当且仅当对每个 $i\in\Z$ 都有
\[
T(\Fil^i(V'))=T(V')\cap\Fil^i(V).
\]
\end{proposition}

\begin{proof}
底层向量空间中，余像和像都可看作 $T(V')$。差别只在过滤。余像继承的第 $i$ 层来自 $\Fil^i(V')$ 的像，即 $T(\Fil^i(V'))$；像作为 $V$ 的子空间继承的第 $i$ 层是 $T(V')\cap\Fil^i(V)$。自然态射在底层上是恒等映射 $T(V')\to T(V')$。它是过滤同构，正是要求两边在每一层给出的子空间完全相等。
\end{proof}

这个判据解释了为什么严格性在过滤线性代数中不可省去：保持过滤只给出包含关系 $T(\Fil^i(V'))\subset T(V')\cap\Fil^i(V)$，严格性要求目标像中的第 $i$ 层没有来自源空间较浅层的额外元素。

\begin{definition}[过滤正合列]
在 $\operatorname{Fil}_F$ 中，一个列 $V'\to V\to V''$ 称为正合，若对每个 $i\in\Z$，普通向量空间列
\[
\Fil^i(V')\longrightarrow \Fil^i(V)\longrightarrow \Fil^i(V'')
\]
都是正合列。
\end{definition}

\begin{proposition}[平移、对偶与张量积]
设 $V,V'$ 是带递减过滤的有限维 $F$-向量空间。
\begin{enumerate}[label=(\arabic*)]
\item 对 $n\in\Z$，定义 $V[n]$ 的底层空间为 $V$，过滤为 $\Fil^i(V[n])=\Fil^{i+n}(V)$。
\item 对偶空间 $V^\vee=\Hom_F(V,F)$ 带过滤
\[
\Fil^i(V^\vee)=\{f\in V^\vee:\Fil^{1-i}(V)\subset\Ker f\}.
\]
在这个约定下有自然过滤同构 $V[n]^\vee\simeq V^\vee[-n]$。
\item 张量积 $V\otimes_FV'$ 带过滤
\[
\Fil^m(V\otimes_FV')
=\sum_{p+q=m}\Fil^p(V)\otimes_F\Fil^q(V').
\]
\end{enumerate}
\end{proposition}

\begin{proof}
(1) 只是在整数指标上整体平移。若 $i+1+n\ge i+n$，则 $\Fil^{i+1+n}(V)\subset\Fil^{i+n}(V)$，所以仍为递减过滤。

(2) 底层对偶空间没有改变，只需比较过滤条件。$f\in\Fil^i(V[n]^\vee)$ 等价于 $\Fil^{1-i}(V[n])\subset\Ker f$，也就是 $\Fil^{1-i+n}(V)\subset\Ker f$。另一方面，$f\in\Fil^i(V^\vee[-n])$ 等价于 $f\in\Fil^{i-n}(V^\vee)$，也就是 $\Fil^{1-(i-n)}(V)=\Fil^{1-i+n}(V)\subset\Ker f$。两者条件一致。

(3) 先验证递减性。若 $p+q=m+1$，则把 $p$ 保持不变并把 $q$ 改为 $q-1$，有 $p+(q-1)=m$，且 $\Fil^q(V')\subset\Fil^{q-1}(V')$。因此每个生成子空间 $\Fil^p(V)\otimes\Fil^q(V')$ 都包含在某个总次数为 $m$ 的生成子空间中，得到 $\Fil^{m+1}\subset\Fil^m$。这个定义使纯张量的层数相加：若 $v$ 处在第 $p$ 层而 $v'$ 处在第 $q$ 层，则 $v\otimes v'$ 处在第 $p+q$ 层。
\end{proof}

\section{无穷扩张}
本节的内容在形式上比前面抽象一些。它的任务不是立刻解决具体数论问题，而是准备两种以后反复出现的语言：第一，无穷对象往往要由有限层的逆极限描述；第二，无穷 Galois 群不是离散群，而是带有自然拓扑的射影有限群。读这一节时可以先抓住一个模型：
\[
\Z_p=\varprojlim_n \Z/p^n\Z.
\]
一个 $p$ 进整数不是某个单独的模 $p^n$ 剩余类，而是一整串彼此相容的模 $p^n$ 剩余类。无穷 Galois 群也是同样思想：它由所有有限 Galois 层的群相容地拼起来。

\subsection{逆极限}

\paragraph{有向集与共尾子集}
一个有向集是一个非空集合 $I$，配备关系 $\le$，满足：
\[
a\le a;
\]
若 $a\le b$ 且 $b\le c$，则 $a\le c$；并且对任意 $a,b\in I$，存在 $c\in I$ 使 $a\le c$ 且 $b\le c$。第三条的意思是：任意两个指标都能被某个更大的指标同时支配。

若 $C\subset I$ 满足：对任意 $a\in I$，存在 $c\in C$ 使 $a\le c$，则称 $C$ 为共尾子集。共尾子集的作用是“没有漏掉足够靠后的信息”。例如用正整数作为指标时，偶数集合是共尾的，因为每个整数之后都有偶数。

\paragraph{拓扑群的逆系统}
以有向集 $I$ 为指标的拓扑群逆系统，是一族拓扑群 $G_i$，以及当 $i\le j$ 时给定连续群同态
\[
\mu_{ji}:G_j\longrightarrow G_i,
\]
满足：
\[
\mu_{ji}\circ\mu_{kj}=\mu_{ki}\qquad(i\le j\le k),
\]
以及
\[
\mu_{ii}=\operatorname{id}_{G_i}.
\]
这些条件说：从更细层 $G_k$ 投到较粗层 $G_i$，不管直接投还是先经过中间层 $G_j$，结果一样。

逆极限 $\varprojlim_i G_i$ 可以显式写成直积中的相容元组集合：
\[
\varprojlim_i G_i
=\left\{(g_i)\in\prod_iG_i:\mu_{ji}(g_j)=g_i\ \text{whenever }i\le j\right\}.
\]
它带有从直积继承的子空间拓扑。投影到第 $i$ 个坐标给出连续同态
\[
\mu_i:\varprojlim_jG_j\longrightarrow G_i.
\]
逆极限的泛性质是：若一个拓扑群 $H$ 同时有相容的连续同态 $\psi_i:H\to G_i$，即 $\mu_{ji}\psi_j=\psi_i$，则存在唯一连续同态
\[
\Psi:H\longrightarrow \varprojlim_iG_i
\]
使得 $\mu_i\Psi=\psi_i$。从元素上看，$\Psi(h)=(\psi_i(h))_i$。

把拓扑群替换为拓扑环、模或其他带结构对象，也能得到相应的逆极限。后面出现的 Witt 向量、完备化、Galois 群、单位群过滤，都可以用这种“相容有限层”的方式理解。

\paragraph{射影有限群}
有限群的逆极限称为射影有限群，也叫 profinite group。它们既有群结构，又有拓扑结构。一个基本事实是：拓扑群 $G$ 是射影有限群，当且仅当 $G$ 是紧的、完全不连通的拓扑群。

若 $G$ 是射影有限群，令 $\{U_i\}$ 遍历 $G$ 的开正规子群且 $[G:U_i]<\infty$。则有自然同构
\[
G\simeq \varprojlim_i G/U_i.
\]
这句话的意思是：射影有限群可以由所有有限商群恢复。元素 $g\in G$ 被所有有限商 $G/U_i$ 中的像共同记录；若两个元素在所有有限商中像相同，则它们相同。

为了描述无限阶，先引入形式无穷积
\[
\prod_p p^{n_p},
\]
其中 $p$ 遍历素数，$n_p$ 是非负整数或 $\infty$。对交换挠群 $A$，定义 $|A|$ 为所有有限子群阶的最小公倍。若 $G$ 是射影有限群、$H$ 是闭子群，则 $[G:H]$ 定义为有限商中
\[
[G/U:H/(H\cap U)]
\]
的最小公倍，其中 $U$ 遍历开正规子群。这个定义把有限指数推广到了无限射影有限情形。

\paragraph{以自然数为指标的逆系统与 ML 条件}
若指标集是非负整数 $\N$，逆系统可写得更具体。设 $\mathcal A$ 是 Abel 范畴，一个逆系统就是一列对象与态射
\[
\cdots\longrightarrow A_{n+1}\xrightarrow{\delta_n}A_n\longrightarrow\cdots\longrightarrow A_0.
\]
若 $i<j$，记复合
\[
\mu_{ji}=\delta_i\circ\delta_{i+1}\circ\cdots\circ\delta_{j-1}:A_j\to A_i.
\]
称逆系统满足 Mittag-Leffler 条件，简称 ML 条件，若对每个固定的 $i$，当 $j$ 足够大时，像
\[
\operatorname{Im}(\mu_{ji}:A_j\to A_i)
\]
稳定下来。称它是 ML 零，若对每个 $i$，存在足够大的 $j$ 使 $\mu_{ji}=0$。

ML 条件的直觉是：向下投影到某一层时，后面越来越远的层不会继续制造新的像。它是控制逆极限右导出函子的技术条件。初学时不必把导出函子细节全部展开，但要知道：没有 ML 条件，逆极限不是精确函子，可能出现 $\varprojlim^1$ 这样的误差项。

这里最容易混淆的是“逆极限存在”和“逆极限保持正合”这两件事。存在性只需要把相容元组放进直积中；正合性则问：如果每一层都有方程可解，能否相容地选出所有层的解。ML 条件正是保证这种相容选择不会在无限下降时失控的一种条件。

\begin{proposition}
设 Abel 范畴 $\mathcal A,\mathcal B$ 有足够内射对象，$F:\mathcal A\to\mathcal B$ 为左正合函子。它诱导函子
\[
F^\N:\mathcal A^\N\longrightarrow \mathcal B^\N.
\]
以 $\varprojlim F$ 记合成 $\varprojlim_n\circ F^\N$。若逆系统 $(A_n,\delta_n)$ 是 ML 零，则
\[
R^p(\varprojlim F)(A_n,\delta_n)=0,\qquad p\ge0.
\]
\end{proposition}

\begin{proof}
这是本章作为同调代数工具调用的消失命题，不在预备章内重证。需要掌握的是各个假设的作用：有足够内射对象使右导出函子 $R^p$ 可以定义；$F$ 左正合保证可以先取内射分解再讨论右导出；ML 零表示每一固定层 $A_i$ 从足够远处接收不到非零信息。于是逆系统对 $\varprojlim$ 以及它的高阶右导出都不给出贡献。后面真正用到此命题时，使用点只会是“某个逆系统满足 ML 零，因此相应导出极限消失”。
\end{proof}

\begin{proposition}
设 $R$ 是环，$\operatorname{Mod}_R$ 为 $R$-模范畴。对逆系统
\[
\cdots\longrightarrow M_{n+1}\xrightarrow{\delta_n}M_n\longrightarrow\cdots\longrightarrow M_0
\]
有正合序列
\[
0\longrightarrow \varprojlim_n M_n
\longrightarrow \prod_n M_n
\xrightarrow{\operatorname{id}-\delta}
\prod_n M_n
\longrightarrow \varprojlim\nolimits^1_n M_n
\longrightarrow 0.
\]
其中
\[
(\operatorname{id}-\delta)((x_n)_n)
=\bigl(x_n-\delta_n(x_{n+1})\bigr)_n.
\]
并且 $p\ge2$ 时 $\varprojlim^p_n M_n=0$。若 $(M_n)$ 满足 ML 条件，则 $\varprojlim^1_n M_n=0$。
\end{proposition}

\begin{proof}
先看正合序列前半段。一个元组 $(x_n)\in\prod_nM_n$ 属于 $\ker(\operatorname{id}-\delta)$，当且仅当
\[
x_n=\delta_n(x_{n+1})\qquad\text{对所有 }n.
\]
这正是逆极限中“相容元组”的定义，所以核就是 $\varprojlim_nM_n$。

后半段的余核定义了 $\varprojlim^1$：若一个元组 $(y_n)$ 不能写成
\[
y_n=x_n-\delta_n(x_{n+1})\qquad(n\ge0),
\]
它就代表一个不能通过相容选择消去的误差类。命题说，对于以自然数为指标的模逆系统，所有高于 $1$ 阶的误差都消失；若满足 ML 条件，连一阶误差也消失。后两句属于导出极限的一般理论，本章把它作为外部输入。读者在后文遇到它时，只需要检查所给逆系统是否满足 ML 条件，随后把 $\varprojlim^1=0$ 当作已经建立的工具使用。
\end{proof}

\subsection{无穷 Galois 扩张}

设 $E/F$ 是代数扩张，次数可以无限。记 $\operatorname{Aut}_F(E)$ 为所有固定 $F$ 的域自同构。若
\[
\{x\in E:\varphi(x)=x,\ \forall\varphi\in\operatorname{Aut}_F(E)\}=F,
\]
则称 $E/F$ 为 Galois 扩张，并记
\[
\Gal(E/F)=G_{E/F}=\operatorname{Aut}_F(E).
\]
这一定义与有限情形兼容：有限时“固定域正好是 $F$”等价于正规可分。

设 $\{K_i:i\in I\}$ 是 $E$ 内所有 $F$ 的有限 Galois 子扩张。它是有向的，因为两个有限 Galois 子扩张 $K_i,K_j$ 的合成域仍是有限 Galois 扩张。并且
\[
E=\bigcup_iK_i=\varinjlim_iK_i.
\]
这里的等式需要解释。若 $x\in E$，由于 $E/F$ 是代数且 Galois，有限扩张 $F(x)/F$ 可嵌入某个有限 Galois 子扩张中，例如取 $x$ 在所有 $F$-嵌入下的共轭并生成的有限正规可分扩张。因此每个元素都会落在某个 $K_i$ 中。
若 $K_i\subset K_j$，限制映射给出投射
\[
\mu_{ji}:\Gal(K_j/F)\longrightarrow \Gal(K_i/F).
\]
于是 $\{\Gal(K_i/F),\mu_{ji}\}$ 构成有限群的逆系统。

\begin{proposition}
设 $E/F$ 是无穷 Galois 扩张，$K_i$ 遍历 $E/F$ 的有限 Galois 子扩张。则
\[
G_{E/F}\simeq \varprojlim_i \Gal(K_i/F).
\]
\end{proposition}

\begin{proof}
对每个 $i$，限制给出同态
\[
G_{E/F}\longrightarrow \Gal(K_i/F).
\]
这些限制彼此相容，因此由逆极限的泛性质得到同态
\[
\Psi:G_{E/F}\longrightarrow \varprojlim_i\Gal(K_i/F).
\]

证明单射。若 $g\ne1$，则存在 $x\in E$ 使 $g(x)\ne x$。取一个有限 Galois 子扩张 $K_i$ 包含 $x$。限制 $g|_{K_i}$ 非平凡，所以 $\Psi(g)$ 的第 $i$ 个坐标非平凡，故 $\Psi(g)\ne1$。

证明满射。取相容元组 $(g_i)_i$，其中 $g_i\in\Gal(K_i/F)$。对 $x\in E$，选 $K_i$ 使 $x\in K_i$，定义
\[
g(x)=g_i(x).
\]
若 $x$ 同时属于另一个 $K_j$，取包含 $K_iK_j$ 的有限 Galois 子扩张 $K_\ell$。相容性保证 $g_\ell$ 限制到 $K_i$ 和 $K_j$ 分别是 $g_i$ 和 $g_j$，所以定义与选择无关。由每个 $g_i$ 都是域自同构得到 $g$ 保持加法、乘法和逆元，并固定 $F$。相容元组的逆元给出 $g$ 的逆映射，因此 $g\in G_{E/F}$，且 $\Psi(g)=(g_i)$。所以 $\Psi$ 满射。
\end{proof}

利用这个同构，把逆极限拓扑转移到 $G_{E/F}$ 上。于是 $G_{E/F}$ 是射影有限群，其单位元的一组邻域基为
\[
\{G_{E/K_i}:K_i/F\text{ 有限 Galois}\}.
\]
这说明无穷 Galois 群中的“小邻域”就是“在某个有限层上看不出来的自同构”。

\begin{theorem}
设 $E/F$ 是 Galois 扩张。令 $\mathcal F$ 为所有中间域 $L$：
\[
F\subset L\subset E,
\]
令 $\mathcal G$ 为 $G_{E/F}$ 的闭子群集合。对 $L\in\mathcal F$，设
\[
A(L)=G_{E/L};
\]
对 $H\in\mathcal G$，设
\[
I(H)=\{x\in E:\varphi(x)=x,\ \forall\varphi\in H\}.
\]
则 $A$ 与 $I$ 给出互反双射
\[
\mathcal F\longleftrightarrow \mathcal G.
\]
\end{theorem}

\begin{proof}
先说明 $A(L)$ 是闭子群。把 $L$ 写成其有限子扩张的并：
\[
L=\bigcup_j L_j.
\]
固定 $L$ 等价于同时固定所有 $L_j$，所以
\[
G_{E/L}=\bigcap_jG_{E/L_j}.
\]
每个有限扩张 $L_j/F$ 都包含在某个有限 Galois 子扩张 $K_j/F$ 中；$G_{E/K_j}$ 是开子群，而 $G_{E/L_j}$ 是包含它的子群，因此也是开子群。拓扑群中的开子群是闭子群，闭子群的交仍闭，所以 $G_{E/L}$ 闭。

接着证明 $I(A(L))=L$。包含 $L\subset I(A(L))$ 来自定义。反向包含的关键是：若 $x\in E\setminus L$，则 $x$ 不可能被所有固定 $L$ 的自同构固定。取包含 $x$ 的有限 Galois 子扩张 $K/F$，令 $M=K\cap L$。由于 $K/F$ 有限 Galois，有限 Galois 理论用于 $K/M$ 给出
\[
K^{\Gal(K/M)}=M.
\]
而 $x\notin M$，所以存在 $\sigma\in\Gal(K/M)$ 使 $\sigma(x)\ne x$。无限 Galois 扩张中的延拓定理允许把 $\sigma$ 延拓为 $E$ 的一个固定 $L$ 的自同构；于是 $x$ 不属于 $I(A(L))$。

最后取闭子群 $H$，令 $T=A(I(H))=G_{E/I(H)}$。按定义 $H\subset T$。要证明相等，只需证明 $H$ 在 $T$ 中稠密，因为 $H$ 已闭。取 $T$ 的任意开正规子群 $V$。令 $L=I(V)$，则 $T/V$ 可看作有限 Galois 扩张 $L/I(H)$ 的 Galois 群。在这个有限层中，$H$ 与 $T$ 有同一个固定域 $I(H)$，有限 Galois 对应给出
\[
T/V=HV/V.
\]
这对每个开正规 $V$ 成立，说明 $H$ 在 $T$ 的每个有限商中都满。射影有限群的拓扑由这些有限商决定，因此 $H$ 在 $T$ 中稠密。因为 $H$ 闭，得到 $H=T$。
\end{proof}

\paragraph{可分闭包与绝对 Galois 群}
选定 $F$ 的代数闭包 $F^{\mathrm{alg}}$。在其中取所有有限可分扩张 $E/F$ 的并，记为
\[
F^{\mathrm{sep}}.
\]
它称为 $F$ 的可分闭包。扩张 $F^{\mathrm{sep}}/F$ 是 Galois 扩张，其 Galois 群
\[
\Gal(F^{\mathrm{sep}}/F)
\]
称为 $F$ 的绝对 Galois 群。它的单位元邻域基为
\[
\{\Gal(F^{\mathrm{sep}}/E):E/F\text{ 有限可分}\}.
\]
因为 $F^{\mathrm{alg}}/F^{\mathrm{sep}}$ 是纯不可分扩张，$F^{\mathrm{sep}}$ 的任一自同构可以唯一延拓到 $F^{\mathrm{alg}}$。因此也可把
\[
\Gal(F^{\mathrm{sep}}/F)
\]
看作 $\operatorname{Aut}_F(F^{\mathrm{alg}})$ 的子群。

若 $E/F$ 是域扩张，取 $E$ 的代数闭包 $E^{\mathrm{alg}}$，可把 $F$ 在其中的代数闭包记作 $F^{\mathrm{alg}}$，于是 $F^{\mathrm{sep}}\subset E^{\mathrm{sep}}$。限制自同构给出连续同态
\[
r:\Gal(E^{\mathrm{sep}}/E)\longrightarrow \Gal(F^{\mathrm{sep}}/F).
\]
其像是闭子群。若 $E/F$ 是代数扩张，则可取同一个代数闭包，$r$ 是单射，于是 $\Gal(E^{\mathrm{sep}}/E)$ 可视为 $\Gal(F^{\mathrm{sep}}/F)$ 的闭子群。若 $E/F$ 是有限扩张，则这个闭子群还是开子群。

\subsection{完备化}

设 $M$ 是交换群，并给定递减子群列
\[
M=M_0\supset M_1\supset M_2\supset\cdots.
\]
以陪集 $x+M_i$ 作为点 $x$ 的邻域基，就在 $M$ 上得到一个群拓扑。加法与减法连续，因为若两个元素分别在足够小的陪集中变化，它们的和与差也落在相应陪集中。这个拓扑是 Hausdorff 的，当且仅当
\[
\bigcap_{i=0}^\infty M_i=0.
\]
直观上，若交不为零，那么非零元素会落在零的所有邻域中，拓扑无法把它和零分开。

序列 $x_n$ 收敛到 $a$，意思是：对每个 $r>0$，存在 $n_r$，使得
\[
n\ge n_r\quad\Longrightarrow\quad x_n-a\in M_r.
\]
序列 $x_n$ 是 Cauchy 序列，意思是：对每个 $r>0$，存在 $n_r$，使得
\[
n,m\ge n_r\quad\Longrightarrow\quad x_n-x_m\in M_r.
\]
称 $M$ 为完备拓扑群，若它是 Hausdorff 的，并且每个 Cauchy 序列都收敛。

\begin{theorem}
给定递减过滤
\[
M=M_0\supset M_1\supset M_2\supset\cdots
\]
且 $\bigcap_iM_i=0$。对 $i\ge j$ 有自然同态 $M/M_i\to M/M_j$。令
\[
\widehat M=\varprojlim_i M/M_i.
\]
则：
\begin{enumerate}[label=(\arabic*)]
\item $\widehat M$ 是完备拓扑群。
\item 自然映射
\[
\varphi:M\longrightarrow \widehat M,\qquad x\longmapsto (x\bmod M_i)_i
\]
是单同态。
\item $\varphi(M)$ 在 $\widehat M$ 中稠密。
\end{enumerate}
\end{theorem}

\begin{proof}
先看单射。若 $\varphi(x)=0$，则 $x\in M_i$ 对所有 $i$ 成立，所以 $x\in\bigcap_iM_i=0$。因此 $\varphi$ 单。

再看稠密性。$\widehat M$ 的一个基本开邻域由有限层同余条件决定。给定 $(\bar x_i)_i\in\widehat M$ 和层数 $r$，取 $x\in M$ 代表 $\bar x_r\in M/M_r$。则 $\varphi(x)$ 与 $(\bar x_i)$ 在第 $r$ 层相同，也就是落在该点的 $r$ 层邻域中。因此 $\varphi(M)$ 与每个基本邻域相交，故稠密。

完备性来自逆极限的构造。一个 $\widehat M$ 中的 Cauchy 序列在每个商群 $M/M_i$ 中最终稳定：这是因为第 $i$ 层邻域的定义正是“第 $i$ 个坐标相同”。把第 $i$ 个坐标的最终稳定值记作 $\bar x_i$。若 $i\ge j$，从足够靠后的项同时读取第 $i$ 与第 $j$ 个坐标，相容性给出 $\bar x_i\bmod M_j=\bar x_j$。于是 $(\bar x_i)_i$ 是逆极限中的元素。原 Cauchy 序列按定义收敛到这个相容元组。Hausdorff 性也由坐标分辨得到：若两个相容元组不同，则它们在某个 $M/M_i$ 中的坐标不同，所以可用第 $i$ 层邻域分开。
\end{proof}

\begin{proposition}
设递减过滤
\[
M=M_0\supset M_1\supset M_2\supset\cdots
\]
在 $M$ 上决定完备 Hausdorff 拓扑，另有
\[
M'=M'_0\supset M'_1\supset M'_2\supset\cdots
\]
也有同样性质。设同态 $u:M\to M'$ 满足 $u(M_n)\subset M'_n$。若对所有 $n$，由 $u$ 决定的同态
\[
u_n:M_n/M_{n+1}\longrightarrow M'_n/M'_{n+1}
\]
都是单射，则 $u$ 是单射；若所有 $u_n$ 都是满射，则 $u$ 是满射。
\end{proposition}

\begin{proof}
先证单射。设 $x\in M$ 且 $u(x)=0$。若 $x\ne0$，由 Hausdorff 性，存在 $r$ 使 $x\notin M_r$。由于 $x\in M_0$，可在有限集合 $\{0,\dots,r-1\}$ 中取最大的 $n$，使 $x\in M_n$。于是 $x\notin M_{n+1}$，所以 $x$ 在 $M_n/M_{n+1}$ 中给出非零类。它经 $u_n$ 映到 $0$，这与 $u_n$ 单射矛盾。因此 $x=0$，$u$ 单射。

再证满射。取 $y\in M'$。先由 $u_0$ 满射，选 $x_0\in M$ 使
\[
y-u(x_0)\in M'_1.
\]
再由 $u_1$ 满射，选 $z_1\in M_1$ 使
\[
y-u(x_0)-u(z_1)\in M'_2.
\]
继续下去，构造 $z_n\in M_n$，使部分和
\[
s_N=x_0+z_1+\cdots+z_N
\]
满足 $y-u(s_N)\in M'_{N+1}$。由于 $z_n\in M_n$，部分和 $s_N$ 是 $M$ 中的 Cauchy 序列；完备性给出极限 $x\in M$。连续性和误差趋于零推出 $u(x)=y$。因此 $u$ 满射。
\end{proof}

\section{特征标}
特征标把群元素映到复数乘法群中，是表示论和 $L$ 函数的共同入口。对交换群，特征标就是一维表示；对一般有限群，特征标是表示矩阵的迹；对 Galois 群，连续特征标会对应有限循环扩张。第十二章的 Artin $L$ 函数正是把 Galois 表示的特征标放进 Euler 乘积中。

\subsection{有限交换群的对偶群}

若群 $G$ 的乘法满足 $xy=yx$，则称 $G$ 为交换群或 Abel 群。本节以 $Z_m$ 记 $m$ 阶循环群，它同构于 $\Z/m\Z$。

\begin{proposition}
设有限交换群 $G$ 的阶有素因子分解
\[
|G|=p_1^{e_1}\cdots p_r^{e_r}.
\]
以 $S(p)$ 记 $G$ 的 $p$-Sylow 子群。则：
\begin{enumerate}[label=(\arabic*)]
\item $G=S(p_1)\times\cdots\times S(p_r)$。
\item 每个 $S(p_i)$ 可分解为循环 $p_i$-群的直积：
\[
S(p_i)\simeq Z_{p_i^{e_{i1}}}\times\cdots\times Z_{p_i^{e_{is}}},
\qquad e_{i1}+\cdots+e_{is}=e_i.
\]
\end{enumerate}
\end{proposition}

\begin{proof}
有限 Abel 群中，不同素数幂阶部分彼此独立。令
\[
S(p)=\{g\in G:g^{p^a}=1\text{ for some }a\}.
\]
它是子群，因为 $G$ 交换。任取 $g\in G$，其阶分解为不同素数幂的乘积。用中国剩余定理可把 $g$ 唯一写成各个 $p_i$-部分的乘积，因此 $G$ 是这些 Sylow 子群的直积。第二点是有限 Abel $p$-群的结构定理：每个有限 Abel $p$-群可唯一分解为若干循环 $p$-幂阶群的直积；阶相乘给出指数和等于 $e_i$。
\end{proof}

绝对值等于 $1$ 的复数组成的交换群记为 $S^1$。设 $G$ 为有限交换群。群同态
\[
\chi:G\longrightarrow S^1
\]
称为 $G$ 的酉特征标。有限群情形下常简称为特征标。特征标的乘法逐点定义：
\[
(\chi\psi)(g)=\chi(g)\psi(g).
\]
于是 $G$ 的所有特征标组成交换群，记为 $G^*$ 或 $\widehat G$，称为特征标群或对偶群。

有限交换群的特征标有两个基础作用。第一，它们能分离群元素：非单位元总能被某个特征标检测出来。第二，它们给出有限群上的 Fourier 分析：普通函数可以按特征标展开。下面先证明第一点，后面的高斯和正是把这种 Fourier 分析用到有限域上。

一个基本例子是 Legendre 符号。设 $p$ 为奇素数，$p\nmid a$，定义
\[
\left(\frac ap\right)=
\begin{cases}
1,&\text{若存在整数 }x\text{ 使 }x^2\equiv a\pmod p,\\
-1,&\text{否则}.
\end{cases}
\]
它给出 $\F_p^\times$ 上的二次特征标。Gauss 二次互反律说：若 $p,q$ 是不同奇素数，则
\[
\left(\frac qp\right)\left(\frac pq\right)
=(-1)^{\frac{p-1}{2}\frac{q-1}{2}}.
\]
这说明特征标不是纯代数记号，它携带深刻的同余信息。

\begin{lemma}
\begin{enumerate}[label=(\arabic*)]
\item 若 $G$ 是有限交换群且 $g\ne1$，则存在非平凡特征标 $\chi\in G^*$ 使 $\chi(g)\ne1$。
\item 若 $A,B$ 是有限交换群，则
\[
(A\times B)^*\simeq A^*\times B^*.
\]
\end{enumerate}
\end{lemma}

\begin{proof}
先证循环群情形。若 $G=Z_m=\langle u\rangle$，特征标由 $\chi(u)$ 决定。由于 $u^m=1$，$\chi(u)$ 必须是 $m$ 次单位根；反过来，任意 $m$ 次单位根 $\zeta$ 都定义一个特征标 $\chi(u^a)=\zeta^a$。若 $g=u^k\ne1$，则 $m\nmid k$。取本原 $m$ 次单位根 $\zeta$，则 $\zeta^k\ne1$，令 $\chi(u)=\zeta$ 便得到 $\chi(g)\ne1$。

一般有限交换群分解为循环群直积。非平凡元素 $g$ 在某个循环因子上的坐标非平凡；在该因子上取能检测它的特征标，在其他因子上取平凡特征标，再相乘得到 $G$ 上的特征标。

第二点中，给定 $A\times B$ 的特征标 $\chi$，限制到 $A\times\{1\}$ 与 $\{1\}\times B$ 分别给出 $A$ 与 $B$ 的特征标；反过来，若 $\alpha\in A^*$、$\beta\in B^*$，则
\[
(a,b)\longmapsto \alpha(a)\beta(b)
\]
是 $A\times B$ 的特征标。这两个构造互逆。
\end{proof}

\begin{proposition}[有限交换群的特征标正交关系]
设 $G$ 是有限交换群。则：
\begin{enumerate}[label=(\arabic*)]
\item 若 $\chi\in G^*$ 非平凡，则 $\sum_{g\in G}\chi(g)=0$。
\item 若 $g\in G$ 且 $g\ne1$，则 $\sum_{\chi\in G^*}\chi(g)=0$；若 $g=1$，该和等于 $|G|$。
\end{enumerate}
\end{proposition}

\begin{proof}
对 (1)，因为 $\chi$ 非平凡，存在 $h\in G$ 使 $\chi(h)\ne1$。记 $S=\sum_{g\in G}\chi(g)$。乘以 $\chi(h)$ 得
\[
\chi(h)S=\sum_{g\in G}\chi(h)\chi(g)=\sum_{g\in G}\chi(hg).
\]
当 $g$ 遍历 $G$ 时，$hg$ 也遍历 $G$，所以右边仍为 $S$。于是 $(\chi(h)-1)S=0$，而 $\chi(h)-1\ne0$，故 $S=0$。

对 (2)，把 $g$ 看成 $\widehat G$ 上的评价特征标 $e_g(\chi)=\chi(g)$。若 $g\ne1$，前一引理说明 $e_g$ 不是平凡特征标，对群 $\widehat G$ 应用 (1) 得 $\sum_{\chi\in G^*}\chi(g)=0$。若 $g=1$，每个 $\chi(1)$ 都等于 $1$，和为 $|G^*|=|G|$。
\end{proof}

取 $g\in G$，定义 $\widehat G$ 上的函数
\[
e_g:\widehat G\longrightarrow S^1,\qquad e_g(\chi)=\chi(g).
\]

\begin{proposition}
设 $G$ 为有限交换群。则：
\begin{enumerate}[label=(\arabic*)]
\item $G\simeq G^*$。
\item Pontryagin 对偶映射
\[
G\longrightarrow G^{**},\qquad g\longmapsto e_g
\]
是自然同构。
\end{enumerate}
\end{proposition}

\begin{proof}
对循环群 $Z_m$，特征标由生成元送到哪个 $m$ 次单位根决定，所以 $Z_m^*\simeq Z_m$。对有限交换群，用结构定理把 $G$ 分解为循环群直积，再用上一引理得到 $G^*$ 也分解为相同阶数的循环群直积，因此 $G\simeq G^*$。

对第二点，映射 $g\mapsto e_g$ 是群同态。若 $g\ne1$，上一引理给出 $\chi$ 使 $\chi(g)\ne1$，于是 $e_g\ne1$，所以该映射单射。由于 $|G|=|G^*|=|G^{**}|$，单射即为同构。
\end{proof}

把以上结果推广到有限交换群的群指数。给定群 $G$，若存在最小正整数 $m$ 使 $g^m=1$ 对所有 $g\in G$ 成立，则称 $m$ 为 $G$ 的群指数。有限群的群指数等于所有元素阶的最小公倍数。

设 $m$ 是有限交换群 $G$ 的群指数，有限循环群 $Z$ 的阶满足 $m\mid |Z|$。则 $\Hom(G,Z)$ 与 $G$ 的特征标群同构。

\begin{proposition}
在上述设定下：
\begin{enumerate}[label=(\arabic*)]
\item $G\simeq \Hom(G,Z)$。
\item $G\simeq \Hom(\Hom(G,Z),Z)$。
\item $\Hom(G,Z)$ 由 $\chi_1,\dots,\chi_r$ 生成，当且仅当：若 $g\in G$ 满足 $\chi_i(g)=1$ 对所有 $i$ 成立，则 $g=1$。
\end{enumerate}
\end{proposition}

\begin{proof}
因为 $Z$ 的阶被 $G$ 的指数整除，$Z$ 中含有足够多的元素来接收每个循环因子的生成元。把 $G$ 分解为循环群直积后，对每个循环因子 $Z_d=\langle u\rangle$，同态 $\varphi:Z_d\to Z$ 由 $\varphi(u)$ 决定，并且 $\varphi(u)$ 必须满足 $\varphi(u)^d=1$。由于 $d$ 整除 $G$ 的指数，从而整除 $|Z|$，循环群 $Z$ 中恰有一个阶为 $d$ 的子群。因此
\[
\Hom(Z_d,Z)\simeq Z_d.
\]
对直积使用 $\Hom(A\times B,Z)\simeq \Hom(A,Z)\times\Hom(B,Z)$，得到 (1)。第二点是对 (1) 再对偶一次。

第 (3) 点是“特征标能否分离点”的判据。若 $\chi_1,\dots,\chi_r$ 生成整个 $\Hom(G,Z)$，而 $g$ 被它们全部送到 $1$，则 $g$ 被所有特征标送到 $1$，由分离性得 $g=1$。

反过来，设这些 $\chi_i$ 生成的子群为 $H\subset\Hom(G,Z)$，并假设共同核为 $\{1\}$。若 $H\ne\Hom(G,Z)$，则商群 $\Hom(G,Z)/H$ 非平凡。有限交换群的特征标能分离点，所以存在非平凡同态 $\Theta:\Hom(G,Z)/H\to Z$。把 $\Theta$ 与商映射合成，得到 $\Hom(G,Z)$ 上一个在 $H$ 上取值 $1$ 的非平凡特征标。由双对偶同构，这个特征标等于某个 $g\in G$ 的评价映射 $\chi\mapsto\chi(g)$。它在所有 $\chi_i$ 上取值 $1$，所以 $g$ 属于共同核，得到 $g=1$；但评价于 $g=1$ 是平凡特征标，与 $\Theta$ 非平凡矛盾。因此 $H=\Hom(G,Z)$。
\end{proof}

\subsection{高斯和}

可以把 $\Z/m\Z$ 上的函数看作 $\Z$ 上周期为 $m$ 的函数。有限 Fourier 展开说：周期函数可以由有限多个指数函数重构。

\begin{proposition}
设 $\eta$ 是 $m$ 次本原单位根，$f(t)$ 是 $\Z$ 上周期为 $m$ 的函数。则
\[
f(t)=\sum_{n=0}^{m-1}c(n)\eta^{nt},
\]
其中
\[
c(n)=\frac1m\sum_{t=0}^{m-1}f(t)\eta^{-nt}.
\]
\end{proposition}

\begin{proof}
把右边展开：
\[
\sum_{n=0}^{m-1}\left(\frac1m\sum_{u=0}^{m-1}f(u)\eta^{-nu}\right)\eta^{nt}
=\frac1m\sum_{u=0}^{m-1}f(u)\sum_{n=0}^{m-1}\eta^{n(t-u)}.
\]
内层和若 $t\equiv u\pmod m$，则等于 $m$；否则是一个公比非 $1$ 的等比和，等于 $0$。因此只剩 $u\equiv t$ 的项，得到 $f(t)$。
\end{proof}

若 $\chi$ 是 $\F_p^\times$ 的特征标，约定 $\chi(0)=0$，把它看成 $\F_p$ 上的函数。取 $p$ 次本原单位根 $\zeta=e^{2\pi i/p}$。称
\[
g_t(\chi)=\sum_{n\in\F_p}\chi(n)\zeta^{nt}
\]
为高斯和，并记 $g(\chi)=g_1(\chi)$。当 $\chi$ 是 Legendre 符号时，常写
\[
g_t=\sum_{n\in\F_p}\left(\frac np\right)\zeta^{nt}.
\]
高斯和把乘法信息 $\chi(n)$ 与加法振荡 $\zeta^{nt}$ 混合在一起，是二次互反律、函数方程和局部常数的早期模型。

\begin{lemma}
若 $\chi$ 是 $\F_p^\times$ 的非平凡特征标，则 $g(\chi)\ne0$，并且 $|g(\chi)|^2=p$。
\end{lemma}

\begin{proof}
因为 $\chi$ 取值在单位圆上，且 $\overline{\chi(b)}=\chi(b)^{-1}$。按约定 $\chi(0)=0$，所以
\[
|g(\chi)|^2
=\sum_{a,b\in\F_p^\times}\chi(a)\chi(b)^{-1}\zeta^{a-b}.
\]
令 $a=tb$，其中 $t\in\F_p^\times$。上式变成
\[
\sum_{t\in\F_p^\times}\chi(t)
\sum_{b\in\F_p^\times}\zeta^{b(t-1)}.
\]
当 $t=1$ 时，内层和为 $p-1$。当 $t\ne1$ 时，$b(t-1)$ 随 $b\in\F_p^\times$ 遍历 $\F_p^\times$，所以内层和为 $\sum_{c\in\F_p^\times}\zeta^c=-1$。因此
\[
|g(\chi)|^2
=(p-1)-\sum_{t\in\F_p^\times,\ t\ne1}\chi(t).
\]
非平凡特征标的总和为 $0$，故 $\sum_{t\ne1}\chi(t)=-1$，于是 $|g(\chi)|^2=p$。特别地 $g(\chi)\ne0$。
\end{proof}

若 $\chi,\lambda$ 是 $\F_p^\times$ 的特征标，定义 Jacobi 和
\[
J(\chi,\lambda)=\sum_{a+b=1}\chi(a)\lambda(b).
\]
\begin{proposition}[Gauss 和与 Jacobi 和的关系]
设 $\chi,\lambda$ 是 $\F_p^\times$ 的非平凡特征标，并按约定在 $0$ 处取值 $0$。当 $\chi\lambda\ne1$ 时，有
\[
J(\chi,\lambda)=\frac{g(\chi)g(\lambda)}{g(\chi\lambda)},
\]
而当 $\lambda=\chi^{-1}$ 时，
\[
J(\chi,\chi^{-1})=-\chi(-1).
\]
\end{proposition}

\begin{proof}
先计算乘积：
\[
g(\chi)g(\lambda)
=\sum_{a,b\in\F_p}\chi(a)\lambda(b)\zeta^{a+b}.
\]
按 $t=a+b$ 分组。因为 $\chi(0)=\lambda(0)=0$，含有 $a=0$ 或 $b=0$ 的项不会贡献。对固定的 $t\in\F_p^\times$，令 $a=tu$、$b=t(1-u)$，则 $u+(1-u)=1$，并且
\[
\chi(a)\lambda(b)=\chi(t)\lambda(t)\chi(u)\lambda(1-u)
=(\chi\lambda)(t)\chi(u)\lambda(1-u).
\]
因此
\[
g(\chi)g(\lambda)
=\left(\sum_{t\in\F_p^\times}(\chi\lambda)(t)\zeta^t\right)
\left(\sum_{u\in\F_p}\chi(u)\lambda(1-u)\right)
=g(\chi\lambda)J(\chi,\lambda),
\]
其中最后一步使用 $t=0$ 项为 $0$。若 $\chi\lambda\ne1$，非平凡特征标的 Gauss 和不为 $0$，便得到第一式。

再设 $\lambda=\chi^{-1}$。由于 $a+b=1$，项 $a=0$ 与 $b=0$ 都为 $0$，其余项可写为
\[
\chi(a)\chi^{-1}(1-a)=\chi\left(\frac{a}{1-a}\right).
\]
当 $a$ 遍历 $\F_p\setminus\{0,1\}$ 时，$a/(1-a)$ 遍历 $\F_p^\times\setminus\{-1\}$。于是
\[
J(\chi,\chi^{-1})
=\sum_{x\in\F_p^\times,\ x\ne-1}\chi(x)
=-\chi(-1),
\]
因为非平凡特征标在整个 $\F_p^\times$ 上的和为 $0$。
\end{proof}

\subsection{Dirichlet 特征标}

以 $R^\times$ 记环 $R$ 的可逆元群。

\begin{example}
以 $(\Z/m\Z)^\times$ 记环 $\Z/m\Z$ 中可逆元组成的乘法群。
\begin{enumerate}[label=(\arabic*)]
\item 若 $m=p_1^{e_1}\cdots p_r^{e_r}$，则由中国剩余定理
\[
(\Z/m\Z)^\times\simeq(\Z/p_1^{e_1}\Z)^\times\times\cdots\times(\Z/p_r^{e_r}\Z)^\times.
\]
\item 若 $p>2$，则 $(\Z/p^e\Z)^\times$ 是 $\varphi(p^e)=p^{e-1}(p-1)$ 阶循环群。
\item $(\Z/2\Z)^\times=\{1\}$，$(\Z/4\Z)^\times\simeq Z_2$，且 $e\ge3$ 时
\[
(\Z/2^e\Z)^\times\simeq Z_2\times Z_{2^{e-2}}.
\]
\end{enumerate}
\end{example}

模 $m$ 的 Dirichlet 特征标是群同态
\[
\chi:(\Z/m\Z)^\times\longrightarrow \C^\times.
\]
因为 $(\Z/m\Z)^\times$ 有 $\varphi(m)$ 个元素，$\chi$ 的像落在 $\varphi(m)$ 次单位根中，所以也可视为酉特征标。单位元称为主特征标，记为 $1$ 或 $\chi_0$。

传统上把 Dirichlet 特征标看成整数上的函数 $\chi:\Z\to\C$，满足：
\[
\chi(a)=0\quad\text{若 }(a,m)\ne1,
\]
\[
\chi(a+km)=\chi(a)\quad\text{若 }(a,m)=1,
\]
以及
\[
\chi(ab)=\chi(a)\chi(b).
\]
这就是把 $(\Z/m\Z)^\times$ 上的群同态延拓为模 $m$ 周期的乘法函数，并在不可逆剩余类上取零。

若存在 $m'<m$、$m'\mid m$ 和模 $m'$ 的 Dirichlet 特征标 $\chi'$，使
\[
\chi=\chi'\circ\rho_{m,m'},
\]
其中 $\rho_{m,m'}:(\Z/m\Z)^\times\to(\Z/m'\Z)^\times$ 是自然投射，则称 $\chi$ 为非原特征标。不是非原特征标的称为原特征标。给定 $\chi$，使它来自更小模数的最小模数称为它的导子。

对应于 $\chi$ 的 Dirichlet $L$-函数定义为
\[
L(s,\chi)=\sum_{n=1}^\infty\frac{\chi(n)}{n^s},\qquad \operatorname{Re}(s)>1.
\]

\begin{theorem}
Dirichlet $L$-函数有以下性质：
\begin{enumerate}[label=(\arabic*)]
\item 在 $\operatorname{Re}(s)>1$ 上有 Euler 乘积
\[
L(s,\chi)=\prod_p(1-\chi(p)p^{-s})^{-1}.
\]
\item 主特征标 $\chi_0$ 的 $L(s,\chi_0)$ 可解析延拓到 $\C\setminus\{1\}$，且 $s=1$ 是唯一极点，留数为 $\varphi(m)/m$。
\item 若 $\chi\ne\chi_0$，则 $L(s,\chi)$ 可解析延拓为整函数，并且 $L(1,\chi)\ne0$。
\item 若 $\chi$ 是模 $m$ 原特征标，令 $\overline\chi(n)=\overline{\chi(n)}$，
\[
\delta(\chi)=\frac{1-\chi(-1)}2,
\]
\[
g(\chi)=\sum_{n=1}^m\chi(n)e^{2\pi in/m},
\]
\[
\Lambda(s,\chi)=\left(\frac m\pi\right)^{(s+\delta)/2}
\Gamma\left(\frac{s+\delta}{2}\right)L(s,\chi),
\qquad
W(\chi)=\frac{g(\chi)}{i^\delta\sqrt m}.
\]
则有函数方程
\[
\Lambda(s,\chi)=W(\chi)\Lambda(1-s,\overline\chi).
\]
\end{enumerate}
\end{theorem}

\begin{proof}
第一点来自完全乘法性。对 $\operatorname{Re}(s)>1$，级数绝对收敛，把每个整数唯一分解成素数幂，得到
\[
\sum_{n\ge1}\frac{\chi(n)}{n^s}
=\prod_p\left(1+\chi(p)p^{-s}+\chi(p)^2p^{-2s}+\cdots\right)
=\prod_p(1-\chi(p)p^{-s})^{-1}.
\]
后面三点属于解析数论核心定理。主特征标情形本质上是删去有限个 Euler 因子的 Riemann zeta 函数，因此继承 $\zeta(s)$ 在 $s=1$ 的单极点。非主特征标的整延拓和 $L(1,\chi)\ne0$ 是 Dirichlet 算术级数定理的关键解析输入。原特征标的函数方程来自 Poisson 求和与高斯和计算；其中 $\Gamma$ 因子补足无穷素位，$W(\chi)$ 是根数。
\end{proof}

这里的 $\Gamma$ 是复变函数。对 $\operatorname{Re}(s)>0$，
\[
\Gamma(s)=\int_0^\infty x^s e^{-x}\frac{dx}{x},
\]
并可延拓为全复平面的亚纯函数，极点位于 $s=0,-1,-2,\dots$。Dirichlet $L$-函数是所有后续 $L$-函数的原型：Euler 乘积、解析延拓、函数方程、特殊值，都会在更一般的情形中重现。

\subsection{有限群的表示}

设有限群 $G$ 在 $n$ 维复向量空间 $V$ 上有表示
\[
\pi:G\longrightarrow \GL(V).
\]
选基后得到矩阵表示 $\pi_e:G\to\GL_n(\C)$。换基只会把矩阵同时共轭：
\[
\pi_{e'}(g)=T\pi_e(g)T^{-1}.
\]
由于矩阵迹在共轭下不变，函数
\[
\chi_\pi(g)=\operatorname{Tr}(\pi_e(g))
\]
与基的选择无关，称为表示 $\pi$ 的特征标。若 $\pi$ 不可约，则称 $\chi_\pi$ 为不可约特征标。

若函数 $\psi:G\to\C$ 满足
\[
\psi(hgh^{-1})=\psi(g),
\]
则称 $\psi$ 为类函数。特征标总是类函数，因为共轭元素的表示矩阵互相共轭。有限群表示论的基本事实是：每个复值类函数都可唯一写成不可约特征标的复线性组合。

从 $G$ 到 $\C$ 的所有函数组成向量空间 $\C^G$。卷积定义为
\[
(x*y)(g)=\sum_{h\in G}x(h)y(h^{-1}g).
\]
有限形式和
\[
\sum_{g\in G}a_g g
\]
组成的代数称为群代数 $\C[G]$。把 $\sum a_gg$ 看作函数 $g\mapsto a_g$，可把群代数嵌入 $\C^G$。在类函数上使用内积
\[
\langle x,y\rangle=\frac1{|G|}\sum_{g\in G}x(g)\overline{y(g)}.
\]
这个内积使不可约特征标成为一组正交基。

若 $\alpha:H\to G$ 是有限群同态，$G$ 上类函数 $\varphi$ 可限制为 $H$ 上类函数：
\[
\alpha^*(\varphi)=\varphi\circ\alpha.
\]
反方向有诱导类函数 $\alpha_*$，由 Frobenius 互反公式刻画：
\[
\langle \varphi,\alpha_*\psi\rangle_G
=\langle \alpha^*\varphi,\psi\rangle_H.
\]
若 $\rho$ 是 $H$ 的表示，则诱导表示 $\operatorname{Ind}_H^G\rho$ 的特征标就是对应的诱导特征标。

\begin{theorem}
设 $\chi$ 是有限群 $G$ 的特征标。则存在 $G$ 的子群 $H_i$、$H_i$ 的次数 $1$ 的特征标 $\psi_i$ 以及整数 $n_i$，使得
\[
\chi=\sum_i n_i\,\operatorname{Ind}_{H_i}^G\psi_i.
\]
\end{theorem}

\begin{proof}
这是有限群表示论中的深定理，这里作为外部工具使用。它的意义是：一般有限群的特征标虽然可能来自高维表示，但在整数线性组合意义下，可由子群的一维特征标诱导生成。后面处理 Artin $L$-函数时，这一定理可以把一般 Artin $L$-函数化归到更接近 Hecke/Dirichlet 型的一维情形。
\end{proof}

\subsection{Galois 群的特征标}

设 $G$ 是拓扑群。连续同态
\[
\chi:G\longrightarrow\C^\times
\]
称为 $G$ 的特征标；若像落在单位圆 $S^1$ 中，称为酉特征标。$G$ 的所有酉特征标组成群，记为 $\widehat G$。若 $\chi(G)=\mu_n$，称 $\chi$ 为 $n$ 阶酉特征标；若 $\chi(G)=1$，称为平凡特征标。

有两个简单但重要的拓扑事实。第一，若酉特征标 $\chi$ 的值总满足 $\operatorname{Re}\chi(g)>0$，则 $\chi$ 平凡，因为单位圆右半圆内没有非平凡子群。第二，紧拓扑群到 $\C^\times$ 的连续特征标必为酉特征标，因为 $\R_{>0}$ 中唯一紧子群是 $\{1\}$。

若 $G$ 的单位元有一组由子群组成的基本邻域，则称 $G$ 为全非连通拓扑群。

\begin{lemma}
\begin{enumerate}[label=(\arabic*)]
\item 若 $G$ 是全非连通局部紧拓扑群，则 $G$ 的特征标必为局部常值。
\item 若 $G$ 是全非连通紧拓扑群，则 $G$ 的特征标必为有限阶酉特征标。
\item 若 $G$ 是紧交换群，并且 $G$ 的每个特征标都是有限阶酉特征标，则 $G$ 是全非连通的。
\end{enumerate}
\end{lemma}

\begin{proof}
对 (1)，连续性说明 $\chi^{-1}$ 把 $1\in\C^\times$ 附近的小邻域拉回为 $G$ 中单位元附近的开邻域。取一个不含非平凡子群的复数邻域，例如足够小的圆盘与 $\C^\times$ 的交。全非连通性允许在原像中取开子群 $U$。于是 $\chi(U)$ 是落在该小邻域内的子群，只能为 $\{1\}$。因此 $U\subset\ker\chi$，而 $\chi$ 在每个陪集 $gU$ 上取常值 $\chi(g)$。

对 (2)，紧性使 $\chi(G)$ 是 $\C^\times$ 的紧子群，所以落在 $S^1$ 中。由 (1)，$\ker\chi$ 是开子群，因此商 $G/\ker\chi$ 是紧离散群，故有限。于是 $\chi$ 有有限阶。

对 (3)，若紧交换群 $G$ 不是全非连通，则其单位连通分支 $G^0$ 非平凡。紧交换群的 Pontryagin 对偶能分离点，所以存在特征标在 $G^0$ 上非平凡；但连通群到有限离散群的连续像必须是平凡的，这与所有特征标有限阶矛盾。
\end{proof}

\begin{proposition}[特征标与循环扩张]
令 $F^{\mathrm{sep}}$ 为域 $F$ 的可分闭包，$G=\Gal(F^{\mathrm{sep}}/F)$。则 $G$ 的每个复特征标都是有限阶酉特征标。若
\[
\chi:G\to\C^\times
\]
是 $n$ 阶酉特征标，令 $H=\ker\chi$。则 $H$ 是正规开子群，商群 $G/H\simeq \mu_n$ 是 $n$ 阶循环群。令
\[
E=(F^{\mathrm{sep}})^H
=\{x\in F^{\mathrm{sep}}:h(x)=x,\ \forall h\in H\}.
\]
则 $E/F$ 是 $n$ 次循环 Galois 扩张，并且 $\Gal(E/F)\simeq\mu_n$。称 $E$ 依附于 $\chi$。

反过来，若 $E/F$ 是 $F^{\mathrm{sep}}/F$ 内的 $n$ 次循环扩张，则它对应一个正规开子群 $H$，且
\[
G/H\simeq \Gal(E/F)\simeq \mu_n.
\]
把投射 $G\to G/H$ 与这个同构合成，就得到一个 $n$ 阶酉特征标 $\chi$，并且 $\ker\chi=H$。这时称 $\chi$ 依附于 $E$。
\end{proposition}

\begin{proof}
群 $G$ 是射影有限群，因而紧且全非连通。上一引理给出每个连续复特征标都是局部常值；紧性又使其像为有限群并落在单位圆中。因此 $\chi$ 有有限阶且为酉特征标。

给定 $n$ 阶 $\chi$，核 $H$ 是连续映射的闭正规子群。由于 $\chi$ 局部常值，存在单位元开邻域落在核中，故 $H$ 开。无限 Galois 对应把开子群对应到有限中间扩张，所以 $E=(F^{\mathrm{sep}})^H$ 是有限 Galois 扩张，并且
\[
\Gal(E/F)\simeq G/H\simeq\chi(G)=\mu_n.
\]
这说明 $E/F$ 是 $n$ 次循环扩张。

反过来，若 $E/F$ 是 $n$ 次循环 Galois 扩张，则 $H=\Gal(F^{\mathrm{sep}}/E)$ 是 $G$ 的正规开子群，且 $G/H\simeq\Gal(E/F)$。选定一个同构 $\Gal(E/F)\simeq\mu_n$，再与自然商映射合成，得到 $G\to\mu_n\subset S^1$ 的连续特征标。它的核正是 $H$。
\end{proof}

\begin{proposition}[特征标只看交换化]
以 $\overline F$ 记 $F$ 的代数闭包。若 $\overline F$ 内有限 Galois 扩张 $E/F$ 的 Galois 群为交换群，则称 $E/F$ 为交换扩张。令 $F^{\mathrm{ab}}$ 为 $\overline F$ 内所有有限交换扩张的并。它是 $F$ 的最大交换扩张。设 $G^{(1)}$ 是对应于 $F^{\mathrm{ab}}$ 的闭正规子群，则 $G^{(1)}$ 是使 $G/G^{(1)}$ 交换的最小闭正规子群。它等于交换子子群
\[
[G,G]=\langle xyx^{-1}y^{-1}:x,y\in G\rangle
\]
的闭包。记
\[
G^{\mathrm{ab}}=G/G^{(1)}.
\]
于是
\[
G^{\mathrm{ab}}\simeq \Gal(F^{\mathrm{ab}}/F).
\]
任何酉特征标 $\chi$ 的核都包含 $G^{(1)}$，所以 $\chi$ 可看作 $G^{\mathrm{ab}}$ 的特征标。反过来，$G^{\mathrm{ab}}$ 的特征标经投射 $G\to G^{\mathrm{ab}}$ 拉回为 $G$ 的特征标。因此
\[
\widehat G=\widehat{G^{\mathrm{ab}}}.
\]
这正是类域论的入口：研究 $G$ 的一维特征标，等价于研究最大交换扩张的 Galois 群。
\end{proposition}

\begin{proof}
有限交换扩张的合成仍是交换扩张，因此它们的并给出最大交换扩张 $F^{\mathrm{ab}}$。无限 Galois 对应把它对应到闭正规子群 $G^{(1)}$。商群 $G/G^{(1)}$ 作用在 $F^{\mathrm{ab}}$ 上，因而同构于 $\Gal(F^{\mathrm{ab}}/F)$，这是交换群。

若 $N$ 是闭正规子群且 $G/N$ 交换，则每个交换子 $xyx^{-1}y^{-1}$ 都落在 $N$ 中；由于 $N$ 闭，交换子子群的闭包也落在 $N$ 中。因此交换子闭包是使商群交换的最小闭正规子群，也就是 $G^{(1)}$。

任取酉特征标 $\chi:G\to S^1$。目标群 $S^1$ 交换，所以 $\chi(xyx^{-1}y^{-1})=1$，从而 $[G,G]$ 及其闭包都包含在 $\ker\chi$ 中。因此 $\chi$ 唯一分解为 $G\to G^{\mathrm{ab}}\to S^1$。反向构造由商映射拉回特征标完成，两者互逆，得到 $\widehat G=\widehat{G^{\mathrm{ab}}}$。
\end{proof}

\section*{习题}
\addcontentsline{toc}{section}{习题}

\begin{exercise}
设 $I=(X^2+1)$、$J=(X^2+X+1)$ 为 $\Q[X]$ 的理想。令 $\sigma:\Q[X]\to\Q[X]/I$、$\tau:\Q[X]\to\Q[X]/J$ 为自然商映射，$\xi=\sigma(X)$，$\omega=\tau(X)$。计算 $p(\xi),q(\xi),\sigma(q(X))$ 以及 $p(\omega),q(\omega),\tau(p(X))$，其中 $p(X)=X^2+1$、$q(X)=X^2+X+1$。
\end{exercise}
\begin{proof}[解答]
状态：可解。在 $\Q[X]/I$ 中有 $\xi^2+1=0$，所以 $\xi^2=-1$。于是 $p(\xi)=0$，且 $q(\xi)=\xi^2+\xi+1=-1+\xi+1=\xi$，因此 $\sigma(q(X))=\xi$。

在 $\Q[X]/J$ 中有 $\omega^2+\omega+1=0$，所以 $\omega^2=-\omega-1$。于是 $q(\omega)=0$，且 $p(\omega)=\omega^2+1=-\omega$，因此 $\tau(p(X))=-\omega$。商环计算的核心是：把定义理想中的多项式关系当作降阶规则。
\end{proof}

\begin{exercise}
令 $\alpha=\sqrt[3]{2}$ 为实根，$\zeta=(-1+\sqrt{-3})/2$ 为本原三次单位根，$K$ 为 $X^3-2$ 在 $\Q$ 上的分裂域。计算 $[\Q(\alpha,\zeta):\Q]$、$[\Q(\zeta\alpha):\Q]$，写出 $K$，并判断 $K/\Q$、$K/\Q(\zeta\alpha)$、$\Q(\zeta\alpha)/\Q$ 是否为 Galois 扩张。
\end{exercise}
\begin{proof}[解答]
状态：可解。多项式 $X^3-2$ 由 Eisenstein 判别法在 $\Q[X]$ 中不可约，所以 $[\Q(\alpha):\Q]=3$。三个根为 $\alpha,\zeta\alpha,\zeta^2\alpha$，故分裂域是 $K=\Q(\alpha,\zeta)$。因为 $\Q(\alpha)\subset\R$，而 $\zeta\notin\R$，所以 $\zeta\notin\Q(\alpha)$，从而 $[\Q(\alpha,\zeta):\Q]=2\cdot3=6$。

元素 $\zeta\alpha$ 也是 $X^3-2$ 的根，故其极小多项式仍为 $X^3-2$，于是 $[\Q(\zeta\alpha):\Q]=3$。扩张 $K/\Q$ 是特征 $0$ 中一个多项式的分裂域，所以是 Galois 扩张。扩张 $K/\Q(\zeta\alpha)$ 是二次扩张，因为 $[K:\Q(\zeta\alpha)]=6/3=2$，特征 $0$ 的二次扩张正规可分，所以是 Galois 扩张。扩张 $\Q(\zeta\alpha)/\Q$ 不是 Galois 扩张，因为它不包含 $X^3-2$ 的所有根，例如不包含非实比值所需的 $\zeta$。
\end{proof}

\begin{exercise}
令 $\xi=1+\sqrt3$。求 $\xi$ 在 $\Q$ 上的极小多项式，并给出 $\Q(\xi)$ 的一组 $\Q$-基。
\end{exercise}
\begin{proof}[解答]
状态：可解。由 $(\xi-1)^2=3$ 得 $\xi^2-2\xi-2=0$。多项式 $X^2-2X-2$ 的判别式为 $12$，不是有理数平方，因此它在 $\Q[X]$ 中不可约。故极小多项式为 $X^2-2X-2$。又 $\sqrt3=\xi-1$，所以 $\Q(\xi)=\Q(\sqrt3)$，一组基可取 $\{1,\xi\}$，也可取 $\{1,\sqrt3\}$。
\end{proof}

\begin{exercise}
设 $F=\{p+qi+r\sqrt5+si\sqrt5:p,q,r,s\in\Q\}\subset\C$，其中 $i^2=-1$。证明表示唯一；对非零元素用共轭乘积构造逆元；证明 $F$ 是一个域。
\end{exercise}
\begin{proof}[解答]
状态：可解。若 $p+qi+r\sqrt5+si\sqrt5=0$，取实部和虚部分别得 $p+r\sqrt5=0$、$q+s\sqrt5=0$。由于 $\sqrt5\notin\Q$，得到 $p=r=0$ 与 $q=s=0$。

令 $\alpha=p+qi+r\sqrt5+si\sqrt5$。把 $\sqrt5$ 和 $i$ 分别取共轭，可得四个共轭元。题中给出的 $\beta$ 正是除 $\alpha$ 外其余三个共轭元的乘积，所以 $c=\alpha\beta$ 是四个共轭元的乘积，落在 $\Q$ 中。若 $\alpha\ne0$，四个共轭元都非零，故 $c\ne0$。于是 $\alpha^{-1}=\beta/c\in F$。集合 $F$ 对加减乘封闭，又含有每个非零元素的逆元，所以是域；事实上 $F=\Q(i,\sqrt5)$。
\end{proof}

\begin{exercise}
证明 $\Q(i,\sqrt5)=\Q(i+\sqrt5)$。
\end{exercise}
\begin{proof}[解答]
状态：可解。令 $\theta=i+\sqrt5$。则 $\theta^2=4+2i\sqrt5$，所以 $i\sqrt5=(\theta^2-4)/2\in\Q(\theta)$。又
\[
\theta(i\sqrt5)=(i+\sqrt5)i\sqrt5=5i-\sqrt5.
\]
把 $\theta=i+\sqrt5$ 与 $\theta(i\sqrt5)=5i-\sqrt5$ 相加，得到 $6i\in\Q(\theta)$，故 $i\in\Q(\theta)$。于是 $\sqrt5=\theta-i\in\Q(\theta)$。这证明 $\Q(i,\sqrt5)\subset\Q(\theta)$，反向包含由 $\theta=i+\sqrt5\in\Q(i,\sqrt5)$ 得到。
\end{proof}

\begin{exercise}
证明 $\sqrt3\notin\Q(\sqrt2)$、$\sqrt5\notin\Q(\sqrt2,\sqrt3)$，并计算 $[\Q(\sqrt2,\sqrt3,\sqrt5):\Q]$。
\end{exercise}
\begin{proof}[解答]
状态：可解。若 $\sqrt3=a+b\sqrt2$，平方并比较 $\sqrt2$ 项得 $2ab=0$。若 $b=0$ 则 $a^2=3$，若 $a=0$ 则 $2b^2=3$，都不可能有有理解。因此 $\sqrt3\notin\Q(\sqrt2)$。

域 $\Q(\sqrt2,\sqrt3)$ 是四次 Galois 扩张，二次子域正好是 $\Q(\sqrt2)$、$\Q(\sqrt3)$、$\Q(\sqrt6)$。若 $\sqrt5$ 属于其中，则 $\Q(\sqrt5)$ 会是这三个二次子域之一，这等价于 $5$ 与 $2,3,6$ 中某个数相差一个有理数平方，矛盾。因此 $\sqrt5\notin\Q(\sqrt2,\sqrt3)$。于是三个平方根逐个使次数乘以 $2$，得到总次数 $2^3=8$。
\end{proof}

\begin{exercise}
设 $f(z)=4z^3-c_2z-c_3$，根为 $e_1,e_2,e_3$。证明 Vieta 关系，并推出判别式
\[
D(f)=(e_1-e_2)^2(e_2-e_3)^2(e_3-e_1)^2=\frac1{16}(c_2^3-27c_3^2).
\]
\end{exercise}
\begin{proof}[解答]
状态：可解。把 $f(z)/4$ 写成 $z^3-(c_2/4)z-(c_3/4)$。Vieta 公式给出
\[
e_1+e_2+e_3=0,\quad e_1e_2+e_2e_3+e_3e_1=-c_2/4,\quad e_1e_2e_3=c_3/4.
\]
由 $e_1+e_2+e_3=0$ 得 $e_2+e_3=-e_1$。于是
\[
e_2^2-e_1e_2=e_2(e_2+e_3)+e_2(e_1-e_3)-e_1e_2
\]
化简时不断用上面的三个对称关系，可把所有二次、三次对称式改写为 $c_2,c_3$。更直接的判别式公式为：三次多项式 $az^3+bz^2+cz+d$ 的判别式是
\[
b^2c^2-4ac^3-4b^3d-27a^2d^2+18abcd.
\]
在这里 $a=4,b=0,c=-c_2,d=-c_3$，所以
\[
\operatorname{disc}(f)=16c_2^3-432c_3^2=16(c_2^3-27c_3^2).
\]
另一方面，判别式也等于
\[
a^{2n-2}\prod_{i<j}(e_i-e_j)^2=4^4D(f)=256D(f).
\]
因此 $D(f)=\operatorname{disc}(f)/256=(c_2^3-27c_3^2)/16$。
\end{proof}

\begin{exercise}
设 $f(z)=4z^3-c_2z-c_3\in\Q[z]$，三个根均不在 $\Q$，且判别式非零。讨论分裂域的 Galois 群；特别计算 $4(z^3-z+1)$ 和 $z^3-3z+1$ 的 Galois 群。
\end{exercise}
\begin{proof}[解答]
状态：可解。若三个根都不在 $\Q$，则 $f/4$ 无有理根。三次多项式可约当且仅当有一次因子，所以 $f$ 不可约，$[\Q(e_1):\Q]=3$。分裂域 $K$ 的自同构忠实作用在三个不同根上，因此 $\Gal(K/\Q)$ 嵌入 $S_3$。

三次不可约多项式的 Galois 群是 $A_3$ 或 $S_3$；判别式在 $\Q$ 中为平方时落在 $A_3$，否则为 $S_3$。对 $4(z^3-z+1)$，有 $c_2=4,c_3=-4$，故 $D=(64-27\cdot16)/16=-23$，不是平方，所以 Galois 群为 $S_3$，其唯一二次子域为 $\Q(\sqrt{-23})$。对 $z^3-3z+1$，判别式为 $-4(-3)^3-27=81$，是平方；该多项式无有理根，所以 Galois 群为 $A_3\simeq Z/3Z$。
\end{proof}

\begin{exercise}
用 Cardano 代换处理 $g(X)=X^3+3HX+G$，并说明三次方程的根可由系数的根式表示；再求 $X^3-6X-9$ 的根。
\end{exercise}
\begin{proof}[解答]
状态：可解。令 $u^3=\xi=(-G+\sqrt{G^2+4H^3})/2$，并取 $v=-H/u$。则 $uv=-H$，且 $v^3=-H^3/u^3$。由于 $\xi$ 满足 $T^2+GT-H^3=0$，有 $\xi-H^3/\xi=-G$，所以 $u^3+v^3=-G$。

对任意三次单位根 $\omega$，三个数 $u+v$、$u\omega+v\omega^2$、$u\omega^2+v\omega$ 都满足
\[
X^3-3uvX-(u^3+v^3)=X^3+3HX+G=0.
\]
因此根落在 $F(G,H)(\omega,\sqrt{G^2+4H^3},\sqrt[3]{(-G+\sqrt{G^2+4H^3})/2})$ 中。

对 $X^3-6X-9$，有 $3H=-6$，即 $H=-2$，且 $G=-9$。判别式项为 $G^2+4H^3=81-32=49$。于是 $u^3=(9+7)/2=8$，可取 $u=2$，$v=-H/u=1$。三个根为 $3$、$2\omega+\omega^2$、$2\omega^2+\omega$。
\end{proof}

\begin{exercise}
在 $R=\Z[\sqrt{-3}]$ 中证明单位群为 $\{\pm1\}$，证明范数为 $4$ 的元素不可约，并用 $2\cdot2=(1+\sqrt{-3})(1-\sqrt{-3})$ 说明 $R$ 不是 UFD。
\end{exercise}
\begin{proof}[解答]
状态：可解。范数为 $N(a+b\sqrt{-3})=a^2+3b^2$。单位必须满足范数为 $1$，故 $a^2+3b^2=1$，只有 $\pm1$。

若 $N(x)=4$ 且 $x=yz$，则 $N(y)N(z)=4$。非单位的范数至少为 $2$。但方程 $a^2+3b^2=2$ 无整数解，所以不存在范数为 $2$ 的元素。因此一个非平凡分解无法存在，$x$ 不可约。

有 $N(2)=4$，$N(1\pm\sqrt{-3})=4$，三者都不可约。等式
\[
2\cdot2=(1+\sqrt{-3})(1-\sqrt{-3})
\]
给出同一元素的两个不可约分解。由于单位只有 $\pm1$，$2$ 与 $1+\sqrt{-3}$ 不相差单位，故唯一分解失败。
\end{proof}

\begin{exercise}
设 $d$ 为平方自由整数，$\Q(\sqrt d)$ 中元素 $r+s\sqrt d$ 满足 $2r\in\Z$、$r^2-s^2d\in\Z$。证明：若 $d\not\equiv1\pmod4$，则 $r,s\in\Z$；若 $d\equiv1\pmod4$，则 $2s\in\Z$ 且 $2r\equiv2s\pmod2$。
\end{exercise}
\begin{proof}[解答]
状态：可解。令 $m=2r\in\Z$。由 $r^2-s^2d\in\Z$ 得 $m^2-4s^2d\in4\Z$，所以 $4s^2d\in\Z$。因为 $d$ 平方自由，分母比较给出 $2s\in\Z$。令 $n=2s\in\Z$，则条件变为 $(m^2-n^2d)/4\in\Z$，即 $m^2\equiv n^2d\pmod4$。

平方模 $4$ 只可能为 $0$ 或 $1$。若 $d\equiv2,3\pmod4$，同余迫使 $m,n$ 都为偶数，因此 $r,s\in\Z$。若 $d\equiv1\pmod4$，同余变为 $m^2\equiv n^2\pmod4$，等价于 $m,n$ 奇偶相同，即 $2r\equiv2s\pmod2$。
\end{proof}

\begin{exercise}
设 $E=\Q(\alpha)$，其中 $\alpha^3+\alpha^2+\alpha+2=0$。把 $(\alpha^2+\alpha+1)(\alpha^2+\alpha)$ 与 $(\alpha-1)^{-1}$ 写成 $a\alpha^2+b\alpha+c$。
\end{exercise}
\begin{proof}[解答]
状态：可解。关系为 $\alpha^3=-\alpha^2-\alpha-2$，从而 $\alpha^4=2-\alpha$。于是
\[
(\alpha^2+\alpha+1)(\alpha^2+\alpha)=\alpha^4+2\alpha^3+2\alpha^2+\alpha=-2\alpha-2.
\]
设 $(\alpha-1)^{-1}=a\alpha^2+b\alpha+c$。比较
\[
(\alpha-1)(a\alpha^2+b\alpha+c)
=(b-2a)\alpha^2+(c-b-a)\alpha+(-2a-c)
\]
与 $1$，得到 $b=2a$、$c=3a$、$-5a=1$。故
\[
(\alpha-1)^{-1}=-\frac15\alpha^2-\frac25\alpha-\frac35.
\]
\end{proof}

\begin{exercise}
证明域的非零元乘法群中的任意有限子群都是循环群。
\end{exercise}
\begin{proof}[解答]
状态：可解。设 $H\subset F^\times$ 有限。它是有限交换群。令 $m$ 为 $H$ 中元素阶的最小公倍数，则每个 $h\in H$ 都满足 $h^m=1$，所以 $H$ 全部落在多项式 $X^m-1$ 的根集中。域上 $m$ 次多项式至多有 $m$ 个根，因此 $|H|\le m$。又 $m$ 是元素阶的最小公倍数，必整除有限交换群 $H$ 的阶，所以 $m\le |H|$。于是 $m=|H|$。有限交换群的指数等于阶时，存在一个元素阶为 $|H|$，所以 $H$ 循环。
\end{proof}

\begin{exercise}
证明 $1,\sqrt2,\sqrt3,\sqrt6$ 在 $\Q$ 上线性无关。
\end{exercise}
\begin{proof}[解答]
状态：可解。若 $a+b\sqrt2+c\sqrt3+d\sqrt6=0$，把式子写成 $(a+b\sqrt2)=-(c+d\sqrt2)\sqrt3$。两边平方可得 $\sqrt2$ 的有理线性表达。更直接地，使用 $\Q(\sqrt2,\sqrt3)$ 的四个自同构，分别令 $\sqrt2,\sqrt3$ 独立变号。对原等式施加四个符号选择并相加，得到 $4a=0$。乘以相应符号后相加，分别得到 $4b\sqrt2=0$、$4c\sqrt3=0$、$4d\sqrt6=0$。因此四个系数全为 $0$。
\end{proof}

\begin{exercise}
研究 $\Q(\sqrt2,\sqrt3)$：求 $\sqrt2+\sqrt3$ 的方程、次数、分裂域性质，并证明它是本原元素。
\end{exercise}
\begin{proof}[解答]
状态：可解。令 $\theta=\sqrt2+\sqrt3$。则 $\theta^2=5+2\sqrt6$，所以 $(\theta^2-5)^2=24$，即 $\theta^4-10\theta^2+1=0$。直接代入可检验这个四次式成立。若题面出现 $X^4-15X^2+38$，代入 $\theta^2=5+2\sqrt6$ 后不为零，计算时应使用由平方消元得到的关系。

由 $\theta(\sqrt3-\sqrt2)=1$ 得 $\sqrt3-\sqrt2=\theta^{-1}\in\Q(\theta)$。因此
\[
\sqrt3=\frac{\theta+\theta^{-1}}2,\qquad \sqrt2=\frac{\theta-\theta^{-1}}2,
\]
所以 $\Q(\theta)=\Q(\sqrt2,\sqrt3)$。由前题，扩张次数为 $4$。该域包含 $\pm\sqrt2,\pm\sqrt3$，所以是 $(X^2-2)(X^2-3)$ 的分裂域。
\end{proof}

\begin{exercise}
令 $\xi=1+\sqrt3$。求极小多项式，证明它不是 $\Q(\sqrt3)$ 中的平方数，并求 $\Q(\xi)$ 的基和次数。
\end{exercise}
\begin{proof}[解答]
状态：可解。极小多项式和基已在习题 3 中得到：$m_\xi(X)=X^2-2X-2$，且 $\Q(\xi)=\Q(\sqrt3)$，基为 $\{1,\sqrt3\}$，次数为 $2$。

若 $1+\sqrt3=(a+b\sqrt3)^2$，比较系数得 $a^2+3b^2=1$、$2ab=1$。由第二式 $b=1/(2a)$。代入第一式并乘以 $4a^2$ 得 $4a^4-4a^2+3=0$。令 $u=a^2$，方程 $4u^2-4u+3=0$ 的判别式为 $-32$，无有理解。因此 $\xi$ 不是平方数。
\end{proof}

\begin{exercise}
若 $\alpha$ 是 $X^3+2X+2$ 的根，$\beta$ 是 $X^3-6X+3$ 的根，证明 $\alpha+\beta$ 满足一个次数至多 $9$ 的 $\Q$ 系数多项式。
\end{exercise}
\begin{proof}[解答]
状态：可解。扩张 $\Q(\alpha)$ 的次数至多 $3$，在其上再加入 $\beta$ 后，$\Q(\alpha,\beta)$ 在 $\Q$ 上次数至多 $9$。元素 $\alpha+\beta$ 属于这个有限维 $\Q$-向量空间。于是 $1,\theta,\theta^2,\dots,\theta^9$（其中 $\theta=\alpha+\beta$）共 $10$ 个元素在维数至多 $9$ 的空间中线性相关。故存在不全为零的 $c_i\in\Q$ 使 $\sum_{i=0}^9c_i\theta^i=0$，这就是所求多项式。
\end{proof}

\begin{exercise}
设 $t$ 是域扩张 $E/F$ 中元素，$f\in F[X]$ 是使 $f(t)=0$ 的最小次数多项式。证明 $f$ 在 $F[X]$ 中不可约；讨论它在 $E[X]$ 中是否不可约；再证明任何不可约且以 $t$ 为根的 $F[X]$ 多项式都与极小多项式相差一个单位因子。
\end{exercise}
\begin{proof}[解答]
状态：可解。若 $f=gh$ 在 $F[X]$ 中非平凡分解，则 $0=f(t)=g(t)h(t)$。域中无零因子，所以 $g(t)=0$ 或 $h(t)=0$。这给出次数更小的消去多项式，矛盾。因此 $f$ 不可约。

在 $E[X]$ 中，$t\in E$，所以 $X-t$ 是 $f$ 的因子。若 $\deg f>1$，则 $f$ 在 $E[X]$ 中可约。最后，设 $m$ 为首一极小多项式，$h\in F[X]$ 不可约且 $h(t)=0$。用除法写 $h=qm+r$，$\deg r<\deg m$。代入 $t$ 得 $r(t)=0$，由最小性 $r=0$，所以 $m\mid h$。因 $h$ 不可约，$h$ 与 $m$ 相差一个非零常数因子。
\end{proof}

\begin{exercise}
证明 Eisenstein 判别法。
\end{exercise}
\begin{proof}[解答]
状态：可解。由 Gauss 引理，只需在 $\Z[X]$ 中排除本原多项式的非平凡分解。设 $f=gh$，且 $g,h$ 都为正次数本原多项式。模 $p$ 后，条件给出 $\overline f=\overline a_nX^n$。在 $\F_p[X]$ 中，$\overline g\overline h$ 是单项式，所以 $\overline g$ 与 $\overline h$ 都是单项式。于是 $g(0)$ 和 $h(0)$ 都被 $p$ 整除。故 $a_0=f(0)=g(0)h(0)$ 被 $p^2$ 整除，违反假设。故 $f$ 在 $\Q[X]$ 中不可约。
\end{proof}

\begin{exercise}
用 Eisenstein 判别法处理 $X^7-2$，证明 $\alpha=\sqrt[7]{2}$ 的幂给出基，并求 $(\alpha^2+1)^{-1}$。
\end{exercise}
\begin{proof}[解答]
状态：可解。对 $X^7-2$ 取 $p=2$，Eisenstein 条件成立，所以它不可约，是 $\alpha$ 的极小多项式。因此 $\{1,\alpha,\dots,\alpha^6\}$ 线性无关并构成 $\Q(\alpha)$ 的基；$\Q(\alpha)=\Q[\alpha]$。任意高次幂都可用关系 $\alpha^7=2$ 降到六次以内，任意非零元素的逆元可由它与极小多项式的扩展 Euclid 等式得到，所以这个张成空间对乘法和逆元封闭。

设 $(\alpha^2+1)^{-1}=A(\alpha)$，其中
\[
A(X)=\frac15(1+2X-X^2-2X^3+X^4+2X^5-X^6).
\]
直接在商环 $\Q[X]/(X^7-2)$ 中计算得 $(X^2+1)A(X)\equiv1$。代入 $X=\alpha$ 得
\[
(\alpha^2+1)^{-1}=\frac15(1+2\alpha-\alpha^2-2\alpha^3+\alpha^4+2\alpha^5-\alpha^6).
\]
\end{proof}

\begin{exercise}
设 $a,b,c$ 是 $X^3+9X^2+3X$ 的根。用基本对称式求 $a^2,b^2,c^2$ 的多项式。
\end{exercise}
\begin{proof}[解答]
状态：可解。Vieta 公式给出 $s_1=a+b+c=-9$、$s_2=ab+ac+bc=3$、$s_3=abc=0$。于是
\[
a^2+b^2+c^2=s_1^2-2s_2=75,
\]
\[
a^2b^2+a^2c^2+b^2c^2=s_2^2-2s_1s_3=9,\qquad a^2b^2c^2=s_3^2=0.
\]
所以 $a^2,b^2,c^2$ 是
\[
Y^3-75Y^2+9Y=0
\]
的根。
\end{proof}

\begin{exercise}
证明 $\sqrt3\notin\Q(\sqrt2)$、$\sqrt5\notin\Q(\sqrt2,\sqrt3)$，并求 $\Q(\sqrt2,\sqrt3,\sqrt5)$ 的次数。
\end{exercise}
\begin{proof}[解答]
状态：可解。此题与习题 6 相同。结论为 $[\Q(\sqrt2,\sqrt3,\sqrt5):\Q]=8$。证明使用多重二次扩张的二次子域分类：$\Q(\sqrt2,\sqrt3)$ 的二次子域只有 $\Q(\sqrt2)$、$\Q(\sqrt3)$、$\Q(\sqrt6)$，不包含 $\Q(\sqrt5)$。
\end{proof}

\begin{exercise}
证明 $(X^2-3)(X^3+1)$ 与 $(X^2-2X-2)(X^2+1)$ 的分裂域均为 $\Q(i,\sqrt3)$；并证明 $X^2-3$ 与 $X^2-2X-2$ 的分裂域均为 $\Q(\sqrt3)$。
\end{exercise}
\begin{proof}[解答]
状态：可解。$X^2-3$ 的根为 $\pm\sqrt3$。$X^3+1=(X+1)(X^2-X+1)$，后二次因子的根为 $(1\pm\sqrt{-3})/2$，它们落在 $\Q(i,\sqrt3)$ 中。因此第一组乘积的分裂域为 $\Q(i,\sqrt3)$。

$X^2-2X-2$ 的根为 $1\pm\sqrt3$，而 $X^2+1$ 的根为 $\pm i$，所以第二组乘积的分裂域同样为 $\Q(i,\sqrt3)$。单独看二次多项式，$X^2-3$ 和 $X^2-2X-2$ 的分裂域分别为 $\Q(\sqrt3)$ 与 $\Q(\sqrt3)$。
\end{proof}

\begin{exercise}
对域自同构群和子域定义 $I(H)$、$A(K)$。证明 $I(A(K))\supset K$、$A(I(H))\supset H$、$A(I(A(K)))=A(K)$、$I(A(I(H)))=I(H)$。
\end{exercise}
\begin{proof}[解答]
状态：可解。若 $x\in K$，则每个固定 $K$ 的自同构固定 $x$，故 $x\in I(A(K))$。若 $\sigma\in H$，则 $\sigma$ 固定所有被 $H$ 固定的元素，故 $\sigma\in A(I(H))$。

由 $I(A(K))\supset K$ 得 $A(I(A(K)))\subset A(K)$。反过来，若 $\sigma\in A(K)$，则 $\sigma$ 属于 $A(K)$ 这个群，因而固定 $I(A(K))$ 中的每个元素，所以 $\sigma\in A(I(A(K)))$。第二个等式同法：由 $A(I(H))\supset H$ 得 $I(A(I(H)))\subset I(H)$；而 $I(H)$ 中的元素按定义被 $A(I(H))$ 全部固定，得反向包含。
\end{proof}

\begin{exercise}
对 $X^3-2$ 的分裂域，求 $[\Q(\sqrt[3]{2},\zeta):\Q]$、所有子域。
\end{exercise}
\begin{proof}[解答]
状态：可解。记 $\alpha=\sqrt[3]{2}$、$\zeta$ 为本原三次单位根，$E=\Q(\alpha,\zeta)$。前面已得 $[E:\Q]=6$ 且 $\Gal(E/\Q)\simeq S_3$。子域由 $S_3$ 的子群给出。全群对应 $\Q$，平凡子群对应 $E$。唯一三阶子群 $A_3$ 对应二次子域 $\Q(\zeta)=\Q(\sqrt{-3})$。三个二阶子群分别对应三个三次子域 $\Q(\alpha)$、$\Q(\zeta\alpha)$、$\Q(\zeta^2\alpha)$。
\end{proof}

\begin{exercise}
证明 $X^6+X^3+1$ 整除 $X^9-1$，并求 $Q(\alpha)\to\C$ 的全部嵌入，其中 $\alpha$ 是 $X^6+X^3+1$ 的根。
\end{exercise}
\begin{proof}[解答]
状态：可解。有 $X^9-1=(X^3-1)(X^6+X^3+1)$。若 $\alpha$ 是后一个多项式的根，则 $\alpha^9=1$ 且 $\alpha^3\ne1$，所以 $\alpha$ 是本原 $9$ 次单位根。于是 $\Q(\alpha)=\Q(\zeta_9)$。嵌入由 $\zeta_9\mapsto\zeta_9^a$ 给出，其中 $a\in(\Z/9\Z)^\times=\{1,2,4,5,7,8\}$。
\end{proof}

\begin{exercise}
令 $\alpha^4=5$ 且 $\alpha>0$。证明 $\Q(i\alpha^2)/\Q$ 正规，$\Q(\alpha+i\alpha)/\Q(i\alpha^2)$ 正规，但 $\Q(\alpha+i\alpha)/\Q$ 不正规。
\end{exercise}
\begin{proof}[解答]
状态：可解。因为 $i\alpha^2=i\sqrt5=\sqrt{-5}$，所以 $\Q(i\alpha^2)=\Q(\sqrt{-5})$ 是二次扩张，故在 $\Q$ 上正规。

令 $\gamma=\alpha+i\alpha=(1+i)\alpha$。则 $\gamma^2=2i\alpha^2\in\Q(i\alpha^2)$，所以 $\gamma$ 在该基域上满足二次多项式 $X^2-2i\alpha^2$；二次可分扩张正规，故 $\Q(\gamma)/\Q(i\alpha^2)$ 正规。

在 $\Q$ 上，$\gamma^4=(1+i)^4\alpha^4=-20$，其共轭根包括 $i\gamma$。若 $\Q(\gamma)/\Q$ 正规，则应包含 $i\gamma$，进而包含 $i$。但 $\Q(\gamma)$ 有一个实结构限制，直接由 $[\Q(\gamma):\Q]=4$ 和其中二次子域为 $\Q(\sqrt{-5})$ 可见 $i\notin\Q(\gamma)$。因此不正规。
\end{proof}

\begin{exercise}
求 $X^4-4X^2-5$ 的分裂域和 Galois 群，并列出真子群和对应不变子域。
\end{exercise}
\begin{proof}[解答]
状态：可解。分解为 $(X^2-5)(X^2+1)$，根为 $\pm\sqrt5,\pm i$，分裂域为 $F=\Q(i,\sqrt5)$。两个二次子域交为 $\Q$，故 $[F:\Q]=4$。自同构由 $i\mapsto\pm i$、$\sqrt5\mapsto\pm\sqrt5$ 独立决定，所以 Galois 群为 $\{I,R,S,T\}\simeq Z/2Z\times Z/2Z$，其中 $R$ 改变 $i$ 符号，$S$ 改变 $\sqrt5$ 符号，$T=RS$。真子群为 $\{I\}$、$\{I,R\}$、$\{I,S\}$、$\{I,T\}$，对应不变域为 $F$、$\Q(\sqrt5)$、$\Q(i)$、$\Q(i\sqrt5)$。
\end{proof}

\begin{exercise}
求 $X^4+1$ 的 Galois 群。
\end{exercise}
\begin{proof}[解答]
状态：可解。根为 $\zeta_8,\zeta_8^3,\zeta_8^5,\zeta_8^7$，分裂域为 $\Q(\zeta_8)=\Q(i,\sqrt2)$。因此
\[
\Gal(\Q(\zeta_8)/\Q)\simeq(\Z/8\Z)^\times\simeq Z/2Z\times Z/2Z,
\]
即 Klein 四元群。
\end{proof}

\begin{exercise}
求 $X^3-X-1$ 在 $\Q$ 和 $\Q(\sqrt{23})$ 上的 Galois 群。
\end{exercise}
\begin{proof}[解答]
状态：可解。多项式无有理根，故在 $\Q$ 上不可约。三次式 $X^3+pX+q$ 的判别式为 $-4p^3-27q^2$；这里 $p=-1,q=-1$，得判别式 $-23$。它不是 $\Q$ 中平方，所以 $\Gal/\Q\simeq S_3$。

在 $\Q(\sqrt{23})$ 上，$-23$ 仍不是平方，因为若 $-23=(a+b\sqrt{23})^2$，比较 $\sqrt{23}$ 系数得 $2ab=0$，再比较有理部分无法成立。因此 Galois 群仍为 $S_3$。若把基域换成判别式域 $\Q(\sqrt{-23})$，则 Galois 群降为 $A_3$。
\end{proof}

\begin{exercise}
设 $p$ 为素数，$F$ 特征 $0$ 且含本原 $p$ 次单位根。若 $E/F$ 是 $p$ 次循环扩张，证明存在 $\xi\in E$ 使 $E=F(\xi)$ 且 $\xi^p\in F$。
\end{exercise}
\begin{proof}[解答]
状态：可解。令 $\Gal(E/F)=\langle\sigma\rangle$，取本原 $p$ 次单位根 $\zeta\in F$。对 $a\in E$ 定义
\[
\xi_a=\sum_{j=0}^{p-1}\zeta^{-j}\sigma^j(a).
\]
并非所有 $\xi_a$ 都为 $0$，否则线性算子 $\sum\zeta^{-j}\sigma^j$ 为零，这会违反 Dedekind 线性无关性。取 $\xi=\xi_a\ne0$。则
\[
\sigma(\xi)=\sum_{j=0}^{p-1}\zeta^{-j}\sigma^{j+1}(a)=\zeta\xi.
\]
所以 $\sigma(\xi^p)=\xi^p$，即 $\xi^p\in F$。又 $\sigma(\xi)\ne\xi$，故 $\xi\notin F$。因为 $[E:F]=p$ 为素数，$F(\xi)$ 是严格中间域，只能等于 $E$。
\end{proof}

\begin{exercise}
令 $\alpha=\sqrt[3]{2}$、$\zeta=(-1+\sqrt{-3})/2$、$\beta=\zeta\alpha$、$F=\Q(\alpha)$、$K=\Q(\beta)$。解答关于 $F,K,FK$、自同构和分裂域 Galois 群的问题。
\end{exercise}
\begin{proof}[解答]
状态：可解。因为 $F\subset\R$ 而 $K$ 含非实元素 $\beta$，故 $K\ne F$。交 $K\cap F$ 不是三次域，否则会等于 $F=K$，所以 $K\cap F=\Q$。合成域
\[
FK=\Q(\alpha,\beta)=\Q(\alpha,\zeta)=\Q(\alpha,\sqrt{-3}),
\]
因为 $\beta/\alpha=\zeta$。因此 $[FK:F]=2$，$[K:\Q]=3$。$K/\Q$ 不是 Galois 扩张，因为不含 $X^3-2$ 的全部根。$\alpha$ 的极小多项式为 $X^3-2$，且 $\{1,\alpha,\alpha^2\}$ 是 $F/\Q$ 的基。

若 $\theta$ 是 $F$ 的 $\Q$-自同构，则 $\theta(\alpha)$ 必须是 $X^3-2$ 的根。域 $F$ 是实域，只含实根 $\alpha$，故 $\theta(\alpha)=\alpha$，自同构群平凡。令 $E=\Q(\alpha,\zeta)$。自同构 $\sigma(\alpha)=\zeta\alpha,\sigma(\zeta)=\zeta$ 与 $\tau(\alpha)=\alpha,\tau(\zeta)=\zeta^2$ 生成 $6$ 个元素 $I,\sigma,\sigma^2,\tau,\sigma\tau,\sigma^2\tau$，满足 $S_3$ 的关系。子群 $H=\{I,\lambda\}$ 的固定域按直接代入检验为 $\Q(\zeta^2\alpha)$。
\end{proof}

\begin{exercise}
若 $p_1,\dots,p_r$ 是互异素数，证明 $\Q(\sqrt{p_1},\dots,\sqrt{p_r})/\Q$ 的 Galois 群为 $(Z/2Z)^r$。
\end{exercise}
\begin{proof}[解答]
状态：可解。每个自同构必须把 $\sqrt{p_i}$ 送到 $\pm\sqrt{p_i}$。这些符号选择可以独立实现，因为平方类 $p_1,\dots,p_r$ 在 $\Q^\times/(\Q^\times)^2$ 中独立。于是扩张次数为 $2^r$，自同构也有 $2^r$ 个，Galois 群为 $(Z/2Z)^r$。
\end{proof}

\begin{exercise}
给定素数 $p$ 和正整数 $r$，构造 $F/\Q$ 使 $\Gal(F/\Q)\simeq Z/p^rZ$。
\end{exercise}
\begin{proof}[解答]
状态：可解，使用 Dirichlet 算术级数定理。取素数 $\ell\equiv1\pmod{p^r}$。分圆扩张 $\Q(\zeta_\ell)/\Q$ 的 Galois 群为 $(\Z/\ell\Z)^\times$，这是阶 $\ell-1$ 的循环群。循环群存在商群 $Z/p^rZ$；令 $F$ 为对应子群的不变域，则 $\Gal(F/\Q)\simeq Z/p^rZ$。
\end{proof}

\begin{exercise}
设 $u,v$ 为独立变量，$\zeta$ 为本原 $n$ 次单位根，$F=\C(u^n+v^n,uv)$。证明 $\C(u,v)/F$ 为 Galois 扩张，且 Galois 群由 $\sigma(u,v)=(v,u)$、$\tau(u,v)=(\zeta u,\zeta^{-1}v)$ 生成，并同构于二面体群 $D_n$。
\end{exercise}
\begin{proof}[解答]
状态：可解。两个变换都固定 $u^n+v^n$ 和 $uv$，故给出 $F$-自同构。它们满足 $\tau^n=1$、$\sigma^2=1$、$\sigma\tau\sigma=\tau^{-1}$，因此生成一个二面体群商。

反向看次数。令 $S=u^n+v^n$、$P=uv$。元素 $u^n$ 与 $v^n$ 是方程 $Z^2-SZ+P^n=0$ 的两个根，所以在 $F$ 上至多二次；在加入 $u^n$ 后，$u$ 至多给出 $n$ 次扩张，且 $v=P/u$。故 $[\C(u,v):F]\le2n$。上述自同构已有 $2n$ 个，因此扩张是 Galois，群为 $D_n$。
\end{proof}

\begin{exercise}
设 $F$ 特征 $0$，$L/F$ 代数，$f\in F[X]$ 的不可约因子均无重根。若 $K,E$ 分别为 $f$ 在 $F,L$ 上的分裂域，证明 $\Gal(E/L)$ 同构于 $\Gal(K/F)$ 的子群。
\end{exercise}
\begin{proof}[解答]
状态：可解。在共同代数闭包中可取 $E=KL$。任意 $\sigma\in\Gal(E/L)$ 固定 $L$，并把 $f$ 的根送到 $f$ 的根，因此限制到 $K$ 给出 $F$-自同构。若 $\sigma|_K$ 为恒等，则 $\sigma$ 在 $KL=E$ 上为恒等，所以限制映射单射 $\Gal(E/L)\hookrightarrow\Gal(K/F)$。
\end{proof}

\begin{exercise}
设 $F$ 特征 $0$。若 $f\in F[X]$ 的 Galois 群可解，证明方程 $f=0$ 可由 $F$ 上根式解出。
\end{exercise}
\begin{proof}[解答]
状态：可解，使用 Galois 可解性判别。令 $K/F$ 为分裂域，$G=\Gal(K/F)$。若 $G$ 可解，则有正规列，其相邻商为素数阶循环群。先在 $F$ 中加入必要的单位根；这一步本身是根式扩张。对每个素数阶循环层，习题 31 说明该层可由添入一个 $p$ 次根得到。沿正规列逐层上升，得到一个由根式生成的扩张，其中包含 $K$，因此 $f$ 的所有根都可由根式表示。
\end{proof}

\begin{exercise}
用尺规作图解释正实数的加减乘除和平方根构造，并讨论能否作出 $\sqrt[3]{2}$。
\end{exercise}
\begin{proof}[解答]
状态：可解。平行线截比定理给出乘法和除法：题设两组平行线构造分别把相似三角形的比例写成 $t/a=b/1$ 与 $s/a=1/b$，所以 $t=ab$、$s=a/b$。加减由在线段上截取长度完成。

平方根构造来自半圆定理。以 $[0,1+a]$ 为直径作半圆，过 $1$ 作垂线交半圆于高 $r$。直角三角形的高定理给出 $r^2=1\cdot a$，所以 $r=\sqrt a$。

尺规可作数从 $\Q$ 出发只允许反复做四则运算和平方根，因此所得扩张次数是 $2$ 的幂。$\sqrt[3]{2}$ 在 $\Q$ 上次数为 $3$，不是 $2$ 的幂，所以不能尺规作出。
\end{proof}

\begin{exercise}
描述 $\F_4$ 与 $\F_8$，并证明 $\F_4$ 不是 $\F_8$ 的子域。
\end{exercise}
\begin{proof}[解答]
状态：可解。在 $\F_2$ 中，$X^4-X=X(X-1)(X^2+X+1)$。取 $X^2+X+1$ 的根 $\omega$，得 $\F_4=\{0,1,\omega,1+\omega\}$，次数为 $2$。

又
\[
X^8-X=X(X-1)(X^3+X+1)(X^3+X^2+1).
\]
取 $\alpha^3+\alpha+1=0$，则
\[
\F_8=\{a+b\alpha+c\alpha^2:a,b,c\in\F_2\},
\]
共 $8$ 个元素，次数为 $3$。有限域 $\F_{2^m}$ 的子域为 $\F_{2^d}$，其中 $d\mid m$。因为 $2\nmid3$，$\F_4$ 不是 $\F_8$ 的子域。
\end{proof}

\begin{exercise}
用特征标和 Jacobi 和计算 $N_p(X^n-a)$、$N_p(X^n+Y^n-1)$，并求 $N_p(X^2+Y^2-1)$。
\end{exercise}
\begin{proof}[解答]
状态：可解。设 $n\mid p-1$，$\chi$ 为 $\F_p^\times$ 的 $n$ 阶特征标。对 $a\in\F_p^\times$，
\[
N_p(X^n-a)=\sum_{i=0}^{n-1}\chi^i(a),
\]
因为右边在 $a$ 为 $n$ 次幂时等于 $n$，否则等于 $0$。于是
\[
N_p(X^n+Y^n-1)=\sum_{a+b=1}N_p(X^n-a)N_p(Y^n-b)
=\sum_{i,j=0}^{n-1}J(\chi^i,\chi^j).
\]
当 $n=2$ 时取二次特征标 $\chi$。逐项计算得
\[
N_p(X^2+Y^2-1)=p-\chi(-1).
\]
若 $p\equiv1\pmod4$，$\chi(-1)=1$，结果为 $p-1$；若 $p\equiv3\pmod4$，$\chi(-1)=-1$，结果为 $p+1$。
\end{proof}

\begin{exercise}
设 $g_t=\sum_n(\frac np)\zeta^{nt}$。证明 $g_t=(\frac tp)g_1$，并计算 $g_t^k$ 的有限 Fourier 系数。
\end{exercise}
\begin{proof}[解答]
状态：可解。若 $t\ne0$，令 $m=nt$，则
\[
g_t=\sum_m\left(\frac{mt^{-1}}p\right)\zeta^m=\left(\frac tp\right)g_1,
\]
因为二次特征标取值为 $\pm1$，其逆等于自身。$t=0$ 时非平凡特征标总和为 $0$。

函数 $t\mapsto g_t^k$ 以 $p$ 为周期。展开乘积：
\[
g_t^k=\sum_{m_1,\dots,m_k}\left(\frac{m_1\cdots m_k}{p}\right)\zeta^{t(m_1+\cdots+m_k)}.
\]
按 $n\equiv m_1+\cdots+m_k\pmod p$ 分组，Fourier 系数为
\[
c_k(n)=\sum_{m_1+\cdots+m_k\equiv n}\left(\frac{m_1\cdots m_k}{p}\right).
\]
\end{proof}

\begin{exercise}
设 $K/F$ 为可分正规代数扩张，给 $G=\Gal(K/F)$ 赋予以 $\Gal(K/E)$ 为基本邻域的拓扑。证明紧 Hausdorff 性、与逆极限拓扑一致、闭包形式的 Galois 对应，以及正规中间扩张给出商拓扑同构。
\end{exercise}
\begin{proof}[解答]
状态：可解。对每个有限 Galois 中间域 $E$，有限群 $\Gal(E/F)$ 离散。限制映射给出
\[
G\simeq\varprojlim_E\Gal(E/F).
\]
有限离散群紧 Hausdorff，直积紧 Hausdorff，闭子空间仍紧 Hausdorff，所以 $G$ 紧 Hausdorff。该逆极限拓扑的单位元基本邻域正是各限制映射的核，即 $\Gal(K/E)$，故与题设拓扑一致。

若 $H\le G$，固定域为 $E=K^H$。固定 $E$ 的自同构在每个有限层上的像，等于 $H$ 的像的闭包；有限层是离散的，闭包就是同一个有限像。因此 $\Gal(K/E)=\overline H$。若 $E/F$ 可分正规，$H=\Gal(K/E)$ 为正规闭子群，限制映射 $G\to\Gal(E/F)$ 连续满射，核为 $H$，由紧群到 Hausdorff 群的连续双射性质得到 $G/H\simeq\Gal(E/F)$ 为拓扑群同构。
\end{proof}

\begin{exercise}
设 $k$ 为特征 $0$ 的代数闭域，$K=k((T))$。证明 $K^{\mathrm{alg}}=\bigcup_{n\ge1}k((T^{1/n}))$，并计算 Galois 群。
\end{exercise}
\begin{proof}[解答]
状态：可解，使用 Newton-Puiseux 定理。该定理说明：在特征 $0$ 且常数域代数闭时，Laurent 级数域的任意代数元可写成某个分数幂变量 $T^{1/n}$ 的 Laurent 级数。因此代数闭包为
\[
K^{\mathrm{alg}}=\bigcup_{n\ge1}k((T^{1/n})).
\]
对固定 $n$，扩张 $k((T^{1/n}))/k((T))$ 是循环 Galois 扩张，自同构由
\[
T^{1/n}\mapsto \zeta_n^aT^{1/n},\qquad a\in\Z/n\Z
\]
给出。若 $m\mid n$，限制映射对应自然投射 $\Z/n\Z\to\Z/m\Z$。取逆极限得到
\[
\Gal(K^{\mathrm{alg}}/K)\simeq\varprojlim_n\Z/n\Z=\widehat{\Z}.
\]
\end{proof}
 \chapter{理想}

\paragraph{第一部分的目标}
第一部分开始真正进入数域论。第零章准备的是域、Galois 群、迹、范、有限域、过滤和特征标；第一章把这些工具放到数域的整数环上，研究理想如何分解。最核心的问题是：若 $E/F$ 是有限 Galois 扩张，整数环扩张
\[
\mathcal O_F\subset \mathcal O_E
\]
中，$\mathcal O_F$ 的素理想在 $\mathcal O_E$ 中会怎样分解？这种分解又怎样被 Galois 群 $\Gal(E/F)$ 控制？

对初学者来说，可以先把第一章理解成一句话：普通整数里有素数分解；数域整数环里元素分解可能失败，但理想分解仍然可靠。代数数论的许多结构，就是建立在这个“把元素换成理想”的动作上。

称有理系数多项式 $f(X)\in\Q[X]$ 的根为代数数。若 $\alpha$ 是代数数，则 $\Q(\alpha)$ 是 $\Q$ 的有限扩张。一般称 $\Q$ 的有限扩张为数域。若 $\mu_m$ 是所有 $m$ 次单位根组成的群，则 $\Q(\mu_m)$ 称为 $m$ 次分圆域；它是后面类域论与 $L$ 函数的重要例子。

本章默认“环”是含乘法单位元的交换环。理想、素理想、极大理想、整元素、整闭包、Noether 条件，是本章的基本语法。

\begin{theorem}
设环 $A$ 有理想 $\mathfrak a_1,\dots,\mathfrak a_n$，并且当 $i\ne j$ 时
\[
\mathfrak a_i+\mathfrak a_j=A.
\]
则：
\begin{enumerate}[label=(\arabic*)]
\item 对任意 $x_1,\dots,x_n\in A$，存在 $x\in A$ 使
\[
x\equiv x_i\pmod{\mathfrak a_i},\qquad 1\le i\le n.
\]
\item 自然映射
\[
A\longrightarrow A/\mathfrak a_1\times\cdots\times A/\mathfrak a_n,\qquad
a\longmapsto(a+\mathfrak a_1,\dots,a+\mathfrak a_n)
\]
诱导环同构
\[
A/\bigcap_{i=1}^n\mathfrak a_i
\simeq
A/\mathfrak a_1\times\cdots\times A/\mathfrak a_n.
\]
\end{enumerate}
\end{theorem}

\begin{proof}
先看 $n=2$。由 $\mathfrak a_1+\mathfrak a_2=A$，可取 $u_1\in\mathfrak a_1$、$u_2\in\mathfrak a_2$ 使
\[
u_1+u_2=1.
\]
给定 $x_1,x_2$，令
\[
x=x_2u_1+x_1u_2.
\]
模 $\mathfrak a_1$ 时，$u_1\equiv0$，且 $u_2=1-u_1\equiv1$，所以 $x\equiv x_1$。模 $\mathfrak a_2$ 时，$u_2\equiv0$，且 $u_1\equiv1$，所以 $x\equiv x_2$。

一般 $n$ 的情形，构造“第 $j$ 个坐标为 $1$、其他坐标为 $0$”的元素。固定 $j$。对每个 $i\ne j$，由 $\mathfrak a_j+\mathfrak a_i=A$，取
\[
u_i\in\mathfrak a_j,\qquad v_i\in\mathfrak a_i,\qquad u_i+v_i=1.
\]
令
\[
y_j=\prod_{i\ne j}v_i.
\]
因为每个 $v_i\equiv0\pmod{\mathfrak a_i}$，所以 $y_j\equiv0\pmod{\mathfrak a_i}$ 对 $i\ne j$ 成立。又 $v_i=1-u_i\equiv1\pmod{\mathfrak a_j}$，所以 $y_j\equiv1\pmod{\mathfrak a_j}$。于是
\[
x=\sum_{j=1}^n x_jy_j
\]
满足所有同余条件。

第二点中，自然映射是环同态。第一点说明它满射。它的核由同时落在所有 $\mathfrak a_i$ 中的元素组成，即
\[
\ker=\bigcap_{i=1}^n\mathfrak a_i.
\]
由环同构定理得到所需同构。
\end{proof}

中国剩余定理后面会反复出现。它的真正作用是把“模一个由互素因子组成的理想”的问题，拆成“分别模每个因子”的问题。理想范数、局部化、有限商环和 Euler 乘积都会用到这个分解思想。

\section{Dedekind 环}
从整数环 $\Z$ 出发，先得到分式域 $\Q=\operatorname{Frac}(\Z)$。若 $F/\Q$ 是有限扩张，我们希望在 $F$ 中找到一个在 $F$ 里扮演整数环角色的环 $R$，满足
\[
\Z\subset R\subset F,\qquad \operatorname{Frac}(R)=F.
\]
这样的 $R$ 有很多，但代数数论需要一个最自然的选择：把 $\Z$ 在 $F$ 中的整闭包作为 $R$。这个环就是数域 $F$ 的整数环 $\mathcal O_F$。

本节的动机可以这样读：整数 $n\in\Z$ 在更大的整数环 $\mathcal O_F$ 中生成理想 $n\mathcal O_F$，它会怎样分解成素理想幂积？
\[
n\mathcal O_F=\mathfrak P_1^{e_1}\cdots\mathfrak P_r^{e_r}.
\]
这就是后面分解、分歧、惯性和 Frobenius 的起点。

\paragraph{整元素与整扩张}
设 $A\subset B$ 是环扩张，$b\in B$。如果存在首一多项式
\[
f(X)=X^n+a_{n-1}X^{n-1}+\cdots+a_0\in A[X]
\]
使 $f(b)=0$，则称 $b$ 在 $A$ 上整。首一条件非常重要：它保证可以把高次幂 $b^n$ 用低次幂线性表示：
\[
b^n=-a_{n-1}b^{n-1}-\cdots-a_0.
\]
由 $B$ 内所有在 $A$ 上整的元素组成的集合称为 $A$ 在 $B$ 中的整闭包。若 $B$ 中每个元素都在 $A$ 上整，则称 $B/A$ 是整环扩张。

\begin{lemma}
设 $A\subset B$ 是环扩张。对 $B$ 中元素 $b_1,\dots,b_n$，以下条件等价：
\begin{enumerate}[label=(\arabic*)]
\item $b_1,\dots,b_n$ 都在 $A$ 上整。
\item $A[b_1,\dots,b_n]$ 是有限生成 $A$-模。
\end{enumerate}
\end{lemma}

\begin{proof}
先证 (1) 推 (2)。对每个 $i$，取首一多项式 $f_i\in A[X]$ 使 $f_i(b_i)=0$，设 $d_i=\deg f_i$。因为 $f_i$ 首一，$b_i^{d_i}$ 可写成 $1,b_i,\dots,b_i^{d_i-1}$ 的 $A$-线性组合。继续乘以 $b_i$ 并归纳，得到任意高次幂 $b_i^m$ 都可写成这些低次幂的 $A$-线性组合。因此任意单项式
\[
b_1^{m_1}\cdots b_n^{m_n}
\]
都可降阶为有限集合
\[
\{b_1^{\ell_1}\cdots b_n^{\ell_n}:0\le \ell_i\le d_i-1\}
\]
的 $A$-线性组合。这个有限集合生成 $A[b_1,\dots,b_n]$，故它是有限生成 $A$-模。

再证 (2) 推 (1)。设
\[
A[b_1,\dots,b_n]=A\alpha_1+\cdots+A\alpha_r.
\]
固定 $b_i$。乘法映射
\[
\mu_{b_i}:A[b_1,\dots,b_n]\to A[b_1,\dots,b_n],\qquad x\mapsto b_ix
\]
把有限生成 $A$-模送到自身。于是存在 $a_{kl}\in A$，使
\[
b_i\alpha_k=\sum_{l=1}^r a_{kl}\alpha_l,\qquad k=1,\dots,r.
\]
记矩阵 $M=(a_{kl})$。把这些等式写成
\[
(b_iI-M)
\begin{pmatrix}
\alpha_1\\ \vdots\\ \alpha_r
\end{pmatrix}=0.
\]
对矩阵 $b_iI-M$ 乘伴随矩阵，得
\[
\det(b_iI-M)\alpha_k=0,\qquad k=1,\dots,r.
\]
因为 $1\in A[b_1,\dots,b_n]$ 可由 $\alpha_k$ 线性表示，推出
\[
\det(b_iI-M)=0.
\]
多项式 $\det(XI-M)\in A[X]$ 是首一多项式，且以 $b_i$ 为根，所以 $b_i$ 在 $A$ 上整。
\end{proof}

\begin{corollary}
若 $A\subset B$ 和 $B\subset C$ 都是整环扩张，则 $A\subset C$ 也是整环扩张。
\end{corollary}

\begin{proof}
取 $c\in C$。因为 $C/B$ 整，存在
\[
c^n+b_{n-1}c^{n-1}+\cdots+b_0=0,\qquad b_i\in B.
\]
因为每个 $b_i$ 在 $A$ 上整，上一个引理给出 $A[b_0,\dots,b_{n-1}]$ 是有限生成 $A$-模。再由上式，$c$ 在 $A[b_0,\dots,b_{n-1}]$ 上整，所以
\[
A[b_0,\dots,b_{n-1},c]
\]
是有限生成 $A[b_0,\dots,b_{n-1}]$-模，从而也是有限生成 $A$-模。由引理反向，$c$ 在 $A$ 上整。
\end{proof}

\paragraph{代数整数环}
有理系数多项式的根称为代数数。设 $F/\Q$ 是数域，令 $\mathcal O_F$ 为 $\Z$ 在 $F$ 中的整闭包。$\mathcal O_F$ 的元素称为 $F$ 中的代数整数，$\mathcal O_F$ 称为 $F$ 的代数整数环或整数环。可以证明
\[
\operatorname{Frac}(\mathcal O_F)=F.
\]
这是因为任意 $\alpha\in F$ 都满足某个有理系数多项式，清分母后可取整数 $m\ne0$ 使 $m\alpha$ 成为代数整数，于是 $\alpha=(m\alpha)/m$ 属于 $\operatorname{Frac}(\mathcal O_F)$。

\paragraph{整环、整闭、素理想与 Noether 条件}
若环 $R$ 中非零元素 $x,y$ 满足 $xy=0$，则称 $x$ 或 $y$ 为零因子。没有零因子且 $1\ne0$ 的环称为整环。等价地，$R$ 是整环当且仅当 $\{0\}$ 是素理想。

若 $R$ 是整环，可构造分式域 $\operatorname{Frac}(R)$。若 $R$ 在 $\operatorname{Frac}(R)$ 中的整闭包仍为 $R$，则称 $R$ 为整闭环。整闭的意思是：分式域里凡是满足首一方程的元素，已经在原环里。

理想 $\mathfrak p\ne R$ 称为素理想，若
\[
ab\in\mathfrak p\quad\Longrightarrow\quad a\in\mathfrak p\text{ 或 }b\in\mathfrak p.
\]
等价地，$R/\mathfrak p$ 是整环。理想 $\mathfrak m$ 称为极大理想，若 $\mathfrak m\ne R$ 且没有严格夹在 $\mathfrak m$ 与 $R$ 之间的理想。等价地，$R/\mathfrak m$ 是域。极大理想一定是素理想。

若 $R$ 的每个理想都是有限生成 $R$-模，则称 $R$ 为 Noether 环。Noether 条件的核心作用是终止性：理想不会出现无限严格上升链。

\begin{definition}
环 $R$ 称为 Dedekind 环，若满足：
\begin{enumerate}[label=(\arabic*)]
\item $R$ 是整环；
\item $R$ 是整闭环；
\item $R$ 是 Noether 环；
\item 每个非零素理想都是极大理想。
\end{enumerate}
\end{definition}

数域整数环确实是 Dedekind 环；源书在理想分解定理之后正式列出这个结论。这里先记住它的意义：即使 $\mathcal O_F$ 中元素唯一分解可能失败，理想仍然有唯一素理想分解。Dedekind 环的四条条件正是为了保证这套理想分解理论能运行。

\section{理想的分解}
在 Dedekind 环中，理想分解取代元素分解。这个替换是代数数论的基本转折：元素可能没有唯一分解，但理想会唯一分解为素理想的乘积。

设 $\mathcal O$ 是 Dedekind 环，$F=\operatorname{Frac}(\mathcal O)$。若 $\mathfrak a,\mathfrak b$ 是 $\mathcal O$ 的理想，说 $\mathfrak a$ 整除 $\mathfrak b$，记为 $\mathfrak a\mid\mathfrak b$，是指
\[
\mathfrak b\subset \mathfrak a.
\]
这个方向和整数整除看起来相反。例如在 $\Z$ 中，$(2)\mid(6)$，而作为集合有 $(6)\subset(2)$。

理想的三个基本运算如下：
\[
\mathfrak a\cap\mathfrak b
\]
是最小公倍，因为它是同时被 $\mathfrak a,\mathfrak b$ 整除的最大集合；
\[
\mathfrak a+\mathfrak b=\{a+b:a\in\mathfrak a,\ b\in\mathfrak b\}
\]
是最大公约，因为它是包含 $\mathfrak a,\mathfrak b$ 的最小理想；乘积定义为
\[
\mathfrak a\mathfrak b
=\left\{\sum_i a_i b_i:\ a_i\in\mathfrak a,\ b_i\in\mathfrak b,\ \text{有限和}\right\}.
\]
若 $\mathfrak p$ 是素理想，并且 $\mathfrak p\mid\mathfrak a\mathfrak b$，也就是 $\mathfrak a\mathfrak b\subset\mathfrak p$，则 $\mathfrak a\subset\mathfrak p$ 或 $\mathfrak b\subset\mathfrak p$。证明是：若两者都不包含于 $\mathfrak p$，取 $a\in\mathfrak a\setminus\mathfrak p$、$b\in\mathfrak b\setminus\mathfrak p$，则 $ab\in\mathfrak a\mathfrak b\subset\mathfrak p$，矛盾于素理想性。

对非零理想 $\mathfrak a\ne0$，定义
\[
\mathfrak a^{-1}=\{x\in F:x\mathfrak a\subset\mathcal O\}.
\]
这是把 $\mathfrak a$ “倒过来”的分式理想。取 $0\ne a_0\in\mathfrak a$，若 $x\in\mathfrak a^{-1}$，则 $xa_0\in\mathcal O$，所以
\[
x\in a_0^{-1}\mathcal O.
\]
因此 $\mathfrak a^{-1}$ 被包含在一个有限生成 $\mathcal O$-模中。因为 $\mathcal O$ Noether，$\mathfrak a^{-1}$ 也是有限生成 $\mathcal O$-模。若 $\mathfrak a\subset\mathfrak b$，则 $x\mathfrak b\subset\mathcal O$ 自动推出 $x\mathfrak a\subset\mathcal O$，所以
\[
\mathfrak a^{-1}\supset\mathfrak b^{-1}.
\]

\begin{lemma}
设 $\mathcal O$ 为 Dedekind 环。
\begin{enumerate}[label=(\arabic*)]
\item 若 $\mathfrak a$ 是 $\mathcal O$ 的非零理想，则存在非零素理想 $\mathfrak p_1,\dots,\mathfrak p_r$，使
\[
\mathfrak a\supset \mathfrak p_1\cdots\mathfrak p_r.
\]
\item 若 $\mathfrak a$ 是非零真理想，则 $\mathfrak a^{-1}\ne\mathcal O$。
\item 若 $\mathfrak a$ 是非零理想，则
\[
\mathfrak a^{-1}\mathfrak a=\mathcal O.
\]
\item 若
\[
\mathfrak p_1\cdots\mathfrak p_n\subset \mathfrak q_1\cdots\mathfrak q_m
\]
其中 $\mathfrak p_i,\mathfrak q_j$ 都是素理想，则 $m\le n$，且集合 $\{\mathfrak q_1,\dots,\mathfrak q_m\}$ 是 $\{\mathfrak p_1,\dots,\mathfrak p_n\}$ 的子集。
\end{enumerate}
\end{lemma}

\begin{proof}
先证 (1)。反设存在不包含任何有限个非零素理想乘积的非零理想。由 Noether 性，这些反例中有一个极大者，记为 $\mathfrak a$。它不是素理想，否则它本身就是一个素理想的乘积。于是存在 $x,y\notin\mathfrak a$ 但 $xy\in\mathfrak a$。理想
\[
\mathfrak a+x\mathcal O,\qquad \mathfrak a+y\mathcal O
\]
都严格大于 $\mathfrak a$，按极大性，它们分别包含某些素理想乘积。于是它们的乘积也包含一个素理想乘积。但
\[
(\mathfrak a+x\mathcal O)(\mathfrak a+y\mathcal O)
\subset \mathfrak a\mathfrak a+x\mathfrak a+y\mathfrak a+xy\mathcal O\subset\mathfrak a.
\]
所以 $\mathfrak a$ 也包含素理想乘积，矛盾。

证 (2)。包含 $\mathcal O\subset\mathfrak a^{-1}$ 总成立。要找一个不在 $\mathcal O$ 中的元素。取 $0\ne a\in\mathfrak a$。由 (1)，可取素理想 $\mathfrak p_1,\dots,\mathfrak p_r$ 使
\[
\mathfrak p_1\cdots\mathfrak p_r\subset a\mathcal O\subset\mathfrak a,
\]
并让 $r$ 最小。取极大理想 $\mathfrak p$ 包含 $\mathfrak a$。由于
\[
\mathfrak p_1\cdots\mathfrak p_r\subset\mathfrak p,
\]
素理想性推出某个 $\mathfrak p_i\subset\mathfrak p$。Dedekind 环中非零素理想极大，所以 $\mathfrak p_i=\mathfrak p$。重排后设 $\mathfrak p_1=\mathfrak p$。

由 $r$ 的最小性，$\mathfrak p_2\cdots\mathfrak p_r$ 不包含于 $a\mathcal O$。取
\[
b\in \mathfrak p_2\cdots\mathfrak p_r,\qquad b\notin a\mathcal O.
\]
令 $x=b/a\in F$。则 $x\notin\mathcal O$，否则 $b=ax\in a\mathcal O$。另一方面
\[
x\mathfrak a\subset x\mathfrak p
=a^{-1}b\mathfrak p
\subset a^{-1}\mathfrak p_1\cdots\mathfrak p_r
\subset \mathcal O.
\]
因此 $x\in\mathfrak a^{-1}\setminus\mathcal O$，故 $\mathfrak a^{-1}\ne\mathcal O$。

证 (3)。按 $\mathfrak a^{-1}$ 的定义，$\mathfrak a^{-1}\mathfrak a\subset\mathcal O$。记
\[
\mathfrak b=\mathfrak a^{-1}\mathfrak a.
\]
若 $\mathfrak b\ne\mathcal O$，由 (2) 得 $\mathfrak b^{-1}\ne\mathcal O$。取 $y\in\mathfrak b^{-1}$。由于
\[
\mathfrak a^{-1}y\mathfrak a
\subset \mathfrak a^{-1}\mathfrak b^{-1}\mathfrak b
\subset \mathfrak a^{-1},
\]
乘以 $\mathfrak a$ 可理解为 $y$ 把有限生成 $\mathcal O$-模 $\mathfrak a^{-1}$ 送到自身。于是 $\mathcal O[y]$ 作用在有限生成模 $\mathfrak a^{-1}$ 上，按整性判别，$y$ 在 $\mathcal O$ 上整。因为 $\mathcal O$ 整闭，$y\in\mathcal O$。这说明 $\mathfrak b^{-1}\subset\mathcal O$，但总有 $\mathcal O\subset\mathfrak b^{-1}$，故 $\mathfrak b^{-1}=\mathcal O$，与 (2) 矛盾。因此 $\mathfrak b=\mathcal O$。

证 (4)。由
\[
\mathfrak p_1\cdots\mathfrak p_n\subset\mathfrak q_1\cdots\mathfrak q_m\subset\mathfrak q_1
\]
和 $\mathfrak q_1$ 的素性，某个 $\mathfrak p_i\subset\mathfrak q_1$。两者都是非零素理想，故都是极大理想，于是 $\mathfrak p_i=\mathfrak q_1$。重排使 $\mathfrak p_1=\mathfrak q_1$。乘以逆理想 $\mathfrak p_1^{-1}$，用 (3) 抵消，得到
\[
\mathfrak p_2\cdots\mathfrak p_n\subset\mathfrak q_2\cdots\mathfrak q_m.
\]
归纳即可得到 $m\le n$，且每个 $\mathfrak q_j$ 都出现在 $\mathfrak p_i$ 中。
\end{proof}

\begin{theorem}
设 $\mathfrak a$ 是 Dedekind 环 $\mathcal O$ 的理想，且 $\mathfrak a\ne(0),(1)$。则 $\mathfrak a$ 可分解为非零素理想的乘积：
\[
\mathfrak a=\mathfrak p_1\cdots\mathfrak p_r.
\]
这些素理想除排列外由 $\mathfrak a$ 唯一决定。
\end{theorem}

\begin{proof}
唯一性由引理 1.5 (4) 直接给出：若同一个理想有两种素理想乘积分解，两边互相包含，应用 (4) 两次得到两边素理想多重集合相同。

证明存在性。由引理 1.5 (1)，$\mathfrak a$ 包含某个素理想乘积。取包含 $\mathfrak a$ 的极大理想 $\mathfrak p$；在 Dedekind 环中 $\mathfrak p$ 是非零素理想。若
\[
\mathfrak p_1\cdots\mathfrak p_r\subset\mathfrak a\subset\mathfrak p,
\]
则由 $\mathfrak p$ 的素性，某个 $\mathfrak p_i\subset\mathfrak p$。非零素理想都是极大理想，所以该 $\mathfrak p_i$ 等于 $\mathfrak p$。重排后可写成
\[
\mathfrak p\mathfrak p_2\cdots\mathfrak p_r\subset\mathfrak a.
\]
乘以 $\mathfrak p^{-1}$ 得
\[
\mathfrak p_2\cdots\mathfrak p_r\subset \mathfrak p^{-1}\mathfrak a\subset\mathcal O.
\]
若 $\mathfrak p^{-1}\mathfrak a=\mathcal O$，则 $\mathfrak a=\mathfrak p$，已经分解完毕。否则 $\mathfrak p^{-1}\mathfrak a$ 是更大的真理想，因为 $\mathfrak a\subset\mathfrak p$ 会推出
\[
\mathfrak a\subset \mathfrak p^{-1}\mathfrak a.
\]
并且它包含的素理想乘积长度比原来少一个。对这个长度归纳，得到 $\mathfrak p^{-1}\mathfrak a$ 是素理想乘积，于是 $\mathfrak a=\mathfrak p(\mathfrak p^{-1}\mathfrak a)$ 也是素理想乘积。
\end{proof}

\begin{corollary}
设 $\mathfrak a,\mathfrak b$ 是 $\mathcal O$ 的理想。
\begin{enumerate}[label=(\arabic*)]
\item $\mathfrak a\subset\mathfrak b$ 当且仅当存在 $\mathcal O$ 的理想 $\mathfrak c$ 使
\[
\mathfrak a=\mathfrak b\mathfrak c.
\]
\item 若
\[
\mathfrak a=\mathfrak p_1^{a_1}\cdots\mathfrak p_n^{a_n},\qquad
\mathfrak b=\mathfrak p_1^{b_1}\cdots\mathfrak p_n^{b_n},
\]
其中 $\mathfrak p_i$ 两两不同，$a_i,b_i\ge0$，则
\[
\mathfrak a+\mathfrak b=\prod_i\mathfrak p_i^{\min(a_i,b_i)},
\]
\[
\mathfrak a\cap\mathfrak b=\prod_i\mathfrak p_i^{\max(a_i,b_i)},
\]
并且
\[
\mathfrak a\mathfrak b=(\mathfrak a+\mathfrak b)(\mathfrak a\cap\mathfrak b).
\]
\end{enumerate}
\end{corollary}

\begin{proof}
把包含关系翻译成素理想指数即可。理想越小，素理想幂指数越大。因此 $\mathfrak a\subset\mathfrak b$ 等价于每个素理想上的指数满足 $a_i\ge b_i$，这又等价于存在非负指数 $c_i=a_i-b_i$ 使 $\mathfrak a=\mathfrak b\prod_i\mathfrak p_i^{c_i}$。

和理想 $\mathfrak a+\mathfrak b$ 是同时包含 $\mathfrak a,\mathfrak b$ 的最小理想，因此在每个素理想上的指数取较小值；交理想 $\mathfrak a\cap\mathfrak b$ 是同时包含在二者中的最大理想，因此指数取较大值。最后一式逐个素理想比较指数：
\[
a_i+b_i=\min(a_i,b_i)+\max(a_i,b_i).
\]
\end{proof}

\begin{proposition}
设 $\mathfrak a,\mathfrak b$ 是 Dedekind 环 $\mathcal O$ 的理想。
\begin{enumerate}[label=(\arabic*)]
\item 存在 $a\in\mathfrak a$ 和 $\mathcal O$ 的理想 $\mathfrak c$，使
\[
\mathfrak b+\mathfrak c=\mathcal O,\qquad \mathfrak a\mathfrak c=a\mathcal O.
\]
\item 加法群 $\mathcal O/\mathfrak b$ 同构于 $\mathfrak a/\mathfrak a\mathfrak b$。
\end{enumerate}
\end{proposition}

\begin{proof}
写
\[
\mathfrak a=\mathfrak p_1^{a_1}\cdots\mathfrak p_n^{a_n},\qquad
\mathfrak b=\mathfrak p_1^{b_1}\cdots\mathfrak p_n^{b_n},
\]
把未出现的指数补成 $0$。对每个 $i$，取
\[
x_i\in\mathfrak p_i^{a_i}\setminus\mathfrak p_i^{a_i+1}.
\]
由中国剩余定理，可取 $a\in\mathcal O$ 使
\[
a\equiv x_i\pmod{\mathfrak p_i^{a_i+1}}
\]
对所有 $i$ 成立。于是 $a$ 在每个 $\mathfrak p_i$ 上的指数正好是 $a_i$，从而
\[
a\mathcal O+\mathfrak a\mathfrak b=\mathfrak a.
\]
由推论 1.7，存在理想 $\mathfrak c$ 使 $a\mathcal O=\mathfrak a\mathfrak c$。把等式
\[
a\mathcal O+\mathfrak a\mathfrak b=\mathfrak a
\]
两边乘以 $\mathfrak a^{-1}$，得
\[
\mathfrak c+\mathfrak b=\mathcal O.
\]
并且 $\mathfrak a\mathfrak c=a\mathcal O$。

对 (2)，定义群同态
\[
\phi:\mathcal O\longrightarrow \mathfrak a/\mathfrak a\mathfrak b,\qquad
x\longmapsto ax+\mathfrak a\mathfrak b.
\]
因为 $a\mathcal O+\mathfrak a\mathfrak b=\mathfrak a$，任意 $\mathfrak a$ 中元素可模 $\mathfrak a\mathfrak b$ 写成 $ax$，所以 $\phi$ 满射。

若 $x\in\ker\phi$，则 $ax\in\mathfrak a\mathfrak b$。乘以 $\mathfrak c$ 得
\[
ax\mathfrak c\subset \mathfrak a\mathfrak b\mathfrak c.
\]
由于 $\mathfrak a\mathfrak c=a\mathcal O$，右边为 $a\mathfrak b$，所以 $ax\mathfrak c\subset a\mathfrak b$。在分式域中可消去非零元 $a$，得到 $x\mathfrak c\subset\mathfrak b$。又 $\mathfrak b+\mathfrak c=\mathcal O$，取 $u\in\mathfrak b$、$v\in\mathfrak c$ 使 $u+v=1$，则
\[
x=xu+xv\in\mathfrak b.
\]
所以 $\ker\phi\subset\mathfrak b$。反向包含由 $x\in\mathfrak b$ 时 $ax\in\mathfrak a\mathfrak b$ 得到。于是 $\ker\phi=\mathfrak b$，同构定理给出
\[
\mathcal O/\mathfrak b\simeq \mathfrak a/\mathfrak a\mathfrak b.
\]
\end{proof}

\begin{corollary}
设 $\mathfrak a$ 是 $\mathcal O$ 的非零理想，$0\ne c\in\mathfrak a$。则存在 $a\in\mathcal O$ 使
\[
\mathfrak a=a\mathcal O+c\mathcal O.
\]
\end{corollary}

\begin{proof}
把命题 1.8 应用于理想 $\mathfrak a$ 与 $c\mathcal O$。得到 $a\in\mathfrak a$ 和理想 $\mathfrak d$，使
\[
c\mathcal O+\mathfrak d=\mathcal O,\qquad \mathfrak a\mathfrak d=a\mathcal O.
\]
于是
\[
a\mathcal O+c\mathcal O\subset\mathfrak a.
\]
反过来，把 $c\mathcal O+\mathfrak d=\mathcal O$ 乘以 $\mathfrak a$，得
\[
\mathfrak a=c\mathfrak a+\mathfrak a\mathfrak d
=c\mathfrak a+a\mathcal O.
\]
由于 $c\mathfrak a\subset c\mathcal O$，右边包含在 $c\mathcal O+a\mathcal O$ 中。因此 $\mathfrak a\subset a\mathcal O+c\mathcal O$，得等号。
\end{proof}

\begin{theorem}
数域 $F$ 的整数环 $\mathcal O_F$ 是 Dedekind 环。
\end{theorem}

\begin{proof}
按定义要验证四点。第一，$\mathcal O_F\subset F$，而 $F$ 是域，所以 $\mathcal O_F$ 没有零因子，是整环。第二，$\mathcal O_F$ 是 $\Z$ 在 $F$ 中的整闭包；若 $x\in F$ 在 $\mathcal O_F$ 上整，则 $\mathcal O_F$ 在 $\Z$ 上整，而整性可传递，所以 $x$ 在 $\Z$ 上整，故 $x\in\mathcal O_F$。因此 $\mathcal O_F$ 在分式域 $F$ 中整闭。

第三，$\mathcal O_F$ 是 Noether 环。标准证明是先取 $F/\Q$ 的一组基，利用迹配对构造一个包含 $\mathcal O_F$ 的有限生成 $\Z$-模，从而 $\mathcal O_F$ 本身是有限生成 $\Z$-模；有限生成 $\Z$-模上的子模 Noether，于是 $\mathcal O_F$ 作为环 Noether。

第四，非零素理想极大。若 $0\ne\mathfrak p\subset\mathcal O_F$ 是素理想，取 $0\ne a\in\mathfrak p$。元素 $a$ 在 $\Z$ 上整，所以它的范数 $N_{F/\Q}(a)$ 是非零整数，并且 $N_{F/\Q}(a)\in a\mathcal O_F\cap\Z\subset\mathfrak p\cap\Z$。因此 $\mathfrak p\cap\Z=(p)$ 为某个素数 $p$。商环 $\mathcal O_F/\mathfrak p$ 是有限整环，因为它是有限 $\Z/(p)$-代数；有限整环是域，所以 $\mathfrak p$ 极大。
\end{proof}

\begin{proposition}
数域 $F$ 的整数环 $\mathcal O_F$ 内任意非零理想 $\mathfrak a$，作为加法群都是秩为 $[F:\Q]$ 的自由交换群。
\end{proposition}

\begin{proof}
整数环 $\mathcal O_F$ 是秩 $n=[F:\Q]$ 的自由 $\Z$-模。理想 $\mathfrak a$ 是 $\mathcal O_F$ 的非零子群，因此是有限生成无挠交换群，故为某个秩 $r$ 的自由交换群。取 $0\ne a\in\mathfrak a$。乘法给出嵌入
\[
\mathcal O_F\longrightarrow \mathfrak a,\qquad x\longmapsto ax,
\]
其像 $a\mathcal O_F\subset\mathfrak a$ 作为加法群秩为 $n$，故 $r\ge n$。另一方面 $\mathfrak a\subset\mathcal O_F$，所以 $r\le n$。因此 $r=n$。
\end{proof}

若
\[
\mathcal O_F=\Z x_1\oplus\cdots\oplus\Z x_n,
\]
则称 $\{x_1,\dots,x_n\}$ 为 $\mathcal O_F$ 或数域 $F$ 的整基。

\paragraph{分式理想}
设 $F$ 是 Dedekind 环 $\mathcal O$ 的分式域。把 $F$ 看作 $\mathcal O$-模。$F$ 内非零有限生成 $\mathcal O$-子模 $\mathfrak f$ 称为 $F$ 的分式理想。对分式理想定义
\[
\mathfrak f^{-1}=\{x\in F:x\mathfrak f\subset\mathcal O\},
\]
并定义乘积
\[
\mathfrak f\mathfrak g
=\left\{\sum_i x_i y_i:\ x_i\in\mathfrak f,\ y_i\in\mathfrak g,\ \text{有限和}\right\}.
\]

\begin{proposition}
设 $F$ 为 Dedekind 环 $\mathcal O$ 的分式域。
\begin{enumerate}[label=(\arabic*)]
\item 全部分式理想按乘法组成交换群，记为 $I_F$，称为分式理想群。
\item 若 $\mathfrak f$ 为分式理想，则 $\mathfrak f^{-1}$ 为分式理想。
\item 若 $\mathfrak f,\mathfrak g$ 为分式理想，则 $\mathfrak f\mathfrak g$ 为分式理想。
\item 若 $\mathfrak f$ 为分式理想，则存在两两不同的素理想 $\mathfrak p_1,\dots,\mathfrak p_n$ 和唯一的整数 $\nu_j\in\Z$，使
\[
\mathfrak f=\mathfrak p_1^{\nu_1}\cdots\mathfrak p_n^{\nu_n}.
\]
\end{enumerate}
\end{proposition}

\begin{proof}
任意分式理想 $\mathfrak f$ 有有限个生成元 $x_1,\dots,x_r\in F$。取共同分母 $0\ne d\in\mathcal O$，使 $dx_i\in\mathcal O$。于是 $d\mathfrak f$ 是 $\mathcal O$ 的非零整理想。由定理 1.6，
\[
d\mathfrak f=\prod_i\mathfrak p_i^{a_i},
\qquad
d\mathcal O=\prod_i\mathfrak p_i^{b_i}.
\]
于是
\[
\mathfrak f=(d\mathfrak f)(d\mathcal O)^{-1}
=\prod_i\mathfrak p_i^{a_i-b_i}.
\]
这给出存在性。唯一性由整理想唯一分解的唯一性给出。分式理想的乘法在指数上就是相加，逆理想在指数上就是取负，因此所有分式理想组成交换群，单位元为 $\mathcal O$，并且 $\mathfrak f^{-1}$ 与 $\mathfrak f\mathfrak g$ 仍为分式理想。
\end{proof}

\section{Dedekind 环扩张}
面对有限域扩张 $E/F$，如果只把 $E$ 看成 $F$-向量空间，我们关心的是基。但在代数数论中，$F$ 往往是 Dedekind 环 $\mathcal O$ 的分式域，$E$ 中有 $\mathcal O$ 的整闭包 $\mathcal O_E$。这时更细的问题是：能否找到 $e_1,\dots,e_n\in\mathcal O_E$，使每个 $y\in\mathcal O_E$ 都能唯一写成
\[
y=a_1e_1+\cdots+a_ne_n,\qquad a_i\in\mathcal O?
\]
若能，则称这组元素为 $\mathcal O_E/\mathcal O$ 的整基。数域相对于 $\Z$ 总有整基；一般 Dedekind 环上则要更谨慎。

\paragraph{迹配对与判别式}
设 $E/F$ 是有限可分扩张，$e_1,\dots,e_n$ 是 $E/F$ 的一组基，$\overline F$ 为 $F$ 的代数闭包，$\iota_1,\dots,\iota_n:E\to\overline F$ 是所有 $F$-嵌入。由第零章迹公式，
\[
\Tr_{E/F}(e_i e_j)=\sum_k \iota_k(e_i)\iota_k(e_j).
\]
因此矩阵
\[
\bigl(\Tr_{E/F}(e_i e_j)\bigr)_{i,j}
\]
可写成嵌入矩阵与其转置的乘积，得到
\[
\det\bigl(\Tr_{E/F}(e_i e_j)\bigr)
=\det\bigl(\iota_i(e_j)\bigr)^2.
\]
这个行列式就是这组基的判别式。它非零意味着迹配对非退化。

\begin{proposition}
设 $E/F$ 为有限可分域扩张，则
\[
(x,y)\longmapsto \Tr_{E/F}(xy)
\]
是 $E/F$ 上的非退化双线性映射。若 $e_1,\dots,e_n$ 是 $E/F$ 的一组基，则
\[
\det\bigl(\Tr_{E/F}(e_i e_j)\bigr)\ne0.
\]
\end{proposition}

\begin{proof}
由本原元素定理，取 $\theta$ 使 $E=F(\theta)$。此时 $1,\theta,\dots,\theta^{n-1}$ 是一组基。若 $\theta_1,\dots,\theta_n$ 是 $\theta$ 的全部共轭，则嵌入矩阵是 Vandermonde 矩阵
\[
(\theta_i^{j-1})_{i,j}.
\]
其行列式为
\[
\prod_{i<j}(\theta_j-\theta_i).
\]
可分性保证这些共轭两两不同，所以该行列式非零。因此迹矩阵的行列式
\[
\prod_{i<j}(\theta_j-\theta_i)^2
\]
非零。换成任意基只会把 Gram 矩阵作合同变换，行列式乘以换基矩阵行列式的平方，仍非零。
\end{proof}

\begin{proposition}
设 $\mathcal O$ 为整闭整环，$F$ 为其分式域，$E/F$ 为有限可分域扩张。以 $\mathcal O_E$ 记 $\mathcal O$ 在 $E$ 内的整闭包。则：
\begin{enumerate}[label=(\arabic*)]
\item 若 $e_i\in\mathcal O_E$ 且 $e_1,\dots,e_n$ 是 $E/F$ 的一组基，令
\[
d=\det(\Tr_{E/F}(e_i e_j)),
\]
则
\[
d\,\mathcal O_E\subset \sum_i\mathcal O e_i.
\]
\item 存在 $E/F$ 的一组基 $x_1,\dots,x_n$，使
\[
\mathcal O_E\subset \sum_i\mathcal O x_i.
\]
\item 若 $\mathcal O$ 是 Noether 整闭整环，则 $\mathcal O_E$ 是有限生成 $\mathcal O$-模。
\item 若 $\mathcal O$ 是主理想整环，则 $E$ 内任意非零有限生成 $\mathcal O_E$-模 $M$，作为 $\mathcal O$-模是秩 $[E:F]$ 的自由模。特别地，$\mathcal O_E/\mathcal O$ 有整基。
\end{enumerate}
\end{proposition}

\begin{proof}
(1) 取 $y\in\mathcal O_E$，写
\[
y=a_1e_1+\cdots+a_ne_n,\qquad a_i\in E.
\]
对每个 $i$ 取迹，得到线性方程组
\[
\Tr_{E/F}(e_i y)=\sum_j \Tr_{E/F}(e_i e_j)a_j.
\]
因为 $e_i,y$ 都在 $\mathcal O_E$ 中，乘积 $e_i y$ 在 $\mathcal O_E$ 中，其迹在 $\mathcal O$ 中。用 Cramer 法则解方程，得
\[
a_j=\frac{c_j}{d},\qquad c_j\in\mathcal O.
\]
所以 $dy=\sum_j c_j e_j\in\sum_j\mathcal O e_j$。

(2) 先说明每个 $x\in E$ 都可通过乘一个非零 $s\in\mathcal O$ 变成整元。设 $x$ 在 $F$ 上的极小多项式为
\[
X^m+c_{m-1}X^{m-1}+\cdots+c_0.
\]
取公分母 $s\in\mathcal O$ 使 $sc_i\in\mathcal O$。把方程改写给 $sx$，可见 $sx$ 满足 $\mathcal O$ 上首一方程，因此 $sx\in\mathcal O_E$。于是可取一组 $E/F$ 的基 $u_1,\dots,u_n$，并使 $u_i\in\mathcal O_E$。

对任意 $x\in\mathcal O_E$，写 $x=\sum_i b_i u_i$。设 $\sigma_1,\dots,\sigma_n$ 是 $E$ 到正规闭包中的全部 $F$-嵌入。由可分性，矩阵 $S=(\sigma_j(u_i))$ 行列式非零。对等式取每个 $\sigma_j$，得到线性方程组
\[
\sigma_j(x)=\sum_i b_i\sigma_j(u_i),\qquad j=1,\dots,n.
\]
用 Cramer 法则解 $b_i$。令 $\delta=\det(S)^2$，则 $\delta b_i$ 可写成 $\sigma_j(x)$ 与 $\sigma_j(u_i)$ 的整系数多项式。由于 $x,u_i$ 都在 $\mathcal O$ 上整，它们的共轭 $\sigma_j(x),\sigma_j(u_i)$ 也在 $\mathcal O$ 上整，所以 $\delta b_i$ 在 $\mathcal O$ 上整。又 $\delta b_i\in F$ 且 $\mathcal O$ 整闭，得到 $\delta b_i\in\mathcal O$。于是
\[
x\in\sum_i\mathcal O\,(u_i/\delta).
\]
令 $x_i=u_i/\delta$ 即得。

(3) 若 $\mathcal O$ Noether，则 $\sum_i\mathcal O x_i$ 是 Noether $\mathcal O$-模。由 (2)，$\mathcal O_E$ 是其子模，所以有限生成。

(4) 主理想整环上的有限生成无挠模自由。由 (1)，取 $E/F$ 的一组整基候选 $e_i\in\mathcal O_E$ 和非零 $d$，有
\[
d\mathcal O_E\subset\sum_i\mathcal O e_i.
\]
这把 $\mathcal O_E$ 控制在一个秩 $n$ 的自由 $\mathcal O$-模中，因此 $\mathcal O_E$ 的秩为 $n$。对任意非零有限生成 $\mathcal O_E$-模 $M$，取生成元并乘共同分母，可把 $M$ 嵌入一个有限秩自由 $\mathcal O$-模中；无挠性和主理想整环结构定理给出 $M$ 自由，秩比较为 $[E:F]$。
\end{proof}

\begin{lemma}
\begin{enumerate}[label=(\arabic*)]
\item 设 $A$ 为整闭整环，$B$ 为整环，且 $B/A$ 是整扩张。若 $\mathfrak P$ 是 $B$ 的非零素理想，则
\[
\mathfrak p=\mathfrak P\cap A
\]
是 $A$ 的非零素理想。
\item 若 $A$ 是域，$B$ 是整环，且 $B/A$ 是整扩张，则 $B$ 是域。
\end{enumerate}
\end{lemma}

\begin{proof}
(1) 先验证 $\mathfrak p=\mathfrak P\cap A$ 是素理想。若 $ab\in\mathfrak p$，则 $ab\in\mathfrak P$，由 $\mathfrak P$ 素得 $a\in\mathfrak P$ 或 $b\in\mathfrak P$。由于 $a,b\in A$，这给出 $a\in\mathfrak p$ 或 $b\in\mathfrak p$。关键是证明 $\mathfrak p\ne0$。

取 $0\ne b\in\mathfrak P$。因为 $b$ 在 $A$ 上整，它满足某个 $A[X]$ 中首一多项式。令 $f(X)$ 为 $b$ 在 $K=\operatorname{Frac}(A)$ 上的极小多项式。由于 $f$ 整除上述首一多项式，$f$ 的根也都在 $A$ 上整；于是 $f$ 的系数作为根的对称多项式也在 $A$ 上整。它们又属于 $K$，而 $A$ 整闭，所以这些系数属于 $A$。

写
\[
f(X)=X^m+a_{m-1}X^{m-1}+\cdots+a_0\in A[X].
\]
因为 $f$ 是极小多项式，$a_0\ne0$。由 $f(b)=0$ 得
\[
a_0=-b(b^{m-1}+a_{m-1}b^{m-2}+\cdots+a_1)\in\mathfrak P.
\]
所以 $0\ne a_0\in\mathfrak P\cap A$。

(2) 若 $B$ 不是域，则存在非零极大理想 $\mathfrak P\subset B$。由 (1)，$\mathfrak P\cap A$ 是 $A$ 的非零素理想。但域没有非零真素理想，只能有 $0$，矛盾。因此 $B$ 是域。
\end{proof}

\begin{theorem}
设 $\mathcal O$ 为 Dedekind 环，$F$ 为其分式域，$E/F$ 为有限可分域扩张。以 $\mathcal O_E$ 记 $\mathcal O$ 在 $E$ 内的整闭包，则 $\mathcal O_E$ 也是 Dedekind 环。
\end{theorem}

\begin{proof}
首先，$\mathcal O_E\subset E$，而 $E$ 是域，所以它是整环。它按定义是整闭包，因此在 $E$ 中整闭。

其次证明 Noether 性。由命题 1.14，$\mathcal O_E$ 是有限生成 $\mathcal O$-模。若 $\mathfrak A$ 是 $\mathcal O_E$ 的理想，则它是有限生成 $\mathcal O$-模的子模。因为 $\mathcal O$ Noether，$\mathfrak A$ 有限生成为 $\mathcal O$-模，从而也有限生成为 $\mathcal O_E$-模。因此 $\mathcal O_E$ Noether。

最后证明非零素理想极大。取 $\mathcal O_E$ 的非零素理想 $\mathfrak P$。由引理 1.15，
\[
\mathfrak p=\mathfrak P\cap\mathcal O
\]
是 $\mathcal O$ 的非零素理想；Dedekind 条件给出 $\mathfrak p$ 极大，所以 $\mathcal O/\mathfrak p$ 是域。商环 $\mathcal O_E/\mathfrak P$ 是 $\mathcal O/\mathfrak p$ 上的整扩张，因为 $\mathcal O_E$ 在 $\mathcal O$ 上整。由引理 1.15 (2)，$\mathcal O_E/\mathfrak P$ 是域，所以 $\mathfrak P$ 极大。
\end{proof}

以下常把 $\mathcal O$ 记为 $\mathcal O_F$。设 $\mathfrak p$ 是 $\mathcal O_F$ 的素理想。它在 $\mathcal O_E$ 中生成的理想为
\[
\mathfrak p\mathcal O_E.
\]
这个理想不是整个 $\mathcal O_E$。若 $\mathfrak p\mathcal O_E=\mathcal O_E$，则模掉 $\mathfrak p$ 后会得到零环；但 $\mathcal O_E/\mathfrak p\mathcal O_E$ 是一个非零有限 $\mathcal O_F/\mathfrak p$-代数。

由于 $\mathcal O_E$ 是 Dedekind 环，$\mathfrak p\mathcal O_E$ 有唯一素理想分解：
\[
\mathfrak p\mathcal O_E=\mathfrak P_1^{e_1}\cdots\mathfrak P_r^{e_r}.
\]
每个 $\mathfrak P_i$ 都满足
\[
\mathfrak P_i\cap\mathcal O_F=\mathfrak p.
\]
称 $\mathfrak P_i$ 为 $\mathfrak p$ 的素除子，记作 $\mathfrak P_i\mid\mathfrak p$。指数 $e_i$ 称为 $\mathfrak P_i$ 在 $\mathfrak p$ 上的分歧指标。包含 $\mathcal O_F\hookrightarrow\mathcal O_E$ 诱导剩余域扩张
\[
\kappa_{\mathfrak p}=\mathcal O_F/\mathfrak p
\subset
\kappa_{\mathfrak P_i}=\mathcal O_E/\mathfrak P_i.
\]
其次数
\[
f_i=[\kappa_{\mathfrak P_i}:\kappa_{\mathfrak p}]
\]
称为惯性次数、剩余域次数或剩余次数。

\begin{theorem}
设 $E/F$ 为有限可分域扩张，$n=[E:F]$。若 $\mathcal O_F$ 的素理想 $\mathfrak p$ 在 $\mathcal O_E$ 中分解为
\[
\mathfrak p\mathcal O_E=\mathfrak P_1^{e_1}\cdots\mathfrak P_r^{e_r},
\]
并令
\[
f_i=[\mathcal O_E/\mathfrak P_i:\mathcal O_F/\mathfrak p],
\]
则
\[
\sum_{i=1}^r e_i f_i=n.
\]
\end{theorem}

\begin{proof}
由中国剩余定理，
\[
\mathcal O_E/\mathfrak p\mathcal O_E
\simeq
\bigoplus_{i=1}^r \mathcal O_E/\mathfrak P_i^{e_i}.
\]
把这些商都看成 $\kappa=\mathcal O_F/\mathfrak p$ 上的向量空间。只需证明两件事：
\[
\dim_\kappa(\mathcal O_E/\mathfrak p\mathcal O_E)=n,
\]
以及
\[
\dim_\kappa(\mathcal O_E/\mathfrak P_i^{e_i})=e_i\,\dim_\kappa(\mathcal O_E/\mathfrak P_i)=e_if_i.
\]

先证第一式。取 $w_1,\dots,w_m\in\mathcal O_E$，使其剩余类构成 $\mathcal O_E/\mathfrak p\mathcal O_E$ 的 $\kappa$-基。证明 $w_1,\dots,w_m$ 是 $E/F$ 的基。

若它们在 $F$ 上线性相关，清分母后有
\[
a_1w_1+\cdots+a_mw_m=0,\qquad a_i\in\mathcal O_F,
\]
且不全为零。令 $\mathfrak a=(a_1,\dots,a_m)$。选 $t\in\mathfrak a^{-1}$ 使 $t\mathfrak a$ 不包含于 $\mathfrak p$；于是 $ta_i$ 的模 $\mathfrak p$ 类不全为零。把上式乘以 $t$ 并模 $\mathfrak p$，得到 $[w_i]$ 的非平凡线性关系，矛盾。因此 $w_i$ 线性无关。

再证生成。令
\[
M=\mathcal O_Fw_1+\cdots+\mathcal O_Fw_m,\qquad N=\mathcal O_E/M.
\]
因为 $[w_i]$ 生成 $\mathcal O_E/\mathfrak p\mathcal O_E$，有
\[
\mathcal O_E=M+\mathfrak p\mathcal O_E,
\]
所以 $N=\mathfrak pN$。模 $N$ 是有限生成 $\mathcal O_F$-模。取生成元 $x_1,\dots,x_s$，由 $N=\mathfrak pN$ 可写
\[
x_i=\sum_j a_{ij}x_j,\qquad a_{ij}\in\mathfrak p.
\]
即 $(A-I)x=0$。取行列式 $d=\det(A-I)$，Cramer 法则给出 $dN=0$。模 $\mathfrak p$ 时 $d\equiv(-1)^s\ne0$，所以 $d\ne0$，在 $F$ 中可逆。于是 $E=Fw_1+\cdots+Fw_m$。故 $m=n$。

再证第二式。固定 $\mathfrak P=\mathfrak P_i$，$e=e_i$。有过滤
\[
\mathcal O_E/\mathfrak P^e\supset
\mathfrak P/\mathfrak P^e\supset\cdots\supset
\mathfrak P^{e-1}/\mathfrak P^e\supset0.
\]
各层商为
\[
\mathfrak P^k/\mathfrak P^{k+1},\qquad k=0,\dots,e-1.
\]
取 $y\in\mathfrak P^k\setminus\mathfrak P^{k+1}$。由 Dedekind 环理想分解，$\mathfrak P^k$ 是 $y\mathcal O_E$ 与 $\mathfrak P^{k+1}$ 的最大公约，于是
\[
\mathfrak P^k=y\mathcal O_E+\mathfrak P^{k+1}.
\]
映射
\[
\mathcal O_E\longrightarrow \mathfrak P^k/\mathfrak P^{k+1},\qquad a\longmapsto ay
\]
满射，核为 $\mathfrak P$，所以
\[
\mathcal O_E/\mathfrak P\simeq \mathfrak P^k/\mathfrak P^{k+1}.
\]
因此每一层的 $\kappa$-维数都是 $f_i$，合计为 $e_if_i$。
\end{proof}

这个公式是第一章的核心之一。它说一个素理想在扩张中分解为若干上方素理想时，总次数由两类数据相加得到：
\[
\text{厚度 }e_i\quad\text{乘以}\quad\text{剩余域扩张次数 }f_i.
\]
若所有 $e_i=f_i=1$ 且 $r=n$，称 $\mathfrak p$ 在 $E$ 中完全分裂。若 $r=1,e_1=1,f_1=n$，称 $\mathfrak p$ 无分解或无分裂。若 $r=1,e_1=n,f_1=1$，称 $\mathfrak p$ 完全分歧。若所有 $e_i=1$，称 $\mathfrak p$ 在 $E/F$ 中无分歧；若某个 $e_i>1$，称 $\mathfrak p$ 分歧。

\section{理想的迹和范数}

上一节把一个素理想 $\mathfrak p\subset\OF$ 在扩张 $\mathcal O_E$ 中怎样分解讲清楚了。接下来要把这种分解变成可以计算的数值。元素有迹 $\Tr_{E/F}$ 和范 $\Nm_{E/F}$，理想也有范数；区别在于元素范数来自线性映射的行列式，理想范数来自商环的元素个数。对初学者来说，最好把理想范数理解为“模掉这个理想后还剩多少个剩余类”。

设 $F/\Q$ 为数域，$\OF$ 为整数环，$\mathfrak a\subset\OF$ 为非零理想。因为 $\OF$ 是秩 $[F:\Q]$ 的自由 $\Z$-模，而非零理想 $\mathfrak a$ 是有限指数子格，所以商 $\OF/\mathfrak a$ 是有限集。

\begin{definition}
非零整理想 $\mathfrak a\subset\OF$ 的绝对范定义为
\[
\Nm(\mathfrak a)=|\OF/\mathfrak a|.
\]
若 $\mathfrak A$ 是扩张数域 $E$ 的非零理想，仍用 $\Nm(\mathfrak A)$ 表示 $|\mathcal O_E/\mathfrak A|$。
\end{definition}

\begin{remark}
这个定义看起来只是在数元素个数，但它把理想和同余方程连接起来。例如在 $\Z$ 中，$\Nm((m))=|\Z/m\Z|=|m|$。因此理想范数是“整数绝对值”的数域推广。
\end{remark}

\begin{proposition}
设 $0\ne a\in\OF$，用 $(a)=a\OF$ 表示主理想。
\begin{enumerate}[label=(\arabic*)]
\item $\Nm((a))=|\Nm_{F/\Q}(a)|$。
\item 若非零理想 $\mathfrak a$ 的素理想分解为
\[
\mathfrak a=\mathfrak p_1^{\nu_1}\cdots\mathfrak p_r^{\nu_r},
\]
则
\[
\Nm(\mathfrak a)=\Nm(\mathfrak p_1)^{\nu_1}\cdots
\Nm(\mathfrak p_r)^{\nu_r}.
\]
\item 对任意非零理想 $\mathfrak a,\mathfrak b\subset\OF$，有
\[
\Nm(\mathfrak a\mathfrak b)=\Nm(\mathfrak a)\Nm(\mathfrak b).
\]
若 $\mathfrak a,\mathfrak b$ 为非零整理想，定义分式理想的范数为
\[
\Nm(\mathfrak a\mathfrak b^{-1})=\Nm(\mathfrak a)\Nm(\mathfrak b)^{-1},
\]
则得到群同态 $\Nm:I_F\to\Q_{>0}^{\times}$。
\end{enumerate}
\end{proposition}

\begin{proof}
(1) 令 $n=[F:\Q]$。取 $\OF$ 的 $\Z$-基 $u_1,\dots,u_n$。乘以 $a$ 是 $\Q$-向量空间 $F$ 上的线性变换；在这组基下写
\[
au_j=\sum_i a_{ij}u_i,\qquad a_{ij}\in\Z.
\]
主理想 $(a)$ 作为 $\Z$-模有基 $au_1,\dots,au_n$。所以 $\OF/(a)$ 的指数等于格变换矩阵 $(a_{ij})$ 的行列式绝对值，即
\[
|\OF/(a)|=|\det(a_{ij})|.
\]
而域范数 $\Nm_{F/\Q}(a)$ 的定义正是“乘以 $a$”这一 $\Q$-线性变换的行列式，因此得到结论。

(2) 先由中国剩余定理得到
\[
\OF/\mathfrak a\simeq
\OF/\mathfrak p_1^{\nu_1}\oplus\cdots\oplus
\OF/\mathfrak p_r^{\nu_r}.
\]
于是商集大小相乘。现在只需证明对非零素理想 $\mathfrak p$ 有
\[
|\OF/\mathfrak p^\nu|=|\OF/\mathfrak p|^\nu.
\]
考虑过滤
\[
\OF/\mathfrak p^\nu\supset
\mathfrak p/\mathfrak p^\nu\supset\cdots\supset
\mathfrak p^{\nu-1}/\mathfrak p^\nu\supset0.
\]
它的相邻商是 $\mathfrak p^i/\mathfrak p^{i+1}$。取
$x\in\mathfrak p^i\setminus\mathfrak p^{i+1}$。因为 $(x)\subset\mathfrak p^i$ 而
$(x)\not\subset\mathfrak p^{i+1}$，主理想 $(x)$ 的素理想分解在 $\mathfrak p$ 处的指数恰为
$i$。对任一非零素理想 $\mathfrak q$，理想和 $(x)+\mathfrak p^{i+1}$ 在
$\mathfrak q$ 处的指数是 $(x)$ 与 $\mathfrak p^{i+1}$ 两个指数的较小值：当
$\mathfrak q=\mathfrak p$ 时得到 $\min(i,i+1)=i$；当
$\mathfrak q\ne\mathfrak p$ 时，$\mathfrak p^{i+1}$ 在 $\mathfrak q$ 处指数为 $0$，较小值为
$0$。于是这个理想和的素理想分解正是
\[
\mathfrak p^i=x\OF+\mathfrak p^{i+1}.
\]
于是映射
\[
\OF\longrightarrow \mathfrak p^i/\mathfrak p^{i+1},\qquad
y\longmapsto xy
\]
满射。其核由满足 $xy\in\mathfrak p^{i+1}$ 的 $y$ 组成。因为 $x\OF$ 在 $\mathfrak p$ 处的指数恰为 $i$，这个条件等价于 $y\in\mathfrak p$。所以
\[
\mathfrak p^i/\mathfrak p^{i+1}\simeq\OF/\mathfrak p.
\]
每一层大小都是 $\Nm(\mathfrak p)$，共有 $\nu$ 层。

(3) 由 (2) 对素理想分解逐项相加指数即可。分式理想情形的定义与乘法性相容，因为同一个分式理想可写成不同的 $\mathfrak a\mathfrak b^{-1}$；若
$\mathfrak a\mathfrak b^{-1}=\mathfrak c\mathfrak d^{-1}$，则
$\mathfrak a\mathfrak d=\mathfrak c\mathfrak b$，整理想乘法性给出
$\Nm(\mathfrak a)\Nm(\mathfrak d)=\Nm(\mathfrak c)\Nm(\mathfrak b)$。
\end{proof}

现在设 $E/F$ 是有限数域扩张。对 $\OF$ 的非零素理想 $\mathfrak p$，写
\[
\mathfrak p\mathcal O_E=\prod_{\mathfrak P\mid\mathfrak p}\mathfrak P^{e_{\mathfrak P|\mathfrak p}}.
\]
这给出理想的扩张同态
\[
i_{E/F}:I_F\longrightarrow I_E,\qquad
i_{E/F}(\mathfrak a)=\mathfrak a\mathcal O_E.
\]
若
\[
\mathfrak A=\prod_{\mathfrak P}\mathfrak P^{n_{\mathfrak P}}
\]
是 $E$ 的分式理想，则定义相对范
\[
\Nm_{E/F}(\mathfrak A)=
\prod_{\mathfrak P}
(\mathfrak P\cap\OF)^{f_{\mathfrak P|(\mathfrak P\cap\OF)}n_{\mathfrak P}},
\]
其中
\[
f_{\mathfrak P|\mathfrak p}
=[\mathcal O_E/\mathfrak P:\OF/\mathfrak p].
\]
特别地，若 $\mathfrak P\cap\OF=\mathfrak p$，则
\[
\Nm_{E/F}(\mathfrak P)=\mathfrak p^{f_{\mathfrak P|\mathfrak p}}.
\]

\begin{proposition}
设 $K\supset E\supset F$ 为有限数域扩张。
\begin{enumerate}[label=(\arabic*)]
\item $i_{K/F}=i_{K/E}\circ i_{E/F}$，且 $\Nm_{K/F}=\Nm_{E/F}\circ\Nm_{K/E}$。
\item 若 $\mathfrak a\in I_F$，则
\[
\Nm_{E/F}(i_{E/F}\mathfrak a)=\mathfrak a^{[E:F]}.
\]
\item 若 $\mathfrak A$ 是 $\mathcal O_E$ 的非零整理想，则
\[
\Nm(\Nm_{E/F}(\mathfrak A))=\Nm(\mathfrak A).
\]
\item 若 $\mathfrak A$ 是 $\mathcal O_E$ 的非零整理想，则 $\Nm_{E/F}(\mathfrak A)$ 是由所有元素范数 $\Nm_{E/F}(x)$，$x\in\mathfrak A$，生成的 $\OF$-理想。
\end{enumerate}
\end{proposition}

\begin{proof}
(1) 对素理想检查即可。设 $\mathfrak p\subset\OF$，$\mathfrak P\subset\mathcal O_E$ 在 $\mathfrak p$ 上方，$\mathfrak Q\subset\mathcal O_K$ 在 $\mathfrak P$ 上方。分歧指标和惯性次数满足塔式公式
\[
e_{\mathfrak Q|\mathfrak p}=e_{\mathfrak Q|\mathfrak P}e_{\mathfrak P|\mathfrak p},
\qquad
f_{\mathfrak Q|\mathfrak p}=f_{\mathfrak Q|\mathfrak P}f_{\mathfrak P|\mathfrak p}.
\]
把这些指数代入扩张理想和相对范的定义，就得到两个复合公式。

(2) 只需令 $\mathfrak a=\mathfrak p$ 为素理想。若
\[
\mathfrak p\mathcal O_E=\prod_{\mathfrak P\mid\mathfrak p}\mathfrak P^{e_{\mathfrak P|\mathfrak p}},
\]
则
\[
\Nm_{E/F}(\mathfrak p\mathcal O_E)
=\prod_{\mathfrak P\mid\mathfrak p}
\mathfrak p^{e_{\mathfrak P|\mathfrak p}f_{\mathfrak P|\mathfrak p}}
=\mathfrak p^{\sum e_{\mathfrak P|\mathfrak p}f_{\mathfrak P|\mathfrak p}}
=\mathfrak p^{[E:F]}.
\]

(3) 对 $\mathfrak A=\mathfrak P$ 为素理想时，
\[
\Nm(\Nm_{E/F}(\mathfrak P))
=\Nm(\mathfrak p^{f_{\mathfrak P|\mathfrak p}})
=|\OF/\mathfrak p|^{f_{\mathfrak P|\mathfrak p}}
=|\mathcal O_E/\mathfrak P|
=\Nm(\mathfrak P).
\]
一般情形由素理想分解和乘法性推出。

(4) 先看包含方向。若 $x\in\mathfrak A$，则主理想 $x\mathcal O_E$ 被 $\mathfrak A$ 整除，即 $x\mathcal O_E\subset\mathfrak A$。把两边写成素理想指数，得到每个 $\mathfrak P$ 上
\[
v_{\mathfrak P}(x\mathcal O_E)\ge v_{\mathfrak P}(\mathfrak A).
\]
取相对范后，每个 $\mathfrak p\subset\OF$ 上的指数满足
\[
v_{\mathfrak p}\bigl(\Nm_{E/F}(x)\mathcal O_F\bigr)
\ge
v_{\mathfrak p}\bigl(\Nm_{E/F}(\mathfrak A)\bigr).
\]
所以所有元素范数生成的理想包含在 $\Nm_{E/F}(\mathfrak A)$ 中。

反向包含使用 Dedekind 环的近似定理：给定有限多个素理想及整数指数，可以取元素使其在这些素理想处具有指定指数，并在预先指定的其他素理想处指数不低于给定下界。把 $\mathfrak A=\prod_{\mathfrak P}\mathfrak P^{n_{\mathfrak P}}$。对固定的 $\mathfrak p$，在每个 $\mathfrak P|\mathfrak p$ 处取元素 $x_{\mathfrak p}\in\mathfrak A$，使 $v_{\mathfrak P}(x_{\mathfrak p})=n_{\mathfrak P}$ 至少对一个达到相对范最小指数的 $\mathfrak P$ 成立，并使其他相关素理想处指数不降低。于是 $\Nm_{E/F}(x_{\mathfrak p})$ 在 $\mathfrak p$ 处达到 $\Nm_{E/F}(\mathfrak A)$ 的指数。对有限多个出现在 $\Nm_{E/F}(\mathfrak A)$ 中的 $\mathfrak p$ 分别取这些元素，所得范数生成的理想在每个素理想处指数都不大于 $\Nm_{E/F}(\mathfrak A)$ 的指数。结合前一段包含关系，二者相等。

这一步的意义是：元素范数和理想范数不是两套互不相关的定义，它们在“由所有元素范数生成的理想”这个层面完全一致。若不调用近似定理，则本断言应作为 Dedekind 理想理论中的标准输入使用。
\end{proof}

\section{判别式}

判别式的作用是检测“迹配对有多退化”。在数域扩张里，迹配对
\[
(x,y)\longmapsto \Tr_{E/F}(xy)
\]
在 $E$ 作为 $F$-向量空间上非退化；这意味着如果 $x\ne0$，总能找到 $y$ 使 $\Tr_{E/F}(xy)\ne0$。但当我们把 $E$ 换成整数环 $B=\mathcal O_E$，配对的取值是否仍落在基环 $A=\OF$，以及落得多“整”，就成为一个精细问题。差别式、余差别式和判别式都是对这个问题的回答。

本节设 $A$ 为 Dedekind 环，$F$ 为其分式域，$E/F$ 为有限可分扩张，$B$ 为 $A$ 在 $E$ 中的整闭包。对 $E$ 的 $A$-子模 $M$，定义迹对偶模
\[
M'=\{x\in E:\Tr_{E/F}(xM)\subset A\}.
\]

\begin{definition}
$B'= \{x\in E:\Tr_{E/F}(xB)\subset A\}$ 称为 $B/A$ 的余差别式，记为 $C_{B/A}$。它的逆分式理想
\[
D_{B/A}=C_{B/A}^{-1}
\]
称为 $B/A$ 的差别式。因为 $1\in B'$ 且 $B\subset B'$，所以 $D_{B/A}\subset B$，即差别式是 $B$ 的整理想。
\end{definition}

\begin{proposition}
设 $\mathfrak a$ 是 $F$ 的分式理想，$\mathfrak b$ 是 $E$ 的分式理想。则
\[
\Tr_{E/F}(\mathfrak b)\subset\mathfrak a
\quad\Longleftrightarrow\quad
\mathfrak b\subset\mathfrak a D_{B/A}^{-1}.
\]
\end{proposition}

\begin{proof}
若 $\mathfrak a=0$，本节实际使用的是非零分式理想情形；零理想情形可从非退化迹配对单独处理。下面设
$\mathfrak a\ne0$。把包含式除以 $\mathfrak a$，得到
\[
\Tr(\mathfrak b)\subset\mathfrak a
\Longleftrightarrow
\Tr(\mathfrak a^{-1}\mathfrak b)\subset A.
\]
按迹对偶模的定义，后一条件等价于
\[
\mathfrak a^{-1}\mathfrak b\subset B'=D_{B/A}^{-1}.
\]
再乘回 $\mathfrak a$，得到所需包含。
\end{proof}

\begin{proposition}
在上述记号下有以下性质。
\begin{enumerate}[label=(\arabic*)]
\item 传递性：若 $K/E$ 也是有限可分扩张，$C$ 是 $A$ 在 $K$ 中的整闭包，则
\[
D_{C/A}=D_{C/B}\,D_{B/A}.
\]
\item 局部化：若 $S\subset A$ 为乘性子集，则
\[
S^{-1}D_{B/A}=D_{S^{-1}B/S^{-1}A}.
\]
\end{enumerate}
\end{proposition}

\begin{proof}
(1) 关键是迹的塔式公式
\[
\Tr_{K/F}=\Tr_{E/F}\circ\Tr_{K/E}.
\]
对 $K$ 的分式理想 $\mathfrak c$，按上一命题反复翻译“迹落在基环”：
\[
\mathfrak c\subset D_{C/B}^{-1}
\Longleftrightarrow
\Tr_{K/E}(\mathfrak c)\subset B,
\]
而
\[
\Tr_{K/E}(\mathfrak c)\subset D_{B/A}^{-1}
\Longleftrightarrow
\Tr_{E/F}(\Tr_{K/E}(\mathfrak c))\subset A.
\]
两者合并后，$K/F$ 的迹对偶正是 $D_{B/A}^{-1}D_{C/B}^{-1}$，取逆得到
$D_{C/A}=D_{C/B}D_{B/A}$。

(2) 先证明余差别式局部化：
\[
S^{-1}B'=(S^{-1}B)'.
\]
若 $x=y/s$，$y\in B'$，则对任意 $b\in B$ 有 $\Tr(xb)=\Tr(yb)/s\in S^{-1}A$，所以 $S^{-1}B'\subset(S^{-1}B)'$。反过来，取 $B$ 的有限 $A$-生成元 $b_1,\dots,b_m$。若 $x\in(S^{-1}B)'$，则每个 $\Tr(xb_i)$ 都属于 $S^{-1}A$。取一个公共分母 $s\in S$，使 $s\Tr(xb_i)\in A$ 对所有 $i$ 成立。由于 $b_i$ 生成 $B$，得到 $sx\in B'$，即 $x\in S^{-1}B'$。取逆理想即得差别式局部化公式。
\end{proof}

\begin{definition}
取 $E/F$ 的所有 $F$-基 $\alpha_1,\dots,\alpha_n$，要求 $\alpha_i\in B$。由行列式
\[
\det\bigl(\Tr_{E/F}(\alpha_i\alpha_j)\bigr)
\]
生成的 $A$-理想称为 $B/A$ 的判别式，记为 $d_{B/A}$；在数域情形也记为 $d_{E/F}$，绝对判别式记为 $d_E$ 或 $d(\mathcal O_E)$。
\end{definition}

\begin{theorem}
判别式与差别式满足
\[
d_{B/A}=\Nm_{E/F}(D_{B/A}).
\]
\end{theorem}

\begin{proof}
这是理想等式，检验它可以逐素理想局部化。由差别式、余差别式和判别式的局部化性质，把
$A$ 局部化到任一非零素理想后，可假设 $A$ 是离散赋值环。此时 $B$ 是有限自由
$A$-模；取 $B$ 的 $A$-基 $\alpha_1,\dots,\alpha_n$，令
\[
T=(\Tr_{E/F}(\alpha_i\alpha_j))_{i,j}.
\]
判别式理想由 $\det T$ 生成。

记 $\alpha_1',\dots,\alpha_n'$ 为迹配对的对偶基：
\[
\Tr_{E/F}(\alpha_i\alpha_j')=\delta_{ij}.
\]
这组对偶基生成 $B'$。换句话说，$B'$ 是“所有与 $B$ 配对后仍落在
$A$ 中的元素”所成的格。由于局部 Dedekind 环上非零分式理想为主理想，存在
$\beta\in E^\times$ 使 $B'=\beta B$，于是 $D_{B/A}=(B')^{-1}=\beta^{-1}B$。

下面把“对偶格”这句话翻成行列式计算。若用 $\alpha_i'$ 作基，则对偶关系说明
$(\Tr(\alpha_i\alpha_j'))$ 是单位矩阵，因此从 $\alpha_i$ 到 $\alpha_i'$ 的换基矩阵正好是
$T^{-1}$；所以 $B'$ 的 Gram 行列式生成 $d_{B/A}^{-1}$。另一方面，因
$B'=\beta B$，也可用 $\beta\alpha_1,\dots,\beta\alpha_n$ 作 $B'$ 的基。其 Gram 行列式为
\[
\det\bigl(\Tr(\beta\alpha_i\,\beta\alpha_j)\bigr)
=\Nm_{E/F}(\beta)^2\det T.
\]
因此在 $A$ 的分式理想群中有
\[
d_{B/A}^{-1}=(\Nm_{E/F}(\beta))^2d_{B/A}.
\]
把两边取赋值。若 $v$ 是 $A$ 的归一化赋值，则
\[
-v(d_{B/A})=2v(\Nm_{E/F}(\beta))+v(d_{B/A}),
\]
故
\[
v(d_{B/A})=-v(\Nm_{E/F}(\beta)).
\]
而 $D_{B/A}=\beta^{-1}B$，所以
\[
v(\Nm_{E/F}(D_{B/A}))=v(\Nm_{E/F}(\beta^{-1}))=-v(\Nm_{E/F}(\beta)).
\]
二者在每个局部化处赋值相同，故全局理想相等：
$d_{B/A}=\Nm_{E/F}(D_{B/A})$。这个证明的核心信息是：差别式先在上层环 $B$
中记录迹对偶格比整数环大多少，判别式再把这份指数信息用范数压回底层 $A$。
\end{proof}

\begin{proposition}
若 $E/F$、$K/E$ 为有限可分扩张，则
\[
d_{K/F}=\Nm_{E/F}(d_{K/E})\,d_{E/F}^{[K:E]}.
\]
若 $S\subset A$ 是乘性子集，则
\[
S^{-1}d_{B/A}=d_{S^{-1}B/S^{-1}A}.
\]
\end{proposition}

\begin{proof}
第一式由差别式传递性和上一条定理推出。设 $C$ 是 $A$ 在 $K$ 中的整闭包。差别式传递性给出
$D_{C/A}=D_{C/B}D_{B/A}C$，于是
\[
d_{K/F}
=\Nm_{K/F}(D_{K/F})
=\Nm_{K/F}(D_{K/E}D_{E/F}\mathcal O_K).
\]
理想范数对乘法和塔相容，因此第一因子为
\[
\Nm_{K/F}(D_{K/E})
=\Nm_{E/F}(\Nm_{K/E}(D_{K/E}))
=\Nm_{E/F}(d_{K/E}).
\]
第二因子来自从 $E$ 拉到 $K$ 的理想。若
$D_{E/F}=\prod\mathfrak p^{a_{\mathfrak p}}$，则
$D_{E/F}\mathcal O_K=\prod_{\mathfrak P|\mathfrak p}\mathfrak P^{a_{\mathfrak p}e_{\mathfrak P|\mathfrak p}}$，取范数后每个
$\mathfrak p$ 的指数变成
$a_{\mathfrak p}\sum_{\mathfrak P|\mathfrak p}e_{\mathfrak P|\mathfrak p}f_{\mathfrak P|\mathfrak p}
=a_{\mathfrak p}[K:E]$，所以
\[
\Nm_{K/F}(D_{E/F}\mathcal O_K)=d_{E/F}^{[K:E]}.
\]
局部化公式由差别式局部化和范数在局部化下的相容性推出：先把
$D_{B/A}$ 局部化，再取相对范；这与先取判别式理想 $d_{B/A}$ 再局部化得到同一个
$S^{-1}A$-理想。
\end{proof}

\begin{definition}
若 $C$ 是 $B$ 的子环且 $A\subset C\subset B$，则
\[
\mathfrak r=\{t\in C:tB\subset C\}
\]
称为 $B$ 在 $C$ 内的导子。导子衡量 $C$ 离整闭环 $B$ 有多远：$\mathfrak r=C$ 等价于 $C=B$。
\end{definition}

\begin{proposition}
设 $n=[E:F]$，$C=A[x]$，并且 $1,x,\dots,x^{n-1}$ 是 $C$ 的 $A$-基。设 $f$ 是 $x$ 在 $F$ 上的特征多项式，$f'$ 为其导数。
\begin{enumerate}[label=(\arabic*)]
\item $f$ 的系数属于 $A$。
\item
\[
1/f'(x),\;x/f'(x),\;\dots,\;x^{n-1}/f'(x)
\]
构成 $C'$ 的一组 $A$-基。
\item 若 $\mathfrak r$ 是 $B$ 在 $C$ 中的导子，则
\[
\mathfrak r=f'(x)D_{B/A}^{-1}.
\]
\item 若 $B=A[x]$，则 $D_{B/A}$ 是由 $f'(x)$ 生成的主理想。
\end{enumerate}
\end{proposition}

\begin{lemma}[Euler 迹公式]
在上一命题的记号下，设 $f(X)=\prod_{k=1}^n(X-x_k)$，其中 $x_1,\dots,x_n$ 是
$x$ 的共轭根。则
\[
\Tr_{E/F}\left(\frac{x^{n-1}}{f'(x)}\right)=1,\qquad
\Tr_{E/F}\left(\frac{x^i}{f'(x)}\right)=0\quad(0\le i\le n-2).
\]
\end{lemma}

\begin{proof}
在分裂域中作部分分式展开：
\[
\frac1{f(X)}
=\sum_{k=1}^n\frac1{f'(x_k)(X-x_k)}.
\]
当 $X$ 很大时，
\[
\frac1{X-x_k}=\frac1X\left(1+\frac{x_k}{X}+\frac{x_k^2}{X^2}+\cdots\right).
\]
比较 $X^{-m-1}$ 的系数，得到
\[
\sum_{k=1}^n\frac{x_k^m}{f'(x_k)}=0\quad(0\le m\le n-2),
\qquad
\sum_{k=1}^n\frac{x_k^{n-1}}{f'(x_k)}=1.
\]
左边正是相应元素的迹，因为迹等于所有共轭值之和。
\end{proof}

\begin{proof}
(1) 元素 $x$ 在 $B$ 中整，所以其特征多项式的系数在 $A$ 上整；这些系数又在 $F$ 中，因 $A$ 整闭而属于 $A$。

为证明 (2)，使用 Euler 迹公式。矩阵
\[
\left(\Tr\left(x^i\frac{x^j}{f'(x)}\right)\right)_{0\le i,j\le n-1}
\]
在 $i+j\le n-2$ 时条目为 $0$，在 $i+j=n-1$ 时条目为 $1$。当
$i+j\ge n$ 时，用特征方程
$x^n+a_{n-1}x^{n-1}+\cdots+a_0=0$ 把 $x^{i+j}$ 降到次数小于 $n$ 的
$A$-线性组合；由已知条目可逐次推出这些高次条目也属于 $A$。所以该矩阵是反对角线上为
$1$、反对角线一侧为 $0$ 的三角型矩阵，行列式为 $\pm1$。这说明列向量
$x^j/f'(x)$ 与基 $1,x,\dots,x^{n-1}$ 在迹配对下给出整数对偶基，从而生成 $C'$。

(3) 对 $t\in E$，逐步翻译条件：
\[
t\in\mathfrak r
\Longleftrightarrow
tB\subset C
\Longleftrightarrow
f'(x)^{-1}tB\subset C'
\Longleftrightarrow
\Tr(f'(x)^{-1}tB)\subset A.
\]
最后一个条件按余差别式定义等价于 $f'(x)^{-1}t\in D_{B/A}^{-1}$，即
$t\in f'(x)D_{B/A}^{-1}$。

(4) 若 $B=C$，则导子 $\mathfrak r=B$。由 (3) 得 $B=f'(x)D_{B/A}^{-1}$，因此 $D_{B/A}=(f'(x))$。
\end{proof}

\begin{theorem}
设 $F$ 为数域，$p$ 为普通素数。则 $p$ 在 $F$ 中分歧，当且仅当 $p$ 整除绝对判别式 $d_F$。
\end{theorem}

\begin{proof}
把上一节的分歧指数和本节的差别式联系起来看。局部化到 $p$ 上方的素理想后，未分歧等价于所有局部扩张的差别式指数为 $0$；有分歧等价于某个局部差别式指数为正。由
\[
d_F=\Nm_{F/\Q}(D_{\OF/\Z})
\]
可把这句话改写为素理想指数的判断：若
$D_{\OF/\Z}=\prod_{\mathfrak p}\mathfrak p^{a_{\mathfrak p}}$，则
\[
d_F=\prod_{\mathfrak p}\bigl(\mathfrak p\cap\Z\bigr)^{a_{\mathfrak p}f_{\mathfrak p|p}}.
\]
所以普通素数 $p$ 整除 $d_F$，等价于存在 $\mathfrak p|p$ 使
$a_{\mathfrak p}>0$。局部判别式理论说明 $a_{\mathfrak p}>0$ 正是局部扩张在
$\mathfrak p$ 处分歧；因此 $p$ 在 $F$ 中分歧与 $p\mid d_F$ 等价。
\end{proof}

\section{Hilbert 分歧理论}

前面的 $e$、$f$、$g$ 公式告诉我们素理想如何分裂。Hilbert 分歧理论进一步说明：当 $E/F$ 是 Galois 扩张时，分解和分歧都能由 Galois 群的子群来读出。这里的思想很重要：一个素理想 $\mathfrak P$ 在自同构作用下可能移动，也可能保持不动；保持不动的自同构构成分解群；在剩余域上还看不出来的自同构构成惯性群。

\subsection{分解群}

\begin{definition}
设 $E/F$ 为有限 Galois 扩张，$G=\Gal(E/F)$，$\mathfrak P$ 为 $\mathcal O_E$ 的素理想。定义
\[
D_{\mathfrak P}=\{\sigma\in G:\sigma\mathfrak P=\mathfrak P\},
\]
称为 $\mathfrak P$ 在 $E/F$ 中的分解群。它的固定域
\[
Z_{\mathfrak P}=E^{D_{\mathfrak P}}
\]
称为 $\mathfrak P$ 的分解域。
\end{definition}

\begin{proposition}
设 $E/F$ 为有限 Galois 扩张，$\mathfrak P\subset\mathcal O_E$，$\mathfrak p=\mathfrak P\cap\OF$。
\begin{enumerate}[label=(\arabic*)]
\item $\mathfrak Q$ 是 $\mathfrak p$ 上方素理想，当且仅当存在 $\sigma\in G$ 使 $\mathfrak Q=\sigma\mathfrak P$。
\item 对任意 $\sigma\in G$，有 $\sigma(\mathcal O_E)=\mathcal O_E$，并诱导同构
\[
\mathcal O_E/\mathfrak P\simeq\mathcal O_E/\sigma\mathfrak P,\qquad
[a]\longmapsto[\sigma a].
\]
因此 $f_{\mathfrak P|\mathfrak p}=f_{\sigma\mathfrak P|\mathfrak p}$。
\item 对任意 $\sigma\in G$，有 $\sigma(\mathfrak p\mathcal O_E)=\mathfrak p\mathcal O_E$，因此
$e_{\mathfrak P|\mathfrak p}=e_{\sigma\mathfrak P|\mathfrak p}$。
\end{enumerate}
\end{proposition}

\begin{proof}
(1) 先说明为什么 Galois 性会给出传递作用。设 $\mathfrak Q$ 也是 $\mathfrak p$ 上方的素理想。局部化到
$A=\mathcal O_{F,\mathfrak p}$，并把 $\mathcal O_E$ 局部化成
$B=S^{-1}\mathcal O_E$，其中 $S=\OF\setminus\mathfrak p$。此时 $B$ 是
$A$ 在 $E$ 中的整闭包，$B$ 的极大理想正好对应于 $\mathfrak p$ 上方的素理想。因为
$E/F$ 是有限 Galois 扩张，$G$ 作用在 $B$ 上，且 $A=B^G$。整扩张中躺在同一个
$A$ 的素理想上方的极大理想，在有限 Galois 群作用下属于同一个轨道；否则把一个轨道中所有极大理想相交，会得到一个 $G$-稳定但不能由 $A$ 中理想拉回的理想，和
$A=B^G$ 以及 lying-over/going-up 的分解性质矛盾。因此存在 $\sigma\in G$ 使
$\sigma\mathfrak P=\mathfrak Q$。反向命题直接成立：若
$\mathfrak Q=\sigma\mathfrak P$，则
$\mathfrak Q\cap\OF=\sigma(\mathfrak P\cap\OF)=\mathfrak p$。

(2) 若 $x\in\mathcal O_E$，则 $x$ 在 $\OF$ 上整，$\sigma x$ 仍满足把系数经 $\sigma$ 固定后的同一个整方程，所以 $\sigma x\in\mathcal O_E$。映射 $[a]\mapsto[\sigma a]$ 良定，因为 $a-b\in\mathfrak P$ 当且仅当 $\sigma a-\sigma b\in\sigma\mathfrak P$。逆映射由 $\sigma^{-1}$ 给出。

(3) 因为 $\sigma$ 固定 $\OF$，故固定 $\mathfrak p$ 的每个元素，从而
\[
\sigma(\mathfrak p\mathcal O_E)=\mathfrak p\mathcal O_E.
\]
把 $\mathfrak p\mathcal O_E$ 的唯一素理想分解施以 $\sigma$，指数保持不变。
\end{proof}

因此若
\[
\mathfrak p\mathcal O_E=\prod_{\mathfrak P_i\mid\mathfrak p}\mathfrak P_i^e
\]
是在 Galois 扩张中的分解式，则所有上方素理想有同一个分歧指数 $e$ 和同一个惯性次数 $f$。上方素理想个数等于 $[G:D_{\mathfrak P}]$，于是
\[
ef[G:D_{\mathfrak P}]=[E:F].
\]
特别地，$\mathfrak p$ 完全分裂当且仅当 $D_{\mathfrak P}=\{1\}$。

\subsection{惯性群}

若 $\sigma\in D_{\mathfrak P}$，它固定理想 $\mathfrak P$，所以诱导剩余域自同构
\[
\overline\sigma:\kappa_{\mathfrak P}\longrightarrow\kappa_{\mathfrak P},
\qquad
[a]\longmapsto[\sigma a],
\]
其中 $\kappa_{\mathfrak P}=\mathcal O_E/\mathfrak P$，$\kappa_{\mathfrak p}=\OF/\mathfrak p$。

\begin{proposition}
扩张 $\kappa_{\mathfrak P}/\kappa_{\mathfrak p}$ 是有限 Galois 扩张，并且
\[
D_{\mathfrak P}\longrightarrow\Gal(\kappa_{\mathfrak P}/\kappa_{\mathfrak p}),
\qquad
\sigma\longmapsto\overline\sigma
\]
是满同态。
\end{proposition}

\begin{proof}
有限性来自 $\mathcal O_E/\mathfrak P$ 是 $\OF/\mathfrak p$ 上有限维向量空间。这个扩张是
Galois 的原因如下。取 $\bar x\in\kappa_{\mathfrak P}$，选取提升
$x\in\mathcal O_E$。$x$ 在 $F$ 上的共轭为 $\sigma x$，其中 $\sigma\in G$；若
$\sigma\notin D_{\mathfrak P}$，则 $\sigma$ 把 $\mathfrak P$ 移到别的上方素理想，不给
$\kappa_{\mathfrak P}$ 的自同构。真正留下来的共轭恰由 $D_{\mathfrak P}$ 给出，它们在模
$\mathfrak P$ 后给出 $\bar x$ 在 $\kappa_{\mathfrak p}$ 上的共轭。因此剩余域扩张是正规且可分的有限域扩张。

还要说明映射满。这里使用的局部事实是：对整闭局部基环上的有限 Galois 扩张，保持给定上方极大理想的自同构群
$D_{\mathfrak P}$ 到剩余域自同构群的约化映射是满射。可把证明想成以下三步。第一，先选取
$\kappa_{\mathfrak P}/\kappa_{\mathfrak p}$ 的一个生成元 $\bar a$，并取提升
$a\in\mathcal O_E$。第二，$\bar a$ 的任意剩余域共轭来自 $a$ 的某个
$F$-共轭在模 $\mathfrak P$ 后的像。第三，若这个共轭最初把 $\mathfrak P$ 送到别的上方素理想，就用上一个命题的传递性在左边复合一个把该素理想送回
$\mathfrak P$ 的自同构。复合后得到的元素属于 $D_{\mathfrak P}$，并在剩余域上给出预定的自同构。
\end{proof}

\begin{definition}
满同态的核
\[
I_{\mathfrak P}=\{\sigma\in D_{\mathfrak P}:\overline\sigma=1\}
\]
称为 $\mathfrak P$ 在 $E/F$ 中的惯性群。其固定域
\[
T_{\mathfrak P}=E^{I_{\mathfrak P}}
\]
称为惯性域。
\end{definition}

\begin{proposition}
设 $E/F$ 为有限 Galois 数域扩张，$\mathfrak P\subset\mathcal O_E$，$\mathfrak p=\mathfrak P\cap\OF$。
\begin{enumerate}[label=(\arabic*)]
\item 有正合列
\[
1\longrightarrow I_{\mathfrak P}\longrightarrow
D_{\mathfrak P}\longrightarrow
\Gal(\kappa_{\mathfrak P}/\kappa_{\mathfrak p})
\longrightarrow1.
\]
\item $|I_{\mathfrak P}|=e_{\mathfrak P|\mathfrak p}$，且 $[D_{\mathfrak P}:I_{\mathfrak P}]=f_{\mathfrak P|\mathfrak p}$。
\item 若 $\mathfrak P_T=\mathfrak P\cap\mathcal O_{T_{\mathfrak P}}$，则
\[
e_{\mathfrak P|\mathfrak P_T}=e_{\mathfrak P|\mathfrak p},\qquad
f_{\mathfrak P|\mathfrak P_T}=1.
\]
\item 若 $\mathfrak P_Z=\mathfrak P\cap\mathcal O_{Z_{\mathfrak P}}$，则
\[
e_{\mathfrak P_T|\mathfrak P_Z}=1,\qquad
f_{\mathfrak P_T|\mathfrak P_Z}=f_{\mathfrak P|\mathfrak p}.
\]
\end{enumerate}
\end{proposition}

\begin{proof}
(1) 是上一命题的核像描述。

(2) 满射给出 $[D_{\mathfrak P}:I_{\mathfrak P}]
=|\Gal(\kappa_{\mathfrak P}/\kappa_{\mathfrak p})|=f_{\mathfrak P|\mathfrak p}$。又由分解群公式
\[
e_{\mathfrak P|\mathfrak p}f_{\mathfrak P|\mathfrak p}
|G:D_{\mathfrak P}|=|G|
\]
和上方素理想个数 $|G:D_{\mathfrak P}|$，得到 $|D_{\mathfrak P}|=ef$，所以 $|I_{\mathfrak P}|=e$。

(3) 扩张 $E/T_{\mathfrak P}$ 的 Galois 群是 $I_{\mathfrak P}$。这一层只有一个上方素理想，因为
$I_{\mathfrak P}$ 是 $\mathfrak P$ 的稳定子在子扩张中的全群。惯性群在剩余域上作用平凡，所以
$f_{\mathfrak P|\mathfrak P_T}=1$。局部次数公式给出
\[
[E:T_{\mathfrak P}]=e_{\mathfrak P|\mathfrak P_T}f_{\mathfrak P|\mathfrak P_T}
=|I_{\mathfrak P}|=e_{\mathfrak P|\mathfrak p},
\]
于是 $e_{\mathfrak P|\mathfrak P_T}=e_{\mathfrak P|\mathfrak p}$。

(4) 扩张 $T_{\mathfrak P}/Z_{\mathfrak P}$ 的 Galois 群为
$D_{\mathfrak P}/I_{\mathfrak P}$，它同构于
$\Gal(\kappa_{\mathfrak P}/\kappa_{\mathfrak p})$，阶为
$f_{\mathfrak P|\mathfrak p}$。这一层的群已经没有惯性核，所以它在
$\mathfrak P_Z$ 处无分歧，即 $e_{\mathfrak P_T|\mathfrak P_Z}=1$；局部次数等于
$f_{\mathfrak P|\mathfrak p}$，故剩余次数也是 $f_{\mathfrak P|\mathfrak p}$。讲解时应把三层关系记成：分解域先把“有几个上方素理想”固定成一个；惯性域再把“真正的分歧”从剩余域扩张中分离出来；惯性群的阶就是分歧指数。
\end{proof}

\subsection{Frobenius 自同构}

设 $\mathfrak p$ 在 $E/F$ 中无分歧。有限域扩张 $\kappa_{\mathfrak P}/\kappa_{\mathfrak p}$ 的 Galois 群由 Frobenius 生成：
\[
[x]\longmapsto [x]^q,\qquad q=|\kappa_{\mathfrak p}|.
\]
因为无分歧时 $I_{\mathfrak P}=1$，分解群 $D_{\mathfrak P}$ 与剩余域 Galois 群同构，于是可以把这个 Frobenius 唯一提升到 $D_{\mathfrak P}$。

\begin{definition}
无分歧情形下，唯一满足
\[
\left[\frac{E/F}{\mathfrak P}\right](x)\equiv x^{\Nm(\mathfrak p)}
\pmod{\mathfrak P},
\qquad x\in\mathcal O_E
\]
的元素称为 $\mathfrak P$ 处的 Frobenius 自同构。
\end{definition}

\begin{proposition}
若 $\tau\in G$，则
\[
\left[\frac{E/F}{\tau\mathfrak P}\right]
=\tau\left[\frac{E/F}{\mathfrak P}\right]\tau^{-1}.
\]
\end{proposition}

\begin{proof}
令 $\sigma=\left[\frac{E/F}{\mathfrak P}\right]$，并取任意
$y\in\mathcal O_E$。写 $y=\tau x$，其中 $x=\tau^{-1}y\in\mathcal O_E$。由
Frobenius 的定义，
$\sigma x\equiv x^q\pmod{\mathfrak P}$。对这个同余施以 $\tau$，得到
\[
\tau\sigma\tau^{-1}(y)\equiv y^q\pmod{\tau\mathfrak P}.
\]
这说明 $\tau\sigma\tau^{-1}$ 满足 $\tau\mathfrak P$ 处 Frobenius 的刻画条件。无分歧时分解群到剩余域 Galois 群是同构，所以满足该同余条件的元素唯一，命题成立。
\end{proof}

因此在非交换 Galois 群中，Frobenius 由共轭类决定；在交换扩张中，共轭不改变元素，所以它只依赖于 $\mathfrak p$，记为
\[
\left(\frac{E/F}{\mathfrak p}\right),
\]
也称 Artin 符号。

\begin{proposition}
设 $\mathfrak p$ 在有限 Galois 扩张 $E/F$ 中无分歧，$\mathfrak P|\mathfrak p$。则 Frobenius 自同构的阶为 $f_{\mathfrak P|\mathfrak p}$，且 $\mathfrak p$ 在 $E$ 中分解为 $[E:F]/f_{\mathfrak P|\mathfrak p}$ 个素理想。特别地，$\mathfrak p$ 完全分裂当且仅当 Frobenius 为恒等元。
\end{proposition}

\begin{proof}
Frobenius 在剩余域扩张 $\kappa_{\mathfrak P}/\kappa_{\mathfrak p}$ 中的阶就是该有限域扩张的次数 $f_{\mathfrak P|\mathfrak p}$。无分歧时 $e=1$，所以公式
\[
efg=[E:F]
\]
给出 $g=[E:F]/f$。完全分裂等价于 $f=1$ 且 $e=1$，也等价于 Frobenius 的阶为 $1$。
\end{proof}

现在设 $E/F$ 为有限交换 Galois 扩张，$S$ 为包含所有分歧素理想的有限集合。由所有 $\mathfrak p\notin S$ 生成的分式理想群子群记为 $I_F^S$。定义
\[
\left(\frac{E/F}{\prod_{\mathfrak p\notin S}\mathfrak p^{n(\mathfrak p)}}\right)
=\prod_{\mathfrak p\notin S}
\left(\frac{E/F}{\mathfrak p}\right)^{n(\mathfrak p)}.
\]
这就是交换扩张的 Artin 同态
\[
\left(\frac{E/F}{\cdot}\right):I_F^S\longrightarrow\Gal(E/F).
\]
它是类域论的核心桥梁：左边是理想，右边是 Galois 群。

\section{理想类群}

理想类群回答的问题很朴素：Dedekind 环中每个理想都能唯一分解为素理想，但它们是否都由一个元素生成？如果不是，失败程度有多大？把所有分式理想组成群，再把主理想看成“平凡”的，剩下的商群就是理想类群。

\begin{definition}
设 $F$ 为数域，$\OF$ 为整数环。非零有限生成 $\OF$-子模 $\mathfrak f\subset F$ 称为分式理想。所有分式理想构成乘法群 $I_F$。主分式理想组成子群
\[
P_F=\{a\OF:a\in F^\times\}.
\]
商群
\[
\Cl_F=I_F/P_F
\]
称为 $F$ 的理想类群，阶 $h_F=|\Cl_F|$ 称为类数。
\end{definition}

\begin{remark}
$h_F=1$ 的意义是：每个理想都是主理想。对 Dedekind 环而言，这等价于 $\OF$ 是主理想环，从而元素唯一分解成立。若 $h_F>1$，元素唯一分解失败，但理想唯一分解仍保留下来；理想类群正是把这种失败压缩成一个群。
\end{remark}

类数是数域的核心不变量。例如 Gauss 的类数问题指出，虚二次域 $\Q(\sqrt{-d})$ 的类数为 $1$ 的正整数 $d$ 恰为
\[
1,2,3,7,11,19,43,67,163.
\]
这不是孤立趣闻，而是说明类数把二次型、素数表示和数域算术联系在一起。

为了走向类域论，需要把普通理想类群推广为射线类群。这里有两个新成分：有限素理想给出通常的同余条件，无穷实嵌入给出正负号条件。

\begin{definition}
设 $\iota_1,\dots,\iota_{r_1}:F\hookrightarrow\R$ 是全部实嵌入。一个模数写成
\[
\mathfrak m=\mathfrak m_\infty\mathfrak m_0,
\]
其中
\[
\mathfrak m_\infty=(\mathfrak p_{\infty,1})^{\nu_1}\cdots
(\mathfrak p_{\infty,r_1})^{\nu_{r_1}},\qquad \nu_j\in\{0,1\},
\]
$\mathfrak m_0$ 是 $\OF$ 的非零理想。若
\[
\mathfrak m_0=\prod_i\mathfrak p_i^{\nu_i},
\]
则 $x\in F^\times$ 满足
\[
x\equiv1\pmod{\mathfrak m}
\]
是指同时满足
\[
\nu_j=1\Rightarrow \iota_j(x)>0,\qquad
v_{\mathfrak p_i}(x-1)\ge\nu_i.
\]
记
\[
F_{\mathfrak m,1}^{\times}
=\{x\in F^\times:x\equiv1\pmod{\mathfrak m}\}.
\]
\end{definition}

\begin{definition}
在上述模数 $\mathfrak m=\mathfrak m_\infty\mathfrak m_0$ 下，元素
$x\in F^\times$ 称为与 $\mathfrak m$ 互素，是指主分式理想 $(x)$ 中不出现任何
$\mathfrak m_0$ 的素因子，即对每个 $\mathfrak p|\mathfrak m_0$ 都有
$v_{\mathfrak p}(x)=0$。记
\[
F_{\mathfrak m}^{\times}
=\{x\in F^\times:x\text{ 与 }\mathfrak m\text{ 互素}\}.
\]
若把 $\OF$ 在所有不整除 $\mathfrak m_0$ 的元素处局部化，记所得环为
$\mathcal O(\mathfrak m)$，则
\[
F_{\mathfrak m}^{\times}=\mathcal O(\mathfrak m)\setminus\{0\}.
\]
这一等式的意思是：分母只允许含有不属于 $\mathfrak m_0$ 的素因子，分子也不能额外带有
$\mathfrak m_0$ 的素因子；这样主理想 $(x)$ 才能落在后面定义的
$I_F^{\mathfrak m}$ 中。
\end{definition}

\begin{definition}
令 $I_F^{\mathfrak m}$ 为由所有不整除 $\mathfrak m_0$ 的素理想生成的分式理想群。定义
\[
S_F^{\mathfrak m}
=\{x\OF:x\in F_{\mathfrak m,1}^{\times}\}\subset I_F^{\mathfrak m}.
\]
商群
\[
I_F^{\mathfrak m}/S_F^{\mathfrak m}
\]
称为模 $\mathfrak m$ 的射线理想类群。
\end{definition}

\begin{remark}
普通理想类群只把任意主理想视为平凡；射线类群只把满足指定同余条件的主理想视为平凡。因此模数越大，被杀掉的主理想越少，射线类群保留的信息越细。
\end{remark}

\begin{definition}
若 $H$ 是分式理想群 $I_F$ 的子群，且
\[
S_F^{\mathfrak m}\subset H\subset I_F^{\mathfrak m},
\]
则称 $H$ 为模 $\mathfrak m$ 定义的理想群。若只使用有限部分
$\mathfrak m_0$ 的同余条件，即
\[
S_{F,\mathrm{fin}}^{\mathfrak m}
=\{x\OF:x\equiv1\pmod{\mathfrak m_0}\},
\]
则商群 $I_F^{\mathfrak m}/S_{F,\mathrm{fin}}^{\mathfrak m}$ 称为模
$\mathfrak m$ 的小射线类群。小射线类群只记录有限素理想同余条件；完整射线类群同时保留有限素理想同余和实嵌入符号。
\end{definition}

\begin{definition}
设 $E/F$ 为有限扩张。令 $I_E^{\mathfrak m}$ 为由所有满足
$\mathfrak P\cap\OF\in I_F^{\mathfrak m}$ 的 $\mathcal O_E$ 素理想 $\mathfrak P$
生成的分式理想群子群。相对范数把它送到 $I_F^{\mathfrak m}$：
\[
\Nm_{E/F}(I_E^{\mathfrak m})\subset I_F^{\mathfrak m}.
\]
这里限制 $\mathfrak P\cap\OF\in I_F^{\mathfrak m}$ 是为了排除模数有限部分中的素理想，使 Artin 符号只在未分歧、允许取 Frobenius 的素理想上定义。
\end{definition}

\begin{theorem}
设 $E/F$ 为有限交换数域扩张。存在 $F$ 的模数 $\mathfrak m$，使得有交换群同构
\[
I_F^{\mathfrak m}/\Nm_{E/F}(I_E^{\mathfrak m})S_F^{\mathfrak m}
\simeq
\Gal(E/F).
\]
\end{theorem}

\begin{remark}
这一定理是后续上同调方法要证明的目标之一。它把 Galois 群完全改写成理想类群的商。讲到这里时不必急着证明，而要抓住方向：类域论说交换扩张可以由底域上的理想同余条件来分类。
\end{remark}

\section{Picard 群}

本节把理想类群放进代数几何语言。若没有学过概形，可以先抓住一句话：可逆模是“代数上的线丛”，Picard 群分类这些线丛；在 Dedekind 环上，可逆模和可逆分式理想是一回事，所以 Picard 群就是理想类群。

先回顾局部化。若 $R$ 是交换环，$S\subset R$ 为乘性子集，则局部化为
$S^{-1}R$。若 $\mathfrak p$ 为素理想，取 $S=R\setminus\mathfrak p$，记
$R_{\mathfrak p}=S^{-1}R$。若 $s\in R$，取
$S=\{1,s,s^2,\dots\}$，则记 $S^{-1}R=R[1/s]$。若 $M$ 是 $R$-模，记
$M_{\mathfrak p}=R_{\mathfrak p}\otimes_RM$。

本节还需要一点仿射概形语言。$\Spec R$ 是 $R$ 的全部素理想组成的空间。对
$S\subset R$，闭集
\[
V(S)=\{\mathfrak p\in\Spec R:S\subset\mathfrak p\},
\]
其补集记为 $D(S)$；特别地，$D(s)$ 是所有不含 $s$ 的素理想。若
$e^2=e$ 是幂等元，则 $D(e)$ 同时开且闭；反过来，$\Spec R$ 的开闭分块由幂等元控制。整环只有
$0,1$ 两个幂等元，所以整环的 $\Spec R$ 连通。

有限生成投射模 $M$ 在每个局部环 $R_{\mathfrak p}$ 上都是自由模，因此可定义局部秩
$\operatorname{rank}(M_{\mathfrak p})$。当 $R$ 是整环时，$\Spec R$ 连通，投射模的秩函数局部常数，所以有限生成投射模的秩在所有素理想处相同。这个事实会在判断“常秩
$1$”时反复使用。

\subsection{可逆模}

\begin{proposition}
设 $R$ 为交换环，$M$ 为 $R$-模，记 $M^*=\Hom_R(M,R)$。以下条件等价。
\begin{enumerate}[label=(\arabic*)]
\item $M$ 是常秩 $1$ 的有限生成投射 $R$-模。
\item 取值映射给出同构
\[
M^*\otimes_R M\simeq R.
\]
\item 存在 $R$-模 $N$，使得
\[
N\otimes_RM\simeq R.
\]
\end{enumerate}
\end{proposition}

\begin{proof}
(1) 推 (2)。考虑取值映射
\[
M^*\otimes_RM\longrightarrow R,\qquad
\alpha\otimes m\longmapsto\alpha(m).
\]
对每个素理想 $\mathfrak p$ 局部化。因为 $M_{\mathfrak p}$ 是秩 $1$ 的自由 $R_{\mathfrak p}$-模，局部取值映射就是
\[
R_{\mathfrak p}\otimes_{R_{\mathfrak p}}R_{\mathfrak p}\to R_{\mathfrak p},
\]
因而是同构。更具体地，若 $e$ 是 $M_{\mathfrak p}$ 的一组基，$e^*$ 是对偶基，则
$e^*\otimes e$ 被送到 $1$，任意 $\alpha\otimes m$ 都是它的倍数。映射的核和余核局部化后为
$0$，而核、余核都是有限生成模，因此它们本身为 $0$；原取值映射为同构。

(2) 推 (3) 可取 $N=M^*$。

(3) 推 (1)。设 $\psi:N\otimes_RM\simeq R$。把逆同构在交换因子后记为
$\varphi:R\to M\otimes_RN$，并写
\[
\varphi(1)=\sum_i m_i\otimes n_i.
\]
定义 $\beta_i:M\to R$ 为 $\beta_i(m)=\psi(n_i\otimes m)$。复合
\[
\alpha:M\simeq R\otimes_RM
\xrightarrow{\varphi\otimes1}
M\otimes_RN\otimes_RM
\xrightarrow{1\otimes\psi}
M\otimes_RR\simeq M
\]
是同构，并且
\[
\alpha(m)=\sum_i\beta_i(m)m_i.
\]
令 $m_i'=\alpha^{-1}(m_i)$，则
\[
m=\sum_i\beta_i(m)m_i'.
\]
这说明有限多个 $m_i'$ 生成 $M$。同时映射
\[
M\longrightarrow R^r,\qquad m\longmapsto(\beta_1(m),\dots,\beta_r(m)),
\]
和
\[
R^r\longrightarrow M,\qquad (a_1,\dots,a_r)\longmapsto\sum_i a_im_i'
\]
的复合为 $\operatorname{id}_M$，故 $M$ 是自由模 $R^r$ 的直和因子，即有限生成投射模。

局部化后 $N_{\mathfrak p}\otimes M_{\mathfrak p}\simeq R_{\mathfrak p}$。在局部环上，有限生成投射模都是自由模；若
$M_{\mathfrak p}$、$N_{\mathfrak p}$ 的秩分别为 $r_{\mathfrak p}$、$s_{\mathfrak p}$，则
$r_{\mathfrak p}s_{\mathfrak p}=1$，所以 $r_{\mathfrak p}=1$。因此 $M$ 常秩
$1$。
\end{proof}

\begin{definition}
满足上一命题条件的 $R$-模称为可逆模。可逆模的同构类在
\[
\langle M\rangle\langle N\rangle=\langle M\otimes_RN\rangle
\]
下构成交换群，称为 $R$ 的 Picard 群，记为 $\Pic(R)$。单位元是 $\langle R\rangle$，逆元是 $\langle M^*\rangle$。
\end{definition}

\subsection{分式理想与投射模}

\begin{proposition}
设 $R$ 为整环，$F$ 为其分式域，$M$ 为有限生成投射 $R$-模。则 $M$ 为可逆模，当且仅当 $M$ 同构于 $F$ 的某个非零 $R$-子模。
\end{proposition}

\begin{proof}
若 $M$ 可逆，则 $F\otimes_RM$ 是一维 $F$-向量空间。取一个
$F$-线性同构 $\varphi:F\otimes_RM\simeq F$，复合自然映射
$M\to F\otimes_RM$ 得到 $R$-线性映射 $M\to F$。这个映射是单射：若
$m$ 被送到 $0$，则 $1\otimes m=0$；因 $M$ 是投射模，故无挠，而整环上无挠模嵌入其泛点局部化，所以 $m=0$。于是 $M$ 同构于
$F$ 的非零 $R$-子模。

反过来，若 $M$ 同构于 $F$ 的某个非零 $R$-子模，则可把 $M$ 看成
$M\subset F$。取 $0\ne m_0\in M$。映射
\[
F\longrightarrow F\otimes_RM,\qquad u/v\longmapsto (u/v)\otimes m_0
\]
非零，且 $F\otimes_RM$ 由 $M$ 的任一非零元素在 $F$ 上生成；因此它是一维
$F$-向量空间。所以 $M$ 在泛点处秩为 $1$。投射模的秩函数在连通分支上局部常数；整环的谱连通，因此 $M$ 常秩 $1$，即为可逆模。
\end{proof}

\begin{definition}
设 $R$ 为整环，$F$ 为分式域。$F$ 的非零 $R$-子模 $M$ 称为分式理想，是指存在非零 $a,b\in F$ 使
\[
aR\subset M\subset bR.
\]
等价地，存在 $R$ 的非零理想 $\mathfrak a$ 和非零元素 $c\in F$ 使
$M=c\mathfrak a$。从第一个条件到第二个条件，可取 $c=b$，则
$b^{-1}M$ 是 $R$ 的非零理想；反向由非零理想 $\mathfrak a$ 有
$0\ne dR\subset\mathfrak a\subset R$ 得到。
对这样的 $M$，记
\[
\overline M=\{x\in F:xM\subset R\}.
\]
\end{definition}

\begin{proposition}
设 $R$ 为整环，$F$ 为其分式域，$M\subset F$ 为非零 $R$-子模。以下条件等价。
\begin{enumerate}[label=(\arabic*)]
\item $M$ 是分式理想，且 $M\overline M=R$。
\item $M$ 是有限生成投射 $R$-模。
\end{enumerate}
在这些条件下，乘法映射
\[
M\otimes_RN\longrightarrow MN
\]
对同类分式理想为同构，并且 $\overline M\simeq M^*$。
\end{proposition}

\begin{proof}
(1) 推 (2)。由 $M\overline M=R$，取
\[
1=m_1n_1+\cdots+m_rn_r,\qquad m_i\in M,\ n_i\in\overline M.
\]
定义
\[
M\to R^r,\quad m\mapsto(mn_1,\dots,mn_r),
\]
以及
\[
R^r\to M,\quad (x_1,\dots,x_r)\mapsto\sum_i x_im_i.
\]
复合为恒等映射，所以 $M$ 是有限生成自由模的直和因子，即有限生成投射模。

(2) 推 (1)。有限生成性给出 $M\subset bR$：若
$u_1,\dots,u_r$ 生成 $M$，把每个 $u_i$ 写成 $a_i/s_i$，取
$b=(s_1\cdots s_r)^{-1}$ 的适当 $R$-倍数即可使所有 $u_i\in bR$。另一方面，取
$0\ne a\in M$ 得 $aR\subset M$，所以 $M$ 是分式理想。

投射性给出有限对偶基：存在 $m_i\in M$、$f_i\in M^*$，使
\[
m=\sum_i f_i(m)m_i.
\]
固定 $0\ne m_0\in M$，把 $M\subset F$ 后可写 $f_i(m)=n_im$，其中 $n_i=f_i(m_0)/m_0\in F$。由 $f_i(M)\subset R$ 得 $n_i\in\overline M$。代入 $m=m_0$ 并除以 $m_0$ 得
\[
1=\sum_i n_im_i,
\]
所以 $M\overline M=R$。

最后，$x\in\overline M$ 给出 $M\to R$，$m\mapsto xm$，从而
$\overline M\to M^*$。上面的对偶基论证也说明每个 $R$-线性映射都来自某个这样的
$x$，所以 $\overline M\simeq M^*$。乘法映射
$M\otimes_RN\to MN$ 在每个素理想处局部化后变为
$R_{\mathfrak p}$ 上两个自由秩一模的乘法映射；选取局部生成元后它就是
$R_{\mathfrak p}\otimes R_{\mathfrak p}\to R_{\mathfrak p}$。局部化处处为同构，故全局为同构。
\end{proof}

\subsection{Dedekind 环的 Picard 群}

\begin{proposition}
设 $\mathcal O$ 为 Dedekind 环，$F$ 为其分式域。则理想类群 $\Cl_F$ 与 Picard 群 $\Pic(\mathcal O)$ 自然同构。
\end{proposition}

\begin{proof}
先把 Picard 群的元素放进分式理想语言。由命题 1.34，每个可逆
$\mathcal O$-模同构于 $F$ 的非零 $\mathcal O$-子模；由命题 1.35，这样的子模正是可逆分式理想。因此每个 Picard 类可用某个非零分式理想表示。定义
\[
\Pic(\mathcal O)\longrightarrow\Cl_F,\qquad
\langle I\rangle\longmapsto[I].
\]
若 $I\simeq J$ 作为 $\mathcal O$-模，则任一同构在张量到 $F$ 后是一维 $F$-向量空间的自同构，故由某个 $a\in F^\times$ 乘法给出，于是 $I=aJ$，二者在类群中相同。映射良定。

它是满射，因为每个理想类有分式理想代表。它是单射，因为若 $[I]$ 为平凡类，则 $I=a\mathcal O$，作为 $\mathcal O$-模同构于 $\mathcal O$，所以 $\langle I\rangle$ 是 $\Pic(\mathcal O)$ 的单位元。乘法相容来自
\[
\langle I\rangle\langle J\rangle
=\langle I\otimes_{\mathcal O}J\rangle
=\langle IJ\rangle.
\]
\end{proof}

\section{Grothendieck 群}

Grothendieck 群把“直和分解”变成加法关系。对初学者来说，可以把 $K_0(R)$ 想成投射模的记账本：一个有限生成投射模 $P$ 记作 $[P]$，直和 $P\oplus Q$ 记作加法 $[P]+[Q]$。本节的目的，是把 $K_0$ 与 Picard 群、理想类群接起来。

\subsection{定义与稳定同构}

\begin{definition}
设 $R$ 为环。令 $\mathcal F$ 为由所有有限生成投射 $R$-模同构类 $\langle P\rangle$ 生成的自由交换群。取关系
\[
\langle P'\rangle+\langle P''\rangle-\langle P'\oplus P''\rangle.
\]
商群称为 $R$ 的 Grothendieck 群，记为 $K_0(R)$。$\langle P\rangle$ 在商群中的像记为 $[P]$，于是
\[
[P\oplus Q]=[P]+[Q].
\]
任意元素都可写成 $[P]-[Q]$ 的形式：因为自由交换群中的元素是有限和
$\sum_i\langle P_i\rangle-\sum_j\langle Q_j\rangle$，到商群后它等于
$[\bigoplus_iP_i]-[\bigoplus_jQ_j]$。这个写法会让“相等”变成稳定同构问题。
\end{definition}

\begin{proposition}
若 $P,Q$ 是有限生成投射 $R$-模，则
\[
[P]=[Q]\quad\text{在 }K_0(R)\text{ 中}
\]
当且仅当存在整数 $n\ge0$，使
\[
P\oplus R^n\simeq Q\oplus R^n.
\]
\end{proposition}

\begin{proof}
若存在这样的稳定同构，则在 $K_0(R)$ 中
\[
[P]+n[R]=[Q]+n[R],
\]
故 $[P]=[Q]$。

反过来，若 $[P]=[Q]$，则在自由交换群中，$\langle P\rangle-\langle Q\rangle$ 是有限个直和关系的整数线性组合：
\[
\langle P\rangle-\langle Q\rangle
=\sum_i(\langle P_i\rangle+\langle Q_i\rangle-\langle P_i\oplus Q_i\rangle)
-\sum_j(\langle P'_j\rangle+\langle Q'_j\rangle-\langle P'_j\oplus Q'_j\rangle).
\]
把负项移到等式另一边，并使用自由交换群的基是同构类这一事实，可得某个有限生成投射模
$X$ 使
\[
P\oplus X\simeq Q\oplus X.
\]
因为 $X$ 投射且有限生成，存在 $Y$ 与 $n$ 使 $X\oplus Y\simeq R^n$。两边再直和 $Y$，得到
\[
P\oplus R^n\simeq Q\oplus R^n.
\]
\end{proof}

\subsection{秩映射与行列式}

若 $R$ 是交换环，则张量积使 $K_0(R)$ 成为交换环：
\[
[P][Q]=[P\otimes_RQ].
\]
对有限生成投射模 $M$，每个局部化 $M_{\mathfrak p}$ 是自由 $R_{\mathfrak p}$-模，定义秩函数
\[
\rho_M:\Spec R\to\Z,\qquad
\mathfrak p\longmapsto\operatorname{rank}_{R_{\mathfrak p}}M_{\mathfrak p}.
\]

\begin{proposition}
设 $R$ 为交换环。
\begin{enumerate}[label=(\arabic*)]
\item $\rho_M$ 是连续函数，其中 $\Z$ 取离散拓扑。
\item
\[
\rho:K_0(R)\to C(\Spec R,\Z),\qquad [M]\mapsto\rho_M
\]
是环同态。
\item 存在环同态
\[
\theta:C(\Spec R,\Z)\to K_0(R)
\]
使 $\rho\circ\theta=1$。
\item 若 $j:\Ker\rho\hookrightarrow K_0(R)$ 为包含映射，定义
\[
p:K_0(R)\to\Ker\rho,\qquad p(x)=x-\theta(\rho(x)),
\]
则 $p\circ j=1$。
\item 因而
\[
K_0(R)\simeq C(\Spec R,\Z)\oplus\Ker\rho.
\]
\end{enumerate}
\end{proposition}

\begin{proof}
(1) 因 $M$ 投射有限生成，存在 $N$ 与 $n$ 使 $M\oplus N\simeq R^n$。集合
\[
\{\mathfrak p:\operatorname{rank}M_{\mathfrak p}\le m\}
\]
可用外幂 $\bigwedge^{m+1}M$ 在 $\mathfrak p$ 处是否为零描述；有限生成模局部化为零的点集是开集。因此 $\rho_M^{-1}(m)$ 可写成有限个开集的交，故为开集。

(2) 直和对应秩相加，张量积对应秩相乘，所以 $\rho$ 是环同态。

(3) 连续函数 $f:\Spec R\to\Z$ 的像有限。每个纤维 $f^{-1}(n_i)$ 是开闭集，对应某个幂等元 $e_i$，并且这些幂等元两两正交、和为 $1$。定义
\[
\theta(f)=\sum_i n_i[Re_i].
\]
若选择了更细的开闭分解，$Re_i$ 相应分解为直和，所以 $K_0$ 中元素不变。加法和乘法由共同加细分解检查。

(4) 先看 $f$ 在开闭分解 $D(e_i)$ 上取常值 $n_i$ 的情形。局部化后，
$Re_i$ 在 $D(e_i)$ 上秩为 $1$，在其他 $D(e_j)$ 上秩为 $0$，所以
$\rho(\sum_i n_i[Re_i])$ 在 $D(e_i)$ 上取值 $n_i$，即 $\rho\circ\theta(f)=f$。

(5) 若 $x\in\Ker\rho$，则 $\rho(x)=0$，所以
\[
p(j(x))=x-\theta(0)=x.
\]
因此 $p\circ j=1$。

(6) 由 $\rho\circ\theta=1$，$\theta$ 给出 $\rho$ 的截面。于是任意 $x\in K_0(R)$ 唯一写成
\[
x=\theta(\rho(x))+\bigl(x-\theta(\rho(x))\bigr),
\]
第二项在 $\Ker\rho$ 中。
\end{proof}

\begin{proposition}
定义
\[
\iota:\Pic(R)\to\Ker\rho,\qquad
\iota(\langle M\rangle)=[M]-[R].
\]
若 $[M]-[R^m]\in\Ker\rho$ 且 $M$ 常秩为 $m$，定义
\[
\lambda([M]-[R^m])=\left\langle\bigwedge^m M\right\rangle.
\]
则 $\lambda:\Ker\rho\to\Pic(R)$ 良定，且 $\lambda\circ\iota=1$。
\end{proposition}

\begin{proof}
若 $M$ 可逆，则 $M$ 常秩 $1$，所以
$\rho([M]-[R])=1-1=0$，$\iota$ 的确落在 $\Ker\rho$。它是群同态，因为
可逆模满足
\[
[M\otimes_RN]-[R]=([M]-[R])+([N]-[R])
\]
等价于 $[M\otimes_RN]+[R]=[M]+[N]$；这可在每个局部化处化为
$R_{\mathfrak p}\oplus R_{\mathfrak p}\simeq R_{\mathfrak p}\oplus R_{\mathfrak p}$，
再由有限生成投射模的局部同构判据得到全局等式。

若
\[
[M]-[R^m]=[N]-[R^n],
\]
则存在 $k$ 使
\[
M\oplus R^{n+k}\simeq N\oplus R^{m+k}.
\]
取最高外幂。因为 $M$ 秩为 $m$，$R^{n+k}$ 秩为 $n+k$，只有一项达到总秩 $m+n+k$，所以
\[
\bigwedge^{m+n+k}(M\oplus R^{n+k})
\simeq
\bigwedge^mM.
\]
对 $N\oplus R^{m+k}$ 作同一最高外幂计算，右边同构于 $\bigwedge^nN$。因此 $\lambda$ 与代表元无关。外幂对直和的最高项公式给出 $\lambda$ 是群同态。若 $M$ 可逆，则秩为 $1$，
\[
\lambda(\iota(\langle M\rangle))
\allowbreak
=\lambda([M]-[R])
=\langle M\rangle.
\]
\end{proof}

由上面的构造可定义行列式映射
\[
\det=\lambda\circ p:K_0(R)\longrightarrow\Pic(R).
\]
它把一个投射模的 $K_0$ 类先去掉秩函数部分，再取最高外幂。直观地说，秩记录“有几维”，行列式记录“最高外幂这条线丛怎样扭曲”。

\subsection{Dedekind 环上的 \texorpdfstring{$K_0$}{K0}}

\begin{lemma}
设 $\mathcal O$ 为 Dedekind 环。任意非零有限生成投射 $\mathcal O$-模同构于若干可逆 $\mathcal O$-模的直和。
\end{lemma}

\begin{proof}
对秩归纳。设 $M\ne0$，且 $M\oplus N\simeq\mathcal O^n$。把 $M$ 嵌入 $\mathcal O^n$，再投影到某个坐标，得到非零同态
\[
\varphi:M\to\mathcal O.
\]
其像 $I=\varphi(M)$ 是 $\mathcal O$ 的非零理想。Dedekind 环中非零理想可逆，因而是投射模。满射 $M\to I$ 分裂，所以
\[
M\simeq I\oplus M'
\]
且 $M'$ 的秩比 $M$ 小 $1$。归纳完成。
\end{proof}

\begin{theorem}
设 $\mathcal O$ 为 Dedekind 环，$F$ 为其分式域。
\begin{enumerate}[label=(\arabic*)]
\item
\[
\iota:\Pic(\mathcal O)\to\Ker\rho,\qquad
\langle M\rangle\mapsto[M]-[\mathcal O]
\]
是群同构。
\item
\[
K_0(\mathcal O)\simeq\Z\oplus\Pic(\mathcal O)
\simeq\Z\oplus\Cl_F.
\]
\end{enumerate}
\end{theorem}

\begin{proof}
因为 $\mathcal O$ 是整环，$\Spec\mathcal O$ 连通，所以
\[
C(\Spec\mathcal O,\Z)\simeq\Z.
\]
于是前面的分解给出
\[
K_0(\mathcal O)\simeq\Z\oplus\Ker\rho.
\]
再证明 $\Ker\rho\simeq\Pic(\mathcal O)$。

首先，$\iota$ 是群同态。对可逆理想 $I,J$，需要比较
\[
([I]-[\mathcal O])+([J]-[\mathcal O])
\quad\text{与}\quad
[IJ]-[\mathcal O].
\]
这等价于
\[
I\oplus J\simeq IJ\oplus\mathcal O.
\]
先处理 $I+J=\mathcal O$ 的情形。取 $a\in I,b\in J$ 使 $a+b=1$。若
$x\in I\cap J$，则 $x=xa+xb$，而 $xa,xb$ 都属于 $IJ$，故
$I\cap J\subset IJ$；反向包含由理想乘积定义得到，所以 $I\cap J=IJ$。加法映射
\[
I\oplus J\longrightarrow\mathcal O,\qquad (u,v)\longmapsto u+v
\]
满射，其核为 $\{(u,-u):u\in I\cap J\}\simeq I\cap J=IJ$。因为
$\mathcal O$ 是自由模，这个短正合列分裂，于是
$I\oplus J\simeq IJ\oplus\mathcal O$。

再处理任意非零理想 $I,J$。取 $0\ne x\in I$。由于 $x\mathcal O\subset I$，在分式理想群中存在理想
$I_1$ 使 $x\mathcal O=II_1$。在商环 $\mathcal O/I_1J$ 中，$I_1/I_1J$
是一个可由单个元素生成的理想；取 $y\in I_1$ 使
\[
I_1=I_1J+y\mathcal O.
\]
两边乘以 $I$，并用 $II_1=x\mathcal O$，得到
\[
x\mathcal O=xJ+yI.
\]
令 $I'=x^{-1}yI$。则 $I'$ 与 $I$ 作为 $\mathcal O$-模同构，且由上式除以
$x$ 得
\[
\mathcal O=J+I'.
\]
所以可把互素情形应用于 $I'$ 与 $J$：
\[
I'\oplus J\simeq I'J\oplus\mathcal O.
\]
乘以非零元素 $x^{-1}y$ 给出 $I'\simeq I$ 和 $I'J\simeq IJ$ 的模同构，故
$I\oplus J\simeq IJ\oplus\mathcal O$。于是 $\iota$ 是群同态。

其次，$\lambda\circ\iota=1$ 已在上一命题证明，所以 $\iota$ 为单射。

最后，设 $x\in\Ker\rho$。按引理，每个投射模都分解成可逆模直和。因此 $x$ 由若干 $[I]-[\mathcal O]$ 生成，落在 $\iota$ 的像中。于是 $\iota$ 满射。结合 Picard 群与理想类群的同构，得到 (2)。
\end{proof}

\section*{习题详解}
\addcontentsline{toc}{section}{习题详解}

本节按第一章习题的题号给出解答。每题先给出可解性状态，再写出完整推理步骤；涉及后文类域论或更深背景的题，会标明所用依赖，并在第一章范围内展开可直接验证的部分。

\begin{exercise}
证明：数域 $F$ 的元素 $\alpha$ 为代数整数，当且仅当 $\alpha$ 在 $\Q$ 上的最小多项式属于 $\Z[X]$。
\end{exercise}

\begin{proof}[解答]
状态：solvable。
若 $\alpha$ 的最小多项式 $m_\alpha(X)$ 属于 $\Z[X]$，并且按最小多项式的约定是首一多项式，则 $\alpha$ 满足一个首一整系数多项式，所以 $\alpha$ 是代数整数。

反过来设 $\alpha$ 是代数整数。则存在首一多项式 $f(X)\in\Z[X]$ 使 $f(\alpha)=0$。令 $m_\alpha(X)\in\Q[X]$ 为 $\alpha$ 在 $\Q$ 上的首一最小多项式。因为 $m_\alpha$ 在 $\Q[X]$ 中整除 $f$，$m_\alpha$ 的全部根都是 $f$ 的根，所以 $\alpha$ 的每个共轭都是代数整数。最小多项式的系数是这些共轭的初等对称多项式，故也是代数整数。另一方面这些系数本来在 $\Q$ 中。若 $c\in\Q$ 同时是代数整数，写 $c=a/b$ 且 $\gcd(a,b)=1$，代入首一整系数方程并清分母得到 $b\mid a^n$，故 $b=1$，所以 $c\in\Z$。因此 $m_\alpha(X)$ 的每个系数都属于 $\Z$，即 $m_\alpha(X)\in\Z[X]$。
\end{proof}

\begin{exercise}
设 $f(X)\in\Z[X]$ 为首一多项式。证明 $f(X)$ 的任一根都是代数整数。
\end{exercise}

\begin{proof}[解答]
状态：solvable。
这正是代数整数定义的直接应用。设 $\alpha$ 是 $f$ 的根，则 $f(\alpha)=0$，且 $f$ 首一、系数在 $\Z$ 中。由于 $\Z\subset\mathcal O_F$ 对任一包含 $\alpha$ 的数域都成立，$\alpha$ 满足一个首一整系数多项式，因此为代数整数。讲解时要强调“首一”不可删去，例如 $2X-1$ 的根 $1/2$ 不是代数整数。
\end{proof}

\begin{exercise}
证明：数域 $F$ 的元素 $\alpha$ 为代数整数，当且仅当 $\Z[\alpha]$ 是有限生成自由 $\Z$-模。
\end{exercise}

\begin{proof}[解答]
状态：solvable。
若 $\alpha$ 是代数整数，设它满足
\[
\alpha^n+c_{n-1}\alpha^{n-1}+\cdots+c_0=0,\qquad c_i\in\Z.
\]
于是
\[
\alpha^n=-(c_{n-1}\alpha^{n-1}+\cdots+c_0).
\]
乘以 $\alpha^r$ 后可把所有 $\alpha^{n+r}$ 写成 $1,\alpha,\dots,\alpha^{n-1}$ 的 $\Z$-线性组合。因此
\[
\Z[\alpha]=\Z+\Z\alpha+\cdots+\Z\alpha^{n-1},
\]
是有限生成 $\Z$-模。它作为 $\Q$-向量空间 $F$ 的子群没有 $\Z$-挠，所以有限生成无挠 $\Z$-模，因而自由。

反过来，若 $\Z[\alpha]$ 是有限生成 $\Z$-模，则乘法映射
\[
m_\alpha:\Z[\alpha]\to\Z[\alpha],\qquad x\mapsto\alpha x
\]
是 $\Z$-模自同态。取有限生成元后，Cayley--Hamilton 定理说明 $m_\alpha$ 满足一个首一整系数多项式。把这个多项式作用在 $1$ 上，就得到 $\alpha$ 满足首一整系数多项式，所以 $\alpha$ 为代数整数。
\end{proof}

\begin{exercise}
设
\[
f(X)=X^m+a_1X^{m-1}+\cdots+a_m,
\]
其中 $a_1,\dots,a_m$ 都是代数整数。若 $f(\beta)=0$，证明 $\beta$ 是代数整数。
\end{exercise}

\begin{proof}[解答]
状态：solvable。
令 $A=\Z[a_1,\dots,a_m]$。代数整数在加法和乘法下封闭，故 $A$ 是由有限个代数整数生成的有限 $\Z$-模。由方程
\[
\beta^m=-(a_1\beta^{m-1}+\cdots+a_m)
\]
这个递推关系说明
\[
A[\beta]=A+A\beta+\cdots+A\beta^{m-1}
\]
是有限生成 $A$-模，从而是有限生成 $\Z$-模。于是 $\Z[\beta]\subset A[\beta]$ 也是有限生成 $\Z$-模，按第 3 题，$\beta$ 是代数整数。
\end{proof}

\begin{exercise}
证明：$r\in\Q$ 是代数整数，当且仅当 $r\in\Z$。
\end{exercise}

\begin{proof}[解答]
状态：solvable。
$r\in\Z$ 时满足首一方程 $X-r=0$，因此是代数整数。反过来，写 $r=a/b$，其中 $a,b\in\Z$，$b>0$，$\gcd(a,b)=1$。若 $r$ 是代数整数，则存在
\[
r^n+c_{n-1}r^{n-1}+\cdots+c_0=0,\qquad c_i\in\Z.
\]
乘以 $b^n$ 得
\[
a^n+c_{n-1}a^{n-1}b+\cdots+c_0b^n=0.
\]
除第一项外每项都被 $b$ 整除，所以 $b\mid a^n$。由 $\gcd(a,b)=1$ 得 $b=1$，故 $r\in\Z$。
\end{proof}

\begin{exercise}
设 $\alpha$ 是非零代数整数，$m(X)$ 是 $\alpha$ 在 $\Q$ 上的最小多项式。证明：$\alpha^{-1}$ 是代数整数，当且仅当 $m(0)=\pm1$。
\end{exercise}

\begin{proof}[解答]
状态：solvable。
写
\[
m(X)=X^n+c_{n-1}X^{n-1}+\cdots+c_1X+c_0\in\Z[X].
\]
因为 $\alpha\ne0$，$c_0=m(0)\ne0$。由 $m(\alpha)=0$ 得
\[
1+c_{n-1}\alpha^{-1}+\cdots+c_1\alpha^{-(n-1)}+c_0\alpha^{-n}=0.
\]
若 $c_0=\pm1$，则
\[
(\alpha^{-1})^n\mp\bigl(1+c_{n-1}\alpha^{-1}+\cdots+c_1\alpha^{-(n-1)}\bigr)=0,
\]
所以 $\alpha^{-1}$ 满足首一整系数多项式，是代数整数。

反过来，若 $\alpha^{-1}$ 是代数整数，则 $\alpha$ 与 $\alpha^{-1}$ 都在整数环中，故 $\alpha$ 是该整数环的单位。元素范数满足
\[
\Nm_{\Q(\alpha)/\Q}(\alpha)\Nm_{\Q(\alpha)/\Q}(\alpha^{-1})=1.
\]
两个范数都是有理代数整数，即整数，所以 $\Nm(\alpha)=\pm1$。但
\[
\Nm_{\Q(\alpha)/\Q}(\alpha)=(-1)^n c_0,
\]
故 $c_0=\pm1$。
\end{proof}

\begin{exercise}
设 $f(X)\in\Z[X]$ 为首一多项式。若有数域元素 $\alpha$ 使 $f(\alpha)$ 为代数整数，证明 $\alpha$ 是代数整数。
\end{exercise}

\begin{proof}[解答]
状态：solvable。
设 $f(X)=X^n+c_{n-1}X^{n-1}+\cdots+c_0$，令 $\gamma=f(\alpha)$。由假设 $\gamma$ 是代数整数。于是 $\alpha$ 是首一多项式
\[
f(X)-\gamma\in \mathcal O_{\Q(\alpha,\gamma)}[X]
\]
的根。第 4 题的证明并没有要求系数都在 $\Z$ 中，只要求它们是代数整数；这里 $c_i$ 与 $-\gamma$ 都是代数整数，所以同一证明适用，得到 $\alpha$ 是代数整数。这里常见误区是以为 $f(\alpha)$ 整就能直接推出 $\alpha$ 整；真正用到的是 $f(X)-f(\alpha)$ 对 $X$ 仍是首一多项式。
\end{proof}

\begin{exercise}
证明环
\[
R=\C[X,Y,Z]/(X^2+Y^2+Z^2-1)
\]
不是唯一分解整环。题中给出的关键恒等式为
\[
(1+Z)(1-Z)=(X+iY)(X-iY).
\]
\end{exercise}

\begin{proof}[解答]
状态：solvable。
在 $R$ 中有
\[
X^2+Y^2=(1-Z)(1+Z),
\]
而 $X^2+Y^2=(X+iY)(X-iY)$，故上述分解成立。

要说明这不是同一个分解，需要检查四个因子都不是单位，并且两边因子之间不只是相差单位。以 $1+Z$ 为例，商环
\[
R/(1+Z)\simeq \C[X,Y]/(X^2+Y^2)
=\C[X,Y]/((X+iY)(X-iY))
\]
不是零环，所以 $1+Z$ 不是单位。对另外三个因子分别取商：$R/(1-Z)$、$R/(X+iY)$、$R/(X-iY)$ 都仍有非零像，例如都可化到某个非零多项式商环中；一个单位生成的商环应为零环，因此 $1-Z$、$X+iY$、$X-iY$ 也都不是单位。

进一步看理想 $(1+Z,X+iY)$。在商环中令 $Z=-1$ 且 $X=-iY$，关系式自动成立，因此商为 $\C[Y]$，不是零环也不是整个环。这说明 $1+Z$ 不整除 $X+iY$。把 $X+iY$ 换成 $X-iY$，或把 $1+Z$ 换成 $1-Z$，可得到相同类型的非整除结论。若 $R$ 是唯一分解整环，则不可约因子分解唯一，左边因子必须与右边某因子相伴，矛盾。因此 $R$ 不是 UFD。
\end{proof}

\begin{exercise}
证明 $\Z[i]$ 是主理想环。
\end{exercise}

\begin{proof}[解答]
状态：solvable。
在 Gaussian 整数环 $\Z[i]$ 上定义范数
\[
N(a+bi)=a^2+b^2.
\]
它满足 $N(\alpha\beta)=N(\alpha)N(\beta)$，且 $N(\alpha)=0$ 当且仅当 $\alpha=0$。

证明 Euclidean 除法。给定 $\alpha,\beta\in\Z[i]$，$\beta\ne0$。在复平面中写 $\alpha/\beta=u+vi$。取整数 $m,n$ 使
\[
|u-m|\le\frac12,\qquad |v-n|\le\frac12.
\]
令 $q=m+ni$，$r=\alpha-q\beta$。则
\[
N(r)=N(\beta)N(\alpha/\beta-q)
\le N(\beta)\left(\frac14+\frac14\right)
<N(\beta),
\]
除非边界等号情形，此时仍可调整最近格点使严格小于。故 $\Z[i]$ 是 Euclidean 环，因此是主理想环。
\end{proof}

\begin{exercise}
以 $\OF$ 记数域 $F$ 的整数环。证明 $\OF$ 的任一理想可用不超过两个元素生成。举例一个数域 $F$ 使 $\OF$ 不是主理想环。
\end{exercise}

\begin{proof}[解答]
状态：solvable。
设 $\mathfrak a\ne0$ 为 $\OF$ 的理想。取 $0\ne x\in\mathfrak a$，则 $(x)\subset\mathfrak a$。把两个理想分解为素理想：
\[
(x)=\prod_{\mathfrak p}\mathfrak p^{m_{\mathfrak p}},\qquad
\mathfrak a=\prod_{\mathfrak p}\mathfrak p^{n_{\mathfrak p}},
\]
其中 $m_{\mathfrak p}\ge n_{\mathfrak p}$。我们要找 $y\in\mathfrak a$，使
\[
(x,y)=\mathfrak a.
\]
这等价于对每个 $\mathfrak p$，有
\[
v_{\mathfrak p}((x,y))
=\min(v_{\mathfrak p}(x),v_{\mathfrak p}(y))
=n_{\mathfrak p}.
\]
只需安排 $v_{\mathfrak p}(y)=n_{\mathfrak p}$ 对有限多个出现在 $(x)$ 或 $\mathfrak a$ 中的 $\mathfrak p$ 成立。由 Dedekind 环的中国剩余定理，可选择 $y$ 满足
\[
y\in\mathfrak p^{n_{\mathfrak p}}\setminus \mathfrak p^{n_{\mathfrak p}+1}
\]
对这些有限个 $\mathfrak p$ 同时成立，并让 $y\in\mathfrak a$。于是 $\mathfrak a=(x,y)$。

非主理想环例子可取 $F=\Q(\sqrt{-5})$。其整数环为 $\Z[\sqrt{-5}]$。在其中
\[
6=2\cdot3=(1+\sqrt{-5})(1-\sqrt{-5})
\]
给出元素唯一分解失败。若 $\mathcal O_F$ 是主理想环，则作为 Dedekind 环它是 UFD，矛盾。因此 $\Z[\sqrt{-5}]$ 不是主理想环。
\end{proof}

\begin{exercise}
设 $d$ 为无平方因子整数，$K=\Q(\sqrt d)$。证明
\[
\mathcal O_K=
\begin{cases}
\Z+\Z\sqrt d,&d\equiv2,3\pmod4,\\
\Z+\Z\frac{1+\sqrt d}{2},&d\equiv1\pmod4.
\end{cases}
\]
\end{exercise}

\begin{proof}[解答]
状态：solvable。
任意 $\alpha\in K$ 可写为 $\alpha=a+b\sqrt d$，$a,b\in\Q$。若 $\alpha$ 是代数整数，则迹和范都是整数：
\[
\Tr(\alpha)=2a\in\Z,\qquad
\Nm(\alpha)=a^2-b^2d\in\Z.
\]
写 $2a=m\in\Z$。又因为 $\alpha$ 满足二次多项式
\[
X^2-\Tr(\alpha)X+\Nm(\alpha)\in\Z[X],
\]
把 $\alpha=(m+2b\sqrt d)/2$ 代入范数公式可见 $4b^2d=m^2-4\Nm(\alpha)\in\Z$。由于 $d$ 无平方因子，分母比较迫使 $2b\in\Z$；写 $2b=n$。于是
\[
\Nm(\alpha)=\frac{m^2-n^2d}{4}\in\Z,
\]
即
\[
m^2\equiv n^2d\pmod4.
\]
平方数模 $4$ 只可能为 $0$ 或 $1$。若 $d\equiv2,3\pmod4$，则同余迫使 $m,n$ 同为偶数，因此 $\alpha\in\Z+\Z\sqrt d$。若 $d\equiv1\pmod4$，则 $m,n$ 同奇或同偶。于是
\[
\alpha=\frac{m+n\sqrt d}{2}
=u+v\frac{1+\sqrt d}{2}
\]
其中 $u,v\in\Z$。反过来，$\sqrt d$ 满足 $X^2-d=0$；当 $d\equiv1\pmod4$ 时，$(1+\sqrt d)/2$ 满足
\[
X^2-X+\frac{1-d}{4}=0.
\]
这些都是首一整系数方程，故列出的元素都整。
\end{proof}

\begin{exercise}
设 $p>2$ 为素数，$p\nmid d$，并令 $K=\Q(\sqrt d)$。用 Legendre 符号判断 $(p)$ 在 $\mathcal O_K$ 中的分解；若 $p$ 整除判别式，则说明分歧情形。
\end{exercise}

\begin{proof}[解答]
状态：solvable。
设 $\delta$ 为二次域判别式：
\[
\delta=
\begin{cases}
4d,&d\equiv2,3\pmod4,\\
d,&d\equiv1\pmod4.
\end{cases}
\]
当 $p\nmid\delta$ 时，$p$ 不分歧。整数环由一个元素 $\omega$ 生成，且其最小多项式模 $p$ 的分解决定 $(p)$ 的分解。

若 $d\equiv2,3\pmod4$，取 $\omega=\sqrt d$，最小多项式为 $X^2-d$。模 $p$ 后：
若 $\left(\frac d p\right)=-1$，$X^2-d$ 在 $\F_p$ 不可约，所以
\[
(p)\ \text{保持素理想}.
\]
若 $\left(\frac d p\right)=1$，取 $a^2\equiv d\pmod p$，则
\[
X^2-d\equiv(X-a)(X+a)\pmod p,
\]
故
\[
(p)=\mathfrak p\mathfrak p',
\quad
\mathfrak p=(p,\sqrt d-a),\quad
\mathfrak p'=(p,\sqrt d+a),
\]
且 $\mathfrak p\ne\mathfrak p'$。

若 $d\equiv1\pmod4$，取 $\omega=(1+\sqrt d)/2$，其判别式仍为 $d$，模 $p$ 是否分裂仍由 $d$ 是否为平方决定。

若 $p\mid\delta$，最小多项式模 $p$ 有重根，于是
\[
(p)=\mathfrak p^2
\]
为分歧情形。这也和 Dedekind 判别式定理一致：素数分歧当且仅当整除判别式。
\end{proof}

\begin{exercise}
在二次域中判断素数 $2$ 的分解：若 $2\nmid\delta$ 且 $d\equiv5\pmod8$，则 $(2)$ 为素理想；若 $2\nmid\delta$ 且 $d\equiv1\pmod8$，则 $(2)$ 分裂；若 $2\mid\delta$，则 $(2)$ 分歧。
\end{exercise}

\begin{proof}[解答]
状态：solvable。
只有 $d\equiv1\pmod4$ 时才可能 $2\nmid\delta$，此时 $\delta=d$，整数环为
\[
\mathcal O_K=\Z[\omega],\qquad \omega=\frac{1+\sqrt d}{2}.
\]
$\omega$ 的最小多项式为
\[
X^2-X+\frac{1-d}{4}.
\]
模 $2$ 后分两种情况。

若 $d\equiv5\pmod8$，则 $(1-d)/4\equiv1\pmod2$，多项式为
\[
X^2+X+1\in\F_2[X],
\]
它在 $\F_2$ 中无根，所以不可约，$(2)$ 保持素理想。

若 $d\equiv1\pmod8$，则 $(1-d)/4\equiv0\pmod2$，多项式为
\[
X^2+X=X(X+1),
\]
有两个不同根，所以
\[
(2)=\mathfrak p\mathfrak p'
\]
分裂为两个不同素理想。

若 $2\mid\delta$，由 Dedekind 判别式定理，$2$ 分歧；在二次扩张中只能是
\[
(2)=\mathfrak p^2.
\]
\end{proof}

\begin{exercise}
设 $d$ 为无平方因子整数，$\delta$ 为 $\Q(\sqrt d)$ 的判别式。证明 $\Q(\zeta_{|\delta|})$ 是包含 $\Q(\sqrt d)$ 的最小分圆域。
\end{exercise}

\begin{proof}[解答]
状态：solvable，使用二次域与二次 Dirichlet 特征标的对应。
所用事实是二次域与二次 Dirichlet 特征标的对应。定义 Kronecker 特征标
\[
\chi_\delta(n)=\left(\frac{\delta}{n}\right),
\]
它的导子正是 $|\delta|$。这里“导子”的含义是：$\chi_\delta$ 可看作模
$|\delta|$ 的乘法特征标，并且不存在更小的正整数 $m$ 使它来自
$(\Z/m\Z)^\times$。分圆域 $\Q(\zeta_m)$ 的 Galois 群为
\[
\Gal(\Q(\zeta_m)/\Q)\simeq(\Z/m\Z)^\times.
\]
一个二次子域对应一个指数 $2$ 子群，也等价于一个二次 Dirichlet 特征标。对 $m=|\delta|$，$\chi_\delta$ 给出满同态
\[
(\Z/|\delta|\Z)^\times\to\{\pm1\}.
\]
其核的固定域是一个二次域。为什么它就是 $\Q(\sqrt d)$？对任意不整除
$\delta$ 的素数 $p$，Frobenius 元在该二次子域中平凡，当且仅当
$\chi_\delta(p)=1$。另一方面，第 12、13 题的二次域分解判据说明，$p$ 在
$\Q(\sqrt d)$ 中分裂，当且仅当同一个符号 $\chi_\delta(p)=1$。两个二次扩张在除有限多个素数外有相同的分裂素数，因此对应同一个二次特征标，固定域就是
$\Q(\sqrt d)$。

再证最小性。若 $\Q(\sqrt d)\subset\Q(\zeta_m)$，则对应的二次特征标必须模
$m$，因为 $\Gal(\Q(\zeta_m)/\Q)\simeq(\Z/m\Z)^\times$ 的指数 $2$ 商正是一个模
$m$ 的二次特征标。这意味着 $\chi_\delta$ 的导子整除 $m$。导子已经是
$|\delta|$，所以 $|\delta|\mid m$。因此最小的 $m$ 是 $|\delta|$，最小分圆域为
$\Q(\zeta_{|\delta|})$。
\end{proof}

\begin{exercise}
证明 $f(X)=X^3-X-1$ 在 $\Q$ 上不可约。设 $\alpha$ 为其根，证明 $\{1,\alpha,\alpha^2\}$ 是 $\Q(\alpha)$ 的一组整基。
\end{exercise}

\begin{proof}[解答]
状态：solvable。
三次多项式在 $\Q$ 上可约当且仅当有有理根。由有理根定理，可能的有理根只有 $\pm1$。直接计算
\[
f(1)=-1,\qquad f(-1)=-1,
\]
所以无有理根，因而不可约。

因为 $f$ 首一且在 $\Z[X]$ 中，$\alpha$ 是代数整数，所以
\[
\Z[\alpha]\subset\mathcal O_{\Q(\alpha)}.
\]
三次多项式 $X^3+pX+q$ 的判别式为 $-4p^3-27q^2$。这里 $p=-1,q=-1$，所以
\[
\operatorname{disc}(1,\alpha,\alpha^2)=4-27=-23.
\]
若 $\mathcal O_K$ 严格大于 $\Z[\alpha]$，则指数
\[
[\mathcal O_K:\Z[\alpha]]
\]
的平方整除 $23$。由于 $23$ 平方自由，指数只能为 $1$。因此 $\mathcal O_K=\Z[\alpha]$，所给三元组为整基。
\end{proof}

\begin{exercise}
设 $\omega=(-1+\sqrt{-3})/2$。证明 $\omega$ 是三次本原单位根，$\Z[\omega]$ 是主理想环，并证明题中关于 $1-\omega$、$3$ 和有理素数 $p$ 的分解断言。
\end{exercise}

\begin{proof}[解答]
状态：solvable。
先有
\[
\omega^2+\omega+1=0,\qquad \omega^3=1,\qquad \omega\ne1,
\]
所以 $\omega$ 是三次本原单位根。

环 $\Z[\omega]$ 的范数为
\[
N(a+b\omega)=(a+b\omega)(a+b\omega^2)=a^2-ab+b^2.
\]
它来自复平面中的六角格。任意复数都可找到最近的 Eisenstein 整数 $q\in\Z[\omega]$，使 $N(z-q)<1$。于是对 $\alpha,\beta\in\Z[\omega]$，$\beta\ne0$，可取 $q$ 使
\[
N(\alpha-q\beta)<N(\beta).
\]
所以 $\Z[\omega]$ 是 Euclidean 环，因而是主理想环。

计算
\[
N(1-\omega)=3.
\]
商环 $\Z[\omega]/(1-\omega)$ 有 $3$ 个元素，因此 $(1-\omega)$ 是极大理想，$1-\omega$ 是素元。又
\[
3=(1-\omega)(1-\omega^2).
\]
因为 $1-\omega^2=(1-\omega)(1+\omega)$，而 $1+\omega=-\omega^2$ 是单位，所以
\[
3=-\omega^2(1-\omega)^2.
\]

若有理素数 $p\equiv2\pmod3$，则 $X^2+X+1$ 在 $\F_p$ 中不可约，所以 $(p)$ 在 $\Z[\omega]$ 中保持素理想，$p$ 是素元。若 $p\equiv1\pmod3$，则 $\F_p^\times$ 中有非平凡三次单位根，$X^2+X+1$ 分裂，于是
\[
(p)=\mathfrak p\bar{\mathfrak p}.
\]
由于 $\Z[\omega]$ 是 PID，$\mathfrak p=(\pi)$，于是
\[
p=u\pi\bar\pi
\]
对某个单位 $u$。吸收单位后可写为 $p=\pi\bar\pi$，且 $\pi,\bar\pi$ 都为素元。
\end{proof}

\begin{exercise}
证明 $\Z$ 的单位为 $\pm1$。再证明 $\Z[\omega]$ 的单位为
\[
\pm1,\ \pm\omega,\ \pm\omega^2.
\]
\end{exercise}

\begin{proof}[解答]
状态：solvable。
若 $n\in\Z$ 是单位，则存在 $m\in\Z$ 使 $mn=1$，故 $n=\pm1$。

对 $\Z[\omega]$，单位 $\varepsilon$ 满足 $N(\varepsilon)=1$。写
\[
\varepsilon=a+b\omega,\qquad N(\varepsilon)=a^2-ab+b^2=1.
\]
把它改写为
\[
4N(\varepsilon)=(2a-b)^2+3b^2=4.
\]
所以 $b=0,\pm1$。逐一代入：$b=0$ 得 $a=\pm1$；$b=1$ 得 $a=0$ 或 $1$，对应 $\omega$ 和 $1+\omega=-\omega^2$；$b=-1$ 得 $a=0$ 或 $-1$，对应 $-\omega$ 和 $\omega^2$。故单位正是
\[
\{\pm1,\pm\omega,\pm\omega^2\}.
\]
\end{proof}

\begin{exercise}
设 $U_d$ 为二次域 $\Q(\sqrt d)$ 的整数环单位群。证明虚二次情形的单位群，并说明实二次情形 $U_d=\{\pm u^n:n\in\Z\}$。
\end{exercise}

\begin{proof}[解答]
状态：solvable。
单位 $\varepsilon$ 必满足
\[
N_{K/\Q}(\varepsilon)=\pm1.
\]
若 $d<0$，把 $K$ 嵌入 $\C$，范数为复绝对值平方：
\[
N(\varepsilon)=|\varepsilon|^2.
\]
单位必须有 $|\varepsilon|=1$。虚二次整数环在复平面中是离散格，与单位圆的交有限。逐一检查只有两个特殊域有额外单位：
\[
U_{-1}=\{\pm1,\pm i\},
\]
\[
U_{-3}=\{\pm1,\pm\omega,\pm\omega^2\}.
\]
当 $d<-3$ 或 $d=-2$ 时，格与单位圆只交于 $\pm1$，故 $U_d=\{\pm1\}$。

若 $d>0$，两个实嵌入给出对数映射
\[
\varepsilon\mapsto(\log|\iota_1(\varepsilon)|,\log|\iota_2(\varepsilon)|).
\]
因为 $|N(\varepsilon)|=1$，两个坐标和为 $0$，单位群映入直线 $x+y=0$。二次情形的 Dirichlet 单位定理断言这个像是该直线中的秩 $1$ 格，也就是某个非零向量的整数倍全体。取对应的正单位 $u>1$ 为最小正生成元；任意单位除以适当的 $u^n$ 后在两个实嵌入下绝对值都为 $1$，只可能是 $\pm1$。于是
\[
U_d=\{\pm u^n:n\in\Z\}.
\]
这一步的外部依赖只用于保证实二次域存在基本单位并且所有单位由它生成。
\end{proof}

\begin{exercise}
设 $R$ 为 Dedekind 环，$F$ 为分式域，$\mathfrak p_1,\dots,\mathfrak p_m$ 互不相同，$n_i\in\Z$。证明存在 $x\in F^\times$ 和与所有 $\mathfrak p_i$ 互素的分式理想 $\mathfrak b$，使
\[
xR=\mathfrak p_1^{n_1}\cdots\mathfrak p_m^{n_m}\mathfrak b.
\]
\end{exercise}

\begin{proof}[解答]
状态：solvable。
令
\[
\mathfrak a=\mathfrak p_1^{n_1}\cdots\mathfrak p_m^{n_m}.
\]
我们要找主理想 $xR$，使 $xR\mathfrak a^{-1}$ 在每个 $\mathfrak p_i$ 处指数为 $0$。选择元素 $x\in F^\times$，满足
\[
v_{\mathfrak p_i}(x)=n_i\qquad(1\le i\le m).
\]
这种元素存在，可由 Dedekind 环中的弱近似得到：先对每个 $i$ 取 $x_i\in\mathfrak p_i^{n_i}\setminus\mathfrak p_i^{n_i+1}$，再用中国剩余定理让一个元素在有限多个局部条件下同时逼近。令
\[
\mathfrak b=xR\,\mathfrak a^{-1}.
\]
则 $v_{\mathfrak p_i}(\mathfrak b)=0$，所以 $\mathfrak b$ 与所有 $\mathfrak p_i$ 互素，且 $xR=\mathfrak a\mathfrak b$。
\end{proof}

\begin{exercise}
设 $R$ 为 Dedekind 环，$\mathfrak a_1,\dots,\mathfrak a_m$ 两两互素。令
\[
\mathfrak b_i=\prod_{j\ne i}\mathfrak a_j.
\]
证明 $\mathfrak b_1+\cdots+\mathfrak b_m=R$。
\end{exercise}

\begin{proof}[解答]
状态：solvable。
若和不是 $R$，则它包含于某个极大理想 $\mathfrak p$。于是每个 $\mathfrak b_i\subset\mathfrak p$。因为 $\mathfrak p$ 是素理想，$\mathfrak b_i=\prod_{j\ne i}\mathfrak a_j\subset\mathfrak p$ 蕴含某个 $j\ne i$ 满足 $\mathfrak a_j\subset\mathfrak p$。

对每个 $i$ 都能找到 $j\ne i$ 使 $\mathfrak a_j\subset\mathfrak p$。当 $m\ge2$ 时，这迫使至少两个不同的 $\mathfrak a_j$ 落在同一个 $\mathfrak p$ 中。但若 $\mathfrak a_r,\mathfrak a_s\subset\mathfrak p$，则
\[
\mathfrak a_r+\mathfrak a_s\subset\mathfrak p\ne R,
\]
与两两互素矛盾。因此和只能是 $R$。
\end{proof}

\begin{exercise}
设 $F/\Q$ 为有限 Galois 扩张，$0\ne a\in\OF$。证明
\[
\sum_{\sigma\in\Gal(F/\Q)}\sigma(a)\overline{\sigma(a)}\ge [F:\Q].
\]
\end{exercise}

\begin{proof}[解答]
状态：solvable。
对每个嵌入值 $z_\sigma=\sigma(a)$，有 $|z_\sigma|^2=\sigma(a)\overline{\sigma(a)}\ge0$。因为 $a$ 是非零代数整数，范数
\[
\prod_\sigma \sigma(a)
\]
是非零整数，故
\[
\prod_\sigma |z_\sigma|^2=|\Nm_{F/\Q}(a)|^2\ge1.
\]
设 $n=[F:\Q]$。对正实数 $|z_\sigma|^2$ 用算术-几何平均不等式：
\[
\frac1n\sum_\sigma |z_\sigma|^2
\ge
\left(\prod_\sigma |z_\sigma|^2\right)^{1/n}
\ge1.
\]
乘以 $n$ 得所需不等式。
\end{proof}

\begin{exercise}
设 $\mathcal O$ 为 Dedekind 环，$F$ 为分式域，$E/F$ 为有限域扩张，不假设可分。证明 $\mathcal O_E$ 仍是 Dedekind 环。
\end{exercise}

\begin{proof}[解答]
状态：solvable，使用一维 Noether 整闭环在有限域扩张中整闭包仍为有限模这一标准有限性定理；本题重点是说明可分性不是 Dedekind 性的必要假设。
记 $\mathcal O_E$ 为 $\mathcal O$ 在 $E$ 中的整闭包。它是域 $E$ 的子环，所以是整环；按定义它在 $E$ 中整闭。非零素理想极大也和可分情形一样：若 $\mathfrak P\subset\mathcal O_E$ 非零素理想，则 $\mathfrak p=\mathfrak P\cap\mathcal O$ 为 $\mathcal O$ 的非零素理想，故极大；商 $\mathcal O_E/\mathfrak P$ 是 $\mathcal O/\mathfrak p$ 上整的整环，因底环为域，商环为域，所以 $\mathfrak P$ 极大。

剩下是 Noether 性。有限性定理给出 $\mathcal O_E$ 是有限 $\mathcal O$-模。为了看清不可分部分为什么不会破坏结论，可把 $E/F$ 分成最大可分子扩张
$E_s/F$ 与纯不可分扩张 $E/E_s$。可分层用迹配对证明
$\mathcal O_{E_s}$ 对 $\mathcal O$ 有限；纯不可分层中，任意 $x\in E$ 有
$x^{p^r}\in E_s$，整元素满足方程 $T^{p^r}-x^{p^r}$，所以整闭包由有限次添入纯不可分根控制。结合有限性定理，$\mathcal O_E$ 是有限
$\mathcal O$-模，从而 Noether。三条性质合并，$\mathcal O_E$ 是 Dedekind 环。
\end{proof}

\begin{exercise}
设 $E/F$ 为有限 Galois 扩张，$\mathfrak p$ 无分歧，$\mathfrak P|\mathfrak p$。证明 Frobenius 共轭类给出的三条分解性质。
\end{exercise}

\begin{proof}[解答]
状态：solvable。
令 $\sigma=\left[\frac{E/F}{\mathfrak P}\right]$，其共轭类记为 $F_{E/F}(\mathfrak p)$。

(1) 设 $\sigma$ 的阶为 $f$。无分歧说明 $e=1$，Frobenius 在剩余域上的阶就是剩余次数 $f_{\mathfrak P|\mathfrak p}$。由基本公式
\[
efg=[E:F]
\]
得
\[
g=[E:F]/f=|\Gal(E/F)|/|\langle\sigma\rangle|
=[\Gal(E/F):\langle\sigma\rangle].
\]

(2) 完全分裂等价于 $g=[E:F]$，等价于 $f=1$，等价于 $\sigma=1$。因为非交换时 Frobenius 只按共轭类确定，故条件为 $F_{E/F}(\mathfrak p)=\{1\}$。

(3) Frobenius 的定义就是：对任意 $a\in\mathcal O_E$，
\[
\sigma(a)\equiv a^{|\OF/\mathfrak p|}\pmod{\mathfrak P}.
\]
而 $|\OF/\mathfrak p|=\Nm(\mathfrak p)$，所以
\[
\left[\frac{E/F}{\mathfrak P}\right]a\equiv a^{\Nm(\mathfrak p)}\pmod{\mathfrak P}.
\]
\end{proof}

\begin{exercise}
设 $\zeta$ 为 $n$ 次本原单位根。证明 $\Q(\zeta)$ 的整数环是 $\Z[\zeta]$。
\end{exercise}

\begin{proof}[解答]
状态：solvable，使用分圆域整数环的标准局部判别式证明。
首先 $\zeta$ 满足 $X^n-1=0$，所以 $\Z[\zeta]$ 中元素都是代数整数，故
\[
\Z[\zeta]\subset\mathcal O_{\Q(\zeta)}.
\]
反向包含要证明任一代数整数都可由 $\zeta$ 整系数表示。证明按素数局部化。先在 $n=p^r$ 为素数幂时证明：令 $\pi=1-\zeta_{p^r}$，有
\[
p=u\pi^{\varphi(p^r)}
\]
其中 $u$ 是单位；这说明 $p$ 在 $\Z[\zeta_{p^r}]$ 中只含一个上方素理想 $(\pi)$，且分歧指数为 $\varphi(p^r)$。由 Eisenstein 多项式
\[
\Phi_{p^r}(X+1)
\]
是 Eisenstein 多项式，所以 $1-\zeta_{p^r}$ 是 $p$ 处的素元。对不等于 $p$ 的素数
$\ell$，分圆多项式的判别式在 $\ell$ 处无分歧，$\Z[\zeta_{p^r}]$ 的局部化已经整闭；在
$p$ 处，Eisenstein 性给出全分歧 DVR，局部整数环也由 $\zeta_{p^r}$ 生成。所有素数处局部化都相等，故
$\mathcal O_{\Q(\zeta_{p^r})}=\Z[\zeta_{p^r}]$。

一般 $n=\prod p_i^{r_i}$ 时，
\[
\Q(\zeta_n)
\]
是这些素数幂分圆域的合成，且不同素数幂分圆域的相对分歧发生在不同的有理素数处。逐素数局部化后，每个局部问题只看见其中一个素数幂部分，其余部分为未分歧合成；因此整环由各个
$\zeta_{p_i^{r_i}}$ 共同生成，得到
\[
\mathcal O_{\Q(\zeta_n)}=\Z[\zeta_{p_1^{r_1}},\dots,\zeta_{p_s^{r_s}}]
=\Z[\zeta_n].
\]
这就是分圆域整数环定理在本题中的应用。
\end{proof}

\begin{exercise}
设 $\zeta_m,\zeta_n$ 分别为 $m,n$ 次本原单位根。证明
\[
\Q(\zeta_m,\zeta_n)=\Q(\zeta_{\operatorname{lcm}(m,n)}).
\]
题目若写成 $\zeta_{mn}$，则该式在 $\gcd(m,n)=1$ 时成立；一般情形应改为最小公倍数。
\end{exercise}

\begin{proof}[解答]
状态：partially solvable：等式若未假设 $\gcd(m,n)=1$ 则为假；下面证明正确形式，并说明互素时得到原式。
令 $\ell=\operatorname{lcm}(m,n)$。因为 $m|\ell$、$n|\ell$，
\[
\zeta_m=\zeta_\ell^{\ell/m},\qquad
\zeta_n=\zeta_\ell^{\ell/n},
\]
所以
\[
\Q(\zeta_m,\zeta_n)\subset\Q(\zeta_\ell).
\]
反过来，设 $m=\prod p^{a_p}$、$n=\prod p^{b_p}$，则
$\ell=\prod p^{\max(a_p,b_p)}$。对每个 $p|\ell$，若 $a_p\ge b_p$，则
$\zeta_m$ 的适当幂给出一个 $p^{a_p}$ 次本原单位根；若 $b_p>a_p$，则
$\zeta_n$ 的适当幂给出一个 $p^{b_p}$ 次本原单位根。因此
$\Q(\zeta_m,\zeta_n)$ 含有所有 $p^{\max(a_p,b_p)}$ 次本原单位根。不同素数幂阶互素，这些单位根的乘积是 $\ell$ 次本原单位根，所以
$\zeta_\ell\in\Q(\zeta_m,\zeta_n)$。于是
$\Q(\zeta_\ell)\subset\Q(\zeta_m,\zeta_n)$。若 $\gcd(m,n)=1$，则 $\ell=mn$，得到题目中的形式；若不互素，例如 $m=n=2$，左边是 $\Q$，而
$\Q(\zeta_{mn})=\Q(i)$，所以原式不能按字面成立。
\end{proof}

\begin{exercise}
设 $p$ 为素数，$\zeta$ 为 $p$ 次本原单位根。证明 $(1-\zeta)$ 是 $\Z[\zeta]$ 的素理想，并且
\[
(p)=(1-\zeta)^{p-1}.
\]
\end{exercise}

\begin{proof}[解答]
状态：solvable。
有
\[
p=\Phi_p(1)=\prod_{a=1}^{p-1}(1-\zeta^a).
\]
而
\[
1-\zeta^a=(1-\zeta)(1+\zeta+\cdots+\zeta^{a-1}).
\]
当 $1\le a\le p-1$ 时，第二个因子在模 $(1-\zeta)$ 后同余于 $a$，不被 $p$ 整除，因此是局部单位。于是所有 $1-\zeta^a$ 与 $1-\zeta$ 相伴，得到
\[
(p)=(1-\zeta)^{p-1}.
\]
商环
\[
\Z[\zeta]/(1-\zeta)\simeq\Z/p\Z
\]
是域，所以 $(1-\zeta)$ 是素理想。
\end{proof}

\begin{exercise}
设 $E/F$ 为数域有限 Galois 扩张，$E=F(\alpha)$，$f(x)\in\OF[x]$ 为 $\alpha$ 的最小多项式。设 $\mathfrak p$ 使 $f\bmod\mathfrak p$ 可分。证明题中关于无分歧、素理想分解和完全分裂的断言。
\end{exercise}

\begin{proof}[解答]
状态：solvable，依赖 Dedekind 分解定理，并要求 $\mathcal O_E=\OF[\alpha]$ 或 $\mathfrak p$ 不整除指数 $[\mathcal O_E:\OF[\alpha]]$；没有这个条件时，单靠 $f\bmod\mathfrak p$ 可分不足以保证结论。
假设 $f\bmod\mathfrak p$ 分解为互异不可约多项式
\[
\bar f=\bar f_1\cdots\bar f_g.
\]
Dedekind 分解定理说明，在 $\mathcal O_E=\OF[\alpha]$ 或 $\mathfrak p$ 不除指数
$[\mathcal O_E:\OF[\alpha]]$ 的条件下，素理想分解由这个模 $\mathfrak p$ 分解给出。题设可分保证没有重复因子，因而分歧指数都为 $1$，所以 $\mathfrak p$ 在 $E$ 中无分歧。若原题省去了指数条件，应把它理解为在本章 Dedekind 分解定理可适用的情形；否则存在由指数带来的例外素数。

定义
\[
\mathfrak P_i=\mathfrak p\mathcal O_E+f_i(\alpha)\mathcal O_E.
\]
商环中有
\[
\mathcal O_E/\mathfrak P_i\simeq
(\OF/\mathfrak p)[X]/(\bar f_i),
\]
这是域，所以 $\mathfrak P_i$ 为素理想。不同 $\bar f_i$ 互素，故 $\mathfrak P_i$ 两两不同。中国剩余定理给出
\[
\mathfrak p\mathcal O_E=\mathfrak P_1\cdots\mathfrak P_g.
\]
最后，$\mathfrak p$ 完全分裂当且仅当每个 $\bar f_i$ 都一次，即 $\bar f$ 在 $\OF/\mathfrak p$ 中完全分裂。这等价于存在足够多根；特别对生成元而言，若存在 $\beta\in\OF$ 使
\[
f(\beta)\equiv0\pmod{\mathfrak p},
\]
则至少有一个一次因子。Galois 条件使一个根模 $\mathfrak p$ 出现时全部共轭根的分解型一致；在无分歧可分情形，这等价于完全分裂。
\end{proof}

\begin{exercise}
在 $\Z[\omega]$ 中定义三次剩余符号，并证明题中四条基本性质。
\end{exercise}

\begin{proof}[解答]
状态：solvable。
设 $\pi$ 是不整除 $3$ 的素元。商环 $\Z[\omega]/(\pi)$ 是有限域，元素个数为 $N(\pi)$。因为 $\pi\nmid3$，有限域中含有三个三次单位根 $1,\omega,\omega^2$。

(1) 多项式
\[
x^3-1=(x-1)(x-\omega)(x-\omega^2)
\]
在 $\Z[\omega]$ 中恒成立，模 $\pi$ 后仍成立。

(2) 若 $\pi\nmid a$，则 $a$ 在有限域 $k=\Z[\omega]/(\pi)$ 中非零。群 $k^\times$ 阶为 $N(\pi)-1$。于是
\[
a^{(N(\pi)-1)/3}
\]
是三次单位根，故同余于 $1,\omega,\omega^2$ 中唯一一个。定义这个唯一值为
\[
\left(\frac a\pi\right)_3.
\]

(3) 若 $a\equiv b\pmod\pi$，则它们在 $k^\times$ 中代表同一元素，幂
\[
a^{(N(\pi)-1)/3},\quad b^{(N(\pi)-1)/3}
\]
相同，所以三次剩余符号相同。

(4) 由有限域乘法可得
\[
(ab)^{(N(\pi)-1)/3}
=a^{(N(\pi)-1)/3}b^{(N(\pi)-1)/3},
\]
故
\[
\left(\frac{ab}{\pi}\right)_3
=\left(\frac a\pi\right)_3\left(\frac b\pi\right)_3.
\]
\end{proof}

\begin{exercise}
设 $F=\Q(\sqrt{-3})$，$E=F(\sqrt[3]{2})$，$\omega=(-1+\sqrt{-3})/2$。证明题中关于 $\OF$、无分歧、Frobenius 和三次剩余符号的公式。
\end{exercise}

\begin{proof}[解答]
状态：solvable，使用 Kummer 扩张在含三次单位根的基域上的基本分歧判据。
(1) 由二次域整数环公式，$-3\equiv1\pmod4$ 对应
\[
\OF=\Z[\omega].
\]
第 16 题已证 $\Z[\omega]$ 是 PID，所以任一素理想可写为 $\mathfrak p=(\pi)$，其中 $\pi$ 为不可约元。

(2) 扩张 $E/F$ 由方程
\[
X^3-2
\]
给出，且 $F$ 已含三次单位根。Kummer 扩张 $F(\sqrt[3]{2})/F$ 的分歧只可能出现在整除 $3$ 或整除 $2$ 的素理想处。也可直接看判别式：$X^3-2$ 的判别式为
$-27\cdot4=-108$，在 $\OF$ 中只含 $2$ 和 $3$ 上方素因子。若
$\pi\nmid6$，则 $\mathfrak p=(\pi)$ 不整除该判别式，并且不属于可能的指数例外素数；Dedekind 判别式定理给出 $\mathfrak p$ 在 $E/F$ 中无分歧。

(3) 设 $\mathfrak P|\mathfrak p=(\pi)$，令 $\alpha=\sqrt[3]{2}$。Frobenius 满足
\[
\Frob_{\mathfrak P}(\alpha)\equiv \alpha^{N(\pi)}\pmod{\mathfrak P}.
\]
因为 $\pi\nmid6$，商域中 $3$ 可逆，并且 $N(\pi)\equiv1\pmod3$；后一事实来自
$\OF/\mathfrak p$ 含有三次单位根 $\omega$ 的像。由 $\alpha^3=2$，
\[
\alpha^{N(\pi)}
=\alpha\cdot(\alpha^3)^{(N(\pi)-1)/3}
=\alpha\cdot2^{(N(\pi)-1)/3}.
\]
这给出题中同余式
$\Frob_{\mathfrak P}(\sqrt[3]{2})\equiv
2^{(N(\pi)-1)/3}\sqrt[3]{2}\pmod{\mathfrak P}$。

(4) 按三次剩余符号定义，
\[
2^{(N(\pi)-1)/3}\equiv
\left(\frac2\pi\right)_3
\pmod\pi.
\]
提升到 $\mathfrak P$ 后，
\[
\Frob_{\mathfrak P}(\alpha)
\equiv
\left(\frac2\pi\right)_3\alpha
\pmod{\mathfrak P}.
\]
由于 $E=F(\alpha)$，一个 $F$-自同构由 $\alpha\mapsto\omega^j\alpha$ 决定。上式说明 Frobenius 与把 $\alpha$ 乘以三次剩余符号的那个自同构相同，因此
\[
\left(\frac{E/F}{\mathfrak p}\right)(\sqrt[3]{2})
=\left(\frac2\pi\right)_3\sqrt[3]{2}.
\]
\end{proof}

\begin{exercise}
设
\[
\delta(K)=\frac{\log|d_{K/\Q}|}{[K:\Q]}.
\]
若 $\gcd(d_{F/\Q},d_{K/\Q})=1$，证明
\[
d_{FK/\Q}=d_{F/\Q}^{[K:\Q]}d_{K/\Q}^{[F:\Q]},
\qquad
\delta(FK)=\delta(F)+\delta(K).
\]
\end{exercise}

\begin{proof}[解答]
状态：solvable。
判别式传递公式给出
\[
d_{FK/\Q}
=\Nm_{F/\Q}(d_{FK/F})\,d_{F/\Q}^{[FK:F]}.
\]
由于 $d_F$ 与 $d_K$ 互素，两个扩张在所有有限素数处没有共同分歧。对 $\mathcal O_F$ 的任一非零素理想局部化：若它位于 $d_F$ 上方，则 $K/\Q$ 在下面的素数处未分歧；若它位于 $d_K$ 上方，则 $F/\Q$ 在该素数处未分歧；其余地方两边都未分歧。因此相对判别式的局部指数与 $d_{K/\Q}\mathcal O_F$ 的局部指数相同，得到
\[
d_{FK/F}=d_{K/\Q}\mathcal O_F.
\]
因此
\[
\Nm_{F/\Q}(d_{FK/F})=d_{K/\Q}^{[F:\Q]}.
\]
同样由判别式互素推出 $F$ 与 $K$ 在 $\Q$ 上线性无交，所以 $[FK:F]=[K:\Q]$，得到第一式。取绝对值、对数并除以
\[
[FK:\Q]=[F:\Q][K:\Q]
\]
即得
\[
\delta(FK)
=\frac{[K:\Q]\log|d_F|+[F:\Q]\log|d_K|}
{[F:\Q][K:\Q]}
=\delta(F)+\delta(K).
\]
\end{proof}

\begin{exercise}
设 $M,N$ 为 $R$-模。证明
\[
\bigwedge^k(M\oplus N)\simeq
\bigoplus_{j=0}^k
\left(\bigwedge^jM\right)\otimes_R
\left(\bigwedge^{k-j}N\right).
\]
\end{exercise}

\begin{proof}[解答]
状态：solvable。
定义从右到左的映射：对纯张量
\[
(m_1\wedge\cdots\wedge m_j)\otimes(n_1\wedge\cdots\wedge n_{k-j})
\]
送到
\[
(m_1,0)\wedge\cdots\wedge(m_j,0)\wedge
(0,n_1)\wedge\cdots\wedge(0,n_{k-j}).
\]
它满足交替关系，所以良定。

反向映射由展开给出。任一生成元
\[
(m_1,n_1)\wedge\cdots\wedge(m_k,n_k)
\]
把每个 $(m_i,n_i)$ 写成 $(m_i,0)+(0,n_i)$ 并按多重线性展开。含有 $j$ 个 $M$-项、$k-j$ 个 $N$-项的部分放入第 $j$ 个直和分量；为把 $M$-项移到前面，需要加入相应置换符号。两边复合在纯生成元上都是恒等，因此为自然同构。
\end{proof}

\begin{exercise}
设 $\mathcal O$ 为 Dedekind 环，$\mathfrak c$ 为非零理想。用中国剩余定理证明 $\mathcal O/\mathfrak c$ 的任意理想是主理想。于是若 $\mathfrak a\supset\mathfrak c$，则存在 $y\in\mathcal O$ 使
\[
\mathfrak a=y\mathcal O+\mathfrak c.
\]
\end{exercise}

\begin{proof}[解答]
状态：solvable。
分解
\[
\mathfrak c=\prod_i\mathfrak p_i^{n_i}.
\]
中国剩余定理给出
\[
\mathcal O/\mathfrak c\simeq
\prod_i\mathcal O/\mathfrak p_i^{n_i}.
\]
因此只需证明每个局部因子 $\mathcal O/\mathfrak p^n$ 的理想都是主理想。在 Dedekind 环局部化 $\mathcal O_{\mathfrak p}$ 中，$\mathfrak p\mathcal O_{\mathfrak p}$ 是主理想，由一个一致化元 $\pi$ 生成。于是
\[
\mathcal O/\mathfrak p^n
\]
的理想正是
\[
(\pi^r)/(\pi^n),\qquad 0\le r\le n,
\]
都是主理想。有限直积中理想逐分量为主，并可用中国剩余定理把各分量生成元拼成一个元素，所以 $\mathcal O/\mathfrak c$ 的任意理想主。

若 $\mathfrak a\supset\mathfrak c$，则 $\mathfrak a/\mathfrak c$ 是 $\mathcal O/\mathfrak c$ 的理想，故由某个 $y+\mathfrak c$ 生成。拉回 $\mathcal O$ 得
\[
\mathfrak a=y\mathcal O+\mathfrak c.
\]
\end{proof}

\begin{exercise}
设 $d$ 为若干互不相同素数的乘积，并按题中方式取虚二次域整数环 $\mathcal O$。证明相关 Dedekind、类群和 $K_0$ 断言。
\end{exercise}

\begin{proof}[解答]
状态：solvable。
题中给出的 $\mathcal O$ 正是虚二次域 $\Q(\sqrt{-d})$ 的整数环：当 $-d\equiv2,3\pmod4$ 时为 $\Z[\sqrt{-d}]$，当 $-d\equiv1\pmod4$ 时为 $\Z[(1+\sqrt{-d})/2]$。数域整数环都是 Dedekind 环，因此 $\mathcal O$ 是 Dedekind 环，分式域为 $\Q(\sqrt{-d})$。

若 $d=3$，则 $\mathcal O=\Z[\omega]$。第 16 题已证明它是 Euclidean 环，因此是 PID，故理想类群为零群。由本章定理
\[
K_0(\mathcal O)\simeq\Z\oplus\Cl(\mathcal O),
\]
得到
\[
K_0(\mathcal O)\simeq\Z.
\]
但子环 $R=\Z[\sqrt{-3}]$ 不是整闭的，因为
\[
\omega=\frac{-1+\sqrt{-3}}2
\]
满足首一整系数多项式 $X^2+X+1=0$，却不在 $R$ 中。Dedekind 环必须整闭，所以 $R$ 不是 Dedekind 环。

若 $d=5$，则 $\mathcal O=\Z[\sqrt{-5}]$。它不是 PID，因为
\[
6=2\cdot3=(1+\sqrt{-5})(1-\sqrt{-5})
\]
给出唯一分解失败：四个因子都不是单位，且范数分别为 $4,9,6,6$，不可能通过单位相伴把左边两个因子配成右边两个因子。另一方面可用 Minkowski 界证明每个理想类都有范数不超过
\[
\frac{2}{\pi}\sqrt{|d_{\mathcal O}|}
=\frac{2}{\pi}\sqrt{20}<3
\]
的代表。因此每个理想类可由范数 $1$ 或 $2$ 的理想代表。范数 $1$ 是主类；范数 $2$ 的素理想只有
\[
\mathfrak p=(2,1+\sqrt{-5}),
\]
因为模掉该理想后有 $2=0$ 且 $\sqrt{-5}\equiv-1$，商环同构于 $\F_2$。计算
\[
\mathfrak p^2=(2),
\]
更精确地，$\mathfrak p^2$ 与 $(2)$ 都是范数 $4$ 的理想且有包含
$\mathfrak p^2\subset(2)$，所以相等。因此 $[\mathfrak p]^2=1$。如果
$\mathfrak p$ 是主理想，设 $\mathfrak p=(a+b\sqrt{-5})$，则
$N(a+b\sqrt{-5})=2$，即 $a^2+5b^2=2$，无整数解。因此
$\mathfrak p$ 非主，类群有两个元素，故
\[
\Cl(\mathcal O)\simeq\Z/2\Z.
\]
于是
\[
K_0(\mathcal O)\simeq\Z\oplus\Z/2\Z.
\]
\end{proof}


 \chapter{格}

本章把数域放进实向量空间来研究。设 $[F:\Q]=n=r_1+2r_2$，其中 $r_1$ 是实嵌入个数，$r_2$ 是共轭复嵌入对数。核心空间是
\[
F_\R=F\otimes_\Q\R\simeq \R^{r_1}\times\C^{r_2}.
\]
本章有三条主线。第一，把 $\OF$ 的理想嵌入 $F_\R$，把它们看成格，然后用体积估计找短向量。第二，用对数映射把单位群放进一个实超平面，得到 Dirichlet 单位定理。第三，用格点计数和 Dirichlet 级数，得到类数公式的残数形式、若干密度结果以及类域论第一不等式。

这一章的思想转换很大：第一章主要在环和理想中做代数分解；第二章开始把同一个对象看成几何图形和解析函数。读的时候不要把“格”“体积”“$L$ 函数”当成外来工具，它们是在回答同一个问题：理想、单位和素理想到底有多少，分布得有多密。

\section{Minkowski 理论}

Minkowski 理论是几何数论的入口。它的基本直觉是：一个格把空间切成许多同样大小的基本区；如果一个对称凸集体积大到超过某个阈值，它不可能完全避开非零格点。

\begin{definition}
设 $V$ 为有限维 $\R$-向量空间，$\Gamma$ 为 $V$ 的加法子群。
\begin{enumerate}[label=(\arabic*)]
\item 若对每个 $\gamma\in\Gamma$，存在 $\gamma$ 的邻域 $U_\gamma$ 使
\[
U_\gamma\cap\Gamma=\{\gamma\},
\]
则称 $\Gamma$ 为离散子群。商空间 $V/\Gamma$ 称为 $\Gamma$ 的轨迹空间。
\item 若存在 $V$ 中线性无关向量 $v_1,\dots,v_m$，使
\[
\Gamma=\Z v_1+\cdots+\Z v_m,
\]
则称 $\Gamma$ 为格。当 $m=\dim_\R V$ 时，称为全格。
\end{enumerate}
\end{definition}

\begin{remark}
“离散”说的是点之间没有无限挤在一起；“格”说的是它由有限个独立方向的整数线性组合生成。命题 2.3 会说明，在有限维实向量空间中，这两种说法本质上等价。
\end{remark}

\begin{lemma}
设 $\Gamma$ 为 $V$ 的格。则 $\Gamma$ 为全格，当且仅当存在有界子集 $D\subset V$，使
\[
V=\bigcup_{\gamma\in\Gamma}(D+\gamma).
\]
\end{lemma}

\begin{proof}
若 $\Gamma$ 为全格，取 $V$ 的基 $v_1,\dots,v_n$ 使
\[
\Gamma=\Z v_1+\cdots+\Z v_n.
\]
令
\[
D=\left\{\sum_{i=1}^n t_iv_i:0\le t_i<1\right\}.
\]
任意 $v=\sum x_iv_i$，写 $x_i=m_i+t_i$，其中 $m_i\in\Z$，$0\le t_i<1$。则
\[
v=\sum m_iv_i+\sum t_iv_i\in \Gamma+D.
\]
所以这些平移覆盖 $V$，且 $D$ 有界。

反过来，设存在有界 $D$ 使 $V=\bigcup_{\gamma\in\Gamma}(D+\gamma)$。令 $W$ 为 $\Gamma$ 张成的实线性子空间。任取 $v\in V$。对每个正整数 $k$，有
\[
kv=d_k+\gamma_k,\qquad d_k\in D,\ \gamma_k\in\Gamma\subset W.
\]
于是
\[
v=\frac{d_k}{k}+\frac{\gamma_k}{k}.
\]
因为 $D$ 有界，$d_k/k\to0$。又 $W$ 是有限维子空间，故闭合，取极限得 $v\in W$。所以 $W=V$，即 $\Gamma$ 的秩等于 $\dim V$，为全格。
\end{proof}

\begin{proposition}
设 $\Gamma$ 为有限维实向量空间 $V$ 的子群。
\begin{enumerate}[label=(\arabic*)]
\item $\Gamma$ 是离散子群，当且仅当 $\Gamma$ 是格。
\item $V/\Gamma$ 紧，当且仅当 $\Gamma$ 是全格。
\end{enumerate}
\end{proposition}

\begin{proof}
若 $\Gamma=\Z v_1+\cdots+\Z v_m$ 且 $v_i$ 线性无关，则在
$W=\sum\R v_i$ 中它与 $\Z^m\subset\R^m$ 同构。标准整数格在原点附近有不含其他格点的小球；平移这个小球即可分离每个格点。把 $W$ 嵌入
$V$ 后，取 $V=W\oplus W'$ 的任一补空间，乘积邻域仍能把格点分开，所以 $\Gamma$ 离散。

反过来，设 $\Gamma$ 离散。先取 $\Gamma$ 中一组极大实线性无关元
$v_1,\dots,v_m$，令 $\Gamma_0=\Z v_1+\cdots+\Z v_m$，$W=\sum_i\R v_i$。极大性说明
$\Gamma\subset W$。考虑由 $v_i$ 张成的半开平行体
\[
P=\left\{\sum_{i=1}^m t_iv_i:0\le t_i<1\right\}.
\]
任意 $\gamma\in\Gamma$ 都可唯一写成 $\gamma=\gamma_0+p$，其中
$\gamma_0\in\Gamma_0$，$p\in P$。于是 $\Gamma/\Gamma_0$ 的每个陪集在有界集
$P$ 中有代表元。因为 $\Gamma$ 离散，$\Gamma\cap\overline P$ 是紧集中的离散子集，故有限。因此 $\Gamma/\Gamma_0$ 有限。取这些有限代表元
$p_1,\dots,p_s$。每个 $p_j$ 在 $W$ 中有有理系数：若某个系数无理，则
$kp_j$ 模 $\Gamma_0$ 的小数部分会给出无限多个落在 $P$ 中的不同
$\Gamma$ 元素，和有限性矛盾。取公共分母 $N$，则
\[
\Gamma\subset \frac1N\Gamma_0.
\]
所以 $\Gamma$ 是有限生成自由 abelian 群，并由 $W$ 中线性无关方向生成，即为格。

对 (2)，若 $\Gamma$ 为全格，前一引理给出有界基本区 $D$，使
$V=D+\Gamma$。取闭包 $\overline D$，它在有限维实向量空间中紧；商映射
$q:V\to V/\Gamma$ 连续，且 $q(\overline D)=V/\Gamma$，所以商空间紧。

反过来，若 $V/\Gamma$ 紧，取原点的一个有界开邻域 $U$。商映射的像
$q(U+\gamma)=q(U)$ 随 $\gamma$ 不变，而 $q(U)$ 的平移覆盖商空间。由紧性，可取有限多个点
$v_1,\dots,v_r\in V$ 使
\[
V/\Gamma=\bigcup_{i=1}^r q(v_i+U).
\]
令 $D=\bigcup_i(v_i+\overline U)$，这是有界集。任意 $v\in V$ 的商类落在某个
$q(v_i+U)$ 中，于是 $v\in v_i+U+\Gamma\subset D+\Gamma$。因此
$V=D+\Gamma$，由前一引理 $\Gamma$ 是全格。
\end{proof}

若 $\Gamma=\Z v_1+\cdots+\Z v_n$ 为 $\R^n$ 中全格，基本平行体
\[
P(v_1,\dots,v_n)=\left\{\sum_{i=1}^n t_iv_i:0\le t_i<1\right\}
\]
的体积为
\[
|\det(v_1,\dots,v_n)|.
\]
若换另一组 $\Z$-基，换基矩阵属于 $\GL_n(\Z)$，行列式为 $\pm1$，所以体积不变。记这个数为 $v(\Gamma)$，称为格的基本区体积或协体积。

\begin{theorem}
设 $\Gamma$ 为 $n$ 维实向量空间 $V$ 的全格。设 $X\subset V$ 满足：若 $x,y\in X$，则 $(x-y)/2\in X$。若
\[
\mu(X)>2^n v(\Gamma),
\]
则 $X\cap\Gamma$ 含有非零点。
\end{theorem}

\begin{proof}
考虑商映射 $q:V\to V/2\Gamma$。格 $2\Gamma$ 的基本区体积为
$2^n v(\Gamma)$。取 $2\Gamma$ 的一个基本平行体 $P$。若 $q$ 在 $X$ 上单射，则可把
$X$ 中每个点唯一平移到 $P$ 中而不改变体积，因此 $\mu(X)\le\mu(P)=2^n v(\Gamma)$。现在
$\mu(X)>2^n v(\Gamma)$，所以 $q$ 在 $X$ 上不能单射；存在不同的 $x,y\in X$，使
\[
x-y\in2\Gamma.
\]
写 $x-y=2\gamma$，$\gamma\in\Gamma$。因为 $x\ne y$，所以 $\gamma\ne0$。由 $X$ 的假设，
\[
\gamma=\frac{x-y}{2}\in X.
\]
故 $X$ 中有非零格点。
\end{proof}

\begin{remark}
最常用的情形是 $X$ 为关于原点对称的凸集。若 $x,y\in X$，则 $x$ 与 $-y$ 都在 $X$，凸性给出 $(x-y)/2=(x+(-y))/2\in X$。因此本定理就是常见的 Minkowski 对称凸体定理。
\end{remark}

\section{加性结构}

本节把数域的加法结构嵌入到实向量空间。设 $F/\Q$ 为有限扩张，$[F:\Q]=n=r_1+2r_2$。取全部嵌入，按
\[
\iota_1,\dots,\iota_{r_1},
\iota_{r_1+1},\dots,\iota_{r_1+r_2},
\overline{\iota}_{r_1+1},\dots,\overline{\iota}_{r_1+r_2}
\]
排列，其中前 $r_1$ 个为实嵌入，后面是复嵌入与其共轭。定义
\[
F_\R=\R^{r_1}\times\C^{r_2},
\]
并定义 Minkowski 嵌入
\[
\iota:F\longrightarrow F_\R,\qquad
a\longmapsto
(\iota_1(a),\dots,\iota_{r_1}(a),
\iota_{r_1+1}(a),\dots,\iota_{r_1+r_2}(a)).
\]

\begin{proposition}
\begin{enumerate}[label=(\arabic*)]
\item 映射
\[
F\otimes_\Q\R\longrightarrow F_\infty,\qquad
a\otimes x\longmapsto x(\iota_j(a))_j
\]
给出 $n$ 维实向量空间同构。
\item $F_\infty$ 与 $F_\R=\R^{r_1}\times\C^{r_2}$ 同构，且 $\iota:F\to F_\R$ 是加法群单射。
\end{enumerate}
\end{proposition}

\begin{proof}
把 $F$ 放进一个包含它的有限 Galois 扩张 $E/\Q$，则从 $F$ 到 $\C$ 的嵌入对应于
$\Gal(E/\Q)$ 对固定 $F$ 的子群的左陪集。不同陪集给出不同嵌入，总数为
$[F:\Q]=n$。这些嵌入把
$F\otimes_\Q\R$ 分解为
\[
\R^{r_1}\times\C^{r_2}
\]
作为实代数的直积：实嵌入贡献一个 $\R$ 因子，一对共轭复嵌入贡献一个
$\C$ 因子。维数为 $r_1+2r_2=n$。若 $\iota(a)=0$，则所有嵌入值为 $0$，特别最小多项式的所有共轭中有 $0$，故 $a=0$；所以 $\iota$ 为单射。
\end{proof}

\begin{theorem}
设 $\mathfrak a$ 为 $\OF$ 的非零理想。则 $\iota(\mathfrak a)$ 是 $F_\R$ 的全格，且其基本区体积为
\[
v(\iota(\mathfrak a))
=2^{-r_2}\Nm(\mathfrak a)\sqrt{|d(\OF)|}.
\]
\end{theorem}

\begin{proof}
取 $\mathfrak a$ 的 $\Z$-基 $\alpha_1,\dots,\alpha_n$。它们也是 $F$ 的 $\Q$-基，所以嵌入后的向量生成 $F_\R$ 的全格。写实坐标矩阵 $M$ 的第 $i$ 行为
\[
\iota_1(\alpha_i),\dots,\iota_{r_1}(\alpha_i),
\Re\iota_{r_1+1}(\alpha_i),\Im\iota_{r_1+1}(\alpha_i),\dots,
\Re\iota_{r_1+r_2}(\alpha_i),\Im\iota_{r_1+r_2}(\alpha_i).
\]
则
\[
v(\iota(\mathfrak a))=|\det M|.
\]

再取复矩阵 $D$，其第 $i$ 行为
\[
\iota_1(\alpha_i),\dots,\iota_{r_1}(\alpha_i),
\iota_{r_1+1}(\alpha_i),\overline{\iota}_{r_1+1}(\alpha_i),\dots,
\iota_{r_1+r_2}(\alpha_i),\overline{\iota}_{r_1+r_2}(\alpha_i).
\]
每一对复列 $(z,\bar z)$ 与实列 $(\Re z,\Im z)$ 的换算行列式贡献绝对值 $2$，因此
\[
|\det M|=2^{-r_2}|\det D|.
\]
而
\[
(\det D)^2=\det(\Tr_{F/\Q}(\alpha_i\alpha_j))=:d(\mathfrak a).
\]
所以
\[
v(\iota(\mathfrak a))=2^{-r_2}\sqrt{|d(\mathfrak a)|}.
\]

若 $\mathfrak a\subset\OF$，则换基指数给出
\[
d(\mathfrak a)=(\OF:\mathfrak a)^2d(\OF)
=\Nm(\mathfrak a)^2d(\OF).
\]
代入即得公式。
\end{proof}

\begin{proposition}
设 $\mathfrak a$ 为 $\OF$ 的非零理想。
\begin{enumerate}[label=(\arabic*)]
\item 存在 $0\ne a\in\mathfrak a$，使
\[
|\Nm_{F/\Q}(a)|
\le
\frac{n!}{n^n}\left(\frac4\pi\right)^{r_2}
\Nm(\mathfrak a)\sqrt{|d(\OF)|}.
\]
\item 每个理想类 $A\in\Cl_F$ 中存在整理想 $\mathfrak b\in A$，使
\[
\Nm(\mathfrak b)
\le
\frac{n!}{n^n}\left(\frac4\pi\right)^{r_2}
\sqrt{|d(\OF)|}.
\]
\end{enumerate}
\end{proposition}

\begin{proof}
(1) 在 $F_\R$ 中取集合
\[
X_t=\left\{(x_i,z_j):\sum_{i=1}^{r_1}|x_i|
 +2\sum_{j=1}^{r_2}|z_j|\le t\right\}.
\]
它是对称凸集。下面计算它的体积。令 $y_i=|x_i|$，$s_j=2|z_j|$。每个实坐标给出符号因子 $2$；每个复坐标用极坐标，面积元为
$2\pi\rho\,d\rho$，而 $s=2\rho$，所以贡献
$(\pi/2)s\,ds$。于是
\[
\mu(X_t)
=2^{r_1}\left(\frac{\pi}{2}\right)^{r_2}
\int_{\sum y_i+\sum s_j\le t}\prod_{j=1}^{r_2}s_j\,dy\,ds.
\]
单纯形积分给出
\[
\int_{\sum y_i+\sum s_j\le t}\prod_{j=1}^{r_2}s_j\,dy\,ds
=\frac{t^n}{n!},
\]
因为 $r_1$ 个变量指数为 $0$，$r_2$ 个变量指数为 $1$，总次数为
$n=r_1+2r_2$。因此
\[
\mu(X_t)=2^{r_1-r_2}\pi^{r_2}\frac{t^n}{n!}.
\]
选择 $t$ 使
\[
\mu(X_t)>2^n v(\iota(\mathfrak a)).
\]
Minkowski 定理给出 $0\ne a\in\mathfrak a$ 且 $\iota(a)\in X_t$。再用算术-几何平均不等式估计
\[
|\Nm_{F/\Q}(a)|
=\prod_{i=1}^{r_1}|\iota_i(a)|
\prod_{j=1}^{r_2}|\iota_{r_1+j}(a)|^2
\]
。这里一共有 $n$ 个非负因子：每个实嵌入给出一个
$|\iota_i(a)|$，每个复嵌入给出两个相同因子
$|\iota_{r_1+j}(a)|$。这些因子的和为
\[
\sum_i|\iota_i(a)|+2\sum_j|\iota_{r_1+j}(a)|\le t.
\]
故
\[
|\Nm_{F/\Q}(a)|\le\left(\frac tn\right)^n.
\]
由体积不等式和
$v(\iota(\mathfrak a))=2^{-r_2}\Nm(\mathfrak a)\sqrt{|d(\OF)|}$ 得
\[
t^n>
\left(\frac4\pi\right)^{r_2}n!\,\Nm(\mathfrak a)\sqrt{|d(\OF)|}.
\]
令 $t$ 向临界值下降，得到
\[
|\Nm_{F/\Q}(a)|
\le
\frac{n!}{n^n}\left(\frac4\pi\right)^{r_2}
\Nm(\mathfrak a)\sqrt{|d(\OF)|}.
\]

(2) 对任一理想类 $A$，取 $\mathfrak a^{-1}\in A^{-1}$ 为整理想。由 (1) 取
$0\ne a\in\mathfrak a^{-1}$，满足上述界。令
\[
\mathfrak b=a(\mathfrak a^{-1})^{-1}=a\mathfrak a.
\]
则 $\mathfrak b$ 是整理想，且属于类 $A$。范数满足
\[
\Nm(\mathfrak b)=\frac{|\Nm_{F/\Q}(a)|}{\Nm(\mathfrak a^{-1})},
\]
代入 (1) 得所需界。
\end{proof}

\begin{theorem}
理想类群 $\Cl_F$ 是有限群。
\end{theorem}

\begin{proof}
上一命题说明每个理想类都有一个范数不超过常数
\[
C_F=\frac{n!}{n^n}\left(\frac4\pi\right)^{r_2}\sqrt{|d(\OF)|}
\]
的整理想代表。对固定整数 $N$，$\OF$ 中范数为 $N$ 的理想只有有限个，因为它们都含有 $N\OF$，对应于有限环 $\OF/N\OF$ 的理想。$N\le C_F$ 只有有限多个可能，所以理想类只有有限多个。
\end{proof}

\section{乘性结构}

加性结构把理想变成格；乘性结构把单位变成对数空间中的格。设
\[
F_\R^\times=(\R^\times)^{r_1}\times(\C^\times)^{r_2}.
\]
定义
\[
\ell:F_\R^\times\to\R^{r_1+r_2}
\]
为
\[
(x_1,\dots,x_{r_1},z_1,\dots,z_{r_2})
\longmapsto
(\log|x_1|,\dots,\log|x_{r_1}|,
2\log|z_1|,\dots,2\log|z_{r_2}|).
\]
对 $a\in F^\times$，记
\[
\lambda(a)=\ell(\iota(a)).
\]

\begin{theorem}
令 $U_F=\OF^\times$。则 $\lambda(U_F)$ 是超平面
\[
H=\left\{(y_1,\dots,y_{r_1+r_2}):y_1+\cdots+y_{r_1+r_2}=0\right\}
\]
中的全格。
\end{theorem}

\begin{proof}
若 $u\in U_F$，则 $\Nm_{F/\Q}(u)=\pm1$，所以
\[
\sum_i \lambda(u)_i=\log|\Nm_{F/\Q}(u)|=0.
\]
故像落在 $H$ 中。离散性来自如下事实：若 $\lambda(u)$ 被限制在一个有界区域内，则所有嵌入 $|\iota(u)|$ 都有上下界；于是 $\iota(u)$ 落在 $F_\R$ 中一个有界集与格 $\iota(\OF)$ 的交中，只可能有有限多个。故 $\lambda(U_F)$ 离散。

满秩性是 Dirichlet 单位定理的核心部分。证明思想是用 Minkowski 理论构造足够多独立单位：在不同无穷嵌入方向上寻找小元素，再取主理想比值，使有限理想部分抵消，留下无穷位对数方向。由此得到 $r_1+r_2-1$ 个线性无关的对数向量，并且它们生成一个使商紧的格。
\end{proof}

\begin{definition}
$F$ 中所有单位根组成的有限群记为 $\mu_F$。它是循环群：因为 $\mu_F$ 是
$\C^\times$ 的有限子群，而复数乘法群的有限子群都由某个本原单位根生成。
\end{definition}

\begin{theorem}
整数环单位群分解为
\[
U_F\simeq \mu_F\times\Z^{r_1+r_2-1}.
\]
\end{theorem}

\begin{proof}
核 $\ker\lambda$ 由所有嵌入绝对值均为 $1$ 的单位组成。它落在有界集与格
$\iota(\OF)$ 的交中，所以有限。若 $u\in\ker\lambda$，则 $u$ 的所有共轭绝对值都为
$1$；它的幂 $u^k$ 也都落在同一个有界集与格的交中，故集合
$\{u^k:k\ge0\}$ 有重复，$u$ 为有限阶元素，即单位根。反过来单位根在所有嵌入下绝对值都为 $1$。因此核为
$\mu_F$。上一条定理说明 $\lambda(U_F)$ 是秩 $r_1+r_2-1$ 的自由交换群。短正合列
\[
1\to\mu_F\to U_F\to \lambda(U_F)\to0
\]
在交换群意义下分裂，于是得到所需直积。
\end{proof}

若 $\mathfrak m=\mathfrak m_\infty\mathfrak m_0$ 为模数，记 $F_{\mathfrak m,1}$ 为满足射线同余 $x\equiv1\pmod{\mathfrak m}$ 的元素，记 $\mu_F^{\mathfrak m}=\mu_F\cap F_{\mathfrak m,1}$。

\begin{theorem}
$U_F\cap F_{\mathfrak m,1}$ 是有限群 $\mu_F^{\mathfrak m}$ 与一个秩 $r_1+r_2-1$ 自由交换群的直积。
\end{theorem}

\begin{proof}
先说明 $U_F\cap F_{\mathfrak m,1}$ 在 $U_F$ 中有限指数。有限部分
$\mathfrak m_0$ 给出映射
\[
U_F\longrightarrow(\OF/\mathfrak m_0)^\times,
\]
其核是满足 $u\equiv1\pmod{\mathfrak m_0}$ 的单位；目标是有限群，所以这个核有限指数。无穷部分只要求若干实嵌入为正，这是到
\[
\{\pm1\}^{r_1}
\]
的符号映射再取核，也只改变有限指数。两种条件合在一起，仍得到
$U_F$ 的有限指数子群。

由 Dirichlet 单位定理，$U_F\simeq\mu_F\times\Z^{r_1+r_2-1}$。有限生成阿贝尔群的有限指数子群与原群有相同自由秩，因此
$U_F\cap F_{\mathfrak m,1}$ 的自由秩仍为 $r_1+r_2-1$。它的扭元素正是同时属于
$\mu_F$ 且满足射线同余条件的单位根，即 $\mu_F^{\mathfrak m}$。于是得到所述直积分解。
\end{proof}

\section{理想估值}

本节的任务是数理想：固定射线理想类 $K\in I_F^{\mathfrak m}/S_F^{\mathfrak m}$，估计该类中范数不超过 $x$ 的整理想个数。思想是把理想计数转化为格点计数，再除以单位群的重复作用。

设
\[
\mathfrak m=\mathfrak m_\infty\mathfrak m_0,\qquad
\mathfrak m_\infty=(\mathfrak p_{\infty,1})^{\nu_1}\cdots
(\mathfrak p_{\infty,r_1})^{\nu_{r_1}},
\]
其中前 $s_{\mathfrak m}$ 个 $\nu_j=1$，其余为 $0$。定义
\[
\Nm(\mathfrak m)=2^{s_{\mathfrak m}}\Nm(\mathfrak m_0).
\]

\begin{lemma}[整理想代表]
每个类 $K\in I_F^{\mathfrak m}/S_F^{\mathfrak m}$ 中都有整理想。
\end{lemma}

\begin{proof}
取类中任意分式理想 $\mathfrak b$。先选 $0\ne c\in\OF$ 使
$c\mathfrak b\subset\OF$，这样 $c\mathfrak b$ 是整理想。还要保证它仍在
$I_F^{\mathfrak m}$ 中，也就是与 $\mathfrak m_0$ 互素。把
$\mathfrak b$ 写成素理想分解；只需在选择 $c$ 时把 $\mathfrak m_0$ 中出现的素理想指数抵消到零，并且在其他素理想处取足够大的非负指数，便得到一个与
$\mathfrak m_0$ 互素的整理想。这个过程只乘了一个主分式理想，所以没有离开
$I_F^{\mathfrak m}/S_F^{\mathfrak m}$ 中给定的类。
\end{proof}

\begin{lemma}[从理想计数到元素计数]
令 $P_x$ 表示类 $K$ 内满足 $\Nm(\mathfrak n)\le x$ 的整理想个数。取整理想
$\mathfrak a\in K^{-1}$。则 $P_x$ 等于满足以下条件的主理想 $(\nu)$ 的个数：
\[
(\nu)\in S_F^{\mathfrak m},\qquad
\nu\in\mathfrak a,\qquad
0<|\Nm_{F/\Q}(\nu)|\le x\Nm(\mathfrak a).
\]
\end{lemma}

\begin{proof}
若 $\mathfrak n\in K$，则 $\mathfrak n\mathfrak a$ 属于
$K\cdot K^{-1}$，即射线主理想类。因此存在
$\nu\in F_{\mathfrak m,1}$ 使
\[
\mathfrak n\mathfrak a=(\nu)\in S_F^{\mathfrak m}.
\]
由于 $\mathfrak n$ 是整理想，$(\nu)\subset\mathfrak a$，故
$\nu\in\mathfrak a$。范数条件变为
\[
0<|\Nm_{F/\Q}(\nu)|\le x\Nm(\mathfrak a).
\]
反过来，若 $\nu$ 满足这些条件，则
\[
\mathfrak n=(\nu)\mathfrak a^{-1}
\]
是整理想，属于 $K$，且范数不超过 $x$。两个元素给出同一个主理想，当且仅当相差一个单位；后面会把这个重复由
$U_F\cap F_{\mathfrak m,1}$ 除掉。
\end{proof}

\begin{lemma}[同余条件的平移格形式]
存在 $\nu_0\in\mathfrak a$，使 $\nu\in\mathfrak a\cap F_{\mathfrak m,1}$ 等价于
\[
\nu\equiv\nu_0\pmod{\mathfrak m_0\mathfrak a}
\]
并满足实嵌入符号条件
$\nu^{(1)},\dots,\nu^{(s_{\mathfrak m})}>0$。
\end{lemma}

\begin{proof}
有限同余条件是
$\nu\equiv1\pmod{\mathfrak m_0}$ 与 $\nu\in\mathfrak a$。因为
$\mathfrak a$ 选自 $I_F^{\mathfrak m}$，它与 $\mathfrak m_0$ 互素，中国剩余定理可在
$\OF$ 中同时解这两个条件，得到一个代表元 $\nu_0$。于是所有解正是
$\nu_0+\mathfrak m_0\mathfrak a$。无穷部分不是同余理想，而是符号条件：若
$\mathfrak p_{\infty,j}$ 出现在 $\mathfrak m_\infty$ 中，则第 $j$ 个实嵌入下必须为正。
\end{proof}

\begin{lemma}[平移格点计数]
设 $S\subset\R^k$ 为有界 Jordan 可测集，$L_v=v+\Z^k$ 为 $\Z^k$ 的一个平移。令
$M(t)=|tS\cap L_v|$。则
\[
\lim_{t\to\infty}\frac{M(t)}{t^k}=\mu(S),
\]
其中 $\mu(S)$ 表示 $S$ 的 $k$ 维体积。
\end{lemma}

\begin{proof}
把问题缩小到网格间距趋向于零的情形。点 $z\in tS\cap L_v$ 等价于
$z/t\in S\cap (t^{-1}L_v)$。以这些点为中心放边长 $1/t$ 的小立方体。完全落在
$S$ 内的小立方体个数记为 $T_1(t)$，与 $S$ 相交的小立方体个数记为
$T_2(t)$。由于 $S$ 的边界体积为零，内近似和外近似的体积都趋向于
$\mu(S)$。并且
\[
T_1(t)\le M(t)\le T_2(t),\qquad
t^{-k}T_1(t)\le \mu(S)\le t^{-k}T_2(t)+o(1).
\]
夹逼后得到所需极限。平移向量 $v$ 只移动网格的相位，不改变网格基本小方块的体积，所以极限仍是同一个体积。
\end{proof}

\begin{definition}
取 $U_F\cap F_{\mathfrak m,1}$ 的自由部分基 $w_1,\dots,w_{r_1+r_2-1}$。令
\[
W_i=\lambda(w_i),\qquad
W=(1,\dots,1,2,\dots,2)\in\R^{r_1+r_2},
\]
其中 $W$ 有 $r_1$ 个 $1$ 和 $r_2$ 个 $2$。模 $\mathfrak m$ 的调控子定义为
\[
R_{\mathfrak m}
=\frac{2^{r_2}}{r_1+2r_2}
\left|\det(W,W_1,\dots,W_{r_1+r_2-1})\right|.
\]
记 $w_{\mathfrak m}=|\mu_F^{\mathfrak m}|$。
\end{definition}

\begin{remark}
这里的向量 $W$ 记录范数的权重。若
$y=(y_1,\dots,y_{r_1},z_1,\dots,z_{r_2})\in F_\R^\times$，则
\[
\log |N_{F_\R/\R}(y)|
=\sum_{i=1}^{r_1}\log |y_i|
 +2\sum_{j=1}^{r_2}\log |z_j|
 = T(\lambda(y)),
\]
其中 $T$ 是与 $W$ 做内积的线性泛函。单位的范数为 $\pm1$，所以
\(\lambda(U_F\cap F_{\mathfrak m,1})\) 落在 $T=0$ 的超平面中。
\end{remark}

\begin{lemma}[单位基本域与理想计数]
设 $\alpha_1,\dots,\alpha_n$ 是 $\mathfrak m_0\mathfrak a$ 的 $\Z$-基，其中
$n=r_1+2r_2$。记 $S$ 为所有 $x=(x_1,\dots,x_n)\in\R^n$ 组成的集合，满足
\[
\nu(x)=\sum_{i=1}^n x_i\alpha_i
\]
在 $F_\R$ 中非零，并且
\begin{enumerate}[label=(\arabic*)]
\item $\lambda(\nu(x))=cW+c_1W_1+\cdots+c_{r_1+r_2-1}W_{r_1+r_2-1}$，其中 $c\in\R$ 且 $0\le c_i<1$；
\item 前 $s_{\mathfrak m}$ 个实嵌入满足 $\nu(x)^{(j)}>0$；
\item $0<|N_{F/\Q}(\nu(x))|\le 1$。
\end{enumerate}
则
\[
\lim_{x\to\infty}\frac{P_x}{x}
=\frac{\mu(S)\Nm(\mathfrak a)}{w_{\mathfrak m}}.
\]
\end{lemma}

\begin{proof}
先把同余条件写成平移格。由前一条引理，存在
$\nu_0\in\mathfrak a$ 使所有有限同余解为
$\nu_0+\mathfrak m_0\mathfrak a$。因为
$\alpha_1,\dots,\alpha_n$ 是 $\mathfrak m_0\mathfrak a$ 的基，存在
$v_1,\dots,v_n\in\Q$ 使
\[
\nu_0=\sum_{i=1}^n v_i\alpha_i .
\]
于是 $L_v=v+\Z^n$ 中的点 $y=(y_1,\dots,y_n)$ 与元素
\[
\nu(y)=\sum_{i=1}^n y_i\alpha_i
\]
一一对应，并且这些元素正好满足有限同余条件
$\nu\equiv\nu_0\pmod{\mathfrak m_0\mathfrak a}$。

现在处理单位造成的重复。若 $\nu$ 给出一个主理想 $(\nu)$，则
$u\nu$ 给出同一个主理想，其中 $u\in U_F\cap F_{\mathfrak m,1}$。反过来，两个满足射线条件的元素给出同一个主理想时，它们的商就是
$U_F\cap F_{\mathfrak m,1}$ 中的单位。单位根
\(\mu_F^{\mathfrak m}\) 中的元素不改变对数向量，数量为 $w_{\mathfrak m}$；自由部分由
$w_1,\dots,w_{r_1+r_2-1}$ 生成。对任意满足条件的 $\nu$，把
\[
\lambda(\nu)=cW+d_1W_1+\cdots+d_{r_1+r_2-1}W_{r_1+r_2-1}
\]
写出后，选整数 $a_i$ 使 $0\le d_i+a_i<1$，并乘以
$w_1^{a_1}\cdots w_{r_1+r_2-1}^{a_{r_1+r_2-1}}$。这样每个主理想在基本带
$0\le c_i<1$ 中恰有 $w_{\mathfrak m}$ 个元素代表。

最后把范数界转化为区域膨胀。若 $t^n=x\Nm(\mathfrak a)$，则
\[
N_{F/\Q}\bigl(\nu(ty)\bigr)=t^nN_{F/\Q}(\nu(y)).
\]
因此 $tS\cap L_v$ 的点数正是 $w_{\mathfrak m}P_x$。由平移格点计数引理，
\[
\mu(S)=\lim_{t\to\infty}\frac{|tS\cap L_v|}{t^n}
=\lim_{x\to\infty}\frac{w_{\mathfrak m}P_x}{x\Nm(\mathfrak a)}.
\]
整理即得公式。
\end{proof}

\begin{lemma}[计数区域的体积]
上述集合 $S$ 的体积为
\[
\mu(S)=
\frac{2^{r_1-s_{\mathfrak m}}(2\pi)^{r_2}R_{\mathfrak m}}
{\sqrt{|d(\OF)|}\,\Nm(\mathfrak m_0\mathfrak a)}.
\]
\end{lemma}

\begin{proof}
第一步，把 $\mathfrak m_0\mathfrak a$ 的坐标变成 Minkowski 嵌入坐标。写
\[
y_i=\sum_j x_j\iota_i(\alpha_j)\quad(1\le i\le r_1),
\]
而对复嵌入写
\[
y_{r_1+j}+\sqrt{-1}\,y_{r_1+r_2+j}
=\sum_k x_k\iota_{r_1+j}(\alpha_k)\quad(1\le j\le r_2).
\]
这个线性变换的 Jacobi 行列式绝对值等于
\[
2^{-r_2}\Nm(\mathfrak m_0\mathfrak a)\sqrt{|d(\OF)|}.
\]
原因是 $\OF$ 在 Minkowski 嵌入下的协体积为
$2^{-r_2}\sqrt{|d(\OF)|}$，而把 $\OF$ 换成理想
$\mathfrak m_0\mathfrak a$ 会把协体积乘以
$\Nm(\mathfrak m_0\mathfrak a)$。有限实符号条件只要求前
$s_{\mathfrak m}$ 个实坐标为正；若先计算所有实坐标均为正的区域 $Y$，则其余
$r_1-s_{\mathfrak m}$ 个实坐标可自由取正负号，产生因子
$2^{r_1-s_{\mathfrak m}}$。

第二步，把复坐标改成极坐标。令
\[
\rho_i=y_i\quad(1\le i\le r_1),\qquad
\rho_{r_1+j}e^{\sqrt{-1}\theta_j}
=y_{r_1+j}+\sqrt{-1}\,y_{r_1+r_2+j}.
\]
角变量独立贡献 $(2\pi)^{r_2}$，径向 Jacobi 因子为
$\rho_{r_1+1}\cdots\rho_{r_1+r_2}$。设
\[
P=\rho_1\cdots\rho_{r_1}\rho_{r_1+1}^2\cdots\rho_{r_1+r_2}^2.
\]
范数条件就是 $0<P\le1$。因为每个 $W_i=\lambda(w_i)$ 的坐标加权和为零，对数条件可写成
\[
\log \rho_i=\frac{\log P}{r_1+2r_2}
 +\sum_{j=1}^{r_1+r_2-1} c_j\log |\iota_i(w_j)|
\]
其中实坐标权重为 $1$，复坐标权重为 $2$，且 $0\le c_j<1$。

第三步，计算从变量
\[
(P,c_1,\dots,c_{r_1+r_2-1})
\]
到 $\rho$ 变量的 Jacobi 行列式。对 $P$ 求导给出
\[
\frac{\partial \rho_i}{\partial P}
=\frac{\rho_i}{(r_1+2r_2)P},
\]
对 $c_j$ 求导给出
\[
\frac{\partial \rho_i}{\partial c_j}
=\rho_i\log |\iota_i(w_j)|.
\]
把这些列合在一起，行列式正是定义 $R_{\mathfrak m}$ 的那个行列式，外加
$\rho_1\cdots\rho_{r_1+r_2}/((r_1+2r_2)P)$。按定义
\[
R_{\mathfrak m}
=\frac{2^{r_2}}{r_1+2r_2}
\left|\det(W,W_1,\dots,W_{r_1+r_2-1})\right|,
\]
所以径向 Jacobi 因子与这一步的行列式相乘后成为常数
$2^{-r_2}R_{\mathfrak m}$。在区域
$0<P\le1$、$0\le c_j<1$ 上积分，体积为
$(2\pi)^{r_2}2^{-r_2}R_{\mathfrak m}$，再乘第一步的线性变换因子与实符号因子，得到
\[
\mu(S)=
\frac{2^{r_1-s_{\mathfrak m}}(2\pi)^{r_2}R_{\mathfrak m}}
{\sqrt{|d(\OF)|}\,\Nm(\mathfrak m_0\mathfrak a)}.
\]
\end{proof}

\begin{lemma}[射线类中的整理想渐近]
在上述记号下，固定射线类 $K$ 中整理想计数满足
\[
\lim_{x\to\infty}\frac{P_x}{x}
=
\frac{2^{r_1}(2\pi)^{r_2}R_{\mathfrak m}}
{\sqrt{|d(\OF)|}\,w_{\mathfrak m}\Nm(\mathfrak m)}.
\]
\end{lemma}

\begin{proof}
把前两条引理相乘即可。由单位基本域与理想计数引理，
\[
\lim_{x\to\infty}\frac{P_x}{x}
=\frac{\mu(S)\Nm(\mathfrak a)}{w_{\mathfrak m}}.
\]
代入体积公式得到
\[
\frac{\Nm(\mathfrak a)}{w_{\mathfrak m}}\cdot
\frac{2^{r_1-s_{\mathfrak m}}(2\pi)^{r_2}R_{\mathfrak m}}
{\sqrt{|d(\OF)|}\,\Nm(\mathfrak m_0\mathfrak a)}
=
\frac{2^{r_1-s_{\mathfrak m}}(2\pi)^{r_2}R_{\mathfrak m}}
{\sqrt{|d(\OF)|}\,w_{\mathfrak m}\Nm(\mathfrak m_0)}.
\]
再用 $\Nm(\mathfrak m)=2^{s_{\mathfrak m}}\Nm(\mathfrak m_0)$，分母中多出的
$2^{s_{\mathfrak m}}$ 移到整体常数里，得到所写形式。
\end{proof}

\begin{theorem}
设 $\mathfrak m$ 为 $F$ 的模，$K\in I_F^{\mathfrak m}/S_F^{\mathfrak m}$。若 $P_x$ 表示 $K$ 内范数不超过 $x$ 的整理想个数，则
\[
\lim_{x\to\infty}\frac{P_x}{x}
=
g_{\mathfrak m}
:=
\frac{2^{r_1}(2\pi)^{r_2}R_{\mathfrak m}}
{\sqrt{|d(\OF)|}\,w_{\mathfrak m}\Nm(\mathfrak m)}.
\]
\end{theorem}

\begin{remark}
重点不是记住常数，而是理解它的来源：$2^{r_1}$ 来自实符号，$(2\pi)^{r_2}$ 来自复极角，$\sqrt{|d|}$ 来自整数环格的协体积，$R_{\mathfrak m}$ 来自单位群基本区，$w_{\mathfrak m}$ 是单位根造成的重复计数。
\end{remark}

\section{\texorpdfstring{$L$}{L} 函数}

本节把上一节的理想计数翻译为 Dirichlet 级数在 $s=1$ 附近的行为。代数信息在这里进入解析函数：一个射线类里有多少个理想，会变成某个 Dirichlet 级数在 $s=1$ 的残数；有限群特征标会把不同射线类分离出来。

\begin{definition}
形如
\[
D(s)=\sum_{n=1}^\infty\frac{a_n}{n^s}
\]
的级数称为 Dirichlet 级数。若存在实数 $\sigma_0$，使级数在
$\Re(s)>\sigma_0$ 收敛并在 $\Re(s)<\sigma_0$ 发散，则称
$\sigma_0$ 为收敛横坐标。
\end{definition}

\begin{theorem}[Dirichlet 级数的基本性质]
设 $D(s)=\sum_{n\ge1}a_nn^{-s}$。
\begin{enumerate}[label=(\arabic*)]
\item 固定 $s_0$。若部分和
\[
S_m(s_0)=\sum_{n=1}^m\frac{a_n}{n^{s_0}}
\]
有界，则对任意 $\delta>0$ 与 $C>0$，级数在
\[
\{s: |s-s_0|\le C,\ \Re(s)\ge\Re(s_0)+\delta\}
\]
上一致收敛，并在 $\Re(s)>\Re(s_0)$ 上给出解析函数。
\item 存在 $\sigma_0\in[-\infty,\infty]$，使 $D(s)$ 在
$\Re(s)>\sigma_0$ 上收敛为解析函数，在 $\Re(s)<\sigma_0$ 上发散。
\item 令 $P_n=\sum_{k\le n}a_k$。若 $|P_n|\le Pn^{\sigma_1}$，其中
$P$ 为常数且 $\sigma_1>0$，则收敛横坐标满足 $\sigma_0\le\sigma_1$。
\item 若
\[
\lim_{n\to\infty}\frac{P_n}{n}=g,
\]
则
\[
\lim_{s\to1,\ \Re(s)>1}(s-1)D(s)=g.
\]
\end{enumerate}
\end{theorem}

\begin{proof}
先证明 Abel 分部求和公式。记 $P_0=0$，则 $a_v=P_v-P_{v-1}$，所以
\[
\sum_{v=N}^{M}a_vv^{-s}
=P_MM^{-s}-P_{N-1}N^{-s}
+\sum_{v=N}^{M-1}P_v\bigl(v^{-s}-(v+1)^{-s}\bigr).
\]
又
\[
v^{-s}-(v+1)^{-s}
=\int_v^{v+1}su^{-s-1}\,du.
\]
因此尾和满足
\[
\left|\sum_{v=N}^{M}a_vv^{-s}\right|
\le |P_M|M^{-\Re(s)}+|P_{N-1}|N^{-\Re(s)}
 +|s|\sum_{v=N}^{M-1}|P_v|\int_v^{v+1}u^{-\Re(s)-1}\,du .
\]

若 $|P_v|\le Pv^{\sigma_1}$，且 $\Re(s)\ge\sigma_1+\epsilon$，右边由
\[
\sum_{v\ge N} v^{-1-\epsilon}
\]
控制，故在每个闭的右半平面 $\Re(s)\ge\sigma_1+\epsilon$ 上一致收敛。紧一致收敛的解析函数列极限仍解析，于是得到 (3)。

(1) 是同一个公式应用在
$b_n=a_nn^{-s_0}$ 与指数 $s-s_0$ 上。部分和 $\sum_{n\le m}b_n$
有界，而 $\Re(s-s_0)\ge\delta$ 给出同样的尾和控制。由此得到右侧区域的一致收敛和解析性。

(2) 来自收敛区域的右半平面性质：若在 $s_1$ 收敛，则由 (1) 得到在
$\Re(s)>\Re(s_1)$ 收敛。把所有收敛点实部的下确界记为 $\sigma_0$，便得到所需的收敛横坐标。

若 $P_n/n\to g$，则把 $D(s)$ 与 $g\zeta(s)$ 相减。设
$Q_n=P_n-gn$，则 $Q_n/n\to0$。对实数 $\sigma>1$，
\[
D(\sigma)-g\zeta(\sigma)
=\sum_{n\ge1}(Q_n-Q_{n-1})n^{-\sigma}.
\]
给定 $\varepsilon>0$，取 $N$ 使 $|Q_n|\le\varepsilon n$ 对 $n\ge N$ 成立。有限的前段乘以
$\sigma-1$ 后趋于 $0$；后段由 Abel 公式被常数倍
\[
\varepsilon\left(N^{1-\sigma}
 +\sigma\sum_{n\ge N}n^{-\sigma}\right)
\]
控制。乘以 $\sigma-1$ 后上界不超过常数倍 $\varepsilon$。令
$\varepsilon$ 趋于 $0$ 得到
$(\sigma-1)(D(\sigma)-g\zeta(\sigma))\to0$。复变量沿
$\Re(s)>1$ 趋于 $1$ 时可在小角形区域中用同一估计；再用
$(s-1)\zeta(s)\to1$，得到 (4)。
\end{proof}

\begin{proposition}
设 $A$ 为有限交换群，$\chi_1,\chi_2:A\to S^1$ 为酉特征标。则
\[
\sum_{a\in A}\chi_1(a)\chi_2(a)=
\begin{cases}
0,&\chi_1\ne\chi_2^{-1},\\
|A|,&\chi_1=\chi_2^{-1}.
\end{cases}
\]
特别地，若 $\chi\ne1$，则 $\sum_{a\in A}\chi(a)=0$。
\end{proposition}

\begin{proof}
若 $\psi=\chi_1\chi_2$ 非平凡，取 $b\in A$ 使 $\psi(b)\ne1$。令 $S=\sum_a\psi(a)$。乘以 $\psi(b)$ 并换元得
\[
\psi(b)S=\sum_a\psi(ba)=\sum_a\psi(a)=S.
\]
因为 $\psi(b)\ne1$，只能有 $S=0$。若 $\psi=1$，每一项都等于
$1$，总和为 $|A|$。取 $\chi_2=1$，便得到非平凡特征标的求和为零。
\end{proof}

\begin{definition}
$I_F^{\mathfrak m}/S_F^{\mathfrak m}$ 的酉特征标称为模 $\mathfrak m$ 的理想特征标。定义
\[
L(s,\chi)=\sum_{(\mathfrak n,\mathfrak m)=1}
\frac{\chi(\mathfrak n)}{\Nm(\mathfrak n)^s},
\qquad \Re(s)>1,
\]
其中求和遍历与 $\mathfrak m$ 互素的整理想。
\end{definition}

\begin{proposition}
令 $\chi$ 为模 $\mathfrak m$ 的理想特征标。
\begin{enumerate}[label=(\arabic*)]
\item $L(s,\chi)$ 在 $\Re(s)>1$ 内紧一致绝对收敛，故解析。
\item 在 $\Re(s)>1$ 有 Euler 乘积
\[
L(s,\chi)=
\prod_{\mathfrak p\nmid\mathfrak m}
\left(1-\frac{\chi(\mathfrak p)}{\Nm(\mathfrak p)^s}\right)^{-1}.
\]
\item 设 $h_{\mathfrak m}=|I_F^{\mathfrak m}/S_F^{\mathfrak m}|$。则
\[
\lim_{s\to1}(s-1)L(s,\chi)
=
\begin{cases}
0,&\chi\ne1,\\
g_{\mathfrak m}h_{\mathfrak m},&\chi=1.
\end{cases}
\]
\end{enumerate}
\end{proposition}

\begin{proof}
(1) 因为 $|\chi(\mathfrak n)|=1$，只需控制
\[
\sum_{(\mathfrak n,\mathfrak m)=1}\Nm(\mathfrak n)^{-\sigma}
\]
在 $\sigma>1$ 时的收敛性。令
\[
A(x)=|\{\mathfrak n:(\mathfrak n,\mathfrak m)=1,\ \Nm(\mathfrak n)\le x\}|.
\]
上一节已经给出每个射线类中的理想数为 $O(x)$，射线类只有有限多个，所以
$A(x)=O(x)$。Abel 求和给出
\[
\sum_{\Nm(\mathfrak n)\le X}\Nm(\mathfrak n)^{-\sigma}
=A(X)X^{-\sigma}
 \sigma\int_1^X A(t)t^{-\sigma-1}\,dt .
\]
当 $\sigma>1$ 时，右边由常数倍
\[
X^{1-\sigma}+\int_1^\infty t^{-\sigma}\,dt
\]
控制。若 $\sigma\ge1+\epsilon$，这个控制在紧集上一致成立，故 $L(s,\chi)$ 在
$\Re(s)>1$ 内紧一致绝对收敛并解析。

(2) 对有限个不除 $\mathfrak m$ 的素理想集合 $T$，有有限乘积展开
\[
\prod_{\mathfrak p\in T}
\left(1-\frac{\chi(\mathfrak p)}{\Nm(\mathfrak p)^s}\right)^{-1}
=
\sum_{\operatorname{supp}(\mathfrak n)\subset T}
\frac{\chi(\mathfrak n)}{\Nm(\mathfrak n)^s}.
\]
这里的等式来自两个事实。第一，每个局部因子都是几何级数
\[
1+\frac{\chi(\mathfrak p)}{\Nm(\mathfrak p)^s}
+\frac{\chi(\mathfrak p)^2}{\Nm(\mathfrak p)^{2s}}+\cdots .
\]
第二，整理想有唯一素理想分解
\(\mathfrak n=\prod \mathfrak p^{e_{\mathfrak p}}\)，并且 $\chi$ 是群同态，所以
\(\chi(\mathfrak n)=\prod\chi(\mathfrak p)^{e_{\mathfrak p}}\)。由 (1) 的绝对收敛，可令
$T$ 递增遍历所有 $\mathfrak p\nmid\mathfrak m$，得到无穷 Euler 乘积。

(3) 对每个射线类 $K$ 定义
\[
\zeta(s,K)=\sum_{\mathfrak n\in K}\frac1{\Nm(\mathfrak n)^s}.
\]
上一节给出该类中理想计数 $P_x\sim g_{\mathfrak m}x$。把
$\zeta(s,K)$ 按整数范数重写为
\[
\zeta(s,K)=\sum_{n\ge1}\frac{a_n(K)}{n^s},\qquad
a_n(K)=|\{\mathfrak n\in K:\Nm(\mathfrak n)=n\}|.
\]
此时部分和 $\sum_{k\le n}a_k(K)$ 正是 $P_n$，由 Dirichlet 级数定理，
\[
\lim_{s\to1}(s-1)\zeta(s,K)=g_{\mathfrak m}.
\]
又
\[
L(s,\chi)=\sum_K\chi(K)\zeta(s,K).
\]
这是因为 $\chi$ 在同一射线类上取同一个值，而全部与 $\mathfrak m$ 互素的整理想按射线类分成有限个互不相交的集合。
于是
\[
\lim_{s\to1}(s-1)L(s,\chi)
=g_{\mathfrak m}\sum_K\chi(K).
\]
由特征标正交性，非平凡特征标时和为 $0$，平凡特征标时和为 $h_{\mathfrak m}$。
\end{proof}

\begin{definition}
数域 $F$ 的 Dedekind $\zeta$ 函数定义为
\[
\zeta_F(s)=\sum_{\mathfrak a\ne0}\Nm(\mathfrak a)^{-s},
\qquad \Re(s)>1,
\]
其中求和遍历 $F$ 的非零整理想。
\end{definition}

由整理想唯一分解，同上一命题的 Euler 乘积证明，在 $\Re(s)>1$ 有
\[
\zeta_F(s)=\prod_{\mathfrak p}(1-\Nm(\mathfrak p)^{-s})^{-1}.
\]
取 $\mathfrak m=1$ 时，射线类群就是通常的理想类群
$\operatorname{Cl}_F=I_F/P_F$。每个非零整理想属于唯一的理想类，因此
\[
\zeta_F(s)=\sum_{K\in\operatorname{Cl}_F}\zeta(s,K).
\]
把前面射线类残数公式在 $\mathfrak m=1$ 的情形求和，得到 Dedekind
$\zeta$ 函数在 $s=1$ 的残数：
\[
\lim_{s\to1}(s-1)\zeta_F(s)
=
\frac{2^{r_1}(2\pi)^{r_2}R_Fh_F}
{w_F\sqrt{|d(\OF)|}}.
\]
这就是类数公式的基本形态。它的含义是：类数 $h_F$、调控子 $R_F$、单位根数
$w_F$ 和判别式 $d(\OF)$ 不是分散的代数不变量，而是共同控制理想计数函数的主项。

\section{密度}

设集合 $S$ 的元素为数域 $F$ 的非零素理想。为了估计这种集合的大小，不能直接数
$|\{\mathfrak p\in S:\Nm(\mathfrak p)\le x\}|$，本节改用解析权重
\(\Nm(\mathfrak p)^{-s}\)。当 $s\to1+$ 时，所有素理想的总和像
\(\log 1/(s-1)\) 一样发散；若 $S$ 占其中固定比例，这个比例就是 Dirichlet 密度。

\begin{definition}
设 $S$ 为 $F$ 的非零素理想集合。若极限
\[
d(S)=
\lim_{s\to1+}
\frac{\sum_{\mathfrak p\in S}\Nm(\mathfrak p)^{-s}}
{\sum_{\mathfrak p}\Nm(\mathfrak p)^{-s}}
\]
存在，则称它为 $S$ 的 Dirichlet 密度。
\end{definition}

若两个函数 $f(s),g(s)$ 的差在 $s=1$ 附近有界且不影响密度计算，本节记作
$f(s)\sim g(s)$。由 Dedekind $\zeta$ 函数的 Euler 乘积，
\[
\log\zeta_F(s)
=\sum_{\mathfrak p}\sum_{m\ge1}\frac1{m\Nm(\mathfrak p)^{ms}}.
\]
其中 $m\ge2$ 的部分在 $s=1$ 附近收敛，因为
\[
\sum_{\mathfrak p}\sum_{m\ge2}\frac1{m\Nm(\mathfrak p)^{m\sigma}}
\le
\sum_{\mathfrak a\ne0}\Nm(\mathfrak a)^{-2\sigma}
\]
在 $\sigma>1/2$ 时有界。另一方面，上一节的类数公式给出
\(\zeta_F(s)\) 在 $s=1$ 有单极点，所以
\[
\sum_{\mathfrak p}\Nm(\mathfrak p)^{-s}\sim \log\frac1{s-1},
\]
因此 Dirichlet 密度也可写成
\[
d(S)=
\lim_{s\to1+}
\frac{\sum_{\mathfrak p\in S}\Nm(\mathfrak p)^{-s}}
{\log\frac1{s-1}}.
\]

\begin{proposition}
数域 $F$ 中剩余次数 $f_{\mathfrak p|p}=1$ 的素理想组成的集合密度为 $1$。
\end{proposition}

\begin{proof}
设 $S$ 为剩余次数为 $1$ 的素理想集合。删去有限多个分歧素理想不会改变密度。若
$\mathfrak p$ 位于普通素数 $p$ 上方且剩余次数大于 $1$，则
\[
\Nm(\mathfrak p)\ge p^2.
\]
对每个 $p$，上方素理想数不超过 $[F:\Q]$，所以这些例外素理想的级数由
\[
[F:\Q]\sum_p p^{-2\sigma},\qquad \sigma=\Re(s),
\]
控制，在 $s=1$ 附近有界。总素理想和发散如 $\log(1/(s-1))$，故例外集合密度为 $0$，所求集合密度为 $1$。
\end{proof}

\begin{proposition}
设 $S_F^{\mathfrak m}\subset H\subset I_F^{\mathfrak m}$，令 $S$ 为 $H$ 中素理想组成的集合。则
\[
d(S)\le \frac1{[I_F^{\mathfrak m}:H]}.
\]
\end{proposition}

\begin{proof}
令 $A=I_F^{\mathfrak m}/H$，$h=|A|$。对 $A$ 的所有特征标用正交关系：
\[
\sum_{\chi\in\widehat A}\chi(\mathfrak p)=
\begin{cases}
h,&\mathfrak p\in H,\\
0,&\mathfrak p\notin H.
\end{cases}
\]
对 Euler 乘积取对数，
\[
\log L(s,\chi)
=\sum_{\mathfrak p\nmid\mathfrak m}\chi(\mathfrak p)\Nm(\mathfrak p)^{-s}+g_\chi(s),
\]
其中
\[
g_\chi(s)=\sum_{\mathfrak p\nmid\mathfrak m}\sum_{r\ge2}
\frac{\chi(\mathfrak p)^r}{r\Nm(\mathfrak p)^{rs}}
\]
在 $s=1$ 附近有界。把所有 $\chi\in\widehat A$ 相加，正交关系把属于
$H$ 的素理想挑出来：
\[
h\sum_{\mathfrak p\in H}\Nm(\mathfrak p)^{-s}
=
\sum_{\chi\in\widehat A}\bigl(\log L(s,\chi)-g_\chi(s)\bigr).
\]
平凡特征标项满足
\[
\log L(s,1)= -\log(s-1)+O(1),
\]
因为上一节已经给出 $(s-1)L(s,1)$ 在 $s=1$ 有非零有限极限。非平凡特征标项
\(\log L(s,\chi)\) 要么在 $s=1$ 附近有界，要么由于 $L(s,\chi)$ 趋于 $0$ 而趋向
$-\infty$；这两种情形都不会提供正的主发散项。

若 $S$ 的密度存在，则
\[
\sum_{\mathfrak p\in S}\Nm(\mathfrak p)^{-s}
=d(S)\log\frac1{s-1}+O(1)
\]
按密度定义理解为主项关系。由于 $S\subset H$，
\[
\sum_{\mathfrak p\in H}\Nm(\mathfrak p)^{-s}
-\sum_{\mathfrak p\in S}\Nm(\mathfrak p)^{-s}\ge0 .
\]
左边的主项系数为 $1/h-d(S)$。若 $d(S)>1/h$，该差在 $s\to1+$ 时会趋向负无穷，与非负性矛盾。因此 $d(S)\le1/h$。
\end{proof}

\begin{lemma}[陪集链与剩余次数]
设 $E/F$ 为有限 Galois 扩张，$G=\Gal(E/F)$，$H$ 为 $G$ 的子群，$K=E^H$。取
$E$ 的一个未分歧素理想 $\mathfrak P$，设 $\mathfrak p=\mathfrak P\cap F$，并令
$\sigma=\operatorname{Frob}_{\mathfrak P}(E/F)$。让 $\sigma$ 在右陪集集合
$H\backslash G$ 上右乘作用。若一个轨道为
\[
H\tau,\ H\tau\sigma,\ \dots,\ H\tau\sigma^{\ell-1},
\qquad H\tau\sigma^\ell=H\tau,
\]
则它对应 $K$ 中一个位于 $\mathfrak p$ 上方的素理想
\(\mathfrak q=\tau(\mathfrak P)\cap K\)，并且
\[
f_{\mathfrak q|\mathfrak p}=\ell.
\]
所有轨道给出的 $\mathfrak q$ 正好是 $K$ 中位于 $\mathfrak p$ 上方的全部素理想。
\end{lemma}

\begin{proof}
先计算一个轨道给出的剩余次数。素理想 $\tau(\mathfrak P)$ 在 $E/K$ 中的分解群为
\[
H_{\tau\mathfrak P}
=\{\eta\in H:\eta\tau(\mathfrak P)=\tau(\mathfrak P)\}
=H\cap \tau G_{\mathfrak P}\tau^{-1}.
\]
因为 $\mathfrak P$ 未分歧，$G_{\mathfrak P}$ 由 Frobenius $\sigma$ 生成。轨道长度
$\ell$ 的条件
\[
H\tau\sigma^\ell=H\tau
\]
等价于 $\tau\sigma^\ell\tau^{-1}\in H$，并且 $\ell$ 是满足这个条件的最小正整数。于是
\[
H_{\tau\mathfrak P}=\langle\tau\sigma^\ell\tau^{-1}\rangle.
\]
由分解群阶数等于剩余次数，
\[
f_{\tau\mathfrak P|\mathfrak q}=|H_{\tau\mathfrak P}|,
\qquad
f_{\tau\mathfrak P|\mathfrak p}=|\langle\sigma\rangle|.
\]
塔公式给出
\[
f_{\mathfrak q|\mathfrak p}
=\frac{f_{\tau\mathfrak P|\mathfrak p}}
{f_{\tau\mathfrak P|\mathfrak q}}
=\frac{|\langle\sigma\rangle|}
{|\langle\tau\sigma^\ell\tau^{-1}\rangle|}
=\ell.
\]

若两个陪集轨道给出同一个 $\mathfrak q$，则两个 $E$ 中素理想都位于 $\mathfrak q$ 上方。由于 $E/K$ 为 Galois 扩张，存在 $\gamma\in H$ 把其中一个送到另一个；这说明两个陪集相差一个 $\sigma$ 的幂，故在同一条轨道中。最后，各轨道长度之和为
 $|H\backslash G|=[K:F]$，而 $K/F$ 中位于 $\mathfrak p$ 上方的素理想满足
\(\sum e_if_i=[K:F]\)。未分歧时 $e_i=1$，前面已经构造出互异素理想且剩余次数之和达到
$[K:F]$，所以没有遗漏。
\end{proof}

\begin{lemma}[循环生成元的共轭计数]
设 $G$ 为有限群，$\sigma\in G$ 的阶为 $n$。定义
\[
\lceil\sigma\rceil
=\{g\sigma^mg^{-1}:g\in G,\ (m,n)=1\}.
\]
则
\[
|\lceil\sigma\rceil|
=\varphi(n)[G:N_G(\langle\sigma\rangle)].
\]
\end{lemma}

\begin{proof}
当 $(m,n)=1$ 时，$\sigma$ 与 $\sigma^m$ 生成同一个循环子群，因此
\[
C_G(\sigma)=C_G(\sigma^m).
\]
固定这样的 $m$，元素 $\sigma^m$ 的共轭类大小为
$[G:C_G(\sigma)]$。当 $m$ 遍历模 $n$ 的可逆剩余类时，一共先得到
\(\varphi(n)[G:C_G(\sigma)]\) 个“带重数”的元素。

现在计算每个元素被重复计算的次数。若 $\sigma^m$ 与 $\sigma^{m'}$ 共轭，则有
$g\langle\sigma\rangle g^{-1}=\langle\sigma\rangle$，所以 $g\in N_G(\langle\sigma\rangle)$。在正规化子中，能把 $\sigma$ 送到某个生成元的不同共轭由
$N_G(\langle\sigma\rangle)/C_G(\sigma)$ 计数。因此每个元素的重数为
\([N_G(\langle\sigma\rangle):C_G(\sigma)]\)。相除得到
\[
\frac{\varphi(n)[G:C_G(\sigma)]}
{[N_G(\langle\sigma\rangle):C_G(\sigma)]}
=\varphi(n)[G:N_G(\langle\sigma\rangle)].
\]
\end{proof}

\begin{theorem}
设 $E/F$ 为有限 Galois 扩张，$\sigma\in G=\Gal(E/F)$。令
\[
\langle\sigma\rangle^\ast
=\{g\sigma^mg^{-1}:g\in G,\ (m,\operatorname{ord}\sigma)=1\}.
\]
若 $S$ 为那些存在 $\mathfrak P|\mathfrak p$ 且 Frobenius 落在 $\langle\sigma\rangle^\ast$ 中的 $F$ 素理想集合，则
\[
d(S)=\frac{|\langle\sigma\rangle^\ast|}{[E:F]}.
\]
\end{theorem}

\begin{proof}
记 $n=\operatorname{ord}(\sigma)$。对 $n$ 作归纳。

当 $n=1$ 时，集合 $S$ 是在 $E$ 中完全分裂的 $F$ 素理想集合。设
$S_E$ 为 $E$ 中位于这些 $\mathfrak p$ 上方的素理想集合。每个完全分裂的
$\mathfrak p$ 在 $E$ 中给出 $[E:F]$ 个素理想，且这些素理想范数都等于
$\Nm_F(\mathfrak p)$。由“剩余次数为 $1$ 的素理想密度为 $1$”应用于数域 $E$，可得
\[
\sum_{\mathfrak P\in S_E}\Nm_E(\mathfrak P)^{-s}
\sim \log\frac1{s-1}.
\]
另一方面左边等于
\[
[E:F]\sum_{\mathfrak p\in S}\Nm_F(\mathfrak p)^{-s}
\]
加上有限分歧素理想造成的有界差。因此
\[
d(S)=\frac1{[E:F]}.
\]

设 $n>1$。对每个 $d|n$，令 $S_d$ 为 Frobenius 落在
\(\lceil\sigma^d\rceil\) 中的素理想集合。当 $d>1$ 时，$\sigma^d$ 的阶小于
$n$，归纳假设给出
\[
d(S_d)=\frac{|\lceil\sigma^d\rceil|}{|G|}.
\]
令 $K=E^{\langle\sigma\rangle}$。由陪集链引理，$F$ 的素理想 $\mathfrak p$ 在
$K$ 中有剩余次数 $1$ 的素理想，当且仅当某个 Frobenius 的共轭落入
\(\langle\sigma\rangle\)，也就是 $\mathfrak p$ 属于某个 $S_d$。若
$\mathfrak p\in S_d$，则剩余次数 $1$ 的上方素理想个数等于
\[
[N_G(\langle\sigma^d\rangle):\langle\sigma\rangle].
\]
把 $K$ 中剩余次数为 $1$ 的素理想按它们在 $F$ 中的收缩分组，并使用前面的密度为
$1$ 的命题，得到主项恒等式
\[
\log\frac1{s-1}
\sim
\sum_{d|n}
[N_G(\langle\sigma^d\rangle):\langle\sigma\rangle]
\sum_{\mathfrak p\in S_d}\Nm(\mathfrak p)^{-s}.
\]
其中 $d>1$ 的项由归纳假设已知，未知项只剩 $S=S_1$。代入
\[
|\langle\sigma\rangle^\ast|
=\varphi(n)[G:N_G(\langle\sigma\rangle)]
\]
以及同一公式对 $\sigma^d$ 的版本，系数化简为
\[
\sum_{\mathfrak p\in S}\Nm(\mathfrak p)^{-s}
\sim
\frac{|\langle\sigma\rangle^\ast|}{|G|}
\log\frac1{s-1}.
\]
这正是所需密度公式。这个证明的核心是：Frobenius 在陪集上的轨道长度就是中间域中的剩余次数。
\end{proof}

\begin{corollary}
有限 Galois 扩张 $E/F$ 中完全分裂的 $F$ 素理想集合的密度为
\[
\frac1{[E:F]}.
\]
\end{corollary}

\begin{proof}
完全分裂等价于 Frobenius 为恒等元。把上一条定理用于 $\sigma=1$，此时
\(|\langle1\rangle^\ast|=1\)，得到密度 $1/[E:F]$。
\end{proof}

\begin{corollary}
设 $E/F$ 为有限 Galois 扩张，$\sigma\in G$。若 $C(\sigma)$ 为 $\sigma$ 的共轭类，则 Frobenius 共轭类等于 $C(\sigma)$ 的未分歧素理想集合密度为
\[
\frac{|C(\sigma)|}{|G|}.
\]
\end{corollary}

\begin{proof}
固定一个未分歧素理想 $\mathfrak p$ 后，不同的 $\mathfrak P|\mathfrak p$ 给出的 Frobenius 元互为共轭。因此“存在 $\mathfrak P|\mathfrak p$ 使 Frobenius 等于
$\sigma$”等价于 Frobenius 共轭类为 $C(\sigma)$。共轭类 $C(\sigma)$ 是若干集合
\(\langle\tau\rangle^\ast\) 的不交并，其中 $\tau$ 可取为该共轭类中不同循环子群的生成元代表。对这些代表应用上一条定理并相加，得到
\[
d=\frac{|C(\sigma)|}{|G|}.
\]
\end{proof}

\begin{theorem}
设 $E/F$ 为有限 Galois 数域扩张，$\mathfrak m$ 为 $F$ 的模。则
\[
[I_F^{\mathfrak m}:\Nm_{E/F}(I_E^{\mathfrak m})S_F^{\mathfrak m}]
\le [E:F].
\]
\end{theorem}

\begin{proof}
令
\[
H=\Nm_{E/F}(I_E^{\mathfrak m})S_F^{\mathfrak m}.
\]
除有限多个分歧素理想和整除 $\mathfrak m$ 的素理想外，若 $\mathfrak p$ 在 $E$ 中完全分裂，则每个上方素理想 $\mathfrak P$ 的范数为 $\mathfrak p$，因此
\[
\mathfrak p=\Nm_{E/F}(\mathfrak P)\in H.
\]
所以完全分裂素理想集合包含于 $H$ 中素理想集合，有限例外不影响密度。由完全分裂密度，
\[
\frac1{[E:F]}\le d(\{\mathfrak p:\mathfrak p\in H\}).
\]
由上一命题，
\[
d(\{\mathfrak p:\mathfrak p\in H\})
\le \frac1{[I_F^{\mathfrak m}:H]}.
\]
合并即得
\[
[I_F^{\mathfrak m}:H]\le [E:F].
\]
\end{proof}

\section*{习题详解}
\addcontentsline{toc}{section}{习题详解}

\begin{exercise}
证明 $\sum_{n\ge1}n^{-s}$ 与 $\prod_p(1-p^{-s})^{-1}$ 在 $\Re(s)>1$ 上解析且相等。
\end{exercise}

\begin{proof}[解答]
\textbf{解答状态：solvable.}
若 $\sigma=\Re(s)>1$，则
\[
\sum_{n=1}^\infty |n^{-s}|=\sum n^{-\sigma}
\]
收敛，并在紧子集上一致收敛，所以定义解析函数。Euler 乘积中
\[
\sum_p |p^{-s}|\le \sum_{n\ge2}n^{-\sigma}<\infty,
\]
故无穷乘积在紧子集上一致收敛。展开每个几何级数：
\[
\prod_p(1-p^{-s})^{-1}
=\prod_p(1+p^{-s}+p^{-2s}+\cdots).
\]
由整数唯一分解，每个正整数 $n=\prod p^{v_p(n)}$ 出现且只出现一次，所得项为 $n^{-s}$。绝对收敛保证可交换求和与乘积，因此相等。
\end{proof}

\begin{exercise}
取模 $m$ 的 Dirichlet 特征标 $\chi$。证明 $L(s,\chi)=\sum_{n\ge1}\chi(n)n^{-s}$ 在 $\Re(s)>1$ 解析，并有 Euler 积。
\end{exercise}

\begin{proof}[解答]
\textbf{解答状态：solvable.}
因为 $|\chi(n)|\le1$，
\[
\sum|\chi(n)n^{-s}|\le\sum n^{-\sigma}
\]
在 $\sigma>1$ 收敛，并且在每个闭半平面 $\sigma\ge1+\epsilon$ 上一致收敛，所以
$L(s,\chi)$ 在 $\Re(s)>1$ 上解析。Dirichlet 特征标完全乘法，即
$\chi(ab)=\chi(a)\chi(b)$；若采用通常的延拓约定，当 $(n,m)>1$ 时
$\chi(n)=0$。于是对每个素数 $p$，
\[
\prod_p\left(1+\frac{\chi(p)}{p^s}
\frac{\chi(p^2)}{p^{2s}}+\cdots\right)
=
\prod_p\left(1+\frac{\chi(p)}{p^s}
\frac{\chi(p)^2}{p^{2s}}+\cdots\right)
=\prod_p\left(1-\frac{\chi(p)}{p^s}\right)^{-1}.
\]
若 $p|m$，则 $\chi(p)=0$，对应的局部因子就是 $1$；若 $p\nmid m$，它是普通几何级数。展开整个乘积时，由唯一分解
\[
n=\prod_p p^{v_p(n)}
\]
得到的项为
\[
\prod_p \frac{\chi(p)^{v_p(n)}}{p^{v_p(n)s}}
=\frac{\chi(n)}{n^s}.
\]
绝对收敛保证无穷乘积与求和交换合法，因此 Euler 积成立。
\end{proof}

\begin{exercise}
证明 Dirichlet 级数的 Abel 分部求和公式，并推出若部分和 $P_n=O(n^{\sigma_1})$，则收敛横坐标 $\sigma_0\le\sigma_1$。
\end{exercise}

\begin{proof}[解答]
\textbf{解答状态：solvable.}
设 $P_n=\sum_{k\le n}a_k$，并约定 $P_0=0$。由于 $a_v=P_v-P_{v-1}$，
\[
\sum_{v=n+1}^{n+m}\frac{a_v}{v^s}
=\sum_{v=n+1}^{n+m}\frac{P_v-P_{v-1}}{v^s}.
\]
把两部分错位：
\[
\sum_{v=n+1}^{n+m}\frac{P_v}{v^s}
-\sum_{v=n+1}^{n+m}\frac{P_{v-1}}{v^s}
=
\sum_{v=n+1}^{n+m}\frac{P_v}{v^s}
-\sum_{u=n}^{n+m-1}\frac{P_u}{(u+1)^s}.
\]
端点项分别为 $P_{n+m}(n+m)^{-s}$ 与 $-P_n(n+1)^{-s}$，中间项为
\[
\sum_{v=n+1}^{n+m-1}P_v
\left(\frac1{v^s}-\frac1{(v+1)^s}\right).
\]
又
\[
\frac1{v^s}-\frac1{(v+1)^s}
=\int_v^{v+1}\frac{s\,du}{u^{s+1}}.
\]
所以得到源公式
\[
\sum_{v=n+1}^{n+m}\frac{a_v}{v^s}
=\frac{P_{n+m}}{(n+m)^s}-\frac{P_n}{(n+1)^s}
+\sum_{v=n+1}^{n+m-1}P_v\int_v^{v+1}\frac{s\,du}{u^{s+1}}.
\]

若 $|P_v|\le Pv^{\sigma_1}$，$\sigma=\Re(s)>\sigma_1$，则
\[
\left|\frac{P_{n+m}}{(n+m)^s}\right|\le \frac{P}{(n+m)^{\sigma-\sigma_1}}
\le \frac{P}{n^{\sigma-\sigma_1}},
\]
并且
\[
\left|\frac{P_n}{(n+1)^s}\right|\le \frac{P}{n^{\sigma-\sigma_1}}.
\]
积分项满足
\[
\left|
\sum_{v=n+1}^{n+m-1}P_v\int_v^{v+1}\frac{s\,du}{u^{s+1}}
\right|
\le
P|s|\sum_{v=n+1}^{n+m-1}\int_v^{v+1}u^{\sigma_1-\sigma-1}\,du.
\]
因为 $\sigma-\sigma_1>0$，
\[
\sum_{v=n+1}^{n+m-1}\int_v^{v+1}u^{\sigma_1-\sigma-1}\,du
\le
\int_n^\infty u^{\sigma_1-\sigma-1}\,du
=\frac1{(\sigma-\sigma_1)n^{\sigma-\sigma_1}}.
\]
因此
\[
\left|\sum_{v=n+1}^{n+m}\frac{a_v}{v^s}\right|
\le
\left(2+\frac{|s|}{\sigma-\sigma_1}\right)
\frac{P}{n^{\sigma-\sigma_1}}.
\]
右端随 $n\to\infty$ 趋于 $0$，且与 $m$ 无关，所以级数在
$\Re(s)>\sigma_1$ 收敛，从而收敛横坐标 $\sigma_0\le\sigma_1$。
\end{proof}

\begin{exercise}
若 $P_n/n\to g$，证明 $\lim_{s\to1,\Re(s)>1}(s-1)D(s)=g$。
\end{exercise}

\begin{proof}[解答]
\textbf{解答状态：solvable.}
令
\[
\varphi(s)=D(s)-g\zeta(s)
=\sum_{n=1}^\infty \frac{a_n-g}{n^s}.
\]
其部分和为 $Q_n=P_n-ng=o(n)$。给定 $\varepsilon>0$，取 $n_1$ 使 $n>n_1$ 时 $|Q_n|\le\varepsilon n$，并令
\[
Q=\max_{1\le n\le n_1}|Q_n|.
\]
对 $\varphi$ 的尾和使用 Abel 公式。前 $n_1$ 项的贡献被
\[
Q|s|\sum_{n=1}^{n_1}\int_n^{n+1}u^{-\sigma-1}\,du
\]
控制；后面的贡献被
\[
\varepsilon |s|\sum_{n>n_1}n\int_n^{n+1}u^{-\sigma-1}\,du
\le
\varepsilon |s|\int_1^\infty u^{-\sigma}\,du
=\frac{\varepsilon |s|}{\sigma-1}
\]
控制。因此
\[
|\varphi(s)|\le |s|\left(Q+\frac{\varepsilon}{\sigma-1}\right)
\]
按源题所要求的形式成立。乘以 $s-1$ 后令 $s\to1$，有限前段消失；尾段的上界由常数倍
$\varepsilon$ 控制。因为 $\varepsilon$ 任意，
\[
\lim_{s\to1}(s-1)\varphi(s)=0.
\]
而 $(s-1)\zeta(s)\to1$，所以 $(s-1)D(s)\to g$。
\end{proof}

\begin{exercise}
证明带单位元有限维代数 $A$ 只有有限个极大理想。设 $\mathfrak r=\cap_i\mathfrak m_i$，证明
\[
A/\mathfrak r\simeq A/\mathfrak m_1\oplus\cdots\oplus A/\mathfrak m_h.
\]
\end{exercise}

\begin{proof}[解答]
\textbf{解答状态：solvable.}
有限维代数是 Artin 环，因此任意降链理想稳定。先证明极大理想有限。若有无限多个不同极大理想
\(\mathfrak m_1,\mathfrak m_2,\dots\)，令
\[
I_r=\mathfrak m_1\cap\cdots\cap\mathfrak m_r.
\]
若 $I_r\subset \mathfrak m_{r+1}$，则由极大理想两两互素可得
\[
A=\mathfrak m_i+\mathfrak m_{r+1}\quad(1\le i\le r),
\]
从而中国剩余思想给出
\[
I_r+\mathfrak m_{r+1}=A,
\]
矛盾。所以 $I_{r+1}=I_r\cap\mathfrak m_{r+1}$ 严格小于 $I_r$。这会构造无限严格下降链，与 Artin 条件矛盾。

不同极大理想两两互素，因为 $\mathfrak m_i+\mathfrak m_j$ 严格包含
\(\mathfrak m_i\)，只能为 $A$。中国剩余定理给出
\[
A/(\mathfrak m_1\cap\cdots\cap\mathfrak m_h)
\simeq
\bigoplus_i A/\mathfrak m_i.
\]
这里交就是 $\mathfrak r$，于是得到所需同构。
\end{proof}

\begin{exercise}
设 $L=E\otimes_FK/\mathfrak m$ 为 $E\otimes_FK$ 的极大理想商。证明 $L$ 是 $E$ 与 $K$ 的域合成。
\end{exercise}

\begin{proof}[解答]
\textbf{解答状态：solvable.}
映射
\[
s:E\to L,\quad x\mapsto x\otimes1+\mathfrak m,\qquad
t:K\to L,\quad y\mapsto1\otimes y+\mathfrak m
\]
是 $F$-代数同态。由于 $E,K$ 是域，非零同态的核只能为 $0$，所以 $s,t$ 都是单射。任意简单张量在商中满足
\[
(x\otimes y)+\mathfrak m=(x\otimes1)(1\otimes y)+\mathfrak m=s(x)t(y).
\]
而 $E\otimes_FK$ 由简单张量线性生成，所以 $L$ 由 $s(E)$ 与 $t(K)$ 作为域生成。于是 $L$ 正是两个嵌入域 $s(E)$、$t(K)$ 在同一上域中的合成域。
\end{proof}

\begin{exercise}
设 $K=F(\theta)$ 为有限可分扩张，$f$ 为 $\theta$ 的最小多项式，在 $E[X]$ 中分解为互异不可约因子 $p_1\cdots p_h$。证明张量积分解和合成域描述。
\end{exercise}

\begin{proof}[解答]
\textbf{解答状态：solvable.}
因为 $K=F(\theta)$ 且 $f$ 是 $\theta$ 的最小多项式，
\[
K\simeq F[X]/(f).
\]
张量到 $E$ 后
\[
E\otimes_FK\simeq E\otimes_FF[X]/(f)\simeq E[X]/(f).
\]
可分性保证 $f$ 在任意扩域上没有重因子，所以不可约因子
$p_i$ 两两互素。中国剩余定理给出
\[
E[X]/(f)\simeq\bigoplus_{i=1}^hE[X]/(p_i).
\]
令 $L_i=E[X]/(p_i)$。因为 $p_i$ 不可约，$L_i$ 是域，并且
\[
[L_i:E]=\deg p_i,\qquad
\sum_i[L_i:E]=\sum_i\deg p_i=\deg f=[K:F].
\]
在 $L_i$ 中，$E$ 由常数嵌入，$K$ 由 $\theta\mapsto X\bmod p_i$ 嵌入；$L_i$ 由
$E$ 和这个 $\theta$ 的像生成，所以它是 $E$ 与 $K$ 的一个合成域。反过来，任一合成域给出
$E\otimes_FK$ 到该域的满同态，其核为极大理想；极大理想在上述直和分解中对应某个
$p_i$，因此合成域同构于某个 $L_i$。
\end{proof}

\begin{exercise}
设 $K\supset E\supset F$，$K/E$ 可分，$E/F$ 不可分，$S/F$ 为 $K/F$ 的极大可分子域。证明自然映射 $E\otimes_FS\to K$ 为同构。
\end{exercise}

\begin{proof}[解答]
\textbf{解答状态：partially solvable.}
这个结论需要把“不可分扩张”理解为“纯不可分扩张”，或者至少需要保证
$E$ 与 $S$ 在 $F$ 上线性无交。若不加这一层条件，$E\cap S$ 可能含有非平凡的
$F$-可分元素，张量积就不一定是域，结论会失去依据。

在纯不可分解释下，证明如下。乘法给出自然 $F$-双线性映射
\[
E\times S\to K,\qquad (e,s)\mapsto es,
\]
从而诱导
\[
E\otimes_FS\to ES\subset K.
\]
纯不可分扩张与可分扩张线性无交：若有限个 $S$ 中元素在 $E$ 上线性相关，取一个最小关系，其系数产生的代数依赖会给出既可分又纯不可分的非平凡元素；这样的元素只能属于
$F$。所以 $E\otimes_FS$ 是域，并嵌入 $K$。

还需证明 $K=ES$。任取 $\alpha\in K$。因为 $K/E$ 可分，$\alpha$ 在
$E$ 上可分；而 $E/F$ 纯不可分意味着存在 $p^r$ 次幂把 $E$ 的有限生成部分送入
$F$。把 $\alpha$ 的最小多项式系数作足够高的 $p$ 次幂，可看出某个由
$\alpha$ 生成的可分部分已经落在 $S$ 中，剩余的纯不可分部分由 $E$ 提供。于是
\(\alpha\in ES\)，故 $K=ES$。因此自然映射是同构。
\end{proof}

\begin{exercise}
证明不可分扩张与可分扩张的张量积为域的两个断言。
\end{exercise}

\begin{proof}[解答]
\textbf{解答状态：partially solvable.}
这里同样需要把第一问中的“不可分扩张”理解为“纯不可分扩张”。若只是非可分扩张，它可能仍含有非平凡可分子扩张，从而未必与给定可分扩张线性无交。

(1) 设 $E/F$ 纯不可分，$K/F$ 可分。先对有限子扩张证明线性无交。若
$K_0/F$ 是有限可分扩张，则 $K_0=F(\alpha)$ 可由一个可分元生成。若
$E\otimes_FK_0$ 不是域，则 $\alpha$ 的最小多项式在 $E$ 上可约，这会产生
$E$ 中某个非平凡的 $F$-可分代数元素；但纯不可分扩张中可分元素只能在
$F$ 里，矛盾。因此 $E\otimes_FK_0$ 是域。对所有有限 $K_0\subset K$ 取并，得到
$E\otimes_FK$ 是域。

(2) 若 $E$ 中所有在 $F$ 上可分的代数元都已在 $F$ 中，则同一证明适用。对有限可分子扩张
$K_0/F$，若 $E\otimes_FK_0$ 不是域，就会产生 $E$ 与 $K_0$ 的非平凡交叠；这个交叠由
$F$ 上可分元素生成，按假设只能是 $F$。所以 $E$ 与 $K_0$ 线性无交，$E\otimes_FK_0$ 是域。再对
$K$ 的有限可分子扩张取并，得到 $E\otimes_FK$ 是域。
\end{proof}

\begin{exercise}
取数域 $F$。证明存在常数 $C$，使任意 $\alpha\in F$ 都有 $\beta\in\OF$ 和非零整数 $t$，满足 $|t|\le C$ 且
\[
|\Nm(t\alpha-\beta)|<1.
\]
\end{exercise}

\begin{proof}[解答]
\textbf{解答状态：solvable.}
把 $\OF$ 通过 Minkowski 嵌入看成 $F_\R$ 中的全格。取一个有界基本区
$D$，并取一个以原点为中心的足够小的开邻域 $U\subset F_\R$，使
\[
u\in U\quad\Longrightarrow\quad |N_{F/\Q}(u)|<1.
\]
这是可能的，因为范数函数在 $F_\R$ 上连续且 $N(0)=0$。

紧群 $F_\R/\iota(\OF)$ 可由有限多个 $U$ 的平移覆盖；设需要 $C$ 个平移。对任意
\(\alpha\in F\)，考察
\[
0,\ \iota(\alpha),\ 2\iota(\alpha),\dots,\ C\iota(\alpha)
\]
在商群中的 $C+1$ 个点。鸽巢原理给出 $0\le i<j\le C$，使两点落入同一个
$U$-小块，于是存在 $\beta\in\OF$ 使
\[
(j-i)\iota(\alpha)-\iota(\beta)\in U.
\]
令 $t=j-i$，则 $1\le |t|\le C$，且由 $U$ 的选择得到
\[
|\Nm(t\alpha-\beta)|<1.
\]
\end{proof}

\begin{exercise}
设 $F=\Q(\sqrt d)$ 为实二次域。证明单位群形如 $\{\pm\varepsilon^n:n\in\Z\}$，并给出理想计数渐近公式。
\end{exercise}

\begin{proof}[解答]
\textbf{解答状态：solvable.}
实二次域有 $r_1=2,r_2=0$。Dirichlet 单位定理给出
\[
U_F\simeq \mu_F\times\Z.
\]
实二次域的单位根只有 $\pm1$，所以 $\mu_F=\{\pm1\}$。取自由部分的一个生成元
$u$。若 $0<u<1$，把它换成 $u^{-1}$；于是可取 $\varepsilon>1$，并且
\[
U_F=\{\pm\varepsilon^n:n\in\Z\}.
\]

对模 $\mathfrak m=1$ 的类数公式残数，
\[
\lim_{s\to1}(s-1)\zeta_F(s)
=\frac{2^{2}R_Fh_F}{w_F\sqrt{|d_F|}}.
\]
这里 $w_F=2$。实二次域的对数嵌入为
\[
u\mapsto(\log|\iota_1(u)|,\log|\iota_2(u)|),
\]
且范数为 $\pm1$ 给出两个坐标和为 $0$。若取生成元 $\varepsilon>1$，调控子为
$R_F=\log\varepsilon$。因此残数为
\[
\frac{2h_F\log\varepsilon}{\sqrt{|d_F|}}.
\]
Dedekind $\zeta$ 函数的系数部分和正是 $N(x;F)$，由本章 Dirichlet 级数残数命题得到
\[
\lim_{x\to\infty}\frac{N(x;F)}x
=\frac{2h_F\log\varepsilon}{\sqrt{|d_F|}}.
\]
\end{proof}

\begin{exercise}
设 $F$ 为虚二次域。证明
\[
\lim_{x\to\infty}\frac{N(x;F)}x
=\frac{2h_F\pi}{w\sqrt{|d_F|}}.
\]
\end{exercise}

\begin{proof}[解答]
\textbf{解答状态：solvable.}
虚二次域有 $r_1=0,r_2=1$。单位群秩为
$r_1+r_2-1=0$，所以没有自由单位方向；调控子按空行列式约定为
$R_F=1$。类数公式残数为
\[
\frac{2^{0}(2\pi)^1R_Fh_F}{w\sqrt{|d_F|}}
=\frac{2\pi h_F}{w\sqrt{|d_F|}}.
\]
另一方面
\[
\zeta_F(s)=\sum_{\mathfrak a\ne0}\Nm(\mathfrak a)^{-s}
\]
的系数部分和就是 $N(x;F)$。由 Dirichlet 级数残数命题，残数等于
$N(x;F)/x$ 的极限，于是得到题中公式。
\end{proof}

\begin{exercise}
设 $E/F$ 为数域有限 Galois 扩张。证明在 $E$ 中完全分裂的 $F$ 素理想集合密度为 $1/[E:F]$。
\end{exercise}

\begin{proof}[解答]
\textbf{解答状态：solvable.}
删去有限多个分歧素理想不会改变 Dirichlet 密度。对未分歧素理想
$\mathfrak p$，它在 $E$ 中完全分裂，当且仅当任取
$\mathfrak P|\mathfrak p$ 都有分解群平凡，也等价于 Frobenius 共轭类为
\(\{1\}\)。Chebotarev 密度定理给出该集合密度为
\[
\frac{|\{1\}|}{|\Gal(E/F)|}=\frac1{[E:F]}.
\]
\end{proof}

\begin{exercise}
固定整数 $d$。证明满足 $\left(\frac d p\right)=1$ 的素数 $p$ 的 Dirichlet 密度为 $1/2$。
\end{exercise}

\begin{proof}[解答]
\textbf{解答状态：partially solvable.}
若 $d$ 是非零非平方整数，则结论成立；若 $d$ 是平方数，则除有限多个素数外
\(\left(\frac dp\right)=1\)，密度为 $1$，所以当前题面需要这个非平方假设。

在 $d$ 非平方时，令 $K=\Q(\sqrt d)$。去掉整除 $2d$ 的素数后，二次互反理论中的基本判别准则给出
\[
\left(\frac d p\right)=1
\]
等价于 $p$ 在 $K$ 中分裂。由于 $K/\Q$ 是二次 Galois 扩张，分裂等价于 Frobenius 为恒等元。完全分裂素数密度为
\[
\frac1{[K:\Q]}=\frac12.
\]
\end{proof}

\begin{exercise}
取素数 $q$。证明满足
\[
p\equiv1\pmod q,\qquad 2^{(p-1)/q}\equiv1\pmod p
\]
的素数 $p$ 的密度为 $1/(q(q-1))$。
\end{exercise}

\begin{proof}[解答]
\textbf{解答状态：solvable.}
令
\[
E=\Q(\zeta_q,\sqrt[q]{2}).
\]
先计算次数。$\Q(\zeta_q)/\Q$ 的次数为 $q-1$。若 $q=2$，则
$E=\Q(\sqrt2)$，次数为 $2=q(q-1)$。若 $q$ 为奇素数，说明
$X^q-2$ 在 $\Q(\zeta_q)$ 上不可约即可。若 $2=a^q$ 对某个
$a\in\Q(\zeta_q)$ 成立，取范数到 $\Q$ 得
\[
2^{q-1}=N_{\Q(\zeta_q)/\Q}(a)^q.
\]
右边是有理数的 $q$ 次幂，而左边素因子 $2$ 的指数为 $q-1$，不被
$q$ 整除，矛盾。Kummer 理论或 Eisenstein--Kummer 判定方法说明
$X^q-2$ 在含 $\zeta_q$ 的域上给出 $q$ 次扩张，所以
\[
[E:\Q]=q(q-1).
\]

对 $p\nmid 2q$，条件 $p\equiv1\pmod q$ 等价于 $p$ 在
\(\Q(\zeta_q)\) 中完全分裂，也就是 $\F_p^\times$ 含全部 $q$ 次单位根。在此条件下，
\[
2^{(p-1)/q}\equiv1\pmod p
\]
等价于 $2$ 在 $\F_p^\times$ 中是 $q$ 次幂；因为 $\F_p^\times$ 是循环群，元素为
$q$ 次幂当且仅当它的 $(p-1)/q$ 次幂为 $1$。这又等价于多项式
$X^q-2$ 在 $\F_p$ 中完全分裂。合起来，两个同余条件等价于 $p$ 在
$E$ 中完全分裂。由完全分裂密度定理，所求密度为
\[
\frac1{[E:\Q]}=\frac1{q(q-1)}.
\]
\end{proof}

\begin{exercise}
设 $f(x)\in\Z[x]$ 是 $n$ 次不可约首一多项式。按当前表述，它要求证明分裂域是 $\Q$ 的 $n$ 次扩张，并推出模 $p$ 有根的素数密度为 $1/n$。
\end{exercise}

\begin{proof}[解答]
\textbf{解答状态：partially solvable.}
按这个题面，命题不成立，因此不能给出证明。反例是
\[
f(x)=x^3-2.
\]
它在 $\Q$ 上不可约，次数 $n=3$，但分裂域为
\[
\Q(\sqrt[3]{2},\zeta_3),
\]
次数为 $6$，不是 $3$。

第二个断言也不按字面成立。一般地，设 $E$ 为 $f$ 的分裂域，$G=\Gal(E/\Q)$。由于 $f$ 不可约，$G$ 传递作用在 $f$ 的 $n$ 个根上。对未分歧素数 $p$，$f\bmod p$ 有根，当且仅当 Frobenius 在这个置换作用中有不动点。因此密度应为
\[
\frac{|\{g\in G:g\text{ 在根集上有不动点}\}|}{|G|},
\]
不一定是 $1/n$。例如 $x^3-2$ 的 Galois 群是 $S_3$，有不动点的元素为恒等元和三个换位，共 $4$ 个，密度为 $4/6=2/3$。

若补充更强假设：分裂域本身就是次数 $n$ 的 Galois 扩张，且 $G$ 在根上正则作用，那么非恒等元素没有不动点，只有恒等 Frobenius 给出根。此时密度才是 $1/|G|=1/n$。
\end{proof}
 \chapter{完备域}

本章从“赋值”进入局部数论。第一章和第二章已经说明：一个数域中的素理想可以分解，理想可以计数，单位群可以放进对数格。现在我们把目光集中在一个素理想附近，问：如果只关心这个素理想所控制的大小，域里会出现怎样的拓扑、极限、完备化和扩张结构？

最基本的例子是 $\Q$ 的 $p$ 进赋值。若
\[
x=p^{v_p(x)}\frac ab,\qquad a,b\in\Z,\quad p\nmid a b,
\]
则 $v_p(x)$ 记录 $x$ 中含有多少个 $p$。令 $|x|_p=p^{-v_p(x)}$，就得到 $p$ 进绝对值。完成化后得到 $\Q_p$。本章的核心就是把这个例子推广到一般赋值域、局部数域和 Lubin--Tate 形式群。

\section{赋值域}

\subsection{赋值}

\begin{definition}
设 $F$ 为域。函数
\[
v:F\longrightarrow \R\cup\{\infty\}
\]
称为 $F$ 的非 Archimede 指数赋值，若满足：
\begin{enumerate}[label=(\arabic*)]
\item $v(x)=\infty$ 当且仅当 $x=0$；
\item $v(xy)=v(x)+v(y)$；
\item $v(x+y)\ge\min\{v(x),v(y)\}$。
\end{enumerate}
带赋值的域 $(F,v)$ 称为赋值域。常把 $v$ 记作 $v_F$ 或 $\operatorname{ord}_F$。
\end{definition}

\begin{remark}
赋值越大，元素越接近 $0$。这和通常绝对值的方向相反：在 $p$ 进情形中，$v_p(p^n)=n$ 越大，$p^n$ 越小。
\end{remark}

\begin{lemma}
若 $v(x)>v(y)$，则
\[
v(x+y)=v(y).
\]
\end{lemma}

\begin{proof}
由赋值三角不等式，
\[
v(x+y)\ge \min\{v(x),v(y)\}=v(y).
\]
另一方面，
\[
y=(x+y)-x,
\]
所以
\[
v(y)\ge \min\{v(x+y),v(x)\}.
\]
若 $v(x+y)>v(y)$，则右边最小值大于 $v(y)$，矛盾。因此只能有 $v(x+y)=v(y)$。
\end{proof}

\begin{example}
设 $\mathcal O$ 为 Dedekind 环，$F=\operatorname{Frac}(\mathcal O)$。对非零素理想 $\mathfrak p$，任意 $x\in F^\times$ 的主分式理想唯一分解为
\[
x\mathcal O=\prod_{\mathfrak q}\mathfrak q^{v_{\mathfrak q}(x)}.
\]
取 $\mathfrak p$ 处的指数 $v_{\mathfrak p}(x)$，并令 $v_{\mathfrak p}(0)=\infty$，就得到 $F$ 的赋值。特别地，$\mathcal O=\Z$ 时，素数 $p$ 给出 $\Q$ 的 $p$ 进赋值。
\end{example}

\begin{proof}[例子的验证]
需要检查三条赋值公理。第一条由定义给出。第二条来自主分式理想乘法：
\[
(xy)\mathcal O=(x\mathcal O)(y\mathcal O)
=\prod_{\mathfrak q}\mathfrak q^{v_{\mathfrak q}(x)+v_{\mathfrak q}(y)},
\]
所以 $\mathfrak p$ 的指数满足
\[
v_{\mathfrak p}(xy)=v_{\mathfrak p}(x)+v_{\mathfrak p}(y).
\]
第三条来自理想包含关系。若
\[
x\mathcal O\subset \mathfrak p^a,\qquad y\mathcal O\subset \mathfrak p^a,
\]
则 $(x+y)\mathcal O\subset \mathfrak p^a$。取
$a=\min\{v_{\mathfrak p}(x),v_{\mathfrak p}(y)\}$，得到
\[
v_{\mathfrak p}(x+y)\ge a.
\]
这说明 Dedekind 环中素理想分解的指数就是一个非 Archimede 赋值。对
\(\Q\) 而言，这正是把非零有理数写成
\[
x=p^{v_p(x)}\frac ab,\qquad p\nmid a,\ p\nmid b,
\]
时得到的 $p$ 进指数。
\end{proof}

\begin{proposition}
设 $(F,v)$ 为赋值域。定义
\[
\mathcal O=\{x\in F:v(x)\ge0\},\qquad
\mathfrak p=\{x\in F:v(x)>0\}.
\]
则：
\begin{enumerate}[label=(\arabic*)]
\item $\mathcal O$ 是整闭整环，且 $F$ 是 $\mathcal O$ 的分式域；
\item $\mathfrak p$ 是 $\mathcal O$ 的唯一极大理想；
\item $\mathcal O^\times=\{x\in F:v(x)=0\}=\mathcal O\setminus\mathfrak p$。
\end{enumerate}
\end{proposition}

\begin{proof}
若 $x,y\in\mathcal O$，则
\[
v(xy)=v(x)+v(y)\ge0,\qquad
v(x+y)\ge\min\{v(x),v(y)\}\ge0,
\]
所以 $\mathcal O$ 是子环。若 $x\in F^\times$，则 $v(x)$ 与 $v(x^{-1})=-v(x)$ 至少有一个非负，所以 $x$ 或 $x^{-1}$ 属于 $\mathcal O$，从而 $F$ 是 $\mathcal O$ 的分式域。

证明整闭性。设 $x\in F$ 在 $\mathcal O$ 上整，即
\[
x^n+a_{n-1}x^{n-1}+\cdots+a_0=0,\qquad a_i\in\mathcal O.
\]
若 $x\notin\mathcal O$，则 $v(x)<0$，所以 $x^{-1}\in\mathcal O$。把方程除以 $x^{n-1}$ 得
\[
x=-(a_{n-1}+a_{n-2}x^{-1}+\cdots+a_0x^{-(n-1)}).
\]
右边在 $\mathcal O$ 中，矛盾。因此 $x\in\mathcal O$。

若 $x\in\mathcal O$ 可逆，存在 $y\in\mathcal O$ 使 $xy=1$，于是
\[
0=v(1)=v(x)+v(y),
\]
且两项都非负，故 $v(x)=0$。反过来若 $v(x)=0$，则 $v(x^{-1})=0$，所以 $x^{-1}\in\mathcal O$。因此单位群就是 $\mathcal O\setminus\mathfrak p$。非单位全体形成唯一极大理想 $\mathfrak p$。

最后把“唯一极大”拆开看。$\mathfrak p$ 对加法封闭：若 $v(x)>0$、$v(y)>0$，则
$v(x+y)\ge\min\{v(x),v(y)\}>0$。若 $a\in\mathcal O$、$x\in\mathfrak p$，则
$v(ax)=v(a)+v(x)>0$，故 $ax\in\mathfrak p$。所以 $\mathfrak p$ 是理想。商环
\(\mathcal O/\mathfrak p\) 中非零类由 $v(x)=0$ 的元素给出，而这样的元素在
\(\mathcal O\) 中可逆，所以商环是域，$\mathfrak p$ 是极大理想。任意极大理想不能含有单位，因此只能由非单位组成；非单位全体正是 $\mathfrak p$，所以所有极大理想都包含于
\(\mathfrak p\)，极大性迫使它等于 $\mathfrak p$。
\end{proof}

\begin{example}
对 $\Q$ 的 $p$ 进赋值，赋值环是
\[
\Z_{(p)}=\left\{\frac ab:a,b\in\Z,\ p\nmid b\right\}.
\]
确实，$v_p(a/b)\ge0$ 等价于分母中没有留下 $p$ 的负指数。它的唯一极大理想是
\[
p\Z_{(p)}=\left\{\frac ab:a,b\in\Z,\ p\mid a,\ p\nmid b\right\},
\]
剩余域是 $\Z_{(p)}/p\Z_{(p)}\simeq\F_p$。
\end{example}

\begin{definition}
$\mathcal O$ 称为 $(F,v)$ 的赋值环，$\mathfrak p$ 称为赋值环的极大理想，剩余域记为
\[
\kappa=\mathcal O/\mathfrak p.
\]
若值群 $v(F^\times)$ 是某个 $r\Z$，则称 $v$ 为离散赋值。若 $v(F^\times)=\Z$，称为标准离散赋值。此时满足 $v(\pi)=1$ 的 $\pi\in\mathcal O$ 称为素元或单化子。
\end{definition}

以下记
\[
U^{(0)}=\mathcal O^\times,\qquad U^{(n)}=1+\mathfrak p^n\quad(n\ge1).
\]

\begin{proposition}
设 $(F,v)$ 为标准离散赋值域，$\pi$ 为素元。
\begin{enumerate}[label=(\arabic*)]
\item 任意 $x\in F^\times$ 可唯一写成
\[
x=u\pi^n,\qquad u\in\mathcal O^\times,\quad n=v(x)\in\Z.
\]
\item $\mathcal O$ 是主理想环，且非零理想全为
\[
\mathfrak p^n=\pi^n\mathcal O=\{x\in F:v(x)\ge n\},\qquad n\ge0.
\]
\item $\mathfrak p^n/\mathfrak p^{n+1}\simeq\mathcal O/\mathfrak p$。
\item $U^{(0)}\supset U^{(1)}\supset U^{(2)}\supset\cdots$。
\item 对 $n\ge1$，
\[
U^{(n)}/U^{(n+1)}\simeq\mathcal O/\mathfrak p,\qquad
U^{(0)}/U^{(n)}\simeq(\mathcal O/\mathfrak p^n)^\times.
\]
\end{enumerate}
\end{proposition}

\begin{proof}
(1) 令 $n=v(x)$。则 $u=x\pi^{-n}$ 满足 $v(u)=0$，所以 $u$ 为单位。若 $x=u\pi^n=u'\pi^{n'}$，取赋值得 $n=n'$，再得 $u=u'$。

(2) 设 $\mathfrak a$ 为非零理想。集合
\[
\{v(x):0\ne x\in\mathfrak a\}\subset\Z_{\ge0}
\]
非空，因此有最小元。取 $x\in\mathfrak a$ 使 $v(x)$ 最小，记 $v(x)=n$。由 (1) 写
$x=u\pi^n$，其中 $u$ 是单位。因为 $\mathfrak a$ 是理想且 $u^{-1}\in\mathcal O$，
\[
\pi^n=u^{-1}x\in\mathfrak a.
\]
于是 $\pi^n\mathcal O\subset\mathfrak a$。反过来，任意 $y\in\mathfrak a$ 有
$v(y)\ge n$，写 $y=w\pi^{v(y)}$，则
\[
y=(w\pi^{v(y)-n})\pi^n\in\pi^n\mathcal O.
\]
所以 $\mathfrak a=\pi^n\mathcal O$。等式
\[
\pi^n\mathcal O=\{x\in F:v(x)\ge n\}
\]
也由同一个分解 $x=u\pi^{v(x)}$ 得到。

(3) 映射
\[
\mathcal O\to \mathfrak p^n/\mathfrak p^{n+1},\qquad
a\mapsto a\pi^n+\mathfrak p^{n+1}
\]
满射，核为 $\mathfrak p$。满射是因为每个 $\mathfrak p^n$ 中的元素都能写成
$a\pi^n$；核为 $\mathfrak p$ 是因为
$a\pi^n\in\mathfrak p^{n+1}$ 等价于 $v(a)\ge1$。第一同构定理给出
\[
\mathfrak p^n/\mathfrak p^{n+1}\simeq\mathcal O/\mathfrak p.
\]

(4) 由 $\mathfrak p^{n+1}\subset\mathfrak p^n$ 直接得到。

(5) 映射
\[
U^{(n)}\to\mathcal O/\mathfrak p,\qquad
1+a\pi^n\mapsto a+\mathfrak p
\]
是从乘法群到加法群的同态，因为
\[
(1+a\pi^n)(1+b\pi^n)
=1+(a+b)\pi^n+ab\pi^{2n},
\]
且 $2n\ge n+1$，最后一项在模 $\mathfrak p^{n+1}$ 后消失。核由
$a\in\mathfrak p$ 描述，正是 $U^{(n+1)}$。

第二个同构由自然映射
\[
\mathcal O^\times\longrightarrow(\mathcal O/\mathfrak p^n)^\times
\]
给出，核正是 $1+\mathfrak p^n$。它满射的理由如下：若
$\bar a\in(\mathcal O/\mathfrak p^n)^\times$，则它模 $\mathfrak p$ 的像非零，所以
$a\notin\mathfrak p$。在赋值环中不在 $\mathfrak p$ 就是单位，故 $a\in\mathcal O^\times$，于是每个可逆剩余类都有单位代表。
\end{proof}

\subsection{绝对值}

赋值适合乘法，绝对值适合拓扑和极限。给定赋值 $v$ 与实数 $q>1$，定义
\[
|0|=0,\qquad |x|=q^{-v(x)}\quad(x\ne0).
\]
则 $|xy|=|x||y|$，且由赋值三角不等式得到更强的非 Archimede 不等式
\[
|x+y|\le\max\{|x|,|y|\}.
\]

\begin{definition}
域 $F$ 上的绝对值是函数 $|\cdot|:F\to\R_{\ge0}$，满足：
\[
|x|=0\Longleftrightarrow x=0,\qquad
|xy|=|x||y|,\qquad
|x+y|\le |x|+|y|.
\]
若满足更强的 $|x+y|\le\max\{|x|,|y|\}$，称为非 Archimede 绝对值。
\end{definition}

用绝对值定义收敛和 Cauchy 序列：$a_n\to a$ 是指 $|a_n-a|\to0$；若任意 $\varepsilon>0$，存在 $N$ 使 $m,n\ge N$ 时 $|a_m-a_n|<\varepsilon$，则称 $\{a_n\}$ 为 Cauchy 序列。若每个 Cauchy 序列都在 $F$ 中收敛，则称 $F$ 完备。

\begin{proposition}
设域 $F$ 带绝对值 $|\cdot|$。
\begin{enumerate}[label=(\arabic*)]
\item 存在完备赋值域 $\widehat F$，使 $F$ 等距嵌入 $\widehat F$ 且像稠密。
\item 若另有 $\widehat F'$ 满足同样性质，则存在唯一保持 $F$ 且保持绝对值的同构 $\widehat F\simeq\widehat F'$。
\end{enumerate}
\end{proposition}

\begin{proof}
令 $C$ 为 $F$ 中所有 Cauchy 序列的集合，$Z$ 为趋于 $0$ 的 Cauchy 序列集合。整个构造的想法是：把“应该有极限但原域中没有极限”的 Cauchy 序列当作新元素；两个序列若差趋于 $0$，就表示同一个新元素。

先记录两个基础事实。第一，每个 Cauchy 序列有界。确切地说，若 $\{a_n\}$ 为 Cauchy 序列，取 $N$ 使 $n\ge N$ 时 $|a_n-a_N|<1$，则
\[
|a_n|\le |a_N|+1\quad(n\ge N),
\]
前面有限多个项再取最大值，便得到统一上界。第二，若 $\{a_n\}$ 不是零序列，则它最终离开 $0$。否则对每个 $\varepsilon>0$ 都有任意大的 $n$ 使 $|a_n|<\varepsilon$，再由 Cauchy 性可推出整个尾部趋于 $0$，与“不属于 $Z$”矛盾。因此存在 $t>0$ 与 $N$，使 $n\ge N$ 时 $|a_n|>t$。

逐项加法和乘法使 $C$ 成为交换环。加法的 Cauchy 性来自
\[
|(a_m+b_m)-(a_n+b_n)|
\le |a_m-a_n|+|b_m-b_n|.
\]
乘法用刚才的有界性：若 $|a_n|\le R$、$|b_n|\le T$，则
\[
|a_mb_m-a_nb_n|
\le |a_m||b_m-b_n|+|b_n||a_m-a_n|
\le R|b_m-b_n|+T|a_m-a_n|,
\]
右边可由 Cauchy 性任意变小。

$Z$ 是理想。若 $\{a_n\}\in Z$、$\{c_n\}\in C$，由 $\{c_n\}$ 有界得
\[
|a_nc_n|\le R|a_n|\to0,
\]
所以乘积仍在 $Z$ 中。若 $\{a_n\}\notin Z$，把前有限项改成 $1$，得到序列
\[
b_n=\begin{cases}
1,&n<N,\\
a_n,&n\ge N,
\end{cases}
\]
则 $\{a_n\}-\{b_n\}\in Z$，且 $|b_n|\ge t$ 对充分大的 $n$ 成立。于是
\[
\left|\frac1{b_m}-\frac1{b_n}\right|
=\frac{|b_m-b_n|}{|b_mb_n|}
\le t^{-2}|b_m-b_n|
\]
对充分大的 $m,n$ 成立，故 $\{b_n^{-1}\}$ 是 Cauchy 序列。于是
\[
\{b_n\}\{b_n^{-1}\}=1
\]
说明 $\{a_n\}$ 在 $C/Z$ 中可逆。任何严格包含 $Z$ 的理想都会含有这样一个非零类，从而含有 $1$，所以 $Z$ 是极大理想，$\widehat F=C/Z$ 是域。

定义
\[
\left|\,[a_n]\,\right|=\lim_{n\to\infty}|a_n|.
\]
这个极限存在，因为
\[
\bigl||a_m|-|a_n|\bigr|\le |a_m-a_n|
\]
使 $\{|a_n|\}$ 成为实数 Cauchy 序列。定义与代表元无关，因为若
$[a_n]=[b_n]$，则 $a_n-b_n\to0$，而
\[
\bigl||a_n|-|b_n|\bigr|\le |a_n-b_n|.
\]
绝对值的三条性质都可逐项取极限得到。例如
\[
|[a_n][b_n]|=\lim |a_nb_n|
=\left(\lim |a_n|\right)\left(\lim |b_n|\right).
\]
三角不等式同样由
$|a_n+b_n|\le |a_n|+|b_n|$ 取极限得到。

常值序列给出嵌入 $F\hookrightarrow C/Z$，且保持绝对值。若
$\alpha=[a_n]\in\widehat F$，把 $a_m$ 看作常值序列，则
\[
|\alpha-a_m|=\lim_{n\to\infty}|a_n-a_m|.
\]
给定 $\varepsilon>0$，由 $\{a_n\}$ 的 Cauchy 性，取 $m$ 足够大可使右边小于
$\varepsilon$。所以 $F$ 在 $\widehat F$ 中稠密。

完备性用对角线法，但要选好代表。设 $\alpha_j\in\widehat F$ 是 Cauchy 序列。因为 $F$ 稠密，可取 $b_j\in F$ 使
\[
|\alpha_j-b_j|<2^{-j}.
\]
则 $\{b_j\}$ 是 $F$ 中的 Cauchy 序列：当 $i,j$ 很大时，
\[
|b_i-b_j|\le |b_i-\alpha_i|+|\alpha_i-\alpha_j|+|\alpha_j-b_j|
\]
三项都可变小。令 $\alpha=[b_j]\in\widehat F$。再由
\[
|\alpha_j-\alpha|\le |\alpha_j-b_j|+|b_j-\alpha|
\]
得到 $\alpha_j\to\alpha$。所以 $\widehat F$ 完备。

唯一性由稠密性给出：若 $x=\lim a_n$，任何等距延拓同构都必须把 $x$ 送到 $\lim a_n$。
\end{proof}

\begin{proposition}
设 $(\widehat F,\widehat v)$ 为赋值域 $(F,v)$ 的完备化。则
\[
\widehat v(\widehat F^\times)=v(F^\times),
\qquad
\kappa_{\widehat F}\simeq\kappa_F.
\]
\end{proposition}

\begin{proof}
若 $x=[a_n]\ne0$，则 Cauchy 序列 $\{a_n\}$ 不趋于 $0$。于是存在 $M$ 与 $N_1$，使 $n\ge N_1$ 时 $v(a_n)<M$；也就是说这些 $a_n$ 不会无限趋向于 $0$。另一方面，Cauchy 性给出 $N_2$，使 $m,n\ge N_2$ 时
\[
v(a_m-a_n)>M.
\]
固定 $N=\max\{N_1,N_2\}$。若 $n\ge N$，则
\[
a_n=(a_n-a_N)+a_N,
\]
且 $v(a_n-a_N)>M>v(a_N)$。由赋值的强三角性质中“较小项决定和”的结论，
\[
v(a_n)=v(a_N).
\]
所以非零 Cauchy 序列的赋值最终稳定。定义
\[
\widehat v([a_n])=\lim_{n\to\infty}v(a_n)
\]
就有意义。

若换一个代表 $b_n$，则 $a_n-b_n\to0$。对充分大的 $n$，
\[
v(a_n-b_n)>\max\{v(a_n),v(b_n)\},
\]
于是同一条“较小项决定和”给出 $v(a_n)=v(b_n)$。所以 $\widehat v$ 与代表无关。它满足赋值公理，因为逐项乘法和逐项加法后再取最终稳定值即可。并且每个
\(\widehat v(x)\) 等于某个充分大的 $v(a_N)$，已经属于 $v(F^\times)$；反向包含由常值序列给出。因此
\[
\widehat v(\widehat F^\times)=v(F^\times).
\]

对剩余域，记 $\widehat{\mathcal O}$、$\widehat{\mathfrak p}$ 为完备化中的赋值环和极大理想。取
$x\in\widehat{\mathcal O}$。由稠密性可取 $a\in F$ 使
\[
\widehat v(x-a)>0.
\]
因为 $\widehat v(x)\ge0$，若 $a\notin\mathcal O_F$，则 $v(a)<0$，从
\[
x=a+(x-a)
\]
会得到 $\widehat v(x)=v(a)<0$，矛盾。所以可取到的这个 $a$ 必在
\(\mathcal O_F\) 中，并且 $x$ 与 $a$ 在剩余域中同类。故自然映射
\[
\mathcal O_F\longrightarrow \widehat{\mathcal O}/\widehat{\mathfrak p}
\]
满射。它的核正是 $\mathfrak p_F$，因为 $a$ 映到 $0$ 等价于
\(\widehat v(a)>0\)，也就是 \(v(a)>0\)。由第一同构定理得到
\[
\kappa_F=\mathcal O_F/\mathfrak p_F\simeq
\widehat{\mathcal O}/\widehat{\mathfrak p}=\kappa_{\widehat F}.
\]
\end{proof}

\begin{proposition}
设 $(F,v)$ 为完备非 Archimede 赋值域，$a_n\in F$。级数
\[
\sum_{n=0}^\infty a_n
\]
在 $F$ 中收敛，当且仅当 $a_n\to0$。
\end{proposition}

\begin{proof}
若级数收敛，部分和
\[
s_N=\sum_{n=0}^N a_n
\]
收敛，所以
\[
a_{N+1}=s_{N+1}-s_N\to0.
\]
反过来，若 $a_n\to0$，则对任意给定的阈值 $M$，存在 $N$ 使 $n\ge N$ 时
$v(a_n)>M$。当 $m>n\ge N$ 时，非 Archimede 不等式给出
\[
v(s_m-s_n)
=v(a_{n+1}+\cdots+a_m)
\ge \min_{n<j\le m}v(a_j)>M.
\]
所以部分和序列是 Cauchy 序列。由于 $F$ 完备，它收敛，级数收敛。
\end{proof}

\begin{proposition}
设 $(F,v)$ 为离散赋值域，$\pi$ 为素元，$R\subset\mathcal O$ 为剩余域 $\kappa$ 的一组代表元且 $0\in R$。若 $\widehat F$ 是完备化，则任意 $x\in\widehat F^\times$ 唯一写成
\[
x=\pi^m(a_0+a_1\pi+a_2\pi^2+\cdots),
\]
其中 $m\in\Z$，$a_i\in R$，且 $a_0\ne0$。
\end{proposition}

\begin{proof}
先写 $x=\pi^m u$，其中 $m=\widehat v(x)$，$u$ 为完备赋值环的单位。这里用到上一命题：
\(\widehat v(\widehat F^\times)=v(F^\times)=\Z\)，所以 $m$ 是整数，$\pi$ 在完备化中仍是素元。

由剩余域同构，存在唯一 $a_0\in R^\times$ 使
\[
u\equiv a_0\pmod{\pi}.
\]
于是 $u-a_0\in\pi\widehat{\mathcal O}$，可写为
\[
u=a_0+\pi b_1,\qquad b_1\in\widehat{\mathcal O}.
\]
再对 $b_1$ 取剩余代表 $a_1\in R$，写
\[
b_1=a_1+\pi b_2.
\]
不断重复得到
\[
u=a_0+a_1\pi+\cdots+a_n\pi^n+\pi^{n+1}b_{n+1}.
\]
因为 $b_{n+1}\in\widehat{\mathcal O}$，所以
\[
\widehat v(\pi^{n+1}b_{n+1})\ge n+1\to\infty,
\]
余项趋于 $0$。由刚才的级数判别，级数收敛，并且
\[
u=\sum_{i=0}^\infty a_i\pi^i.
\]
乘回 $\pi^m$ 得到所需表达式。

唯一性也逐位证明。若
\[
u=\sum_{i\ge0}a_i\pi^i=\sum_{i\ge0}a_i'\pi^i,
\]
先模 $\pi$ 得 $a_0\equiv a_0'\pmod{\pi}$。两者都在代表元集合 $R$ 中，所以
$a_0=a_0'$。相减后除以 $\pi$，再模 $\pi$ 得 $a_1=a_1'$。如此归纳，全部系数相同。指数 $m$ 的唯一性由取赋值直接得到。
\end{proof}

\begin{example}
对 $(\Q,v_p)$，取 $R=\{0,1,\dots,p-1\}$，完备化为 $\Q_p$。每个非零 $x\in\Q_p$ 可唯一写成
\[
x=p^m(a_0+a_1p+a_2p^2+\cdots),\qquad a_0\ne0.
\]
\end{example}

\begin{proposition}
设 $|\cdot|_1,|\cdot|_2$ 是域 $F$ 的两个绝对值。以下条件等价：
\begin{enumerate}[label=(\arabic*)]
\item 它们定义相同拓扑；
\item 对任意 $a\in F$，$|a|_1<1$ 当且仅当 $|a|_2<1$；
\item 存在实数 $r>0$，使 $|a|_1=|a|_2^r$ 对所有 $a\in F$ 成立。
\end{enumerate}
\end{proposition}

\begin{proof}
若拓扑相同，则 $a^n\to0$ 在一个绝对值下成立当且仅当在另一个下成立；而 $a^n\to0$ 等价于 $|a|<1$。故 (1) 推 (2)。

设 (2) 成立。先说明它也保持 $=1$ 和 $>1$。若 $|a|_1>1$，则
$|a^{-1}|_1<1$，所以 $|a^{-1}|_2<1$，即 $|a|_2>1$。交换两个绝对值的角色，可得到
$|a|_2>1$ 时也有 $|a|_1>1$。若
$|a|_1=1$，既不是 $<1$ 也不是 $>1$，所以 $|a|_2=1$。

固定 $b$ 使 $|b|_1\ne1$，必要时把 $b$ 换成 $b^{-1}$，可设
$|b|_1>1$，于是 $|b|_2>1$。对任意 $a\in F^\times$ 和正整数 $m,n$，
\[
|a|_1^m<|b|_1^n
\]
等价于
\[
|a^m b^{-n}|_1<1.
\]
由 (2)，这又等价于
\[
|a^m b^{-n}|_2<1,
\]
也就是
\[
|a|_2^m<|b|_2^n.
\]
因此有理数 $\frac nm$ 同时位于两个实数
\[
\frac{\log|a|_1}{\log|b|_1},\qquad
\frac{\log|a|_2}{\log|b|_2}
\]
的同一侧。由有理数在实数中的稠密性，这两个实数相等：
\[
\frac{\log|a|_1}{\log|b|_1}
=
\frac{\log|a|_2}{\log|b|_2}.
\]
令
\[
r=\frac{\log|b|_1}{\log|b|_2}>0,
\]
便有 $|a|_1=|a|_2^r$。所以 (2) 推 (3)。

若 (3) 成立，则开球
\[
\{x:|x-a|_1<\varepsilon\}
\]
等于
\[
\{x:|x-a|_2<\varepsilon^{1/r}\},
\]
所以两个绝对值定义同一拓扑。
\end{proof}

\begin{theorem}
设 $|\cdot|_1,\dots,|\cdot|_n$ 是域 $F$ 上两两不等价的绝对值。给定 $a_1,\dots,a_n\in F$ 与 $\varepsilon>0$，存在 $x\in F$，使
\[
|x-a_i|_i<\varepsilon,\qquad 1\le i\le n.
\]
\end{theorem}

\begin{proof}
证明思想和中国剩余定理相同，只是理想同余换成拓扑小邻域。先说明关键构造。对每个 $i$，由于 $|\cdot|_i$ 与 $|\cdot|_j$ 不等价，可以找到元素在第 $i$ 个绝对值下相对很大、在第 $j$ 个绝对值下相对很小。经过有限次相乘，可构造 $c_i\in F$，使
\[
|c_i|_i>1,\qquad |c_i|_j<1\quad(j\ne i).
\]
取高次幂后，
\[
d_i=c_i^N
\]
满足 $|d_i|_i$ 很大而 $|d_i|_j$ 很小。令
\[
b_i=\frac{d_i}{1+d_i}.
\]
当 $N$ 足够大时，在第 $i$ 个绝对值下 $d_i$ 很大，所以 $b_i$ 接近 $1$；在其他绝对值下
$d_i$ 很小，所以 $b_i$ 接近 $0$。

更精确地，可使
\[
|b_i-1|_i<\delta,\qquad |b_i|_j<\delta\quad(j\ne i)
\]
同时成立。然后取
\[
x=\sum_i a_i b_i
\]
并令 $\delta$ 足够小。对第 $k$ 个绝对值，
\[
x-a_k=a_k(b_k-1)+\sum_{i\ne k}a_i b_i.
\]
右边每一项都可由 $\delta$ 控制，所以可使
\[
|x-a_k|_k<\varepsilon.
\]
这就完成同时逼近。
\end{proof}

\section{赋值域扩张}

若 $(E,w)/(F,v)$ 满足 $F\subset E$ 且 $w|_F=v$，称为赋值域扩张。赋值延拓的核心问题是：给定 $F$ 上的 $v$，它在代数扩张 $E$ 上有多少个延拓？

设 $F_v$ 为 $(F,v)$ 的完备化，取代数闭包 $\overline{F_v}$，并把 $v$ 唯一延拓到 $\overline{F_v}$。若 $E/F$ 为代数扩张，每个 $F$-嵌入 $\tau:E\to\overline{F_v}$ 给出 $E$ 上赋值
\[
w=v\circ\tau.
\]
不同嵌入可能给出同一个赋值；差异由 $F_v$-自同构控制。

\begin{proposition}
设 $(F,v)$ 为赋值域，$E/F$ 为代数扩张。
\begin{enumerate}[label=(\arabic*)]
\item $v$ 可延拓为 $E$ 上的赋值。
\item 若 $w$ 是 $E$ 上延拓 $v$ 的赋值，则存在 $F$-嵌入 $\tau:E\to\overline{F_v}$，使 $w=\bar v\circ\tau$。
\item 对 $F$-嵌入 $\tau,\tau':E\to\overline{F_v}$，有 $\bar v\circ\tau=\bar v\circ\tau'$ 当且仅当存在 $F_v$-自同构 $\sigma$，使 $\tau'=\sigma\circ\tau$。
\item 若 $E=F(\theta)$ 为有限可分扩张，$\theta$ 的最小多项式 $f$ 在 $F_v$ 上分解为不可约因子
\[
f=f_1\cdots f_r,
\]
则每个 $f_j$ 对应一个 $v$ 在 $E$ 上的延拓。
\end{enumerate}
\end{proposition}

\begin{proof}
(1) 先把 $E$ 放进一个代数闭包中。取 $F_v$ 的代数闭包
\(\overline{F_v}\)，并把 $v$ 延拓到其上的赋值 $\bar v$。任取一个
$F$-嵌入
\[
\tau:E\hookrightarrow\overline{F_v}.
\]
定义
\[
w(x)=\bar v(\tau x).
\]
因为 $\tau$ 是域同态，而 $\bar v$ 满足赋值公理，$w$ 也满足赋值公理。若
$a\in F$，则 $\tau(a)=a$，故 $w(a)=v(a)$。这给出 $v$ 在 $E$ 上的延拓。

(2) 反过来，设 $w$ 是 $E$ 上延拓 $v$ 的赋值。先完成 $E$ 得到 $E_w$。由于
$E_w$ 是含有 $F$ 的完备赋值域，$F$ 在其中的闭包可视为 $F_v$ 的一个副本。把
$E_w$ 嵌入 $\overline{F_v}$，再限制到 $E$，得到一个 $F$-嵌入
\[
\tau:E\to\overline{F_v}.
\]
完备域上的赋值延拓唯一，所以 $\bar v\circ\tau$ 在 $E$ 上等于原来的 $w$。

(3) 若 $\tau'=\sigma\tau$，由于 $\bar v$ 在 $\overline{F_v}$ 上对 $F_v$-自同构不变，二者给出同一赋值。反过来，若赋值相同，先在 $\tau(E)$ 与 $\tau'(E)$ 之间定义 $\sigma(\tau x)=\tau'x$。等距性让它连续延拓到 $\tau(E)F_v$，再由代数闭包延拓定理延拓为 $\overline{F_v}$ 的 $F_v$-自同构。

(3) 中的“等距性”具体如下。若 $z=\tau(x)$，则
\[
\bar v(z)=\bar v(\tau x)=\bar v(\tau' x),
\]
所以 $\tau x\mapsto\tau'x$ 保持赋值，也保持由赋值得到的拓扑。由于 $F$ 在
$F_v$ 中稠密，这个映射能唯一连续延拓到合成域
\(\tau(E)F_v\to\tau'(E)F_v\)。代数闭包中的嵌入延拓定理再把它延拓为
\(\overline{F_v}\) 的 $F_v$-自同构。

(4) $F$-嵌入 $\tau:E=F(\theta)\to\overline{F_v}$ 由 $\theta$ 的像决定，而
\(\tau(\theta)\) 必是 $f$ 在 \(\overline{F_v}\) 中的根。若两个根属于同一个不可约因子
$f_j$，则它们在 $F_v$ 上共轭，存在 $F_v$-自同构把一个根送到另一个根，于是由 (3) 给出同一个赋值延拓。若两个根属于不同不可约因子，则不存在这样的
$F_v$-自同构，因而给出不同延拓。于是不可约因子 $f_j$ 与赋值延拓一一对应。
\end{proof}

\begin{definition}
若 $E/F$ 有限，$w|v$，记 $E_w$ 为 $E$ 关于 $w$ 的完备化。定义
\[
e_{w|v}=[w(E^\times):v(F^\times)]
\]
为分歧指标，
\[
f_{w|v}=[\kappa_E:\kappa_F]
\]
为剩余次数。
\end{definition}

\begin{proposition}
设 $(F,v)$ 为赋值域，$E/F$ 为有限可分扩张。
\begin{enumerate}[label=(\arabic*)]
\item 有 $F_v$-代数同构
\[
E\otimes_FF_v\simeq\prod_{w|v}E_w.
\]
\item
\[
[E:F]=\sum_{w|v}[E_w:F_v].
\]
\item 对 $x\in E$，
\[
\Nm_{E/F}(x)=\prod_{w|v}\Nm_{E_w/F_v}(x),
\qquad
\Tr_{E/F}(x)=\sum_{w|v}\Tr_{E_w/F_v}(x).
\]
\item 若 $v$ 离散，则
\[
[E:F]=\sum_{w|v}e_{w|v}f_{w|v}.
\]
\end{enumerate}
\end{proposition}

\begin{proof}
写 $E=F(\theta)$，$\theta$ 的最小多项式为 $f$。可分性保证 $f$ 在
$F_v[X]$ 中没有重因子。首先有自然同构
\[
E\otimes_FF_v\simeq F_v[X]/(f).
\]
这个同构把 $\theta\otimes1$ 对应到 $X$ 的剩余类；它成立的原因是
$E\simeq F[X]/(f)$，张量到 $F_v$ 后得到
\[
F_v\otimes_FF[X]/(f)\simeq F_v[X]/(f).
\]

把 $f$ 在 $F_v[X]$ 中分解为互素不可约因子
\[
f=\prod_{w|v} f_w.
\]
上一命题说明这些不可约因子正好对应 $v$ 在 $E$ 上的延拓。中国剩余定理给出
\[
F_v[X]/(f)\simeq\prod_w F_v[X]/(f_w).
\]
每个因子 \(F_v[X]/(f_w)\) 是由 $\theta$ 在该因子中的像生成的完备赋值域，它含有 $E$ 的像，并对应该赋值 $w$；按完备化的唯一性，它就是 $E_w$。于是得到 (1)。

(2) 对 (1) 两边取 $F_v$-维数：
\[
\dim_{F_v}(E\otimes_FF_v)=\dim_F E=[E:F],
\]
而右边直积的维数为
\[
\sum_{w|v}[E_w:F_v].
\]
所以得到维数公式。

(3) 设 $m_x$ 是 $E$ 作为 $F$-向量空间上“乘以 $x$”的线性变换。张量到
$F_v$ 后，$m_x\otimes1$ 在
$E\otimes_FF_v$ 上的矩阵与原矩阵相同，只是把系数看作 $F_v$ 中元素，所以它的行列式和迹仍分别是
\(\Nm_{E/F}(x)\) 与 \(\Tr_{E/F}(x)\)。在直积分解
\[
E\otimes_FF_v\simeq\prod_{w|v}E_w
\]
下，这个线性变换分解为各个 $E_w$ 上的“乘以 $x$”。直积上线性变换的行列式为各分量行列式的乘积，迹为各分量迹的和，因此
\[
\Nm_{E/F}(x)=\prod_{w|v}\Nm_{E_w/F_v}(x),
\qquad
\Tr_{E/F}(x)=\sum_{w|v}\Tr_{E_w/F_v}(x).
\]

(4) 若 $v$ 离散，则每个 $E_w/F_v$ 是完备离散赋值域的有限扩张。完备离散赋值域有限扩张满足基本等式
\[
[E_w:F_v]=e_{w|v}f_{w|v}.
\]
代入 (2)，得到
\[
[E:F]=\sum_{w|v}e_{w|v}f_{w|v}.
\]
\end{proof}

\section{完备域扩张}

完备性让代数方程可以由模 $\mathfrak p$ 的近似解提升为真正解，这就是 Hensel 引理。它是 $p$ 进分析的基本工具。

\begin{lemma}
设 $(F,v)$ 为完备赋值域，$\mathcal O$ 为赋值环，$\mathfrak p$ 为极大理想。若
\[
f(X)\in\mathcal O[X],\qquad
\bar f(X)=\bar g(X)\bar h(X)\in\kappa[X],
\]
其中 $\bar g,\bar h$ 互素，则存在 $g,h\in\mathcal O[X]$，使
\[
f=gh,\qquad g\equiv\bar g\pmod{\mathfrak p},\qquad
h\equiv\bar h\pmod{\mathfrak p},
\]
且 $\deg g=\deg\bar g$。
\end{lemma}

\begin{proof}
把 $g,h$ 的系数逐次近似构造。先取 $\bar g,\bar h$ 的任意提升
$g_1,h_1\in\mathcal O[X]$，并固定 $\deg g_1=\deg\bar g$。因为
\(\bar g,\bar h\) 互素，存在 $\bar A,\bar B\in\kappa[X]$ 使
\[
\bar A\bar g+\bar B\bar h=1.
\]
这句话是 Hensel 提升的代数核心：它说明模 $\mathfrak p$ 后，任意误差都可写成
\(\bar h\cdot(\text{修正 }g)+\bar g\cdot(\text{修正 }h)\) 的形式。

假设已经模 $\mathfrak p^n$ 得到
\[
f\equiv g_nh_n\pmod{\mathfrak p^n},
\]
并且 $g_n\equiv g_1$、$h_n\equiv h_1\pmod{\mathfrak p}$。要修正到模
$\mathfrak p^{n+1}$，写
\[
g_{n+1}=g_n+\Delta g,\qquad h_{n+1}=h_n+\Delta h,
\]
其中修正项系数在 $\mathfrak p^n$ 中。线性化后要解
\[
f-g_nh_n\equiv h_n\Delta g+g_n\Delta h\pmod{\mathfrak p^{n+1}}.
\]
二次项 $\Delta g\Delta h$ 的系数在 $\mathfrak p^{2n}\subset\mathfrak p^{n+1}$ 中，所以在模
$\mathfrak p^{n+1}$ 的计算中为零。将
\[
\frac{f-g_nh_n}{\mathfrak p^n}
\]
的系数模 $\mathfrak p$ 后得到一个 $\kappa[X]$ 中的误差多项式。Bezout 恒等式把它写成
\[
\bar h\cdot \overline{\Delta g/\mathfrak p^n}
+\bar g\cdot \overline{\Delta h/\mathfrak p^n}.
\]
提升这些系数，就得到所需的 $\Delta g,\Delta h$。同时可通过除以 $g_n$ 后取余的方式控制
\(\deg g_{n+1}=\deg\bar g\)。

这样得到的系数序列在 $\mathfrak p$-进拓扑下逐级稳定，因此是 Cauchy 序列。完备性保证系数收敛到 $g,h\in\mathcal O[X]$。乘法连续，故
$f=gh$；逐级构造又保留了
$g\equiv\bar g$、$h\equiv\bar h\pmod{\mathfrak p}$ 与 $\deg g=\deg\bar g$。
\end{proof}

\begin{corollary}
设 $f(X)=a_0+\cdots+a_nX^n\in\mathcal O[X]$，令 $v(f)=\min_i v(a_i)$。
\begin{enumerate}[label=(\arabic*)]
\item 若 $v(f)=0$、$v(a_n)>0$，且存在 $0<i<n$ 使 $v(a_i)=0$，则 $f$ 可约。
\item 若 $v(f)=0$、$f$ 不可约且 $v(a_n)>0$，则所有 $i<n$ 都满足 $v(a_i)>0$。
\item 若 $f\in\mathcal O[X]$ 不可约，则 $\bar f$ 是某个不可约多项式的幂。
\end{enumerate}
\end{corollary}

\begin{proof}
(1) 令 $j<n$ 为最大的 $v(a_j)=0$ 的指标。模 $\mathfrak p$ 后，$\bar f$ 的次数为 $j$，并且
\[
\bar f=X^r\bar q(X)
\]
其中 $\bar q(0)\ne0$，而 $r\ge1$ 来自最高次数项落入 $\mathfrak p$ 后丢失的那一段。于是
$X^r$ 与 $\bar q(X)$ 互素。Hensel 引理把这个互素分解提升成
\[
f=A(X)B(X),
\]
其中 $A,B$ 都有正次数，所以 $f$ 可约。

(2) 是 (1) 的逆否命题。

(3) 把 $\bar f$ 在 $\kappa[X]$ 中分解为不可约因子幂的乘积。若出现两个不同不可约因子，则把其中一个因子幂与剩余因子幂分成两块，它们互素。Hensel 引理会把这个互素分解提升为
$f$ 的非平凡分解，矛盾。因此 $\bar f$ 只能是某一个不可约多项式的幂。
\end{proof}

\begin{proposition}
设 $F$ 为完备赋值域，$f\in\mathcal O[X]$。若存在 $a_0\in\mathcal O$，使
\[
|f(a_0)|<|f'(a_0)|^2,
\]
则 Newton 迭代
\[
a_{n+1}=a_n-\frac{f(a_n)}{f'(a_n)}
\]
收敛到 $f$ 的根 $a\in\mathcal O$，且
\[
|a-a_0|\le \left|\frac{f(a_0)}{f'(a_0)}\right|.
\]
\end{proposition}

\begin{proof}
令
\[
c=\left|\frac{f(a_0)}{f'(a_0)^2}\right|<1.
\]
先说明迭代合法。由假设 $f'(a_0)\ne0$。若
\[
|a_n-a_0|\le \left|\frac{f(a_0)}{f'(a_0)}\right|,
\]
则 Taylor 展开给出
\[
f'(a_n)=f'(a_0)+(a_n-a_0)\gamma_n,\qquad \gamma_n\in\mathcal O
\]
并且第二项绝对值严格小于 $|f'(a_0)|$，所以
\[
|f'(a_n)|=|f'(a_0)|.
\]
因此每一步的分母都非零。

Taylor 展开给出
\[
f(a_{n+1})
=f(a_n)-f'(a_n)\frac{f(a_n)}{f'(a_n)}
\allowbreak
\beta_n\left(\frac{f(a_n)}{f'(a_n)}\right)^2,
\]
故
\[
|f(a_{n+1})|\le
\left|\frac{f(a_n)}{f'(a_n)}\right|^2.
\]
同时 $f'(a_{n+1})$ 与 $f'(a_n)$ 绝对值相同。于是
\[
\left|\frac{f(a_{n+1})}{f'(a_{n+1})^2}\right|
\le
\left|\frac{f(a_n)}{f'(a_n)^2}\right|^2.
\]
归纳得到误差按平方下降。又
\[
|a_{n+1}-a_n|
=\left|\frac{f(a_n)}{f'(a_n)}\right|
\]
趋于 $0$，在非 Archimede 完备域中，级数
\[
a_0+\sum_{n\ge0}(a_{n+1}-a_n)
\]
收敛，所以 $\{a_n\}$ 收敛。设极限为 $a$。多项式函数连续，故
$f(a)=0$。非 Archimede 三角不等式给出
\[
|a-a_0|\le \max_n |a_{n+1}-a_n|
=\left|\frac{f(a_0)}{f'(a_0)}\right|.
\]
\end{proof}

\begin{proposition}
设 $(F,v)$ 为完备赋值域，$E/F$ 为代数扩张。则 $v$ 唯一延拓为 $E$ 上的赋值。若 $E/F$ 有限，则 $E$ 也是完备赋值域，且对 $x\in E$，
\[
w(x)=\frac1{[E:F]}v(\Nm_{E/F}(x)).
\]
\end{proposition}

\begin{proof}
先看有限扩张。若 $x\in E^\times$，范数 $N_{E/F}(x)$ 属于 $F^\times$，定义
\[
w(x)=\frac1{[E:F]}v(N_{E/F}(x)),\qquad w(0)=\infty.
\]
乘法性来自范数乘法性。若 $a\in F^\times$，则
\[
N_{E/F}(a)=a^{[E:F]},
\]
所以 $w(a)=v(a)$。

三角不等式可通过完备性得到。把 $E$ 看成有限维 $F$-向量空间，取一个
$F$-基。所有由 $F$ 上绝对值诱导的有限维向量空间范数给出同一拓扑；乘法
$E\times E\to E$ 连续，且 $E^\times$ 中求逆在远离 $0$ 的地方连续。由标准的完备赋值域有限维线性代数，$E$ 上的绝对值延拓唯一，所以上式给出的 $w$ 是唯一延拓。有限维完备性同样说明：$E$ 中 Cauchy 序列逐坐标为 $F$ 中 Cauchy 序列，因 $F$ 完备而收敛，所以 $E$ 完备。

若 $E/F$ 是任意代数扩张，则任一元素落在某个有限子扩张中。对有限子扩张按上面唯一延拓，若两个有限子扩张相交于更小扩张，唯一性保证定义一致。于是可在整个
$E$ 上拼接出唯一延拓。

若 $\sigma:E\to E$ 是 $F$-同构，则 $w\circ\sigma$ 也是 $v$ 的延拓；唯一性给出
$w\circ\sigma=w$。因此 $\sigma$ 保持 $w\ge0$ 与 $w>0$ 的条件，即
\[
\sigma(\mathcal O_E)=\mathcal O_E,\qquad
\sigma(\mathfrak p_E)=\mathfrak p_E.
\]
\end{proof}

\begin{proposition}
设 $(E,w)/(F,v)$ 为完备赋值域有限扩张，分歧指标为 $e$，剩余次数为 $f$。则
\[
ef\le [E:F].
\]
若 $v$ 离散，则
\[
ef=[E:F].
\]
\end{proposition}

\begin{proof}
取 $x_1,\dots,x_f\in\mathcal O_E$，其剩余类构成 $\kappa_E/\kappa_F$ 的基。再取 $\pi_1,\dots,\pi_e\in E^\times$，使 $w(\pi_i)$ 代表 $w(E^\times)/v(F^\times)$ 的不同陪集。若
\[
\sum_{i,j}c_{ij}x_j\pi_i=0,\qquad c_{ij}\in F,
\]
比较各项赋值。对固定 $i$，若某个 $c_{ij}$ 非零，令
$m_i=\min_j v(c_{ij})$，除以一个达到最小值的系数后看剩余域。因为
\(\bar x_1,\dots,\bar x_f\) 线性无关，最小赋值项不会在剩余域中相消，所以
\[
w\left(\sum_jc_{ij}x_j\right)=m_i.
\]
于是第 $i$ 块
\[
\pi_i\sum_jc_{ij}x_j
\]
的赋值属于陪集 $w(\pi_i)+v(F^\times)$。不同 $i$ 给出不同陪集，所以这些块的赋值互不相同。赋值互不相同的若干项相加时，和的赋值等于最小赋值，不可能等于
$\infty$。因此若总和为 $0$，每一块都必须为 $0$，再由剩余域基的线性无关性得所有
$c_{ij}=0$。所以 $ef$ 个元素 $x_j\pi_i$ 在 $F$ 上线性无关，得到
$ef\le [E:F]$。

若 $v$ 离散，取 $E$ 的素元 $\Pi$，$F$ 的素元 $\pi$。有
\[
\pi=u\Pi^e,\qquad u\in\mathcal O_E^\times.
\]
任意 $E$ 元素可乘以 $F$ 中适当的 $\pi$ 幂，把它的赋值调整到代表
$0,1,\dots,e-1$ 的区间。单位部分再由剩余域基 $x_1,\dots,x_f$ 模
\(\mathfrak p_E\) 展开，逐次减去近似项后赋值升高。完备性保证这个逐次展开收敛，于是
$E$ 由
\[
x_j\Pi^i,\qquad 1\le j\le f,\ 0\le i<e
\]
张成。故 $[E:F]\le ef$。与前面的不等式合并得到等式。
\end{proof}

\begin{proposition}
设 $(E,w)/(F,v)$ 为完备离散赋值域有限扩张，$\kappa_E/\kappa_F$ 可分。则存在 $x\in\mathcal O_E$，使
\[
\mathcal O_E=\mathcal O_F[x],
\]
并且 $1,x,\dots,x^{n-1}$ 是 $\mathcal O_E$ 的 $\mathcal O_F$-基。
\end{proposition}

\begin{proof}
选 $\bar x\in\kappa_E$ 生成可分扩张 $\kappa_E/\kappa_F$，取提升
$x\in\mathcal O_E$。设 $g(X)\in\mathcal O_F[X]$ 是 $\bar x$ 的最小多项式的首一提升。因为
\(\bar x\) 可分，
\[
\bar g'(\bar x)\ne0,
\]
所以 $g'(x)$ 是 $\mathcal O_E$ 的单位。

若 $g(x)$ 已经是 $E$ 的素元，就保留这个 $x$。若 $w(g(x))\ge2$，取
$h\in\mathcal O_E$ 使 $w(h)=1$。Taylor 展开给出
\[
g(x+h)=g(x)+hg'(x)+h^2r,\qquad r\in\mathcal O_E.
\]
其中 $w(g(x))\ge2$，$w(hg'(x))=1$，$w(h^2r)\ge2$，所以
\[
w(g(x+h))=1.
\]
把 $x$ 换成 $x+h$，仍然有同一个剩余类 $\bar x$，并且 $\Pi=g(x)$ 是
$E$ 的素元。

由上一命题，元素
\[
x^i\Pi^j,\qquad 0\le i<f,\ 0\le j<e
\]
构成 $\mathcal O_E$ 的 $\mathcal O_F$-基。因为 $\Pi=g(x)$ 是 $x$ 的多项式，所有
$x^i\Pi^j$ 都属于 $\mathcal O_F[x]$，所以
\[
\mathcal O_E=\mathcal O_F[x].
\]
设 $n=[E:F]=ef$，则 $1,x,\dots,x^{n-1}$ 也可取为由该生成元给出的
\(\mathcal O_F\)-基。
\end{proof}

\begin{lemma}
设 $F$ 为完备赋值域，$\alpha,\beta$ 为 $F$ 上代数元。若对任意非恒等 $F$-嵌入 $\sigma:F(\alpha)\to F^{\mathrm{alg}}$，都有
\[
v(\alpha-\beta)>v(\alpha-\sigma\alpha),
\]
则
\[
F(\alpha)\subset F(\beta).
\]
\end{lemma}

\begin{proof}
取固定 $F(\beta)$ 的嵌入 $\tau:F(\alpha,\beta)\to F^{\mathrm{alg}}$。若 $\tau(\alpha)\ne\alpha$，则
\[
v(\tau\alpha-\beta)=v(\tau(\alpha-\beta))=v(\alpha-\beta),
\]
而
\[
\tau\alpha-\alpha=(\tau\alpha-\beta)+(\beta-\alpha).
\]
右边两项的赋值都等于 $v(\alpha-\beta)$。假设给出
\[
v(\alpha-\beta)>v(\alpha-\tau\alpha).
\]
但非 Archimede 不等式说明右边和的赋值至少为
$v(\alpha-\beta)$，即
\[
v(\tau\alpha-\alpha)\ge v(\alpha-\beta),
\]
矛盾。因此任何固定 $F(\beta)$ 的嵌入都必须固定 $\alpha$。这说明
\(\alpha\) 属于所有这类嵌入的固定域，也就是 \(F(\beta)\)，故
\[
F(\alpha)\subset F(\beta).
\]
\end{proof}

\begin{proposition}
设 $F$ 为完备赋值数域，$f\in F[X]$ 为首一不可约 $n$ 次多项式。存在常数 $M(f)$，使得任意首一 $n$ 次多项式 $g$ 若满足 $v(g-f)\ge M(f)$，则 $g$ 不可约，且 $g$ 的一个根生成与 $f$ 的根相同的扩张。
\end{proposition}

\begin{proof}
设 $f$ 的根为 $\alpha_1,\dots,\alpha_n$，令
\[
r=\min_{i\ne j}v(\alpha_i-\alpha_j).
\]
若 $g$ 与 $f$ 的系数足够接近，则 $g(\alpha_i)$ 的乘积
\[
\prod_i g(\alpha_i)=\prod_{i,j}(\alpha_i-\beta_j)
\]
赋值足够大，其中 $\beta_j$ 为 $g$ 的根。因此存在 $i,j$ 使
\[
v(\alpha_i-\beta_j)>r.
\]
Krasner 引理给出 $F(\alpha_i)\subset F(\beta_j)$。另一方面
$\beta_j$ 是 $n$ 次多项式 $g$ 的根，所以
\[
[F(\beta_j):F]\le n.
\]
而 $f$ 不可约给出 $[F(\alpha_i):F]=n$。于是
\[
n\le [F(\beta_j):F]\le n,
\]
故两者相等，$F(\alpha_i)=F(\beta_j)$。这迫使 $\beta_j$ 在 $F$ 上的最小多项式次数为
$n$，所以 $g$ 不可约，并且它的一个根生成与 $f$ 的根相同的扩张。
\end{proof}

\begin{definition}
完备赋值域有限扩张 $E/F$ 称为无分歧扩张，若
\[
e(E/F)=1,\qquad f(E/F)=[E:F],
\]
且剩余域扩张可分。若最大无分歧子扩张为 $F$，称 $E/F$ 为全分歧扩张。若剩余特征为 $p$，且去掉最大无分歧部分后分歧指数与 $p$ 互素，称为顺分歧扩张；否则出现狂分歧。
\end{definition}

\section{局部数域}

\subsection{\texorpdfstring{$\Q_p$}{Qp} 的拓扑}

局部数域的直觉来自 $\Q_p$。在实数里，两个距离都小的数相加仍然只是“较小”；
在 $p$ 进世界里，强三角不等式
\[
|x+ y|_p\le \max\{|x|_p,|y|_p\}
\]
使空间呈现树状结构：两个点只要落在同一个闭球内，便可以互相充当球心。记
\[
\begin{aligned}
U_a(r)&=\{x\in\Q_p:|x-a|_p<r\},\\
B_a(r)&=\{x\in\Q_p:|x-a|_p\le r\},\\
S_a(r)&=\{x\in\Q_p:|x-a|_p=r\}.
\end{aligned}
\]
这里 $U_a(r)$ 是开球，$B_a(r)$ 是闭球，$S_a(r)$ 是球面。下面的命题不是几何趣闻，
而是以后所有“邻域由理想给出”“单位群由过滤给出”的拓扑基础。

\begin{proposition}
对 $a\in\Q_p$、$r>0$：
\begin{enumerate}[label=(\arabic*)]
\item 若 $b\in B_a(r)$，则 $B_a(r)=B_b(r)$；
\item $B_a(r)$ 是开集；
\item $U_a(r)$ 是闭集；
\item $S_a(r)=\bigcup_{x\in S_a(r)}U_x(r)$。
\end{enumerate}
\end{proposition}

\begin{proof}
首先证明 (1)。若 $x\in B_b(r)$，则
\[
|x-a|_p\le\max\{|x-b|_p,|b-a|_p\}\le r,
\]
故 $x\in B_a(r)$，于是 $B_b(r)\subset B_a(r)$。又因 $b\in B_a(r)$ 等价于
 $|a-b|_p\le r$，所以 $a\in B_b(r)$。将刚才的包含证明用于中心 $b$ 和点 $a$，
得到 $B_a(r)\subset B_b(r)$。两个包含合在一起给出 $B_a(r)=B_b(r)$。

证明 (2)。取任意 $b\in B_a(r)$。由 (1) 得 $B_b(r)=B_a(r)$，并且
$U_b(r)\subset B_b(r)$。于是每个 $b$ 都有一个开球 $U_b(r)$ 完全落在 $B_a(r)$ 中。
这正是 $B_a(r)$ 为开集的定义。

证明 (3)。定义关系 $x\sim y$ 当且仅当 $|x-y|_p<r$。自反性来自 $|x-x|_p=0$，
对称性来自 $|x-y|_p=|y-x|_p$。若 $x\sim y$ 且 $y\sim z$，则
\[
|x-z|_p\le \max\{|x-y|_p,|y-z|_p\}<r,
\]
所以 $x\sim z$。因此这是等价关系。等价类正是开球 $U_x(r)$，所以每个等价类开。
$U_a(r)$ 的补集是所有其他等价类的并，故补集开，因而 $U_a(r)$ 闭。

证明 (4)。若 $x\in S_a(r)$ 且 $y\in U_x(r)$，则 $|x-a|_p=r$ 且 $|y-x|_p<r$。在非
Archimede 绝对值中，当两个加数的绝对值不相等时，和的绝对值等于较大者。于是
\[
|y-a|_p=|(y-x)+(x-a)|_p=r,
\]
所以 $y\in S_a(r)$。这证明 $U_x(r)\subset S_a(r)$。反向包含由 $x\in U_x(r)$ 给出：
每个球面上的点 $x$ 属于右边并集中的 $U_x(r)$。因此
$S_a(r)=\bigcup_{x\in S_a(r)}U_x(r)$。
\end{proof}

\begin{example}
以 $\Q_3$ 为例，$B_0(1)=\Z_3$。若 $b\in\Z_3$，则 $|b|_3\le1$。对任意 $x$，
$|x-b|_3\le1$ 当且仅当 $|x|_3\le1$：一个方向由强三角不等式给出，另一个方向把
$x=(x-b)+b$ 代回去。于是 $\Z_3=B_b(1)$。这说明在 $3$ 进闭单位球内，每个点都是
闭单位球的中心。
\end{example}

\subsection{局部数域的拓扑}

\begin{definition}
特征 $0$ 的完备离散赋值域 $F$ 若剩余域 $\kappa_F$ 有限，则称为局部数域。若
$|\kappa_F|=q$，常取素元 $\pi$，记赋值环为 $\mathcal O$，极大理想为
$\mathfrak p=(\pi)$，并记
\[
U^{(0)}=\mathcal O^\times,\qquad U^{(n)}=1+\mathfrak p^n\quad(n\ge1).
\]
这里 $U^{(n)}$ 是单位群中越来越小的邻域，它们后来会控制局部互反律中的范群。
\end{definition}

\begin{proposition}
设 $(F,v)$ 为完备离散赋值域，$\mathcal O$ 为赋值环，$\mathfrak p$ 为极大理想。
\begin{enumerate}[label=(\arabic*)]
\item 度量 $d(x,y)=q^{-v(x-y)}$ 使 $F$ 成为拓扑域，$F^\times$ 成为拓扑群。
\item
\[
\mathcal O\supset\mathfrak p\supset\mathfrak p^2\supset\cdots
\]
是 $0$ 的邻域基。
\item
\[
U^{(0)}\supset U^{(1)}\supset U^{(2)}\supset\cdots
\]
是 $F^\times$ 中 $1$ 的邻域基。
\item 自然映射给出拓扑环同构
\[
\mathcal O\simeq\varprojlim_n\mathcal O/\mathfrak p^n.
\]
\item 自然映射给出拓扑群同构
\[
U^{(0)}\simeq\varprojlim_n U^{(0)}/U^{(n)}.
\]
\end{enumerate}
\end{proposition}

\begin{proof}
把 $v$ 归一化为 $v(F^\times)=\Z$。对 $x,y,z\in F$，
$d(x,z)=q^{-v(x-z)}\le\max\{d(x,y),d(y,z)\}$，所以 $d$ 是非 Archimede 度量。
加法连续的原因是
\[
d((x+h),(y+k))=q^{-v((x-y)+(h-k))}
\le \max\{d(x,y),d(h,k)\}.
\]
乘法连续可由
\[
(x+h)(y+k)-xy=xk+yh+hk
\]
和强三角不等式得到；当 $x\ne0$ 且 $h$ 足够小时，$x+h=x(1+h/x)$ 仍非零，且
\[
(x+h)^{-1}-x^{-1}=-\frac{h}{x(x+h)}
\]
趋于 $0$，所以取逆连续。这证明 $F$ 是拓扑域，$F^\times$ 是拓扑群。

对 (2)，有 $\mathfrak p^n=\{x\in F:v(x)\ge n\}$。度量球
$\{x:d(0,x)<q^{-N}\}$ 等于 $\{x:v(x)>N\}$，也就是某个 $\mathfrak p^n$。
反过来每个 $\mathfrak p^n$ 本身就是 $0$ 的一个球。因此
$\mathcal O\supset\mathfrak p\supset\mathfrak p^2\supset\cdots$ 给出 $0$ 的邻域基。
对 (3)，在 $F^\times$ 中 $1$ 的小邻域正是满足 $v(u-1)\ge n$ 的单位，即
$U^{(n)}=1+\mathfrak p^n$。

证明 (4)。自然映射
\[
\Psi:\mathcal O\longrightarrow\varprojlim_n\mathcal O/\mathfrak p^n,\qquad
x\longmapsto (x\bmod\mathfrak p^n)_n
\]
是环同态。若 $\Psi(x)=0$，则 $x\in\bigcap_n\mathfrak p^n$。对离散赋值而言，非零
$x$ 的赋值是某个整数，不可能同时大于等于所有 $n$，所以 $x=0$，即 $\Psi$ 单。

再证满。给定兼容系 $(x_n)_n$，其中 $x_n\in\mathcal O/\mathfrak p^n$ 且
$x_{n+1}\equiv x_n\pmod{\mathfrak p^n}$。逐步选择代表元 $a_n\in\mathcal O$。兼容性
说明 $a_{n+1}-a_n\in\mathfrak p^n$，于是 $\{a_n\}$ 是 Cauchy 序列。完备性给出
$a_n\to a\in F$。因为 $\mathcal O$ 是闭球，它在完备度量中闭，故 $a\in\mathcal O$。
固定 $m$ 时，所有 $n\ge m$ 都满足 $a_n\equiv a_m\pmod{\mathfrak p^m}$，取极限得
$a\equiv a_m\pmod{\mathfrak p^m}$。于是 $\Psi(a)=(x_n)_n$。

拓扑同构性也来自同一组核：$\Psi^{-1}$ 把逆极限中投射到
$\mathcal O/\mathfrak p^n$ 为 $0$ 的基本邻域送回 $\mathfrak p^n$。

证明 (5)。单位元在逆极限中逐层判定：$x\in\mathcal O$ 是单位，当且仅当它模
$\mathfrak p$ 非零；这等价于每个像在 $\mathcal O/\mathfrak p^n$ 中都是单位。因此由
(4) 得
\[
\mathcal O^\times\simeq(\varprojlim\mathcal O/\mathfrak p^n)^\times
\simeq\varprojlim(\mathcal O/\mathfrak p^n)^\times
\simeq\varprojlim U^{(0)}/U^{(n)}.
\]
这些同构保持由核 $U^{(n)}$ 定义的邻域，故也是拓扑群同构。
\end{proof}

\begin{lemma}
自然包含 $\Z\subset\Z_{(p)}\subset\Z_p$ 对每个 $n\ge1$ 诱导环同构
\[
\Z/p^n\Z\simeq \Z_{(p)}/p^n\Z_{(p)}\simeq \Z_p/p^n\Z_p,
\]
其中 $\Z_{(p)}=\{a/b\in\Q:p\nmid b\}$。
\end{lemma}

\begin{proof}
先看 $\Z\to\Z_{(p)}$。若 $a\in\Z$ 在 $\Z_{(p)}/p^n\Z_{(p)}$ 中为 $0$，则
$a=p^n b/c$，其中 $p\nmid c$。于是 $ca=p^n b$。由于 $c$ 与 $p$ 互素，$p^n\mid a$，
所以核正是 $p^n\Z$。满性来自分母可逆：若 $b/c\in\Z_{(p)}$ 且 $p\nmid c$，则存在
$d\in\Z$ 使 $cd\equiv1\pmod {p^n}$，从而
$b/c\equiv bd\pmod {p^n\Z_{(p)}}$。

再看 $\Z\to\Z_p$。若整数 $a$ 在 $\Z_p/p^n\Z_p$ 中为 $0$，则 $v_p(a)\ge n$，也就是
$p^n\mid a$。给定 $\bar x\in\Z_p/p^n\Z_p$，由 $\Q$ 在 $\Q_p$ 中稠密，取
$r\in\Q$ 使 $|x-r|_p\le p^{-n}$。又 $x\in\Z_p$，可把 $r$ 改写为 $b/c$ 且 $p\nmid c$。
按上一段选择 $d$ 使 $cd\equiv1\pmod {p^n}$，则 $bd\in\Z$ 与 $r$ 模 $p^n$ 相同，
从而与 $x$ 模 $p^n\Z_p$ 相同。故 $\Z/p^n\Z\simeq\Z_p/p^n\Z_p$。
\end{proof}

\begin{proposition}
局部数域 $F$ 必为某个 $\Q_p$ 的有限扩张，且其赋值环 $\mathcal O_F$ 是 $F$ 的唯一极大紧子环。
\end{proposition}

\begin{proof}
特征 $0$ 给出嵌入 $\Q\subset F$。把 $F$ 的绝对值限制到 $\Q$，得到 $\Q$ 上非平凡
非 Archimede 绝对值；由 Ostrowski 定理，它等价于某个 $p$ 进绝对值。由于 $F$
完备，$\Q$ 在 $F$ 中的闭包同构于 $\Q_p$，所以 $F$ 是 $\Q_p$ 的扩域。

剩余域有限使 $\mathcal O_F/\mathfrak p^n$ 对所有 $n$ 都有限。上面的逆极限描述给出
\[
\mathcal O_F\simeq\varprojlim\mathcal O_F/\mathfrak p^n
\]
为有限离散集的逆极限，故 $\mathcal O_F$ 紧。选 $\mathcal O_F$ 的一组提升
$a_1,\dots,a_f$，其剩余类构成 $\kappa_F/\F_p$ 的基。再取素元 $\pi$，设
$p\mathcal O_F=\mathfrak p^e$。每个 $\mathcal O_F$ 元素都可逐次写成
$\sum_{0\le i<e,1\le j\le f} c_{ij}a_j\pi^i$ 的 $\Z_p$-线性极限，因此
$\mathcal O_F$ 是有限生成 $\Z_p$-模。这推出 $F$ 是有限维 $\Q_p$-向量空间。

若 $A\subset F$ 是紧子环而含有某个 $x$ 满足 $v(x)<0$，则 $x^n\in A$ 且
$v(x^n)=nv(x)\to-\infty$。这组元素没有收敛子列，因为不同 $n$ 的元素越来越远离
$\mathcal O_F$ 的有界部分，和紧性矛盾。因此任意紧子环都包含在
$\mathcal O_F$ 中。既然 $\mathcal O_F$ 自身紧，它就是唯一极大紧子环。
\end{proof}

\begin{theorem}
设 $F$ 为局部数域，剩余域 $\kappa_F$ 有 $q=p^m$ 个元素。若 $E/F$ 有限，则：
\begin{enumerate}[label=(\arabic*)]
\item $E$ 有唯一极大无分歧子扩张 $T/F$，且 $[T:F]=f(E/F)$；$T$ 由一个本原 $(q^{f(E/F)}-1)$ 次单位根生成。
\item $E/T$ 全分歧，$[E:T]=e(E/F)$。若 $\Pi$ 是 $E$ 的素元，则 $E=T(\Pi)$，且 $\Pi$ 在 $T$ 上的最小多项式是 Eisenstein 多项式。
\end{enumerate}
\end{theorem}

\begin{proof}
先解释无分歧部分。有限域 $\F_q$ 的 $f$ 次扩张唯一，记为 $\F_{q^f}$，其乘法群是
阶 $q^f-1$ 的循环群。取 $\bar\zeta\in\F_{q^f}^\times$ 为生成元，它满足
$\bar\zeta^{q^f-1}=1$，且导数 $(q^f-1)X^{q^f-2}$ 在剩余特征 $p$ 下不为零，因为
$p\nmid(q^f-1)$。Hensel 引理把 $\bar\zeta$ 唯一提升为一个 $(q^f-1)$ 次单位根
$\zeta$。令 $T=F(\zeta)$。它的剩余域包含 $\F_q(\bar\zeta)=\F_{q^f}$，而
$\zeta$ 满足次数至多 $q^f-1$ 的单位根方程；由无分歧扩张的基本等式可得到
$[T:F]=f$ 且 $e(T/F)=1$。所以 $T/F$ 是无分歧扩张。

在给定有限扩张 $E/F$ 内，取其剩余域扩张 $\kappa_E/\kappa_F$ 的可分部分。有限域
都是完美域，所以该剩余域扩张本身可分。按上一段，存在唯一无分歧扩张 $T/F$
使 $\kappa_T=\kappa_E$，并可嵌入 $E$。唯一性来自 Hensel 提升的唯一性：剩余域中
的同一个本原单位根只有一个提升。于是 $T$ 是 $E$ 中最大无分歧子扩张，且
$[T:F]=[\kappa_E:\kappa_F]=f(E/F)$。

接着看 $E/T$。由于 $\kappa_T=\kappa_E$，这个扩张的剩余次数为 $1$，所以全部次数
都由分歧指数承担，$E/T$ 全分歧，$[E:T]=e(E/F)$。若 $\Pi$ 是 $E$ 的素元，则
$w(\Pi)$ 生成 $w(E^\times)/w(T^\times)$，剩余域又没有扩大。用
$1,\Pi,\dots,\Pi^{e-1}$ 对值群代表，再用 $\mathcal O_T$ 的单位近似逐层修正，可得
$E=T(\Pi)$。设 $\Pi$ 在 $T$ 上的最小多项式为
$X^e+a_{e-1}X^{e-1}+\cdots+a_0$。所有共轭根都有正赋值，故 $a_i\in\mathfrak p_T$
对 $i<e$；常数项是根的乘积，赋值恰为 $1$，因而不在 $\mathfrak p_T^2$。这正是
Eisenstein 条件。
\end{proof}

\begin{proposition}
对正整数 $m$，当 $n$ 足够大时，幂映射
\[
U^{(n)}\longrightarrow U^{(n+v(m))},\qquad x\longmapsto x^m
\]
是群同构。
\end{proposition}

\begin{proof}
令 $\pi$ 为素元，写 $m=u\pi^{v(m)}$，其中 $u\in\mathcal O^\times$。若
$x=1+a\pi^n\in U^{(n)}$，二项式展开给出
\[
x^m=1+ma\pi^n+\binom m2a^2\pi^{2n}+\cdots.
\]
第一非平凡项落在 $\mathfrak p^{n+v(m)}$ 中；当 $n$ 大于所有高次项可能干扰的界限时，
后面的项赋值更高。因此 $x^m\in U^{(n+v(m))}$。

证明满性。给定 $1+a\pi^{n+v(m)}\in U^{(n+v(m))}$，要求
$(1+y\pi^n)^m=1+a\pi^{n+v(m)}$。展开后等价于
\[
-a+uy+\pi^{n-v(m)}f(y)=0
\]
其中 $f(Y)\in\mathcal O[Y]$。当 $n>v(m)$ 时，这个多项式模 $\mathfrak p$ 化为
$-\,\bar a+\bar uY$，导数为单位 $\bar u$。Hensel 引理给出 $y\in\mathcal O$ 的解。

证明单性。核由 $U^{(n)}$ 中的 $m$ 次单位根组成。局部域中的单位根有限；当 $n$
增大时，$\bigcap_nU^{(n)}=\{1\}$，所以对足够大的 $n$，核只剩 $1$。满性和单性合在一起
给出同构。
\end{proof}

\subsection{单位群}

\begin{lemma}
若正整数 $n$ 的 $p$ 进展开为 $n=a_0+a_1p+\cdots+a_sp^s$，$0\le a_i<p$，令
$s_p(n)=a_0+\cdots+a_s$，则
\[
v_p(n!)=\frac{n-s_p(n)}{p-1}.
\]
\end{lemma}

\begin{proof}
$n!$ 中含有的因子 $p$ 的个数为
\[
v_p(n!)=\sum_{j\ge1}\left\lfloor\frac n{p^j}\right\rfloor.
\]
当 $j>s$ 时该项为 $0$。由 $p$ 进展开，
\[
\left\lfloor\frac n{p^j}\right\rfloor
=a_j+a_{j+1}p+\cdots+a_sp^{s-j}.
\]
求和并交换求和次序得到
\[
\sum_{j=1}^s\left\lfloor\frac n{p^j}\right\rfloor
=\sum_{i=1}^s a_i(1+p+\cdots+p^{i-1})
=\sum_{i=1}^s a_i\frac{p^i-1}{p-1}
=\frac{n-s_p(n)}{p-1}.
\]
\end{proof}

\begin{proposition}
设 $F$ 为 $p$ 局部数域，绝对分歧指数 $e$ 由 $p\mathcal O_F=\mathfrak p^e$ 决定。
\begin{enumerate}[label=(\arabic*)]
\item 若 $v_p(x)>1/(p-1)$，则
\[
\exp x=1+x+\frac{x^2}{2!}+\cdots
\]
收敛。
\item 若 $x\in\mathfrak p$，则
\[
\log(1+x)=x-\frac{x^2}{2}+\frac{x^3}{3}-\cdots
\]
收敛。
\item 当 $m>e/(p-1)$ 时，
\[
\exp:\mathfrak p^m\longrightarrow U^{(m)},\qquad
\log:U^{(m)}\longrightarrow\mathfrak p^m
\]
互为同胚群同构。
\end{enumerate}
\end{proposition}

\begin{proof}
把 $\Q_p$ 上的 $v_p$ 延拓到 $F$，并仍记为 $v_p$。若 $v$ 是归一化到
$v(\pi)=1$ 的赋值，则 $v(p)=e$，因而 $v_p(x)=v(x)/e$。指数级数的第 $n$ 项满足
\[
v_p(x^n/n!)=n v_p(x)-v_p(n!)
=n\left(v_p(x)-\frac1{p-1}\right)+\frac{s_p(n)}{p-1}.
\]
若 $v_p(x)>1/(p-1)$，右边趋于 $+\infty$，所以第 $n$ 项趋于 $0$。在非
Archimede 完备域中，级数收敛等价于通项趋于 $0$，指数级数收敛。

对数级数中
\[
v_p(x^n/n)=n v_p(x)-v_p(n).
\]
若 $x\in\mathfrak p$，则 $v_p(x)>0$。而 $v_p(n)\le\log_p n$，所以
$n v_p(x)-v_p(n)\to+\infty$，对数级数收敛。

最后证明互逆和同构。形式幂级数恒等式
\[
\exp(\log(1+z))=1+z,\qquad \log(\exp x)=x
\]
在上述收敛区域内可逐项代入。若 $m>e/(p-1)$ 且 $x\in\mathfrak p^m$，则
$v_p(x)>1/(p-1)$。指数展开中除首项 $x$ 外，其余项赋值更高，故
$v(\exp x-1)=v(x)$，于是 $\exp$ 把 $\mathfrak p^m$ 映入 $U^{(m)}$。同样，
若 $1+z\in U^{(m)}$，则 $\log(1+z)$ 的首项为 $z$，其余项赋值更高，所以
$v(\log(1+z))=v(z)$。两者互逆，且形式恒等式
$\log((1+x)(1+y))=\log(1+x)+\log(1+y)$ 给出群同态性。连续性来自级数在这些邻域中
一致收敛。
\end{proof}

\begin{theorem}
设局部数域 $F$ 的剩余域有 $q=p^f$ 个元素，$d=[F:\Q_p]$。
\begin{enumerate}[label=(\arabic*)]
\item
\[
F^\times\simeq \langle\pi\rangle\times\mathcal O^\times.
\]
\item
\[
\mathcal O^\times\simeq \mu_{q-1}\times U^{(1)}.
\]
\item $U^{(1)}$ 是有限 $p$ 幂单位根群与 $\Z_p^d$ 的直积。
\item 若 $\mu(F)$ 为 $F$ 中全部单位根，则
\[
F^\times\simeq
\Z\oplus \Z/(q-1)\Z\oplus \Z/p^t\Z\oplus \Z_p^d
\]
对某个 $t\ge0$ 成立。
\item 对正整数 $n$，
\[
[F^\times:F^{\times n}]
=|\mu_n(F)|\, n\, p^{d\,v_p(n)}.
\]
\end{enumerate}
\end{theorem}

\begin{proof}
(1) 对 $x\in F^\times$，令 $r=v(x)$。则 $u=x\pi^{-r}$ 满足 $v(u)=0$，所以
$u\in\mathcal O^\times$，且 $x=\pi^r u$。若 $\pi^r u=1$，取赋值得 $r=0$，再得
$u=1$，所以这是直积。

(2) 剩余映射 $\mathcal O^\times\to\kappa^\times$ 满射，核为 $U^{(1)}$。有限域
$\kappa^\times$ 是阶 $q-1$ 的循环群，它的元素正是 $X^{q-1}-1$ 的根。由于
$p\nmid(q-1)$，导数 $(q-1)X^{q-2}$ 在每个非零剩余根处不为零。Hensel 引理把每个
剩余根唯一提升为 $\mathcal O^\times$ 中的 $(q-1)$ 次单位根，得到正合列的分裂：
\[
1\to U^{(1)}\to\mathcal O^\times\to\kappa^\times\to1
\]
且 $\mathcal O^\times\simeq\mu_{q-1}\times U^{(1)}$。

(3) 当 $m>e/(p-1)$ 时，$\log:U^{(m)}\to\mathfrak p^m$ 是同胚群同构。作为
$\Z_p$-模，$\mathcal O$ 是秩 $d=[F:\Q_p]$ 的自由模，而 $\mathfrak p^m$ 是它的有限
指数子模，所以 $\mathfrak p^m\simeq\Z_p^d$。因此 $U^{(m)}\simeq\Z_p^d$。又
$U^{(m)}$ 在 $U^{(1)}$ 中有限指数，因为每个商 $U^{(i)}/U^{(i+1)}$ 都同构于有限加群
$\mathfrak p^i/\mathfrak p^{i+1}$。于是 $U^{(1)}$ 是有限生成 $\Z_p$-模。有限生成
$\Z_p$-模分解为有限挠子部分和自由部分；这里的挠子元素正是 $p$ 幂阶单位根，
自由部分秩为 $d$。

(4) 把 (1)--(3) 合并。阶与 $p$ 互素的单位根来自 $\mu_{q-1}$，$p$ 幂阶单位根来自
$U^{(1)}$ 的挠子部分，所以
$\mu(F)=\mu_{q-1}\times\mu_{p^t}$，并得到所列分解。

(5) 在直积分解中逐项计算 $n$ 次幂映射的余核大小。$\Z$ 上乘以 $n$ 的余核大小是
$n$。若 $n=p^rs$ 且 $(s,p)=1$，则乘以 $s$ 在 $\Z_p$ 上可逆，乘以 $p^r$ 的余核大小为
$p^r$，所以 $\Z_p^d$ 贡献 $p^{d v_p(n)}$。有限单位根部分由正合列
\[
1\to\mu_n(F)\to\mu(F)\xrightarrow{x\mapsto x^n}\mu(F)\to\mu(F)/\mu(F)^n\to1
\]
给出贡献 $|\mu_n(F)|$。三项相乘得到公式。
\end{proof}

\begin{lemma}
在多项式环 $\Z[T]$ 中令 $I=(p,T)$。对 $n\ge0$，
\[
(1+T)^{p^n}-1\in T I^n.
\]
\end{lemma}

\begin{proof}
用归纳法。$n=0$ 时，$(1+T)-1=T\in TI^0$。设
$(1+T)^{p^n}=1+Tu$ 且 $u\in I^n$。则
\[
(1+T)^{p^{n+1}}=(1+Tu)^p
=1+pTu+\binom p2T^2u^2+\cdots+T^pu^p.
\]
除第一项外，中间项都含有因子 $pT^2u^2$ 或更高次；它们属于 $TI^{n+1}$，因为
$p u\in I^{n+1}$ 且 $T u\in I^{n+1}$。末项 $T^pu^p=T\cdot T^{p-1}u^p$ 也属于
$TI^{n+1}$。第一项 $pTu=T\cdot pu$ 属于 $TI^{n+1}$。因此
$(1+T)^{p^{n+1}}-1\in TI^{n+1}$。
\end{proof}

\begin{proposition}
设 $|\cdot|_p$ 是 $\Q_p$ 的绝对值，并延拓到 $F$。
\begin{enumerate}[label=(\arabic*)]
\item 若 $t\in F$、$|t|_p\le1$，则
\[
|(1+t)^{p^n}-1|_p
\le |t|_p\max\{|t|_p,p^{-1}\}^n.
\]
特别地，若 $|t|_p<1$，则 $(1+t)^{p^n}\to1$。
\item 投射 $\omega:\mathcal O^\times=\mu_{q-1}\times U^{(1)}\to\mu_{q-1}$ 可由
\[
\omega(x)=\lim_{n\to\infty}x^{p^{n!}}
\]
给出。
\item 对 $a\in F$，$|a-1|_p<1$ 当且仅当 $a^{p^n}\to1$。
\item 若 $\zeta$ 是 $p^t$ 阶本原单位根，则
\[
|\zeta-1|_p=p^{-1/\varphi(p^t)}.
\]
\end{enumerate}
\end{proposition}

\begin{proof}
(1) 由上一引理，$(1+T)^{p^n}-1=T f(T)$，且 $f(T)\in I^n$。$I^n$ 的每一项都是
$p^iT^{n-i}$ 的整系数组合。代入 $T=t$ 后，每一项的绝对值不超过
$p^{-i}|t|_p^{n-i}$。再乘前面的 $|t|_p$，得到所需估计。若 $|t|_p<1$，则
\(\max\{|t|_p,p^{-1}\}<1\)，右边趋于 $0$。

(2) 写 $x=\zeta u$，其中 $\zeta\in\mu_{q-1}$、$u\in U^{(1)}$。因为 $p$ 与 $q-1$ 互素，
当 $n$ 足够大时，Euler 定理给出 $p^{n!}\equiv1\pmod {q-1}$，故
$\zeta^{p^{n!}}=\zeta$。另一方面 $u=1+t$ 且 $|t|_p<1$，由 (1) 得
$u^{p^N}\to1$，取 $N=n!$ 仍收敛到 $1$。因此
$x^{p^{n!}}=\zeta^{p^{n!}}u^{p^{n!}}\to\zeta$。这正是投射 $\omega(x)$。

(3) 若 $|a-1|_p<1$，由 (1) 得 $a^{p^n}\to1$。反过来若 $a^{p^n}\to1$，则 $|a|_p=1$。
若 $|a-1|_p>1$，则 $|a|_p=|a-1|_p>1$，与 $|a^{p^n}|\to1$ 矛盾。于是只需排除
$|a-1|_p=1$。在这种情况下二项式展开给出：对 $1\le j\le p-1$，
$|\binom pj a^{p-j}(-1)^j|_p<1$，而端点项的差保持绝对值 $1$；当 $p=2$ 时再用
$|a^2-1|_p=|(a^2+1)-2|_p=1$。故 $|a^p-1|_p=1$，迭代得到
$|a^{p^n}-1|_p=1$，与收敛到 $1$ 矛盾。

(4) 先令 $t=1$。若 $\zeta^p=1$ 且 $\zeta\ne1$，写 $\zeta=1+\xi$。由 (3) 知
 $|\xi|_p<1$。方程
\[
0=\frac{(1+\xi)^p-1}{\xi}
=p+\binom p2\xi+\cdots+\xi^{p-1}
\]
可写成 $p(1+\xi y)+\xi^{p-1}=0$，其中 $|y|_p\le1$。因为 $|\xi|_p<1$，
$|1+\xi y|_p=1$，所以 $|\xi|_p^{p-1}=|p|_p=p^{-1}$，即
$|\zeta-1|_p=p^{-1/(p-1)}$。

若 $\zeta$ 是 $p^t$ 阶本原单位根，则 $\zeta^{p^{t-1}}$ 是 $p$ 阶本原单位根。写
$\zeta=1+\theta$。展开 $\zeta^{p^{t-1}}-1$ 时，主项为 $\theta^{p^{t-1}}$，其余项赋值
更高，因此
\[
|\theta|_p^{p^{t-1}}=p^{-1/(p-1)}.
\]
于是 $|\zeta-1|_p=p^{-1/(p^{t-1}(p-1))}=p^{-1/\varphi(p^t)}$。
\end{proof}

\begin{proposition}
存在唯一连续同态
\[
\log:F^\times\to F
\]
使 $\log\pi=0$，且对 $1+x\in U^{(1)}$，
\[
\log(1+x)=x-\frac{x^2}{2}+\frac{x^3}{3}-\cdots.
\]
\end{proposition}

\begin{proof}
由分解
\[
F^\times=\langle\pi\rangle\times\mu_{q-1}\times U^{(1)}
\]
定义 $\log(\pi)=0$，对 $\mu_{q-1}$ 上的元素定义 $\log(\zeta)=0$，在 $U^{(1)}$ 上用
收敛的对数级数。由于 $U^{(1)}$ 中可能含有 $p$ 幂阶单位根，需要说明该定义仍然为
群同态。对 $U^{(m)}$，$m>e/(p-1)$，前面的指数与对数已经给出群同构；而
$U^{(m)}$ 是 $U^{(1)}$ 的开子群，任意 $u\in U^{(1)}$ 的足够高 $p$ 次幂落入
$U^{(m)}$。连续性把小邻域上的同态性质推广到整个 $U^{(1)}$。

唯一性如下。连续同态若在 $\pi$ 上取 $0$，并且在 $U^{(1)}$ 上满足给定级数，则它在
$\langle\pi\rangle$ 和 $U^{(1)}$ 上都已被决定。阶与 $p$ 互素的单位根 $\zeta$ 满足
$\zeta^{q-1}=1$，把同态作用后得 $(q-1)\log\zeta=0$。域 $F$ 的加法群没有非零有限阶
元素，所以 $\log\zeta=0$。因此整个 $F^\times$ 上唯一。
\end{proof}

\begin{proposition}
若 $1+x\in U^{(1)}$，则存在唯一连续映射
\[
\Z_p\longrightarrow U^{(1)},\qquad z\longmapsto (1+x)^z
\]
延拓整数幂，并且
\[
(1+x)^z=\sum_{n=0}^\infty \binom{z}{n}x^n.
\]
\end{proposition}

\begin{proof}
对非负整数 $z$，右边就是二项式定理。若 $z\in\Z_p$，二项式系数
\[
\binom{z}{n}=\frac{z(z-1)\cdots(z-n+1)}{n!}
\]
属于 $\Z_p$；这是整数值多项式在 $\Z_p$ 上连续延拓的标准性质。因为
$x\in\mathfrak p$，有 $|x^n|_p\to0$，而 $|\binom{z}{n}|_p\le1$，所以级数收敛到
$U^{(1)}$ 中的元素。

连续性由级数一致收敛得到。若 $z_k\in\Z_{\ge0}$ 收敛到 $z\in\Z_p$，则
$(1+x)^{z_k}$ 按上式收敛到定义的 $(1+x)^z$，所以它确实是整数幂映射的连续延拓。
乘法律来自多项式恒等式
\[
(1+X)^{a+b}=(1+X)^a(1+X)^b
\]
先对非负整数 $a,b$ 成立，再由 $\Z_{\ge0}$ 在 $\Z_p$ 中稠密和两边连续，推广到
$a,b\in\Z_p$。
\end{proof}

\subsection{局部域的无分歧扩张}

设 $F$ 为局部数域，剩余域有 $q=p^m$ 个元素。单位群分解中的
$\mu_{q-1}$ 是剩余域乘法群的 Teichmuller 提升：投射
$\mathcal O\to\kappa_F$ 限制到 $\mu_{q-1}\cup\{0\}$ 是双射。这个事实使无分歧扩张
可以用阶与 $p$ 互素的单位根完全描述。

\begin{proposition}
对每个 $n\ge1$，存在唯一 $n$ 次无分歧扩张 $F_n/F$。它由一个本原 $(q^n-1)$ 次单位根生成。极大无分歧扩张为
\[
F^{\mathrm{nr}}=\bigcup_{n\ge1}F_n.
\]
其 Galois 群由 Frobenius 拓扑生成，Frobenius 在剩余域上作用为
\[
x\longmapsto x^q.
\]
\end{proposition}

\begin{proof}
有限域 $\F_q$ 有唯一 $n$ 次扩张 $\F_{q^n}$，其乘法群循环，阶为 $q^n-1$。取
$\bar\zeta$ 为本原 $(q^n-1)$ 次单位根。由于 $p\nmid(q^n-1)$，Hensel 引理把
$\bar\zeta$ 唯一提升为某个 $\zeta$，令 $F_n=F(\zeta)$。剩余域包含
$\F_q(\bar\zeta)=\F_{q^n}$，而无分歧构造给出 $[F_n:F]=n$。因此 $F_n/F$ 是 $n$ 次
无分歧扩张。

若 $E/F$ 是任意 $n$ 次无分歧扩张，则 $\kappa_E/\kappa_F$ 是 $n$ 次有限域扩张，
只能是 $\F_{q^n}/\F_q$。其中本原 $(q^n-1)$ 次单位根的 Teichmuller 提升属于
$E$，并生成同一个无分歧扩张，所以 $E=F_n$。这给出唯一性。

把所有 $F_n$ 合起来得到 $F^{\mathrm{nr}}$。它的剩余域是
\[
\bigcup_{n\ge1}\F_{q^n},
\]
即 $\kappa_F$ 的可分闭包。每个有限层的 Galois 群同构于
$\operatorname{Gal}(\F_{q^n}/\F_q)$，由 $x\mapsto x^q$ 生成。有限层上的这些生成元
彼此兼容，因而在逆极限上给出一个拓扑生成元，称为 Frobenius 自同构。它在
$\mu_{(p)}=\bigcup_{(r,p)=1}\mu_r$ 上由
\[
\zeta\longmapsto \zeta^q
\]
给出；模 $\mathfrak p$ 后正是有限域上的 $q$ 次幂 Frobenius。
\end{proof}

\begin{remark}
后面局部类域论把 $F^\times$ 的素元部分送到无分歧 Galois 群的 Frobenius 方向，
把单位部分送到全分歧方向。第三章此处的结构分解
$F^\times=\langle\pi\rangle\times\mathcal O^\times$，正是后续互反映射的拓扑蓝图。
\end{remark}

\section{形式群}

本节进入 Lubin--Tate 理论的入口。前一节说明，无分歧扩张由剩余域和
Frobenius 控制；剩下的困难是全分歧交换扩张。Lubin--Tate 的想法是：不直接
猜全分歧扩张的方程，而是在局部域的极大理想上制造一个“加法规律”，再研究这个
形式加法下被 $\pi^n$ 杀掉的点。这样得到的除点域正好给出局部类域论中单位群
对应的全分歧部分。

\subsection{形式模}

\begin{definition}
设 $\mathcal O$ 为环。$\mathcal O$ 上的交换一维形式群是幂级数
\[
F(X,Y)\in\mathcal O[[X,Y]]
\]
满足：
\begin{enumerate}[label=(\arabic*)]
\item $F(X,Y)\equiv X+Y\pmod{\deg2}$；
\item $F(X,F(Y,Z))=F(F(X,Y),Z)$；
\item $F(X,Y)=F(Y,X)$。
\end{enumerate}
形式群同态 $f:F\to G$ 是无常数项幂级数 $f(X)$，满足
\[
f(F(X,Y))=G(f(X),f(Y)).
\]
\end{definition}

\begin{definition}
若有环同态
\[
\mathcal O\to\operatorname{End}_{\mathcal O}(F),\qquad a\mapsto [a]_F
\]
且
\[
[a]_F(X)\equiv aX\pmod{\deg2},
\]
则称 $F$ 为形式 $\mathcal O$-模。
\end{definition}

\begin{remark}
形式群中的 $F(X,Y)$ 不是普通函数，而是“形式加法”。当 $x,y$ 落在局部域的极大
理想中时，幂级数收敛，$F(x,y)$ 才变成真正的加法运算。最基本的两个例子是
\[
F_a(X,Y)=X+Y,\qquad F_m(X,Y)=X+Y+XY.
\]
第二个例子来自乘法群，因为 $(1+X)(1+Y)-1=X+Y+XY$。Lubin--Tate 理论把这个
乘法形式群推广到任意局部域和任意素元 $\pi$。
\end{remark}

\begin{proof}[乘法形式群的检验]
对 $F_m(X,Y)=X+Y+XY$，线性部分为 $X+Y$。结合律来自
\[
1+F_m(X,F_m(Y,Z))=(1+X)(1+Y)(1+Z)
=1+F_m(F_m(X,Y),Z),
\]
交换律来自乘法交换律。若 $n\in\Z$，则
\[
[n]_{F_m}(X)=(1+X)^n-1
\]
给出自同态，因为
$1+[n]_{F_m}(F_m(X,Y))=((1+X)(1+Y))^n$。这个例子说明后面的
$f(X)=(1+X)^p-1$ 为什么产生分圆扩张。
\end{proof}

\subsection{Lubin--Tate 模}

设 $K$ 为局部域，$\pi$ 为素元，剩余域大小为 $q$。令 $\mathcal F_\pi$ 为所有满足
\[
f(X)\equiv \pi X\pmod{\deg2},\qquad
f(X)\equiv X^q\pmod{\pi}
\]
的 $f(X)\in\mathcal O_K[[X]]$ 的集合，称其中元素为 $\pi$ 的 Frobenius 级数。
第一个同余规定了切空间上的乘 $\pi$；第二个同余规定了剩余域上的 Frobenius。
也就是说，$f$ 同时记住“素元方向”和“剩余域 Frobenius 方向”。

\begin{lemma}
取 $f,g\in\mathcal F_\pi$，以及线性形式
\[
L(X_1,\dots,X_n)=a_1X_1+\cdots+a_nX_n.
\]
存在唯一 $F(X_1,\dots,X_n)\in\mathcal O_K[[X_1,\dots,X_n]]$，使
\[
F\equiv L\pmod{\deg2},\qquad
f(F(X_1,\dots,X_n))=F(g(X_1),\dots,g(X_n)).
\]
\end{lemma}

\begin{proof}
记 $X=(X_1,\dots,X_n)$ 与 $g(X)=(g(X_1),\dots,g(X_n))$。我们按总次数递推构造
多项式 $F_r(X)$，使
\[
F_r(X)\equiv L(X)\pmod{\deg2},\qquad
f(F_r(X))\equiv F_r(g(X))\pmod{\deg(r+1)}.
\]
$r=1$ 时取 $F_1=L$。设 $F_r$ 已经构造。要把同余推进到次数 $r+2$，令
\[
F_{r+1}=F_r+H_{r+1},
\]
其中 $H_{r+1}$ 是 $(r+1)$ 次齐次多项式。由于 $f(T)=\pi T+\text{高次项}$，在模
次数 $r+2$ 下有
\[
f(F_r+H_{r+1})\equiv f(F_r)+\pi H_{r+1}.
\]
又因为 $g(X_i)\equiv \pi X_i\pmod{\deg2}$，齐次多项式满足
\[
H_{r+1}(g(X))\equiv \pi^{r+1}H_{r+1}(X)\pmod{\deg(r+2)}.
\]
因此误差的
$(r+1)$ 次部分从 $C_{r+1}$ 变为
\[
C_{r+1}+(\pi-\pi^{r+1})H_{r+1}.
\]
需要取
\[
H_{r+1}=-\frac{C_{r+1}}{\pi-\pi^{r+1}}.
\]
这里 $\pi-\pi^{r+1}=\pi(1-\pi^r)$，且 $1-\pi^r$ 是单位。还要知道
$C_{r+1}$ 的系数都被 $\pi$ 整除。模 $\pi$ 后，$f(T)\equiv T^q$ 与
$g(X_i)\equiv X_i^q$，于是
\[
f(F_r(X))-F_r(g(X))\equiv F_r(X)^q-F_r(X_1^q,\dots,X_n^q)\equiv0\pmod\pi,
\]
最后一个同余来自剩余特征 $p$ 下的 Frobenius 公式。故 $H_{r+1}$ 的系数仍在
$\mathcal O_K$ 中。递推得到兼容的 $F_r$，取形式极限
$F=\lim_rF_r$。唯一性由同一个递推式给出：每一阶修正都被当前误差决定。
\end{proof}

\begin{proposition}
对任意 $f\in\mathcal F_\pi$：
\begin{enumerate}[label=(\arabic*)]
\item 存在唯一形式群 $F_f$，使 $f\in\operatorname{End}(F_f)$。
\item 存在唯一环同态 $\mathcal O_K\to\operatorname{End}(F_f)$，记为 $a\mapsto[a]_f$，满足 $[\pi]_f=f$ 且 $[a]_f(X)\equiv aX\pmod{\deg2}$。
\item 若 $f,g\in\mathcal F_\pi$，则对每个 $a\in\mathcal O_K$ 存在唯一同态
$[a]_{f,g}:F_g\to F_f$，满足
\[
[a+b]_{f,g}=[a]_{f,g}+_{F_f}[b]_{f,g},\quad
[a]_{f,g}\circ[b]_{g,h}=[ab]_{f,h},\quad
f\circ[a]_{f,g}=[a]_{f,g}\circ g.
\]
\end{enumerate}
\end{proposition}

\begin{proof}
把上一引理用于 $g=f$ 和 $L=X+Y$，得到唯一的幂级数
$F_f(X,Y)$，满足线性部分为 $X+Y$ 且
\[
f(F_f(X,Y))=F_f(f(X),f(Y)).
\]
结合律的检验使用唯一性：幂级数 $F_f(X,F_f(Y,Z))$ 与
$F_f(F_f(X,Y),Z)$ 都有线性部分 $X+Y+Z$，并且都满足与 $f$ 交换的同一函数方程；
由上一引理的唯一性二者相等。交换律同样由 $F_f(X,Y)$ 与 $F_f(Y,X)$ 的同一线性项
和同一方程得到。故 $F_f$ 是形式群，且 $f$ 是它的自同态。

对 $a\in\mathcal O_K$，把上一引理用于线性形式 $L=aX$，得到唯一
$[a]_f(X)$，满足 $[a]_f(X)\equiv aX\pmod{\deg2}$ 且
$f\circ[a]_f=[a]_f\circ f$。令 $a=\pi$ 时，$f$ 本身满足这些条件，所以
$[\pi]_f=f$。加法和乘法相容性仍由唯一性推出：例如
$[a+b]_f$ 与 $[a]_f+_{F_f}[b]_f$ 有同一线性项 $(a+b)X$，并且都与 $f$ 交换；
组合 $[a]_f\circ[b]_f$ 与 $[ab]_f$ 也有同一线性项 $abX$，并且都与 $f$ 交换。
唯一性迫使它们相等。这给出环同态。

若 $f,g$ 不同，则用上一引理中的两个 Frobenius 级数 $f,g$ 与线性形式 $aX$，得到
$[a]_{f,g}$。同态性质和三条相容公式也由“线性项相同且满足同一交换方程”的唯一性
逐项推出。
\end{proof}

\subsection{\texorpdfstring{$\pi^n$}{pi n} 除点}

\begin{definition}
设 $K^{\mathrm{alg}}$ 为代数闭包。定义
\[
F_f(n)=\{x\in\mathfrak p_{K^{\mathrm{alg}}}: [\pi^n]_f(x)=0\},
\]
称为 Lubin--Tate 模的 $\pi^n$ 除点。
\end{definition}

\begin{proposition}
\begin{enumerate}[label=(\arabic*)]
\item $F_f(n)$ 是 $\mathcal O_K$-模。
\item 若 $a\in\mathcal O_K$ 且 $f,g\in\mathcal F_\pi$，则
$x\mapsto[a]_{g,f}(x)$ 给出 $\mathcal O_K$-模同态 $F_f(n)\to F_g(n)$；若 $a$ 是单位，
则它是同构。
\item $F_f(n)\simeq\mathcal O_K/\pi^n\mathcal O_K$。
\item 由单位 $u\in U_K$ 给出的自同构 $x\mapsto[u]_f(x)$ 诱导
\[
\operatorname{Aut}_{\mathcal O_K}(F_f(n))\simeq U_K/U_K^{(n)}.
\]
\end{enumerate}
\end{proposition}

\begin{proof}
(1) 若 $[\pi^n]_f(x)=0$，则
\[
[\pi^n]_f([a]_f(x))=[a]_f([\pi^n]_f(x))=0,
\]
所以 $[a]_f(x)$ 仍是除点。形式加法 $F_f(x,y)$ 也保持除点，因为
\[
[\pi^n]_f(F_f(x,y))=F_f([\pi^n]_f(x),[\pi^n]_f(y))=0.
\]
因此 $F_f(n)$ 是
$\mathcal O_K$-模。

(2) 相容式给出
\[
[\pi^n]_g([a]_{g,f}(x))=[a]_{g,f}([\pi^n]_f(x))=0,
\]
所以确实映入 $F_g(n)$。若 $a$ 是单位，取 $a^{-1}\in\mathcal O_K^\times$，则
\[
[a^{-1}]_{f,g}\circ[a]_{g,f}=[1]_{f,f},\qquad
[a]_{g,f}\circ[a^{-1}]_{f,g}=[1]_{g,g}.
\]
两边复合都是恒等同态，故 $[a]_{g,f}$ 为同构。

(3) 先看 $n=1$。因为 $f(X)\equiv X^q\pmod\pi$ 且线性项为 $\pi X$，多项式
$f(X)/X$ 是 Eisenstein 型多项式，非零根连同 $0$ 一共给出 $q$ 个 $\pi$ 除点，正好
对应 $\mathcal O_K/\pi$ 的元素。递推地，$[\pi^n]_f=f\circ[\pi^{n-1}]_f$。每个
$\pi^{n-1}$ 除点的原像由一个平移后的 Eisenstein 方程给出，新增一层大小为 $q$，
故 $|F_f(n)|=q^n$。

取一个不属于 $F_f(n-1)$ 的元素 $\lambda_n\in F_f(n)$，称为本原 $\pi^n$ 除点。映射
\[
\mathcal O_K/\pi^n\to F_f(n),\qquad
a\mapsto [a]_f(\lambda_n)
\]
是 $\mathcal O_K$-模同态。若 $[a]_f(\lambda_n)=0$，则 $a$ 杀掉本原
$\pi^n$ 除点，这迫使 $a\in\pi^n\mathcal O_K$；若 $v(a)=r<n$，则
$[a]_f(\lambda_n)$ 落到本原 $\pi^{n-r}$ 除点而非 $0$。所以该映射单。两边都有
$q^n$ 个元素，故为同构。

(4) 在同构 $F_f(n)\simeq\mathcal O_K/\pi^n$ 下，$[u]_f$ 对应于乘以单位 $u$。
全部 $\mathcal O_K$-模自同构都由乘以单位给出，且 $u$ 与 $u'$ 给出同一自同构当且
仅当 $u\equiv u'\pmod{\pi^n}$。因此
\[
\operatorname{Aut}_{\mathcal O_K}(F_f(n))\simeq
(\mathcal O_K/\pi^n)^\times\simeq U_K/U_K^{(n)}.
\]
\end{proof}

\begin{proposition}
若 $f,g\in\mathcal F_\pi$，则 $K(F_f(n))=K(F_g(n))$。
\end{proposition}

\begin{proof}
由上一命题，$[1]_{g,f}$ 给出 $F_f(n)\to F_g(n)$ 的同构，且其幂级数系数在
$\mathcal O_K$ 中。因此 $F_g(n)\subset K(F_f(n))$。交换 $f,g$ 得反向包含，故两域
相等。
\end{proof}

\begin{theorem}
令 $E_{\pi,n}=K(F_f(n))$。此域与 $f\in\mathcal F_\pi$ 的选择无关。记
$G_{\pi,n}=\Gal(E_{\pi,n}/K)$，则
\[
G_{\pi,n}\simeq U_K/U_K^{(n)}.
\]
此外 $E_{\pi,n}/K$ 是全分歧交换扩张，
\[
[E_{\pi,n}:K]=q^{n-1}(q-1).
\]
若 $f^n$ 表示 $f$ 的 $n$ 次复合，令
\[
\Phi_n(X)=(f^{n-1}(X))^{q-1}+\pi,
\]
则存在 $\lambda\in E_{\pi,n}$ 满足 $\Phi_n(\lambda)=0$、
$E_{\pi,n}=K(\lambda)$，并且
\[
N_{E_{\pi,n}/K}(-\lambda)=\pi.
\]
\end{theorem}

\begin{proof}
域与 $f$ 无关已经由上一命题证明。对 $u\in U_K$，自同态 $[u]_f$ 作用在
$F_f(n)$ 上，因此诱导 $K$-自同构 $\sigma_u$：
\[
\sigma_u(x)=[u]_f(x)\qquad(x\in F_f(n)).
\]
若 $u\in U_K^{(n)}$，则它在 $F_f(n)\simeq\mathcal O_K/\pi^n$ 上作用为恒等；若它
作用为恒等，则 $u\equiv1\pmod{\pi^n}$。所以得到单同态
\[
U_K/U_K^{(n)}\hookrightarrow G_{\pi,n}.
\]

取本原除点 $\lambda$。任意 $\sigma\in G_{\pi,n}$ 保持方程 $[\pi^n]_f(X)=0$，所以
$\sigma(\lambda)$ 仍是 $F_f(n)$ 的本原元素。按 $F_f(n)\simeq\mathcal O_K/\pi^n$，
本原元素正是 $[u]_f(\lambda)$，其中 $u$ 模 $\pi^n$ 为单位。于是 $\sigma=\sigma_u$。
故 $G_{\pi,n}\simeq U_K/U_K^{(n)}$，扩张交换，且阶数为
\[
|(\mathcal O_K/\pi^n)^\times|=q^{n-1}(q-1).
\]

再看 $\Phi_n$。本原 $\pi^n$ 除点 $\lambda$ 的像 $f^{n-1}(\lambda)$ 是非零
$\pi$ 除点。非零 $\pi$ 除点满足 $f(Y)=0$ 且 $Y\ne0$；由
$f(Y)/Y\equiv Y^{q-1}+\pi$ 的 Eisenstein 形状，可选取本原除点使
\[
\Phi_n(\lambda)=(f^{n-1}(\lambda))^{q-1}+\pi=0.
\]
这个方程是 Eisenstein 型方程，故 $K(\lambda)/K$ 全分歧，并且 $\lambda$ 是素元。
由于 $\lambda$ 的所有共轭为 $[u]_f(\lambda)$，$u$ 遍历
$(\mathcal O_K/\pi^n)^\times$，最小多项式的常数项等于这些共轭的乘积。Eisenstein
方程给出这个乘积为 $(-1)^{[E_{\pi,n}:K]}\pi$。因此
$N_{E_{\pi,n}/K}(-\lambda)=\pi$。全分歧性也可从 $\lambda$ 是素元且剩余域未扩大
得到。
\end{proof}

\begin{example}
当 $K=\Q_p$，取
\[
f(X)=(1+X)^p-1
\]
时，形式群是乘法形式群。因为 $[p^n]_{F_m}(X)=(1+X)^{p^n}-1$，所以
$\pi^n$ 除点为
\[
\{\zeta-1:\zeta^{p^n}=1\}.
\]
因此
\[
E_{\pi,n}=\Q_p(\zeta_{p^n}).
\]
这个例子说明 Lubin--Tate 除点域把分圆扩张推广到了任意局部域。
\end{example}

\subsection{除点的变换}

若 $\pi'=u\pi$ 是另一个素元，$E_{\pi,n}$ 与 $E_{\pi',n}$ 未必在 $K$ 上相等。源文
接下来证明：与极大无分歧扩张合成后，它们相等。这一步的意义是把“选哪个素元”
造成的差异放入无分歧方向中消去。

令 $K^{\mathrm{nr}}$ 为极大无分歧扩张，$\widehat K$ 为其完备化，$\varphi$ 为
Frobenius 在 $\widehat K$ 上的连续延拓。它在剩余域上作用为 $x\mapsto x^q$。

\begin{lemma}
对 $c\in\mathcal O_{\widehat K}$，存在 $x\in\mathcal O_{\widehat K}$ 使
$\varphi(x)-x=c$。对 $c\in U_{\widehat K}$，存在 $x\in U_{\widehat K}$ 使
\[
\frac{\varphi(x)}x=c.
\]
\end{lemma}

\begin{proof}
先证乘法方程。把 $c$ 模极大理想化为 $\bar c$。剩余域是 $\kappa_K$ 的代数闭包，
所以存在 $\bar x_1\ne0$ 满足 $\bar x_1^q/\bar x_1=\bar c$。提升为
$x_1\in U_{\widehat K}$，则
\[
c=a_1\frac{\varphi(x_1)}{x_1},\qquad a_1\in U_{\widehat K}^{(1)}.
\]
接着在商 $U^{(r)}/U^{(r+1)}\simeq\bar\kappa$ 上重复同样步骤，逐次取
$x_r\in U_{\widehat K}^{(r-1)}$，使余项提高到 $U^{(r)}$。无限乘积
$x=x_1x_2\cdots$ 收敛，因为 $x_r\to1$，并满足 $\varphi(x)/x=c$。

加法方程使用加法过滤 $\mathfrak p^r/\mathfrak p^{r+1}\simeq\bar\kappa$。在每一层上
需要解 $y^q-y=\bar c$，这在代数闭剩余域中有解。逐层提升并取收敛和，得到
$\varphi(x)-x=c$。
\end{proof}

\begin{lemma}
设 $\pi'=u\pi$，$u\in U_K$，取 $f\in\mathcal F_\pi$、$f'\in\mathcal F_{\pi'}$。存在
\[
\theta(X)=\varepsilon X+\cdots\in\mathcal O_{\widehat K}[[X]],\qquad
\varepsilon\in U_{\widehat K},
\]
使
\[
\theta^\varphi(X)=\theta([u]_f(X)),
\]
并且
\[
\theta(F_f(X,Y))=F_{f'}(\theta(X),\theta(Y)),\qquad
\theta([a]_f(X))=[a]_{f'}(\theta(X))
\]
对所有 $a\in\mathcal O_K$ 成立。于是 $\theta$ 给出
$F_f(n)\simeq F_{f'}(n)$ 的 $\mathcal O_K$-模同构。
\end{lemma}

\begin{proof}
第一步构造 $\alpha(X)=\varepsilon X+\cdots$ 满足
$\alpha^\varphi=\alpha\circ[u]_f$。先选 $\varepsilon$ 使
$\varphi(\varepsilon)/\varepsilon=u$，这是上一引理的乘法方程。设已构造
$\alpha_r$ 使
\[
\alpha_r^\varphi(X)\equiv \alpha_r([u]_f(X))\pmod{X^{r+1}}.
\]
若误差的 $X^{r+1}$ 系数为 $c$，把
$\alpha_{r+1}=\alpha_r+a\varepsilon^{r+1}X^{r+1}$ 代入，新的系数要求满足
$a-\varphi(a)=c/(\varphi\varepsilon)^{r+1}$。这由上一引理的加法方程可解。递推取
极限得到 $\alpha$。

第二步把 $\alpha$ 修正为形式模同构。令
$h=\alpha^\varphi\circ f\circ\alpha^{-1}$。由
$\alpha^\varphi=\alpha\circ[u]_f$ 和 $\pi'=u\pi$，可检查 $h$ 的系数固定于
$K$，且 $h(X)\equiv\pi'X\pmod{\deg2}$、$h(X)\equiv X^q\pmod{\pi'}$，即
$h\in\mathcal F_{\pi'}$。由前面命题，存在唯一同态 $[1]_{f',h}$。取
\[
\theta=[1]_{f',h}\circ\alpha.
\]
它满足列出的三条相容关系。若 $x\in F_f(n)$，则
$[\pi^n]_{f'}(\theta(x))=\theta([\pi^n]_f(x))=0$，所以 $\theta(x)\in F_{f'}(n)$。
线性项 $\varepsilon$ 是单位，故 $\theta$ 只在 $0$ 处取 $0$；两边除点集大小相同，于是
该映射为同构。
\end{proof}

\begin{proposition}
若 $\pi'$ 也是 $K$ 的素元，则
\[
K^{\mathrm{nr}}E_{\pi,n}=K^{\mathrm{nr}}E_{\pi',n}.
\]
\end{proposition}

\begin{proof}
上一引理给出 $F_{f'}(n)=\theta(F_f(n))$，其中 $\theta$ 的系数在 $\widehat K$ 中。
故 $F_{f'}(n)$ 包含在 $K^{\mathrm{nr}}E_{\pi,n}$ 的完备化中。它们是代数元，因而实际
落在代数域 $K^{\mathrm{nr}}E_{\pi,n}$ 中。于是
$K^{\mathrm{nr}}E_{\pi',n}\subset K^{\mathrm{nr}}E_{\pi,n}$。交换 $\pi$ 与 $\pi'$ 得反向
包含。
\end{proof}

\begin{theorem}
设 $\omega_\pi:K^\times\to\Gal(K^{\mathrm{nr}}E_{\pi,n}/K)$ 由
\[
a=u\pi^m\mapsto
\begin{cases}
\varphi^m&\text{在 }K^{\mathrm{nr}}\text{ 上},\\
[u^{-1}]_f&\text{在 }F_f(n)\text{ 上}
\end{cases}
\]
定义。若 $\pi'$ 也是素元，则
\[
\omega_\pi(\pi')|_{E_{\pi',n}}=\operatorname{id}.
\]
\end{theorem}

\begin{proof}
写 $\pi'=u\pi$。令 $\theta$ 为上一引理中的除点变换。对 $x\in F_f(n)$，由定义
$\omega_\pi(\pi)$ 在 $K^{\mathrm{nr}}$ 上为 $\varphi$，在 $E_{\pi,n}$ 上为恒等；
$\omega_\pi(u)$ 在 $K^{\mathrm{nr}}$ 上为恒等，在除点上为 $[u^{-1}]_f$。因此先作用
$\omega_\pi(\pi)$ 会把 $\theta(x)$ 变为 $\theta^\varphi(x)$，再作用 $\omega_\pi(u)$ 得
\[
\omega_\pi(\pi')\theta(x)
=\omega_\pi(u)\omega_\pi(\pi)\theta(x)
=\theta^\varphi([u^{-1}]_f(x)).
\]
由 $\theta^\varphi=\theta\circ[u]_f$，这等于
\[
\theta([u]_f([u^{-1}]_f(x)))=\theta(x).
\]
这些 $\theta(x)$ 生成 $E_{\pi',n}$，所以 $\omega_\pi(\pi')$ 在 $E_{\pi',n}$ 上限制为
恒等。
\end{proof}

\section{数域的赋值}

现在把局部理论放回数域。局部域不是孤立出现的：数域 $F$ 的每个非零素理想
$\mathfrak p\subset\OF$ 都给出一个局部视角。若
\[
x\OF=\mathfrak p^{v_{\mathfrak p}(x)}\mathfrak a,\qquad
(\mathfrak p,\mathfrak a)=1,
\]
则 $v_{\mathfrak p}$ 是 $F$ 的标准离散赋值。它的完备化就是通常记作 $F_{\mathfrak p}$
的局部域。若 $E/F$ 是有限扩张，$\mathcal O_E$ 中位于 $\mathfrak p$ 上方的素理想
$\mathfrak P$ 给出 $E$ 上的标准赋值 $v_{\mathfrak P}$。本节说明：数域上的有限赋值
正是这些素理想赋值；Galois 群在这些延拓赋值上按传递方式作用；稳定一个延拓赋值
的子群就是分解群。

\begin{proposition}
设 $F$ 为数域，$\nu$ 为 $F$ 的离散赋值。则存在 $F$ 的非零素理想 $\mathfrak p$，使
$\nu$ 与 $v_{\mathfrak p}$ 等价。
\end{proposition}

\begin{proof}
令
\[
\mathcal O_\nu=\{x\in F:\nu(x)\ge0\},\qquad
\mathfrak m_\nu=\{x\in F:\nu(x)>0\}.
\]
把它与整数环相交，得到
\[
\mathfrak p=\mathfrak m_\nu\cap\OF.
\]
这是 $\OF$ 的素理想：若 $ab\in\mathfrak p$，则 $\nu(a)+\nu(b)>0$。对代数整数
$a,b$ 有 $\nu(a),\nu(b)\ge0$；因此至少一个赋值为正，对应元素落入 $\mathfrak p$。
它不是零理想，因为非平凡离散赋值在数域上会使某些代数整数具有正赋值；若交为零，
则 $\OF$ 的所有非零元素都在 $\mathcal O_\nu^\times$ 中，分式域中的每个元素都赋值
为 $0$，与非平凡性矛盾。

局部化 $\mathcal O_{F,\mathfrak p}$ 嵌入 $\mathcal O_\nu$。若 $a\notin\mathfrak p$，则
$\nu(a)=0$，所以 $a$ 在 $\mathcal O_\nu$ 中可逆。由局部化的泛性质得到
$\mathcal O_{F,\mathfrak p}\subset\mathcal O_\nu$。另一方面，$\mathcal O_{F,\mathfrak p}$ 是离散赋值环，
其分式域为 $F$，其赋值就是 $v_{\mathfrak p}$。两个离散赋值环在同一分式域中具有同一
中心 $\mathfrak p$，于是它们给出同一拓扑；等价地，存在正整数 $c$ 使
\[
\nu(x)=c\,v_{\mathfrak p}(x)\qquad(x\in F^\times).
\]
所以 $\nu$ 与 $v_{\mathfrak p}$ 等价。
\end{proof}

\begin{proposition}
设 $E/F$ 为有限可分扩张，且
\[
\mathfrak p\mathcal O_E=\mathfrak P_1^{e_1}\cdots\mathfrak P_g^{e_g}.
\]
则 $v_{\mathfrak p}$ 在 $E$ 上的全部延拓为
\[
w_i=\frac1{e_i}v_{\mathfrak P_i}\qquad(1\le i\le g).
\]
\end{proposition}

\begin{proof}
先验证 $w_i$ 确实延拓 $v_{\mathfrak p}$。若 $x\in F^\times$，则
\[
x\mathcal O_E=(x\OF)\mathcal O_E.
\]
若 $x\OF=\mathfrak p^{v_{\mathfrak p}(x)}\mathfrak a$ 且
$(\mathfrak p,\mathfrak a)=1$，扩张到 $\mathcal O_E$ 后，$\mathfrak P_i$ 在
$x\mathcal O_E$ 中的指数为 $e_i v_{\mathfrak p}(x)$。因此
\[
w_i(x)=\frac1{e_i}v_{\mathfrak P_i}(x)=v_{\mathfrak p}(x).
\]

反过来，设 $w$ 是 $E$ 上延拓 $v_{\mathfrak p}$ 的离散赋值。由上一命题，$w$ 等价于
某个素理想 $\mathfrak P\subset\mathcal O_E$ 给出的标准赋值。限制到 $F$ 后，其中心
就是 $\mathfrak P\cap\OF$。因为 $w|_F$ 与 $v_{\mathfrak p}$ 等价，这个中心必须是
$\mathfrak p$。所以 $\mathfrak P$ 是某个 $\mathfrak P_i$。归一化条件
$w|_F=v_{\mathfrak p}$ 决定比例因子只能是 $1/e_i$，故 $w=w_i$。
\end{proof}

\begin{proposition}
设 $(F,v)$ 为赋值域，$E/F$ 为 Galois 扩张，$w,w'$ 是 $v$ 在 $E$ 上的两个延拓。则
存在 $\sigma\in\Gal(E/F)$，使
\[
w'=w\circ\sigma.
\]
\end{proposition}

\begin{proof}
先设 $E/F$ 有限，记 $G=\Gal(E/F)$。若不存在这样的 $\sigma$，则两族赋值
\[
\{w\circ\sigma:\sigma\in G\},\qquad
\{w'\circ\sigma:\sigma\in G\}
\]
互不相交。由逼近定理，可取 $x\in E$，使对每个 $\sigma\in G$ 都有
\[
|\sigma x|_w<1,\qquad |\sigma x|_{w'}>1.
\]
令
\[
\alpha=N_{E/F}(x)=\prod_{\sigma\in G}\sigma x\in F.
\]
第一组不等式给出 $|\alpha|_v<1$，因为 $w$ 在 $F$ 上限制为 $v$；第二组不等式给出
$|\alpha|_v>1$。同一个 $F$ 中元素不可能同时满足这两个严格不等式，矛盾。故有限
Galois 情形成立。

若 $E/F$ 是无限 Galois 扩张，令 $M$ 遍历 $E/F$ 的有限 Galois 中间扩张，并设
\[
X_M=\{\sigma\in G:w'|_M=(w\circ\sigma)|_M\}.
\]
有限情形说明每个 $X_M$ 非空。集合 $X_M$ 是闭集：若 $\sigma\notin X_M$，则存在
$a\in M$ 使两边赋值不同，所有落在同一陪集 $\sigma\Gal(E/M)$ 中的元素在 $M$ 上
限制相同，所以该陪集包含在补集中。有限交性质也成立：对
$M_1,\dots,M_r$，取合成域 $M=M_1\cdots M_r$，则
$X_M\subset X_{M_1}\cap\cdots\cap X_{M_r}$。Galois 群 $G$ 是紧的射影极限，闭集族
具有有限交性质，故 $\bigcap_MX_M$ 非空。取其中的 $\sigma$，便有 $w'=w\circ\sigma$
在整个 $E$ 上成立。
\end{proof}

\begin{definition}
设 $(E,w)/(F,v)$ 为赋值 Galois 扩张。定义分解群
\[
D_{w|v}=\{\sigma\in\Gal(E/F):w\circ\sigma=w\}.
\]
它是保持所选延拓赋值 $w$ 的 Galois 元素集合。定义更高分歧群
\[
G_i=\{\sigma\in\Gal(E/F):w(\sigma a-a)\ge i+1,\ \forall a\in\mathcal O_E\}.
\]
其中 $\mathcal O_E=\{x\in E:w(x)\ge0\}$。$G_0$ 称为惯性群，记为 $I_{w|v}$。
\end{definition}

\begin{proposition}
设 $(E,w)/(F,v)$ 为赋值域 Galois 扩张，剩余域分别为 $\kappa_E,\kappa_F$。则
$\kappa_E/\kappa_F$ 为 Galois 扩张，并有正合列
\[
1\longrightarrow I_{w|v}\longrightarrow D_{w|v}
\longrightarrow \Gal(\kappa_E/\kappa_F)\longrightarrow1.
\]
\end{proposition}

\begin{proof}
对 $\sigma\in D_{w|v}$，若 $a\in\mathcal O_E$，则
$w(\sigma a)=w(a)\ge0$，所以 $\sigma$ 保持 $\mathcal O_E$。若 $a\in\mathfrak p_E$，
则 $w(\sigma a)>0$，所以 $\sigma$ 也保持极大理想。于是 $\sigma$ 在剩余域上诱导
$\kappa_F$-自同构，得到群同态
\[
D_{w|v}\longrightarrow\Gal(\kappa_E/\kappa_F).
\]
它的核正是对剩余域作用为恒等的元素，也就是 $I_{w|v}$。

剩余域扩张为 Galois 的原因来自 Galois 正规性：若 $\bar a\in\kappa_E$ 的最小多项式
在 $\kappa_F$ 上有一个根，则可提升到 $\mathcal O_E$ 中的元素并用 $E/F$ 的正规性
把共轭带回 $E$，再模极大理想得到所有剩余共轭。

有限扩张时，满射就是第一章分解群到剩余域 Galois 群的标准满射。无限扩张时，令
$Z_w$ 为 $D_{w|v}$ 的固定域。$Z_w$ 的剩余域为 $\kappa_F$。给定
$\bar\sigma\in\Gal(\kappa_E/\kappa_F)$ 和一个有限 Galois 子扩张
$\mu/\kappa_F$，取 $E/Z_w$ 的有限中间扩张 $M/Z_w$，使 $\kappa_M$ 包含 $\mu$。
有限情形给出
\[
\Gal(M/Z_w)\longrightarrow\Gal(\kappa_M/\kappa_F)
\]
满射，因此可取 $\tau\in D_{w|v}$，使它在 $\mu$ 上与 $\bar\sigma$ 相同。这样的
$\tau$ 在所有有限剩余子扩张上逼近 $\bar\sigma$，所以像是稠密子群。又
$D_{w|v}$ 紧，连续像紧从而闭；既稠密又闭，像就是整个
$\Gal(\kappa_E/\kappa_F)$。
\end{proof}

\begin{proposition}
设 $(E,w)/(F,v)$ 为赋值域有限 Galois 扩张，$F_v,E_w$ 分别为完备化。则
$E_w/F_v$ 为 Galois 扩张，并有自然同构
\[
D_{w|v}\simeq\Gal(E_w/F_v).
\]
\end{proposition}

\begin{proof}
若 $\sigma\in D_{w|v}$，则 $w(\sigma x-\sigma y)=w(x-y)$，所以 $\sigma$ 对 $w$-拓扑
连续，并唯一延拓为 $E_w$ 的连续自同构。它固定 $F$，故固定 $F_v$，于是得到同态
\[
D_{w|v}\to\Gal(E_w/F_v).
\]
单射来自 $E$ 在 $E_w$ 中稠密：若延拓后为恒等，则在稠密子域 $E$ 上已经为恒等。

证明满射。任取 $\tau\in\Gal(E_w/F_v)$。由于 $E/F$ 有限，完备化张量分解给出
\[
E\otimes_FF_v\simeq\prod_{w_i|v}E_{w_i}.
\]
其中因子由 $v$ 在 $E$ 上的延拓赋值编号。固定所选因子 $E_w$，$F_v$-自同构
$\tau$ 保持该因子，并把稠密的 $E$ 的像带回同一个因子。限制到 $E$ 后得到一个
$F$-嵌入 $E\to E_w$。因为 $E/F$ 是 Galois，这个嵌入来自某个
$\sigma\in\Gal(E/F)$。它保持所选完备化因子，等价于 $w\circ\sigma=w$，即
$\sigma\in D_{w|v}$。这个 $\sigma$ 的连续延拓就是 $\tau$。因此同态满。

既然 $E_w$ 是有限 Galois 扩张 $E/F$ 在一个延拓赋值处的完备化，其自同构群由上面的
分解群给出，故 $E_w/F_v$ 是 Galois 扩张。
\end{proof}

\section*{习题详解}
\addcontentsline{toc}{section}{习题详解}

\begin{exercise}
设 $\widehat F$ 是赋值域 $F$ 的完备化。证明：
\[
v(\widehat F^\times)=v(F^\times),\qquad
\kappa_{\widehat F}\simeq\kappa_F .
\]
\end{exercise}

\begin{proof}[解答]
状态：solvable（可解）。

记 $F$ 的赋值环和极大理想为
\[
\mathcal O_F=\{x:v(x)\ge0\},\qquad
\mathfrak p_F=\{x:v(x)>0\}.
\]
完备化 $\widehat F$ 中的元素由 $F$ 中 Cauchy 序列表示。先证明值群没有变大。取
$x\in\widehat F^\times$，选 $x_n\in F$ 且 $x_n\to x$。因为 $x\ne0$，存在 $N_0$
使 $x_n\ne0$ 且 $v(x_n-x)>v(x)$ 对 $n\ge N_0$ 成立。于是
\[
v(x_n)=v\bigl(x+(x_n-x)\bigr)=v(x)
\]
对 $n\ge N_0$ 成立；这里用的是非 Archimede 赋值的基本事实：若 $v(y)>v(x)$，
则 $v(x+y)=v(x)$。所以 $v(x)=v(x_n)\in v(F^\times)$。另一方向来自嵌入
$F^\times\subset\widehat F^\times$，因此两个值群相等。

再证明剩余域不变。令 $\mathcal O_{\widehat F}$、$\mathfrak p_{\widehat F}$ 为
完备化的赋值环和极大理想。嵌入 $\mathcal O_F\to\mathcal O_{\widehat F}$ 诱导
\[
\mathcal O_F/\mathfrak p_F\longrightarrow
\mathcal O_{\widehat F}/\mathfrak p_{\widehat F}.
\]
它的核正是 $\mathfrak p_F$。只需证明满射。取
$\bar x\in\mathcal O_{\widehat F}/\mathfrak p_{\widehat F}$，选代表
$x\in\mathcal O_{\widehat F}$。由 $F$ 在 $\widehat F$ 中稠密，可取 $a\in F$
满足 $v(x-a)>0$。因为 $v(x)\ge0$ 且 $v(x-a)>0$，若 $v(a)<0$，则
$v(x-a)=\min\{v(x),v(a)\}<0$，矛盾；故 $a\in\mathcal O_F$。又
$x-a\in\mathfrak p_{\widehat F}$，所以 $x$ 与 $a$ 有同一剩余类。诱导映射满射，
从而得到 $\kappa_F\simeq\kappa_{\widehat F}$。
\end{proof}

\begin{exercise}
设不可约多项式
\[
f(X)=X^n+a_{n-1}X^{n-1}+\cdots+a_0\in\Q_p[X],
\]
且 $a_0\in\Z_p$。证明 $a_{n-1},\dots,a_0\in\Z_p$。
\end{exercise}

\begin{proof}[解答]
状态：solvable（可解）。

令 $\alpha$ 为 $f$ 的一个根，$L=\Q_p(\alpha)$。因为 $f$ 在 $\Q_p[X]$ 中首一且不可约，
它是 $\alpha$ 在 $\Q_p$ 上的最小多项式。完备离散赋值在有限扩张中唯一延拓，所以
$f$ 的所有共轭根在代数闭包中的赋值相同。记这个共同赋值为 $t$。

由根与常数项的关系，
\[
a_0=(-1)^n\prod_{i=1}^n\alpha_i,
\]
其中 $\alpha_i$ 是 $f$ 的全部共轭根。因此
\[
v_p(a_0)=\sum_{i=1}^n v_p(\alpha_i)=nt .
\]
题设 $a_0\in\Z_p$ 给出 $v_p(a_0)\ge0$，于是 $t\ge0$。这说明每个共轭根
$\alpha_i$ 都是代数整数，也就是属于相应有限扩张的整数环。

系数 $a_{n-j}$ 是带符号的 $j$ 次初等对称多项式：
\[
a_{n-j}=(-1)^j\sum_{1\le i_1<\cdots<i_j\le n}
\alpha_{i_1}\cdots\alpha_{i_j}.
\]
整数环在加法和乘法下封闭，所以每个 $a_{n-j}$ 是 $\Q_p$ 中的整元素。$\Q_p$ 中的
整元素正是 $\Z_p$，故 $a_{n-1},\ldots,a_0\in\Z_p$。
\end{proof}

\begin{exercise}
设
\[
f(X)=a_nX^n+\cdots+a_1X+a_0\in\Z_p[X],
\]
满足 $|a_n|=1$，$|a_i|<1$ 对 $0<i<n$ 成立，且 $|a_0|=1/p$。证明 $f$ 在 $\Q_p$ 上不可约。
\end{exercise}

\begin{proof}[解答]
状态：solvable（可解）。

把绝对值翻译成 $p$ 进赋值。$|a_n|=1$ 表示 $a_n$ 是 $\Z_p$ 的单位；
$|a_i|<1$ 对 $0<i<n$ 成立表示 $a_i\in p\Z_p$；$|a_0|=1/p$ 表示
$a_0\in p\Z_p$ 但 $a_0\notin p^2\Z_p$。

假设 $f=gh$ 在 $\Q_p[X]$ 中有非平凡分解。由于 $a_n$ 是单位，用 Gauss 引理可把
$g,h$ 取在 $\Z_p[X]$ 中，并且二者首项系数都是单位。模 $p$ 后，
\[
\bar f(X)=\bar a_nX^n
\]
是一个非零常数乘 $X^n$。因此 $\bar g$ 与 $\bar h$ 都只能是非零常数乘 $X$ 的幂；
由于分解非平凡，两个多项式的次数都正，所以这两个幂的指数都正。于是
$g(0),h(0)\in p\Z_p$，从而
\[
a_0=f(0)=g(0)h(0)
\]
属于 $p^2\Z_p$。这与 $a_0\notin p^2\Z_p$ 矛盾，所以 $f$ 在 $\Q_p$ 上不可约。
\end{proof}

\begin{exercise}
证明：对每个 $n$，存在唯一无分歧扩张 $F_n/\Q_p$，满足 $[F_n:\Q_p]=n$，且 $F_n=\Q_p(\zeta)$，其中 $\zeta$ 是 $p^n-1$ 次本原单位根。并证明 $\Q_5(\sqrt2)=\Q_5(\zeta)$，其中 $\zeta$ 为 $24$ 次本原单位根。
\end{exercise}

\begin{proof}[解答]
状态：solvable（可解）。

先说明存在性。$\Q_p$ 的剩余域为 $\F_p$，有限域在每个次数 $n$ 上有唯一扩张
$\F_{p^n}$。取 $\F_{p^n}^\times$ 的生成元 $\bar\zeta$，它的阶为 $p^n-1$。多项式
$X^{p^n-1}-1$ 在 $\bar\zeta$ 处的导数为
\[
(p^n-1)\bar\zeta^{p^n-2},
\]
该元素在特征 $p$ 的剩余域中非零，所以 Hensel 引理把 $\bar\zeta$ 唯一提升为一个
$(p^n-1)$ 次单位根 $\zeta$。由 $\bar\zeta$ 生成剩余域可得
\[
F_n=\Q_p(\zeta),\qquad [F_n:\Q_p]=n,
\]
且该扩张的剩余域就是 $\F_{p^n}$，分歧指数为 $1$，所以它无分歧。

再说明唯一性。若 $E/\Q_p$ 是 $n$ 次无分歧扩张，则剩余域为 $\F_{p^n}$。其乘法群中
的本原元可唯一 Teichmuller 提升到 $E$ 中，所以 $E$ 含有上述 $\zeta$，并且
$\Q_p(\zeta)$ 已经有次数 $n$。因此 $E=\Q_p(\zeta)$。

对 $p=5,n=2$，唯一二次无分歧扩张为 $\Q_5(\zeta_{24})$。在 $\F_5$ 中平方剩余为
$1,4$，所以 $X^2-2$ 没有根。它模 $5$ 不可约，且判别式 $8$ 不被 $5$ 整除，因此
$\Q_5(\sqrt2)/\Q_5$ 的剩余域次数为 $2$、分歧指数为 $1$。它也是唯一的二次无分歧扩张，
故
\[
\Q_5(\sqrt2)=\Q_5(\zeta_{24}).
\]
\end{proof}

\begin{exercise}
在 $5$ 进情形中证明题中关于 $\sqrt2,\sqrt3,\sqrt6$ 和 $24$ 次单位根的断言。
\end{exercise}

\begin{proof}[解答]
状态：solvable（可解）。

(1) 在 $\F_5$ 中非零平方为 $1^2=1$、$2^2=4$、$3^2=4$、$4^2=1$。因此
$2,3$ 都不是平方，$X^2-2$ 与 $X^2-3$ 在 $\F_5[X]$ 中不可约。两个多项式的判别式
分别为 $8$ 和 $12$，都不是 $5$ 的倍数，所以相应二次扩张没有分歧。于是
$\Q_5(\sqrt2)$ 与 $\Q_5(\sqrt3)$ 都是 $\Q_5$ 的二次无分歧扩张。

(2) 二次无分歧扩张唯一，故
\[
\Q_5(\sqrt2)=\Q_5(\sqrt3).
\]
在这个共同的域中取
\[
\alpha=\sqrt2\sqrt3.
\]
则 $\alpha^2=6$。又
\[
\frac{\alpha\sqrt3}{3}
=\frac{\sqrt2(\sqrt3)^2}{3}
=\sqrt2.
\]
若把 $\sqrt3$ 换成另一个根，等式右端换成负号，所以题中的 $\pm$ 来自平方根选择。

(3) 设 $\beta=\sqrt2$，则整数环为 $\Z_5[\beta]$，剩余域为
\[
\F_5[\bar\beta]\simeq \F_5[T]/(T^2-2).
\]
取
\[
\bar\xi=1+2\bar\beta .
\]
在关系 $\bar\beta^2=2$ 下直接计算：
\[
\bar\xi^2=4+4\bar\beta,\quad
\bar\xi^3=2\bar\beta,\quad
\bar\xi^6=3,\quad
\bar\xi^{12}=4=-1,\quad
\bar\xi^{24}=1.
\]
由于 $\bar\xi^{12}\ne1$，并且 $\bar\xi^6\ne1$、$\bar\xi^8=2+2\bar\beta\ne1$，它的阶不是
$1,2,3,4,6,8,12$，只能是 $24$。

现在把 $\bar\xi$ 提升为 $X^{24}-1$ 的根。导数 $24X^{23}$ 在模 $5$ 后非零，所以 Hensel
提升唯一。用 Newton 迭代
\[
x_{r+1}=x_r-\frac{x_r^{24}-1}{24x_r^{23}}
\]
从 $x_0=1+2\beta$ 出发，在环 $(\Z/125\Z)[\beta]/(\beta^2-2)$ 中得到
\[
\xi\equiv 26+77\beta \pmod{125}.
\]
写成 $5$ 进前三位就是
\[
\xi\equiv (1+25)+(2+3\cdot25)\beta \pmod{5^3}.
\]
直接代回可验算 $(26+77\beta)^{24}\equiv1\pmod{125}$，且它模 $5$ 的剩余类已经是
$24$ 阶元，所以该提升是 $24$ 次本原单位根。
\end{proof}

\begin{exercise}
按题中较弱定义讨论绝对值的等价关系，并证明等价绝对值相差正实数幂，且定义相同拓扑。
\end{exercise}

\begin{proof}[解答]
状态：solvable（可解）。

(1) 对非平凡绝对值定义 $|\cdot|_1\equiv|\cdot|_2$ 为
\[
|x|_1<1\Longrightarrow |x|_2<1.
\]
自反性由定义直接得到。为了证明对称性和传递性，先证明一个比较引理：若
$|x|_1<1\Rightarrow |x|_2<1$，则存在 $\alpha>0$ 使
$|x|_2=|x|_1^\alpha$ 对所有 $x\in F^\times$ 成立。

取 $a$ 使 $|a|_1<1$；由假设 $|a|_2<1$。对任意 $x$ 和正整数 $m,n$，若
\[
|x|_1^m<|a|_1^n
\Longleftrightarrow
|x^m a^{-n}|_1<1,
\]
则
\[
|x|_2^m<|a|_2^n.
\]
若 $|x|_1^m>|a|_1^n$，则 $|a^n x^{-m}|_1<1$，从而
$|a|_2^n<|x|_2^m$。所以两个实数
\[
\frac{\log |x|_1}{\log |a|_1}
\quad\text{与}\quad
\frac{\log |x|_2}{\log |a|_2}
\]
与每个有理数 $n/m$ 的大小关系一致。实数由有理数的截断确定，所以二者相等。令
\[
\alpha=\frac{\log |a|_2}{\log |a|_1}>0,
\]
得到
\[
\log |x|_2=\alpha\log |x|_1,
\]
也就是 $|x|_2=|x|_1^\alpha$。当 $|x|_1=1$ 时两边都等于 $1$，同样成立。
这个幂关系立刻给出反向蕴含，因此关系是对称的。若
$|\cdot|_1\equiv|\cdot|_2$ 且 $|\cdot|_2\equiv|\cdot|_3$，蕴含合成给出
$|\cdot|_1\equiv|\cdot|_3$，所以它也是传递的。

(2) 上面的比较引理已经给出：若 $|\cdot|_2$ 等价于 $|\cdot|_1$，则存在
$\alpha>0$，使 $|x|_2=|x|_1^\alpha$。

(3) 若 $|x|_2=|x|_1^\alpha$，则 $|x-a|_2<\varepsilon$ 等价于 $|x-a|_1<\varepsilon^{1/\alpha}$。因此基本邻域相同，拓扑相同。
\end{proof}

\begin{exercise}
设 $F$ 有非 Archimede 绝对值。证明强三角不等式的基本后果、Cauchy 判别、级数判别和双重级数换序。
\end{exercise}

\begin{proof}[解答]
状态：solvable（可解）。

(1) 题设中的“非亚绝对值”就是满足常数 $c=1$ 的情形，所以
\[
|a+b|\le\max\{|a|,|b|\}.
\]

(2) 若 $|a|<|b|$，则
\[
|b|=|(a+b)-a|\le\max\{|a+b|,|a|\}.
\]
由于 $|a|<|b|$，只能有 $|a+b|\ge |b|$。另一边由强三角不等式 $|a+b|\le |b|$，故相等。

(3) 若 $\{a_n\}$ Cauchy，则 $|a_{n+1}-a_n|\to0$。反过来，若相邻差趋于 $0$，则对 $m>n$，
\[
|a_m-a_n|
\le\max_{n\le i<m}|a_{i+1}-a_i|,
\]
右边随 $n\to\infty$ 变为 $0$，故 Cauchy。

(4) 设部分和 $s_N=\sum_{n=0}^N a_n$。相邻差为 $s_{N+1}-s_N=a_{N+1}$。由 (3)，
部分和 Cauchy 当且仅当 $a_n\to0$。在完备情形中这等价于级数收敛；若未假定完备，
则结论应理解为“部分和 Cauchy 当且仅当项趋于零”。当级数收敛时，有限和满足
\[
\left|\sum_{n=0}^N a_n\right|\le\max_{0\le n\le N}|a_n|.
\]
令 $N\to\infty$ 得
\[
\left|\sum_{n=0}^{\infty}a_n\right|
\le\max_n|a_n|
\]
。

(5) 设
\[
S_{M,N}=\sum_{i=0}^M\sum_{j=0}^N b_{ij}.
\]
要证明两个累次和相等，只需证明矩形部分和有同一个极限。给定 $\varepsilon>0$。由
“对 $j$ 一致趋于零”可取 $I$，使 $i\ge I$ 时对所有 $j$ 都有 $|b_{ij}|<\varepsilon$。
对有限多个 $i<I$，条件 $\lim_{j\to\infty}b_{ij}=0$ 给出 $J$，使 $j\ge J$ 且
$i<I$ 时 $|b_{ij}|<\varepsilon$。于是只要 $(i,j)$ 在矩形
$\{0,\ldots,I-1\}\times\{0,\ldots,J-1\}$ 外，就有 $|b_{ij}|<\varepsilon$。
强三角不等式说明任意有限尾部和的绝对值不超过尾部各项绝对值的最大值，因此矩形部分和
是 Cauchy。先令 $N\to\infty$ 再令 $M\to\infty$，或先令 $M\to\infty$ 再令
$N\to\infty$，都得到这个矩形极限，所以两种累次求和收敛并相等。
\end{proof}

\begin{exercise}
设非 Archimede 域上幂级数 $f(X)=\sum a_nX^n$，令 $1/r=\limsup\sqrt[n]{|a_n|}$。证明收敛半径判别和恒等定理。
\end{exercise}

\begin{proof}[解答]
状态：solvable（可解）。

在非 Archimede 域中，幂级数在给定 $x$ 处收敛等价于项 $a_nx^n$ 趋于 $0$。令
\[
\rho=\limsup_{n\to\infty}\sqrt[n]{|a_n|}=1/r.
\]
则
\[
\limsup\sqrt[n]{|a_nx^n|}
=|x|\limsup\sqrt[n]{|a_n|}
=|x|/r.
\]
若 $|x|<r$，取 $\lambda<1$ 使 $|x|\rho<\lambda$，则从某项起
$|a_nx^n|\le \lambda^n$，故项趋于 $0$。若 $|x|>r$，则有无穷多个 $n$ 满足
$\sqrt[n]{|a_nx^n|}>1$，项不趋于 $0$。当 $|x|=r$ 时，项的绝对值就是
$|a_n|r^n$，所以边界是否收敛完全由 $\lim |a_n|r^n$ 是否为 $0$ 决定。
这分别给出题中 (1)--(4)，其中 $r=0$ 和 $r=\infty$ 是同一判别的两个端点。

证明 (5)。令 $h=f-g=\sum c_nX^n$。若 $h\ne0$，取最小 $N$ 使 $c_N\ne0$，则
\[
h(X)=X^N(c_N+c_{N+1}X+\cdots).
\]
括号内的级数在 $X$ 足够接近 $0$ 时收敛，并且尾部
$\sum_{k>N}c_kX^{k-N}$ 的绝对值小于 $|c_N|$。强三角不等式给出括号内绝对值等于
$|c_N|$，所以它不为零。由于 $a_m\to0$ 且 $a_m\ne0$，充分大的 $m$ 会满足
$h(a_m)\ne0$，这与 $f(a_m)=g(a_m)$ 矛盾。因此 $h=0$，即 $f(X)=g(X)$。
\end{proof}

\begin{exercise}
证明 $\Q$ 的任一非平凡绝对值等价于通常绝对值或某个 $p$ 进绝对值。
\end{exercise}

\begin{proof}[解答]
状态：solvable（可解）。

先证明 $|\cdot|_{p,c}\equiv|\cdot|_{p,c'}$。若 $r=p^vb$ 且 $p\nmid b$，则
$|r|_{p,c}=c^v$、$|r|_{p,c'}=(c')^v$。因此
\[
|r|_{p,c'}=|r|_{p,c}^{\log c'/\log c},
\]
二者相差正实数幂，故等价。

下面证明 $\Q$ 上的分类。若存在整数 $b\ge2$ 使 $|b|>1$，取任意正整数 $m$，写成
$b$ 进展开
\[
m=c_0+c_1b+\cdots+c_kb^k,\qquad 0\le c_i<b.
\]
令 $C=\max_{0\le c<b}|c|$。由题中绝对值条件反复使用加法估计，可得
$|m|\le C'(k+1)^A |b|^k$，其中常数只与该绝对值和 $b$ 有关。因为
$b^k\le m<b^{k+1}$，取对数并让 $m$ 的幂次放大，可推出
\[
|m|=m^\alpha
\]
对某个 $\alpha>0$ 成立。于是绝对值等价于通常绝对值。

若所有整数 $n$ 都满足 $|n|\le1$，则由加法条件得到非 Archimede 三角不等式。非平凡性
保证存在素数 $p$ 使 $|p|<1$。若另一个素数 $\ell$ 也满足 $|\ell|<1$，取
$ap^r+b\ell^s=1$，当 $r,s$ 很大时右边两项绝对值都小于 $1$，强三角不等式会给出
$|1|<1$，矛盾。因此只有一个素数 $p$ 的绝对值小于 $1$，其他素数绝对值都等于 $1$。
对有理数 $x=p^v a/b$ 且 $p\nmid ab$，得到
\[
|x|=|p|^v,
\]
所以它等价于某个 $|\,\cdot\,|_{p,c}$。
\end{proof}

\begin{exercise}
设 $(F,|\cdot|)$ 为完备赋值域，且 $|\cdot|$ 为 Archimede 绝对值。证明 $F$ 是 $\R$ 或 $\C$。
\end{exercise}

\begin{proof}[解答]
状态：solvable（可解）。

由上一题，$|\cdot|$ 在 $\Q$ 上等价于通常绝对值。把绝对值取一个正幂后，可设它在
$\Q$ 上就是通常绝对值。完备性给出一个等距嵌入 $\R\hookrightarrow F$。

若 $F$ 含有 $i$，则把 $F$ 看作复 Banach 代数。对任意 $a\in F$，若
$a-\lambda$ 对每个 $\lambda\in\C$ 都非零，则
\[
R(\lambda)=(a-\lambda)^{-1}
\]
是定义在整个复平面上的 $F$ 值解析函数。当 $|\lambda|>2|a|$ 时，
\[
(a-\lambda)^{-1}
=-\lambda^{-1}\bigl(1-a/\lambda\bigr)^{-1}
=-\lambda^{-1}\sum_{n\ge0}(a/\lambda)^n,
\]
因此 $R(\lambda)\to0$。对任意连续复线性泛函 $\ell:F\to\C$，函数
$\ell(R(\lambda))$ 是有界整函数，所以为常数，并且由无穷远处极限为 $0$ 得它恒为
$0$。Hahn--Banach 定理使连续线性泛函能分离点，于是 $R(\lambda)=0$，矛盾。
因此存在 $\lambda\in\C$ 使 $a=\lambda$。每个 $a$ 都在 $\C$ 中，所以 $F=\C$。

若 $F$ 不含 $i$，构造 $F(i)=F[X]/(X^2+1)$。它仍是完备 Archimede 赋值代数；并且
对 $x+iy\ne0$ 有逆元
\[
(x+iy)^{-1}=\frac{x-iy}{x^2+y^2},
\]
因为 $x^2+y^2=0$ 会推出 $i=x/y\in F$。于是 $F(i)$ 是含 $i$ 的完备 Archimede 域，
由上一段 $F(i)\simeq\C$。复共轭的固定域是 $\R$，而 $F$ 正是 $F(i)$ 中共轭固定的
部分，所以 $F\simeq\R$。
\end{proof}

\begin{exercise}
证明分圆域的基本事实：本原根判别、次数、Galois 群、分圆多项式公式和根差乘积。
\end{exercise}

\begin{proof}[解答]
状态：solvable（可解）。

(1) 若 $\zeta$ 为 $m$ 次本原单位根，则 $\zeta^a$ 的阶为 $m/\gcd(a,m)$。所以它仍为 $m$ 次本原根当且仅当 $(a,m)=1$。

(2) 所有与 $m$ 互素的 $a$ 给出本原根 $\zeta^a$，共有 $\varphi(m)$ 个。定义
\[
\Phi_m(X)=\prod_{(a,m)=1}(X-\zeta^a).
\]
它在 Galois 群作用下不变，故属于 $\Q[X]$；又它是整系数首一多项式。$\zeta$ 的
最小多项式的根必须是这些本原根，反过来每个本原根由 $\Q$-嵌入 $\zeta\mapsto\zeta^a$
得到，所以最小多项式就是 $\Phi_m$，次数为 $\varphi(m)$。

(3) 任意 $\sigma\in\Gal(\Q(\zeta)/\Q)$ 把 $\zeta$ 送到另一个本原根 $\zeta^a$，$(a,m)=1$。映射 $\sigma\mapsto a\bmod m$ 给出
\[
\Gal(\Q(\zeta)/\Q)\simeq(\Z/m\Z)^\times.
\]

(4)
\[
\Phi_m(X)=\prod_{(a,m)=1}(X-\zeta^a).
\]
所有 $m$ 次单位根按精确阶 $d|m$ 分组，故
\[
X^m-1=\prod_{d|m}\Phi_d(X).
\]
Möbius 反演给出
\[
\Phi_m(X)=\prod_{d|m}(X^d-1)^{\mu(m/d)}.
\]

(5) 对 $X^m-1$ 求导得 $mX^{m-1}$。在根 $\xi_i$ 处，
\[
\prod_{j\ne i}(\xi_i-\xi_j)=m\xi_i^{m-1}.
\]
对全部 $i$ 相乘：
\[
\prod_i\prod_{j\ne i}(\xi_i-\xi_j)
=m^m\left(\prod_i\xi_i\right)^{m-1}.
\]
多项式 $X^m-1$ 的根之积为 $(-1)^{m+1}$，所以
\[
\left(\prod_i\xi_i\right)^{m-1}=(-1)^{(m+1)(m-1)}=(-1)^{m-1}.
\]
因此
\[
\prod_i\prod_{j\ne i}(\xi_i-\xi_j)=(-1)^{m-1}m^m.
\]
\end{proof}

\begin{exercise}
证明 profinite 整数环
\[
\widehat\Z=\varprojlim_n\Z/n\Z
\]
同构于 $\prod_p\Z_p$，并证明相关闭包和闭理想断言。
\end{exercise}

\begin{proof}[解答]
状态：solvable（可解）。

由中国剩余定理，若 $n=\prod_p p^{v_p(n)}$，则
\[
\Z/n\Z\simeq\prod_{p|n}\Z/p^{v_p(n)}\Z.
\]
对 $n$ 取逆极限，得到
\[
\widehat\Z\simeq\prod_p\varprojlim_r\Z/p^r\Z
=\prod_p\Z_p.
\]
模 $n$ 后，
\[
\widehat\Z/n\widehat\Z\simeq\Z/n\Z.
\]
逆极限拓扑中，$n\widehat\Z$ 是 $0$ 的开邻域，且所有这样的集合构成邻域基。对角嵌入
$\Z\to\widehat\Z$ 在每个有限商 $\Z/n\Z$ 上满射，所以 $\Z$ 在 $\widehat\Z$ 中稠密。

交集 $\Z\cap m\widehat\Z=m\Z$：若整数 $a$ 在每个 $\Z_p$ 中都属于 $m\Z_p$，则
$v_p(a)\ge v_p(m)$ 对所有 $p$ 成立，故 $m|a$；反向包含由定义给出。

最后证明闭理想是主理想。把闭理想 $I$ 投影到各分量 $\Z_p$，得闭理想 $I_p$。因为
$\Z_p$ 是离散赋值环，闭理想只能是 $p^{r_p}\Z_p$，其中
$r_p\in\Z_{\ge0}\cup\{\infty\}$，约定 $p^\infty\Z_p=\{0\}$。于是
\[
I=\prod_p p^{r_p}\Z_p.
\]
取 $a=(a_p)_p\in\prod_p\Z_p$，其中 $v_p(a_p)=r_p$；若 $r_p=\infty$ 就取
$a_p=0$。则 $I=a\widehat\Z$，所以 $I$ 为主理想。
\end{proof}

\begin{exercise}
设 Hausdorff 拓扑群 $G$ 的单位元有有限指数开子群邻域基。证明 $n\mapsto\sigma^n$ 可延拓为 $\widehat\Z\to G$ 的连续同态，并描述有限群情形的核。
\end{exercise}

\begin{proof}[解答]
状态：solvable（可解）。

令 $\phi:\Z\to G$，$\phi(n)=\sigma^n$。若 $H_\iota$ 是单位元邻域基中的开子群，则
\[
\phi^{-1}(H_\iota)=\{n:\sigma^n\in H_\iota\}
\]
是 $\Z$ 的子群；又 $\phi(\Z)H_\iota/H_\iota$ 在有限集合 $G/H_\iota$ 中，所以
$\phi^{-1}(H_\iota)$ 有有限指数。这正是 $\Z$ 的 profinite 拓扑中的开子群，所以
$\phi$ 连续。由 $\widehat\Z$ 的泛性质，$\phi$ 唯一延拓为连续同态
\[
\widehat\phi:\widehat\Z\to G.
\]
像 $\widehat\phi(\widehat\Z)$ 是 $\phi(\Z)$ 的闭包，即 $\langle\sigma\rangle$。

若 $G$ 有限，设 $\sigma$ 的阶为 $d$。原映射 $\Z\to G$ 的核是 $d\Z$。延拓到
$\widehat\Z$ 后，核是 $d\Z$ 在 $\widehat\Z$ 中的闭包，即 $d\widehat\Z$。像为
有限循环群 $\langle\sigma\rangle$。
\end{proof}

\begin{exercise}
证明 $\Z_p^\times$ 的结构，并描述 $\widehat\Z^\times$ 的挠子闭包。
\end{exercise}

\begin{proof}[解答]
状态：solvable（可解）。

若 $p\ne2$，有分解
\[
\Z_p^\times\simeq \mu_{p-1}\times(1+p\Z_p).
\]
这里 $\mu_{p-1}$ 是 Teichmuller 根部分。对数级数在 $1+p\Z_p$ 上收敛，并给出拓扑群同构
$1+p\Z_p\simeq p\Z_p$。作为 $\Z_p$ 加法群，$p\Z_p\simeq\Z_p$，所以
\[
\Z_p^\times\simeq \Z/(p-1)\Z\times\Z_p.
\]
当 $p=2$，
\[
\Z_2^\times\simeq\{\pm1\}\times(1+4\Z_2)\simeq \Z/2\Z\times\Z_2.
\]
由 $\widehat\Z^\times\simeq\prod_p\Z_p^\times$，逐个素数取有限根部分，得到子群
\[
\Z/2\Z\oplus\bigoplus_{p\ne2}\Z/(p-1)\Z
\]
嵌入 $\widehat\Z^\times$。记有限根部分的乘积为
\[
\widehat T=\Z/2\Z\times\prod_{p\ne2}\Z/(p-1)\Z.
\]
每个挠子元素在 $p$ 分量中必须落入 $\Z_p^\times$ 的有限根部分，因为 $\Z_p$ 加法部分
无非平凡挠元。反过来，有限个分量非平凡的根部分元素是挠元。因此挠子子群的闭包正是
$\widehat T$。
\end{proof}

\begin{exercise}
设 $\Omega=\bigcup_n\Q(\mu_n)$。证明 $\Gal(\Omega/\Q)\simeq\widehat\Z^\times$，并描述题中各固定域的 Galois 群。
\end{exercise}

\begin{proof}[解答]
状态：solvable（可解）。

对每个 $n$，
\[
\Gal(\Q(\mu_n)/\Q)\simeq(\Z/n\Z)^\times.
\]
对 $n$ 取逆极限得
\[
\Gal(\Omega/\Q)\simeq\varprojlim_n(\Z/n\Z)^\times=\widehat\Z^\times.
\]
若 $\widetilde\Q$ 为挠子子群固定域，则
\[
\Gal(\widetilde\Q/\Q)
\simeq \widehat\Z^\times/\overline{(\widehat\Z^\times)_{\mathrm{tor}}}.
\]
由上一题，商去有限根部分后，每个 $p$ 分量留下一个 $\Z_p$，因此
\[
\Gal(\widetilde\Q/\Q)\simeq\prod_p\Z_p=\widehat\Z.
\]

对 $\Omega_p=\bigcup_n\Q(\mu_{p^n})$，
\[
\Gal(\Omega_p/\Q)\simeq\Z_p^\times.
\]
其挠子部分为 $\mu_{p-1}$（当 $p=2$ 时为 $\{\pm1\}$，按低层根统一处理），固定域
$\widetilde\Q^{(p)}$ 的 Galois 群为
\[
\Gal(\widetilde\Q^{(p)}/\Q)\simeq\Z_p.
\]
所有 $p$ 分量的合成对应于直积分解
$\widehat\Z^\times\simeq\prod_p\Z_p^\times$，所以 $\widetilde\Q$ 是各
$\widetilde\Q^{(p)}$ 的合成域。
\end{proof}

\begin{exercise}
证明
\[
U^{(0)}\simeq\varprojlim_i U^{(0)}/U^{(i)}.
\]
\end{exercise}

\begin{proof}[解答]
状态：solvable（可解）。

这是 $\mathcal O^\times$ 完备性的单位群版本。自然映射
\[
U^{(0)}\to\varprojlim_iU^{(0)}/U^{(i)}
\]
的核为 $\cap_iU^{(i)}=\{1\}$，因为若 $u-1$ 落在所有 $\mathfrak p^i$ 中，则
$u-1=0$。给定兼容系 $(u_i)$，它在 $\mathcal O/\mathfrak p^i$ 中给出兼容单位类。由
\[
\mathcal O\simeq\varprojlim_i\mathcal O/\mathfrak p^i
\]
得到极限 $u\in\mathcal O$。由于 $u$ 模 $\mathfrak p$ 的像是单位，$u\notin\mathfrak p$，
所以 $u\in\mathcal O^\times$。这证明满射。商群
$U^{(0)}/U^{(i)}$ 的核邻域正是 $U^{(i)}$，因此该双射也保持拓扑。
\end{proof}

\begin{exercise}
设 $F$ 为局部数域，剩余特征为 $p$。证明对 $n\ge0$，
\[
(1+\mathfrak p)^{p^n}\subset 1+\mathfrak p^{n+1}.
\]
\end{exercise}

\begin{proof}[解答]
状态：solvable（可解）。

先证一步包含：对任意 $i\ge1$，
\[
(1+\mathfrak p^i)^p\subset 1+\mathfrak p^{i+1}.
\]
若 $x\in\mathfrak p^i$，则
\[
(1+x)^p-1=px+\binom p2x^2+\cdots+x^p.
\]
因为剩余特征为 $p$，有 $p\in\mathfrak p$，故 $px\in\mathfrak p^{i+1}$。中间项
$\binom pkx^k$ 至少含有一个 $p$ 和 $x^k$，也落在 $\mathfrak p^{i+1}$ 中。最后
$x^p\in\mathfrak p^{pi}\subset\mathfrak p^{i+1}$。所以一步包含成立。

从 $i=1$ 开始迭代 $n$ 次，得到
\[
(1+\mathfrak p)^{p^n}\subset1+\mathfrak p^{n+1}.
\]
\end{proof}

\begin{exercise}
设 $S$ 为数域 $F$ 的有限素理想集，$p_0\notin S$。给定 $a_{\mathfrak p}\in F$，证明存在 $x\in F$，使 $\mathfrak p\in S$ 时 $|x-a_{\mathfrak p}|_{\mathfrak p}<\varepsilon$，且对 $\mathfrak p\notin S\cup\{p_0\}$ 有 $|x|_{\mathfrak p}\le1$。
\end{exercise}

\begin{proof}[解答]
状态：partially solvable（部分可解；题面按严格不等号逐字读取时不成立）。

先说明阻塞点。若要求对所有 $\mathfrak p\notin S\cup\{\mathfrak p_0\}$ 都有
$|x|_{\mathfrak p}<1$，则 $x$ 必须属于无限多个素理想。非零数域元素只在有限多个素理想处
有正赋值，所以这会迫使 $x=0$。但 $x=0$ 一般不能同时满足
$|x-a_{\mathfrak p}|_{\mathfrak p}<\varepsilon$。因此严格不等号版本不是普遍真命题。

标准可解版本是把外部条件改成 $|x|_{\mathfrak p}\le1$。证明如下。取整数环中足够高的幂
$\mathfrak p^{N_{\mathfrak p}}$，使同余
\[
x\equiv a_{\mathfrak p}\pmod{\mathfrak p^{N_{\mathfrak p}}}
\]
推出 $|x-a_{\mathfrak p}|_{\mathfrak p}<\varepsilon$。在 Dedekind 环的局部化
$\OF[1/\mathfrak p_0]$ 中，有限个理想
$\mathfrak p^{N_{\mathfrak p}}\OF[1/\mathfrak p_0]$ 两两互素。中国剩余定理给出
$x\in\OF[1/\mathfrak p_0]\subset F$，使所有这些同余同时成立。由于 $x$ 的分母只允许含
$\mathfrak p_0$，对 $\mathfrak p\notin S\cup\{\mathfrak p_0\}$ 有
$v_{\mathfrak p}(x)\ge0$，即 $|x|_{\mathfrak p}\le1$。
\end{proof}

\begin{exercise}
设 $(F,v)$ 为局部数域。证明 $F$ 没有非空连通开集。
\end{exercise}

\begin{proof}[解答]
状态：solvable（可解）。

任取非空开集 $U$ 和 $a\in U$。存在整数 $N$ 使
\[
a+\mathfrak p^N\subset U.
\]
再取 $M>N$。子球 $a+\mathfrak p^M$ 在 $F$ 中既开又闭：开性来自它本身是球；闭性来自
其补集是其他同半径球的并。并且
\[
a+\mathfrak p^M\subsetneq a+\mathfrak p^N\subset U.
\]
在相对拓扑中，$a+\mathfrak p^M$ 是 $U$ 的非空开闭真子集。因此任意非空开集都不连通，
所以 $F$ 没有非空连通开集。
\end{proof}

\begin{exercise}
设 $f(x)$ 在 $x\in\mathfrak p^\mu$ 上收敛，$g(y)=\sum b_ny^n$ 在 $\eta$ 处收敛且 $v(b_n\eta^n)\ge\mu$。证明 $F(x)=f(g(x))$ 在 $x=\eta$ 收敛且 $F(\eta)=f(g(\eta))$。
\end{exercise}

\begin{proof}[解答]
状态：solvable（可解）。

写 $f(X)=\sum_{m\ge0}a_mX^m$。题设说当 $z\in\mathfrak p^\mu$ 时
$\sum_m a_mz^m$ 收敛。又 $g(\eta)=\sum_{n\ge1}b_n\eta^n$ 收敛，并且每一项
$b_n\eta^n$ 都属于 $\mathfrak p^\mu$。强三角不等式给出有限部分和仍在
$\mathfrak p^\mu$ 中；闭性给出极限也在 $\mathfrak p^\mu$ 中，即
\[
g(\eta)\in\mathfrak p^\mu.
\]
所以 $f(g(\eta))$ 有意义。

形式复合为
\[
F(X)=\sum_{m\ge0}a_m\bigl(g(X)\bigr)^m.
\]
在 $X=\eta$ 处，第 $m$ 项为 $a_m(g(\eta))^m$，它正是 $f(g(\eta))$ 的第 $m$ 项。
需要确认把 $g(\eta)$ 作为级数代入不会改变和。对固定 $m$，$g(\eta)^m$ 是 $m$ 重级数；
因为 $b_n\eta^n\to0$，且所有项都在同一理想 $\mathfrak p^\mu$ 中，任意固定 $m$ 的展开
收敛到通常乘积。再对 $m$ 求和时，$a_m(g(\eta))^m\to0$ 来自 $f$ 在
$g(\eta)\in\mathfrak p^\mu$ 处收敛。第 7 题的双重级数换序判别适用，因此复合级数在
$\eta$ 处收敛，并且
\[
F(\eta)=\sum_{m\ge0}a_m(g(\eta))^m=f(g(\eta)).
\]
\end{proof}

\begin{exercise}
设 $(F,v)$ 为局部数域，$\mathcal O$ 为赋值环，$\pi$ 为素元。设 $f(x)\in\mathcal O[x]$，$a\in\mathcal O$，$v(f'(a))=N$，且
\[
f(a)\equiv0\pmod{\pi^{2N+1}}.
\]
证明 Hensel 提升、单位群 $p$ 次幂包含、指数对数双射等结论。
\end{exercise}

\begin{proof}[解答]
状态：solvable（可解）。

(1) 条件等价于
\[
v(f(a))\ge2N+1>2N=2v(f'(a)).
\]
Newton 型 Hensel 引理给出根 $b\in\mathcal O$，且
\[
v(b-a)\ge N+1,
\]
即 $b\equiv a\pmod{\pi^{N+1}}$。

(2) 令 $e=v(p)$。若 $u=1+x\in U^{(e+1)}$，要求 $y^p=u$。考虑
\[
F(Y)=Y^p-u
\]
在 $Y=1$ 处。$F(1)=1-u=-x$ 的赋值至少 $e+1$，$F'(1)=p$ 的赋值为 $e$。对更深的 $U^{(2e+1)}$，有 $v(F(1))\ge2e+1$，满足 Hensel 条件，得 $p$ 次根在 $U^{(e+1)}$ 中。所以
\[
U^{(2e+1)}\subset (U^{(e+1)})^p.
\]

(3) 令 $\kappa=\lfloor e/(p-1)\rfloor+1$。若 $v(x)\ge\kappa$，则
\[
v(x^n/n!)=nv(x)-v(n!)
\]
。Legendre 公式给出
\[
v_p(n!)=\sum_{r\ge1}\left\lfloor\frac n{p^r}\right\rfloor\le\frac{n-1}{p-1}.
\]
由于 $v(p)=e$，所以 $v(n!)\le e(n-1)/(p-1)$。于是
\[
v(x^n/n!)\ge n\kappa-\frac{e(n-1)}{p-1}
=\kappa+(n-1)\left(\kappa-\frac e{p-1}\right)\ge\kappa.
\]

(4) 由 (3)，对 $x\in\mathfrak p^\kappa$，
\[
\exp x=\sum_{n\ge0}\frac{x^n}{n!}
\]
收敛，且 $\exp x-1\in\mathfrak p^\kappa$。对 $u=1+y\in U^{(\kappa)}$，
\[
\log u=\sum_{n\ge1}(-1)^{n+1}\frac{y^n}{n}
\]
也收敛并落在 $\mathfrak p^\kappa$；估计方式与 (3) 相同但更简单，因为 $v(n)\le v(n!)$。
形式恒等式 $\log(\exp x)=x$ 与 $\exp(\log(1+y))=1+y$ 在有理形式幂级数中成立，
而两边在当前收敛范围内逐项合法，所以 $\exp:\mathfrak p^\kappa\to U^{(\kappa)}$
与 $\log$ 互为逆映射。故 $\xi\mapsto\exp\xi$ 是双射。
\end{proof}

\begin{exercise}
设 $F$ 为局部数域。证明 $U^{(1)}$ 的 $p$ 幂包含关系，并构造连续 $\Z_p$-模结构，说明 $U^{(1)}$ 是有限生成 $\Z_p$-模，秩为 $[F:\Q_p]$。
\end{exercise}

\begin{proof}[解答]
状态：solvable（可解）。

(1) 若 $x=1+a$，$a\in\mathfrak p^i$，则
\[
x^p-1=pa+\binom p2a^2+\cdots+a^p.
\]
由于 $p\in\mathfrak p$，第一项在 $\mathfrak p^{i+1}$ 中；中间项含有 $p$ 且含有
$a^2$；最后一项在 $\mathfrak p^{pi}\subset\mathfrak p^{i+1}$ 中。因此
\[
U^{(i)p}\subset U^{(i+1)}.
\]
迭代得 $U^{(i)p^j}\subset U^{(i+j)}$。

(2) 对固定 $i$，若整数 $\mu,\nu$ 满足 $\mu\equiv\nu\pmod{p^i}$，则
$\mu-\nu=p^iq$。由 (1) 得 $x^{p^i}\in U^{(i+1)}\subset U^{(i)}$，所以
$x^\mu$ 与 $x^\nu$ 在 $U^{(1)}/U^{(i)}$ 中同类。于是得到良定作用
\[
\Z/p^i\Z\times U^{(1)}/U^{(i)}\to U^{(1)}/U^{(i)}.
\]
对 $i$ 取逆极限，得到连续作用
\[
\Z_p\times U^{(1)}\to U^{(1)}.
\]

取 $m>\frac e{p-1}$，则对数给出拓扑群同构
\[
\log:U^{(m)}\xrightarrow{\sim}\mathfrak p^m.
\]
右边是 $F$ 的一个 $\Z_p$-格，作为 $\Z_p$-模自由秩为 $[F:\Q_p]$。商
$U^{(1)}/U^{(m)}$ 有限，因此 $U^{(1)}$ 是一个有限生成 $\Z_p$-模；有限商不改变
$\Z_p$-秩，所以
\[
\operatorname{rank}_{\Z_p}U^{(1)}=[F:\Q_p].
\]
\end{proof}
 \chapter{类群}

第三章已经把一个数域 $F$ 在每个素位处的完备化讲清楚了。第四章的核心问题是：
怎样把这些局部域重新组织成一个全局对象，使它既保留所有局部信息，又有足够好的拓扑性质，
可以用紧性、离散性和商群来研究理想类群。加元环和理元群就是为这个目的服务的。

\begin{notation}
设 $F$ 为数域。若两个赋值在 $F$ 上诱导同一个拓扑，就把它们看作同一个素位
（prime place）。有限素位来自非零素理想；无穷素位来自 $F$ 到 $\R$ 或 $\C$ 的嵌入。
对素位 $v$，记 $F_v$ 为 $F$ 在 $v$ 处的完备化。若 $v$ 是有限素位，记
$\mathcal O_v$ 为 $F_v$ 的赋值环，$\mathfrak p_v$ 为其极大理想，$U_v=\mathcal O_v^\times$
为单位群。
\end{notation}

这里要特别分清两件事。第一，有限素位处的 $F_v$ 是第三章意义下的局部数域；
它带有离散赋值、紧开子环 $\mathcal O_v$ 和单位群 $U_v$。第二，无穷素位处的完备化是
$\R$ 或 $\C$，它们也是局部紧域，但本书把“局部数域”这个词保留给 $\Q_p$ 的有限扩张。
这不是语言细节，而是后面定义限制直积时必须区别的地方：有限处有自然的紧开整数环
$\mathcal O_v$，无穷处没有对应的“整数环”，所以无穷处总要全部放进有限集合 $S$ 中。

把所有无穷素位的完备化放在一起，可写作
\[
F_\infty=\prod_{v\mid\infty}F_v\simeq F\otimes_\Q\R.
\]
这个同构的含义是：把 $F$ 的所有实嵌入和复嵌入同时记录下来，就得到一个有限维实代数。
有限素位则对应于 $\operatorname{Spec}\mathcal O_F$ 的闭点。于是“素位”这个词把代数整数环的
素理想和 Archimede 嵌入放进同一套语言里。这样做的好处是，后面所有乘积公式、逼近定理、
理想映射都可以用一个统一的指标 $v$ 来写。

本章有两条主线。第一条是加法主线：构造加元环 $\A_F$，证明 $F$ 在其中离散且
$\A_F/F$ 紧，并由此推出强逼近。第二条是乘法主线：构造理元群 $\A_F^\times$，
定义理元类群，并把它与通常的理想类群、射线类群联系起来。直观地说，加元是
“additive idele”，用于加法拓扑和逼近；理元是“ideal element”的乘法版本，用来记录
每个有限素位处的理想指数。后面类域论中的同调计算正是建立在这些拓扑群上。

\section{加元环}

本章的主题叫“类群”，但开头并不直接从理想类群讲起，而是先把数域 $F$ 的所有完备化
$F_v$ 同时放在一个拓扑环里。这样做的动机是：数域上的算术问题常常可以在每个素位
$v$ 上局部化；可是只看一个 $F_v$ 会丢掉全局条件，只看笛卡尔积 $\prod_vF_v$ 又太大，
不再局部紧。加元环就是两者之间的折中：它允许有限多个地方有分母，同时要求几乎所有
有限地方保持整性。

\textbf{定义（加元环）.}
设 $S_\infty$ 为 $F$ 的全体无穷素位。若 $S$ 是包含 $S_\infty$ 的有限素位集，定义
\[
 \A_F^S=\prod_{v\in S}F_v\times\prod_{v\notin S}\mathcal O_v.
\]
这里第二个乘积只真正限制有限素位，因为无穷素位已经全在 $S$ 中。定义
\[
 \A_F=\bigcup_{S\supset S_\infty,\;S\text{ finite}}\A_F^S.
\]
称 $\A_F$ 为 $F$ 的加元环（adele ring），元素称为加元（adele）。因此，一个加元就是
一个族 $a=(a_v)_v$，其中 $a_v\in F_v$，并且除有限多个有限素位外都有
$a_v\in\mathcal O_v$。

这个定义里的“除有限多个外”不能省。若允许无穷多个有限素位同时出现负赋值，就会失去
自然的紧开部分；若要求所有有限素位都整，又不能容纳一个一般的全局元素在有限多个素理想
处的分母。加元环正好模仿分式理想的行为：一个数或理想只在有限多个素理想处有非平凡指数。

逐分量的加法和乘法是良定的。若 $a\in\A_F^S$、$b\in\A_F^{S'}$，则
$a,b\in\A_F^{S\cup S'}$；对 $v\notin S\cup S'$，有 $a_v,b_v\in\mathcal O_v$，而
$\mathcal O_v$ 是子环，所以 $a_v+b_v$ 与 $a_vb_v$ 仍在 $\mathcal O_v$ 中。因此
\[
 (a+b)_v=a_v+b_v,\qquad (ab)_v=a_vb_v
\]
使 $\A_F$ 成为环。

\textbf{拓扑（限制直积拓扑）.}
每个 $\A_F^S$ 取积拓扑，并规定它是 $\A_F$ 的开子环。等价地，$0$ 的一组邻域基为
\[
 \prod_vN_v,
\]
其中 $N_v$ 是 $F_v$ 中 $0$ 的邻域，并且除有限多个 $v$ 外 $N_v=\mathcal O_v$。
这个拓扑称为 $F_v$ 关于 $\mathcal O_v$ 的限制直积拓扑。它的局部紧性来自两点：
有限多个被放大的分量是局部紧空间，剩下的 $\mathcal O_v$ 是紧空间；有限乘积和任意紧
乘积保持局部紧所需的局部模型。以后所有“加元开集”的计算，实际都是在控制有限多个分量，
其余分量保持在 $\mathcal O_v$ 中。

\textbf{主加元.}
每个 $a\in F$ 可对角嵌入所有 $F_v$。若 $a\ne0$，则分式理想 $(a)$ 的素理想分解只含
有限多个素理想；等价地，除有限多个有限素位外 $v(a)\ge0$，即 $a\in\mathcal O_v$。
因此
\[
 a\longmapsto (a)_v
\]
确实给出单射 $F\hookrightarrow\A_F$。这样的加元称为主加元。这里的单射性来自任意一个
完备化嵌入 $F\to F_v$ 都是单射；若一个全局元素在所有完备化中都为零，它本身就为零。
以后我们常把 $F$ 看成 $\A_F$ 的子环。

\begin{proposition}
设 $E/F$ 为数域的可分扩张，则有自然的线性代数同构和拓扑同构
\[
 \A_F\otimes_FE\simeq \A_E,
\]
并且它把对角嵌入的 $E=F\otimes_FE$ 送到 $\A_E$ 中的主加元子域 $E$。
\end{proposition}

\begin{proof}
取 $E/F$ 的一组 $F$-基 $\xi_1,\ldots,\xi_n$。作为 $F$-向量空间，
\[
 \A_F\otimes_FE=\A_F\xi_1\oplus\cdots\oplus\A_F\xi_n.
\]
在每个素位 $v$ 上，第三章的完备化张量积分解给出
\[
 F_v\otimes_FE\simeq\prod_{w|v}E_w,
\]
其中 $w$ 遍历 $E$ 中位于 $v$ 上方的素位。把这些局部同构对所有 $v$ 拼合起来，就得到
\[
 \prod_v'(F_v\otimes_FE)\simeq
 \prod_v'\prod_{w|v}E_w\simeq \prod_w'E_w=\A_E.
\]
这里撇号表示限制直积。需要检查限制条件相容。对固定的有限素位 $v$，
$\mathcal O_v\xi_1+\cdots+\mathcal O_v\xi_n$ 是 $F_v\otimes_FE$ 中的一个
$\mathcal O_v$-格；在分解
\[
F_v\otimes_FE\simeq\prod_{w|v}E_w
\]
下，它变成 $\prod_{w|v}E_w$ 中的一个开紧 $\mathcal O_v$-子模。另一方面，
$\prod_{w|v}\mathcal O_w$ 也是同一有限维 $F_v$-向量空间中的开紧 $\mathcal O_v$-格。
两个格在局部域上线性等价，因而彼此包含在有限倍放缩之中。并且除有限多个 $v$ 外，
基 $\xi_i$ 可作为整基来使用，这时这两个格相同。于是
\[
 \mathcal O_v\xi_1+\cdots+\mathcal O_v\xi_n
\]
与 $\prod_{w|v}\mathcal O_w$ 给出的限制条件只在有限多个地方不同；改变有限多个地方不改变限制直积的拓扑类型。于是上面的逐点局部同构给出全局拓扑同构。

若 $x\in E$ 写成 $x=\sum_i a_i\xi_i$，其中 $a_i\in F$，则 $\A_F\otimes_FE$ 中的 $x$ 在每个 $v$ 处分量为同一个元素 $x$，经局部分解后正是 $\A_E$ 中的对角主加元。因此主子域也相容。
\end{proof}

命题 4.1 的意义是：加元不是只适用于一个固定数域的孤立构造。数域扩张后，加元环按张量积自然变换，这使后面把局部信息、Galois 作用和上同调放在同一个对象上成为可能。

\begin{theorem}
$F$ 是 $\A_F$ 的离散子环。把 $\A_F$ 看作加法拓扑群，则商群 $\A_F/F$ 是紧群。
\end{theorem}

\begin{proof}
先解释为什么可以把问题化到 $F=\Q$。若 $\xi_1,\ldots,\xi_n$ 是 $F/\Q$ 的一组 $\Q$-基，则命题 4.1 给出加法拓扑群同构
\[
 \A_F\simeq \A_\Q\xi_1\oplus\cdots\oplus \A_\Q\xi_n.
\]
在这个同构下，一个元素 $x\in F$ 唯一写成 $x=\sum_iq_i\xi_i$，其中 $q_i\in\Q$，
所以对角嵌入的 $F$ 对应于 $\Q^n\subset\A_\Q^n$。因此
\[
\A_F/F\simeq(\A_\Q/\Q)^n
\]
作为加法拓扑群同构。若 $\Q$ 在 $\A_\Q$ 中离散，则 $\Q^n$ 在 $\A_\Q^n$ 中离散；
若 $\A_\Q/\Q$ 紧，则有限乘积 $(\A_\Q/\Q)^n$ 紧。所以只需证明 $F=\Q$ 的情形。

证明 $\Q$ 离散。取 $\A_\Q$ 中 $0$ 的邻域
\[
 U=\{x=(x_\infty,x_2,x_3,\ldots): |x_\infty|<1,\ |x_p|_p\le1\text{ 对所有 }p\}.
\]
若 $q\in\Q\cap U$，则 $|q|_p\le1$ 对所有素数 $p$，这说明 $q\in\Z$；同时 $|q|_\infty<1$，故整数 $q$ 只能为 $0$。所以 $U\cap\Q=\{0\}$，即 $\Q$ 离散。

再证明商紧。对素数 $p$，记
\[
 \Q_{(p)}=\left\{\frac{a}{p^m}:a\in\Z,\ m\ge0\right\}.
\]
任意 $x_p\in\Q_p$ 都可写成
\[
 x_p=\xi_p+u_p,\qquad \xi_p\in\Q_{(p)},\quad u_p\in\Z_p,
\]
因为取 $p$-进展开中负幂部分即可。若 $x=(x_v)\in\A_\Q$，只有有限多个 $p$ 满足 $x_p\notin\Z_p$。对这些 $p$ 取上述 $\xi_p$，其他 $p$ 取 $\xi_p=0$，令
\[
 \xi=\sum_p\xi_p\in\Q.
\]
对每个 $p$，把 $\xi$ 写作 $\xi=\xi_p+\eta_p$，其中 $\eta_p=\sum_{q\ne p}\xi_q$。由于 $\eta_p$ 的分母不含 $p$，故 $\eta_p\in\Z_p$。于是
\[
 x_p-\xi=(x_p-\xi_p)-\eta_p\in\Z_p.
\]
这说明每个加元模去某个有理数后，其所有有限分量都落在 $\Z_p$ 中。

最后处理实分量。任意实数 $r$ 可写成 $r=n+s$，其中 $n\in\Z$ 且 $-\frac12\le s\le\frac12$。因此
\[
 C=\left[-\frac12,\frac12\right]\times\prod_p\Z_p
\]
是 $\A_\Q$ 中的紧集，并且 $\A_\Q=\Q+C$。紧集 $C$ 在商群 $\A_\Q/\Q$ 中的像覆盖全商，所以 $\A_\Q/\Q$ 紧。
\end{proof}

从学习角度看，定理 4.2 是加元环的第一个核心事实。它说 $F$ 在所有局部完备化的限制直积中不是稠密的，而是一个离散格；并且余商是紧的。这和几何数论中“格在实向量空间中离散且商紧”的图像完全平行，只是这里实向量空间被加元环替代。

\begin{lemma}
存在只依赖于 $F$ 的常数 $C>0$，使得若 $a=(a_v)\in\A_F$ 满足
\[
 \prod_v |a_v|_v>C,
\]
则存在 $\xi\in F^\times$，使得对所有素位 $v$ 都有
\[
 |\xi|_v\le |a_v|_v.
\]
\end{lemma}

\begin{proof}
定理 4.2 已给出 $\A_F/F$ 紧。取 Haar 测度 $\mu$，并取一个标准紧邻域
\[
 B=\{x=(x_v)\in\A_F: |x_v|_v\le 1/10\text{ 若 }v|\infty,\ |x_v|_v\le1\text{ 若 }v<\infty\}.
\]
它的测度正且有限。设 $\dot\mu(\A_F/F)$ 为商群的 Haar 测度。选择
\[
 C>\frac{\dot\mu(\A_F/F)}{\mu(B)}.
\]
对给定 $a=(a_v)$，令
\[
 T=\{t=(t_v): |t_v|_v\le |a_v|_v/10\text{ 若 }v|\infty,\ |t_v|_v\le |a_v|_v\text{ 若 }v<\infty\}.
\]
Haar 测度在逐分量放缩下按因子 $\prod_v|a_v|_v$ 缩放，所以当 $\prod_v|a_v|_v>C$ 时有
\[
 \mu(T)>\dot\mu(\A_F/F).
\]
将 $T$ 投到紧商 $\A_F/F$。若投影是单射，则投影测度等于 $\mu(T)$，不可能大于整个商空间的测度。因此存在不同的 $t',t''\in T$，使得 $t'-t''=\xi\in F$。

由于 $t'\ne t''$，有 $\xi\ne0$。在非 Archimede 素位，超距不等式给出
\[
 |\xi|_v=|t'_v-t''_v|_v\le \max\{|t'_v|_v,|t''_v|_v\}\le |a_v|_v.
\]
在 Archimede 素位，两点都落在半径 $|a_v|_v/10$ 的区间或圆盘中，因此差的绝对值至多 $|a_v|_v/5$，特别不超过 $|a_v|_v$。于是得到所需的 $\xi$。
\end{proof}

\begin{lemma}
有以下两个逼近结论。

\begin{enumerate}[label=(\arabic*)]
\item $\A_F$ 中存在集合 $W=\{x=(x_v): |x_v|_v\le\delta_v\}$，其中除有限多个 $v$ 外 $\delta_v=1$，使得任意 $\alpha\in\A_F$ 都能写成
\[
 \alpha=\theta+\gamma,\qquad \theta\in W,\quad \gamma\in F.
\]
\item 固定一个素位 $v_0$。若对所有 $v\ne v_0$ 给定 $\delta_v>0$，且除有限多个 $v$ 外 $\delta_v=1$，则存在 $\beta\in F^\times$，使得
\[
 |\beta|_v\le\delta_v\qquad(v\ne v_0).
\]
\end{enumerate}
\end{lemma}

\begin{proof}
(1) 定理 4.2 说明 $\A_F/F$ 紧，因此存在紧集 $K\subset\A_F$ 的像覆盖整个商。紧集在
限制直积中有一个重要后果：每个坐标投影 $K_v\subset F_v$ 都有界，并且除有限多个有限素位外
$K_v\subset\mathcal O_v$。于是可以选择实数或赋值半径 $\delta_v>0$，使
$K_v\subset\{x_v:|x_v|_v\le\delta_v\}$，且除有限多个 $v$ 外 $\delta_v=1$。令
\[
\ W=\{x=(x_v): |x_v|_v\le\delta_v\}.
\]
因为 $K\subset W$，而 $K$ 的像已经覆盖 $\A_F/F$，任意 $\alpha\in\A_F$ 都能写成
$\alpha=\theta+\gamma$，其中 $\theta\in W$、$\gamma\in F$。

(2) 对每个 $v\ne v_0$ 取 $\alpha_v\in F_v$，满足 $0<|\alpha_v|_v\le\delta_v$，并且当 $\delta_v=1$ 时取 $|\alpha_v|_v=1$。接着选择 $v_0$ 处分量 $\alpha_{v_0}$ 足够大，使得
\[
 \prod_v|\alpha_v|_v>C,
\]
其中 $C$ 是引理 4.3 中的常数。由引理 4.3，存在 $\beta\in F^\times$，使 $|\beta|_v\le|\alpha_v|_v$ 对所有 $v$ 成立。于是对 $v\ne v_0$ 有 $|\beta|_v\le\delta_v$。
\end{proof}

\begin{theorem}
固定素位 $v_0$。记
\[
 \A_F^{v_0}=\prod_{v\ne v_0}'F_v
\]
为省去 $v_0$ 后、关于 $\mathcal O_v$ 的限制直积。则 $F$ 在 $\A_F^{v_0}$ 中稠密。
\end{theorem}

\begin{proof}
稠密性的邻域表述如下：给定有限素位集 $S$，其中不含 $v_0$，给定 $\alpha_v\in F_v$ 和误差 $\varepsilon_v>0$，要求找 $\beta\in F$，使
\[
 |\beta-\alpha_v|_v<\varepsilon_v\quad(v\in S),
\]
并且对 $v\notin S\cup\{v_0\}$ 有 $\beta\in\mathcal O_v$。

由引理 4.4(1)，取坐标盒 $W=\{|x_v|_v\le\delta_v\}$，使 $\A_F=F+W$。由引理 4.4(2)，存在 $\lambda\in F^\times$，满足
\[
 |\lambda|_v\delta_v<\varepsilon_v\quad(v\in S),\qquad
 |\lambda|_v\delta_v\le1\quad(v\notin S\cup\{v_0\}).
\]
在 $\A_F$ 中取一个元素 $\alpha$，使它在 $v\in S$ 的分量为给定的 $\alpha_v$，其他分量取 $0$。因为 $\A_F=F+\lambda W$，存在 $\beta\in F$ 与 $\theta\in W$，使
\[
 \alpha=\beta+\lambda\theta.
\]
于是 $v\in S$ 时
\[
 |\beta-\alpha_v|_v=|\lambda\theta_v|_v\le|\lambda|_v\delta_v<\varepsilon_v.
\]
若 $v\notin S\cup\{v_0\}$，则 $\alpha_v=0$，故
\[
 |\beta|_v=|\lambda\theta_v|_v\le|\lambda|_v\delta_v\le1,
\]
即 $\beta\in\mathcal O_v$。这正是所需的稠密性。
\end{proof}

强逼近定理可以这样记：如果允许在一个地方 $v_0$ 不受控制，那么全局元素 $F$ 可以同时满足任意有限多个局部近似条件，并在其他有限地方保持整性。这是后面处理 $S$-整数、射线同余和理想类的基本工具。

\section{理元群}

加元环是加法和环结构的限制直积。类域论真正使用更多的是乘法版本，也就是理元群。直观上，理元把每个局部域的非零元素放在一起，但仍要求除有限多个地方外是局部单位。
如果直接取 $\prod_vF_v^\times$，会得到一个过大的乘法群：无穷多个有限素位都可以带有
非零理想指数，和分式理想只含有限多个素因子的事实不相符。理元群把这个问题修正为
“几乎处处为单位”。这也是 Chevalley 在类域论中引入理元的关键点：它既是局部紧拓扑群，
又能同时承载所有局部乘法群 $F_v^\times$ 的信息。

\begin{definition}
设 $S$ 是包含 $S_\infty$ 的有限素位集。定义
\[
 \A_F^{\times S}=\prod_{v\in S}F_v^\times\times\prod_{v\notin S}\mathcal O_v^\times.
\]
定义 $F$ 的理元群为
\[
 \A_F^\times=\bigcup_{S\supset S_\infty,\;S\text{ finite}}\A_F^{\times S}.
\]
其元素称为理元。也就是说，理元是族 $a=(a_v)_v$，满足 $a_v\in F_v^\times$，并且除有限多个有限素位外 $a_v\in\mathcal O_v^\times$。
\end{definition}

拓扑仍用限制直积定义：每个 $\A_F^{\times S}$ 取积拓扑，并规定它们是开子群。单位元 $1$ 的邻域基为
\[
 \prod_vN_v,
\]
其中 $N_v$ 是 $F_v^\times$ 中 $1$ 的邻域，且除有限多个 $v$ 外 $N_v=\mathcal O_v^\times$。这样 $\A_F^\times$ 是局部紧拓扑群。

注意一个细节：$\A_F^\times$ 作为集合等于加元环 $\A_F$ 的可逆元，但它的拓扑不是简单地从 $\A_F$ 继承的子空间拓扑。原因是一般拓扑环中，取逆映射在子空间拓扑下不一定连续；理元群必须让取逆连续。

\begin{proposition}
有以下结论。
\begin{enumerate}[label=(\arabic*)]
\item 加元环 $\A_F$ 的可逆元集合等于理元群 $\A_F^\times$。
\item 映射
\[
 \iota:\A_F^\times\longrightarrow \A_F\times\A_F,\qquad
 x\longmapsto (x,x^{-1})
\]
是单射。
\item 若 $\iota(\A_F^\times)$ 取 $\A_F\times\A_F$ 的子空间拓扑，而 $\A_F^\times$ 取限制直积拓扑，则 $\iota$ 是拓扑同构。
\end{enumerate}
\end{proposition}

\begin{proof}
(1) 若 $a=(a_v)\in\A_F^\times$，则每个 $a_v\ne0$，且除有限多个有限素位外 $a_v\in\mathcal O_v^\times$。于是 $a^{-1}=(a_v^{-1})$ 仍满足加元条件，所以 $a$ 是 $\A_F$ 的可逆元。

反过来，若 $a$ 是 $\A_F$ 的可逆元，则存在 $b=(b_v)\in\A_F$ 使 $ab=1$。逐分量看，$a_vb_v=1$，所以 $a_v\in F_v^\times$。又因为 $a,b$ 都是加元，除有限多个有限素位外 $a_v,b_v\in\mathcal O_v$；在这些地方 $a_vb_v=1$，故 $a_v,b_v$ 都是 $\mathcal O_v^\times$ 中的单位。因此 $a\in\A_F^\times$。

(2) 单射来自第一分量：若 $(x,x^{-1})=(y,y^{-1})$，则 $x=y$。

(3) 对固定有限集 $S$，$\A_F^{\times S}$ 中的元素及其逆元在 $\A_F^S$ 中都有受控分量。映射 $x\mapsto (x,x^{-1})$ 在每个局部域 $F_v^\times$ 上是连续的，在 $\mathcal O_v^\times$ 上也是连续的。有限多个非单位地方取普通积拓扑，几乎所有地方取紧单位群拓扑，于是限制直积拓扑恰好等于 $\iota$ 诱导的拓扑。反向连续性由投影到第一分量得到。
更具体地说，$\iota(\A_F^\times)$ 在 $\A_F\times\A_F$ 中的一个基本邻域同时要求
$x_v$ 接近给定值并要求 $x_v^{-1}$ 接近给定逆元。有限多个指定分量上的这些要求正是
$F_v^\times$ 的通常拓扑；未指定的有限素位上，$x_v$ 与 $x_v^{-1}$ 都在
$\mathcal O_v$ 中等价于 $x_v\in\mathcal O_v^\times$。因此诱导拓扑的基本开集和
限制直积拓扑的基本开集逐项一致。
\end{proof}

\begin{proposition}
对任意 $a\in F^\times$，有
\[
 \prod_v |a|_v=1.
\]
\end{proposition}

\begin{proof}
先看 $F=\Q$。若
\[
 a=\pm\prod_pp^{n_p}
\]
是有理数的素因子分解，其中只有有限多个 $n_p$ 非零，则
\[
 |a|_\infty=\prod_pp^{n_p},\qquad |a|_p=p^{-n_p}.
\]
因此
\[
 |a|_\infty\prod_p|a|_p
 =\prod_pp^{n_p}\prod_pp^{-n_p}=1.
\]

一般数域情形，把有限素位按其下方的有理素数 $p$ 分组。局部绝对值的规范化满足
\[
 \prod_{v|p}|a|_v=|\Nm_{F/\Q}(a)|_p,
\]
这是因为 $p$ 上方的所有完备化范数相乘，正好给出全局范数在 $\Q_p$ 中的绝对值；
有限扩张中局部范数的乘积等于全局范数。无穷素位的乘积给出
$|\Nm_{F/\Q}(a)|_\infty$。于是
\[
 \prod_v|a|_v
 =|\Nm_{F/\Q}(a)|_\infty\prod_p|\Nm_{F/\Q}(a)|_p=1,
\]
最后一步是 $\Q$ 上的乘积公式应用于非零有理数 $\Nm_{F/\Q}(a)$。
\end{proof}

把 $a\in F^\times$ 对角嵌入 $\A_F^\times$，得到主理元。乘积公式说明主理元的所有局部大小相乘总是 $1$。因此理元范数
\[
 |x|_\A=\prod_v|x_v|_v
\]
在主理元上恒等于 $1$。

\begin{proposition}
$F^\times$ 是 $\A_F^\times$ 的离散闭子群。
\end{proposition}

\begin{proof}
先证明离散。取 $1$ 的邻域
\[
 U=\{x=(x_v)\in\A_F^\times:
 |x_v|_v=1\text{ 若 }v<\infty,\quad |x_v-1|_v<1\text{ 若 }v|\infty\}.
\]
若 $x\in F^\times\cap U$ 且 $x\ne1$，则对有限素位有 $|x|_v=1$，所以
\[
 |x-1|_v\le\max\{|x|_v,1\}=1.
\]
对无穷素位有 $|x-1|_v<1$。将乘积公式用于 $x-1\in F^\times$，得到
\[
 1=\prod_v|x-1|_v
 =\prod_{v<\infty}|x-1|_v\prod_{v|\infty}|x-1|_v<1,
\]
矛盾。因此 $F^\times\cap U=\{1\}$，$F^\times$ 离散。

再证闭性。设 $y\in\A_F^\times$ 不属于 $F^\times$。若每个 $F^\times$ 的邻域都与
$y$ 相交，则 $1$ 的每个邻域都与 $y^{-1}F^\times$ 相交。取上面构造的邻域 $U$，并取
一个对称邻域 $V$，满足 $VV^{-1}\subset U$。若存在两个不同的主理元
$a,b\in F^\times\cap yV$，则 $a^{-1}b\in F^\times\cap V^{-1}V\subset F^\times\cap U$，
所以 $a^{-1}b=1$，即 $a=b$。因此 $yV$ 中至多有一个主理元。若 $y$ 在闭包中，则
可以取 $a\in F^\times\cap yV$；再把 $V$ 缩小而仍包含 $y^{-1}a$，闭包性会迫使
同一个 $a$ 落在所有足够小的 $yV$ 中，故 $y^{-1}a$ 落在所有 $1$ 的邻域交中。理元群是
Hausdorff 群，邻域交为 $\{1\}$，所以 $y=a\in F^\times$，矛盾。故补集开，
$F^\times$ 闭。
\end{proof}

定义
\[
 \A_F^1=\{x\in\A_F^\times: |x|_\A=1\}.
\]
它是理元群中“总大小为 $1$”的部分，主理元 $F^\times$ 全部包含在其中。

\begin{lemma}
$\A_F^1$ 作为 $\A_F$ 的子集是闭子集；并且 $\A_F^\times$ 与 $\A_F$ 在 $\A_F^1$ 上诱导相同的拓扑。
\end{lemma}

\begin{proof}
先证闭性。若 $a\in\A_F\setminus\A_F^1$，分三种情况。

若某个分量 $a_{v_0}=0$，取有限集 $S$ 包含 $v_0$、所有无穷素位以及所有满足
 $|a_v|_v>1$ 的有限素位。对 $v\in S$ 选 $a_v$ 的足够小邻域，使所有
$x_v$ 的有限乘积绝对值小于 $1$；这可以做到，因为 $v_0$ 分量可以被限制到任意小。
在 $S$ 外要求 $x_v\in\mathcal O_v$，则外部乘积至多为 $1$。所以该加元邻域中的
元素不可能有总范数 $1$，它与 $\A_F^1$ 不交。

若 $a\in\A_F^\times$ 且 $|a|_\A<1$，取有限集 $S$ 包含所有满足 $|a_v|_v\ne1$ 的地方以及所有无穷素位。令 $x$ 在 $S$ 上足够接近 $a$，在 $S$ 外满足 $|x_v|_v\le1$。由于有限多个连续函数的乘积连续，可使
\[
 \prod_{v\in S}|x_v|_v<1,
\]
而 $S$ 外乘积不超过 $1$，故 $|x|_\A<1$。这个邻域不遇到 $\A_F^1$。

若 $a\in\A_F^\times$ 且 $|a|_\A>1$，记 $c=|a|_\A$。取有限集 $S$ 包含所有满足
$|a_v|_v\ne1$ 的地方、所有无穷素位，以及所有剩余域大小不超过 $2c$ 的有限素位。
在有限集 $S$ 上缩小邻域，使
\[
\qquad 1<\prod_{v\in S}|x_v|_v<2c .
\]
在 $S$ 外只要求 $x_v\in\mathcal O_v$。若所有 $S$ 外分量都是单位，则外部乘积为
$1$，总范数大于 $1$。若某个 $S$ 外分量不是单位，则因该素位不在 $S$，它的剩余域大小
大于 $2c$，故 $|x_v|_v\le(2c)^{-1}$；其他外部分量绝对值至多 $1$，于是总范数小于
$2c\cdot(2c)^{-1}=1$。两种情形都不可能得到总范数 $1$。于是 $\A_F^1$ 闭。

再比较拓扑。理元拓扑比加元子空间拓扑强，因为包含映射
$\A_F^\times\to\A_F$ 连续。反过来，取 $a\in\A_F^1$ 的理元邻域。它由有限集
$S$ 上的条件 $x_v\in N_v$ 和 $S$ 外条件 $x_v\in\mathcal O_v^\times$ 给出。
由于 $a\in\A_F^1$，可把 $S$ 扩大到包含所有 $|a_v|_v\ne1$ 的地方。再把 $S$ 扩大，
使 $S$ 外任一非单位的绝对值都小于一个预先给定的数 $\rho<1$。把 $S$ 上的加元邻域
缩小，使 $\prod_{v\in S}|x_v|_v$ 始终落在
$(1,\rho^{-1})$ 或其对应的足够窄区间内，并且满足原理元邻域在 $S$ 上的要求。若
$x\in\A_F^1$ 且 $S$ 外只有 $x_v\in\mathcal O_v$，一旦某个 $S$ 外分量不是单位，
外部乘积至多为 $\rho$，总范数便不能等于 $1$；因此 $S$ 外自动有
$x_v\in\mathcal O_v^\times$。这说明每个理元邻域在 $\A_F^1$ 上都含有一个加元子空间
邻域，两种诱导拓扑一致。
\end{proof}

\begin{theorem}
商群 $\A_F^1/F^\times$ 是紧拓扑群。
\end{theorem}

\begin{proof}
由引理 4.10，可在加元环的拓扑中寻找紧集代表。取引理 4.3 的常数 $C$，再取一个加元 $a=(a_v)$，使
\[
 \prod_v|a_v|_v>C.
\]
例如只在一个无穷素位把绝对值取足够大，其他地方取 $1$。

令
\[
 W=\{x=(x_v)\in\A_F: |x_v|_v\le |a_v|_v\text{ 对所有 }v\}.
\]
这是 $\A_F$ 中的紧集：有限多个地方是有界闭集，几乎所有有限地方是 $\mathcal O_v$。

取任意 $b=(b_v)\in\A_F^1$。对加元 $b^{-1}a$，有
\[
 \prod_v|b_v^{-1}a_v|_v
 =|b|_\A^{-1}\prod_v|a_v|_v
 =\prod_v|a_v|_v>C.
\]
由引理 4.3，存在 $\eta\in F^\times$，使
\[
 |\eta|_v\le |b_v^{-1}a_v|_v.
\]
于是
\[
 |\eta b_v|_v\le |a_v|_v,
\]
即 $\eta b\in W\cap\A_F^1$。这说明每个 $\A_F^1/F^\times$ 的陪集都有代表落在紧集 $W\cap\A_F^1$ 中。由于 $\A_F^1$ 闭，$W\cap\A_F^1$ 紧；它在商中的像覆盖全商，所以 $\A_F^1/F^\times$ 紧。
\end{proof}

\section{理元类群}

理元类群是理元群模去主理元所得的商群。它是类域论的主要对象，因为它把“所有局部乘法数据”按“来自一个全局元素的数据”取商。直观地说，理元记录每个地方的局部倍数；主理元表示这些倍数本来就来自同一个全局数。商掉主理元后，留下的就是局部数据不能全局化的障碍。

设 $F$ 的模写作
\[
 \mathfrak m=\mathfrak m_\infty\mathfrak m_0,\qquad
 \mathfrak m_0=\prod_{v<\infty}\mathfrak p_v^{n_v}.
\]
这里 $n_v\ge0$ 且除有限多个外为 $0$；$\mathfrak m_\infty$ 由部分实素位组成。若 $v$ 是有限素位，定义
\[
 U_v^{(0)}=\mathcal O_v^\times,\qquad
 U_v^{(n)}=1+\mathfrak p_v^n\quad(n>0).
\]
若 $v$ 是实素位并且出现在 $\mathfrak m_\infty$ 中，取 $U_v^{(1)}=\R_+^\times$；若不出现，取 $U_v^{(0)}=\R^\times$。复素位处取 $U_v^{(0)}=\C^\times$。
实素位处出现 $\R_+^\times$ 的原因是射线类群不仅记录有限素位同余，还记录实嵌入下的
符号。条件“在指定实嵌入下为正”不能用有限素理想同余表达，所以必须把它放在
$\mathfrak m_\infty$ 中。复素位没有正负之分，因此不产生额外条件。

定义模 $\mathfrak m$ 的同余子群
\[
 \A_F^{\times\mathfrak m}=\prod_v U_v^{(n_v)},
\]
其中无穷处按上面的正性条件解释。还记
\[
 \widehat F_{\mathfrak m_0}=\prod_{v<\infty}U_v^{(n_v)},\qquad
 F_\infty^\times=\prod_{v|\infty}F_v^\times.
\]
于是
\[
 \A_F^{\times\mathfrak m_0}
 =F_\infty^\times\times\widehat F_{\mathfrak m_0}.
\]

\begin{definition}
定义理元类群
\[
 C_F=\A_F^\times/F^\times.
\]
同时定义
\[
 C(\mathfrak m_0)=\A_F^\times/\widehat F_{\mathfrak m_0}F^\times,
\]
\[
 C_F^{\mathfrak m}=\A_F^{\times\mathfrak m}F^\times/F^\times,
\qquad
 C_F/C_F^{\mathfrak m}
 =\A_F^\times/\A_F^{\times\mathfrak m}F^\times.
\]
\end{definition}

这些定义初看符号很多，核心只有一句：不同的开子群给出不同精细程度的类群。$C_F$ 是最基本的理元类群；加入模 $\mathfrak m$ 后，有限处的同余条件和实处的正性条件共同形成射线类结构。后面的类域论会把这些商群与 Abel 扩张的 Galois 群联系起来。
读这些商群时可以按“先限制，再取商”理解。$C_F^{\mathfrak m}$ 是理元类群中能用一个满足
模 $\mathfrak m$ 局部条件的理元代表表示的子群；商 $C_F/C_F^{\mathfrak m}$ 则把理元类按
模 $\mathfrak m$ 的同余精度区分。$C(\mathfrak m_0)$ 只除去有限处的同余单位，不处理无穷处
正性，因此它是连接理元语言和传统模 $\mathfrak m_0$ 射线类群的中间对象。

\section{理想}

本节把理元语言和第一章的理想语言连接起来。有限素位 $v$ 对应 $\mathcal O_F$ 的非零素理想 $\mathfrak p_v$。若 $a=(a_v)\in\A_F^\times$，在有限处定义
\[
 i(a)=\prod_{v<\infty}\mathfrak p_v^{v(a_v)}.
\]
因为理元除有限多个地方外是单位，只有有限多个指数非零，所以这是一个分式理想。于是得到群同态
\[
 i:\A_F^\times\longrightarrow I_F,
\]
其中 $I_F$ 是 $F$ 的分式理想群。

同态 $i$ 的核正是
\[
 \A_F^{\times S_\infty}
 =\prod_{v|\infty}F_v^\times\times\prod_{v<\infty}\mathcal O_v^\times.
\]
另一方面，主理元 $a\in F^\times$ 的像是主分式理想 $(a)$。因此 $i$ 诱导正合列
\[
 1\longrightarrow
 \A_F^{\times S_\infty}F^\times/F^\times
 \longrightarrow C_F
 \longrightarrow \Cl_F
 \longrightarrow1.
\]
这个正合列说明理想类群是理元类群的一个离散影子：理元类群含有连续的局部单位信息，映到理想类群时只保留有限素位的赋值指数。

\begin{proposition}
当有限素位集 $S$ 充分大并包含 $S_\infty$ 时，有
\[
 \A_F^\times=\A_F^{\times S}F^\times.
\]
\end{proposition}

\begin{proof}
理想类群 $\Cl_F$ 有限，取每个理想类的一个代表分式理想
\[
 \mathfrak a_1,\ldots,\mathfrak a_h.
\]
令 $S$ 包含 $S_\infty$，并包含所有出现在这些 $\mathfrak a_i$ 的素理想分解中的有限素位。

任取理元 $a=(a_v)\in\A_F^\times$。其理想 $i(a)$ 属于某个理想类，所以存在 $b\in F^\times$ 与某个 $i$，使
\[
 i(a)=\mathfrak a_i(b).
\]
等价地，
\[
 i(ab^{-1})=\mathfrak a_i.
\]
由于 $\mathfrak a_i$ 只在 $S$ 中的有限素位有非零指数，对 $v\notin S$ 有
\[
 v((ab^{-1})_v)=0,
\]
即 $(ab^{-1})_v\in\mathcal O_v^\times$。因此 $ab^{-1}\in\A_F^{\times S}$，从而
\[
 a\in \A_F^{\times S}F^\times.
\]
\end{proof}

现在整理射线类群的符号。设
\[
 \mathfrak m=\mathfrak m_\infty\mathfrak m_0,\qquad
 \mathfrak m_0=\prod_{v<\infty}\mathfrak p_v^{n_v}.
\]
记 $I_F^{\mathfrak m}$ 为由所有不整除 $\mathfrak m_0$ 的素理想生成的分式理想群。模 $\mathfrak m$ 的射线主理想群 $S_F^{\mathfrak m}$ 由主理想 $(a)$ 组成，其中 $a\in F^\times$ 满足
\[
 a\equiv1\pmod{\mathfrak m_0},
\]
并且在 $\mathfrak m_\infty$ 中指定的所有实嵌入下为正。若不要求无穷处正性，只记相应群为 $S_F^{\mathfrak m_0}$。
为写第三个正合列，记
\[
F_{\mathfrak m_0}^\times
=\{x\in F^\times: x\in U_v^{(n_v)}\text{ 对所有 }v<\infty\}.
\]
这是“有限处模 $\mathfrak m_0$ 同余于 $1$”的全局元素群。通过无穷嵌入，它自然看作
$F_\infty^\times$ 的子群。

\begin{proposition}
同态 $i:\A_F^\times\to I_F$ 给出以下正合序列：
\[
 1\to \A_F^{\times S_\infty}F^\times/F^\times
 \to C_F\to \Cl_F\to1,
\]
\[
 1\to C_F^{\mathfrak m}\to C_F\to \Cl_F^{\mathfrak m}\to1,
\]
\[
 1\to
 F_\infty^\times/F_{\mathfrak m_0}^\times
 \to C(\mathfrak m_0)
 \to I_F^{\mathfrak m}/S_F^{\mathfrak m_0}
 \to1.
\]
\end{proposition}

\begin{proof}
第一列已经在本节开头说明：$i$ 的核是 $\A_F^{\times S_\infty}$，而主理元映到主理想。
满射性来自每个分式理想都可写成有限素理想幂的乘积；给定
$\prod_{v<\infty}\mathfrak p_v^{n_v}$，取理元在这些 $v$ 处分量为素元的 $n_v$ 次幂，
其他有限处取单位，无穷处取 $1$，便映到该分式理想。模去主理元后，余核正是
$\Cl_F$。

第二列中，商映射由 $i$ 诱导。若理元 $a$ 在 $\Cl_F^{\mathfrak m}$ 中映为平凡，表示 $i(a)$ 是一个射线主理想 $(x)$，其中 $x$ 满足模 $\mathfrak m$ 的同余与正性条件。于是
\[
 i(ax^{-1})=1,
\]
并且 $ax^{-1}$ 在模 $\mathfrak m$ 要求的各局部单位子群中，所以 $a$ 在 $C_F$ 中落入
$C_F^{\mathfrak m}$。反过来，若 $a$ 的类属于 $C_F^{\mathfrak m}$，则可写成
$a=ux$，其中 $u\in\A_F^{\times\mathfrak m}$、$x\in F^\times$。有限处同余条件说明
$i(u)$ 在 $I_F^{\mathfrak m}$ 中不给出非平凡射线类，而 $i(x)=(x)$ 是满足模
$\mathfrak m$ 条件的射线主理想；所以 $a$ 映到射线类群单位元。因此核正是
$C_F^{\mathfrak m}$。满射性与第一列相同：任何与 $\mathfrak m_0$ 互素的分式理想都可由
有限处理元的赋值指数实现。

第三列只保留有限同余子群 $\widehat F_{\mathfrak m_0}$。映射
\[
C(\mathfrak m_0)=\A_F^\times/\widehat F_{\mathfrak m_0}F^\times
\longrightarrow I_F^{\mathfrak m}/S_F^{\mathfrak m_0}
\]
仍由 $i$ 给出。有限处落在 $\widehat F_{\mathfrak m_0}$ 的理元在所有允许的有限素位处
赋值为 $0$，不会改变 $I_F^{\mathfrak m}$ 中的类；主理元会把理想乘上一个模
$\mathfrak m_0$ 的射线主理想。因此映射良定。

求核。若 $a$ 映到单位元，则存在 $x\in F^\times$ 使 $i(a)=(x)$，并且 $x$ 在有限处满足
模 $\mathfrak m_0$ 的同余条件。于是 $ax^{-1}$ 的有限处分量都落在
$\widehat F_{\mathfrak m_0}$ 中，故在 $C(\mathfrak m_0)$ 中，$a$ 的类可由一个只剩无穷
处分量的理元表示。反过来，纯无穷理元没有有限赋值贡献，映到理想群中为单位类。
两个无穷理元表示同一个核元素，当且仅当它们相差一个全局元素 $x\in F^\times$，且这个
$x$ 的有限处分量属于 $\widehat F_{\mathfrak m_0}$，也就是 $x\in F_{\mathfrak m_0}^\times$。
所以核为
\[
 F_\infty^\times/F_{\mathfrak m_0}^\times.
\]
满射性仍由给定理想的有限赋值指数构造理元得到。这给出正合性。
\end{proof}

\section*{习题详解}
\addcontentsline{toc}{section}{习题详解}

\begin{exercise}
对 $z=(z_v)\in\A_\Q^\times$，设
\[
 r(z)=\operatorname{sgn}(z_\infty)\prod_p|z_p|_p^{-1}.
\]
证明：
\[
 \ker r=\R_+^\times\times\prod_p\Z_p^\times,\qquad
 \A_\Q^\times=\Q^\times\ker r,\qquad
 \A_\Q^\times=\Q^\times\times\R_+^\times\times\prod_p\Z_p^\times.
\]
\end{exercise}

\begin{proof}[解答]
状态：solvable（可解）。

关键是把有限处赋值和无穷处符号分开。

先求核。因为 $z\in\A_\Q^\times$，除有限多个 $p$ 外 $z_p\in\Z_p^\times$，所以乘积是有限乘积。写
\[
 |z_p|_p=p^{-n_p},\qquad n_p\in\Z,
\]
且只有有限多个 $n_p\ne0$。于是
\[
 \prod_p|z_p|_p^{-1}=\prod_pp^{n_p}
\]
是正有理数。若 $r(z)=1$，则 $\operatorname{sgn}(z_\infty)=1$ 且 $\prod_pp^{n_p}=1$。素因子分解唯一性给出所有 $n_p=0$，所以 $z_p\in\Z_p^\times$，且 $z_\infty>0$。因此
\[
 \ker r=\R_+^\times\times\prod_p\Z_p^\times.
\]

再证 $\A_\Q^\times=\Q^\times\ker r$。对 $q\in\Q^\times$ 的主理元，乘积公式给出
\[
 r(q)=\operatorname{sgn}(q)\prod_p|q|_p^{-1}
 =\operatorname{sgn}(q)|q|_\infty=q.
\]
对任意 $z\in\A_\Q^\times$，令 $q=r(z)\in\Q^\times$。则
\[
 r(q^{-1}z)=r(q)^{-1}r(z)=1,
\]
所以 $q^{-1}z\in\ker r$，即 $z=q(q^{-1}z)\in\Q^\times\ker r$。

最后证直积性。若 $q\in\Q^\times\cap\ker r$，则 $q=r(q)=1$。因此 $\Q^\times$ 与 $\ker r$ 交平凡，乘积表示唯一。代入核的描述，得到
\[
 \A_\Q^\times\simeq
 \Q^\times\times\R_+^\times\times\prod_p\Z_p^\times.
\]
\end{proof}

\begin{exercise}
证明：
\begin{enumerate}[label=(\arabic*)]
\item $F^\times$ 是 $\A_F^1$ 的离散子群。
\item 有拓扑群同构
\[
 \A_F^\times/F^\times\simeq \A_F^1/F^\times\times\R_+^\times.
\]
\end{enumerate}
\end{exercise}

\begin{proof}[解答]
状态：solvable（可解）。

(1) 命题 4.9 已证明 $F^\times$ 在 $\A_F^\times$ 中离散。由于乘积公式给出 $F^\times\subset\A_F^1$，而 $\A_F^1$ 取子空间拓扑，$F^\times$ 在 $\A_F^1$ 中仍离散。

(2) 理元范数
\[
 |\cdot|_\A:\A_F^\times\to\R_+^\times
\]
是连续群同态，核为 $\A_F^1$。它在主理元 $F^\times$ 上为 $1$，因此诱导
\[
 C_F=\A_F^\times/F^\times\to\R_+^\times.
\]
这个同态满射：取一个无穷素位 $v_0$。若 $v_0$ 为实素位，对 $t>0$ 取理元在 $v_0$ 分量为 $t$、其他分量为 $1$；若 $v_0$ 为复素位，取分量为正实数 $\sqrt t$，在规范化绝对值下其局部绝对值为 $t$。这样得到范数为 $t$ 的理元。

选择上述截面 $s:\R_+^\times\to\A_F^\times$。任意 $a\in\A_F^\times$ 可唯一写为
\[
 a=(a\,s(|a|_\A)^{-1})\,s(|a|_\A),
\]
其中第一因子在 $\A_F^1$。模去 $F^\times$ 后得到
\[
 \A_F^\times/F^\times\simeq \A_F^1/F^\times\times\R_+^\times.
\]
连续性来自范数映射和截面的连续性。
\end{proof}

\begin{exercise}
设 $S$ 是包含所有无穷素位的有限素位集。证明：
\begin{enumerate}[label=(\arabic*)]
\item $\A_F^\times/F^\times\A_F^{\times S}$ 是有限集。
\item 存在这样的 $S$，使 $\A_F^\times=F^\times\A_F^{\times S}$。
\end{enumerate}
\end{exercise}

\begin{proof}[解答]
状态：solvable（可解）。

用同态 $i:\A_F^\times\to I_F$。因为 $\A_F^{\times S}$ 在 $S$ 外的有限素位都是局部单位，所以商
\[
 \A_F^\times/F^\times\A_F^{\times S}
\]
只记录 $S$ 外素理想指数所决定的理想类。更精确地，$i$ 诱导到理想类群 $\Cl_F$ 的映射。两个理元若相差主理元和 $\A_F^{\times S}$ 中元素，则其理想只差一个主理想以及 $S$ 中素理想的乘积。因此该商是 $\Cl_F$ 的某个商集，故有限。

第二问即命题 4.13。取理想类群每个类的代表分式理想，再把这些代表中出现的有限素位全部加入 $S$。对任意理元 $a$，理想 $i(a)$ 与某个代表只差主理想。乘以相应主理元后，其 $S$ 外赋值全为零，所以落入 $\A_F^{\times S}$。于是
\[
 \A_F^\times=F^\times\A_F^{\times S}.
\]
\end{proof}

\begin{exercise}
证明：
\begin{enumerate}[label=(\arabic*)]
\item 对 $\xi\in F$，除有限多个 $v$ 外有 $|\xi|_v\le1$。
\item 设 $M$ 为 $F^\times$ 中所有满足 $|\xi|_v\le1$ 对每个 $v$ 都成立的元素。证明 $M$ 是 $F$ 中全体单位根组成的有限循环群。
\end{enumerate}
\end{exercise}

\begin{proof}[解答]
状态：solvable（可解）。

(1) 若 $\xi=0$，结论直接成立。若 $\xi\ne0$，主理想 $(\xi)$ 的素理想分解只含有限多个素理想。对不出现的有限素位 $v$，有 $v(\xi)\ge0$，即 $\xi\in\mathcal O_v$，所以 $|\xi|_v\le1$。无穷素位本来只有有限多个。

(2) 先证 $M$ 中元素都是单位根。若 $\xi\in M$，则 $\xi\ne0$ 且
\[
 |\xi|_v\le1\quad\text{对所有 }v.
\]
乘积公式给出
\[
 1=\prod_v|\xi|_v.
\]
所有因子都不超过 $1$，因此每个因子都等于 $1$。特别是所有无穷嵌入下 $|\xi|_v=1$。

把 $F$ 通过 Minkowski 嵌入到
\[
 F_\R=\prod_{v|\infty}F_v.
\]
条件 $|\xi|_v=1$ 说明 $\xi$ 的所有幂 $\xi^n$ 在 $F_\R$ 中都落在一个紧集内；同时有限处 $|\xi|_v=1$ 说明 $\xi$ 是 $\mathcal O_F^\times$。代数整数环在 Minkowski 空间中是离散格，所以一个紧集里只含有限多个代数整数。因此序列
\[
 1,\xi,\xi^2,\ldots
\]
只能取有限多个值，存在 $m>n$ 使 $\xi^m=\xi^n$，从而 $\xi^{m-n}=1$。所以 $\xi$ 是单位根。

反过来，若 $\xi$ 是单位根，则 $\xi^N=1$。任意绝对值下都有
\[
 |\xi|_v^N=1,
\]
故 $|\xi|_v=1\le1$。因此 $M$ 正是 $F$ 中单位根群。域中的有限乘法子群是循环群，所以 $M$ 是有限循环群。
\end{proof}

\begin{exercise}
设 $F$ 为数域。有限素位 $v$ 处，有限维 $F_v$-向量空间的 $F_v$-格指紧开 $\mathcal O_v$-子模；无穷素位处，有限维 $F_v$-向量空间的格指离散子群 $\Gamma$，使商空间紧。设 $V$ 为有限维 $F$-向量空间，$F$-格指有限生成 $\mathcal O_F$-子模 $L\subset V$，并且 $L$ 含有 $V$ 的一组 $F$-基。证明：
\begin{enumerate}[label=(\arabic*)]
\item 对 $V$ 的 $F$-格 $L$，$L$ 在 $V\otimes_FF_v$ 中的闭包 $L_v$ 等于 $L$ 生成的 $\mathcal O_v$-模。
\item 给定 $V$ 中的 $F$-格 $L$，并对每个有限素位 $v$ 给定 $V\otimes_FF_v$ 的 $F_v$-格 $N_v$。存在 $F$-格 $M$ 使 $M$ 的闭包为各 $N_v$，当且仅当除有限多个 $v$ 外 $N_v=L_v$。此时 $M=\bigcap_{v<\infty}(V\cap N_v)$。
\item 若 $L\supset M$ 是两个 $F$-格，则自然映射
\[
 L/M\longrightarrow\prod_{v<\infty}L_v/M_v
\]
是同构，并且
\[
 [L:M]=\prod_{v<\infty}[L_v:M_v].
\]
\end{enumerate}
\end{exercise}

\begin{proof}[解答]
状态：solvable（可解，第二问按有限素位闭包理解）。

原题第二问中闭包应理解为有限素位闭包；若写成无穷素位闭包，则与已给的 $N_v$ 索引不一致。

(1) 取 $L$ 的有限生成元 $e_1,\ldots,e_m$。它们在 $V_v=V\otimes_FF_v$ 中生成一个 $\mathcal O_v$-子模
\[
 \mathcal O_ve_1+\cdots+\mathcal O_ve_m.
\]
因为 $\mathcal O_v$ 紧，有限生成 $\mathcal O_v$-模在有限维 $F_v$-空间中是闭的。另一方面，$F$ 在 $F_v$ 中稠密，$\mathcal O_F$ 在 $\mathcal O_v$ 中稠密，所以 $L$ 的闭包包含所有 $\mathcal O_v$-线性组合。于是闭包恰为该 $\mathcal O_v$-模。

(2) 必要性：若 $M$ 是一个 $F$-格，则 $M$ 与 $L$ 都是有限生成 $\mathcal O_F$-模，并含有同一个 $F$-向量空间的基。存在非零 $a,b\in\mathcal O_F$ 使
\[
 aL\subset M\subset b^{-1}L.
\]
只有有限多个素理想整除 $ab$，所以除这些素位外，局部化后 $M_v=L_v$。因此 $N_v=L_v$ 除有限多个外成立。

充分性：设 $T$ 是有限素位集，使 $N_v=L_v$ 对 $v\notin T$ 成立。定义
\[
 M=\bigcap_{v<\infty}(V\cap N_v).
\]
因为只有 $T$ 中的条件真正改变 $L$，可在 $T$ 内分别选择整数幂 $\mathfrak p_v^{r_v}$，使
\[
 \mathfrak p_v^{r_v}L_v\subset N_v\subset \mathfrak p_v^{-r_v}L_v.
\]
把这些有限条件合并成非零分式理想界限，得到
\[
 aL\subset M\subset b^{-1}L
\]
对某些非零 $a,b\in\mathcal O_F$ 成立。于是 $M$ 是有限生成 $\mathcal O_F$-模，并含有 $V$ 的 $F$-基，所以是 $F$-格。

对每个有限 $v$，$M$ 的局部闭包包含在 $N_v$ 中由定义直接得到。反向包含来自强逼近或中国剩余定理：给定 $x\in N_v$，只需在有限多个素位同时逼近 $x$ 并满足其他 $N_w$ 的整性条件，就能找到 $M$ 中元素趋近于 $x$。因此闭包为 $N_v$。

唯一性：若两个 $F$-格有相同的所有有限局部闭包，则一个元素属于其中一个格，当且仅当它在每个有限完成中属于对应闭包；所以二者都等于 $\bigcap_v(V\cap N_v)$。

(3) 首先只有有限多个 $v$ 满足 $L_v\ne M_v$。定义
\[
 \Phi:L/M\to\prod_{v<\infty}L_v/M_v.
\]
若 $x\in L$ 在所有 $L_v/M_v$ 中为零，则 $x\in M_v$ 对所有有限 $v$ 成立。由上一问的交描述，$x\in M$，所以 $\Phi$ 单射。

证明满射。右侧只有有限多个非零因子。对这些素位 $v_1,\ldots,v_r$，选择代表 $x_i\in L_{v_i}$。由于 $L/M$ 是有限 $\mathcal O_F$-模，可用中国剩余定理在有限商中同时指定各局部分量，得到 $x\in L$，使其在每个 $L_{v_i}/M_{v_i}$ 中等于给定类，在其他地方为零。于是 $\Phi$ 满射。

同构给出阶数相等：
\[
 [L:M]=|L/M|
 =\prod_{v<\infty}|L_v/M_v|
 =\prod_{v<\infty}[L_v:M_v].
\]
\end{proof}

\begin{exercise}
设 $F$ 为数域，$\mathcal O_F$ 为整数环，$F_v$ 为 $F$ 在有限素位 $v$ 处的完备化，$\mathcal O_v$ 为赋值环，$\mathfrak m_v$ 为其极大理想。设
\[
 \mathfrak p_v=\mathcal O_F\cap\mathfrak m_v.
\]
证明：
\begin{enumerate}[label=(\arabic*)]
\item $v\mapsto\mathfrak p_v$ 是有限素位集合到 $\mathcal O_F$ 非零素理想集合的双射。
\item 自然同态 $\mathcal O_F/\mathfrak p_v\to\mathcal O_v/\mathfrak m_v$ 是同构。
\end{enumerate}
\end{exercise}

\begin{proof}[解答]
状态：solvable（可解）。

(1) 有限素位 $v$ 给出 $F$ 上的非 Archimede 绝对值，也给出离散赋值。其赋值环与 $F$ 的交为
\[
 \{x\in F: |x|_v\le1\},
\]
极大理想与 $\mathcal O_F$ 的交为
\[
 \{x\in\mathcal O_F: |x|_v<1\}.
\]
这是 $\mathcal O_F$ 的非零素理想。反过来，每个非零素理想 $\mathfrak p$ 给出离散赋值
\[
 v_\mathfrak p(x)=\operatorname{ord}_\mathfrak p(x),
\]
从而给出完备化 $F_\mathfrak p$ 和有限素位。两种构造互逆，所以得到双射。

(2) 映射良定，因为 $\mathfrak p_v=\mathcal O_F\cap\mathfrak m_v$。它是单射。证明满射：$\mathcal O_F$ 在 $\mathcal O_v$ 中稠密，而 $\mathfrak m_v$ 是开理想，所以任意 $y\in\mathcal O_v$ 可由某个 $x\in\mathcal O_F$ 逼近到 $y-x\in\mathfrak m_v$。于是 $y$ 的剩余类来自 $x$。因此
\[
 \mathcal O_F/\mathfrak p_v\simeq\mathcal O_v/\mathfrak m_v.
\]
\end{proof}

\begin{exercise}
证明以下范数与指数公式。
\begin{enumerate}[label=(\arabic*)]
\item 若 $\mathfrak a,\mathfrak b$ 是分式理想且 $\mathfrak a\supset\mathfrak b$，设
\[
 \mathfrak a^{-1}\mathfrak b=\prod_v\mathfrak p_v^{n(v)}.
\]
证明
\[
 [\mathfrak a:\mathfrak b]=\prod_v[\mathcal O_F:\mathfrak p_v]^{n(v)}.
\]
\item 存在群同态 $N:I_F\to\Q_+^\times$，使 $N(\mathfrak p)=[\mathcal O_F:\mathfrak p]$。
\item 若 $\mathfrak a\supset\mathfrak b$，证明 $[\mathfrak a:\mathfrak b]=N(\mathfrak b)/N(\mathfrak a)$。
\item 若 $a=(a_v)\in\A_F^\times$，证明
\[
 N(i(a))=\prod_{v<\infty}|a_v|_v^{-1}.
\]
\item 对 $\xi\in F^\times$，设有 $\rho$ 个实素位 $w$ 满足 $\xi$ 在 $F_w\simeq\R$ 中为负。证明
\[
 N(i(\xi))=(-1)^\rho N_{F/\Q}(\xi).
\]
\end{enumerate}
\end{exercise}

\begin{proof}[解答]
状态：solvable（可解）。

(1) 先在 $\mathfrak a=\mathcal O_F$ 的情形证明。若
\[
 \mathfrak b=\prod_v\mathfrak p_v^{n(v)}
\]
是整理想，则 Dedekind 环中互素素理想幂的中国剩余分解给出
\[
 \mathcal O_F/\mathfrak b\simeq
 \prod_v \mathcal O_F/\mathfrak p_v^{n(v)}.
\]
并且
\[
 |\mathcal O_F/\mathfrak p_v^{n}|
 =|\mathcal O_F/\mathfrak p_v|^n.
\]
故
\[
 [\mathcal O_F:\mathfrak b]
 =\prod_v[\mathcal O_F:\mathfrak p_v]^{n(v)}.
\]
一般情形可乘以某个非零 $c\in F^\times$，使 $c\mathfrak a$ 和 $c\mathfrak b$ 为整理想。乘以 $c$ 不改变指数，并且
\[
 (c\mathfrak a)^{-1}(c\mathfrak b)=\mathfrak a^{-1}\mathfrak b.
\]
于是得到公式。

(2) 分式理想群 $I_F$ 是由非零素理想自由生成的 Abel 群。规定
\[
 N\left(\prod_v\mathfrak p_v^{m(v)}\right)
 =\prod_v[\mathcal O_F:\mathfrak p_v]^{m(v)}
\]
即可得到唯一的群同态 $I_F\to\Q_+^\times$。

(3) 由 (1) 和 (2)，
\[
 [\mathfrak a:\mathfrak b]
 =N(\mathfrak a^{-1}\mathfrak b)
 =N(\mathfrak b)/N(\mathfrak a).
\]

(4) 若
\[
 i(a)=\prod_{v<\infty}\mathfrak p_v^{v(a_v)},
\]
则
\[
 N(i(a))=\prod_{v<\infty}N(\mathfrak p_v)^{v(a_v)}.
\]
规范化绝对值满足
\[
 |a_v|_v=N(\mathfrak p_v)^{-v(a_v)}.
\]
因此
\[
 N(i(a))=\prod_{v<\infty}|a_v|_v^{-1}.
\]

(5) 将 (4) 应用于主理元 $\xi$：
\[
 N(i(\xi))=\prod_{v<\infty}|\xi|_v^{-1}.
\]
由乘积公式，
\[
 \prod_{v<\infty}|\xi|_v^{-1}
 =\prod_{v|\infty}|\xi|_v.
\]
无穷处乘积等于 $|N_{F/\Q}(\xi)|$。若有 $\rho$ 个实嵌入下 $\xi<0$，复嵌入贡献总为正，故
\[
 |N_{F/\Q}(\xi)|=(-1)^\rho N_{F/\Q}(\xi).
\]
这给出结论。
\end{proof}

\begin{exercise}
证明与模 $\mathfrak m$ 相关的理元类群正合序列：
\[
1\to
\A_F^{\times\mathfrak m}F^\times/\widehat F_{\mathfrak m}F^\times
\to
\A_F^\times/\widehat F_{\mathfrak m}F^\times
\to
\A_F^\times/\A_F^{\times\mathfrak m}F^\times
\to1,
\]
\[
1\to
(\A_F^{\times\mathfrak m}\cap F^\times)/(\widehat F_{\mathfrak m}\cap F^\times)
\to
\A_F^{\times\mathfrak m}/\widehat F_{\mathfrak m}
\to
\A_F^{\times\mathfrak m}F^\times/\widehat F_{\mathfrak m}F^\times
\to1,
\]
并推出
\[
1\to F_\infty^\times/O_F^{\mathfrak m}
\to C(\mathfrak m)\to I_F^{\mathfrak m}/S_F^{\mathfrak m}\to1.
\]
\end{exercise}

\begin{proof}[解答]
状态：solvable（可解）。

为避免符号混乱，记 $\widehat F_{\mathfrak m}$ 为有限处同余子群，$O_F^{\mathfrak m}$ 为满足模 $\mathfrak m$ 同余和无穷正性条件的全局单位型子群。

第一列是一般群论事实。若 $H\subset K\subset G$，则
\[
 1\to K/H\to G/H\to G/K\to1
\]
正合。这里取
\[
 G=\A_F^\times,\quad
 K=\A_F^{\times\mathfrak m}F^\times,\quad
 H=\widehat F_{\mathfrak m}F^\times.
\]
包含关系由 $\widehat F_{\mathfrak m}\subset\A_F^{\times\mathfrak m}$ 给出。

第二列来自乘法映射
\[
 \A_F^{\times\mathfrak m}\to
 \A_F^{\times\mathfrak m}F^\times/\widehat F_{\mathfrak m}F^\times.
\]
它是满射，因为目标群的每个陪集都有形如 $uf$ 的代表，其中
$u\in\A_F^{\times\mathfrak m}$、$f\in F^\times$，而在模
$\widehat F_{\mathfrak m}F^\times$ 后这个代表等于 $u$ 的像。其核由那些
$x\in\A_F^{\times\mathfrak m}$ 满足 $x\in\widehat F_{\mathfrak m}F^\times$ 的元素组成，也就是
\[
 \A_F^{\times\mathfrak m}\cap \widehat F_{\mathfrak m}F^\times.
\]
模去 $\widehat F_{\mathfrak m}$ 后，核正是
\[
 (\A_F^{\times\mathfrak m}\cap F^\times)/(\widehat F_{\mathfrak m}\cap F^\times).
\]
这给出第二个正合列。

接着计算各项。有限处同余子群没有无穷分量，因此
\[
 \A_F^{\times\mathfrak m}/\widehat F_{\mathfrak m}\simeq F_\infty^\times.
\]
全局元素同时落在 $\A_F^{\times\mathfrak m}$ 中，正是满足模 $\mathfrak m$ 的有限同余和无穷正性条件的元素；在无穷处分量中形成子群 $O_F^{\mathfrak m}$。因此
\[
 \A_F^{\times\mathfrak m}F^\times/\widehat F_{\mathfrak m}F^\times
 \simeq F_\infty^\times/O_F^{\mathfrak m}.
\]

最后，用 $i:\A_F^\times\to I_F$。在模 $\mathfrak m$ 的商中，有限处属于 $\A_F^{\times\mathfrak m}$ 的理元给出的理想只含被排除的同余平凡部分；主理元给出射线主理想。因此
\[
 \A_F^\times/\A_F^{\times\mathfrak m}F^\times
 \simeq I_F^{\mathfrak m}/S_F^{\mathfrak m}.
\]
把这些识别代入第一正合列，得到
\[
1\to F_\infty^\times/O_F^{\mathfrak m}
\to C(\mathfrak m)\to I_F^{\mathfrak m}/S_F^{\mathfrak m}\to1.
\]
\end{proof}

\begin{exercise}
设 $S$ 是包含所有无穷素位的有限素位集，定义
\[
 \mathcal O_S=F\cap\A_F^S.
\]
讨论并证明 $S$-整数环的基本性质：$\mathcal O_S$ 的逼近性质、Dedekind 性、分式域以及 $S$-单位群的有限生成性和秩。
\end{exercise}

\begin{proof}[解答]
状态：partially solvable（部分可解；题面第一问的全空间稠密性需按强逼近形式修正）。

先说明题面中的稠密性需要作精确理解。若原命题写成“$\mathcal O_S$ 在 $\A_F^S$ 中稠密”，当 $S=S_\infty$ 时不成立。例如 $F=\Q$、$S=\{\infty\}$ 时，$\mathcal O_S=\Z$，它在 $\R\times\prod_p\Z_p$ 中的实分量不是稠密的。正确的强逼近形式是：固定一个 $v_0\in S$ 后，$\mathcal O_S$ 在省去 $v_0$ 的相应限制直积中稠密；若 $S$ 含有至少一个有限素位，也可利用该有限素位允许分母来逼近其余指定地方。

先证明 $\mathcal O_S$ 的描述。由定义，$x\in F$ 属于 $\mathcal O_S$ 当且仅当对所有有限 $v\notin S$ 有 $x\in\mathcal O_v$，即
\[
 v_\mathfrak p(x)\ge0\qquad(\mathfrak p\notin S).
\]
所以
\[
 \mathcal O_S=\mathcal O_F[\mathfrak p^{-1}:\mathfrak p\in S,\ \mathfrak p<\infty],
\]
也就是允许只在 $S$ 中有限素位出现分母的元素环。

逼近性质来自定理 4.5。固定 $v_0\in S$。给定有限多个 $v\ne v_0$ 的局部目标，并要求 $v\notin S$ 时保持整性，强逼近给出 $\beta\in F$ 同时满足这些条件。由于它在所有 $v\notin S$ 处整，$\beta\in\mathcal O_S$。因此 $\mathcal O_S$ 在省去 $v_0$ 后的限制直积中稠密。

再证 $\mathcal O_S$ 是 Dedekind 环。它是 Dedekind 环 $\mathcal O_F$ 对乘法集
\[
 T=\mathcal O_F\setminus\bigcup_{\mathfrak p\in S_f}\mathfrak p
\]
的局部化，其中 $S_f$ 是 $S$ 中有限素位对应的素理想集合。Dedekind 环的局部化仍是一维 Noether 整闭整环；非零素理想正是 $\mathcal O_F$ 中不属于 $S_f$ 的非零素理想扩张而来。因此 $\mathcal O_S$ 是 Dedekind 环。它的分式域仍是 $F$，因为局部化不改变原整环的分式域。

最后证明 $S$-单位群结构。由上面对 $\mathcal O_S$ 的描述，
\[
 \mathcal O_S^\times
 =\{x\in F^\times: v_\mathfrak p(x)=0\text{ 对所有 }\mathfrak p\notin S_f\}.
\]
定义赋值映射
\[
 \lambda:\mathcal O_S^\times\to\Z^{S_f},\qquad
 x\mapsto (v_\mathfrak p(x))_{\mathfrak p\in S_f}.
\]
其核是 $\mathcal O_F^\times$。Dirichlet 单位定理给出
\[
 \mathcal O_F^\times\simeq\mu(F)\times\Z^{r_1+r_2-1}.
\]
映射 $\lambda$ 的像由主理想在 $S_f$ 中的指数向量组成。理想类群有限说明这些 $S_f$ 中素理想类张成的子群有限指数地控制允许的指数关系；取适当幂后，每个 $\mathfrak p\in S_f$ 的类成为主类，从而得到足够多的 $S$-单位使像在一个秩为 $|S_f|$ 的格中有有限指数。于是
\[
 \operatorname{rank}\mathcal O_S^\times
 =(r_1+r_2-1)+|S_f|=|S|-1.
\]
其中 $|S|=r_1+r_2+|S_f|$ 按 Archimede 素位和有限素位一起计数。挠子部分仍是 $\mu(F)$，所以 $\mathcal O_S^\times$ 是有限生成交换群，秩为 $|S|-1$。
\end{proof}
 \chapter{上同调群}

第二部分的任务，是把类域论中反复出现的“范数、核、余核、局部到整体的障碍”放进一个统一的群上同调框架中。第五章先不碰具体数域，而是建立有限群和射影有限群的同调工具。后面第六章计算局部域，第七章计算理元类群时，真正反复调用的正是本章的连接同态、诱模、维数移位、杯积、Tate 定理和类成原则。

本章要学的不是一堆符号，而是三条主线。第一，$H^q(G,A)$ 记录 $G$-模 $A$ 中“局部兼容但不一定全局平凡”的数据。第二，Tate 上同调把同调和上同调接成一条双向无穷的链，使循环群情形出现周期二。第三，类成原则说：若某个 Galois 模 $A$ 的上同调像整数模 $\Z$ 那样受一个基本类控制，就能从上同调构造互反映射。

把它放回代数数论的主线中，设 $E/F$ 是有限 Galois 扩张，$G=\operatorname{Gal}(E/F)$。
第二部分反复计算的是 $H^q(G,A)$，其中 $A$ 会依次取为四类对象：
\[
E^\times,\qquad
\prod_{w|v}E_w^\times\text{ 及局部单位群},\qquad
\A_E^\times,\qquad
C_E=\A_E^\times/E^\times .
\]
这四类对象正好对应前面几章的全局域、局部域、理元群和理元类群。它们之间有短正合列，
例如
\[
1\longrightarrow E^\times\longrightarrow \A_E^\times\longrightarrow C_E\longrightarrow1.
\]
上同调的长正合列会把这些对象的计算连接起来。对初学者来说，可以先把本章理解成
“为这些短正合列准备计算工具”：连接同态负责在长正合列中移动，诱模和维数移位负责把
难次数化到低次数，Herbrand 商负责控制有限性，Tate 定理负责给出周期二。

类域论中的互反律不是突然出现的公式。它来自这样一个思想：若某个类成对象的上同调由一个
基本类控制，那么杯积就把整数系数的上同调和该对象的上同调联系起来；这种联系在次数上
表现为周期二，在 Galois 群的交换化上表现为互反映射。因此本章虽然看起来抽象，实际是在
为后面“局部互反”和“整体互反”搭建代数机器。

\section{有限群的同调群}

设 $G$ 为有限群。若交换群 $A$ 上有群同态
\[
 \rho:G\longrightarrow \operatorname{Aut}(A),
\]
则称 $A$ 为 $G$-模。通常写 $ga$ 代替 $\rho(g)(a)$。若所有 $g$ 都平凡作用，则称作用平凡。等价地，$A$ 是群环 $\Z[G]$ 上的模。

初学者最容易误会的一点是：群上同调不是凭空给群造新不变量，而是研究“带有群作用的对象”中哪些元素、映射、扩张可以下降到不变量层。类域论中的 $G=\operatorname{Gal}(E/F)$，而 $A$ 常常是 $E^\times$、局部单位群、理元群或理元类群。

\subsection{标准分解}

为了定义同调，必须先有一个标准的自由分解。这里“分解”的意思是：用一串自由
$\Z[G]$-模去解析平凡 $G$-模 $\Z$，再把这个解析拿去对任意 $G$-模 $A$ 做
$\operatorname{Hom}_G(-,A)$ 或张量运算。标准分解的优点是完全显式，缺点是公式看起来长。
学习时不要先背公式，而要先记住它的来源：每一项都在记录“把若干群元素相乘时，哪一个
相邻乘法被合并，或者哪一个端点被删去”。

对 $n\ge1$，令 $G^n$ 为 $n$ 个 $G$ 的直积，以 $X_n$ 表示由符号
$(g_1,\ldots,g_n)$ 生成的自由 $\Z[G]$-模：
\[
 X_n=\bigoplus_{(g_1,\ldots,g_n)\in G^n}\Z[G](g_1,\ldots,g_n).
\]
令 $X_0=X_{-1}=\Z[G]\cdot1$，并约定 $X_{-n-1}=X_n$。正指标部分使用非齐性边界。若 $n>1$，
\[
\begin{aligned}
d_n(g_1,\ldots,g_n)
={}&g_1(g_2,\ldots,g_n)
 +\sum_{i=1}^{n-1}(-1)^i(g_1,\ldots,g_ig_{i+1},\ldots,g_n)\\
&+(-1)^n(g_1,\ldots,g_{n-1}).
\end{aligned}
\]
并定义
\[
 d_1(g)=g-1.
\]
例如
\[
d_2(g,h)=g(h)-(gh)+(g).
\]
这个式子正是 1-上链边界里“先作用 $g$，再合并 $gh$，最后保留 $g$”三种贡献的交替和。
直接计算
\[
d_1d_2(g,h)=g(h-1)-(gh-1)+(g-1)=0,
\]
说明相邻两个边界复合为零在低阶上确实成立。

负指标部分是为了把同调也接进同一条链。记
\[
 N_G=\sum_{g\in G}g\in \Z[G]
\]
为范数元，增广映射和上增广映射分别为
\[
 \varepsilon:\Z[G]\to\Z,\qquad \sum n_gg\mapsto \sum n_g,
\]
\[
 \mu:\Z\to\Z[G],\qquad n\mapsto nN_G.
\]
令 $d_0=\mu\varepsilon$，所以 $d_0(1)=N_G$。再令
\[
 d_{-1}:X_{-1}\to X_{-2},\qquad
 1\mapsto \sum_{g\in G}g^{-1}(g)-(g).
\]
一般地，若 $n\ge1$，则 $d_{-n-1}:X_{-n-1}\to X_{-n-2}$ 在生成元
$(g_1,\ldots,g_n)$ 上定义为
\[
\begin{aligned}
d_{-n-1}(g_1,\ldots,g_n)
={}&\sum_{g\in G}g^{-1}(g,g_1,\ldots,g_n)\\
&+\sum_{g\in G}\sum_{i=1}^n(-1)^i
(g_1,\ldots,g_ig,g^{-1},g_{i+1},\ldots,g_n)\\
&+\sum_{g\in G}(-1)^{n+1}(g_1,\ldots,g_n,g).
\end{aligned}
\]
这些求和项看起来比正指标边界复杂，是因为负半边要同时编码同调的链复形和范数元
$N_G$。低阶处可以看到它怎样衔接：$d_0(1)=N_G$，而
\[
d_{-1}N_G
=\sum_{h\in G}h\left(\sum_{g\in G}g^{-1}(g)-(g)\right)=0,
\]
因为第一重求和经变量替换后与第二重求和相同。所以 $d_{-1}d_0=0$。

这些映射组成正合的 $G$-模链
\[
 \cdots\leftarrow X_{-2}\leftarrow X_{-1}\leftarrow X_0
 \leftarrow X_1\leftarrow X_2\leftarrow\cdots.
\]
称为 $G$ 的非齐性标准分解。
正合性是标准分解的核心性质。直观上，正指标部分是平凡模 $\Z$ 的 bar resolution；
它有显式收缩同伦，所以除了增广处以外没有同调。负指标部分是把对偶的 bar resolution
接到范数映射上。后面使用时通常不需要重新证明整条正合性，但必须知道：正合性保证
由它定义的上同调群只依赖于 $G$ 和 $A$，而不是某个任意选择的链复形。

\begin{remark}
这里的“非齐性”意味着 $n$-上链可看成 $G^n\to A$ 的映射。齐性分解则把 $n$-上链看成 $G^{n+1}\to A$ 的等变映射。两种写法等价：齐性写法对对称性更自然，非齐性写法对实际计算更短。
\end{remark}

若 $M$ 用乘法记号，非齐性 $2$-上闭链，也称因子集，是映射
\[
 \varphi:G\times G\to M
\]
满足
\[
 g_1\varphi(g_2,g_3)\,\varphi(g_1g_2,g_3)^{-1}
 \varphi(g_1,g_2g_3)\,\varphi(g_1,g_2)^{-1}=1.
\]
这个公式的意义是：先合成 $g_2,g_3$ 再合成 $g_1$，与先合成 $g_1,g_2$ 再合成 $g_3$，两条路径给出的误差相同。中心单代数、域扩张中的因子集都用这个条件表达结合律。

非齐性 $2$-上边界，也称平凡因子集，是由一个映射 $\psi:G\to M$ 给出的
\[
\varphi(g_1,g_2)
=g_1\psi(g_2)\,\psi(g_1)\,\psi(g_1g_2)^{-1}.
\]
它表示因子集来自“换一套截面”造成的误差，因此在上同调中应当被视为零。若改用齐性分解，
则 $2$-因子集是映射 $f:G^3\to M$，满足等变性
\[
f(sg_0,sg_1,sg_2)=s f(g_0,g_1,g_2)
\]
以及闭链条件
\[
f(g_1,g_2,g_3)f(g_0,g_2,g_3)^{-1}
f(g_0,g_1,g_3)f(g_0,g_1,g_2)^{-1}=1.
\]
齐性平凡因子集形如
\[
f(g_0,g_1,g_2)=z(g_0,g_1)z(g_1,g_2)z(g_0,g_2)^{-1},
\]
其中 $z(g_0g_1,g_0g_2)=g_0z(g_1,g_2)$。两种语言互相转换；非齐性公式适合计算，
齐性公式适合看出 $G$-等变性。

\subsection{Tate 上同调群}

设 $A$ 为 $G$-模。定义
\[
 A^n=\operatorname{Hom}_G(X_n,A).
\]
因为上一小节的 $X_\bullet$ 是一个链复形，向任意 $G$-模 $A$ 取
$G$-模同态后，箭头方向反过来，得到上链复形
\[
 \cdots\to A^{-2}\xrightarrow{\partial_{-1}}A^{-1}
 \xrightarrow{\partial_0}A^0\xrightarrow{\partial_1}A^1
 \xrightarrow{\partial_2}A^2\to\cdots,
\]
其中
\[
 (\partial_nx)(\sigma)=x(d_n\sigma).
\]
这里的下标容易混淆：$x\in A^{n-1}$ 是定义在 $X_{n-1}$ 上的等变同态，而
$d_n\sigma\in X_{n-1}$，所以 $x(d_n\sigma)$ 定义了一个 $X_n\to A$ 的等变同态。
上一小节的 $d_nd_{n+1}=0$ 保证 $\partial_{n+1}\partial_n=0$。

\begin{definition}
定义
\[
 Z^n(G,A)=\ker\partial_{n+1},\qquad
 B^n(G,A)=\operatorname{im}\partial_n,
\]
并定义 Tate 上同调群
\[
 H^n(G,A)=Z^n(G,A)/B^n(G,A).
\]
\end{definition}

当 $n\ge1$ 时，$H^n(G,A)$ 是通常的群上同调；当 $n\le -2$ 时，它对应通常的群同调，
更具体地说 $H^{-n-1}(G,A)$ 对应 $H_n(G,A)$。$H^{-1}$ 和 $H^0$ 是 Tate 理论中特别插入的
两项，它们把“同调半边”和“上同调半边”接起来。记
\[
 I_G=(g-1:g\in G)\subset \Z[G],
\]
\[
 A^G=\{a\in A:ga=a\text{ 对所有 }g\in G\},
\qquad
{}_{N_G}A=\{a\in A:N_Ga=0\}.
\]
下面把两个低阶公式算出来。因为 $X_0=\Z[G]$，一个 $G$-等变同态
$X_0\to A$ 由 $1\mapsto a$ 唯一决定，所以 $A^0\simeq A$。由于
$X_{-1}=\Z[G]$ 也是秩一自由 $\Z[G]$-模，一个 $G$-等变同态
$X_{-1}\to A$ 同样由 $1$ 的像唯一决定，因此 $A^{-1}\simeq A$。在这些识别下，$\partial_0:A^{-1}\to A^0$ 由
$d_0(1)=N_G$ 给出，因此
\[
\partial_0(a)=N_Ga=\sum_{g\in G}ga.
\]
而 $\partial_1:A^0\to A^1$ 由 $d_1(g)=g-1$ 给出，因此
\[
(\partial_1a)(g)=ga-a.
\]
于是 $\ker\partial_1=A^G$，$\operatorname{im}\partial_0=N_GA$，故
\[
H^0(G,A)=A^G/N_GA.
\]

再看 $H^{-1}$。$\ker\partial_0$ 正是
${}_{N_G}A=\{a:N_Ga=0\}$。$\operatorname{im}\partial_{-1}$ 由 $d_{-1}$ 给出；
在 $A^{-2}\simeq\operatorname{Hom}_G(X_{-2},A)$ 的识别下，它生成的正是所有
$(g-1)a$ 的和，也就是 $I_GA$。因此
\[
 H^{-1}(G,A)={}_{N_G}A/I_GA,
\qquad
 H^0(G,A)=A^G/N_GA.
\]
因此 $H^0$ 的意思是“不变量模去范数”，$H^{-1}$ 的意思是“范数为零的元素模去由群作用差
$(g-1)a$ 生成的差分”。这两个公式在局部类域论中非常重要：$H^0$ 衡量不变量是否都是
范数，$H^{-1}$ 衡量范数为零的关系是否都来自群作用差。

此外
\[
 Z^1(G,A)=\{x:G\to A:x(gh)=x(g)+gx(h)\},
\]
\[
 B^1(G,A)=\{x:g\mapsto ga-a,\ a\in A\}.
\]
所以 $H^1$ 是交叉同态模去主交叉同态。若作用平凡，它就是 $\operatorname{Hom}(G,A)$。
这里 $x(gh)=x(g)+gx(h)$ 是 1-上闭链条件；它说 $x$ 不是普通同态，而是被 $G$ 的作用
扭曲后的同态。主交叉同态 $g\mapsto ga-a$ 表示只由选择一个元素 $a\in A$ 并比较它和
它的 $G$-共轭得到的平凡类。Hilbert 90 的常见形式正是说，在某些乘法群上这样的
$H^1$ 为零。

\subsection{连接同态}

先说函子性。若 $f:A\to A'$ 是 $G$-模同态，则对每个 $n$ 都有
\[
f_*:\operatorname{Hom}_G(X_n,A)\longrightarrow \operatorname{Hom}_G(X_n,A'),
\qquad x\longmapsto f\circ x.
\]
因为 $f$ 与 $G$-作用相容，所以 $f\circ x$ 仍是 $G$-等变同态；又因为
$f_*\partial=\partial f_*$，它给出上同调群同态
\[
\bar f_n:H^n(G,A)\longrightarrow H^n(G,A').
\]
这说明 Tate 上同调不是只给每个模一个群，还把 $G$-模同态变成上同调群同态。

若
\[
 0\to A'\to A\to A''\to0
\]
是 $G$-模短正合列，则逐项取
$\operatorname{Hom}_G(X_n,-)$。因为 $X_n$ 是自由 $\Z[G]$-模，特别是射影模，所以得到
上链复形的短正合列。由蛇引理得到连接同态
\[
 \delta_n:H^n(G,A'')\to H^{n+1}(G,A').
\]
它们组成长正合列：
\[
\cdots\to H^n(G,A')\to H^n(G,A)\to H^n(G,A'')
\xrightarrow{\delta_n}H^{n+1}(G,A')\to\cdots.
\]

连接同态可按“提升、取边界、落回子模”三步理解。取 $[\alpha]\in H^n(G,A'')$，其中
$\alpha\in Z^n(G,A'')$。先把 $\alpha$ 提升为某个 $n$-上链
$\widetilde\alpha\in A^n$。由于 $\alpha$ 是闭链，$\partial\widetilde\alpha$ 在
$A''$ 中的像为 $0$，所以 $\partial\widetilde\alpha$ 实际来自 $A'$，记其原像为
$\beta\in A'^{\,n+1}$。再检查 $\partial\beta=0$，于是 $\beta$ 给出
$H^{n+1}(G,A')$ 中的类。不同提升只会改变一个上边界，所以得到良定同态
\[
\delta_n([\alpha])=[\beta].
\]
这就是“提升失败就是下一维边界”的精确含义。

连接同态还满足自然性。若有短正合列之间的交换图
\[
\begin{tikzcd}
0 \arrow[r] & A' \arrow[r] \arrow[d,"f'"] & A \arrow[r] \arrow[d,"f"] &
A'' \arrow[r] \arrow[d,"f''"] & 0\\
0 \arrow[r] & C' \arrow[r] & C \arrow[r] & C'' \arrow[r] & 0,
\end{tikzcd}
\]
则相应长正合列中的方块交换，特别是
\[
\bar f'_{n+1}\circ\delta_n=\delta_n\circ \bar f''_n.
\]
证明直接来自上面的三步构造：先提升再取边界，或先通过 $f''$ 再提升，得到的上链在
$C'$ 中只差一个边界。后面 Hilbert 90、Kummer 序列、Brauer 群和类成基本类都反复使用
这种自然性。

\subsection{诱模}

\begin{definition}
称 $G$-模 $A$ 为 $G$-诱模，如果存在 $A$ 的子群 $S$，使
\[
 A=\bigoplus_{g\in G}gS.
\]
\end{definition}

最基本的例子是 $\Z[G]$。若 $S$ 是交换群且 $G$ 平凡作用在 $S$ 上，则
\[
 \Z[G]\otimes S=\bigoplus_{g\in G}g(1\otimes S)
\]
也是诱模。更一般地，若 $A=\bigoplus_{g\in G}gS$ 是诱模，$B$ 是任意 $G$-模，则
\[
A\otimes B
=\left(\bigoplus_{g\in G}gS\right)\otimes B
=\bigoplus_{g\in G}(gS\otimes gB)
=\bigoplus_{g\in G}g(S\otimes B),
\]
所以 $A\otimes B$ 仍是诱模。这里第二个等号使用 $gB=B$ 作为集合，但保留 $g$ 是为了
标出 $G$-作用怎样把一份 $S\otimes B$ 移到另一份。诱模的作用是提供“上同调为零”的对象，
好让我们用分解去计算一般模。

\begin{proposition}
设 $A$ 是 $G$-诱模。
\begin{enumerate}[label=(\arabic*)]
\item 若 $H$ 为 $G$ 的子群，则 $A$ 作为 $H$-模仍是 $H$-诱模。
\item 若 $H$ 是 $G$ 的正规子群，则 $A^H$ 是 $G/H$-诱模。
\end{enumerate}
\end{proposition}

\begin{proof}
设 $A=\bigoplus_{g\in G}gS$。若取左陪集分解 $G=\coprod Hx$，则
\[
 A=\bigoplus_{h\in H}h\left(\bigoplus_xxS\right),
\]
括号内的 $\bigoplus_xxS$ 是一个普通子群，$H$ 作用把它的各个 $h$-平移分量排成直和。
这正符合 $H$-诱模的定义，所以 $A$ 作为 $H$-模是诱模。

若 $H\triangleleft G$，取 $a\in A^H$。把 $a$ 唯一写成
\[
a=\sum_{g\in G}gs_g,\qquad s_g\in S.
\]
对任意 $h\in H$，因为 $ha=a$，又
\[
ha=\sum_{g\in G}hgs_g.
\]
将指标 $g$ 换成 $h^{-1}g$，得到
\[
ha=\sum_{g\in G}g\,s_{h^{-1}g}.
\]
直和表达唯一，因此 $s_g=s_{h^{-1}g}$，也就是同一个左 $H$-陪集上的系数相同。取
$G=\coprod_y yH$，则
\[
 a=\sum_y y\sum_{h\in H}hs_y=\sum_y yN_Hs_y.
\]
反向包含也成立。若 $H$ 正规，则 $yN_Hy^{-1}=N_H$，所以对任意 $h_0\in H$，
\[
h_0(yN_Hs)=y(y^{-1}h_0y)N_Hs=yN_Hs.
\]
因此每个 $yN_Hs$ 都被 $H$ 固定。于是
\[
 A^H=\bigoplus_{y\in G/H}yN_HS,
\]
这是 $G/H$-诱模。
\end{proof}

\begin{definition}
设 $A$ 为 $G$-模。
\begin{enumerate}[label=(\arabic*)]
\item 若 $H^n(G,A)=0$ 对所有 $n>0$ 成立，则称 $A$ 为零调。
\item 若对 $G$ 的所有子群 $H$ 和所有整数 $q$ 都有 $H^q(H,A)=0$，则称 $A$ 有平凡上同调。
\end{enumerate}
\end{definition}

\begin{proposition}
诱模有平凡上同调。
\end{proposition}

\begin{proof}
先设 $A=\bigoplus_{g\in G}gS$。对任意自由 $\Z[G]$-模 $X$，一个 $G$-模同态 $X\to A$ 由它投影到 $S$ 的普通 $\Z$-模同态决定，因此
\[
 \operatorname{Hom}_G(X,A)\simeq \operatorname{Hom}_\Z(X,S).
\]
标准分解在 $\Z$-模意义下不仅正合，而且是分裂正合：忘掉 $G$-作用后，bar resolution
有显式收缩同伦。对分裂正合列取 $\operatorname{Hom}_\Z(-,S)$ 仍然正合。因此
\[
\cdots\to \operatorname{Hom}_G(X_q,A)\to
\operatorname{Hom}_G(X_{q+1},A)\to\cdots
\]
正合，对应的 Tate 上同调全为零。对子群的结论由命题 5.2(1) 得到。
\end{proof}

\begin{proposition}
若 $G$-模 $A$ 有零调分解
\[
 0\to A\to Y^0\to Y^1\to Y^2\to\cdots,
\]
则
\[
 H^n(G,A)\simeq H^n\bigl((Y^\bullet)^G\bigr).
\]
\end{proposition}

\begin{proof}
记 $C^p=\ker(Y^p\to Y^{p+1})$。由短正合列
\[
 0\to C^p\to Y^p\to C^{p+1}\to0
\]
和 $Y^p$ 的零调性，连接同态给出逐步同构
\[
 H^n(G,A)\simeq H^{n-1}(G,C^1)\simeq\cdots\simeq H^1(G,C^{n-1}).
\]
另一方面，$(Y^\bullet)^G$ 的第 $n$ 个同调正是
\[
 \ker((Y^n)^G\to(Y^{n+1})^G)/
 \operatorname{im}((Y^{n-1})^G\to(Y^n)^G),
\]
它由同一串短正合列的 $H^0$ 部分给出。两边通过相同的连接同态识别。
\end{proof}

现在给出诱导模的常用形式。若 $H\le G$，$A$ 是 $H$-模，定义
\[
 \operatorname{Ind}_H^G(A)=
 \{f:G\to A:f(st)=sf(t)\text{ 对所有 }s\in H,t\in G\},
\]
并令 $G$ 通过右乘作用：
\[
 (t\cdot f)(s)=f(st).
\]

\begin{lemma}
设 $H$ 是 $G$ 的子群，$A$ 是 $H$-模。则对 $n\ge0$ 有自然同构
\[
 H^n(G,\operatorname{Ind}_H^G A)\simeq H^n(H,A).
\]
\end{lemma}

\begin{proof}
先看 $n=0$。若 $f\in(\operatorname{Ind}_H^G A)^G$，则 $f$ 由 $f(1)$ 决定，而且 $f(1)\in A^H$。故 $H^0$ 相同。

一般 $n$ 用齐性复形。对 $\operatorname{Ind}_H^G A$ 的 $G$-等变上链，取在 $1$ 处的值，得到 $H$-等变上链；反向把 $H$-等变函数 $u$ 送到
\[
 f(g_0,\ldots,g_n)(s)=u(sg_0,\ldots,sg_n).
\]
两个构造互逆：从 $f$ 取 $f(\cdots)(1)$ 再按上式恢复，会利用诱导模条件
$f(st)=sf(t)$；从 $u$ 构造 $f$ 再取 $s=1$ 则直接回到 $u$。删去某个变量的微分操作与
在 $s=1$ 取值相容，所以这是复形同构。因为齐性复形由诱模组成，命题 5.5 给出上同调同构。
\end{proof}

\begin{proposition}
设 $A$ 为 $G$-模。Shapiro 同构与上限制同态相容：由自然映射进入诱导模所得的同态，与先 Shapiro 再 corestriction 所得同态一致。
\end{proposition}

\begin{proof}
两边都是从 $H$-上同调到 $G$-上同调的自然 $\delta$-函子态射。它们在 $H^0$ 上都把 $a\in A^H$ 送到陪集求和
\[
 \sum_{s\in G/H}sa\in A^G.
\]
泛 $\delta$-函子的唯一性推出所有次数相同。
\end{proof}

\subsection{维数移位}

维数移位的目标是把任意次数的上同调转成 $H^0$ 或 $H^1$ 来算。记
\[
 I_G=\ker(\varepsilon:\Z[G]\to\Z),\qquad
 J_G=\Z[G]/\Z N_G.
\]
有短正合列
\[
 0\to I_G\to \Z[G]\to\Z\to0,
\qquad
 0\to\Z\to\Z[G]\to J_G\to0.
\]

\begin{definition}
对 $G$-模 $A$ 定义
\[
A^{[m]}=
\begin{cases}
J_G^{\otimes m}\otimes A,&m\ge0,\\
I_G^{\otimes(-m)}\otimes A,&m\le0.
\end{cases}
\]
当 $m=0$ 时约定空张量积为 $\Z$，所以 $A^{[0]}=A$。正的移位反复使用
$J_G$，负的移位反复使用 $I_G$。
\end{definition}

\begin{theorem}
对所有整数 $q,m$ 和 $G$ 的所有子群 $H$，有同构
\[
 \delta^m:H^{q-m}(H,A^{[m]})\simeq H^q(H,A).
\]
\end{theorem}

\begin{proof}
把上面两个短正合列张量 $A$。这里张量后仍正合，是因为这两个短正合列作为
$\Z$-模短正合列是分裂的：$\Z[G]\to\Z$ 可由 $1\mapsto 1_G$ 给出一个
$\Z$-线性截面，$\Z\to\Z[G]$ 的像 $\Z N_G$ 是自由循环子群，也可在自由交换群
$\Z[G]$ 中取补。于是得到
\[
 0\to I_G\otimes A\to \Z[G]\otimes A\to A\to0,
\]
\[
 0\to A\to \Z[G]\otimes A\to J_G\otimes A\to0.
\]
中间项 $\Z[G]\otimes A$ 是诱模，对所有子群都有平凡上同调。长正合列中的连接同态于是给出
\[
 H^{q-1}(H,J_G\otimes A)\simeq H^q(H,A),
\]
\[
 H^{q+1}(H,I_G\otimes A)\simeq H^q(H,A).
\]
反复合成这些同构，就得到任意 $m$ 的结论。
\end{proof}

\begin{lemma}
设 $G=\langle\sigma\rangle$ 是 $n$ 阶循环群。
\begin{enumerate}[label=(\arabic*)]
\item $\Z[G]=\oplus_{i=0}^{n-1}\Z\sigma^i$，$N_G=1+\sigma+\cdots+\sigma^{n-1}$，且 $I_G=(\sigma-1)\Z[G]$。
\item 映射 $Z^1(G,A)\to Z^{-1}(G,A)$，$x\mapsto x(\sigma)$ 是同构。
\end{enumerate}
\end{lemma}

\begin{proof}
(1) 群环基底来自循环群元素。若 $k\ge1$，则
\[
 \sigma^k-1=(\sigma-1)(1+\sigma+\cdots+\sigma^{k-1}),
\]
所以所有增广为零的元素都由 $\sigma-1$ 生成。

(2) 若 $x$ 是 $1$-上闭链，则
\[
 x(\rho\tau)=x(\rho)+\rho x(\tau).
\]
由此 $x(1)=0$，并且
\[
 x(\sigma^k)=\sum_{i=0}^{k-1}\sigma^ix(\sigma).
\]
令 $k=n$ 得 $N_Gx(\sigma)=x(1)=0$。反过来，若 $a\in A$ 满足 $N_Ga=0$，定义
\[
 x(\sigma^k)=\sum_{i=0}^{k-1}\sigma^ia
\]
即可得到唯一的 $1$-上闭链，且 $x(\sigma)=a$。
\end{proof}

\begin{theorem}
若 $G$ 是循环群，则对所有 $q\in\Z$ 有自然同构
\[
 H^q(G,A)\simeq H^{q+2}(G,A).
\]
\end{theorem}

\begin{proof}
引理 5.10 把 $Z^1$ 与 $Z^{-1}$ 对应起来。边界也对应：$1$-边界由 $a\mapsto \sigma a-a$ 给出，$-1$ 次边界正是 $I_GA=(\sigma-1)A$。因此
\[
 H^1(G,A)\simeq H^{-1}(G,A).
\]
再用维数移位把这个同构搬到任意次数，得到周期二。
\end{proof}

\begin{proposition}
设 $|G|=n$。则
\[
 n\cdot H^q(G,A)=0
\]
对所有 $q$ 成立。若 $A$ 是有限生成 $\Z$-模，则各 $H^q(G,A)$ 都是有限群。
\end{proposition}

\begin{proof}
在 $H^0$ 中，$A^G/N_GA$ 被 $n$ 杀掉，因为对 $a\in A^G$ 有 $N_Ga=na$。由维数移位可推广到所有 $q$。若 $A$ 有限生成，则每个有限次数的上链群也是有限生成，且上同调被 $n$ 杀掉，所以为有限群。
\end{proof}

\begin{proposition}
若 $A$ 是唯一可除群，即对任意正整数 $n$ 和 $a\in A$，方程 $nx=a$ 有唯一解，则 $A$ 有平凡上同调。
\end{proposition}

\begin{proof}
对任意子群 $H$，命题 5.12 说明 $H^q(H,A)$ 被 $|H|$ 杀掉。唯一可除性使乘以 $|H|$ 在 $A$ 上是同构，并诱导上同调上的同构。因此一个同时被 $|H|$ 杀掉、又被乘以 $|H|$ 自同构的群只能为零。
\end{proof}

\begin{proposition}
设有限群 $G$ 平凡作用在 $\Z$ 上，且 $n=|G|$。则
\[
H^{-2}(G,\Z)\simeq G/[G,G],\qquad H^{-1}(G,\Z)=0,
\]
\[
H^0(G,\Z)\simeq \Z/n\Z,\qquad H^1(G,\Z)=0,
\]
并且
\[
H^2(G,\Z)\simeq H^1(G,\Q/\Z)\simeq\operatorname{Hom}(G,\Q/\Z).
\]
\end{proposition}

\begin{proof}
先算 $H^0$ 与 $H^{-1}$。平凡作用下 $A^G=\Z$，范数 $N_G$ 是乘以 $n$，所以
\[
H^0(G,\Z)=\Z/n\Z.
\]
同时
\[
{}_{N_G}\Z=\{a\in\Z:na=0\}=0,
\]
故 $H^{-1}(G,\Z)=0$。

再看 $H^1$。平凡作用下，$1$-上闭链就是群同态 $G\to\Z$。有限群到无挠群 $\Z$ 的同态
只能为零，所以 $H^1(G,\Z)=0$。

计算 $H^{-2}$ 用增广理想。由短正合列
\[
0\to I_G\to \Z[G]\to\Z\to0
\]
和 $\Z[G]$ 的平凡上同调，连接同态给出
\[
H^{-2}(G,\Z)\simeq H^{-1}(G,I_G).
\]
而
\[
H^{-1}(G,I_G)=I_G/I_G^2.
\]
映射
\[
G\longrightarrow I_G/I_G^2,\qquad g\longmapsto (g-1)\bmod I_G^2
\]
是群同态，因为
\[
gh-1=(g-1)+(h-1)+(g-1)(h-1),
\]
最后一项在 $I_G^2$ 中。它的核包含换位子群，并且 $I_G$ 由所有 $g-1$ 生成，所以诱导同构
$G/[G,G]\simeq I_G/I_G^2$。

最后，用短正合列
\[
0\to\Z\to\Q\to\Q/\Z\to0.
\]
$\Q$ 是唯一可除群，上一命题给出它有平凡上同调。因此连接同态给出
\[
H^2(G,\Z)\simeq H^1(G,\Q/\Z).
\]
平凡作用下，$\Q/\Z$ 上的 $1$-上闭链是群同态 $G\to\Q/\Z$，而 $1$-上边界为零，所以
$H^1(G,\Q/\Z)=\operatorname{Hom}(G,\Q/\Z)$。
\end{proof}

\begin{proposition}
设 $H\le G$，$G=\coprod_{i=1}^m s_iH$，令
\[
 B=\bigoplus_{i=1}^m s_iA.
\]
则 $B$ 是从 $H$ 到 $G$ 的诱导模，且
\[
 H^q(G,B)\simeq H^q(H,A).
\]
\end{proposition}

\begin{proof}
这是 Shapiro 引理在陪集代表下的显式写法。诱导模中的函数由各陪集代表上的值决定：
若 $f\in\operatorname{Ind}_H^G(A)$，则 $f(hs_i)=hf(s_i)$，所以给出 $f$ 等价于给出
每个 $f(s_i)\in A$。这把诱导模识别为
\[
B=\bigoplus_{i=1}^m s_iA.
\]
在这个识别下，投影到 $i=1$ 的分量就是源书中的 $\pi:B\to A$。Shapiro 同构说明
\[
\bar\pi\circ\operatorname{Res}:H^q(G,B)\to H^q(H,A)
\]
是同构。$q=0$ 时可以直接看出：若 $b\in A^H$，则
\[
\nu(b)=\sum_{i=1}^m s_i b
\]
是 $G$-不变量，且投影限制后回到 $b$；反方向也由不变量条件逐陪集恢复。一般次数由
维数移位或 Shapiro 引理的复形证明推出。
\end{proof}

\subsection{Herbrand 商}

前面已经把循环群的 Tate 上同调化成两个算子：
\[
 \sigma-1,\qquad N_G=1+\sigma+\cdots+\sigma^{n-1}.
\]
这两个算子互相复合为零，因为
\[
(\sigma-1)N_G=\sigma^n-1=0,\qquad N_G(\sigma-1)=\sigma^n-1=0.
\]
因此
\[
H^0(G,A)=\ker(\sigma-1)/\operatorname{im}N_G,\qquad
H^{-1}(G,A)=\ker N_G/\operatorname{im}(\sigma-1).
\]
Herbrand 商要记录的不是单个上同调群，而是这两个相邻 Tate 上同调群的阶数比例。
这个比例在短正合列中会相乘；在有限模上等于 $1$；在数论应用里，它能把单位群、
idele 群、idele 类群的 cohomology size 计算压缩成一个可传递的数。

先给出不带群作用的一般定义。设 $A$ 是交换群，$\varphi,\psi:A\to A$ 是自同态，并满足
\[
\varphi\psi=\psi\varphi=0.
\]
于是 $\operatorname{im}\psi\subset\ker\varphi$ 且
$\operatorname{im}\varphi\subset\ker\psi$，所以下面的两个商群有定义。若指数有限，定义
\[
 h_{\varphi,\psi}(A)=
 \frac{|\ker\varphi/\operatorname{im}\psi|}
 {|\ker\psi/\operatorname{im}\varphi|}.
\]
如果 $\psi$ 是乘以正整数 $n$ 的映射，$\varphi=0$，则记
\[
h_{0,n}(A)=\frac{|A/nA|}{|{}_nA|},\qquad
{}_nA=\{a\in A:na=0\}.
\]
这里分子测量 ``模掉 $n$ 倍以后还剩多少''，分母测量 ``本来已经被 $n$ 杀掉多少''。
当 $A\simeq \Z^r\oplus A_{\mathrm{tor}}$ 且 $A_{\mathrm{tor}}$ 有限时，有限挠子对这个比例
的贡献抵消，故
\[
h_{0,n}(A)=n^r.
\]
证明是直接的：$A_{\mathrm{tor}}/nA_{\mathrm{tor}}$ 与 ${}_nA_{\mathrm{tor}}$ 有相同阶数，
而 $\Z^r/n\Z^r$ 的阶为 $n^r$，${}_n\Z^r=0$。

若 $G=\langle\sigma\rangle$ 是 $n$ 阶循环群，取
\[
\varphi=\sigma-1,\qquad \psi=N_G,
\]
并且相应两个商群有限，则把 $h_{\sigma-1,N_G}(A)$ 记作 $h(A)$。

\begin{proposition}
若 $A$ 为有限群，则 $h_{\varphi,\psi}(A)=1$。
\end{proposition}

\begin{proof}
因为 $A$ 有限，任一自同态 $\alpha:A\to A$ 都满足
\[
|A|=|\ker\alpha|\,|\operatorname{im}\alpha|.
\]
把 $\alpha$ 分别取为 $\varphi$ 与 $\psi$，得到
\[
|\operatorname{im}\varphi|=\frac{|A|}{|\ker\varphi|},\qquad
|\operatorname{im}\psi|=\frac{|A|}{|\ker\psi|}.
\]
又因为 $\operatorname{im}\psi\subset\ker\varphi$ 与
$\operatorname{im}\varphi\subset\ker\psi$，所以
\[
|\ker\varphi/\operatorname{im}\psi|
=\frac{|\ker\varphi|}{|\operatorname{im}\psi|},\qquad
|\ker\psi/\operatorname{im}\varphi|
=\frac{|\ker\psi|}{|\operatorname{im}\varphi|}.
\]
代入定义得
\[
h_{\varphi,\psi}(A)
=\frac{|\ker\varphi|\,|\operatorname{im}\varphi|}
{|\ker\psi|\,|\operatorname{im}\psi|}
=\frac{|A|}{|A|}=1.
\]
这说明 Herbrand 商只看有限生成群的自由秩部分，有限挠子部分不会改变它。
\end{proof}

\begin{proposition}
设 $G=\langle\sigma\rangle$ 是 $n$ 阶循环群，$A$ 为有限生成 $G$-模，记 $h(A)=h_{\sigma-1,N_G}(A)$。则
\[
 h(A)=\frac{|H^0(G,A)|}{|H^{-1}(G,A)|}
 =\frac{|H^2(G,A)|}{|H^1(G,A)|}.
\]
若 $0\to A\to B\to C\to0$ 是正合列，则
\[
 h(B)=h(A)h(C).
\]
若 $C$ 有限，则 $h(A)=h(B)$。
\end{proposition}

\begin{proof}
循环群的 Tate 上同调有周期 $2$，且上一节已经得到
\[
 H^0=A^G/N_GA,\qquad H^{-1}={}_{N_G}A/I_GA,
\]
其中 $I_GA=(\sigma-1)A$。这正是
\[
h(A)=\frac{|H^0(G,A)|}{|H^{-1}(G,A)|}.
\]
再用周期性 $H^{q+2}(G,A)\simeq H^q(G,A)$，得到
\[
\frac{|H^2(G,A)|}{|H^1(G,A)|}
=\frac{|H^0(G,A)|}{|H^{-1}(G,A)|}.
\]

现在证明乘积性。短正合列给出长正合列。截取周期中的六个相邻项，得到循环正合序列
\[
\begin{tikzcd}[column sep=small]
H^{-1}(G,A)\arrow[r,"\alpha_1"]&
H^{-1}(G,B)\arrow[r,"\alpha_2"]&
H^{-1}(G,C)\arrow[d,"\alpha_3"]\\
H^{0}(G,C)\arrow[u,"\alpha_6"]&
H^{0}(G,B)\arrow[l,"\alpha_5"']&
H^{0}(G,A)\arrow[l,"\alpha_4"']
\end{tikzcd}
\]
这里的箭头方向只表示这六项在长正合列中的相邻关系。令
$F_i=|\operatorname{im}\alpha_i|$。正合性说明每一项的阶等于进入该项的像阶乘以离开该项的
像阶：
\[
\begin{aligned}
|H^{-1}(G,A)|&=F_6F_1,&
|H^{-1}(G,B)|&=F_1F_2,&
|H^{-1}(G,C)|&=F_2F_3,\\
|H^0(G,A)|&=F_3F_4,&
|H^0(G,B)|&=F_4F_5,&
|H^0(G,C)|&=F_5F_6.
\end{aligned}
\]
于是
\[
h(A)h(C)
=\frac{F_3F_4}{F_6F_1}\cdot
\frac{F_5F_6}{F_2F_3}
=\frac{F_4F_5}{F_1F_2}
=h(B).
\]
若 $C$ 有限，则前一命题给出 $h(C)=1$，乘积公式化为 $h(B)=h(A)$。
\end{proof}

\begin{lemma}
\begin{enumerate}[label=(\arabic*)]
\item 若 $G$ 是 $n$ 阶循环群并平凡作用在 $A$ 上，则 $h(A)=h_{0,n}(A)$。
\item 若 $0\to A\to B\to C\to0$ 是交换群正合列，则 $h_{0,n}(B)=h_{0,n}(A)h_{0,n}(C)$。
\item 若 $f,g$ 是交换群 $A$ 的可交换自同态，则 $h_{0,gf}(A)=h_{0,g}(A)h_{0,f}(A)$。
\end{enumerate}
\end{lemma}

\begin{proof}
(1) 平凡作用下每个 $\sigma^i$ 都是恒等映射，所以
\[
\sigma-1=0,\qquad N_G=1+\sigma+\cdots+\sigma^{n-1}=n.
\]
代入循环群 Herbrand 商定义得到结论。

(2) 把 $A,B,C$ 都看成平凡 $G$-模，其中 $|G|=n$。由上一命题的乘积性得到
\[
h_{0,n}(B)=h(B)=h(A)h(C)=h_{0,n}(A)h_{0,n}(C).
\]

(3) 因为 $fg=gf$，先看由 $g$ 与 $f$ 产生的短正合列
\[
0\to \frac{\ker f}{g(A)\cap\ker f}
\to \frac{A}{g(A)}
\to \frac{f(A)}{fg(A)}
\to 0.
\]
中间箭头把 $a+g(A)$ 送到 $f(a)+fg(A)$。它的核是
$(\ker f+g(A))/g(A)\simeq \ker f/(g(A)\cap\ker f)$，像是 $f(A)/fg(A)$。
取阶得
\[
\frac{[A:fg(A)]}{[A:f(A)]}
=\frac{[A:g(A)]\,|g(A)\cap\ker f|}{|\ker f|}.
\]
另一方面，$g$ 诱导同构
\[
\frac{\ker(fg)}{\ker g}\simeq g(A)\cap\ker f,
\qquad a+\ker g\mapsto g(a),
\]
所以
\[
|\ker(fg)|=|\ker g|\,|g(A)\cap\ker f|.
\]
把这两个等式相除，得到
\[
 \frac{[A:gf(A)]}{|\ker(gf)|}
 =
 \frac{[A:g(A)]}{|\ker g|}\,
 \frac{[A:f(A)]}{|\ker f|}.
\]
这正是 $h_{0,gf}(A)=h_{0,g}(A)h_{0,f}(A)$。
\end{proof}

\begin{theorem}
设 $G$ 是素数 $p$ 阶循环群，$A$ 是有限生成 $G$-模。记 $\operatorname{rk}(A)$ 为 $A$ 的秩，则
\[
 h(A)^{p-1}=\frac{h_{0,p}(A^G)^p}{h_{0,p}(A)}
\]
并且
\[
 h(A)=p^{(p\,\operatorname{rk}(A^G)-\operatorname{rk}(A))/(p-1)}.
\]
\end{theorem}

\begin{proof}
令 $G=\langle\sigma\rangle$。第一步把 $A$ 分成不动部分和增广理想作用出来的部分。映射
$A\to I_GA$ 为 $a\mapsto(\sigma-1)a$，其核是 $A^G$，故有正合列
\[
 0\to A^G\to A\to I_GA\to0
\]
和 Herbrand 商乘积性
\[
h(A)=h(A^G)h(I_GA).
\]
在 $A^G$ 上作用平凡，所以
\[
h(A^G)=h_{0,p}(A^G).
\]

第二步计算 $I_GA$ 上的 $p$ 倍映射。因为 $N_G(\sigma-1)=0$，
$I_GA$ 自然是 $\Z[G]/\Z N_G$-模。取 $p$ 次本原单位根 $\zeta_p$，有
\[
 \Z[G]/\Z N_G\simeq
 \Z[X]/(1+X+\cdots+X^{p-1})\simeq \Z[\zeta_p],
\qquad \sigma\mapsto \zeta_p.
\]
在 $\Z[\zeta_p]$ 中，理想 $(p)$ 等于 $(\zeta_p-1)^{p-1}$ 乘以一个单位的主理想。
换回 $\Z[G]/\Z N_G$，存在一个单位 $\varepsilon$，使得在 $I_GA$ 上
\[
p=(\sigma-1)^{p-1}\varepsilon.
\]
单位给出的自同态是同构，所以它的 $h_{0,\varepsilon}$ 为 $1$。由上一引理，
\[
h_{0,p}(I_GA)
=h_{0,(\sigma-1)^{p-1}}(I_GA)
=h_{0,\sigma-1}(I_GA)^{p-1}.
\]
在 $I_GA$ 上 $N_G=0$，因此
\[
h_{0,\sigma-1}(I_GA)
=\frac{|I_GA/(\sigma-1)I_GA|}{|\ker(\sigma-1)|}
=\frac{1}{h_{\sigma-1,N_G}(I_GA)}
=\frac{1}{h(I_GA)}.
\]
于是
\[
h_{0,p}(I_GA)=\frac{1}{h(I_GA)^{p-1}}.
\]

再把平凡模的乘积公式用于
$0\to A^G\to A\to I_GA\to0$，得
\[
h_{0,p}(A)=h_{0,p}(A^G)h_{0,p}(I_GA)
=\frac{h_{0,p}(A^G)}{h(I_GA)^{p-1}}.
\]
也就是
\[
h(I_GA)^{p-1}=\frac{h_{0,p}(A^G)}{h_{0,p}(A)}.
\]
与 $h(A)=h_{0,p}(A^G)h(I_GA)$ 合并，得到
\[
h(A)^{p-1}
=h_{0,p}(A^G)^{p-1}
\frac{h_{0,p}(A^G)}{h_{0,p}(A)}
=\frac{h_{0,p}(A^G)^p}{h_{0,p}(A)}.
\]

最后把有限生成交换群写成自由部分加有限挠子。有限挠子对 $h$ 与 $h_{0,p}$ 的贡献都是 $1$，
所以只需看自由秩。若 $M\simeq\Z^r$，则
\[
h_{0,p}(M)=|M/pM|=p^r.
\]
因此
\[
h(A)^{p-1}
=\frac{p^{p\,\operatorname{rk}(A^G)}}{p^{\operatorname{rk}(A)}}
=p^{p\,\operatorname{rk}(A^G)-\operatorname{rk}(A)},
\]
开 $(p-1)$ 次方得到所需公式。
\end{proof}

\subsection{限制同态}

上同调不仅随模变，也随群变。数论中经常出现塔
$F\subset E\subset L$，于是有子群
\[
\operatorname{Gal}(L/E)\le \operatorname{Gal}(L/F),
\]
也有商群
\[
\operatorname{Gal}(L/F)\twoheadrightarrow \operatorname{Gal}(E/F).
\]
为了在这些群之间比较上同调类，需要把群同态转化为上同调群之间的同态。本节先处理
子群方向，这就是限制同态；再处理反方向的上限制同态。

设 $f:G\to G'$ 是群同态，$A$ 为 $G'$-模。通过
\[
g\cdot a=f(g)a
\]
可以把 $A$ 看成 $G$-模。非齐性标准分解的第 $q$ 项由 $G^q$ 上的函数描述，群同态 $f$
把
\[
(g_1,\ldots,g_q)\in G^q
\]
送到
\[
(f(g_1),\ldots,f(g_q))\in (G')^q.
\]
这给出标准分解之间的链映射，因此给出复形同态
\[
\operatorname{Hom}_{G'}(X_\bullet(G'),A)\to
\operatorname{Hom}_{G}(X_\bullet(G),A).
\]
取上同调得到换群同态
\[
f^\ast:H^n(G',A)\to H^n(G,A).
\]
在非齐性 $n$-上链语言中，这个映射就是预合成：
\[
(f^\ast c)(g_1,\ldots,g_n)=c(f(g_1),\ldots,f(g_n)).
\]
这个公式说明边界会被送到边界，上闭链会被送到上闭链。

还有一个与极限有关的基本兼容性。若有限群形成逆系统
\[
G=\varprojlim_i G_i
\]
而离散模形成直系统
\[
A=\varinjlim_i A_i
\]
且作用相容，则连续上链只依赖某个有限商 $G_i$，值也落在某个 $A_i$ 中。因此
\[
H^n\!\left(\varprojlim_iG_i,\varinjlim_iA_i\right)
\simeq
\varinjlim_i H^n(G_i,A_i).
\]
这个同构的用处是：以后处理射影有限群时，可以先在有限商上构造同态，再通过极限得到
连续上同调的同态。

当 $H\le G$ 且 $f:H\hookrightarrow G$ 时，称为限制同态：
\[
\operatorname{Res}^n:H^n(G,A)\to H^n(H,A).
\]
它的非齐性公式是把 $G^n$ 上的上链限制到 $H^n$：
\[
(\operatorname{Res}^n c)(h_1,\ldots,h_n)=c(h_1,\ldots,h_n).
\]
在 $n=0$ 时，Tate 上同调的元素写作 $a+N_GA$，其中 $a\in A^G$。因为
$A^G\subset A^H$，所以
\[
\operatorname{Res}^0(a+N_GA)=a+N_HA.
\]
这个公式良定义：若 $a$ 改变一个 $G$-范数 $N_Gb$，则
\[
N_Gb=\sum_{sH\in G/H}s\,N_Hb
\]
在限制到 $H$ 的 Tate 商中属于 $N_HA$ 的像。

限制同态还与连接同态交换。对短正合列
\[
0\to A'\to A\to A''\to0
\]
有交换图
\[
\begin{tikzcd}
H^n(G,A'')\arrow[r,"\delta"]\arrow[d,"\operatorname{Res}^n"']&
H^{n+1}(G,A')\arrow[d,"\operatorname{Res}^{n+1}"]\\
H^n(H,A'')\arrow[r,"\delta"']&
H^{n+1}(H,A').
\end{tikzcd}
\]
这是因为连接同态由 ``提升上链、取边界、落回子模'' 构造，而限制到子群与这三步逐项相容。
反过来，$\operatorname{Res}^0$ 的公式和这个交换图唯一决定所有次数的限制同态：
用维数移位把任意次数逐步移到 $0$ 次，交换图保证每一步得到同一个候选映射。

反方向还有上限制同态
\[
\operatorname{Cor}^n:H^n(H,A)\to H^n(G,A).
\]
它不是简单地把函数从 $H$ 延拓到 $G$，而是把 $H$-不变量沿陪集求和，转成
$G$-不变量。先在 $0$ 次说明定义。取左陪集代表
\[
G/H=\{s_1H,\ldots,s_mH\}.
\]
对 $a\in A^H$ 定义
\[
N_{G/H}(a)=\sum_i s_ia.
\]
这个元素与代表元的选择无关：若把 $s_i$ 换成 $s_ih_i$，其中 $h_i\in H$，则
$s_ih_ia=s_ia$。它属于 $A^G$，因为任取 $g\in G$，集合
\[
\{gs_1H,\ldots,gs_mH\}
\]
仍是全部左陪集，求和只发生重排。于是定义
\[
\operatorname{Cor}^0(a+N_HA)=N_{G/H}(a)+N_GA.
\]
若 $a$ 改变一个 $H$-范数 $N_Hb$，则
\[
N_{G/H}(N_Hb)=N_Gb,
\]
所以上式在商群中良定义。

高次数的上限制同态由两条性质唯一刻画：第一，零次就是刚才的陪集求和；第二，对任一短正合列
\[
0\to A'\to A\to A''\to0
\]
有交换图
\[
\begin{tikzcd}
H^n(H,A'')\arrow[r,"\delta"]\arrow[d,"\operatorname{Cor}^n"']&
H^{n+1}(H,A')\arrow[d,"\operatorname{Cor}^{n+1}"]\\
H^n(G,A'')\arrow[r,"\delta"']&
H^{n+1}(G,A').
\end{tikzcd}
\]
存在性可以用诱导模和 Shapiro 引理构造；唯一性同样来自维数移位。

限制与上限制的关键关系是
\[
\operatorname{Cor}^n\circ\operatorname{Res}^n=[G:H]\operatorname{id}.
\]
先看 $n=0$。若 $a\in A^G$，则
\[
\operatorname{Cor}^0\operatorname{Res}^0(a+N_GA)
=N_{G/H}(a)+N_GA
=\sum_i s_ia+N_GA
=[G:H]a+N_GA.
\]
这里用到 $a$ 被整个 $G$ 固定。一般次数由维数移位得到。更具体地说，把 $A$ 嵌入诱导模
$A^1=\Z[G]\otimes A$，再令余商为 $A^{[1]}$；连接同态把 $H^0(G,A^{[n]})$ 同构到
$H^n(G,A)$。限制、上限制都与连接同态交换，所以零次的乘指数公式沿着维数移位传到所有
次数。

这个关系对有限群上同调的检测非常有用。由于有限群 $G$ 的 Tate 上同调在每个次数都是
被 $|G|$ 杀掉的挠群，它可以分解为各个素数的准素部分：
\[
H^n(G,A)=\bigoplus_p H^n(G,A)_p.
\]

\begin{proposition}
设 $p$ 为素数，$G_p$ 为 $G$ 的 $p$-Sylow 子群。则
\[
\operatorname{Res}^n:H^n(G,A)_p\to H^n(G_p,A)
\]
是单射，且
\[
\operatorname{Cor}^n:H^n(G_p,A)\to H^n(G,A)_p
\]
是满射。
\end{proposition}

\begin{proof}
设 $m=[G:G_p]$。Sylow 定理给出 $p\nmid m$。在 $p$-准素群 $H^n(G,A)_p$ 上，乘以
$m$ 是自同构：如果某个元素阶为 $p^r$，则存在整数 $u,v$ 使 $um+vp^r=1$，从而
乘以 $u$ 是乘以 $m$ 的逆映射。由上面的乘指数关系，
\[
\operatorname{Cor}^n\operatorname{Res}^n=m\operatorname{id}
\]
在 $H^n(G,A)_p$ 上是自同构。

若 $x\in H^n(G,A)_p$ 且 $\operatorname{Res}^n x=0$，则
\[
mx=\operatorname{Cor}^n\operatorname{Res}^n x=0.
\]
乘以 $m$ 可逆，所以 $x=0$，限制同态在 $p$-准素部分上单射。

再取任意 $y\in H^n(G,A)_p$。因为乘以 $m$ 可逆，可取 $x$ 使 $mx=y$。于是
\[
y=mx=\operatorname{Cor}^n(\operatorname{Res}^n x),
\]
说明 $y$ 属于上限制同态的像。又 $H^n(G_p,A)$ 被 $|G_p|$ 杀掉，故其像落在
$p$-准素部分中，上限制同态到 $H^n(G,A)_p$ 是满射。
\end{proof}

\begin{proposition}
若对所有素数 $p$ 和 $G$ 的 $p$-Sylow 子群 $G_p$ 都有
\[
H^n(G_p,A)=0,
\]
则
\[
H^n(G,A)=0.
\]
\end{proposition}

\begin{proof}
对每个素数 $p$，上一命题给出单射
\[
H^n(G,A)_p\hookrightarrow H^n(G_p,A).
\]
假设右端为零，所以 $H^n(G,A)_p=0$。有限群上同调是被 $|G|$ 杀掉的挠群，因此它等于其
各准素部分的直和。所有准素部分都为零，整个 $H^n(G,A)$ 也为零。
\end{proof}

\subsection{膨胀同态}

限制同态把 $G$ 的上同调类拿到子群 $H$ 上。若 $H$ 是正规子群，还可以反过来从商群
$G/H$ 的上同调类制造 $G$ 的上同调类。这个操作叫膨胀：商群上的上链本来只看陪集，
把它沿投射 $G\to G/H$ 拉回后，就变成 $G$ 上的上链。

设 $H\triangleleft G$。若 $a\in A^H$，则对 $g\in G$ 与 $h\in H$ 有
\[
h(ga)=g(g^{-1}hg)a=ga,
\]
因为 $g^{-1}hg\in H$。所以 $ga\in A^H$，从而 $A^H$ 是 $G/H$-模。商映射
$\pi:G\to G/H$ 给出换群同态
\[
\pi^\ast:H^n(G/H,A^H)\to H^n(G,A^H),
\]
再与包含 $i:A^H\hookrightarrow A$ 诱导的同态复合，得到膨胀同态
\[
 \operatorname{Inf}:H^n(G/H,A^H)\to H^n(G,A).
\]
在非齐性上链上，它的公式是
\[
(\operatorname{Inf}c)(g_1,\ldots,g_n)
=c(g_1H,\ldots,g_nH),
\]
再把取值从 $A^H$ 看作 $A$ 中的元素。

\begin{proposition}
设 $A$ 是 $G$-模，$H\triangleleft G$。
\begin{enumerate}[label=(\arabic*)]
\item 有正合列
\[
 0\to H^1(G/H,A^H)\xrightarrow{\operatorname{Inf}}
 H^1(G,A)\xrightarrow{\operatorname{Res}}H^1(H,A).
\]
\item 若 $H^i(H,A)=0$ 对 $i=0,\ldots,q-1$ 成立，则有正合列
\[
 0\to H^q(G/H,A^H)\xrightarrow{\operatorname{Inf}}
 H^q(G,A)\xrightarrow{\operatorname{Res}}H^q(H,A).
\]
\end{enumerate}
\end{proposition}

\begin{proof}
(1) 先证明膨胀在 $H^1(G/H,A^H)$ 上单射。设
$x:G/H\to A^H$ 是 $1$-上闭链，且膨胀后的类在 $H^1(G,A)$ 中为零。于是存在
$a\in A$，使得对每个 $g\in G$，
\[
 x(gH)=ga-a,
\]
其中右端按 $G$ 对 $A$ 的作用理解。取 $h\in H$，由于 $hH=H$，且 $1$-上闭链满足
$x(H)=0$，得到
\[
0=x(hH)=ha-a.
\]
所以 $a\in A^H$。因此 $x(gH)=g(a)-a$ 已经是商群 $G/H$ 上、取值于 $A^H$ 的边界，
原类为零，膨胀单射。

接着证明像等于限制的核。若 $x:G/H\to A^H$ 是 $1$-上闭链，则对 $h\in H$，
\[
(\operatorname{Res}\operatorname{Inf}x)(h)=x(hH)=x(H)=0.
\]
所以
\[
\operatorname{im}\operatorname{Inf}\subset\ker\operatorname{Res}.
\]

反过来，设 $x:G\to A$ 是 $1$-上闭链，且其限制到 $H$ 后为零类。于是存在 $a\in A$，
使得
\[
x(h)=ha-a\qquad(h\in H).
\]
令 $y(g)=ga-a$，它是一个 $1$-边界；把 $x$ 改成同调的代表
\[
x'=x-y.
\]
则 $x'(h)=0$ 对所有 $h\in H$ 成立。由 $1$-上闭链公式
\[
x'(st)=x'(s)+s\,x'(t)
\]
得到两条关系：
\[
x'(gh)=x'(g)+g\,x'(h)=x'(g),
\]
以及
\[
x'(hg)=x'(h)+h\,x'(g)=h\,x'(g).
\]
由于 $H$ 正规，$hg=g(g^{-1}hg)$，且 $g^{-1}hg\in H$，第一条关系又给出
\[
x'(hg)=x'(g).
\]
与第二条关系比较，得到 $h\,x'(g)=x'(g)$。因此 $x'(g)\in A^H$。第一条关系还说明
$x'(g)$ 只依赖右陪集 $gH$。于是可以定义
\[
z:G/H\to A^H,\qquad z(gH)=x'(g).
\]
该定义良好，且把 $1$-上闭链公式降到商群上：
\[
z(g_1g_2H)
=x'(g_1g_2)
=x'(g_1)+g_1x'(g_2)
=z(g_1H)+g_1H\,z(g_2H).
\]
所以 $z$ 是 $G/H$ 上的 $1$-上闭链，并且 $\operatorname{Inf}z=x'$。由于 $x'$ 与 $x$
相差一个边界，$[x]$ 属于膨胀同态的像。

(2) 用维数移位把高次命题降到刚证明的一次命题。令
\[
B=\Z[G]\otimes A,\qquad C=J_G\otimes A,
\]
其中 $A\to B$ 的映射为 $a\mapsto\sum_{g\in G}g\otimes g^{-1}a$，而 $C$ 是余商。得到短正合列
\[
0\to A\to B\to C\to0.
\]
$B$ 作为 $G$-模是诱导模；限制到 $H$ 后仍是诱导模；取 $H$-不变量后，$B^H$ 作为
$G/H$-模也是诱导模。因此三种上同调
\[
H^i(G,B),\qquad H^i(H,B),\qquad H^i(G/H,B^H)
\]
在非零需要的位置都消失。由假设 $H^i(H,A)=0$ 对 $0\le i\le q-1$ 成立，长正合列给出
\[
H^i(H,C)\simeq H^{i+1}(H,A)=0\qquad(0\le i\le q-2).
\]
同时 $0\to A^H\to B^H\to C^H\to H^1(H,A)=0$ 保证取 $H$-不变量后仍有足够的正合性。

连接同态把关于 $C$ 的 $(q-1)$ 次正合列同构到关于 $A$ 的 $q$ 次正合列。更具体地，有交换图
\[
\begin{tikzcd}[column sep=small]
0\arrow[r]&
H^{q-1}(G/H,C^H)\arrow[r,"\operatorname{Inf}"]\arrow[d,"\delta"']&
H^{q-1}(G,C)\arrow[r,"\operatorname{Res}"]\arrow[d,"\delta"]&
H^{q-1}(H,C)\arrow[d,"\delta"]\\
0\arrow[r]&
H^q(G/H,A^H)\arrow[r,"\operatorname{Inf}"']&
H^q(G,A)\arrow[r,"\operatorname{Res}"']&
H^q(H,A).
\end{tikzcd}
\]
当 $q=1$ 时，第二行正合已由 (1) 证明。若 $q>1$，归纳假设给出第一行正合；三条竖直连接
同态在所需位置为同构，于是第二行也正合。这完成高次情形。
\end{proof}

\section{张量积}

这一节的目标，是把 ``上同调类可以相乘'' 这件事做成一个严格、可追踪、能和限制/
上限制/连接同态兼容的运算。它的数论意义很直接：后面构造类成基本类、证明 Tate 定理、
比较局部和整体类域论时，都要反复把一个上同调类与另一个配对，再看结果落到哪里。

设 $A,B$ 为 $G$-模。张量积 $A\otimes B$ 取对角作用：
\[
 g(a\otimes b)=ga\otimes gb.
\]
这不是随便选的作用，而是唯一把两边的群作用同时保留下来的自然选择：$g$ 先作用在
左因子，再作用在右因子，最后得到同一个张量项。

若 $a\in A^G$、$b\in B^G$，则 $a\otimes b\in(A\otimes B)^G$。在 Tate 的 $H^0$ 中得到
\[
 H^0(G,A)\times H^0(G,B)\to H^0(G,A\otimes B),
\qquad
 (\bar a,\bar b)\mapsto \overline{a\otimes b}.
\]
这就是杯积在零维的起点。

\begin{theorem}
存在唯一一组双线性映射
\[
 \cup:H^p(G,A)\times H^q(G,B)\to H^{p+q}(G,A\otimes B)
\]
满足：
\begin{enumerate}[label=(\arabic*)]
\item $p=q=0$ 时由张量积诱导。
\item 对 $A$ 变量的短正合列，杯积与连接同态相容：
\[
 \delta(a''\cup b)=\delta a''\cup b.
\]
\item 对 $B$ 变量的短正合列，杯积与连接同态满足符号规则：
\[
 \delta(a\cup b'')=(-1)^p a\cup\delta b''.
\]
\end{enumerate}
\end{theorem}

\begin{proof}
先在一个变量次数为 $0$ 时直接定义。一般次数用维数移位：把 $a\in H^p(G,A)$ 拉回到 $H^0(G,A^{[p]})$，把 $b\in H^q(G,B)$ 拉回到 $H^0(G,B^{[q]})$，在零维取张量积，再通过维数移位送到 $H^{p+q}(G,A\otimes B)$。符号 $(-1)^p$ 来自两个连接同态交换时的反交换规则。唯一性来自每个次数都能由 $H^0$ 经过维数移位得到。
\end{proof}

\begin{proposition}
杯积满足以下性质：
\begin{enumerate}[label=(\arabic*)]
\item 交换律带符号：
\[
 a^p\cup b^q=(-1)^{pq}b^q\cup a^p.
\]
\item 结合律：
\[
 (a^p\cup b^q)\cup c^r=a^p\cup(b^q\cup c^r).
\]
\item 对 $G$-模同态函子性相容。
\item 对子群限制和上限制相容：
\[
 \operatorname{Res}(a\cup b)=\operatorname{Res}a\cup\operatorname{Res}b,
\]
\[
 \operatorname{Cor}(\operatorname{Res}a\cup b)=a\cup\operatorname{Cor}b.
\]
\end{enumerate}
\end{proposition}

\begin{proof}
先解释为什么要有这些性质。零维时，$a\otimes b$ 的定义应当只依赖不变量本身，而不能依赖
代表元，所以它必须对称地接受张量积的自然性质。高次时，杯积是把两个上链先拉到零维，
在零维做张量，再搬回去，因此所有兼容性都应由维数移位继承。

交换律的符号来自链复形的反交换性。若 $a$ 的次数是 $p$，$b$ 的次数是 $q$，把它们拉到
零维时要分别经过 $p$ 次和 $q$ 次连接同态。交换这两段连接同态时，每跨过一次会积出一个
$-1$，总共有 $pq$ 次跨越，所以得到
\[
 a^p\cup b^q=(-1)^{pq}b^q\cup a^p.
\]
结合律同样先在零维由张量积的结合律得到，再由维数移位传回一般次数。

函子性是说：先对系数做模同态，再做杯积，和先做杯积再把结果推过去，得到同一个类。
这是因为整个构造只用了链映射与张量积，而它们本身就是函子性的。

限制相容性最直接：把代表上链限制到子群，张量后再限制，和先张量再限制没有区别。
上限制相容性稍微更像“求和交换”：
\[
\operatorname{Cor}(\operatorname{Res}a\cup b)=a\cup\operatorname{Cor}b.
\]
它的本质是陪集求和与张量乘积可交换。
\end{proof}

\begin{proposition}
设 $a^p\in H^p(G,A)$。在同构 $G^{\operatorname{ab}}\simeq H^{-2}(G,\Z)$ 下，把 $\sigma$ 的像记为 $\bar\sigma$。则
\[
 a^0\cup b^q=a_0\otimes b_q,\qquad
 a^p\cup b^0=a_p\otimes b_0,
\]
\[
 a^1\cup b^{-1}=\sum_{\tau\in G}a_1(\tau)\otimes \tau b_{-1},
\]
\[
 a^1\cup\bar\sigma=a_1(\sigma),
\qquad
 a^2\cup\bar\sigma=\sum_{\tau\in G}a_2(\tau,\sigma).
\]
\end{proposition}

\begin{proof}
这一组公式的意义，是把抽象的杯积和具体的链级表达式对上号。后面在类成理论里，基本类、
Frobenius、范数和特征标的计算都靠这些低维公式。

第一式是定义。第二式把 $a^1$ 通过正合列
\[
 0\to A\to \Z[G]\otimes A\to J_G\otimes A\to0
\]
降到零维，计算连接同态后得到求和公式。更具体地说，$a^1$ 由一个 1-上闭链 $a_1$ 表示，
而 $b^{-1}$ 在 $H^{-1}$ 中可以看成某个被范数杀掉的元素。杯积后在零维出现的就是把
每个群元素的作用逐项加起来，所以得到
\[
 a^1\cup b^{-1}=\sum_{\tau\in G}a_1(\tau)\otimes \tau b_{-1}.
\]

后两式把 $\bar\sigma$ 看成 $H^{-2}(G,\Z)$ 中的类，再用增广理想的连接同态计算。
具体地说，$\sigma$ 在 $G^{\mathrm{ab}}$ 中的像对应于 $H^{-2}(G,\Z)$ 的一个生成元；
把它和 $a^1$、$a^2$ 杯积，就是把 $\sigma$ 代入相应的 1-上闭链或 2-上闭链：
\[
 a^1\cup\bar\sigma=a_1(\sigma),\qquad
 a^2\cup\bar\sigma=\sum_{\tau\in G}a_2(\tau,\sigma).
\]
这类公式后面会不断出现，所以这里把它们写成可直接代入的链级形式。
\end{proof}

\begin{proposition}
任意 $G$-模配对 $A\times B\to C$ 都诱导杯积
\[
 H^p(G,A)\times H^q(G,B)\to H^{p+q}(G,C),
\]
并且满足函子性、连接同态相容性和符号规则。
\end{proposition}

\begin{proof}
配对 $A\times B\to C$ 等价于 $G$-模同态 $A\otimes B\to C$。先用定理 5.23 得到
$A\otimes B$ 中的杯积，再复合同态 $A\otimes B\to C$。这样做的好处是：一旦配对被视作
张量积上的模同态，所有相容性就不必重新证明，而是从定理 5.23 的三条性质直接继承。
\end{proof}

\begin{proposition}
设 $G$ 为有限群，$A$ 为 $G$-模。则配对
\[
 \operatorname{Hom}(A,\Q/\Z)\times A\to\Q/\Z
\]
诱导同构
\[
 H^i(G,\operatorname{Hom}(A,\Q/\Z))
 \simeq \operatorname{Hom}(H^{-i-1}(G,A),\Q/\Z).
\]
若 $A$ 是自由 $\Z$-模，则配对
\[
 \operatorname{Hom}(A,\Z)\times A\to\Z
\]
诱导
\[
 H^i(G,\operatorname{Hom}(A,\Z))
 \simeq \operatorname{Hom}(H^{-i}(G,A),\Q/\Z).
\]
\end{proposition}

\begin{proof}
第一式在 $i=0$ 可直接由 $H^0$ 和 $H^{-1}$ 的定义验证。一般 $i$ 用维数移位：
把 $A$ 先移成 $A^{[i]}$，再把杯积推到零维，最后再移回来。因为对偶配对
\[
 \operatorname{Hom}(A,\Q/\Z)\times A\to\Q/\Z
\]
把左边元素看成把右边元素送入 $\Q/\Z$ 的线性函数，所以
$H^i(G,\operatorname{Hom}(A,\Q/\Z))$ 的元素正好对应于
$H^{-i-1}(G,A)$ 的对偶函数。

第二式来自短正合列
\[
 0\to\operatorname{Hom}(A,\Z)\to\operatorname{Hom}(A,\Q)
 \to\operatorname{Hom}(A,\Q/\Z)\to0.
\]
由于 $\operatorname{Hom}(A,\Q)$ 是唯一可除群，命题 5.13 给出其上同调为零，连接同态把第二式
化为第一式。直观上，这是因为 $\Q$ 是“无障碍”的中介，$\Q/\Z$ 只记录有限阶信息。
\end{proof}

\section{Tate 定理}

\begin{lemma}
若对所有子群 $H\subset G$ 都有
\[
 H^1(H,A)=0=H^2(H,A),
\]
则对所有子群 $H$ 都有
\[
 H^0(H,A)=0=H^3(H,A).
\]
\end{lemma}

\begin{proof}
这一引理是 Tate 定理的起手式：把“低次两个相邻群都消失”推进到“再往前再往后各消失
两个”。它告诉我们，一旦某个二次条件成立，就能把同调的整个时间轴往两边推开。

先对 $|G|$ 归纳。若 $G$ 不是 $p$-群，则取任意 Sylow 子群 $G_p$。假设所有真子群都已知
满足结论，则 $G_p$ 的上同调满足 $H^1(G_p,A)=H^2(G_p,A)=0$。由命题 5.21，
全群的 $H^0$ 与 $H^3$ 都为零。

若 $G$ 是 $p$-群，则存在正规子群 $H\triangleleft G$ 使 $G/H$ 为素数阶循环群。
对 $H$ 代入假设，得到
\[
H^1(H,A)=H^2(H,A)=0.
\]
膨胀-限制正合列给出
\[
0\to H^1(G/H,A^H)\to H^1(G,A)\to H^1(H,A)=0,
\]
所以 $H^1(G/H,A^H)\simeq H^1(G,A)=0$。
再看 $H^2$，同样得到
\[
H^2(G/H,A^H)\simeq H^2(G,A)=0.
\]
现在对循环群 $G/H$ 作用在 $A^H$ 上应用周期性：
\[
H^3(G/H,A^H)\simeq H^1(G/H,A^H)=0,\qquad
H^0(G/H,A^H)\simeq H^2(G/H,A^H)=0.
\]
把这两项通过膨胀拉回 $G$，得到 $H^3(G,A)=0$ 和 $H^0(G,A)=0$。
\end{proof}

\begin{proposition}
若存在 $q_0$，使对所有子群 $H\subset G$ 都有
\[
 H^{q_0}(H,A)=0=H^{q_0+1}(H,A),
\]
则对所有子群 $H$ 和所有整数 $q$，都有
\[
 H^q(H,A)=0.
\]
\end{proposition}

\begin{proof}
这一步是“消失区间外推”的纯形式版本。把 $A$ 维数移位到 $A^{[q_0-1]}$，则假设
\[
H^1(H,A^{[q_0-1]})\simeq H^{q_0}(H,A)=0,\qquad
H^2(H,A^{[q_0-1]})\simeq H^{q_0+1}(H,A)=0.
\]
由引理 5.28，得到
\[
H^0(H,A^{[q_0-1]})=0,\qquad H^3(H,A^{[q_0-1]})=0.
\]
再翻回原来的次数，便是
\[
H^{q_0-1}(H,A)=0,\qquad H^{q_0+2}(H,A)=0.
\]
把这个步骤重复应用：一旦知道某两个相邻次数消失，就能把零区间向左、向右各推进一格。
于是所有整数次数都被扫空。
\end{proof}

\begin{proposition}
设 $u\in H^2(G,A)$。若对所有子群 $H\subset G$，杯积
\[
 u_H\cup -:H^q(H,\Z)\to H^{q+2}(H,A)
\]
在 $q=0,1$ 时是同构，则它在所有 $q$ 时都是同构。
\end{proposition}

\begin{proof}
令
\[
f_q=u_H\cup-:H^q(H,\Z)\to H^{q+2}(H,A).
\]
把它的“失效部分”记为 $B$，也就是把源和目标之间的映射锥取上同调后留下来的差对象。
如果 $q=0,1$ 时 $f_q$ 是同构，那么这个差对象在两个相邻次数上为零。命题 5.29 说：
只要对所有子群在两个相邻次数上同时消失，就会推出所有次数都消失。所以差对象全为零，
也就说明所有 $f_q$ 都是同构。
\end{proof}

\begin{theorem}
设 $A$ 为 $G$-模。若存在 $u\in H^2(G,A)$，使得对 $G$ 的任意子群 $H$：
\[
 |H|\cdot H^0(H,A)=0,\qquad
 H^1(H,A)=0,
\]
并且 $u_H$ 诱导
\[
 H^0(H,\Z)\simeq H^2(H,A),
\]
则对所有 $q$ 都有同构
\[
 H^q(H,\Z)\xrightarrow{\;u_H\cup-\;}H^{q+2}(H,A).
\]
\end{theorem}

\begin{proof}
这个定理把前面所有抽象命题压缩成类成理论的最终形式。条件
\[
|H|\cdot H^0(H,A)=0,\qquad H^1(H,A)=0
\]
正是命题 5.30 的低维假设：零次说明范数商是循环的，一级说明没有一阶障碍。
而 $u_H$ 在 $H^2(H,A)$ 中给出的基本类把二次信息固定下来。

因此命题 5.30 直接告诉我们，杯积
\[
u_H\cup-:H^q(H,\Z)\to H^{q+2}(H,A)
\]
在 $q=0,1$ 已经是同构，就会自动延伸到全部整数 $q$。这正是 Tate 定理在类成模上的
抽象版本。
\end{proof}

\section{射影有限群的同调群}

有限 Galois 理论里，真正出现的群往往不是单个有限群，而是许多有限群组成的逆极限。
绝对 Galois 群就是这样的对象。它的每个有限商都对应一个有限 Galois 扩张，而连续上同调
要做的，就是把这些有限层上的 cohomology 拼成一个稳定的极限对象。

射影有限群是有限群的逆极限。Galois 群 $\operatorname{Gal}(F^{\operatorname{sep}}/F)$ 通常无限，但它的有限商正是有限 Galois 扩张的 Galois 群。因此连续上同调必须允许“通过有限商计算”。

\begin{proposition}
设 $G$ 为射影有限群，$A$ 为离散 $G$-模。以下条件等价：
\begin{enumerate}[label=(\arabic*)]
\item 作用映射 $G\times A\to A$ 连续。
\item 对每个 $a\in A$，稳定子 $G_a$ 是开子群。
\item $A=\bigcup_U A^U$，其中 $U$ 遍历 $G$ 的开正规子群。
\end{enumerate}
\end{proposition}

\begin{proof}
离散拓扑下，集合开当且仅当它是任意子集的并，所以映射连续等价于每个点的原像开。
对固定 $a\in A$，映射
\[
G\to A,\qquad \sigma\mapsto \sigma a
\]
的原像 $\{ \sigma:\sigma a=a\}$ 正是稳定子 $G_a$。因此 (1)$\Leftrightarrow$(2)。

若 (2) 成立，则对每个 $a$，$G_a$ 开，故含某个开正规子群 $U$，于是 $a\in A^U$，
从而 $A=\bigcup_U A^U$。反过来，若 $A$ 是这些不动子群的并，则每个 $a$ 都被某个开正规子群
固定，因此稳定子开。这样 (2)$\Leftrightarrow$(3)。
\end{proof}

\begin{definition}
设 $G$ 为射影有限群，$A$ 为连续离散 $G$-模。定义连续上同调
\[
 H^q(G,A)=\varinjlim_U H^q(G/U,A^U),
\]
其中 $U$ 遍历 $G$ 的开正规子群，转移映射由膨胀同态给出。
\end{definition}

\begin{proposition}
连续 $G$-模 $A$ 有平凡上同调，当且仅当对所有开正规子群 $U$，有限群 $G/U$-模 $A^U$ 有平凡上同调。
\end{proposition}

\begin{proof}
定义把连续上同调写成有限层上同调的直极限。若每一有限层都为零，直极限自然为零。
反过来，给定某个连续类，它只出现在某个有限层里，因为上链取值离散且连续，
所以它会下降到某个 $G/U$ 上的类。若所有连续类都消失，那么这个有限层类也必须消失，
于是每个有限层都平凡。
\end{proof}

\section{类成}

现在把 Tate 定理翻译成类域论需要的语言。这里的关键不是再造一个新理论，而是把前面
关于有限群、循环群、杯积、基本类和维数移位的所有结果，装进 Galois 群的逆极限框架里。
这样一来，范数、包含、投射和共轭就都可以用同一个杯积公式来读出。

设 $G_F$ 是某个域 $F$ 的绝对 Galois 群，$G_K\triangleleft G_F$ 对应有限 Galois 扩张 $K/F$。令
\[
 G(K/F)=G_F/G_K.
\]

\begin{definition}
称连续 $G_F$-模 $A$ 满足类成原则，如果对每个有限 Galois 扩张 $K/F$，存在同构
\[
 \operatorname{inv}_{K|F}:H^2(G(K/F),A^{G_K})
 \simeq \frac1{[K:F]}\Z/\Z,
\]
并且这些同构与中间域的限制、上限制、膨胀相容；同时 $H^1(G(K/F),A^{G_K})=0$。
\end{definition}

这个定义的意思是：$A$ 的二阶上同调由一个规范的循环群控制，而且一阶障碍消失。
一旦 $H^1=0$，二阶类就不再有“非平凡扭曲”；一旦 $H^2$ 被 $\Z/\Z$ 型商控制，
就可以把一个基本类选出来，作为所有计算的统一起点。局部类域论中 $A=\bar F^\times$，
整体类域论中 $A$ 是理元类群的极限版本。

\begin{theorem}
若连续 $G_F$-模 $A$ 满足类成原则，且 $G_K\triangleleft G_F$，则基本类
\[
 u_{K|F}\in H^2(G(K/F),A^{G_K})
\]
决定同构
\[
 H^q(G(K/F),\Z)\simeq H^{q+2}(G(K/F),A^{G_K})
\]
对所有 $q$ 成立。
\end{theorem}

\begin{proof}
取 $u_{K|F}$ 为在不变量同构下对应 $1/[K:F]+\Z$ 的类。类成原则给出 Tate 定理 5.31 的低维条件：
\[
H^1=0,\qquad H^2\cong \frac1{[K:F]}\Z/\Z
\]
并且这些类与中间域相容。于是命题 5.31 中的基本类条件成立，杯积
\[
u_{K|F}\cup-
\]
把所有次数的群都推到下一层，给出全部次数的同构。
\end{proof}

由定理 5.36，取 $q=-2$ 得
\[
 H^{-2}(G(K/F),\Z)\simeq H^0(G(K/F),A^{G_K}).
\]
而
\[
 H^{-2}(G,\Z)\simeq G^{\operatorname{ab}},
\]
所以得到从 $A^{G_F}$ 到 $G(K/F)^{\operatorname{ab}}$ 的映射。这就是范剩符号的来源。

\begin{proposition}
设 $A$ 满足类成原则。若 $a\in A^{G_F}$，$\chi\in H^1(G(K/F),\Q/\Z)$，则
\[
 \chi((a,K|F))
 =
 \operatorname{inv}_{K|F}\bigl(\bar a\cup\delta\chi\bigr),
\]
其中 $\delta:H^1(G(K/F),\Q/\Z)\to H^2(G(K/F),\Z)$ 来自
\[
 0\to\Z\to\Q\to\Q/\Z\to0.
\]
若 $E/F$、$K/F$ 为相容的正规扩张，则范剩符号与自然投射相容。
\end{proposition}

\begin{proof}
由定理 5.36，$\bar a$ 可写成
\[
 \bar a=u_{K|F}\cup \bar\sigma_a
\]
其中 $\bar\sigma_a\in H^{-2}(G(K/F),\Z)$ 对应 $\sigma_a=(a,K|F)$。于是
\[
 \bar a\cup\delta\chi
 =
 u_{K|F}\cup(\bar\sigma_a\cup\delta\chi).
\]
低维杯积公式给出
\[
 \bar\sigma_a\cup\chi=\chi(\sigma_a)\in H^{-1}(G,\Q/\Z).
\]
再连接到 $H^0(G,\Z)$，并用 $\operatorname{inv}(u_{K|F})=1/[K:F]+\Z$，得到公式。自然性来自杯积、
膨胀和限制同态的相容性。
\end{proof}

\begin{proposition}
范剩符号对中间域、范映射、包含映射和共轭作用相容。
\end{proposition}

\begin{proof}
四种相容性都来自同一个机制：基本类在限制、上限制、膨胀下相容，而杯积与这些同态相容。
对中间域 $K/E/F$，上限制对应范映射 $N_{E/F}$，限制对应群的迁移映射，膨胀对应 Galois 群
自然投射，共轭对应把扩张整体搬到共轭扩张。把这些同态放入 Tate 同构方块，交换性给出四个图。
\end{proof}

\begin{proposition}
域 $F$ 的范剩符号
\[
 (\ ,F):A^{G_F}\to G_F^{\operatorname{ab}}
\]
的核为
\[
 \bigcap_E N_{E|F}A^{G_E},
\]
像是 $G_F^{\operatorname{ab}}$ 的稠密子群。
\end{proposition}

\begin{proof}
元素 $a$ 映到 $1$ 当且仅当它在每个有限正规层 $E/F$ 的范剩符号中为 $1$。有限层的核是范数像 $N_{E|F}A^{G_E}$，交起来得到核。

稠密性按 profinite 拓扑检验。$G_F^{\operatorname{ab}}$ 的基本邻域由有限商 $G(E/F)^{\operatorname{ab}}$ 的纤维给出。每个有限层的范剩符号由 Tate 同构得到满射，所以任意有限商中的目标都可由某个 $a$ 达到。因此全局像在所有有限商中满，故稠密。
\end{proof}

\section{域的上同调}

这一节把前面关于有限群的结论重新放回到具体域扩张里。加法群上的正规基解释了为什么
有限 Galois 扩张的上同调会消失；乘法群上的 Hilbert 90 则告诉我们，真正非平凡的部分
不是 1-上闭链本身，而是它们是否能写成一个全局元素的商。

\begin{proposition}
设 $E/F$ 是有限 Galois 扩张。则对所有 $q$，
\[
 H^q(\operatorname{Gal}(E/F),E)=0,
\]
其中 $E$ 按加法群看作 Galois 模。
\end{proposition}

\begin{proof}
正规基定理给出 $a\in E$，使 $\{ga:g\in G\}$ 是 $F$-基。因此
\[
 E=\bigoplus_{g\in G}F\cdot ga
\]
是诱模。命题 5.4 说明诱模有平凡上同调。
\end{proof}

\begin{theorem}
设 $E/F$ 是有限 Galois 扩张。则
\[
 H^1(\operatorname{Gal}(E/F),E^\times)=1.
\]
\end{theorem}

\begin{proof}
这一结论是 Hilbert 90 的上同调形式。它的意思是：乘法 1-上闭链没有真正的新障碍，
每个上闭链都能写成一个“坐标商” $b/\tau(b)$。这正是以后处理循环扩张和 Kummer 理论时
最常用的工具。

记 $G=\operatorname{Gal}(E/F)$。取 $1$-上闭链 $a_\sigma\in E^\times$，即
\[
 a_{\sigma\tau}=a_\sigma\,\sigma(a_\tau).
\]
要证明它是边界。对 $c\in E$ 定义
\[
 b=\sum_{\sigma\in G}a_\sigma\,\sigma(c).
\]
由 Dedekind 自同构线性无关性，可选 $c$ 使 $b\ne0$。对 $\tau\in G$，
\[
 \tau(b)=\sum_\sigma \tau(a_\sigma)\,\tau\sigma(c).
\]
利用上闭链关系 $a_{\tau\sigma}=a_\tau\,\tau(a_\sigma)$，得
\[
 \tau(a_\sigma)=a_\tau^{-1}a_{\tau\sigma}.
\]
代入并重排求和，得到
\[
 \tau(b)=a_\tau^{-1}b.
\]
于是
\[
 a_\tau=b/\tau(b),
\]
这正是一个乘法边界。故 $H^1(G,E^\times)=1$。
\end{proof}

\begin{corollary}
若 $E/F$ 是循环有限扩张，$\sigma$ 生成 $G$，且 $x\in E^\times$ 满足
\[
 N_{E/F}(x)=1,
\]
则存在 $c\in E^\times$ 使
\[
 x=\sigma(c)/c.
\]
\end{corollary}

\begin{proof}
这里的条件 $N_{E/F}(x)=1$ 正好是 $H^{-1}$ 的定义中“被范数杀掉”的部分。
循环群周期性给出
\[
 H^{-1}(G,E^\times)\simeq H^1(G,E^\times)=1.
\]
而
\[
 H^{-1}(G,E^\times)
 =
 \{x:N_{E/F}x=1\}/\{ \sigma(c)c^{-1}:c\in E^\times\}.
\]
故范数为 $1$ 的元素必为 $\sigma(c)/c$。
\end{proof}

\begin{corollary}
若 $E/F$ 是有限域扩张，则对所有 $q$，
\[
 H^q(\operatorname{Gal}(E/F),E^\times)=1.
\]
\end{corollary}

\begin{proof}
有限域扩张的 Galois 群是循环群，$E^\times$ 是有限群。Hilbert 90 给出 $H^1=1$。Herbrand 商为 $1$，所以 $H^2=1$。由循环群周期性，所有次数都为 $1$。
\end{proof}

\section{Kummer 扩张}

Kummer 理论的角色很清楚：当域里已经有足够多的单位根时，开 $m$ 次根的扩张就完全由
Galois 群的特征标和子群结构控制。它是后面 Brauer 群、循环代数和类域论之间最直接的
桥梁之一。

设 $F$ 的特征不整除正整数 $m$，并且 $F$ 含有全体 $m$ 次单位根 $\mu_m$。记
\[
 E^{\times(m)}=\{x^m:x\in E^\times\}.
\]

\begin{definition}
称域扩张 $E/F$ 为 $m$-Kummer 扩张，如果：
\begin{enumerate}[label=(\arabic*)]
\item $\operatorname{Gal}(E/F)$ 是指数整除 $m$ 的 Abel 群；
\item $F$ 含 $\mu_m$；
\item $\operatorname{char}F\nmid m$。
\end{enumerate}
\end{definition}

\begin{proposition}
设 $F$ 含 $\mu_m$，$m'|m$，$E/F$ 为有限 $m'$-Kummer 扩张。则：
\begin{enumerate}[label=(\arabic*)]
\item 对 $c^m=a\in F^\times$，定义
\[
 \chi_c(\sigma)=\sigma(c)/c\in\mu_m.
\]
得到正合列
\[
 1\to F^\times\to \sqrt[m]{F}\cap E^\times
 \xrightarrow{\chi}\operatorname{Hom}(\operatorname{Gal}(E/F),\mu_m)\to1.
\]
\item 有同构
\[
 (E^{\times(m)}\cap F^\times)/F^{\times(m)}
 \simeq \operatorname{Gal}(E/F)^\vee
 \simeq \operatorname{Gal}(E/F).
\]
\item $E=F(\sqrt[m]{a_1},\ldots,\sqrt[m]{a_r})$ 当且仅当这些根在
\[
 (\sqrt[m]{F}\cap E^\times)/F^\times
\]
中生成。
\item 若 $a_iF^{\times(m)}$ 生成 $(E^{\times(m)}\cap F^\times)/F^{\times(m)}$，则 $E$ 是
\[
 \prod_i(X^m-a_i)
\]
的分裂域。
\end{enumerate}
\end{proposition}

\begin{proof}
Kummer 理论的本质是把“根的集合”换成“特征标的集合”。每一个 $m$ 次根都贡献一个
$\mu_m$-值特征标，反过来特征标也都来自某个根。这就是根与群之间的双射。

(1) 若 $c^m\in F^\times$，则
\[
 \sigma(c)^m=c^m,
\]
故 $\sigma(c)/c\in\mu_m$。因为 $\mu_m\subset F$，该映射是特征标。若特征标平凡，则 $c$ 被所有 $\sigma$ 固定，故 $c\in F^\times$，核正是 $F^\times$。

满射用 Hilbert 90。给定 $\chi:G\to\mu_m$，它是 $E^\times$ 中的 $1$-上闭链。Hilbert 90 给出 $c\in E^\times$，使
\[
 \chi(\sigma)=\sigma(c)/c.
\]
于是 $\sigma(c^m)=c^m$，故 $c^m\in F^\times$。

(2) 映射 $c\mapsto c^mF^{\times(m)}$ 的核为 $F^\times$，所以第一步给出同构到特征标群。有限 Abel 群与其 $\mu_m$-值特征标群非典范同构。

(3) 若 $E$ 由这些根生成，则固定所有根的 Galois 元素只能是恒等元；因此对应特征标生成整个特征标群。由 (1) 得根类生成。反向，若根类生成，则任何固定这些根的 Galois 元素被所有特征标杀掉，故为恒等元，扩张域就是 $E$。

(4) 由 (2)(3)，选择这些 $a_i$ 的 $m$ 次根生成 $E$。因为 $\mu_m\subset F$，$X^m-a_i$ 的全部根都在加入一个根后得到，所以 $E$ 正是乘积多项式的分裂域。
\end{proof}

\begin{theorem}
设 $F$ 含 $\mu_m$。令 $\mathcal G$ 为满足
\[
 F^{\times(m)}\subset\Gamma\subset F^\times,\qquad
 |\Gamma/F^{\times(m)}|<\infty
\]
的子群集合；令 $\mathcal F$ 为所有 $m'$-Kummer 扩张 $E/F$，其中 $m'|m$。则
\[
 \Gamma\mapsto F(\sqrt[m]{\Gamma}),
\qquad
 E\mapsto E^{\times(m)}\cap F^\times
\]
互为逆映射。
\end{theorem}

\begin{proof}
这一定理是 Kummer 理论的分类版：有限指数子群对应有限开根号扩张。
取 $\Gamma$，选 $a_1,\ldots,a_r$ 使其模 $F^{\times(m)}$ 生成 $\Gamma/F^{\times(m)}$。由于 $F$ 含 $\mu_m$ 且特征不整除 $m$，多项式
\[
 \prod_i(X^m-a_i)
\]
可分，分裂域是 Galois 扩张。任一 Galois 元素只把 $\sqrt[m]{a_i}$ 乘以 $m$ 次单位根，因此 Galois 群 Abel 且指数整除 $m$。所以得到 Kummer 扩张。

从 Kummer 扩张 $E$ 出发，命题 5.45 说明 $E$ 由 $E^{\times(m)}\cap F^\times$ 中元素的 $m$ 次根生成，因此回到同一个 $E$。从 $\Gamma$ 出发，命题 5.45(2) 又说明分裂域中落回 $F^\times$ 的 $m$ 次幂恰为 $\Gamma$。故两映射互逆。
\end{proof}

\section*{习题详解}
\addcontentsline{toc}{section}{习题详解}

\begin{exercise}
设 $A,B$ 为 $G$-模，$H$ 为 $G$ 的子群。取 $a\in A^G$、$b\in B^H$。证明：
\[
 (\operatorname{Res}\bar a)\cup\bar b
 =a\otimes b+N_H(A\otimes B),
\]
\[
 \operatorname{Cor}((\operatorname{Res}\bar a)\cup\bar b)
 =
 \sum_{y\in G/H}a\otimes yb+N_G(A\otimes B)
 =
 \bar a\cup\operatorname{Cor}\bar b.
\]
\end{exercise}

\begin{proof}[解答]
本题可解。第一式只是在 $H^0$ 中展开杯积定义。限制后 $\operatorname{Res}\bar a$ 由同一个 $a\in A^H$ 表示，故
\[
 (\operatorname{Res}\bar a)\cup\bar b
 =a\otimes b+N_H(A\otimes B).
\]

第二式用 $H^0$ 中上限制的定义。若 $G/H$ 的代表为 $y$，则
\[
 \operatorname{Cor}(a\otimes b+N_H(A\otimes B))
 =
 \sum_{y\in G/H}y(a\otimes b)+N_G(A\otimes B).
\]
因为 $a\in A^G$，$y(a\otimes b)=a\otimes yb$。另一方面
\[
 \operatorname{Cor}\bar b=\sum_{y\in G/H}yb+N_GB,
\]
再与 $\bar a$ 做零维杯积，得到同一个元素。
\end{proof}

\begin{exercise}
证明范剩符号
\[
 (\ ,F):A^{G_F}\to G_F^{\operatorname{ab}}
\]
的核为
\[
 \bigcap_{G_E\triangleleft G_F}N_{E|F}A^{G_E},
\]
像是 $G_F^{\operatorname{ab}}$ 的稠密子集。
\end{exercise}

\begin{proof}[解答]
这是命题 5.39 的展开。若 $a$ 在全局范剩符号下为 $1$，则它在每个有限正规商
\[
 G(E/F)^{\operatorname{ab}}
\]
中的像都为 $1$。有限层类成同构告诉我们，有限层范剩符号的核正是 $N_{E|F}A^{G_E}$。所以 $a$ 属于所有这些范数像的交。

反过来，若 $a$ 属于所有 $N_{E|F}A^{G_E}$，则每个有限层范剩符号都为 $1$，故在逆极限 $G_F^{\operatorname{ab}}$ 中也为 $1$。

证明稠密性。profinite 群中，像稠密等价于在每个有限连续商中像为满。对任一有限正规扩张 $E/F$，有限层范剩符号
\[
 A^{G_F}\to G(E/F)^{\operatorname{ab}}
\]
由 Tate 同构给出满射。因此全局像在所有有限商上满，故稠密。
\end{proof}

\begin{exercise}
设 $F$ 为域，$G=\operatorname{Gal}(F^{\operatorname{sep}}/F)$。
\begin{enumerate}[label=(\arabic*)]
\item 证明 $\operatorname{Gal}(F^{\operatorname{sep}}/F)\simeq \operatorname{Aut}_F(F^{\operatorname{alg}})$。
\item 证明从 $G^m$ 到离散集 $S$ 的映射连续，当且仅当它局部正则。
\item 用连续因子集描述
\[
 H^2(G,F^{\operatorname{sep}\times}).
\]
\end{enumerate}
\end{exercise}

\begin{proof}[解答]
这一组题把连续 Galois 群、局部正则和连续因子集三件事串起来，正是后面看 Brauer 群
和二阶上同调时要用的语言。

(1) $F^{\operatorname{alg}}/F^{\operatorname{sep}}$ 是纯不可分扩张。纯不可分扩张中，每个元素只有唯一的共轭根，所以 $F^{\operatorname{sep}}$ 的 $F$-自同构至多有一种方式延拓到 $F^{\operatorname{alg}}$。存在性由代数闭包延拓定理给出，唯一性由纯不可分性给出。因此限制映射
\[
 \operatorname{Aut}_F(F^{\operatorname{alg}})
 \to \operatorname{Gal}(F^{\operatorname{sep}}/F)
\]
是同构。

(2) 若 $f:G^m\to S$ 连续，$S$ 离散，则对每一点 $(\sigma_1,\ldots,\sigma_m)$，纤维
\[
 f^{-1}(f(\sigma_1,\ldots,\sigma_m))
\]
是开集。profinite 群的开邻域含有形如
\[
 H\sigma_1\times\cdots\times H\sigma_m
\]
的基本邻域，其中 $H$ 是某个开正规子群。故 $f$ 在这些集合上取常值，即局部正则。反过来，若局部正则，则每一点有这样的开邻域落在一个纤维中，所以连续。

(3) 连续 $2$-上链可用齐性写法表示为局部正则映射
\[
 f:G^3\to F^{\operatorname{sep}\times}
\]
满足等变条件和上闭链条件。题中给出的等式
\[
 f(\rho,\sigma,\tau)f(\nu,\rho,\tau)
 =
 f(\nu,\sigma,\tau)f(\nu,\rho,\sigma)
\]
就是齐性 $2$-cocycle 条件的乘法写法。若 $z:G^2\to F^{\operatorname{sep}\times}$ 连续，则
\[
 \partial z(\rho,\sigma,\tau)=z(\rho,\sigma)z(\sigma,\tau)z(\rho,\tau)^{-1}
\]
是平凡因子集。边界的集合 $B^2$ 是 $Z^2$ 的子群，因为 $\partial$ 是群同态且 $\partial^2=1$。于是连续上同调为
\[
 H^2(G,F^{\operatorname{sep}\times})=Z^2/B^2.
\]
\end{proof}

\begin{proof}[补充]
为了把第 (3) 问讲得更完整，可以再补一句：因子集对应的是扩张在乘法上的扭曲，
而边界 $B^2$ 对应重选截面后得到的等价扭曲。也就是说，$H^2$ 不是把所有扭曲都记住，
而是把本质不同的扭曲分出来。
\end{proof}

\begin{exercise}
证明中心单代数的基本性质，并构造 Brauer 群 $\operatorname{Br}(k)$。
\end{exercise}

\begin{proof}[解答]
本题可解，核心是 Artin--Wedderburn 定理和 Skolem--Noether 定理。这里真正的意思是：
中心单代数的“本质部分”只有一个除代数，矩阵块只是重复；而相似类把这种重复淡化掉，
留下可做群运算的对象。

(1) 若 $A$ 是有限维中心单 $k$-代数，Artin--Wedderburn 定理给出
\[
 A\simeq M_n(D)
\]
其中 $D$ 是有限维除 $k$-代数。$D$ 由 $A$ 唯一决定到同构，因为 $D$ 是简单左 $A$-模的自同态除环的反环。可分闭包 $k^{\operatorname{sep}}$ 上没有非平凡有限维中心除代数，故
\[
 A\otimes_k k^{\operatorname{sep}}\simeq M_N(k^{\operatorname{sep}}).
\]

(2) Skolem--Noether 定理说明中心单代数的每个 $k$-自同构都是内自同构，即存在 $a\in A^\times$ 使
\[
 \varphi(x)=a^{-1}xa.
\]
证明思路是把 $A$ 看作简单 $A$-模，比较原作用和经 $\varphi$ 扭曲的作用；两个简单表示同构，给出共轭元。

(3) 若 $A,A'$ 中心单，则 $A\otimes_kA'$ 的中心为 $k$，双边理想由单性推出只能为 $0$ 或全体，所以仍中心单。若
\[
 A\simeq M_n(D),\quad A'\simeq M_{n'}(D'),\quad
 D\otimes_kD'\simeq M_m(D''),
\]
则
\[
 A\otimes_kA'\simeq M_{nn'm}(D'').
\]
定义 $A\sim A'$ 当且仅当其除代数部分同构。这一关系自反、对称、传递；每个等价类含唯一除代数代表。

(4) 反代数 $A^\circ$ 与 $A$ 作为向量空间相同，乘法反向。映射
\[
 A\otimes_kA^\circ\to \operatorname{End}_k(A),\qquad
 a\otimes b\mapsto(x\mapsto axb)
\]
是代数同态。两边维数都为 $(\dim_kA)^2$，且 $A$ 中心单保证核为 $0$，故同构。因此
\[
 A\otimes_kA^\circ\simeq M_{\dim_kA}(k),
\]
所以 $cl(A^\circ)$ 是 $cl(A)$ 的逆。

(5) 令 $\operatorname{Br}(k)$ 为中心单代数的相似类。乘法定义为
\[
 cl(A)cl(A')=cl(A\otimes_kA').
\]
单位元是 $cl(k)$，逆元为 $cl(A^\circ)$。结合律来自张量积结合律。因此 $\operatorname{Br}(k)$ 是群。
\end{proof}

\begin{exercise}
设 $A$ 是 $F$ 上中心单代数，$\dim_FA=n^2$，$E/F$ 为域扩张，$\varphi:A\to M_n(E)$ 为 $F$-线性映射。证明扩张
\[
 \varphi_E:A\otimes_FE\to M_n(E)
\]
是同构，当且仅当 $\varphi$ 是保单位的 $F$-代数同态。
\end{exercise}

\begin{proof}[解答]
若 $\varphi_E$ 是代数同构，则限制到 $A$ 后满足
\[
 \varphi(1_A)=1_n,\qquad \varphi(ab)=\varphi(a)\varphi(b).
\]

反过来，若 $\varphi$ 是保单位代数同态，则 $\varphi_E$ 是 $E$-代数同态。由于 $A$ 中心单，$A\otimes_FE$ 是中心单 $E$-代数或至少半单中心单的基变换；其维数为 $n^2$，与 $M_n(E)$ 相同。保单位同态的核是双边理想，不能是全体，因此为 $0$。维数相同，所以单射即满射，故为同构。
\end{proof}

\begin{exercise}
证明中心单代数给出同构
\[
 \operatorname{Br}(F)\simeq H^2(G,F^{\operatorname{sep}\times}),
\qquad
 G=\operatorname{Gal}(F^{\operatorname{sep}}/F).
\]
\end{exercise}

\begin{proof}[解答]
这是 Brauer 群与二阶 Galois 上同调的标准构造。给定中心单 $F$-代数 $A$，选有限 Galois 分裂域 $E/F$ 和同构
\[
 \varphi:A\otimes_FE\simeq M_n(E).
\]
对 $\rho,\sigma\in G$，两个嵌入 $\varphi^\rho,\varphi^\sigma$ 都把 $A$ 表示为矩阵代数中的中心单子代数。由 Skolem--Noether 定理，存在
\[
 Y(\rho,\sigma)\in \operatorname{GL}_n(F^{\operatorname{sep}})
\]
使
\[
 Y(\rho,\sigma)\varphi^\sigma=\varphi^\rho Y(\rho,\sigma).
\]
局部常值性来自 $Y$ 在某个有限分裂域上定义。结合三重比较可得
\[
 f_{A,\varphi}(\rho,\sigma,\tau)
 =
 Y(\rho,\sigma)Y(\sigma,\tau)Y(\rho,\tau)^{-1}
 \in F^{\operatorname{sep}\times},
\]
且它满足 $2$-因子集条件。

更换 $Y$ 或扩大分裂域只会把 $f$ 乘以一个边界，因此得到良定的类
\[
 [f_{A,\varphi}]\in H^2(G,F^{\operatorname{sep}\times}).
\]
张量积代数对应因子集相乘，反代数对应取逆，所以这是群同态。

满射可由给定因子集构造交叉积代数得到；单射来自因子集为边界时，交叉积裂为矩阵代数。
于是得到同构。换句话说，Brauer 类就是“扭曲矩阵代数”的扭曲方式，而这个扭曲方式
正好由 $H^2$ 记录。
\end{proof}

\begin{exercise}
给定 $\theta\in F^\times$ 和特征标 $\chi:G\to\C^\times$，取局部常值函数 $\Phi:G\to\R$ 使
\[
 \chi(\sigma)=e^{2\pi i\Phi(\sigma)}.
\]
定义相应因子集并证明得到类 $\{\chi,\theta\}$；证明 $\chi\mapsto\{\chi,\theta\}$ 是群同态，且 $\theta\mapsto\{\chi,\theta\}$ 的核是范数群。
\end{exercise}

\begin{proof}[解答]
这一题是第 5 章把字符和 Brauer 类联系起来的关键模型。它说明：一个一维特征标和一个
标量 $\theta$ 就足以生成一个二阶类。

设
\[
 f(\rho,\sigma,\tau)
 =
 \theta^{\Phi(\sigma\rho^{-1})+\Phi(\tau\sigma^{-1})-\Phi(\tau\rho^{-1})}.
\]
因为 $\chi$ 是同态，指数中的加法误差总为整数；所以 $f$ 的取值落在由 $\theta$ 生成的离散子群中，并且局部常值。把四个点 $\nu,\rho,\sigma,\tau$ 代入 $2$-cocycle 条件，指数项逐项抵消，得到因子集。

若换 $\Phi'$，令
\[
 z(\rho,\sigma)=
 \theta^{(\Phi'-\Phi)(\sigma\rho^{-1})},
\]
直接代入边界公式可得
\[
 f'f^{-1}=\partial z.
\]
故上同调类与 $\Phi$ 的选择无关，记为 $\{\chi,\theta\}$。

若 $\chi_1,\chi_2$ 对应 $\Phi_1,\Phi_2$，则 $\chi_1\chi_2$ 可取提升 $\Phi_1+\Phi_2$，相应因子集是两个因子集的乘积，所以
\[
 \{\chi_1\chi_2,\theta\}=\{\chi_1,\theta\}+\{\chi_2,\theta\}.
\]

设 $\chi$ 对应有限循环扩张 $E/F$。类 $\{\chi,\theta\}$ 正是循环代数的 Brauer 类。由 Hilbert 90 和循环代数分裂判别，循环代数分裂当且仅当
\[
 \theta\in N_{E/F}E^\times.
\]
因此 $\theta\mapsto\{\chi,\theta\}$ 的核是 $N_{E/F}E^\times$。
\end{proof}

\begin{exercise}
设 $E/F$ 是 $n$ 次循环扩张，$\alpha$ 生成 $\operatorname{Gal}(E/F)$，特征标 $\chi(\alpha)=e^{2\pi i/n}$，取 $\theta\in F^\times$。构造循环代数
\[
 [E/F;\chi,\theta]
\]
并证明其 Brauer 类对应 $\{\chi,\theta\}$。
\end{exercise}

\begin{proof}[解答]
这题只是把上一题的抽象因子集换成最经典的循环代数模型。一个基元 $u_1$ 负责把
$E$ 上的 Galois 作用编码进代数乘法，关系 $u_1^n=\theta$ 则把二阶扭曲固定下来。

令
\[
 X=\bigoplus_{i=0}^{n-1}Eu_i
\]
为左 $E$-向量空间。定义乘法规则：
\[
 u_0=1,\qquad u_1^n=\theta u_0,\qquad u_i=u_1^i,
\]
\[
 u_1\xi=\alpha(\xi)u_1\qquad(\xi\in E).
\]
这些规则唯一决定双 $F$-线性乘法。任意元素写为 $\sum_i\xi_iu_i$，乘法由分配律和上面关系化为标准形。结合律只需检查生成元 $\xi\in E$ 与 $u_1$：
\[
 u_1(\xi\eta)=\alpha(\xi\eta)u_1
 =\alpha(\xi)\alpha(\eta)u_1
 =(u_1\xi)\eta,
\]
以及 $u_1^n=\theta$ 与 $F$ 中元素可交换。故 $X$ 成为 $F$-代数，称为循环代数。

在 $E$ 上基变换后，循环代数裂为矩阵代数；$u_1$ 的作用给出伴随矩阵，Galois 元素作用时产生的因子正是第 7 题构造的 $\theta$-因子集。把该分裂同构代入 Brauer 群到 $H^2(G,F^{\operatorname{sep}\times})$ 的构造，得到的上同调类就是
\[
 \{\chi,\theta\}.
\]
因此
\[
 cl([E/F;\chi,\theta])\longmapsto \{\chi,\theta\}.
\]
\end{proof}
 \chapter{局部域的上同调群}

这一章把第五章的上同调工具真正用于局部域。前面已经学过局部数域、无分歧扩张、形式群和 Lubin--Tate 除点；这里要把这些对象放进类成原则中，得到局部互反律。对初学者来说，最容易迷路的地方有两个。第一，$H^2$ 不是凭空出现的抽象群，它通过赋值
\[
0\longrightarrow U_E\longrightarrow E^\times\stackrel{v_E}{\longrightarrow}\Z\longrightarrow0
\]
把乘法群的单位部分和整数赋值联系起来。第二，互反映射不是一个临时定义的同态，而是 Tate 定理把基本类 $u_{E|F}\in H^2(\Gal(E/F),E^\times)$ 与低维 Tate 上同调配对后产生的同构。

\section{无分歧扩张}

设 $F$ 为特征 $0$ 的局部域，$\mathcal O_F$ 为整数环，$\mathfrak p_F$ 为极大理想，$\kappa_F$ 为有限剩余域。无分歧扩张 $E/F$ 的特点是分歧指数 $e(E/F)=1$，所以素元不变，全部扩张度来自剩余域扩张：$[E:F]=[\kappa_E:\kappa_F]$。这使得无分歧扩张成为局部类域论中最容易计算的一部分。

若 $E/F$ 是有限无分歧 Galois 扩张，剩余域扩张 $\kappa_E/\kappa_F$ 也是有限 Galois 扩张，并由 Frobenius 生成。记
\[
\varphi_{E/F}\in\Gal(E/F)
\]
为提升剩余域 Frobenius $x\mapsto x^{q_F}$ 的自同构，也就是说
\[
\overline{\varphi_{E/F}(a)}=\bar a^{q_F}\qquad (a\in\mathcal O_E).
\]
这里 $q_F=|\kappa_F|$。这个定义把 Galois 群中的元素和剩余域上的显式幂映射连接起来。

\begin{proposition}
设 $F\subset K\subset E$，且 $E/F$、$K/F$ 都是有限无分歧扩张。则 Frobenius 满足：
\[
\varphi_{K/F}=\varphi_{E/F}|_K=\varphi_{E/K}\Gal(E/K)\in\Gal(K/F),
\]
并且
\[
\varphi_{E/K}=\varphi_{E/F}^{[K:F]}.
\]
\end{proposition}

\begin{proof}
先看限制。对任意 $a\in\mathcal O_K$，把它也看成 $\mathcal O_E$ 中元素，则
\[
\overline{\varphi_{E/F}(a)}=\bar a^{q_F}.
\]
右边属于 $\kappa_K$，所以 $\varphi_{E/F}(a)$ 的剩余类在 $\kappa_K$ 中。由于 $K/F$ 是无分歧扩张，$\mathcal O_K$ 是 $\mathcal O_F$ 的非分歧整闭扩张；剩余域自同构决定唯一的无分歧提升，因此 $\varphi_{E/F}$ 把 $K$ 保持到自身，并且限制后的自同构就是 $\varphi_{K/F}$。

把 $\varphi_{E/F}$ 在商群 $\Gal(E/F)/\Gal(E/K)\simeq\Gal(K/F)$ 中取像，得到的正是它在 $K$ 上的限制，所以第一式的商群表述也成立。

再看第二式。设 $q_K=|\kappa_K|$。因为 $K/F$ 无分歧，$q_K=q_F^{[K:F]}$。对任意 $a\in\mathcal O_E$，
\[
\overline{\varphi_{E/F}^{[K:F]}(a)}=\bar a^{q_F^{[K:F]}}=\bar a^{q_K}.
\]
这正是 $E/K$ 的 Frobenius 在剩余域上的作用。无分歧提升唯一，所以 $\varphi_{E/K}=\varphi_{E/F}^{[K:F]}$。
\end{proof}

\begin{proposition}
设 $F$ 为特征 $0$ 的局部域，$E/F$ 为有限无分歧扩张，$G=\Gal(E/F)$。则对所有整数 $q$，
\[
\widehat H^q(G,U_E)=1.
\]
特别地，
\[
U_F=N_{E/F}U_E.
\]
\end{proposition}

\begin{proof}
单位群有主单位过滤
\[
U_E=U_E^{(0)}\supset U_E^{(1)}\supset U_E^{(2)}\supset\cdots,\qquad U_E^{(n)}=1+\mathfrak p_E^n.
\]
无分歧意味着 $\mathfrak p_E=\pi\mathcal O_E$ 可取同一个素元 $\pi\in F$，并且
\[
U_E^{(n)}/U_E^{(n+1)}\simeq \kappa_E
\]
作为加法群和 $G$-模，$G$ 的作用就是剩余域扩张的 Galois 作用。

先处理有限层。有限域扩张 $\kappa_E/\kappa_F$ 的加法群上同调为零：这是有限域加法群的正常基定理的结果，因为 $\kappa_E$ 作为 $\kappa_F[G]$-模是自由秩 $1$，从而是诱导模；第五章已经说明诱导模 Tate 上同调全为零。因此每个商 $U_E^{(n)}/U_E^{(n+1)}$ 的 Tate 上同调都为零。由短正合列
\[
1\longrightarrow U_E^{(n+1)}\longrightarrow U_E^{(n)}
\longrightarrow \kappa_E\longrightarrow0
\]
得到
\[
\widehat H^q(G,U_E^{(n+1)})\simeq \widehat H^q(G,U_E^{(n)})
\]
对所有 $q$ 成立。

再把过滤推进到足够深处。令 $m=|G|$。当 $n$ 充分大时，$p$ 进对数或二项式展开给出同构
\[
U_E^{(n)}\longrightarrow U_E^{(n+v_E(m))},\qquad x\longmapsto x^m.
\]
于是 $\widehat H^q(G,U_E^{(n)})$ 上的乘 $m$ 映射是同构。另一方面，有限群 $G$ 的 Tate 上同调总被 $|G|=m$ 杀掉。一个群若同时被 $m$ 杀掉而乘 $m$ 又为同构，只能为零。由前面的同构链回推，$\widehat H^q(G,U_E)=0$。

取 $q=0$。按 Tate 零维上同调定义，
\[
\widehat H^0(G,U_E)=U_E^G/N_{E/F}U_E=U_F/N_{E/F}U_E.
\]
该群为零，故 $U_F=N_{E/F}U_E$。
\end{proof}

设 $F_0$ 为特征 $0$ 的局部域，$F_0^{\mathrm{nr}}$ 为其最大无分歧扩张。它是所有有限无分歧扩张的并，且
\[
\Gal(F_0^{\mathrm{nr}}/F_0)\simeq \widehat{\Z},
\]
拓扑生成元由 Frobenius 给出。下面要说明无分歧塔已经满足类成原则的形式要求。

\begin{theorem}
设 $F_0$ 为特征 $0$ 的局部域，$F_0^{\mathrm{nr}}$ 为最大无分歧扩张。对任意有限无分歧扩张 $E/F$，其中 $F_0\subset F\subset E\subset F_0^{\mathrm{nr}}$，有：
\[
H^1(\Gal(E/F),E^\times)=1,
\]
并且存在同构
\[
\operatorname{inv}_{E|F}:H^2(\Gal(E/F),E^\times)\longrightarrow \frac1{[E:F]}\Z/\Z.
\]
这些同构关于中间域的膨胀、限制和扩张次数相容。
\end{theorem}

\begin{proof}
记 $G=\Gal(E/F)$。Hilbert 90 给出 $H^1(G,E^\times)=1$。

为了定义不变量映射，使用赋值短正合列
\[
1\longrightarrow U_E\longrightarrow E^\times\stackrel{v_E}{\longrightarrow}\Z\longrightarrow0.
\]
由上一命题，$U_E$ 的 Tate 上同调为零，所以
\[
H^2(G,E^\times)\simeq H^2(G,\Z).
\]
再用短正合列
\[
0\longrightarrow \Z\longrightarrow\Q\longrightarrow\Q/\Z\longrightarrow0.
\]
因为有限群 $G$ 在唯一可除群 $\Q$ 上的高阶上同调为零，连接同态给出
\[
H^1(G,\Q/\Z)\simeq H^2(G,\Z).
\]
于是一个 $2$ 阶上同调类先经赋值变成 $H^2(G,\Z)$ 中的类，再经连接同态逆映到一个特征标 $\chi:G\to\Q/\Z$。定义
\[
\operatorname{inv}_{E|F}(c)=\chi(\varphi_{E/F}).
\]
若 $[E:F]=n$，则 $G$ 是由 $\varphi_{E/F}$ 生成的 $n$ 阶循环群，因此 $\chi(\varphi_{E/F})$ 恰落在 $\frac1n\Z/\Z$ 中，而且每个这样的值都由某个特征标得到，所以这是同构。

相容性来自三个事实：赋值映射与域嵌入相容；连接同态对上同调函子自然；Frobenius 在塔中的行为由上一命题给出。若 $F\subset K\subset E$ 且 $K/F$ 无分歧，膨胀不改变在 $\varphi$ 上的取值。若 $E/K$ 无分歧，则限制时 Frobenius 由 $\varphi_{E/K}=\varphi_{E/F}^{[K:F]}$ 给出，所以不变量乘以 $[K:F]$。
\end{proof}

\begin{proposition}
设 $F_0\subset F\subset F_0^{\mathrm{nr}}$ 为有限无分歧扩张。则
\[
H^2(\Gal(F_0^{\mathrm{nr}}/F),F_0^{\mathrm{nr}\times})\simeq \Q/\Z.
\]
\end{proposition}

\begin{proof}
连续上同调按有限子扩张取直接极限：
\[
H^2(\Gal(F_0^{\mathrm{nr}}/F),F_0^{\mathrm{nr}\times})
\simeq \varinjlim_E H^2(\Gal(E/F),E^\times),
\]
其中 $E/F$ 遍历有限无分歧扩张。对每个次数 $n=[E:F]$，上一定理 给出
\[
H^2(\Gal(E/F),E^\times)\simeq \frac1n\Z/\Z.
\]
有限无分歧扩张对每个 $n\ge1$ 唯一存在，且当 $E\subset E'$ 时不变量映射与包含
\[
\frac1n\Z/\Z\subset \frac1{n'}\Z/\Z
\]
相容。因此直接极限为
\[
\bigcup_{n\ge1}\frac1n\Z/\Z=\Q/\Z.
\]
\end{proof}

\section{局部互反律}

本节从无分歧扩张进入任意有限 Galois 扩张。核心困难是：任意扩张不再由剩余域控制，单位群上同调也不再消失。但是局部域的 Brauer 群有一个不变量映射
\[
\operatorname{inv}_F:\Br(F)\longrightarrow \Q/\Z,
\]
它把所有有限扩张的 $H^2$ 统一起来。局部互反律就是把这个 $H^2$ 中的基本类转换成乘法群商和 Abel Galois 群的同构。

\subsection{\texorpdfstring{$H^2$}{H2}}

\begin{theorem}
设 $E/F$ 为有限正规局部域扩张，$E'/F$ 为无分歧扩张，且 $[E':F]=[E:F]$。则通过局部 Brauer 群中的自然嵌入，可以把
\[
H^2(\Gal(E/F),E^\times)
\]
与
\[
H^2(\Gal(E'/F),E'^\times)
\]
看作 $\Br(F)$ 中同一个阶数受 $[E:F]$ 控制的子群。
\end{theorem}

\begin{proof}
每个有限 Galois 扩张 $L/F$ 都给出相对 Brauer 群
\[
\Br(L/F)=\ker\bigl(\Br(F)\to\Br(L)\bigr)
\simeq H^2(\Gal(L/F),L^\times).
\]
这个同构来自交叉积代数：一个 $2$-余循环记录乘法关系，构造出的中心单代数在 $L$ 上分裂；反过来，任何在 $L$ 上分裂的中心单代数都由这样的因子集给出。

局部域的关键事实是：$\Br(F)$ 的元素可以用无分歧扩张计算。具体地，若
\[
c\in H^2(\Gal(E/F),E^\times),
\]
把它看作 $\Br(F)$ 中元素。取包含 $E$ 与 $E'$ 的有限正规扩张 $L/F$。限制到 $L$ 后它分裂，因而 $c$ 的阶整除 $[E:F]$。无分歧扩张 $E'/F$ 的相对 Brauer 群由上一节的不变量映射识别为
\[
\frac1{[E:F]}\Z/\Z.
\]
因此 $\Br(F)$ 中所有阶数整除 $[E:F]$ 的元素都已经由 $E'/F$ 的无分歧相对 Brauer 群表示。

换句话说，任意正规扩张的 $H^2$ 类，在局部 Brauer 群中可以用同次数无分歧扩张的 $H^2$ 类来度量。这个结论正是后面定义一般不变量映射所需要的桥梁：困难扩张的 $H^2$ 先进入 $\Br(F)$，再用无分歧模型读取它的值。
\end{proof}

\begin{definition}
设 $E/F$ 为有限正规局部域扩张。取无分歧扩张 $E'/F$，使 $[E':F]=[E:F]$。对
\[
c\in H^2(\Gal(E/F),E^\times)
\]
先把 $c$ 看作 $\Br(F)$ 中元素，再用上一定理 把它与 $E'/F$ 的相对 Brauer 群相比较，定义
\[
\operatorname{inv}_{E|F}(c)=\operatorname{inv}_{E'|F}(c)\in \frac1{[E:F]}\Z/\Z.
\]
\end{definition}

这个定义的意思是：一般扩张的 $H^2$ 不直接由 Frobenius 计算，但它在 Brauer 群中有一个无分歧替身；不变量映射就是在无分歧替身上读数。

\begin{proposition}
若 $F$ 是 $p$ 进局部域，则
\[
\Br(F)\simeq \Q/\Z.
\]
\end{proposition}

\begin{proof}
Brauer 群是相对 Brauer 群沿所有有限 Galois 扩张的直接极限：
\[
\Br(F)=\varinjlim_E H^2(\Gal(E/F),E^\times).
\]
上一定理 说明每个有限层中的类都可由同次数无分歧层读取不变量。有限无分歧层给出的子群为 $\frac1n\Z/\Z$，其中 $n$ 是扩张次数。局部域对每个 $n$ 有唯一的 $n$ 次无分歧扩张，所以所有这些子群的并为
\[
\bigcup_{n\ge1}\frac1n\Z/\Z=\Q/\Z.
\]
这给出同构 $\operatorname{inv}_F:\Br(F)\to\Q/\Z$。它的单射性来自相对 Brauer 群的无分歧代表：若不变量为 $0$，则在某个无分歧层中对应的 $H^2$ 类为零，因而原 Brauer 类也为零。满射性由有限无分歧扩张的基本类给出。
\end{proof}

\subsection{类成原则}

\begin{theorem}
设 $F_0$ 为特征 $0$ 的局部域，$\bar F_0$ 为代数闭包。则
\[
(\Gal(\bar F_0/F_0),\bar F_0^\times)
\]
满足类成原则：对任意有限正规扩张 $E/F$，其中 $F_0\subset F\subset E\subset\bar F_0$，
\[
H^1(\Gal(E/F),E^\times)=1,
\]
并且存在与塔相容的不变量同构
\[
\operatorname{inv}_{E|F}:H^2(\Gal(E/F),E^\times)\longrightarrow \frac1{[E:F]}\Z/\Z.
\]
\end{theorem}

\begin{proof}
第一条是 Hilbert 90。第二条由上一小节的 Brauer 群不变量构造。需要检查的相容性有两种。

若 $F\subset K\subset E$ 且 $K/F$ 正规，则膨胀映射对应于把在 $K$ 上分裂的 Brauer 类看作在 $E$ 上也分裂的 Brauer 类；不变量不变。若 $E/K$ 正规，则限制映射对应于底域从 $F$ 改为 $K$。局部不变量在有限扩张下满足
\[
\operatorname{inv}_K(\operatorname{Res}_{K/F}c)=[K:F]\operatorname{inv}_F(c).
\]
这在无分歧情形已经由 Frobenius 计算验证；一般情形由上一小节把类比较到同次数无分歧扩张后得到。于是类成原则的两条相容性成立。
\end{proof}

对有限正规扩张 $E/F$，取基本类
\[
u_{E|F}\in H^2(\Gal(E/F),E^\times)
\]
为满足
\[
\operatorname{inv}_{E|F}(u_{E|F})=\frac1{[E:F]}+\Z
\]
的元素。

\begin{proposition}
设 $E/F$ 为特征 $0$ 的局部域有限正规扩张，$G=\Gal(E/F)$。

\begin{enumerate}
\item 杯积基本类给出 Tate 同构
\[
u_{E|F}\cup -:\widehat H^q(G,\Z)\longrightarrow \widehat H^{q+2}(G,E^\times)
\]
对所有整数 $q$ 成立。
\item 取 $q=-2$，得到同构
\[
r_{E|F}:G^{\operatorname{ab}}\longrightarrow F^\times/N_{E/F}E^\times.
\]
它称为局部互反映射。其逆映射与自然投影复合，给出范剩符号
\[
(\ ,E|F):F^\times\longrightarrow G^{\operatorname{ab}}.
\]
\end{enumerate}
\end{proposition}

\begin{proof}
类成原则正是 Tate 定理的输入。对每个子群 $H\subset G$，有限扩张 $E/E^H$ 也满足
\[
H^1(H,E^\times)=1,\qquad H^2(H,E^\times)\simeq \frac1{|H|}\Z/\Z,
\]
并且基本类限制到 $H$ 后仍是 $E/E^H$ 的基本类。因此第五章的 Tate 定理适用，杯积给出所有次数的同构。

当 $q=-2$ 时，
\[
\widehat H^{-2}(G,\Z)\simeq G^{\operatorname{ab}},
\]
而
\[
\widehat H^0(G,E^\times)=F^\times/N_{E/F}E^\times.
\]
杯积同构于是给出
\[
G^{\operatorname{ab}}\simeq F^\times/N_{E/F}E^\times.
\]
本讲义把这个方向记为 $r_{E|F}$；范剩符号则取它的逆，再从 $F^\times$ 投到范数商。
\end{proof}

\begin{proposition}
设 $K\supset E\supset F$ 为局部域有限扩张，$K/F$ 正规。范剩符号满足以下相容性。
\begin{enumerate}
\item 若 $E/F$ 正规，则从 $\Gal(K/F)^{\operatorname{ab}}$ 到 $\Gal(E/F)^{\operatorname{ab}}$ 的投影与 $F^\times$ 上同一个元素的范剩符号相容。
\item 包含 $F^\times\hookrightarrow E^\times$ 对应于 Galois 群上的迁移映射
\[
\operatorname{Ver}:\Gal(K/F)^{\operatorname{ab}}\to \Gal(K/E)^{\operatorname{ab}}.
\]
\item 范数 $N_{E/F}:E^\times\to F^\times$ 对应于包含
\[
\Gal(K/E)\hookrightarrow \Gal(K/F)
\]
诱导的 Abel 化映射。
\item 对任意 $F$-嵌入的共轭作用，范剩符号与共轭后的域扩张相容。
\end{enumerate}
\end{proposition}

\begin{proof}
这些相容性都来自同一个机制：基本类在限制、膨胀和芯限制下相容，杯积又与这些上同调操作相容。

第一条对应于膨胀。若 $E/F$ 正规，则 $G_{K/F}\to G_{E/F}$ 的投影在 $H^{-2}(-,\Z)$ 上就是 Abel 化投影；在 $H^0(-,K^\times)$ 上就是把 $F^\times/N_{K/F}K^\times$ 映到 $F^\times/N_{E/F}E^\times$。杯积方块交换，所以范剩符号投影相容。

第二条对应于限制。负二阶 Tate 上同调中的限制正是群论中的迁移映射；零阶上同调中的限制是把 $F^\times$ 看成 $E^\times$ 的元素。杯积与限制相容，因此包含映射和迁移相配。

第三条对应于芯限制。零阶上同调中的芯限制由范数给出；负二阶中的芯限制由子群包含诱导。杯积与芯限制相容，于是范数和包含相配。

第四条是函子性。共轭同构把扩张 $K/F$ 变为 $\sigma K/\sigma F$，也把基本类带到对应基本类；杯积构造不依赖坐标选择，所以范剩符号在共轭下保持相同形式。
\end{proof}

\begin{proposition}
若 $E/F$ 是有限无分歧扩张，则对任意 $a\in F^\times$，
\[
(a,E|F)=\varphi_{E/F}^{v_F(a)}.
\]
\end{proposition}

\begin{proof}
设 $G=\Gal(E/F)$，取特征标 $\chi\in H^1(G,\Q/\Z)=\operatorname{Hom}(G,\Q/\Z)$。范剩符号由基本类的杯积定义，因此
\[
\chi((a,E|F))=\operatorname{inv}_{E|F}(\bar a\cup\delta\chi),
\]
其中 $\bar a$ 是 $a$ 在 $F^\times/N_{E/F}E^\times$ 中的类，$\delta\chi\in H^2(G,\Z)$ 是由 $0\to\Z\to\Q\to\Q/\Z\to0$ 得到的连接像。

无分歧情形中，赋值把 $\bar a$ 送到 $v_F(a)\in\Z/n\Z$，单位部分因 $U_F=N_{E/F}U_E$ 没有贡献。不变量映射的定义正是把 $\delta\chi$ 在 Frobenius 上读数，所以
\[
\chi((a,E|F))=v_F(a)\chi(\varphi_{E/F})
=\chi(\varphi_{E/F}^{v_F(a)}).
\]
有限 Abel 群中的元素由所有特征标分离，因此
\[
(a,E|F)=\varphi_{E/F}^{v_F(a)}.
\]
\end{proof}

\begin{proposition}
设 $F^{\mathrm{nr}}/F$ 为最大无分歧扩张，$\varphi_F=\varprojlim_E\varphi_{E/F}$。则
\[
(a,F^{\mathrm{nr}}|F)=\varphi_F^{v_F(a)}.
\]
特别地，
\[
(a,F^{\mathrm{nr}}|F)=1\quad\Longleftrightarrow\quad a\in U_F.
\]
\end{proposition}

\begin{proof}
对任意有限无分歧子扩张 $E/F$，把等式投影到 $\Gal(E/F)$ 后，由上一命题得到
\[
\operatorname{pr}_E(a,F^{\mathrm{nr}}|F)=(a,E|F)=\varphi_{E/F}^{v_F(a)}
=\operatorname{pr}_E(\varphi_F^{v_F(a)}).
\]
所有有限投影相同，所以射影极限中的元素相同。又 $\varphi_F$ 是 $\widehat{\Z}$ 的拓扑生成元，$\varphi_F^m=1$ 当且仅当 $m=0$。因此核正是 $v_F(a)=0$ 的元素，也就是单位群 $U_F$。
\end{proof}

\begin{proposition}
全局定义的局部范剩符号
\[
(\ ,F):F^\times\longrightarrow G_F^{\operatorname{ab}}
\]
是单射。
\end{proposition}

\begin{proof}
其核为所有有限 Abel 扩张范数群的交：
\[
\ker(\ ,F)=\bigcap_{E/F\ \mathrm{finite\ abelian}}N_{E/F}E^\times.
\]
我们证明这个交只有 $1$。

写
\[
F^\times=\langle\pi\rangle\times U_F.
\]
无分歧扩张已经检测赋值部分。若 $a=\pi^m u$ 在核中，那么投到最大无分歧扩张得到
\[
\varphi_F^{m}=1,
\]
所以 $m=0$，即 $a\in U_F$。

对单位部分，局部类域论的存在定理说明每个开有限指数子群都是某个有限 Abel 扩张的范数群。单位群有过滤
\[
U_F\supset U_F^{(1)}\supset U_F^{(2)}\supset\cdots,\qquad \bigcap_n U_F^{(n)}=\{1\}.
\]
对每个 $n$，存在有限 Abel 扩张 $E_n/F$ 使 $N_{E_n/F}E_n^\times$ 包含在 $U_F^{(n)}$ 与适当赋值条件的乘积中。若 $u$ 属于所有范数群，则 $u\in U_F^{(n)}$ 对所有 $n$ 成立，故 $u=1$。因此核为 $\{1\}$。
\end{proof}

\subsection{除点}

这一小节把局部互反律与第三章的 Lubin--Tate 除点联系起来。设 $K$ 为局部域，$\pi$ 为素元，$\mathcal O_K$ 为整数环，剩余域大小为 $q$。取 Lubin--Tate 幂级数 $f\in\mathcal F_\pi$，即
\[
f(X)\equiv \pi X\pmod{X^2},\qquad f(X)\equiv X^q\pmod{\pi}.
\]
对应的一维形式群记为 $F_f$。令 $F_f[\pi^n]$ 为 $\pi^n$-除点，设
\[
E_{\pi,n}=K(F_f[\pi^n]).
\]
第三章已经构造了同态
\[
\omega_\pi:K^\times\longrightarrow \Gal(K^{\mathrm{nr}}E_{\pi,n}/K),
\]
它在无分歧部分把 $\pi$ 送到 Frobenius，在全分歧 Lubin--Tate 部分由单位作用给出。

\begin{proposition}
对任意 $a\in K^\times$，
\[
\omega_\pi(a)=(a,K)|_{K^{\mathrm{nr}}E_{\pi,n}}.
\]
\end{proposition}

\begin{proof}
先验证生成元。乘法群由素元和单位生成。

在 $K^{\mathrm{nr}}$ 上，$\omega_\pi(\pi)$ 按定义是 Frobenius；上一命题说明
\[
(\pi,K^{\mathrm{nr}}|K)=\varphi_K,
\]
所以两者在无分歧部分一致。在 $E_{\pi,n}$ 上，Lubin--Tate 理论给出 $\pi$ 是 $E_{\pi,n}/K$ 的范数，因此
\[
(\pi,E_{\pi,n}|K)=1,
\]
而 $\omega_\pi(\pi)$ 在除点扩张上也按构造平凡。

再看单位 $u\in U_K$。在无分歧部分 $v_K(u)=0$，所以范剩符号平凡；$\omega_\pi(u)$ 的无分歧限制也平凡。在 $E_{\pi,n}$ 上，Lubin--Tate 的单位作用把 $u$ 对应到形式群自同态 $[u]_f$。局部互反律的相容性说明范剩符号在全分歧 Lubin--Tate 扩张上的作用正是这一单位作用。于是两者在 $K^{\mathrm{nr}}$ 和 $E_{\pi,n}$ 上都一致，故在合成域上一致。
\end{proof}

\begin{proposition}
设 $a=u\pi^m$，其中 $u\in U_K$，$m\in\Z$。对 $x\in F_f[\pi^n]$，
\[
(a,E_{\pi,n}|K)(x)=[u^{-1}]_f(x).
\]
并且
\[
N_{E_{\pi,n}/K}E_{\pi,n}^\times=U_K^{(n)}\times\langle\pi\rangle.
\]
\end{proposition}

\begin{proof}
上一命题把范剩符号与 $\omega_\pi$ 识别。Lubin--Tate 约定中，互反映射方向与单位作用互为逆，所以 $u$ 在 Galois 群上作用为 $[u^{-1}]_f$。素元 $\pi$ 在 $E_{\pi,n}$ 上作用平凡，因为 $\pi$ 是该全分歧 Lubin--Tate 扩张的范数。因此
\[
(u\pi^m,E_{\pi,n}|K)(x)=[u^{-1}]_f(x).
\]

范数群等于范剩符号的核。上式说明 $u\pi^m$ 在 $E_{\pi,n}$ 上作用平凡，当且仅当 $[u^{-1}]_f$ 在 $F_f[\pi^n]$ 上平凡。Lubin--Tate 理论中
\[
[u]_f\equiv X\pmod{\pi^n}
\]
当且仅当 $u\in U_K^{(n)}$。因此核正是
\[
U_K^{(n)}\times\langle\pi\rangle.
\]
\end{proof}

\begin{proposition}
设 $K/F$ 为有限无分歧扩张，$q=|\kappa_F|$，$\varphi_{K/F}$ 为 Frobenius。则局部互反映射在无分歧部分满足
\[
(\pi,K|F)=\varphi_{K/F},\qquad (u,K|F)=1\ (u\in U_F),
\]
并且
\[
(a,K|F)=\varphi_{K/F}^{v_F(a)}\qquad(a\in F^\times).
\]
\end{proposition}

\begin{proof}
前两式分别是素元和单位的情形。素元的像由无分歧 Frobenius 规定，单位的像为零，因为无分歧扩张的单位范数群就是全体单位群。乘法群分解
\[
F^\times=\langle\pi\rangle\times U_F
\]
后，任意 $a\in F^\times$ 都可写成 $a=\pi^{v_F(a)}u$。范剩符号是群同态，因此
\[
(a,K|F)=(\pi,K|F)^{v_F(a)}(u,K|F)=\varphi_{K/F}^{v_F(a)}.
\]
\end{proof}

\begin{proposition}
设 $a=up^m\in\Q_p^\times$，$u\in\Z_p^\times$，$\zeta$ 为 $p^n$ 次本原单位根。取正整数 $r$ 满足
\[
ur\equiv1\pmod{p^n}.
\]
则
\[
(a,\Q_p(\zeta)|\Q_p)(\zeta)=\zeta^r.
\]
\end{proposition}

\begin{proof}
对 $\Q_p$ 取 Lubin--Tate 幂级数
\[
f(X)=(1+X)^p-1.
\]
若 $x=\zeta-1$，则 $x$ 是形式乘法群的 $p^n$-除点。上一命题给出
\[
(a,\Q_p(\zeta)|\Q_p)(x)=[u^{-1}]_f(x).
\]
形式乘法群中 $[b]_f(X)=(1+X)^b-1$。取 $r\equiv u^{-1}\pmod{p^n}$，则在 $p^n$-除点上
\[
[u^{-1}]_f(x)=[r]_f(x)=(1+x)^r-1=\zeta^r-1.
\]
于是
\[
(a,\Q_p(\zeta)|\Q_p)(\zeta)
=1+(a,\Q_p(\zeta)|\Q_p)(x)
=\zeta^r.
\]
\end{proof}

\section{分圆域}

分圆域把局部互反律变成可以计算的公式。对 $\Q$ 来说，单位根扩张给出最经典的 Abel 扩张；对 $\Q_p$ 来说，加入 $\ell^n$ 次单位根时，$\ell\ne p$ 与 $\ell=p$ 的行为完全不同：前者无分歧，后者全分歧。

\begin{proposition}
设 $\zeta$ 为 $m$ 次本原单位根。则：
\begin{enumerate}
\item $\zeta^a$ 是 $m$ 次本原单位根，当且仅当 $(a,m)=1$。
\item $[\Q(\zeta):\Q]=\varphi(m)$。
\item $\Gal(\Q(\zeta)/\Q)\simeq(\Z/m\Z)^\times$，其中 $a$ 对应 $\zeta\mapsto\zeta^a$。
\item 第 $m$ 个分圆多项式为
\[
\Phi_m(X)=\prod_{(a,m)=1}(X-\zeta^a),
\]
并满足
\[
X^m-1=\prod_{d\mid m}\Phi_d(X).
\]
\item 若 $\mu_m=\{\xi_1,\dots,\xi_m\}$，则
\[
\prod_{i\ne j}(\xi_i-\xi_j)=(-1)^{m-1}m^m.
\]
\end{enumerate}
\end{proposition}

\begin{proof}
第一条来自阶的计算：$\zeta^a$ 的阶为 $m/(a,m)$，所以它是本原根当且仅当 $(a,m)=1$。

第二、三条一起证明。每个 $\Q$-嵌入由 $\zeta$ 的像决定，而 $\zeta$ 的像必须仍是 $m$ 次本原单位根；可选像正好有 $\varphi(m)$ 个。映射
\[
(\Z/m\Z)^\times\to\Gal(\Q(\zeta)/\Q),\qquad a\mapsto(\zeta\mapsto\zeta^a)
\]
是群同态，且由上一句知双射。

第四条中，$X^m-1$ 的根正是所有阶整除 $m$ 的单位根。按单位根的精确阶分类，得到
\[
X^m-1=\prod_{d\mid m}\Phi_d(X).
\]

第五条是 $X^m-1$ 的判别式计算。设 $f(X)=X^m-1=\prod_i(X-\xi_i)$。则
\[
f'(\xi_i)=m\xi_i^{m-1}=\prod_{j\ne i}(\xi_i-\xi_j).
\]
把 $i$ 全部相乘：
\[
\prod_i f'(\xi_i)=m^m\prod_i\xi_i^{m-1}.
\]
而 $\prod_i\xi_i=(-1)^{m+1}$，整理符号得到所给公式。
\end{proof}

\begin{proposition}
设 $F$ 为数域。
\begin{enumerate}
\item 对任意正整数 $n$，存在循环分圆扩张 $E/F$，使 $n\mid [E:F]$。
\item 给定有限素位集 $S$ 与正整数 $m$，可取循环分圆扩张 $E/F$，使每个有限素位 $v\in S$ 的局部次数 $[E_w:F_v]$ 被 $m$ 整除；若 $v$ 为实无穷素位，还可要求局部次数为 $2$。
\end{enumerate}
\end{proposition}

\begin{proof}
思想是用大阶单位根制造循环 Galois 群，再通过选择素数幂阶数控制局部次数。

先看 $F=\Q$。若取素数 $\ell$ 和足够大的 $r$，则
\[
\Gal(\Q(\zeta_{\ell^r})/\Q)\simeq(\Z/\ell^r\Z)^\times
\]
含有很大的循环 $\ell$-幂部分；当 $\ell\ne2$ 时阶为 $\ell^{r-1}(\ell-1)$，当 $\ell=2$ 时也有可控的 $2$-幂循环子扩张。取相应循环子域即可使扩张次数被指定整数整除。一般 $F$ 的情形，把这样的 $\Q$ 上循环分圆扩张与 $F$ 合成；取阶数足够大，可避开与 $F$ 的有限交，从而保持所需整除性。

第二条使用同一构造的局部版本。对有限多个 $v\in S$，选择素数幂阶的分圆扩张，使在每个 $F_v$ 上的局部扩张次数含有所需的素数幂因子。不同素数幂部分线性无关地合成，局部次数相乘。若 $v$ 为实素位，加入一个全虚二次分圆部分即可把实嵌入局部化为 $\C/\R$，局部次数为 $2$。由于 $S$ 有限，逐一满足后合成即可。
\end{proof}

\begin{proposition}
取素数 $p,\ell$ 与正整数 $n$，当 $\ell=2$ 时设 $n\ge2$。令 $\zeta$ 为 $\ell^n$ 次本原单位根。则：
\begin{enumerate}
\item 若 $p\ne\ell$，则 $\Q_p(\zeta)/\Q_p$ 为无分歧扩张。
\item 若 $p=\ell$，则 $\Q_p(\zeta)/\Q_p$ 为全分歧扩张，且
\[
[\Q_p(\zeta):\Q_p]=p^{n-1}(p-1),
\]
$\zeta-1$ 是 $\Q_p(\zeta)$ 的素元，
\[
N_{\Q_p(\zeta)/\Q_p}(1-\zeta)=p,
\]
并且
\[
N_{\Q_p(\zeta)/\Q_p}\Q_p(\zeta)^\times
=\langle p\rangle\times U_{\Q_p}^{(n)}.
\]
\end{enumerate}
\end{proposition}

\begin{proof}
若 $p\ne\ell$，则 $\ell^n$ 次单位根的阶与剩余特征互素。它们可在未分歧扩张的剩余域中出现：只需取 $f$ 使 $\ell^n\mid p^f-1$，则 $\F_{p^f}^\times$ 含有 $\ell^n$ 次本原单位根。Hensel 提升把剩余域中的单位根唯一提升到无分歧扩张中，所以 $\Q_p(\zeta)/\Q_p$ 无分歧。

若 $p=\ell$，分圆多项式
\[
\Phi_{p^n}(X)=\frac{X^{p^n}-1}{X^{p^{n-1}}-1}
\]
在 $X=Y+1$ 后成为 Eisenstein 多项式。于是 $Y=\zeta-1$ 是素元，扩张全分歧，次数为 $\varphi(p^n)=p^{n-1}(p-1)$。范数公式来自常数项：
\[
N(1-\zeta)=\Phi_{p^n}(1)=p.
\]
最后的范数群公式是上一小节 Lubin--Tate 计算在 $\Q_p$、$f(X)=(1+X)^p-1$ 下的特例：
\[
N\Q_p(\zeta)^\times=\langle p\rangle\times U_{\Q_p}^{(n)}.
\]
\end{proof}

\begin{proposition}
设 $\zeta$ 为 $\ell^n$ 次本原单位根，当 $\ell=2$ 时取 $n\ge2$。对任意 $a\in\Q^\times$，有乘积公式
\[
\prod_v (a,\Q_v(\zeta)/\Q_v)=1,
\]
其中 $v$ 遍历 $\Q$ 的所有素位，实素位处按 $\C/\R$ 的范剩符号理解，且除有限多个 $v$ 外因子为 $1$。
\end{proposition}

\begin{proof}
只需看 $\zeta$ 上的作用。若 $v=p\ne \ell,\infty$，扩张无分歧，故
\[
(a,\Q_p(\zeta)/\Q_p)(\zeta)=\zeta^{p^{v_p(a)}}.
\]
若 $p=\ell$，写 $a=up^{v_p(a)}$，由前一节显式公式，作用为
\[
\zeta\longmapsto \zeta^r,\qquad r\equiv u^{-1}\pmod{\ell^n}.
\]
若 $v=\infty$，当 $a>0$ 时实处作用平凡，当 $a<0$ 时为复共轭，即 $\zeta\mapsto \zeta^{-1}$。

把所有局部作用相乘，指数为
\[
\operatorname{sgn}(a)\cdot
\prod_{p\ne\ell}p^{v_p(a)}\cdot r
\quad\text{在 }(\Z/\ell^n\Z)^\times\text{ 中。}
\]
由有理数分解
\[
a=\operatorname{sgn}(a)\prod_p p^{v_p(a)}
\]
以及 $r\equiv (a\ell^{-v_\ell(a)})^{-1}\pmod{\ell^n}$，上式同余于 $1$。因此所有局部范剩符号的乘积在 $\zeta$ 上作用平凡。分圆扩张由 $\zeta$ 生成，所以乘积为恒等自同构。
\end{proof}

\begin{proposition}
若 $a\in\Q^\times$，则它在所有有限素位和无穷素位上的局部范剩符号乘积为 $1$，这正是分圆域上的全局乘积公式。
\end{proposition}

\begin{proof}
这是上一命题的直接重述，但把结论单独提出来更方便后面引用。局部互反律把每个素位的作用写成对单位根的幂作用；乘积公式说这些局部幂指数在全局上总和为零，因此总作用为恒等。换言之，局部互反映射并不是彼此孤立的，它们在全局分圆扩张中彼此抵消。
\end{proof}

\section*{习题详解}
\addcontentsline{toc}{section}{习题详解}

\paragraph{习题1.}
设 $F$ 为 $p$ 进局部域。称 $F^\times$ 的子群 $N$ 为范子群，若存在有限正规扩张 $E/F$ 使
\[
N=N_{E/F}E^\times.
\]
证明：范子群正好是开且有限指数的子群；有限 Abel 扩张 $E/F$ 与范子群 $N_{E/F}E^\times$ 一一对应；并比较开、闭、有限指数条件。

\emph{解答。}
先证明范子群是开且有限指数。局部互反律给出
\[
F^\times/N_{E/F}E^\times\simeq \Gal(E/F)^{\operatorname{ab}}.
\]
右边有限，所以指数有限。范数映射连续，$E^\times$ 中的单位过滤在范数下包含某个 $U_F^{(n)}$；等价地，有限扩张的范数群含有 $F^\times$ 的一个开邻域。因此范数群开。

反过来，设 $N\subset F^\times$ 开且有限指数。局部类域论的存在定理断言存在唯一有限 Abel 扩张 $E/F$，使
\[
N=N_{E/F}E^\times.
\]
这里唯一性也可从互反律看出：全局局部互反映射
\[
F^\times\to G_F^{\operatorname{ab}}
\]
把开有限指数子群 $N$ 对应到开子群，开子群的固定域就是有限 Abel 扩张 $E$；该扩张的范数群正是 $N$。

双射性如下。若 $E_1,E_2$ 是有限 Abel 扩张且范数群相同，则互反律给出
\[
\Gal(E_i/F)\simeq F^\times/N_{E_i/F}E_i^\times.
\]
更精确地，两个扩张对应于 $G_F^{\operatorname{ab}}$ 的同一个开子群，所以固定域相同。若 $N_1\subset N_2$，对应的扩张满足 $E_2\subset E_1$，这是 Galois 对应的反向性。

最后比较拓扑条件。局部域乘法群
\[
F^\times\simeq \langle\pi\rangle\times U_F
\]
是局部紧群。任意开子群自动闭且指数有限的结论对一般局部紧群不总成立，但在这里成立：开子群在离散因子 $\Z$ 中给出 $d\Z$，在紧因子 $U_F$ 中给出开子群，紧群的开子群指数有限。闭且有限指数的子群也是开，因为商群有限离散，商映射连续，单位元的逆像开。因此在 $F^\times$ 中，开有限指数、闭有限指数、有限指数且含某个 $U_F^{(n)}$ 是同一类对象，也就是范子群。

\paragraph{习题2.}
设 $\widehat{\Z}=\varprojlim_n\Z/n\Z$，$F^{\mathrm{nr}}$ 为局部数域 $F$ 的最大无分歧扩张。说明
\[
\Gal(F^{\mathrm{nr}}/F)\simeq\widehat{\Z},
\]
并解释 Frobenius 的提升如何给出互反映射在无分歧部分的公式。

\emph{解答。}
对每个 $n\ge1$，$F$ 有唯一的 $n$ 次无分歧扩张 $F_n$。若 $m\mid n$，则 $F_m\subset F_n$。于是
\[
F^{\mathrm{nr}}=\bigcup_n F_n,\qquad
\Gal(F^{\mathrm{nr}}/F)=\varprojlim_n\Gal(F_n/F).
\]
每个 $\Gal(F_n/F)$ 都是由 Frobenius $\varphi_n$ 生成的循环群，阶为 $n$，因此
\[
\Gal(F_n/F)\simeq\Z/n\Z.
\]
限制映射 $\Gal(F_n/F)\to\Gal(F_m/F)$ 对应于自然投影 $\Z/n\Z\to\Z/m\Z$，故极限为 $\widehat{\Z}$。

令 $\varphi_F=(\varphi_n)_n$。对 $a\in F^\times$，无分歧范剩符号在每个 $F_n$ 上为
\[
(a,F_n|F)=\varphi_n^{v_F(a)}.
\]
取射影极限得到
\[
(a,F^{\mathrm{nr}}|F)=\varphi_F^{v_F(a)}.
\]
因此单位群 $U_F$ 正是无分歧互反映射的核，而素元 $\pi$ 映到 Frobenius。

\paragraph{习题3.}
设 $\widehat{F^{\mathrm{nr}}}$ 为最大无分歧扩张的完备化，$E/F$ 为有限全分歧 Abel 扩张。证明全分歧部分的互反映射可在 $\widehat{F^{\mathrm{nr}}}E$ 中用 Frobenius 的提升来描述，并说明元素 $b$ 满足的范数方程如何产生给定自同构。

\emph{解答。}
记 $L=\widehat{F^{\mathrm{nr}}}$。扩张 $LE/L$ 仍为有限全分歧 Abel 扩张。由于 $L$ 的剩余域是代数闭包的完备化，$L$ 的无分歧部分已经消失；剩下的 Galois 信息全部来自全分歧扩张。

取 $\sigma\in\Gal(E/F)$。把 $F$ 上 Frobenius 的提升记为 $\Phi$，它作用在 $L$ 上，并可延拓到 $LE$。寻找 $b\in (LE)^\times$ 使
\[
\frac{\Phi(b)}{b}
\]
落在与 $\sigma$ 对应的 $1$-余循环中。Hilbert 90 说明这样的余循环若范数为 $1$，就可写成 $\Phi(b)/b$ 的形式。这里的范数条件来自 $\sigma$ 在有限 Abel 群中的关系：沿一个完整轨道相乘必须回到 $1$。

一旦找到 $b$，定义
\[
a=N_{LE/L}(b)
\]
或在相应有限层取范数下降到 $F^\times$。局部互反律的相容性说明 $a$ 的范剩符号就是 $\sigma$。证明中每一步都使用同一个原则：在完备无分歧基变换后，Frobenius 控制下降数据；Hilbert 90 把一维余循环写成商 $\Phi(b)/b$；范数把这个下降数据压回 $F^\times$。

\paragraph{习题4.}
设 $K$ 为完备离散赋值域，$\chi:\Gal(K^{\mathrm{sep}}/K)\to\Q/\Z$ 为有限阶非分歧特征标，$\varpi$ 为素元。说明循环代数
\[
\{\chi,\varpi\}
\]
的 Brauer 不变量，并判定何时为除代数。

\emph{解答。}
非分歧特征标 $\chi$ 只依赖于
\[
\Gal(K^{\mathrm{nr}}/K)\simeq\widehat{\Z}
\]
上的 Frobenius。若 $\chi$ 的阶为 $n$，则它对应于 $n$ 次无分歧扩张 $L/K$ 的一个生成特征标，满足
\[
\chi(\varphi_{L/K})=\frac{r}{n}+\Z
\]
且 $(r,n)=1$ 当 $\chi$ 阶正好为 $n$。

循环代数 $\{\chi,\varpi\}$ 对应于杯积
\[
\bar\varpi\cup \delta\chi\in H^2(\Gal(L/K),L^\times).
\]
按无分歧不变量的定义，$\bar\varpi$ 的赋值为 $1$，所以
\[
\operatorname{inv}_K\{\chi,\varpi\}=\chi(\varphi_{L/K})=\frac{r}{n}+\Z.
\]
中心单代数在局部域上由 Brauer 不变量决定。其指数等于不变量在 $\Q/\Z$ 中的阶。若 $\chi$ 阶为 $n$ 且 $(r,n)=1$，则不变量阶为 $n$，循环代数的指数等于次数 $n$，因此它是除代数。若不变量阶严格小于次数，则代数分裂成矩阵代数上的较小除代数，不是次数 $n$ 的除代数。

\paragraph{习题5.}
设 $K$ 为局部域，$\mu$ 为 $K^{\mathrm{sep}}$ 中所有单位根，$\eta:\Br(K)\to\Q/\Z$ 为局部不变量映射。对有限阶特征标 $\chi$ 与 $K^\times$ 的元素 $\theta$，定义配对
\[
(\chi,\theta)_K=\eta(\chi\cup \theta).
\]
证明该配对与局部互反映射互为对偶，并判定非分歧特征标。

\emph{解答。}
把 $\chi$ 看作
\[
H^1(G_K,\Q/\Z)=\operatorname{Hom}_{\mathrm{cont}}(G_K,\Q/\Z)
\]
中的类，把 $\theta\in K^\times$ 通过 Kummer 或乘法群的零维上同调看作与 $H^0(G_K,K^{\mathrm{sep}\times})$ 相关的类。杯积
\[
\chi\cup \theta\in H^2(G_K,K^{\mathrm{sep}\times})=\Br(K)
\]
再经 $\eta$ 得到 $\Q/\Z$ 中元素。

局部互反映射给出连续同态
\[
\operatorname{rec}_K:K^\times\longrightarrow G_K^{\operatorname{ab}}.
\]
Tate 对偶性和基本类定义说明
\[
(\chi,\theta)_K=\chi(\operatorname{rec}_K(\theta)).
\]
这就是“互为对偶”的精确含义：固定 $\theta$ 时，配对在所有特征标上的取值等于评价 $\operatorname{rec}_K(\theta)$；固定 $\chi$ 时，它是 $K^\times$ 上的连续特征。

现在判定非分歧特征标。$\chi$ 非分歧当且仅当它在惯性群上为零，也就是它只通过
\[
G_K^{\operatorname{ab}}\twoheadrightarrow \Gal(K^{\mathrm{nr}}/K)\simeq\widehat{\Z}
\]
分解。由无分歧互反公式，单位群 $U_K$ 映到惯性部分，且在最大无分歧商中为核。因此
\[
\chi\ \text{非分歧}
\quad\Longleftrightarrow\quad
(\chi,u)_K=0\ \text{对所有 }u\in U_K.
\]
若 $\varpi$ 是素元，则
\[
(\chi,\varpi)_K=\chi(\varphi_K),
\]
所以非分歧特征标完全由它在素元上的配对值决定。这也说明映射
\[
\chi\longmapsto \bigl(\theta\mapsto(\chi,\theta)_K\bigr)
\]
把有限阶 Galois 特征标嵌入 $K^\times$ 的 Pontryagin 对偶；局部类域论进一步说明它的像正是连续有限阶特征标。
 \chapter{理元类的上同调群}

第六章完成了局部域的类成原则。本章把这些局部结果粘合到数域上，目标是证明理元类群
\[
C_F=\A_F^\times/F^\times
\]
满足类成原则，从而得到整体互反律、类域存在性定理、Artin 互反律和 Hilbert 类域定理。这里的学习重点不是背诵图表，而是看清楚三层结构：
\[
\text{局部乘法群}\quad\longrightarrow\quad \text{理元群}\quad\longrightarrow\quad \text{理元类群}.
\]
局部类域论告诉我们每个完备化 $F_v$ 上怎么互反；理元群把所有 $v$ 的信息同时记录；理元类群再除去全局主元，正好留下全局 Abel 扩张的信息。

\section{理元的上同调群}

设 $E/F$ 为有限 Galois 数域扩张，$G=\Gal(E/F)$。若 $a=(a_P)_P\in\A_E^\times$，定义 $G$ 在理元群上的作用为
\[
(\sigma a)_P=\sigma\bigl(a_{\sigma^{-1}P}\bigr).
\]
这个公式只是说：$\sigma$ 既移动素位，也作用在对应局部域的元素上。若 $a=(a_v)_v\in\A_F^\times$，则把它嵌入 $\A_E^\times$ 时，在每个 $P|v$ 处取同一个局部分量 $a_v$。

对 $F$ 的素位 $v$，选定 $E$ 的一个素位 $P|v$，记分解群为
\[
G_P=\{\sigma\in G:\sigma P=P\}.
\]
它自然同构于局部扩张 $E_P/F_v$ 的 Galois 群。写
\[
E_v^\times=\prod_{P|v}E_P^\times,\qquad U_{E,v}=\prod_{P|v}U_P.
\]
若 $S$ 是包含所有无穷素位的有限素位集，则
\[
\A_{E,S}^\times=\prod_{v\in S}E_v^\times\times\prod_{v\notin S}U_{E,v}.
\]
这些群随 $S$ 增大而并起来就是 $\A_E^\times$。

\begin{proposition}
设 $E/F$ 为有限 Galois 数域扩张。则
\[
(\A_E^\times)^G=\A_F^\times,\qquad E^\times\cap \A_F^\times=F^\times.
\]
\end{proposition}

\begin{proof}
若 $a=(a_P)_P\in(\A_E^\times)^G$，则对每个 $v$ 和每个 $P|v$，分量 $a_P$ 被 $G_P$ 固定，所以
\[
a_P\in E_P^{G_P}=F_v.
\]
若 $P'=\sigma P$，由 $G$-不变性得到 $a_{P'}=\sigma(a_P)$；而 $a_P\in F_v$ 时 $\sigma(a_P)$ 正是同一个 $F_v$ 元素在 $E_{P'}$ 中的像。因此所有 $P|v$ 的分量来自同一个 $a_v\in F_v^\times$，且限制直积条件也下降到 $F$，所以 $a\in\A_F^\times$。反向包含由定义立即得到。

第二个等式中，若 $x\in E^\times\cap\A_F^\times$，则 $x$ 在所有局部嵌入 $E_P$ 中都落在 $F_v$。特别地，对每个 $\sigma\in G$，$x$ 和 $\sigma x$ 在所有完备化中的像相同，因此 $x=\sigma x$。故 $x\in E^G=F$。反向包含来自对角嵌入：若 $x\in F^\times$，则它作为 $E^\times$ 的元素当然属于 $E^\times$，而它在每个 $F_v$ 中的局部分量也给出 $\A_F^\times$ 的元素，所以 $F^\times\subset E^\times\cap\A_F^\times$。
\end{proof}

\begin{proposition}
设 $E/F$ 为有限 Galois 数域扩张，$v$ 为 $F$ 的素位，$P|v$。则
\[
\widehat H^q(G,E_v^\times)\simeq \widehat H^q(G_P,E_P^\times).
\]
若 $v$ 为有限素位且 $E_P/F_v$ 无分歧，则
\[
\widehat H^q(G,U_{E,v})=0
\]
对所有整数 $q$ 成立。
\end{proposition}

\begin{proof}
所有 $P|v$ 构成 $G/G_P$。因此
\[
E_v^\times=\prod_{\sigma\in G/G_P}E_{\sigma P}^\times
\]
作为 $G$-模就是从 $G_P$-模 $E_P^\times$ 诱导到 $G$ 的模。Shapiro 引理给出第一式。

同样把单位部分写成诱导模，
\[
U_{E,v}=\prod_{\sigma\in G/G_P}U_{\sigma P}
\]
是由 $U_P$ 诱导而来，所以
\[
\widehat H^q(G,U_{E,v})\simeq \widehat H^q(G_P,U_P).
\]
若 $E_P/F_v$ 无分歧，第六章已经证明右边为零。
\end{proof}

\begin{proposition}
设 $S$ 为 $F$ 的有限素位集，包含所有无穷素位和所有在 $E/F$ 中分歧的有限素位。对每个 $v$ 选定 $P|v$。则：
\[
\widehat H^q(G,\A_{E,S}^\times)\simeq\prod_{v\in S}\widehat H^q(G_P,E_P^\times).
\]
特别地，
\[
\widehat H^1(G,\A_{E,S}^\times)=0,
\]
并且
\[
\widehat H^0(G,\A_{E,S}^\times)\simeq\prod_{v\in S}G_P^{\operatorname{ab}}.
\]
令 $S$ 增大取极限，有
\[
\widehat H^q(G,\A_E^\times)\simeq\prod_v \widehat H^q(G_P,E_P^\times),
\qquad
\widehat H^1(G,\A_E^\times)=0.
\]
\end{proposition}

\begin{proof}
由定义，
\[
\A_{E,S}^\times=\prod_{v\in S}E_v^\times\times\prod_{v\notin S}U_{E,v}.
\]
对 $v\notin S$，扩张在 $v$ 处无分歧，所以上一命题给出单位因子的 Tate 上同调为零。剩下的因子是有限直积，Tate 上同调与有限直积相容，得到第一式。

局部 Hilbert 90 给出
\[
H^1(G_P,E_P^\times)=0,
\]
所以 $\widehat H^1(G,\A_{E,S}^\times)=0$。局部类域论给出
\[
\widehat H^0(G_P,E_P^\times)=F_v^\times/N_{E_P/F_v}E_P^\times
\simeq G_P^{\operatorname{ab}},
\]
于是得到零维公式。

最后，$\A_E^\times=\varinjlim_S\A_{E,S}^\times$。在有限群作用下，上同调可与这种滤过直极限交换；每个素位最终进入 $S$，所以极限给出全体素位的乘积公式。
\end{proof}

\begin{proposition}
设 $E/F$ 为有限 Galois 数域扩张，$a=(a_v)\in\A_F^\times$。存在 $b\in\A_E^\times$ 使
\[
N_{E/F}b=a
\]
当且仅当对每个素位 $v$，存在 $b_P\in E_P^\times$ 使
\[
N_{E_P/F_v}b_P=a_v.
\]
\end{proposition}

\begin{proof}
必要性直接来自范数的逐分量定义。反过来，若每个局部分量都是局部范数，则 $a$ 在
\[
\widehat H^0(G,\A_E^\times)=\A_F^\times/N_{E/F}\A_E^\times
\]
中的类在每个局部商
\[
F_v^\times/N_{E_P/F_v}E_P^\times
\]
中都为零。上一命题把全局零维 Tate 上同调识别为这些局部零维群的乘积，所以 $a$ 的全局类为零，即 $a=N_{E/F}b$。
\end{proof}

若 $F_0\subset F$，并在固定代数闭包中让 $E/F$ 遍历所有有限正规扩张，则膨胀同态
\[
H^2(\Gal(E/F),\A_E^\times)\to H^2(\Gal(K/F),\A_K^\times)
\]
在 $E\subset K$ 时是单射。这来自膨胀-限制正合列和刚才得到的 $H^1=0$。因此可以把
\[
H^2(\Gal(F_0^{\mathrm{alg}}/F),\A_{F_0^{\mathrm{alg}}}^\times)
\]
看作所有有限层 $H^2(\Gal(E/F),\A_E^\times)$ 的并。

\begin{proposition}
设 $F$ 为数域。则
\[
\bigcup_{E/F\ \mathrm{finite\ normal}}H^2(\Gal(E/F),E^\times)
\]
等于只让 $E/F$ 遍历有限循环分圆扩张时得到的并；同样，
\[
H^2(\Gal(F^{\mathrm{alg}}/F),\A_{F^{\mathrm{alg}}}^\times)
\]
也可由有限循环分圆扩张的理元上同调类生成。
\end{proposition}

\begin{proof}
给定一个二阶类 $z$，它已经来自某个有限正规扩张 $L/F$。令 $m$ 为 $z$ 的阶。只有有限多个素位 $v$ 的局部分量 $z_v$ 非零。第六章的分圆构造给出一个有限循环分圆扩张 $E/F$，使得这些有限素位处的局部次数都被 $m$ 整除，实素位处需要时局部次数为 $2$。

设 $K=LE$。若把 $z$ 膨胀到 $K/F$，再限制到 $K/E$，则在每个局部位置，不变量都会乘以相应局部次数。我们选择 $E$ 的方式保证这些乘积为零，所以限制为零。膨胀-限制正合列说明 $z$ 已经来自 $H^2(\Gal(E/F),\A_E^\times)$。这证明了理元版本。把 $\A_E^\times$ 换成 $E^\times$ 时，同样使用局部不变量和分圆扩张的局部次数控制，得到域乘法群版本。
\end{proof}

对 Abel 扩张 $L/F$，定义理元范剩符号
\[
(a,L|F)=\prod_v(a_v,L_w|F_v)\in \Gal(L/F),\qquad a=(a_v)\in\A_F^\times,
\]
其中 $w|v$ 任取。除有限多个 $v$ 外，$a_v$ 是单位且 $L_w/F_v$ 无分歧，所以局部符号为 $1$；乘积有限。由于 $\Gal(L/F)$ 交换，乘积与素位和 $w$ 的选择无关。

还要定义理元不变量映射。对 $c\in H^2(G,\A_E^\times)$，由局部分解得到 $c_v\in H^2(G_P,E_P^\times)$，令
\[
\operatorname{inv}_{E/F}(c)=\sum_v\operatorname{inv}_{E_P/F_v}(c_v)\in\frac1{[E:F]}\Z/\Z.
\]
这是有限和，并且与 $P|v$ 的选择无关。

\begin{proposition}
若 $K\supset E\supset F$ 且 $K/F$ 正规，则理元不变量满足：
\[
\operatorname{inv}_{K/F}(c)=\operatorname{inv}_{E/F}(c)
\]
对 $c\in H^2(\Gal(E/F),\A_E^\times)$ 成立；并且
\[
\operatorname{inv}_{K/E}(\operatorname{Res}c)=[E:F]\operatorname{inv}_{K/F}(c),
\]
\[
\operatorname{inv}_{K/F}(\operatorname{Cor}d)=\operatorname{inv}_{K/E}(d).
\]
\end{proposition}

\begin{proof}
三条公式都是局部类域论中不变量相容性的逐素位求和。第一条中，膨胀只把局部类从 $E_P/F_v$ 看成 $K_Q/F_v$ 的类，不变量不变。第二条中，底域从 $F_v$ 换成 $E_P$，局部不变量乘以 $[E_P:F_v]$；对所有 $P|v$ 求和后，总乘数为
\[
\sum_{P|v}[E_P:F_v]=[E:F].
\]
第三条中，芯限制在局部上对应局部芯限制，而局部不变量与芯限制相容；对所有素位求和即得。
\end{proof}

\begin{lemma}
设 $L/F$ 为有限 Abel 扩张，$a\in\A_F^\times$，$\chi\in\operatorname{Hom}(\Gal(L/F),\Q/\Z)$。若 $[a]$ 是 $a$ 在
\[
\widehat H^0(\Gal(L/F),\A_L^\times)
\]
中的类，$\delta\chi\in H^2(\Gal(L/F),\Z)$ 为连接像，则
\[
\chi((a,L|F))=\operatorname{inv}_{L/F}([a]\cup\delta\chi).
\]
\end{lemma}

\begin{proof}
把公式投到每个素位 $v$。局部类域论给出
\[
\chi_v((a_v,L_w|F_v))
=\operatorname{inv}_{L_w/F_v}([a_v]\cup\delta\chi_v),
\]
其中 $\chi_v$ 是 $\chi$ 在分解群上的限制。理元范剩符号是局部符号的乘积，特征标把乘积化为和；理元不变量也是局部不变量之和。因此逐项相加得到公式。
\end{proof}

\begin{theorem}
若 $E/F$ 为有限正规数域扩张，则
\[
\operatorname{inv}_{E/F}\bigl(H^2(\Gal(E/F),E^\times)\bigr)=0.
\]
\end{theorem}

\begin{proof}
先把问题化到 $\Q$ 上的正规扩张。给定 $E/F$，取包含 $E$ 的有限正规扩张 $N/\Q$。由膨胀和芯限制相容性，若定理对 $N/\Q$ 成立，则对 $E/F$ 成立。

再由前面的分圆化命题，把 $H^2$ 类放到某个循环分圆扩张 $L/\Q$ 上。设 $G=\Gal(L/\Q)$，取生成特征标 $\chi\in H^1(G,\Q/\Z)$。循环群情形中，Tate 同构说明任意
\[
c\in H^2(G,L^\times)
\]
可写成
\[
c=[a]\cup\delta\chi
\]
的形式，其中 $a\in\Q^\times$。

由上一引理，
\[
\operatorname{inv}_{L/\Q}(c)=\chi((a,L|\Q)).
\]
第六章分圆域乘积公式说明全局理元范剩符号中主理元 $a$ 的局部符号乘积为 $1$，即 $(a,L|\Q)=1$。于是 $\operatorname{inv}_{L/\Q}(c)=0$。回到原来的 $E/F$，相容性给出结论。
\end{proof}

\section{计算一阶上同调}

这一节的目标是证明
\[
H^1(\Gal(E/F),C_E)=1
\]
对任意有限 Galois 数域扩张成立。它是整体类成原则的一半。证明分两步：先得到第一不等式，再用 Kummer 扩张证明第二不等式，最后用群论归纳推广到任意正规扩张。

\subsection{第一不等式}

\begin{theorem}
设 $E/F$ 为素数次数 $p$ 的循环扩张。则
\[
\frac{|\widehat H^0(\Gal(E/F),C_E)|}{|H^1(\Gal(E/F),C_E)|}=[E:F].
\]
特别地，
\[
[C_F:N_{E/F}C_E]\ge [E:F].
\]
\end{theorem}

\begin{proof}
取 $S$ 包含所有无穷素位和所有分歧素位，并足够大，使
\[
\A_F^\times=\A_{F,S}^\times F^\times,\qquad
\A_E^\times=\A_{E,\widetilde S}^\times E^\times.
\]
记 $S$ 中在 $E$ 中不分裂的素位数为 $n_0$。因为 $[E:F]=p$，每个 $v\in S$ 要么完全分裂，要么只有一个上方素位，后者分解群阶为 $p$。上一节理元上同调计算给出
\[
h(\A_{E,\widetilde S}^\times)=p^{n_0}.
\]

令
\[
E_{\widetilde S}=E^\times\cap\A_{E,\widetilde S}^\times.
\]
这是 $S$-单位群。Dirichlet 单位定理给出它的秩，从而 Chevalley 的 Herbrand 商公式给出
\[
h(E_{\widetilde S})=p^{n_0-1}.
\]
现在
\[
C_E=\A_E^\times/E^\times\simeq \A_{E,\widetilde S}^\times/E_{\widetilde S},
\]
Herbrand 商在短正合列中相乘，所以
\[
h(C_E)=\frac{h(\A_{E,\widetilde S}^\times)}{h(E_{\widetilde S})}=p.
\]
这就是第一式。又
\[
\widehat H^0(G,C_E)=C_F/N_{E/F}C_E,
\]
所以分子大小至少为 $p=[E:F]$。
\end{proof}

\begin{proposition}
若 $E/F$ 为 $p^r$ 次循环扩张，则 $F$ 中有无穷多个素位在 $E$ 中不完全分裂。
\end{proposition}

\begin{proof}
先设 $[E:F]=p$。若只有有限多个不分裂素位，记它们为 $S$。给定任意类 $\bar a\in C_F$，用逼近定理选择 $x\in F^\times$，使 $ax^{-1}$ 在所有 $v\in S$ 处落入局部 $p$ 次幂或局部范数群。对 $v\notin S$，扩张完全分裂，局部范数映射满。于是由理元范数判别命题，$ax^{-1}$ 是全局理元范数，故 $C_F=N_{E/F}C_E$。这与第一不等式矛盾。

一般 $p^r$ 次循环扩张含唯一 $p$ 次子扩张 $L/F$。若 $E/F$ 只有有限多个不完全分裂素位，则 $L/F$ 也只有有限多个不完全分裂素位，和上一步矛盾。
\end{proof}

\subsection{第二不等式}

\begin{lemma}
设 $F$ 含 $p$ 次单位根，$x\in F^\times$，$E=F(\sqrt[p]{x})$。若有限素位 $v\nmid p$，则：
\begin{enumerate}
\item $v$ 在 $E$ 中无分歧，当且仅当 $x\in U_v(F_v^\times)^p$；
\item $v$ 在 $E$ 中完全分裂，当且仅当 $x\in(F_v^\times)^p$。
\end{enumerate}
\end{lemma}

\begin{proof}
局部化到 $F_v$。因为 $v\nmid p$ 且 $F$ 含 $p$ 次单位根，Kummer 扩张由
\[
X^p-x
\]
控制。把 $x$ 写为
\[
x=u y^p\pi^m,\qquad u\in U_v.
\]
若 $m$ 不被 $p$ 整除，则扩张有分歧。若 $m$ 被 $p$ 整除，则把 $\pi^m$ 吸收到 $y^p$ 中，扩张由单位 $u$ 决定。此时剩余多项式 $X^p-\bar u$ 在剩余域中不可约时给出无分歧扩张，在剩余域中完全分裂时由 Hensel 引理提升为局部完全分裂。于是无分歧等价于 $x\in U_v(F_v^\times)^p$，完全分裂等价于 $x\in(F_v^\times)^p$。
\end{proof}

\begin{lemma}
设 $F$ 含 $p$ 次单位根，$S$ 为包含无穷素位的有限素位集，$F_S$ 为 $S$-单位群。给定 $x_0\in F_S$，令 $E=F(\sqrt[p]{x_0})$。若 $n=|S|$，则存在 $q_1,\dots,q_{n-1}\notin S$，使：
\begin{enumerate}
\item 每个 $q_i$ 在 $E$ 中完全分裂；
\item 若 $x\in F^\times$，且 $S$ 中所有素位在 $F(\sqrt[p]{x})$ 中完全分裂、除 $q_1,\dots,q_{n-1}$ 外的所有 $S$ 外素位无分歧，则 $x\in(F^\times)^p$。
\end{enumerate}
\end{lemma}

\begin{proof}
考虑 Kummer 扩张
\[
K=F(\sqrt[p]{F_S}).
\]
由 Kummer 对应，$\Gal(K/F)$ 是若干个 $p$ 阶循环群的乘积，维数由 $S$-单位模 $p$ 次幂决定。利用上一命题得到的“有无穷多个不分裂素位”，可以在合适的 $p$ 次子扩张中选取素位 $q_i$，使它们的 Frobenius 对偶地检测 $F_S/(F_S)^p$ 的各个坐标，同时又在给定的 $E$ 中完全分裂。

具体地，选择 $q_i$ 使评价映射
\[
F_S/(F_S)^p\longrightarrow\prod_{i=1}^{n-1}U_{q_i}/(U_{q_i})^p
\]
为同构。单射性来自上一局部判别：若 $x$ 在每个 $q_i$ 处是 $p$ 次幂，则 $F(\sqrt[p]{x})$ 在这些地方完全分裂，从而固定域必须退化为 $F$。两边阶相同，故为同构。

若 $x$ 满足第二条条件，先用 $S$-单位部分把它调整到各 $q_i$ 处为局部 $p$ 次幂，再用第一条局部判别与第一不等式的推论排除非平凡 Kummer 扩张。于是 $x$ 是全局 $p$ 次幂。
\end{proof}

\begin{theorem}
设 $E/F$ 为素数次数 $p$ 的循环扩张，且 $F$ 含 $p$ 次单位根。则
\[
[C_F:N_{E/F}C_E]\le [E:F].
\]
\end{theorem}

\begin{proof}
要构造一个子群 $J\subset C_F$，使
\[
J\subset N_{E/F}C_E,\qquad [C_F:J]=p.
\]
取足够大的 $S$ 包含无穷素位、$p$ 上方素位和分歧素位，再取上一引理给出的辅助素位集合 $S'=\{q_1,\dots,q_{n-1}\}$。定义理元子群 $I$：在 $S$ 中的分量取局部 $p$ 次幂，在 $S'$ 中的分量任意，在其他有限素位取单位；令
\[
J=IF^\times/F^\times\subset C_F.
\]

指数计算分两部分。分子来自 $S$ 中局部商
\[
\prod_{v\in S}F_v^\times/(F_v^\times)^p,
\]
由局部单位群结构和乘积公式计算。分母来自 $S\cup S'$-单位群模 $p$ 次幂。辅助素位的选择使额外交集恰好被消去，最后得到
\[
[C_F:J]=p.
\]

还要证明 $J\subset N_{E/F}C_E$。取 $a\in I$。对每个 $v$，逐一检查 $a_v$ 是局部范数。若 $v$ 在 $E$ 中完全分裂，则
\[
E_P\simeq F_v
\]
在每个上方素位上成立，局部范数就是各分量相乘；取一个分量为 $a_v$、其余分量为 $1$，范数便是 $a_v$。若 $v\notin S$ 且无分歧，则 $a_v$ 是单位，而第六章已经证明无分歧扩张中单位全是范数。若 $v\in S$，按 $I$ 的定义 $a_v$ 属于 $(F_v^\times)^p$；对 $p$ 次循环局部扩张，局部范数群含有 $p$ 次幂，因为
\[
N_{E_P/F_v}(y)=y^p\qquad(y\in F_v^\times)
\]
当 $y$ 看作嵌入 $E_P^\times$ 的基域元素时成立。于是所有局部分量都是局部范数。由理元范数判别命题，$a$ 是全局理元范数，所以 $J\subset N_{E/F}C_E$。因此指数至多为 $p$。
\end{proof}

\begin{theorem}
若 $E/F$ 为素数次数 $p$ 的循环扩张且 $F$ 含 $p$ 次单位根，则
\[
H^1(\Gal(E/F),C_E)=1.
\]
\end{theorem}

\begin{proof}
第一不等式给出
\[
[C_F:N_{E/F}C_E]\ge p,
\]
第二不等式给出反向不等式。因此
\[
|\widehat H^0(G,C_E)|=p.
\]
而第一不等式中的 Herbrand 商公式为
\[
\frac{|\widehat H^0(G,C_E)|}{|H^1(G,C_E)|}=p.
\]
代入可得 $|H^1(G,C_E)|=1$。
\end{proof}

\begin{theorem}
对任意有限 Galois 数域扩张 $E/F$，
\[
H^1(\Gal(E/F),C_E)=1.
\]
\end{theorem}

\begin{proof}
对 $|G|$ 归纳。若 $G$ 不是 $p$-群，则对每个 Sylow 子群 $G_p$，归纳假设给出
\[
H^1(G_p,C_E)=1.
\]
由 Sylow 检测，$H^1(G,C_E)=1$。

若 $G$ 是 $p$-群且 $|G|>p$，取正规子群 $H$，使 $G/H$ 阶为 $p$。令 $K=E^H$。归纳假设给出
\[
H^1(H,C_E)=1,\qquad H^1(G/H,C_K)=1.
\]
膨胀-限制正合列于是给出 $H^1(G,C_E)=1$。

剩下 $|G|=p$。若 $F$ 不含 $p$ 次单位根，取 $F'=F(\mu_p)$，令 $E'=EF'$。由于 $[F':F]$ 整除 $p-1$，它比 $p$ 小；对 $F'/F$ 使用归纳，对 $E'/F'$ 使用上一定理，再用膨胀-限制正合列把结论下降到 $E/F$。若 $F$ 已含 $\mu_p$，上一定理 已经处理。
\end{proof}

\section{计算二阶上同调}

本节证明理元类群的二阶上同调由不变量映射完全刻画：
\[
H^2(\Gal(E/F),C_E)\simeq \frac1{[E:F]}\Z/\Z.
\]
这就是整体类成原则的另一半。

\begin{proposition}
若 $E/F$ 为有限 Galois 数域扩张，则
\[
|H^2(\Gal(E/F),C_E)|\mid [E:F].
\]
\end{proposition}

\begin{proof}
对 $|G|$ 归纳。若 $G$ 不是素数幂阶群，则限制到 Sylow 子群后，限制映射在对应 $p$-部分上单射；归纳假设说明每个 $p$-部分的阶整除相应 Sylow 阶，故总阶整除 $|G|$。

若 $G$ 是 $p$-群，取正规子群 $g$ 使 $G/g$ 阶为 $p$，令 $L=E^g$。由上一节 $H^1(g,C_E)=1$，膨胀-限制给出正合列
\[
1\to H^2(G/g,C_L)\to H^2(G,C_E)\to H^2(g,C_E).
\]
归纳假设给出右端阶整除 $|g|$，而左端在素数阶情形阶整除 $p$。故中间群阶整除 $p|g|=|G|$。
\end{proof}

和理元群一样，由膨胀单射可以定义
\[
H^2(\Gal(F^{\mathrm{alg}}/F),C_{F^{\mathrm{alg}}})
=\bigcup_E H^2(\Gal(E/F),C_E).
\]

\begin{proposition}
对循环扩张 $E/F$，序列
\[
1\to H^2(G,E^\times)\to H^2(G,\A_E^\times)
\stackrel{\operatorname{inv}_{E/F}}{\longrightarrow}
\frac1{[E:F]}\Z/\Z\to0
\]
正合。
\end{proposition}

\begin{proof}
满射性先在 $p^r$ 次循环扩张上证明。由前面“有无穷多个不分裂素位”，取一个素位 $v_0$ 使局部次数
\[
[E_{P_0}:F_{v_0}]=[E:F].
\]
在该局部位置取一个局部 $H^2$ 类，其局部不变量为 $1/[E:F]$，其他位置取零。由理元上同调分解，这给出理元上同调类，其总不变量为 $1/[E:F]$。一般循环扩张分解为素数幂部分，再用 Bezout 组合得到满射。

核包含 $H^2(G,E^\times)$，这是上一节定理给出的全局主元不变量为零。反过来，$H^1(G,C_E)=1$ 使短正合列
\[
1\to E^\times\to\A_E^\times\to C_E\to1
\]
给出
\[
1\to H^2(G,E^\times)\to H^2(G,\A_E^\times)\to H^2(G,C_E).
\]
上一步满射和阶数整除 $[E:F]$ 的命题一起说明不变量的核不能比 $H^2(G,E^\times)$ 更大，因此正合。
\end{proof}

\begin{proposition}
若 $E/F$ 为循环扩张，则自然映射
\[
H^2(G,\A_E^\times)\to H^2(G,C_E)
\]
满射，且
\[
|H^2(G,C_E)|=[E:F].
\]
\end{proposition}

\begin{proof}
长正合列中有
\[
H^2(G,\A_E^\times)\to H^2(G,C_E)\to H^3(G,E^\times).
\]
循环群 Tate 上同调周期为 $2$，而 Hilbert 90 给出 $H^1(G,E^\times)=1$，所以
\[
H^3(G,E^\times)\simeq H^1(G,E^\times)=1.
\]
因此映射满射。上一命题把商
\[
H^2(G,\A_E^\times)/H^2(G,E^\times)
\]
识别为 $\frac1{[E:F]}\Z/\Z$，故 $H^2(G,C_E)$ 的阶为 $[E:F]$。
\end{proof}

\begin{theorem}
若 $E/F$ 为有限 Galois 数域扩张，且存在循环扩张 $E'/F$ 满足 $[E':F]=[E:F]$，则
\[
H^2(\Gal(E/F),C_E)\simeq H^2(\Gal(E'/F),C_{E'}).
\]
并且
\[
H^2(\Gal(F^{\mathrm{alg}}/F),C_{F^{\mathrm{alg}}})
\]
由有限循环扩张的二阶上同调类生成。
\end{theorem}

\begin{proof}
取合成域 $K=EE'$。由膨胀-限制正合列，相关的二阶类可放到 $H^2(\Gal(K/F),C_K)$ 中考察。
来自 $E/F$ 的那些类，正是限制到 $K/E$ 后为零的类；对 $E'/F$ 也同样。

循环扩张 $E'/F$ 的二阶类群由上一命题完全控制，阶为 $[E':F]=[E:F]$。若把 $E'/F$ 的类膨胀到 $K/F$，再限制到 $K/E$，其不变量被乘以 $[E:F]$，所以为零；因此它落在来自 $E/F$ 的子群中。这给出单射
\[
H^2(\Gal(E'/F),C_{E'})\hookrightarrow H^2(\Gal(E/F),C_E).
\]
右端阶数整除 $[E:F]$，左端阶数等于 $[E:F]$，故为同构。

第二条由第六章的循环分圆扩张构造给出：任何有限正规层的类都可在同次数循环扩张中比较，所以整个直接极限由循环层生成。
\end{proof}

现在定义 $C_F$ 的不变量映射。先在理元层已有
\[
\operatorname{inv}_F:H^2(\Gal(F^{\mathrm{alg}}/F),\A_{F^{\mathrm{alg}}}^\times)\to \Q/\Z.
\]
短正合列
\[
1\to (F^{\mathrm{alg}})^\times\to\A_{F^{\mathrm{alg}}}^\times\to C_{F^{\mathrm{alg}}}\to1
\]
诱导
\[
j:H^2(\Gal(F^{\mathrm{alg}}/F),\A_{F^{\mathrm{alg}}}^\times)
\to H^2(\Gal(F^{\mathrm{alg}}/F),C_{F^{\mathrm{alg}}}).
\]

\begin{theorem}
映射 $j$ 满射，并且
\[
\operatorname{inv}_F(\ker j)=0.
\]
\end{theorem}

\begin{proof}
给定 $C$-层二阶类，由上一定理 可把它放到某个循环扩张 $E/F$ 上。循环情形中
\[
H^2(G,\A_E^\times)\to H^2(G,C_E)
\]
满射，所以它有理元层提升。这证明 $j$ 满射。

若 $c\in\ker j$，则在某个有限正规扩张 $E/F$ 上，$c\in H^2(G,\A_E^\times)$ 且在 $H^2(G,C_E)$ 中为零。长正合列说明 $c$ 来自 $H^2(G,E^\times)$。上一节的主元不变量为零定理给出
\[
\operatorname{inv}_F(c)=0.
\]
\end{proof}

因此可以定义
\[
\operatorname{inv}_F:H^2(\Gal(F^{\mathrm{alg}}/F),C_{F^{\mathrm{alg}}})\to\Q/\Z
\]
为：先取任意理元上同调提升，再取理元不变量。上一定理 保证定义与提升无关。有限层上记为
\[
\operatorname{inv}_{E/F}:H^2(\Gal(E/F),C_E)\to\frac1{[E:F]}\Z/\Z.
\]

\begin{theorem}
映射
\[
\operatorname{inv}_{E/F}:H^2(\Gal(E/F),C_E)\longrightarrow \frac1{[E:F]}\Z/\Z
\]
是同构。取直接极限得到
\[
\operatorname{inv}_F:H^2(\Gal(F^{\mathrm{alg}}/F),C_{F^{\mathrm{alg}}})\longrightarrow \Q/\Z
\]
也是同构。
\end{theorem}

\begin{proof}
取同次数循环扩张 $E'/F$。上一节已经把 $H^2(\Gal(E/F),C_E)$ 与循环情形识别。循环情形中，二阶群阶为 $[E:F]$，且不变量映射满到 $\frac1{[E:F]}\Z/\Z$；因此它是同构。直接极限版本由
\[
\Q/\Z=\bigcup_n\frac1n\Z/\Z
\]
得到。

相容性也来自理元层的不变量相容：若 $K\supset E\supset F$ 且 $K/F$ 正规，则
\[
\operatorname{inv}_{K/E}(\operatorname{Res}c)=[E:F]\operatorname{inv}_{K/F}(c).
\]
\end{proof}

\section{整体互反律}

\subsection{类成原则}

\begin{theorem}
设 $F_0$ 为数域，$F_0^{\mathrm{alg}}$ 为代数闭包。令
\[
G=\Gal(F_0^{\mathrm{alg}}/F_0).
\]
则 $(G,C_{F_0^{\mathrm{alg}}})$ 满足类成原则：对任意有限正规扩张 $E/F$，
\[
H^1(\Gal(E/F),C_E)=1,
\]
并且
\[
\operatorname{inv}_{E/F}:H^2(\Gal(E/F),C_E)\simeq \frac1{[E:F]}\Z/\Z
\]
与膨胀、限制相容。
\end{theorem}

\begin{proof}
第一条是第二节的 $H^1$ 消失定理。第二条是第三节的不变量同构。相容性已经在理元不变量相容性和 $C_E$ 层不变量定义中验证。
\end{proof}

\subsection{范剩符号}

对有限正规扩张 $E/F$，取基本类
\[
u_{E|F}\in H^2(\Gal(E/F),C_E)
\]
使
\[
\operatorname{inv}_{E/F}(u_{E|F})=\frac1{[E:F]}+\Z.
\]
由 Tate 定理，杯积给出同构
\[
u_{E|F}\cup-:\widehat H^q(\Gal(E/F),\Z)\to
\widehat H^{q+2}(\Gal(E/F),C_E).
\]
取 $q=-2$，得到
\[
r_{E|F}:\Gal(E/F)^{\operatorname{ab}}\longrightarrow C_F/N_{E/F}C_E.
\]
其逆复合自然投影定义整体范剩符号
\[
(\ ,E|F):C_F\longrightarrow \Gal(E/F)^{\operatorname{ab}}.
\]

\begin{proposition}
若 $K\supset E\supset F$ 且 $K/F$ 正规，则整体范剩符号满足投影、迁移、范数和共轭相容性。具体地：
\begin{enumerate}
\item 若 $E/F$ 正规，$\Gal(K/F)^{\operatorname{ab}}\to\Gal(E/F)^{\operatorname{ab}}$ 与 $C_F$ 上符号相容；
\item 包含 $C_F\to C_E$ 对应群论迁移映射；
\item 范数 $N_{E/F}:C_E\to C_F$ 对应 $\Gal(K/E)\to\Gal(K/F)$；
\item 域同构 $\sigma$ 下，符号与共轭扩张相容。
\end{enumerate}
\end{proposition}

\begin{proof}
这是类成原则下基本类的函子性。膨胀、限制、芯限制分别对应投影、迁移和范数；共轭相容来自 Galois 上同调的自然性。杯积与这些操作相容，所以由杯积定义的互反映射和范剩符号也相容。
\end{proof}

\subsection{类域}

把数域 $F$ 的模写成
\[
\mathfrak m=\prod_v v^{n_v},
\]
其中只有有限多个 $n_v>0$，无穷实素位的指数为 $0$ 或 $1$。若 $a=(a_v)\in\A_F^\times$，写
\[
a\equiv1\pmod{\mathfrak m}
\]
表示：有限处 $n_v>0$ 时 $a_v\in1+\mathfrak p_v^{n_v}$；有限处 $n_v=0$ 时 $a_v\in U_v$；实处 $n_v=1$ 时 $a_v>0$。定义
\[
\A_F^{\mathfrak m}=\{a\in\A_F^\times:a\equiv1\pmod{\mathfrak m}\},
\qquad
C_F^{\mathfrak m}=\A_F^{\mathfrak m}F^\times/F^\times.
\]
商 $C_F/C_F^{\mathfrak m}$ 是模 $\mathfrak m$ 的射线类群。

\begin{theorem}
有限 Abel 扩张 $E/F$ 与 $C_F$ 的开有限指数子群 $N$ 反向一一对应：
\[
E\longmapsto N_{E/F}C_E.
\]
并且
\[
\Gal(E/F)\simeq C_F/N_{E/F}C_E.
\]
对每个模 $\mathfrak m$，与 $C_F^{\mathfrak m}$ 对应的扩张称为模 $\mathfrak m$ 的射线类域，满足
\[
\Gal(F^{\mathfrak m}/F)\simeq C_F/C_F^{\mathfrak m}.
\]
任意有限 Abel 扩张都包含在某个射线类域中。
\end{theorem}

\begin{proof}
对有限 Abel 扩张，整体互反律给出满同态
\[
C_F\to\Gal(E/F),
\]
其核为 $N_{E/F}C_E$，所以得到开有限指数子群。

反过来，若 $N\subset C_F$ 开且有限指数，则 $C_F/N$ 是有限 Abel 群。取足够小的模同余子群 $C_F^{\mathfrak m}\subset N$，则 $C_F/N$ 是射线类群 $C_F/C_F^{\mathfrak m}$ 的商。射线类域存在性给出 $F^{\mathfrak m}/F$，取对应商群的固定域即可得到所需 $E$。唯一性由 Galois 对应和互反律给出。

包含关系反向：$E_1\subset E_2$ 时，$N_{E_2/F}C_{E_2}\subset N_{E_1/F}C_{E_1}$。合成和交对应核的交与乘积，这是有限 Abel 群 Galois 对应的直接翻译。
\end{proof}

\begin{theorem}
设 $E/F$ 为有限 Abel 扩张。存在最小的模 $\mathfrak f$，使
\[
C_F^{\mathfrak f}\subset N_{E/F}C_E.
\]
它称为 $E/F$ 的导子。素理想 $\mathfrak p$ 在 $E$ 中分歧，当且仅当 $\mathfrak p|\mathfrak f$。
\end{theorem}

\begin{proof}
所有满足 $C_F^{\mathfrak m}\subset N_{E/F}C_E$ 的模在整除关系下非空，因为上一定理 说明 $E$ 包含在某个射线类域中。两个这样的模取最大公因子仍满足条件：有限处这是局部单位过滤与局部导子的交性质，实处是正性条件的交性质。因此存在最小模 $\mathfrak f$。

若 $\mathfrak p\nmid\mathfrak f$，则在 $\mathfrak p$ 处没有主单位同余条件，局部互反映射的惯性部分落入核，故局部扩张无分歧。若 $\mathfrak p|\mathfrak f$，去掉 $\mathfrak p$ 的指数后同余子群不再落入范数核，说明 $\mathfrak p$ 处的某个单位惯性信息非平凡，故 $\mathfrak p$ 分歧。
\end{proof}

\subsection{Artin 互反律}

设 $E/F$ 为有限 Abel 扩张，取模 $\mathfrak m$ 使 $E\subset F^{\mathfrak m}$。若 $\mathfrak p\nmid\mathfrak m$，则 $\mathfrak p$ 在 $E$ 中无分歧，可以定义 Frobenius
\[
\left(\frac{E/F}{\mathfrak p}\right)\in\Gal(E/F).
\]
对与 $\mathfrak m$ 互素的分式理想群 $I_F^{\mathfrak m}$，定义 Artin 映射为
\[
\left(\frac{E/F}{\cdot}\right):I_F^{\mathfrak m}\to\Gal(E/F),
\qquad
\mathfrak p\mapsto\left(\frac{E/F}{\mathfrak p}\right).
\]

\begin{theorem}
令
\[
i:\A_F^\times\to I_F,\qquad (a_v)\mapsto\prod_{v<\infty}v^{v(a_v)}.
\]
则整体范剩符号与 Artin 映射相容，并诱导同构
\[
I_F^{\mathfrak m}/N_{E/F}(I_E^{\mathfrak m})S_F^{\mathfrak m}\simeq \Gal(E/F),
\]
其中 $S_F^{\mathfrak m}$ 是由 $x\equiv1\pmod{\mathfrak m}$ 的主理想生成的子群。
\end{theorem}

\begin{proof}
对素理想 $\mathfrak p\nmid\mathfrak m$，取理元 $a$，使 $a_{\mathfrak p}$ 为素元，其他有限处为单位。局部互反律在无分歧扩张上给出
\[
(a_{\mathfrak p},E_P|F_{\mathfrak p})=\Frob_{\mathfrak p},
\]
其他有限处贡献为 $1$，无穷处由模条件控制。因此整体范剩符号的像正是 Artin Frobenius。乘法性把素理想情形推广到所有与 $\mathfrak m$ 互素的分式理想。

核由两部分组成：来自 $E$ 的范理想，以及满足模同余的主理想。前者在范剩符号中为零，后者对应 $C_F^{\mathfrak m}$，而 $E\subset F^{\mathfrak m}$ 保证它也在核中。反过来，若理想在 Artin 映射下为零，选择对应理元，它在 $C_F/N_{E/F}C_E$ 中为零，故理想类属于 $N_{E/F}(I_E^{\mathfrak m})S_F^{\mathfrak m}$。于是得到同构。
\end{proof}

记 $I_F$ 为分式理想群，$P_F$ 为主理想群，$\Cl_F=I_F/P_F$。记 $P_F^+$ 为由全正元素生成的主理想群，窄理想类群为
\[
\Cl_F^+=I_F/P_F^+.
\]

\begin{theorem}
数域 $F$ 的窄 Hilbert 类域是 $F$ 的最大 Abel 扩张，其中所有有限素位无分歧。Hilbert 类域是 $F$ 的最大 Abel 扩张，其中所有素位，包括实无穷素位，都无分歧。并且任意理想在 Hilbert 类域中延拓后成为主理想。
\end{theorem}

\begin{proof}
有限素位无分歧对应导子的有限部分为 $1$；若允许实素位分歧，则得到窄类群 $\Cl_F^+$，对应窄 Hilbert 类域。若还要求实素位也无分歧，就除去正性区别，得到普通类群 $\Cl_F$ 和 Hilbert 类域。

Artin 互反律给出
\[
\Cl_F\simeq \Gal(H_F/F).
\]
若 $\mathfrak a$ 是 $F$ 的理想，它在 $H_F$ 中的延拓的 Artin 符号为限制后的恒等元，因为在 $H_F$ 上对应理想已经来自范数核。类群同构说明其理想类为零，所以 $\mathfrak a\mathcal O_{H_F}$ 为主理想。
\end{proof}

\section{Weil 群}

Weil 群把互反律包装成一个群扩张对象。它不是简单地等于绝对 Galois 群，而是在 Abel 化处刻意放入类群 $C_F$，从而让 Frobenius、互反映射和表示论能同时出现。
这一点对初学者很重要：绝对 Galois 群的 Abel 化只看有限 Abel 扩张，而类域论给出的自然源头是 $C_F$。Weil 群的设计目标，就是造一个映到绝对 Galois 群的群 $W_F$，使每个有限层的 Abel 化看起来像对应的类群：
\[
W_E^{\operatorname{ab}}\simeq C_E.
\]
这样一来，局部和整体互反律不再只是许多同构，而被统一成一个群扩张的结构。

若 $G$ 是拓扑群，记 $G^c$ 为交换子群闭包，$G^{\operatorname{ab}}=G/G^c$。若 $H\subset G$ 为有限指数闭子群，迁移映射
\[
\operatorname{Ver}:G^{\operatorname{ab}}\to H^{\operatorname{ab}}
\]
由选取陪集截面并把共轭修正项相乘得到；它是群论中从大群 Abel 化到小群 Abel 化的规范映射。

统一记
\[
C_F=\begin{cases}
\A_F^\times/F^\times,& F\text{ 为数域},\\
F^\times,& F\text{ 为局部域}.
\end{cases}
\]
固定可分闭包 $F^{\mathrm{sep}}$。

\begin{definition}
一个 $F^{\mathrm{sep}}/F$ 的 Weil 群由拓扑群 $W_F$、连续同态
\[
\alpha_F:W_F\to\Gal(F^{\mathrm{sep}}/F)
\]
以及对每个有限扩张 $E/F$ 的同构
\[
r_E:C_E\to W_E^{\operatorname{ab}},\qquad
W_E=\alpha_F^{-1}(\Gal(F^{\mathrm{sep}}/E))
\]
组成，满足以下条件：
\begin{enumerate}
\item $\alpha_F(W_F)$ 在绝对 Galois 群中稠密；
\item 若 $E/F$ 有限 Galois，则 $W_F/W_E\simeq\Gal(E/F)$；
\item $r_E$ 与类域论互反映射相容；
\item 对共轭和有限扩张塔，$r_E$ 与共轭、范数和迁移映射相容。
\end{enumerate}
\end{definition}

\begin{proposition}
Weil 群存在，当且仅当对每个有限 Galois 扩张 $E/F$，存在
\[
\alpha_{E/F}\in H^2(\Gal(E/F),C_E)
\]
使杯积给出所有次数的 Tate 同构
\[
\widehat H^n(\Gal(E/F),\Z)\simeq \widehat H^{n+2}(\Gal(E/F),C_E),
\]
并且这些类在膨胀和限制下满足
\[
\operatorname{Inf}\alpha_{E'/F}=[E:E']\alpha_{E/F},\qquad
\operatorname{Res}\alpha_{E/F}=\alpha_{E/E'}.
\]
\end{proposition}

\begin{proof}
若 $W_F$ 已存在，则对有限 Galois 扩张 $E/F$ 有群扩张
\[
1\to C_E\to W_{E/F}\to\Gal(E/F)\to1,
\]
其中 $W_{E/F}=W_F/W_E^c$。这个扩张对应一个二阶上同调类 $\alpha_{E/F}$。由于 $C_E$ 满足类成原则，这个类就是基本类，杯积给出 Tate 同构；塔相容性由 Weil 群定义中的迁移和范数相容性给出。

反过来，若这样的基本类系统存在，则每个 $\alpha_{E/F}$ 定义一个群扩张 $W_{E/F}$。这里用到群扩张和二阶上同调的标准对应：给定代表余循环 $c(\sigma,\tau)\in C_E$，可把集合 $C_E\times\Gal(E/F)$ 赋予乘法
\[
(x,\sigma)(y,\tau)=(x+\sigma y+c(\sigma,\tau),\sigma\tau),
\]
所得群的扩张类就是 $\alpha_{E/F}$。换代表余循环只会得到同构的扩张。相容性保证当 $E$ 增大时这些扩张之间有投影。取射影极限
\[
W_F=\varprojlim_E W_{E/F}
\]
并赋予极限拓扑，即得到满足 Weil 群公理的对象。
\end{proof}

\begin{proposition}
Weil 群在同构意义下唯一。若 $W_F$ 与 $W'_F$ 都是 Weil 群，则存在同构
\[
\theta:W_F\to W'_F
\]
与到绝对 Galois 群的映射及所有 $r_E$ 相容。
\end{proposition}

\begin{proof}
对每个有限 Galois 扩张 $E/F$，两个有限层扩张
\[
1\to C_E\to W_{E/F}\to\Gal(E/F)\to1
\]
和
\[
1\to C_E\to W'_{E/F}\to\Gal(E/F)\to1
\]
对应同一个基本类，所以存在扩张同构。任意两个这样的同构相差 $C_E/C_F$ 的内共轭。有限层同构组成一个紧的主齐性系统，塔相容映射连续，因此可在射影极限中选出相容族。这个相容族诱导所需的 $\theta$。
\end{proof}

局部 Weil 群可更具体地描述。设 $F$ 为局部数域，$F^{\mathrm{nr}}$ 为最大无分歧扩张。几何 Frobenius 记为 $\Phi_F$。定义
\[
W_F=\{\sigma\in\Gal(F^{\mathrm{sep}}/F):\sigma|_{F^{\mathrm{nr}}}=\Phi_F^n\text{ for some }n\in\Z\}.
\]
它是绝对 Galois 群的稠密子群。惯性群 $I_F=\Gal(F^{\mathrm{sep}}/F^{\mathrm{nr}})$ 作为开子群带原来的拓扑，而 $W_F/I_F\simeq\Z$。

\begin{proposition}
上述局部 $W_F$ 是 $F^{\mathrm{sep}}/F$ 的 Weil 群。若 $E/F$ 有限，则
\[
W_E\simeq W_F\cap\Gal(F^{\mathrm{sep}}/E),
\]
且有限扩张 $E/F$ 与 $W_F$ 的开有限指数子群一一对应。
\end{proposition}

\begin{proof}
限制映射给出
\[
W_F/I_F\simeq\Z,
\]
惯性群的拓扑保证 $W_F\to\Gal(F^{\mathrm{sep}}/F)$ 连续且像稠密。局部互反律给出
\[
F^\times\simeq W_F^{\operatorname{ab}},
\]
并且对有限扩张 $E/F$，同样有 $E^\times\simeq W_E^{\operatorname{ab}}$。范数和迁移相容性正是局部互反律的相容性。

若 $E/F$ 有限，则 $\Gal(F^{\mathrm{sep}}/E)$ 与 $W_F$ 的交是有限指数开子群，且它在 $W_F$ 中的商对应 $E/F$ 的嵌入集合；当 $E/F$ Galois 时商为 $\Gal(E/F)$。反过来，任意开有限指数子群在绝对 Galois 群中的闭包是有限指数闭子群，对应一个有限扩张 $E/F$，该开子群就是对应的 $W_E$。
\end{proof}

\section{注记}

整体互反律可以看作 $GL(1)$ 的 Langlands 对应。若
\[
\rho:\Gal(F^{\mathrm{alg}}/F)\to\C^\times
\]
是一维连续表示，则它通过 Abel 化分解。与范剩符号复合得到
\[
\chi:C_F\to\C^\times.
\]
这就是 Hecke 特征标。互反律保证 Artin $L$-函数与对应 Hecke $L$-函数一致。把 $GL(1)$ 换成 $GL(n)$，问题就变成从 $n$ 维 Galois 表示到自守表示的对应，这正是 Langlands 纲领的入口。

\section*{习题详解}
\addcontentsline{toc}{section}{习题详解}

\paragraph{习题1.}
证明 $\Q(\zeta_m)$ 是 $\Q$ 的适当模的射线类域，并推出 Kronecker--Weber 定理。

\emph{解答。}
对 $\Q$，取模 $\mathfrak m=m\cdot\infty$。射线类群可由
\[
(\Z/m\Z)^\times
\]
描述：与 $m$ 互素的有理素数 $p$ 的理想类对应 $p\bmod m$。分圆域中
\[
\Gal(\Q(\zeta_m)/\Q)\simeq(\Z/m\Z)^\times,\qquad
p\mapsto(\zeta_m\mapsto\zeta_m^p).
\]
这正是 Artin 映射，所以 $\Q(\zeta_m)$ 是模 $m\infty$ 的射线类域。

若 $F/\Q$ 为有限 Abel 扩张，类域存在性定理说明 $F$ 包含在某个 $\Q$ 的射线类域中。刚才已识别射线类域为分圆域，故存在 $m$ 使
\[
F\subset \Q(\zeta_m).
\]
这就是 Kronecker--Weber 定理。

\paragraph{习题2.}
证明整体理元范剩符号等于局部范剩符号的乘积：
\[
((a_v),E|F)=\prod_v(a_v,E_w|F_v).
\]

\emph{解答。}
整体范剩符号由基本类杯积定义。对任意有限特征标 $\chi$，第七章的杯积公式给出
\[
\chi(((a_v),E|F))
=\operatorname{inv}_{E/F}([(a_v)]\cup\delta\chi).
\]
右边按定义是局部不变量之和：
\[
\sum_v\operatorname{inv}_{E_w/F_v}([a_v]\cup\delta\chi_v)
=\sum_v\chi_v((a_v,E_w|F_v)).
\]
这等于 $\chi$ 作用在局部符号乘积上的值。有限 Abel 群的特征标分离元素，所以两个 Galois 元素相等。

\paragraph{习题3.}
证明 Artin 互反律的交换图，并说明素理想情形
\[
\left(\frac{E/F}{\mathfrak p}\right)=(\pi,E_P|F_\mathfrak p).
\]

\emph{解答。}
第一步，理元到理想的映射 $i$ 只记录有限分量的赋值。若理元 $a$ 满足 $a\equiv1\pmod{\mathfrak m}$，则它在 Artin 映射中落入射线主理想子群 $S_F^\mathfrak m$。若 $a=N_{E/F}b$，则 $i(a)=N_{E/F}i(b)$。因此 $i$ 把整体范剩符号的核送入 Artin 映射的核，得到交换图。

第二步，若 $\mathfrak p\nmid\mathfrak m$，取理元 $a$ 在 $\mathfrak p$ 处为素元 $\pi$，其他有限处为单位。局部符号乘积公式中，除 $\mathfrak p$ 外的有限处贡献为 $1$，所以整体符号为
\[
(\pi,E_P|F_\mathfrak p).
\]
无分歧局部互反律说明这就是 Frobenius，因此得到素理想的 Artin 符号。

第三步，由素理想生成与乘法性，Artin 映射在 $I_F^\mathfrak m$ 上满射，核为 $N_{E/F}(I_E^\mathfrak m)S_F^\mathfrak m$，于是得到 Artin 互反律。

\paragraph{习题4.}
证明 Artin 符号关于限制、范数和共轭的相容性。

\emph{解答。}
这些都是局部 Frobenius 的基本相容性。若 $K\supset E\supset F$，且 $\mathfrak a$ 在 $K$ 中无分歧，则 $K/F$ 的 Frobenius 限制到 $E$ 就是 $E/F$ 的 Frobenius，因此
\[
\operatorname{Res}_{E}\left(\frac{K/F}{\mathfrak a}\right)
=\left(\frac{E/F}{\mathfrak a}\right).
\]
若 $E/F$ 有限，$L/F$ Abel，$\mathfrak q| \mathfrak p$，剩余次数为 $f$，则局部剩余域 Frobenius 的幂关系给出
\[
\operatorname{Res}_{L}\left(\frac{LE/E}{\mathfrak q}\right)
=\left(\frac{L/F}{\mathfrak p}\right)^f.
\]
对范数理想，把 $\mathfrak q$ 的范数写成 $\mathfrak p^f$ 即得交换图。共轭相容性来自
\[
\sigma\Frob_\mathfrak p\sigma^{-1}=\Frob_{\sigma\mathfrak p}.
\]

\paragraph{习题5.}
证明 Artin 映射满射，并说明导子的判别条件。

\emph{解答。}
给定有限 Abel 扩张 $E/F$，选模 $\mathfrak m$ 包含所有分歧有限素位和必要的无穷条件。整体互反律给出
\[
C_F/N_{E/F}C_E\simeq\Gal(E/F).
\]
通过理元到理想的映射，$I_F^\mathfrak m$ 的像在 $C_F/C_F^\mathfrak m$ 中生成整个射线类商。具体地，一个与 $\mathfrak m$ 互素的分式理想
\[
\mathfrak a=\prod_{\mathfrak p\nmid\mathfrak m}\mathfrak p^{n_\mathfrak p}
\]
可由理元 $a=(a_v)$ 表示，其中 $a_\mathfrak p$ 取素元的 $n_\mathfrak p$ 次幂，其他有限处取单位，无穷处按模的符号条件取正分量。这个理元在 $C_F/C_F^\mathfrak m$ 中的类只依赖于 $\mathfrak a$ 的射线类。由于整体互反律在 $C_F/N_{E/F}C_E$ 上是满射，且 $C_F^\mathfrak m$ 已被选入范数核，Artin 映射在 $I_F^\mathfrak m$ 上满射。

若增大模的有限指数，使 $S_F^\mathfrak m$ 落入 Artin 核，就得到
\[
S_F^\mathfrak m\subset\ker\left(\frac{E/F}{\cdot}\right)\subset I_F^\mathfrak m.
\]
满足该条件的模有最大公因子，定义为导子 $\mathfrak f$。局部互反律说明 $\mathfrak p$ 的导子指数为 $0$ 当且仅当局部扩张无分歧；因此 $\mathfrak p$ 分歧当且仅当 $\mathfrak p|\mathfrak f$。

\paragraph{习题6.}
设 $S_F^\mathfrak m\subset H\subset I_F^\mathfrak m$。证明存在 Abel 扩张 $E/F$，使 Artin 映射核为 $H$，并推出 Hilbert 类域的分裂判别。

\emph{解答。}
商 $I_F^\mathfrak m/H$ 是有限 Abel 群。通过射线类同构
\[
I_F^\mathfrak m/S_F^\mathfrak m\simeq C_F/C_F^\mathfrak m,
\]
$H$ 对应 $C_F$ 的开有限指数子群。类域存在性定理给出唯一有限 Abel 扩张 $E/F$，其范数核对应这个子群。Artin 互反律说明
\[
I_F^\mathfrak m/H\simeq\Gal(E/F).
\]
分歧素位都整除 $\mathfrak m$，因为 $\mathfrak m$ 外的素理想在射线类群中有 Frobenius 元素，局部互反律在这些素位上只通过无分歧商读出素元的赋值；主单位群落入核，惯性作用为零。

取 $\mathfrak m=1$，$S_F^1=P_F$，得到
\[
I_F/P_F=\Cl_F\simeq\Gal(H_F/F).
\]
所以 $H_F$ 是最大无分歧 Abel 扩张。素理想 $\mathfrak p$ 在 $H_F$ 中完全分裂，当且仅当它的 Frobenius 为 $1$。在上面的同构下，$\Frob_\mathfrak p$ 正是 $\mathfrak p$ 在类群中的类，因此完全分裂等价于 $[\mathfrak p]=1$，也就是 $\mathfrak p$ 为主理想。

\paragraph{习题7.}
设 $n$ 无平方因子且 $n\not\equiv3\pmod4$，令 $K=\Q(\sqrt{-n})$，$H_n$ 为 Hilbert 类域。证明奇素数 $p\nmid n$ 可写成
\[
p=x^2+ny^2
\]
当且仅当 $p$ 在 $H_n$ 中完全分裂，并给出最小多项式判别。

\emph{解答。}
形式 $x^2+ny^2$ 是二次域 $K$ 中元素 $x+y\sqrt{-n}$ 的范数。若
\[
p=x^2+ny^2,
\]
则 $p$ 在 $K$ 中分裂，且上方素理想
\[
\mathfrak p=(p,x+y\sqrt{-n})
\]
为主理想。Hilbert 类域分裂判别说明主理想类的素理想在 $H_n$ 中完全分裂。反过来，若 $p$ 在 $H_n$ 中完全分裂，则它先在 $K$ 中分裂，并且上方素理想在类群中平凡，因此为主理想；取生成元即可得到范数表示。

若 $H_n=K(\alpha)$，$f_n$ 为 $\alpha$ 的最小多项式，则 $p$ 在 $H_n$ 中完全分裂等价于：$p$ 在 $K$ 中分裂，且某个上方素理想在 $H_n$ 中剩余次数 $1$。前者等价于 Kronecker 符号
\[
\left(\frac{-n}{p}\right)=1,
\]
后者等价于 $f_n(X)\equiv0\pmod p$ 有根，并排除判别式整除 $p$ 的坏约化情况。于是得到题中判别。

\paragraph{习题8.}
证明 $\Q(\sqrt{-14})$ 的 Hilbert 类域为
\[
\Q(\sqrt{-14},\sqrt{2\sqrt2-1}).
\]

\emph{解答。}
令 $K=\Q(\sqrt{-14})$，候选域
\[
H=K(\sqrt{2\sqrt2-1})
\]
是 $K$ 的二次扩张。先检查它只在 $K$ 的无穷处和有限分歧候选处可能分歧；由于 $K$ 为虚二次域无实处，重点是有限素位。用二次扩张判别式或局部平方判别可验证 $H/K$ 在所有有限素位无分歧。另一方面，二次型或 Minkowski 界计算给出 $K$ 的类数为 $2$。因此最大无分歧 Abel 扩张次数为 $2$，候选的非平凡无分歧二次扩张就是 Hilbert 类域。

\paragraph{习题9.}
证明 $\Q(\sqrt3)$ 的普通类数为 $1$，窄类数为 $2$，并确定 Hilbert 类域和窄 Hilbert 类域。

\emph{解答。}
实二次域 $F=\Q(\sqrt3)$ 的整数环为 $\Z[\sqrt3]$，判别式为 $12$。Minkowski 界很小，可直接检查每个理想类都有范数 $1$ 的代表，因此普通类数 $h_F=1$。

单位 $2+\sqrt3$ 的两个实嵌入分别为正数和负数的倒数，所以存在范数 $1$ 但不是全正调整所需的单位结构；普通主理想和全正主理想相差一个符号类，故窄类数为 $2$。普通 Hilbert 类域就是 $F$ 自身。窄 Hilbert 类域是唯一只在无穷处允许分歧的二次 Abel 扩张，即
\[
F(i)=\Q(\sqrt3,i).
\]

\paragraph{习题10.}
设全局范剩符号核为 $D_F$，像为 $S_F$。证明 $S_F$ 稠密，并描述 $D_F$。

\emph{解答。}
对每个有限 Abel 扩张 $E/F$，范剩符号诱导满射
\[
C_F\to\Gal(E/F).
\]
取所有有限 Abel 扩张的射影极限，得到像在
\[
G_F^{\operatorname{ab}}
\]
的每个有限商中都满射，因此像稠密。

核为所有有限 Abel 扩张范数群的交：
\[
D_F=\bigcap_E N_{E/F}C_E.
\]
类域存在性说明开有限指数子群都来自范数群，所以 $D_F$ 是 $C_F$ 中落入所有开有限指数商核的部分。它等于单位元连通分支；在数域情形中，这主要来自无穷实正连通和模长方向。由于连通可除群在有限离散商中为零，$D_F$ 是无限可除群，并满足 $D_F^m=D_F$。

\paragraph{习题11.}
证明：给定整数 $a$ 与 $r>1$，若 $q$ 为素数，则存在素数 $p$ 使
\[
a^{q^r}\equiv1\pmod p
\]
且可按需要让阶含有指定 $q$-幂因子。

\emph{解答。}
考虑分圆扩张
\[
\Q(\zeta_{q^r},a^{1/q^r})/\Q.
\]
在合适的 Kummer 扩张中选取 Frobenius，使它在 $q^r$ 次根上产生非平凡指定作用，同时在单位根部分满足所需同余。Chebotarev 密度定理保证存在无穷多个未分歧素数具有这个 Frobenius。对这样的 $p$，$a$ 在 $\F_p^\times$ 中的阶被指定的 $q$-幂控制，从而得到所需同余。

\paragraph{习题12.}
计算 $\Br(\R)$ 与 $\Br(\C)$，并解释实处 Hasse 不变量和符号。

\emph{解答。}
$\C$ 代数闭，所以
\[
\Br(\C)=0.
\]
对 $\R$，唯一有限 Galois 扩张是 $\C/\R$，群为二阶。相对 Brauer 群
\[
H^2(\Gal(\C/\R),\C^\times)
\]
有两个元素：平凡类和四元数代数类 $\mathbb H$。因此
\[
\Br(\R)=\{\mathbb R,\mathbb H\}.
\]
循环代数 $[\C/\R;\chi,-1]$ 正是四元数代数，因为生成元 $j$ 满足 $j^2=-1$ 且 $jz=\bar z j$。

定义 Hasse 不变量时，平凡类取 $1$，四元数类取 $-1$。实 Hilbert 符号满足
\[
(\chi,\theta)_\R=
\begin{cases}
1,&\theta>0,\\
-1,&\theta<0.
\end{cases}
\]
这与 $\R^\times/N_{\C/\R}\C^\times\simeq\{\pm1\}$ 完全一致。

\paragraph{习题13.}
设 $k$ 为数域。说明 Brauer 群局部化、Hasse 不变量以及乘积公式。

\emph{解答。}
中心单代数 $A/k$ 经张量积给出局部中心单代数
\[
A_v=A\otimes_k k_v,
\]
从而有同态
\[
\Br(k)\to\Br(k_v).
\]
局部类域论给出
\[
\Br(k_v)\simeq \Q/\Z
\]
或乘法记号中的单位根群。定义 $h_v(A)$ 为 $A_v$ 的局部不变量。

若 $A$ 是循环代数 $[L/k;\chi,\theta]$，则局部化后是
\[
[L_w/k_v;\chi\circ\rho_v,\theta],
\]
所以
\[
h_v(A)=(\chi\circ\rho_v,\theta)_{k_v}.
\]
全局主元乘积公式和第七章理元主元不变量为零给出
\[
\prod_v h_v(A)=1.
\]
这是 Hasse--Brauer--Noether 局部-整体理论的核心公式。

\paragraph{习题14.}
设 $\chi$ 为 $G_k$ 的有限阶特征标，定义局部特征标 $\chi_v$。证明几乎处处无分歧，并定义全局配对
\[
(\chi,z)_k=\prod_v(\chi_v,z_v)_{k_v}.
\]

\emph{解答。}
有限阶特征标通过某个有限 Abel 扩张 $L/k$ 分解。除有限多个分歧素位外，$L/k$ 无分歧，所以 $\chi_v$ 无分歧；若几乎处处 $\chi_v$ 平凡，则所有未分歧素位的 Frobenius 在 $L/k$ 中都被 $\chi$ 杀掉。Chebotarev 密度定理说明这些 Frobenius 稠密于 $\Gal(L/k)$，故 $\chi$ 平凡。

对 $z=(z_v)\in\A_k^\times$，除有限多个 $v$ 外，$z_v$ 是单位且 $\chi_v$ 无分歧，所以局部符号为 $1$。因此乘积有限，配对良定义。若 $z\in k^\times$ 为主理元，则第六章和第七章的乘积公式给出
\[
(\chi,z)_k=1.
\]
所以配对下降为
\[
G_k^\vee\times \A_k^\times/k^\times\to \mu.
\]
双线性和连续性逐局部成立，有限乘积后仍成立。若配对对所有理元类都为 $1$，则 $\chi$ 在全局互反像上为 $1$；互反像稠密，故 $\chi=1$。对固定理元类的消失与存在有限阶特征标检测有限商相同。

\paragraph{习题15.}
用上题配对定义
\[
a:\A_k^\times\to\Gal(k^{\operatorname{ab}}/k)
\]
并证明它与局部互反、稠密性和核相容。

\emph{解答。}
对 $z\in\A_k^\times$，定义 $a(z)$ 为唯一满足
\[
\chi(a(z))=(\chi,z)_k
\]
的 Abel Galois 元素。有限阶特征标分离 profinite Abel 群元素，因此 $a(z)$ 唯一。若 $j_v:k_v^\times\to\A_k^\times$ 是单处分量嵌入，则
\[
a\circ j_v=\rho_v\circ a_v
\]
由局部符号定义立即得到。

像的闭包在每个有限 Abel 商 $\Gal(E/k)$ 中都是满的。
原因是 Artin 互反律给出同构
\[
C_k/N_{E/k}C_E\simeq\Gal(E/k).
\]
故闭包为整个 $\Gal(k^{\operatorname{ab}}/k)$。

核由所有有限阶特征标同时为 $1$ 的理元组成，即所有有限 Abel 扩张范数群的交。类域论把它识别为 $k^\times$ 与 $\A_k^\times$ 单位元连通分支的闭包。最后，连续有限阶特征标群
\[
\Gal(k^{\operatorname{ab}}/k)^\vee
\]
通过复合 $a$ 与在 $k^\times$ 上平凡的有限阶理元类特征标群同构，这正是 Pontryagin 对偶形式的整体互反律。
 \chapter{对偶定理}

这一章的任务，是把前面已经出现的有限群上同调、Tate 上同调、局部类域论、整体类域论，组织成一套能够同时处理局部和整体 Galois 群的对偶机器。初学时最容易卡住的地方，是看见许多群
\[
H^i(G,A),\qquad H^{2-i}(G,A'),\qquad P^i(G_S,A),\qquad \Sha^i(G_S,A)
\]
却不知道它们为什么要排成这些形状。应当先抓住一条主线：上同调群记录某种方程或粘合问题的障碍；对偶定理说明这些障碍不是散乱出现的，而是可以由另一个上同调群通过配对完全检测出来。

在这一章中，若 $A$ 是交换群，记
\[
A^*=\Hom(A,\Q/\Z).
\]
若 $A$ 是有限交换群，也常写
\[
A'=\Hom(A,\mu),
\]
其中 $\mu$ 是所处代数闭包中的所有单位根。若 $A$ 是局部紧交换群，则 Pontryagin 对偶是
\[
A^\vee=\Hom_{\mathrm{cont}}(A,\R/\Z).
\]
有限群的情形中 $A^*$ 与 $A^\vee$ 没有实质差别；局部紧群中必须保留拓扑条件。

\section{有限群的同调群}

本节把有限群的同调放到 Tate 上同调的同一张坐标系中。前面我们主要使用上同调 $H^i(G,A)$；现在要补上同调 $H_i(G,A)$，因为负次数 Tate 上同调正是同调群的另一种记法。

设 $G$ 是有限群，$A$ 是左 $G$-模。令 $I_G\subset \Z[G]$ 为增广理想，令
\[
A_G=A/I_GA
\]
为余不变量群。直观上，$A^G$ 是把 $G$ 的作用强制为平凡以后留下的共同固定部分；$A_G$ 是把所有 $ga-a$ 都视作零以后得到的最大平凡商。一个是子对象，一个是商对象，二者在 Tate 理论中由范映射连接。

为了定义同调，取自由分解
\[
P_n=\Z[G^{n+1}],\qquad n\geq 0,
\]
其边界映射为
\[
\partial_n(g_0,\ldots,g_n)
=\sum_{i=0}^n(-1)^i(g_0,\ldots,\widehat g_i,\ldots,g_n).
\]
增广映射 $\partial_0:P_0\to \Z$ 把 $\sum a_g g$ 送到 $\sum a_g$。于是得到正合列
\[
\cdots\longrightarrow P_2\longrightarrow P_1\longrightarrow P_0\longrightarrow \Z\longrightarrow 0.
\]
把它与 $A$ 张量，然后取 $G$-余不变量，定义
\[
H_n(G,A)=H_n\bigl((P_\bullet\otimes A)_G\bigr),\qquad n\geq 0.
\]
这个定义的含义可以用最低次数来读。$H_0(G,A)=A_G$，因为第零个同调就是把作用差 $ga-a$ 全部压成零。$H_1(G,A)$ 记录关系之间的关系，和上同调 $H^1$ 记录扭曲同态相对。

Tate 上同调的负次数与同调的关系是
\[
\widehat H^n(G,A)\simeq H_{-n-1}(G,A),\qquad n<-1.
\]
特殊次数为
\[
\widehat H^{-1}(G,A)=\Ker(N_G:A_G\to A^G),\qquad
\widehat H^0(G,A)=A^G/N_GA,
\]
其中
\[
N_Ga=\sum_{g\in G}ga.
\]
这两个式子很重要：负一次与零次是同调和上同调交界处的接缝，许多长正合列和杯积都要经过这里。

若 $U$ 是 $G$ 的正规子群，则投影 $G\to G/U$ 与商 $A\to A_U$ 给出映射
\[
\operatorname{coinf}:H_n(G,A)\longrightarrow H_n(G/U,A_U),
\]
称为上膨胀同态。它的作用是先在大群上写出同调类，再把 $U$ 的作用压掉，使之成为商群的同调类。后面定义射影有限群的负次数 Tate 上同调时，正是用这种映射来取逆极限。

\section{射影有限群的上同调群}

设
\[
G=\varprojlim_U G/U
\]
是射影有限群，$U$ 遍历开正规子群，$A$ 是离散连续 $G$-模。离散连续的意思是每个元素 $a\in A$ 都被某个开子群固定。因此任何有限个元素都已经生活在某个有限商 $G/U$ 上。这个事实是把无限 Galois 群问题化成有限商问题的入口。

\subsection{Tate 上同调群}

对正次数，上同调用膨胀映射形成正向系统：
\[
\widehat H^n(G,A)=\varinjlim_U \widehat H^n(G/U,A^U),\qquad n>0.
\]
对非正次数，Tate 上同调要用上节的上膨胀或范映射来形成逆向系统：
\[
\widehat H^n(G,A)=\varprojlim_U \widehat H^n(G/U,A^U),\qquad n\leq 0.
\]
这样定义看似不对称，实则是由正次数上同调和负次数同调的本性决定的：上同调类可以沿有限商膨胀到更细层，同调类则要沿商映射向下传递。

泛范群定义为
\[
N_GA=\bigcap_U N_{G/U}(A^U)\subset A^G.
\]
它由所有有限层范映射共同约束出来。局部和整体类域论中，“成为所有有限层的范”是非常强的条件，因此泛范群常常可除，甚至是对偶定理里消失的部分。

\subsection{杯积}

杯积是对偶定理的乘法结构。有限群情形有配对
\[
\widehat H^i(G,A)\times H^{q-i}(G,B)\longrightarrow H^q(G,A\otimes B),
\]
在通过有限商取极限时，要检查膨胀、限制、上膨胀、范映射彼此相容。相容性的核心是：先把一个类传到商群再作杯积，和先作杯积再传到商群，得到同一个结果。

\begin{proposition}
设 $G$ 是有限群，$U$ 是正规子群，$A,B$ 是 $G$-模，$i\leq 0$，$q\geq 1$。由杯积、上膨胀和膨胀给出的自然方块交换。
\end{proposition}

\begin{proof}
把 $\widehat H^i$ 在 $i\leq 0$ 的部分改写为同调群。上膨胀由链复形
\[
(P_\bullet(G)\otimes A)_G\longrightarrow (P_\bullet(G/U)\otimes A_U)_{G/U}
\]
诱导。上同调的膨胀由把 $G/U$-上链复形拉回到 $G$-上链复形诱导。杯积在链层面由 Alexander--Whitney 型公式给出：它先取前若干个变量进入一个上链，再取后若干个变量进入另一个上链，然后把两个值在 $A\otimes B$ 中相乘。

需要验证的不是某个抽象符号，而是两个操作对代表元的结果相同。若先取商，所有出现的 $g\in G$ 都被替换为其像 $\bar g\in G/U$，而 $A$ 的元素也取到 $A_U$。若先作杯积再取商，张量中的作用差
\[
ga\otimes gb-a\otimes b
\]
在余不变量中被压成零。两条路径都把同一个链代表送到由 $\bar g$ 与 $A_U,B^U$ 决定的代表，因此方块交换。
\end{proof}

于是对射影有限群也有杯积
\[
\widehat H^i(G,A)\times \widehat H^{q-i}(G,B)
\longrightarrow H^q(G,A\otimes B),
\]
其中 $i\leq 0$，$q\geq 1$，且 $A$ 通常要求为有限生成离散模，以保证固定子模和 Hom 与有限层极限协调。

\subsection{级紧模}

对偶定理需要把离散模与紧群相互转换。设 $G$ 是射影有限群，$A$ 是离散 $G$-模。若对每个开子群 $U$，固定子群 $A^U$ 是紧群，并且当 $V\subset \sigma U\sigma^{-1}$ 时，映射
\[
A^U\longrightarrow A^V,\qquad a\longmapsto \sigma a
\]
连续，则称 $A$ 为级紧 $G$-模。这个条件的作用是保证每个有限层都有拓扑上的良好性，并且不同有限层之间的转移映射连续。

\begin{proposition}
设
\[
0\longrightarrow A'\longrightarrow A\longrightarrow A''\longrightarrow 0
\]
是级紧 $G$-模正合列。若相应的范映射在需要处满射，则 Tate 上同调给出长正合列
\[
\cdots\to \widehat H^m(G,A')\to \widehat H^m(G,A)\to
\widehat H^m(G,A'')\to \widehat H^{m+1}(G,A')\to\cdots .
\]
\end{proposition}

\begin{proof}
在每个有限商 $G/U$ 上，Tate 上同调已经有长正合列。关键是让这些有限层长正合列在取极限后仍然正合。正次数处使用正向极限，交换群范畴中的正向极限保持正合。非正次数处使用逆向极限，逆向极限一般只左正合，因此需要范映射满射或 Mittag--Leffler 型条件来消去一阶导出极限。级紧性正是保证这些转移映射有足够拓扑紧性的条件。

具体地说，若 $x\in \widehat H^m(G,A)$ 在 $\widehat H^m(G,A'')$ 中为零，则在足够小的有限层上，$x$ 的代表已经落在 $A'$ 的像中。有限层长正合列给出提升；转移映射相容使这些提升构成一个极限元素，因此 $x$ 来自 $\widehat H^m(G,A')$。其他位置的核像相等按相同的有限层验证与极限传递完成。
\end{proof}

\subsection{上同调维数}

射影有限群 $G$ 的上同调维数衡量高阶上同调何时消失。若对所有离散挠 $G$-模 $A$ 都有
\[
H^q(G,A)=0,\qquad q>n,
\]
则称 $\operatorname{cd}(G)\leq n$。若对所有离散 $G$-模都成立，则称严格上同调维数 $\operatorname{scd}(G)\leq n$。若只考察 $p$-挠模，则得到 $\operatorname{cd}_p(G)$。

这个概念决定对偶中的数字。例如局部 $p$ 进域的绝对 Galois 群有上同调维数 $2$，所以局部对偶把 $H^i$ 与 $H^{2-i}$ 配成对。若某个群的维数是 $n$，则最高非零维通常是 $n$，迹映射也落在 $H^n$。

\begin{proposition}
设 $U$ 是 $G$ 的开子群。
\begin{enumerate}[label=(\arabic*)]
\item 若 $n=\operatorname{scd}(G)$，则对任意离散 $G$-模 $A$，余限制映射
\[
\operatorname{cor}:H^n(U,A)\longrightarrow H^n(G,A)
\]
满射。
\item 若 $U$ 是开正规子群，则有同构
\[
H^n(U,A)_{G/U}\simeq H^n(G,A).
\]
\end{enumerate}
\end{proposition}

\begin{proof}
取 $A$ 的消解
\[
0\to A\to X^0\to X^1\to\cdots
\]
使 $X^r$ 在上同调计算上平凡。令 $Y^n=\Ker(X^n\to X^{n+1})$。由维数假设，$H^q(G,Y^n)=0$ 与 $H^q(U,Y^n)=0$ 在 $q>0$ 处成立，因此 $H^n(G,A)$ 与 $Y^n$ 的不变量商等价，$H^n(U,A)$ 也有同样的描述。

当 $U$ 正规时，$G/U$ 作用在 $H^n(U,A)$ 上。余限制在链层面由范元
\[
N_{G/U}=\sum_{\bar g\in G/U}\bar g
\]
给出；取余不变量以后，范映射正好把 $U$-层的最高维类送到 $G$-层的最高维类。由于高于 $n$ 的上同调为零，连接同态没有额外障碍，得到同构。非正规开子群的满射可通过取开正规子群 $V\subset U$，再把
\[
H^n(V,A)\to H^n(U,A)\to H^n(G,A)
\]
分解来得到。
\end{proof}

\section{谱序列}

谱序列是本章最节省力气的工具。它的意思不是一次写出答案，而是把复杂对象过滤成许多层，先算每层，再看层与层之间的微分如何修正结果。
如果还没有熟悉谱序列，可以先把它当作“带页码的长正合列”。第一页或第二页记录容易计算的局部信息；随着页数增加，微分逐步消去不该存活的部分；最后留下的 $E_\infty$ 页不是直接等于目标上同调，而是目标上同调某个过滤的分次商。也就是说，谱序列真正告诉我们的是：目标对象可以被一层一层拼起来，每一层由 $E_\infty^{p,q}$ 给出。

设 $A^\bullet$ 是带下降过滤 $F^pA^\bullet$ 的复形，且每一总次数上过滤有限。记
\[
E_1^{p,q}=H^{p+q}(F^pA^\bullet/F^{p+1}A^\bullet).
\]
谱序列写作
\[
E_1^{p,q}\Rightarrow H^{p+q}(A^\bullet).
\]
若来自双复形 $K^{p,q}$，常用的形式是
\[
E_2^{p,q}=H^p\bigl(H^q(K^{\bullet,\bullet})\bigr)
\Rightarrow H^{p+q}(\operatorname{Tot}K^{\bullet,\bullet}).
\]
读这个符号时，$E_2^{p,q}$ 是“先沿竖方向取上同调，再沿横方向取上同调”；极限项给出总复形上同调的分次部分。
本章只需要两种读法。第一，若某一页只剩一列或一行，后续微分没有地方可去，谱序列退化，目标上同调就由这一列或这一行给出。第二，低次数项会给出五项正合列，用来比较 $H^1$、$H^2$ 以及子群和商群的上同调。

为了把这句话落到可操作的层面，先说明“退化”和“边缘映射”在证明中怎样使用。若谱序列位于第一象限，微分的形状为
\[
d_r:E_r^{p,q}\longrightarrow E_r^{p+r,q-r+1}.
\]
因此从低次数项出发时，许多微分因为指标变成负数而不存在。比如总次数 $1$ 的部分只涉及 $E_r^{1,0}$ 和 $E_r^{0,1}$；总次数 $2$ 的部分只涉及 $E_r^{2,0},E_r^{1,1},E_r^{0,2}$。五项正合列正是把这些低次数项中仍然可能出现的微分写出来：
\[
0\to E_2^{1,0}\to H^1\to E_2^{0,1}\xrightarrow{d_2}E_2^{2,0}\to H^2.
\]
这里的 $H^m$ 表示谱序列收敛到的总上同调。这个公式的意义是：$E_2^{1,0}$ 一定能嵌入目标的一次上同调；目标的一次上同调映到 $E_2^{0,1}$ 后，其像正是被 $d_2$ 杀掉的部分；而 $d_2$ 的像正是 $E_2^{2,0}$ 中被送到目标二次上同调前已经消失的部分。

在本章后面的应用中，谱序列不是为了计算某个具体数值，而是为了证明某些自然映射确实是同构。证明路线通常如下：先写出一个谱序列；再证明除一条行、列或对角线外，其余项都为零；最后把唯一剩下的边缘映射解释成杯积、迹映射或局部化映射。这样做的好处是，许多看上去需要逐链追踪的相容性，可以被收敛性质和低次正合列统一处理。

\begin{proposition}
设 $K^{p,q}$ 是第一象限双复形，且对每个 $q$，横向复形
\[
K^{0,q}\to K^{1,q}\to K^{2,q}\to\cdots
\]
在 $K^{1,q}$ 以后正合。令
\[
B^\bullet=\Ker(K^{0,\bullet}\to K^{1,\bullet}).
\]
则总复形的上同调自然同构于 $H^\bullet(B^\bullet)$。
\end{proposition}

\begin{proof}
把双复形转置，先沿横向取上同调。由于横向复形在正次数处正合，横向上同调只剩第零列，也就是
\[
H^0_{\mathrm{hor}}(K^{\bullet,q})=\Ker(K^{0,q}\to K^{1,q})=B^q.
\]
因此与该过滤相伴的谱序列在 $E_2$ 页只剩一列。只剩一列时，不再有非零微分可以进出这些项，谱序列退化，总上同调就是这列复形的上同调，即 $H^\bullet(B^\bullet)$。

还要检查这个同构不是只在抽象层面存在。总复形的 $m$ 次项是
\[
\operatorname{Tot}^mK=\bigoplus_{p+q=m}K^{p,q}.
\]
若一个总闭链写作 $x_0+x_1+\cdots+x_m$，其中 $x_p\in K^{p,m-p}$，横向正合性允许我们从最高的 $p$ 开始逐步减去总边界，把所有 $p>0$ 的分量消去。剩下的代表元位于 $K^{0,m}$，且其横向微分为零，因此它属于 $B^m$。若两个 $B^\bullet$ 中的闭链在总复形中相差一个总边界，同样按横向正合性整理边界，可把这个边界替换为来自 $B^{m-1}$ 的边界。于是闭链和边界两个层面都与 $B^\bullet$ 一致，得到自然同构。
\end{proof}

\subsection{Hochschild--Serre 谱序列}

设 $G$ 是射影有限群，$K$ 是闭正规子群，$A$ 是离散 $G$-模。直觉上，先看 $K$ 的上同调，再让商群 $G/K$ 作用；这应当恢复 $G$ 的上同调。

\begin{theorem}
设 $G$ 是射影有限群，$K$ 是闭正规子群，$A$ 是离散 $G$-模。存在第一象限谱序列
\[
E_2^{p,q}=H^p(G/K,H^q(K,A))\Rightarrow H^{p+q}(G,A).
\]
\end{theorem}

\begin{proof}
取 $A$ 的连续上链消解 $X^\bullet(G,A)$。先取 $K$-不变量得到复形 $(X^\bullet)^K$，再对商群 $G/K$ 取连续上链。由此形成双复形
\[
C^{p,q}=C^p(G/K,(X^q)^K).
\]
沿竖方向先取上同调，得到
\[
H^q((X^\bullet)^K)=H^q(K,A),
\]
于是
\[
E_2^{p,q}=H^p(G/K,H^q(K,A)).
\]
沿横方向先取上同调时，$X^q$ 是诱导型或相对平凡的 $G$-模，故 $(X^q)^K$ 对 $G/K$ 的高阶上同调消失。于是总复形上同调由横向零次列给出，而零次列是
\[
C^0(G/K,(X^\bullet)^K)^{G/K}=(X^\bullet)^G.
\]
它的上同调正是 $H^\bullet(G,A)$。两种计算同一个总复形，得到所需谱序列。
\end{proof}

这个谱序列最常用的低次部分是
\[
0\to H^1(G/K,A^K)\to H^1(G,A)\to H^1(K,A)^{G/K}
\to H^2(G/K,A^K)\to H^2(G,A).
\]
它解释了许多“先限制到子群，再看商群不变量”的操作。

\subsection{Tate 谱序列}

Tate 谱序列把“最高维上同调的对偶”表示成某个新模的上同调。设 $G$ 是射影有限群，$A$ 是离散 $G$-模，且 $\operatorname{cd}(G,A)\leq n$。对闭子群 $H$ 定义
\[
D_j(H,A)=\varinjlim_{U\subset H}\Hom(H^j(U,A),\Q/\Z),
\]
其中 $U$ 遍历 $H$ 的开子群，转移映射来自余限制的对偶。

\begin{theorem}
设 $G$ 是射影有限群，$H$ 是闭子群，$A$ 是离散 $G$-模，且 $\operatorname{cd}(G,A)\leq n$。则存在谱序列
\[
E_2^{p,q}=H^p(G/H,D_{n-q}(H,A))
\Rightarrow \Hom(H^{n-(p+q)}(G,A),\Q/\Z).
\]
\end{theorem}

\begin{proof}
取 $A$ 的消解 $A\to X^\bullet$，截断到长度 $n$：
\[
0\to X^0\to X^1\to\cdots\to X^{n-1}\to Z^n\to 0.
\]
维数假设给出 $H^r(H,Z^n)=0$ 对 $r>0$ 成立，因此这个截断复形在每个开子群上足以计算直到最高维的对偶信息。

对开正规子群 $U$ 取 $U$-不变量，再取 $\Hom(-,\Q/\Z)$，得到有限长度复形 $Y^\bullet_U$。构造双复形
\[
C^{p,q}=C^p(G/U,Y^q_U).
\]
先沿竖方向取上同调，得到 $H^{n-q}(U,A)^*$；再对 $U$ 取正向极限，得到 $D_{n-q}(H,A)$。这给出 $E_2$ 页。

先沿横方向取上同调时，由所选复形的平凡性，正次数横向上同调消失，总上同调化为
\[
H^{p+q}(H^0(G,Y^\bullet_U)).
\]
把 $U$ 的极限传进去，正好得到 $H^{n-(p+q)}(G,A)^*$。这说明该双复形的两种计算分别给出谱序列的 $E_2$ 页和收敛目标。
\end{proof}

当 $H=1$ 时，谱序列的边缘映射给出
\[
H^p(G,D_n(A))\longrightarrow H^{n-p}(G,A)^*.
\]
在 $p=n$ 且 $A=\Z$ 时得到迹映射
\[
\operatorname{Tr}:H^n(G,D_n(\Z))\longrightarrow \Q/\Z.
\]
迹映射是所有后面对偶配对落到 $\Q/\Z$ 的出口。

\begin{theorem}
设 $\operatorname{cd}(G,\Z)\leq n$，$A$ 是有限生成离散 $G$-模，且 $\operatorname{cd}(G,A)\leq n$。由杯积与迹映射得到的映射
\[
H^p(G,D_n(A))\longrightarrow H^{n-p}(G,A)^*
\]
与 Tate 谱序列的边缘映射相同。
\end{theorem}

\begin{proof}
两种映射都来自同一个双复形。边缘映射是在谱序列中从 $E_2^{p,0}$ 走向极限项的自然映射；杯积映射则先用
\[
D_n(A)\times A\longrightarrow D_n(\Z)
\]
把系数相乘，再用迹映射送到 $\Q/\Z$。在链层面，若 $\phi$ 代表 $D_n(A)$ 的类，$a$ 代表 $A$ 的类，则杯积把 $(\phi,a)$ 送到函数 $x\mapsto \phi(ax)$。这正是构造双复形时对偶化截断消解所用的评价映射。

因此两条路线在代表元上给出同一个 $\Q/\Z$ 值。由于上同调类由代表元模去边界决定，而评价映射与微分相容，等式在上同调层成立。
\end{proof}

\section{成对偶模}

本节把上节的构造抽象成一句话：如果存在一个 $G$-模 $D_n$ 使
\[
H^i(G,A)^*\simeq H^{n-i}(G,\Hom(A,D_n)),
\]
那么 $D_n$ 就是维数 $n$ 的成对偶模。它把“取上同调再取对偶”的操作改写成“先取 Hom 到 $D_n$ 再取上同调”。

这个定义背后的动机是表示问题。函子
\[
A\longmapsto H^n(G,A)^*
\]
把一个系数模 $A$ 送到最高维上同调的对偶。若它能写成
\[
A\longmapsto \Hom_G(A,D)
\]
的形式，那么所有关于最高维上同调的问题都可以改成研究一个固定对象 $D$ 中的 $G$-等变映射。换言之，$D$ 收集了最高维上同调能检测到的全部信息。局部域中这个 $D$ 是单位根群 $\mu$；整体情形中它会由 ideles 类群方向的对象替代。

这里的星号 $(-)^*=\Hom(-,\Q/\Z)$ 也有实际作用。$\Q/\Z$ 是内射群，取 $\Hom(-,\Q/\Z)$ 会把有限群变成同阶的对偶群，并把许多右正合信息变成左正合信息。对初学者来说，可以把 $H^n(G,A)^*$ 看成“所有能用 $\Q/\Z$ 测量 $H^n(G,A)$ 的线性函数”。成对偶模定理说，这些测量函数不是凭空来的，它们都由 $A$ 到某个固定模 $D$ 的等变映射给出。

对射影有限群，定义
\[
D_i(A)=\varinjlim_U H^i(U,A)^*.
\]
若 $D_i(\Z)=0$ 且对所有杀死 $A$ 中挠元的整数 $m$，乘 $m$ 在 $D_{i+1}(\Z)$ 上满射，则 $D_i(A)=0$。原因是：把 $A$ 分成自由部分和挠部分。自由部分由 $D_i(\Z)$ 控制；挠部分由正合列
\[
0\to \Z\xrightarrow{m}\Z\to \Z/m\Z\to 0
\]
控制，连接映射显示
\[
D_i(\Z/m\Z)\simeq \operatorname{coker}\bigl(D_{i+1}(\Z)\xrightarrow{m}D_{i+1}(\Z)\bigr),
\]
而该余核在满射条件下为零。

\begin{lemma}
设 $D_i(\Z)=0$。令 $A$ 为有限生成离散 $G$-模。若每个杀死 $A_{\mathrm{tor}}$ 的整数 $m$ 都使
\[
D_{i+1}(\Z)\xrightarrow{m}D_{i+1}(\Z)
\]
满射，则 $D_i(A)=0$。
\end{lemma}

\begin{proof}
取开子群使 $A$ 的有限个生成元都被固定，因而可先按平凡作用理解 $A$。把 $A$ 分解为自由部分与有限挠部分。自由部分是若干个 $\Z$ 的直和，故 $D_i$ 为零。

有限挠部分可由若干 $\Z/m\Z$ 的扩张得到。对
\[
0\to \Z\xrightarrow{m}\Z\to \Z/m\Z\to 0
\]
应用 $D_\bullet$ 的长正合列，在 $D_i(\Z)=0$ 处得到
\[
D_{i+1}(\Z)\xrightarrow{m}D_{i+1}(\Z)\to D_i(\Z/m\Z)\to 0.
\]
乘 $m$ 满射，所以 $D_i(\Z/m\Z)=0$。再用有限次扩张的长正合列，推出 $D_i(A)=0$。
\end{proof}

\begin{theorem}
设 $G$ 是射影有限群，$\operatorname{scd}(G)=n<\infty$。令
\[
D=D_n(\Z).
\]
则对有限生成离散 $G$-模 $A$ 有自然同构
\[
H^n(G,A)^*\simeq \Hom_G(A,D).
\]
若还满足 $D_k(\Z)=0$ 对 $k<n$ 成立，且 $D$ 对相关素数可除，则对所有 $i$ 有自然同构
\[
H^{n-i}(G,A)^*\simeq H^i(G,\Hom(A,D)).
\]
\end{theorem}

\begin{proof}
先处理 $A=\Z[G/U]$。由 Shapiro 引理，
\[
H^n(G,\Z[G/U])\simeq H^n(U,\Z).
\]
取对偶并对开子群极限，得到
\[
H^n(U,\Z)^*\simeq D^U\simeq \Hom_G(\Z[G/U],D).
\]
这里最后一个同构需要展开。一个 $G$-等变映射 $f:\Z[G/U]\to D$ 由 $f(1\cdot U)$ 决定；等变性要求对每个 $u\in U$ 有
\[
u f(1\cdot U)=f(u\cdot U)=f(1\cdot U),
\]
所以 $f(1\cdot U)\in D^U$。反过来，任取 $d\in D^U$，令 $f(gU)=gd$，由于 $d$ 被 $U$ 固定，这个定义与代表元选择无关。于是 $\Hom_G(\Z[G/U],D)\simeq D^U$。这说明置换模情形中，定理只是 Shapiro 引理加上 $D$ 的定义。

一般有限生成 $A$ 可用有限个置换模或自由 $\Z[G/U]$-模给出表示
\[
0\to A\to F_0\to F_1.
\]
对 $H^n(-)^*$ 与 $\Hom_G(-,D)$ 写出两行左正合列。前一步已知 $F_0,F_1$ 上为同构，由五引理的短形式得到 $A$ 上的同构。
这里要注意方向。短正合列在系数变量 $A$ 中给出上同调长正合列；由于 $n$ 是严格上同调维数，$H^{n+1}$ 消失，所以在最高维取对偶后得到左正合片段。另一方面 $\Hom_G(-,D)$ 对第一个变量反变，也给出左正合片段。两行的前两个位置已经由置换模情形识别，交换性来自杯积评价的自然性，因此核也被识别，得到一般 $A$ 的同构。

第二个同构来自 Tate 谱序列。若 $D_k(A)=0$ 在 $k<n$ 处成立，谱序列只剩一条相关对角线，边缘映射即为同构。上一引理把 $D_k(\Z)=0$ 与 $D$ 的可除性转化为 $D_k(A)=0$，所以杯积加迹映射给出所有次数的同构。
更具体地说，Tate 谱序列的 $E_2$ 项形如
\[
E_2^{p,q}=H^p(G,D_{n-q}(A)).
\]
若 $q>0$ 时 $D_{n-q}(A)=0$，则只有 $q=0$ 那一行可能非零，所有进入或离开该行的高页微分都没有落点。因此
\[
H^p(G,D_n(A))\simeq H^{n-p}(G,A)^*.
\]
再用 $D_n(A)\simeq\Hom(A,D_n(\Z))$，就得到定理中写出的形式。这个过程解释了为什么定理不是单独证明每个次数，而是一次性由谱序列给出全部次数。
\end{proof}

\section{类成对偶}

类形成的抽象定理是本章的核心桥梁。它说：若 $G$-模 $C$ 满足类形成条件，特别是对每个有限商都有
\[
\widehat H^2(G,C)\simeq \frac1{|G|}\Z/\Z
\]
并且低维 Tate 上同调具有规定消失，则 $C$ 扮演维数 $2$ 的对偶模。

这里的“维数 $2$”来自类域论的基本类。有限 Galois 扩张 $L/K$ 的类形成把一个特殊的二上同调类
\[
u_{L/K}\in H^2(\Gal(L/K),C_L)
\]
放在中心位置。这个类可以看成扩张 $L/K$ 的“全局或局部互反律压缩包”：与它作杯积，等于把普通的上同调类转换成对偶类。若没有这个基本类，$H^i$ 与 $H^{2-i}$ 之间不会自然相遇；有了基本类，二者的次数刚好加到 $2$，再通过不变量映射落入 $\Q/\Z$。

为理解这个结论，先看有限群 $G$。取 $u\in H^2(G,C)$ 为基本类。构造 $C(u)=C\oplus B$，其中 $B$ 由符号 $b_g$ 生成，并让 $G$ 的作用由代表 $u$ 的二上链控制。这样做的目的，是把基本类 $u$ 在扩大模 $C(u)$ 中变成上同调零类。于是 $C(u)$ 具有平凡 Tate 上同调，长正合列便把 $C$ 的上同调和理想的移位联系起来。

\begin{theorem}
设 $G$ 是有限群，$C$ 是满足 Tate 类形成条件的 $G$-模。令 $u$ 是 $H^2(G,C)$ 中对应 $\frac1{|G|}\Z/\Z$ 的基本类。若 $A$ 是自由有限生成 $\Z$-模，则杯积
\[
\widehat H^i(G,\Hom(A,C))\times \widehat H^{2-i}(G,A)
\longrightarrow \widehat H^2(G,C)\simeq \frac1{|G|}\Z/\Z
\]
诱导同构
\[
\widehat H^i(G,\Hom(A,C))\simeq \widehat H^{2-i}(G,A)^*.
\]
\end{theorem}

\begin{proof}
用基本类 $u$ 构造 $C(u)$。有短正合列
\[
0\to C\to C(u)\to I_G\to 0,
\]
其中 $I_G$ 是增广理想。由构造，$C(u)$ 的 Tate 上同调为零。于是长正合列给出
\[
\widehat H^i(G,\Hom(A,C))\simeq
\widehat H^{i-1}(G,\Hom(A,I_G)).
\]
又由
\[
0\to I_G\to \Z[G]\to \Z\to 0
\]
和 $\Z[G]$ 的诱导性，得到
\[
\widehat H^{i-1}(G,\Hom(A,I_G))
\simeq \widehat H^{i-2}(G,\Hom(A,\Z)).
\]

现在 $\Hom(A,\Z)$ 与 $A$ 之间有自然评价配对。有限生成自由性保证没有 Ext 挠项干扰。通过两次连接同态，上述配对正好变成与基本类 $u$ 的杯积。有限群 Tate 上同调的标准完美性给出目标同构。
把“通过两次连接同态”说得更具体一些。第一条短正合列把 $C$ 替换为 $I_G$，这对应把次数降低一；第二条短正合列把 $I_G$ 替换为 $\Z$，又把次数降低一。两次降低后，原来的
\[
\widehat H^i(G,\Hom(A,C))
\]
被识别到
\[
\widehat H^{i-2}(G,\Hom(A,\Z)).
\]
此时与 $\widehat H^{2-i}(G,A)$ 的杯积落在 $\widehat H^0(G,\Z)$，再由增广映射取值。沿两条短正合列反向追踪，$\widehat H^0(G,\Z)$ 中的这个值正是先与基本类 $u$ 杯积再进入 $\widehat H^2(G,C)$ 的值。因而连接同态没有引入新的配对，只是把同一个评价配对移位了两次。
\end{proof}

\begin{theorem}
设 $G$ 是射影有限群，$C$ 是满足类形成原则的 $G$-模。设 $A$ 为有限生成离散 $G$-模。则对所有整数 $i$，杯积配对
\[
\widehat H^i(G,\Hom(A,C))\times \widehat H^{2-i}(G,A)
\longrightarrow \widehat H^2(G,C)
\]
诱导拓扑同构
\[
\widehat H^i(G,\Hom(A,C))\simeq \widehat H^{2-i}(G,A)^*.
\]
\end{theorem}

\begin{proof}
取开正规子群 $U$ 使 $A=A^U$。在每个有限商 $G/U$ 上应用上一定理。若 $V\subset U$，有限层同构与膨胀、限制、范映射相容；这是杯积与转移映射交换的结果。于是可对有限层取极限。

正次数处是正向极限，非正次数处是逆向极限。类形成原则保证基本类在有限层之间按不变量映射相容，所以右侧的 $\widehat H^2$ 始终识别为相应的 $\frac1{|G/U|}\Z/\Z$ 的极限。由此得到射影有限群层面的拓扑同构。
拓扑同构的含义也要分清。正次数 Tate 上同调按膨胀取正向极限，通常带离散拓扑；非正次数按通缩取逆向极限，通常是紧的射影有限群。Pontryagin 对偶会交换这两种拓扑：离散有限层的正向极限对偶为有限对偶的逆向极限，反之亦然。有限层配对已经是完美配对，并且转移映射相容，所以极限后的配对仍把左侧识别为右侧的连续对偶，而不是只给出集合层面的双射。
\end{proof}

\begin{theorem}
设 $G,C$ 如上，且 $A$ 是有限生成离散 $G$-模。若把 $C$ 的泛范部分记入拓扑结构，则 $C$ 在维数 $2$ 的意义下是成对偶模：对合适的 $A$ 有
\[
\widehat H^i(G,\Hom(A,C))\simeq \widehat H^{2-i}(G,A)^*.
\]
\end{theorem}

\begin{proof}
这是前一定理的成对偶模语言重写。需要说明的是拓扑：左侧由有限层 Tate 上同调取极限，右侧取离散或紧拓扑的 Pontryagin 对偶。类形成原则给出泛范群可除性，因而可除部分在有限对偶中不被检测，剩下的有限层配对成为完美配对。故该同构不仅是抽象群同构，也是拓扑同构。
\end{proof}

\section{局部对偶}

设 $F$ 是局部域，$G_F=\Gal(\overline F/F)$。若 $F$ 是 $p$ 进局部数域，则
\[
\operatorname{cd}(G_F)=2
\]
在有限模范围内成立。局部类域论给出
\[
H^2(G_F,\mu)\simeq \Q/\Z,
\]
其中 $\mu$ 是 $\overline F$ 中所有单位根。于是 $\mu$ 是局部维数 $2$ 的对偶模。

\begin{theorem}
设 $F$ 是局部数域，$A$ 是有限 $G_F$-模。对 $0\leq i\leq 2$，杯积
\[
H^i(G_F,\Hom(A,\mu))\times H^{2-i}(G_F,A)
\longrightarrow H^2(G_F,\mu)\simeq \Q/\Z
\]
诱导有限交换群的同构
\[
H^i(G_F,\Hom(A,\mu))\simeq H^{2-i}(G_F,A)^*.
\]
\end{theorem}

\begin{proof}
由局部类域论，$G_F$ 与 $\overline F^\times$ 形成局部类形成。上一节的类成对偶定理可直接用于 $C=\overline F^\times$，基本类就是局部基本类。Kummer 正合列
\[
1\to \mu_n\to \overline F^\times\xrightarrow{n}\overline F^\times\to 1
\]
把有限系数的对偶模识别为单位根群。由于 $\operatorname{cd}(G_F)=2$，高于二次的上同调消失，配对只发生在 $i=0,1,2$。

有限性来自局部 Galois 上同调的有限性定理。非退化性由类形成对偶给出；具体说，若左侧的类与所有右侧类杯积为零，则它在每个有限层与基本类配对为零，有限群层面的完美性迫使该类为零。若右侧的类与所有左侧类杯积为零，也把它限制到每个有限层，有限层完美配对迫使这些限制全为零，而连续上同调类由某个有限层膨胀而来，所以该右侧类本身为零。
\end{proof}

无限素位也有相应结论。若 $F=\C$，则 $G_F=1$，上同调只剩零次。若 $F=\R$，则 $G_F\simeq \Z/2\Z$，Tate 上同调用二阶循环群计算。

\begin{theorem}
令 $G=\Gal(\C/\R)$，$A$ 为有限生成 $G$-模。杯积
\[
\widehat H^i(G,\Hom(A,\C^\times))\times
\widehat H^{2-i}(G,A)\longrightarrow \widehat H^2(G,\C^\times)
\]
诱导同构
\[
\widehat H^i(G,\Hom(A,\C^\times))\simeq
\widehat H^{2-i}(G,A)^*.
\]
\end{theorem}

\begin{proof}
$G$ 是二阶循环群。对循环群，Tate 上同调以两个算子控制：
\[
N=1+\sigma,\qquad I=1-\sigma.
\]
因此
\[
\widehat H^0(G,M)=M^G/NM,\qquad
\widehat H^{-1}(G,M)=\Ker N/IM,
\]
并且二周期。$\C^\times$ 在复共轭作用下满足
\[
H^2(G,\C^\times)\simeq \R^\times/N_{\C/\R}\C^\times\simeq \{\pm1\}\simeq \frac12\Z/\Z.
\]
对 $A=\Z$ 与符号模 $\Z(1)$ 可直接计算配对非退化。一般有限生成 $G$-模由这些基本模和有限二挠模通过有限次扩张得到；长正合列与五引理把非退化性传递到一般 $A$。
\end{proof}

局部有限模还有 Euler--Poincare 特征公式。若
\[
h^i(F,A)=|H^i(G_F,A)|,
\]
定义
\[
\chi(F,A)=\frac{h^0(F,A)h^2(F,A)}{h^1(F,A)}.
\]
对 $p$ 进局部数域，Tate 特征公式写作
\[
\chi(F,A)=|A|_F^{-1},
\]
其中右侧按 $F$ 的归一化绝对值理解。这个公式说明：局部三个上同调群的大小不是自由变化的，它们被局部域的绝对值精确约束。

\section{整体对偶}

设 $F$ 是数域，$S$ 是包含所有无限素位的素位集。记 $F_S$ 为在 $S$ 外无分歧的最大扩张，
\[
G_S=\Gal(F_S/F).
\]
本节的目标，是把局部对偶拼成整体对偶。拼接并不自动成立，因为局部可解的数据可能没有整体解；这种失败由后面的 $\Sha$ 群记录。

先回顾整体 ideles。令 $\A_F^\times$ 为 ideles 群，$C_F=\A_F^\times/F^\times$ 为 ideles 类群。绝对值给出
\[
|\cdot|:C_F\to \R_{>0}^\times,
\]
其核的单位连通部分记作 $C_F^0$。整体类域论给出正合列
\[
1\to C_F^0\to C_F\longrightarrow \Gal(F^{\mathrm{ab}}/F)\to 1.
\]
这说明 ideles 类群的可见商正是整体 Abel Galois 群，而 $C_F^0$ 是被互反映射杀掉的泛范部分。

在 $S$-限制情形中，定义 $C_S(F)$ 为只保留 $S$ 中局部乘法分量并模去 $S$-单位所得的类群。相应有
\[
1\to C_S(F)^0\to C_S(F)\longrightarrow \Gal(F_S/F)^{\mathrm{ab}}\to 1.
\]
这里 $C_S(F)^0$ 仍然是泛范部分，并且可除。

\begin{theorem}
设 $A$ 是有限生成离散 $G_S$-模，且其挠部分的阶只含 $S$ 中允许的素数。
\begin{enumerate}[label=(\arabic*)]
\item $C_S^1$ 是级紧 $G_S$-模，满足类形成原则，其泛范群可除。
\item 映射 $C_S^1\to C_S$ 诱导
\[
\widehat H^i(G_S,\Hom(A,C_S^1))
\simeq
\widehat H^i(G_S,\Hom(A,C_S)).
\]
\item 对 $i\leq 0$，杯积诱导拓扑同构
\[
\widehat H^i(G_S,\Hom(A,C_S^1))
\simeq \widehat H^{2-i}(G_S,A)^*.
\]
\item 因而有
\[
\widehat H^0(G_S,\Hom(A,C_S))\simeq \widehat H^2(G_S,A)^*.
\]
\end{enumerate}
\end{theorem}

\begin{proof}
第一点来自整体类域论的类形成。有限层 $E/F$ 上，$C_S(E)^1$ 的范像是 $C_S(F)^0$，而这个群可除；类形成的基本类由整体基本类给出，局部不变量和为零保证有限层之间相容。

第二点看短正合列
\[
0\to C_S^1\to C_S\to C_S/C_S^1\to 0.
\]
商 $C_S/C_S^1$ 是实向量群或其有限维商，因而对有限生成离散 $A$，$\Hom(A,C_S/C_S^1)$ 是可除模。可除模在 Tate 上同调中贡献为零，所以 $C_S^1$ 与 $C_S$ 给出同一上同调。

第三点把上一节的类成对偶应用于 $G_S$ 与 $C_S^1$。由于泛范群可除，有限层非退化配对通过极限保留下来。第四点取 $i=0$，再用第二点把 $C_S^1$ 换成 $C_S$。
\end{proof}

整体对偶的含义是：$G_S$ 的二次上同调由 ideles 类群方向的零次上同调完全检测。它是 Poitou--Tate 序列的核心输入。

这一点可以按“检测器”来读。一个类 $x\in H^2(G_S,A)$ 若在整体对偶下被所有
\[
\Hom(A,C_S)^{G_S}
\]
中的元素配对为零，那么整体对偶定理说 $x=0$。所以 ideles 类群不是额外装饰，而是用来测试整体二次上同调的全部探针。Poitou--Tate 序列后面把这些探针再限制到各个局部域，从而把“整体为零”和“每个局部为零”之间的差异精确写成 $\Sha$ 群。

\section[Pi 与 Sha]{$P^i$ 与 $\Sha^i$}

要把局部与整体放进同一条长正合列，必须引入两个对象。第一个是局部直积 $P^i$，它收集所有素位处的局部上同调。第二个是 $\Sha^i$，它是局部化映射的核，记录整体类在所有局部处都消失的现象。

对 $G_S$-模 $A$，定义局部化映射
\[
\lambda^i:H^i(G_S,A)\longrightarrow P^i(G_S,A).
\]
这里
\[
P^i(G_S,A)=\prod_{v\in S}' H^i(F_v,A)
\]
是限制直积，限制条件是在几乎所有有限素位处落入无分歧部分
\[
H^i_{\mathrm{nr}}(F_v,A)=\operatorname{im}\bigl(H^i(\Gal(F_v^{\mathrm{nr}}/F_v),A)
\to H^i(F_v,A)\bigr).
\]
定义
\[
\Sha^i(G_S,A)=\Ker\lambda^i.
\]
这就是“局部处看不见的整体障碍”。例如 $\Sha^1$ 表示某个一上同调类在所有局部域中都是零，但整体上可能不是零。

\begin{proposition}
设 $A$ 是有限生成 $G_S$-模，且满足 $S$ 的挠条件。若 $M=A$ 或 $M=A'=\Hom(A,\mathcal O_S^\times)$，则局部化映射
\[
H^i(G_S,M)\longrightarrow \prod_{v\in S}H^i(F_v,M)
\]
的像包含在 $P^i(G_S,M)$ 中。
\end{proposition}

\begin{proof}
取 $x\in H^i(G_S,M)$。连续上同调类由有限数据表示，因此存在有限 Galois 扩张 $E/F$，$E\subset F_S$，使得 $x$ 从
\[
H^i(\Gal(E/F),M^{\Gal(F_S/E)})
\]
膨胀而来。除有限多个素位外，$E/F$ 在该素位无分歧，且 $M$ 的局部作用也无分歧。于是这些素位处的局部限制类来自未分歧商
\[
\Gal(F_v^{\mathrm{nr}}/F_v),
\]
即落在 $H^i_{\mathrm{nr}}(F_v,M)$ 中。故局部化像确实位于限制直积 $P^i$。
\end{proof}

边界映射
\[
\Hom(A,C_S)^{G_S}\longrightarrow H^1(G_S,A')
\]
的像是有限的，且存在开正规子群 $U$ 使范像被边界映射杀掉。这个事实保证 $\Sha^1$ 是有限群。直观上，这是整体类域论的有限性与局部化条件共同作用的结果：虽然 ideles 类群巨大，但真正进入上同调障碍的部分只有有限多。

\begin{theorem}
设 $A\in \operatorname{Mod}_S(G_S)$。边界映射
\[
\Hom(A,C_S)^{G_S}\longrightarrow H^1(G_S,A')
\]
的像有限，并且 $\Sha^i(G_S,A')$ 是有限群。
\end{theorem}

\begin{proof}
把 $C_S$ 写成有限层 $C_S(E)$ 的极限。类域论给出：对足够大的有限扩张 $E/F$，使 $A$ 与 $A'$ 的作用都在 $E$ 上平凡，并且使 $S$ 中相关素位在适当扩张中完全分裂。此时从更大层到 $E$ 层的范映射在类群方向上把边界像压成零。

因此边界映射通过一个紧商
\[
\Hom(A,C_S)^{G_S}/N_{G_S/U}\Hom(A,C_S)^U
\]
因子化。这个商在当前有限生成假设下是有限群，所以边界像有限。$\Sha^1$ 是该有限像的子商，故有限。其他出现的 $\Sha^i$ 通过局部对偶和整体对偶与它配对或嵌入有限群，由此也为有限群。
\end{proof}

\section{Poitou--Tate 序列}

现在把前面的局部化、局部对偶、整体对偶和 $\Sha$ 群装配成一条长正合列。设 $A$ 是有限 $G_S$-模，且 $|A|\in N(S)$。令
\[
A'=\Hom(A,\mathcal O_S^\times).
\]
局部 Tate 对偶给出
\[
P^i(G_S,A)\simeq P^{2-i}(G_S,A')^*,\qquad 0\leq i\leq 2.
\]
因此可以把局部化映射的对偶接到整体上同调后面。

\begin{lemma}
对 $i=0,1$，由 Shapiro 映射、整体对偶和局部对偶组成的方块交换。
\end{lemma}

\begin{proof}
所有箭头都由同一个杯积配对诱导。沿上方路径，先把整体 ideles 方向的类限制到各个局部域，再与局部上同调类作局部杯积，最后用局部不变量映射
\[
\operatorname{inv}_v:H^2(F_v,\overline F_v^\times)\to \Q/\Z
\]
取值。沿下方路径，先作整体杯积得到 $H^2(G_S,C_S)$，再由整体不变量映射取值。整体不变量就是所有局部不变量的和，而互反律保证该和与限制后的局部杯积一致。因此方块交换。
若把一个整体类记为 $x$，一个局部族记为 $(y_v)_v$，方块交换就是说
\[
\operatorname{inv}\bigl(x\cup y\bigr)=\sum_{v\in S}\operatorname{inv}_v(x_v\cup y_v).
\]
右侧只有有限多个非零项，因为限制直积的定义要求几乎所有有限素位处落在未分歧部分，而未分歧部分与整数单位方向的杯积不产生新的局部不变量。这个有限性保证上式不仅形式上有意义，而且能作为从局部配对到整体配对的桥梁。
\end{proof}

\begin{lemma}
边界映射
\[
\delta:\Hom(A,C_S)^{G_S}\to H^1(G_S,A')
\]
可通过 $H^2(G_S,A)^*$ 分解。
\end{lemma}

\begin{proof}
由整体对偶定理，
\[
\widehat H^0(G_S,\Hom(A,C_S))\simeq H^2(G_S,A)^*.
\]
从 $G_S$-不变量到 Tate 零次上同调的自然商映射把范子群压成零。上一节说明边界映射在某个范子群上为零，因此 $\delta$ 必定通过该商映射因子化。于是存在映射
\[
H^2(G_S,A)^*\longrightarrow H^1(G_S,A')
\]
使得原边界映射等于二者复合。
\end{proof}

\begin{theorem}
设 $A$ 是有限 $G_S$-模，且 $|A|\in N(S)$。杯积诱导有限群的完美配对
\[
\Sha^1(G_S,A')\times \Sha^2(G_S,A)\longrightarrow \Q/\Z.
\]
若 $A$ 是无挠且属于相应的 $S$-模范畴，也有配对
\[
\Sha^1(G_S,A)\times \Sha^2(G_S,A')\longrightarrow \Q/\Z.
\]
\end{theorem}

\begin{proof}
先证明有限模情形。由上一引理，得到正合片段
\[
P^0(G_S,A')\to H^2(G_S,A)^*\to \Sha^1(G_S,A')\to 0.
\]
对偶后得到
\[
0\to \Sha^1(G_S,A')^*\to H^2(G_S,A)\to P^0(G_S,A')^*.
\]
再用局部对偶 $P^0(G_S,A')^*\simeq P^2(G_S,A)$，可见 $\Sha^1(G_S,A')^*$ 正好是局部化
\[
H^2(G_S,A)\to P^2(G_S,A)
\]
的核，也就是 $\Sha^2(G_S,A)$。这给出
\[
\Sha^2(G_S,A)\simeq \Sha^1(G_S,A')^*,
\]
等价于所述完美配对。

无挠情形把 $A$ 与 $A'$ 的角色互换，并使用同样的局部对偶和整体对偶。有限性由上一节已得的 $\Sha$ 有限性保证。
\end{proof}

\begin{theorem}
设 $A$ 是有限 $G_S$-模，且 $|A|\in N(S)$。存在正合序列
\[
0\to H^0(G_S,A)\to P^0(G_S,A)\to H^2(G_S,A')^*
\]
\[
\to H^1(G_S,A)\to P^1(G_S,A)\to H^1(G_S,A')^*
\]
\[
\to H^2(G_S,A)\to P^2(G_S,A)\to H^0(G_S,A')^*\to 0.
\]
\end{theorem}

\begin{proof}
序列中的第一、第四、第七个箭头是局部化映射。第三、六、九个位置来自局部对偶后取对偶的映射。中间的连接映射由整体对偶与边界映射给出。

除中间
\[
H^1(G_S,A)\to P^1(G_S,A)\to H^1(G_S,A')^*
\]
处外，正合性由上一引理、局部对偶和整体对偶直接给出。中间处需要检验局部化像的正交补。局部 Tate 对偶说明 $P^1(G_S,A)$ 与 $P^1(G_S,A')$ 完美配对；整体杯积的局部不变量和为零，说明 $H^1(G_S,A)$ 的像包含在 $H^1(G_S,A')$ 像的正交补中。

反向包含用阶数计算完成。设有限复形 $C^\bullet$ 为上述九项序列。除了中间点，已知正合。局部 Euler--Poincare 公式和整体有限性公式给出该复形的交错阶数为 $1$。因此中间位置的核与像也必须具有同一阶数，包含关系升级为相等，正合性成立。
把最后一步展开如下。设中间处的核为
\[
K=\Ker\bigl(P^1(G_S,A)\to H^1(G_S,A')^*\bigr),
\]
而整体局部化的像为
\[
I=\operatorname{im}\bigl(H^1(G_S,A)\to P^1(G_S,A)\bigr).
\]
前面已经由局部不变量和为零证明 $I\subset K$。九项复形除这一处外都正合，交错阶数为 $1$ 等价于
\[
\frac{|K|}{|I|}=1.
\]
因为 $I$ 是 $K$ 的子群且二者阶数相等，只能有 $I=K$。这就补上了全序列唯一没有由图追踪直接给出的正合点。
\end{proof}

Poitou--Tate 序列的读法是：整体上同调进入所有局部上同调；局部配对的正交补再回到对偶模的整体上同调。它不是三行公式的拼接，而是一条“整体--局部--对偶整体”的闭合回路。
在应用时，常按下面三步读它。先看
\[
H^i(G_S,A)\to P^i(G_S,A)
\]
确定一个整体类在所有局部域中的影子。再用局部 Tate 对偶把 $P^i(G_S,A)$ 的正交补解释成 $P^{2-i}(G_S,A')$ 的条件。最后回到 $H^{2-i}(G_S,A')^*$，判断这些局部条件能否由对偶模的整体类解释。$\Sha$ 群正是这条回路闭合失败的部分。

例如在 $i=1$ 的位置，序列中
\[
H^1(G_S,A)\to P^1(G_S,A)\to H^1(G_S,A')^*
\]
的正合性可以读成一句非常实用的话：一个局部族 $(\xi_v)_v\in P^1(G_S,A)$ 来自整体一上同调类，当且仅当它与所有对偶模 $A'$ 的整体一上同调类的局部化都正交。这里的“正交”是局部 Tate 配对后再把所有局部不变量相加。这样，Poitou--Tate 把求整体解的问题转化为检查所有对偶整体类给出的线性条件。

\section{上同调理论和数论}

本节是概念性收束：为什么代数数论要承受这些上同调技术？答案不是“为了抽象”，而是因为数论对象经常由局部数据、整体约束和对偶障碍同时控制。上同调提供统一语言，能够把这些关系写成可计算的正合列与配对。

\subsection[L-函数]{$L$-函数}

$L$-函数通常是带 Euler 积的 Dirichlet 级数
\[
L(s)=\sum_{n\geq1}\frac{a_n}{n^s}=\prod_p L_p(p^{-s})^{-1}.
\]
最基本的例子是 Riemann zeta 函数
\[
\zeta(s)=\sum_{n\geq1}\frac1{n^s}=\prod_p(1-p^{-s})^{-1}.
\]
标准问题包括解析延拓、函数方程、零点和极点、特殊值、系数 $a_n$ 的渐近。上同调进入这些问题，是因为许多 $a_n$ 来自 Frobenius 在某个上同调空间上的迹。

如果 $X/\Q$ 是光滑射影簇，则 $l$ 进上同调
\[
H^i_{\mathrm{et}}(X_{\overline\Q},\Q_l)
\]
携带 $G_\Q$ 的连续表示。好素数处的局部因子由 Frobenius 的特征多项式给出。这样，几何对象的点数、Frobenius 迹和 $L$-函数被同一个上同调空间连接起来。

\subsection{表示相容系统}

固定维数 $d$。对每个素数 $l$，设有连续表示
\[
\rho_l:G_\Q\to \GL(V_l),
\]
其中 $V_l$ 是 $d$ 维 $\Q_l$-向量空间。若存在有限素位集 $S$，使得对 $q\notin S$ 且 $q\ne l$，表示在 $q$ 处无分歧，并且
\[
\det(1-\rho_l(\Frob_q)T\mid V_l)
\]
属于 $\Q[T]$ 且与 $l$ 无关，则称这是一族 $l$ 进表示相容系统。

几何来源是：取 $X/\Q$ 光滑射影，令
\[
V_l=H^i_{\mathrm{et}}(X_{\overline\Q},\Q_l).
\]
不同 $l$ 的上同调空间不同，但好素数处 Frobenius 多项式相同，因为它们都由同一个几何点数公式控制。这是上同调为 $L$-函数提供算术系数的基本机制。

\subsection[p 进 L-函数]{$p$ 进 $L$-函数}

$p$ 进 $L$-函数把复特殊值插值成 $p$ 进解析函数。设 $\Gamma\simeq \Z_p$，其 Iwasawa 代数是
\[
\Lambda=\varprojlim_n \Z_p[\Gamma/\Gamma^{p^n}]\simeq \Z_p[[T]].
\]
若有相容系统 $V$，理想图像是存在一个 $p$ 进解析函数
\[
\chi\longmapsto L_p(\chi,V),
\]
在有限阶特征 $\chi$ 处给出经周期和 Euler 因子修正后的复值 $L(1,\chi,V)$。

对初学者来说，关键不是立即掌握构造，而是理解“插值”二字：有限阶特征在 $p$ 进单位圆盘中稠密，若一组特殊值满足足够强的同余关系，就可能由一个 $p$ 进解析函数统一表示。

\subsection{岩泽理论}

岩泽理论问：$p$ 进 $L$-函数给出的解析量，是否等于某个由 Galois 上同调定义的代数量。设 $Q_\infty/\Q$ 是 $\Z_p$-扩张，$\Gamma=\Gal(Q_\infty/\Q)$，$\Lambda=\Z_p[[\Gamma]]$。给定 $p$ 进表示 $V_p$ 及格 $T_p$，令 $A_p=V_p/T_p$。在
\[
H^1(\Gal(\overline\Q/Q_\infty),A_p)
\]
中加入局部条件得到 Selmer 群 $\operatorname{Sel}(A_p)$。其 Pontryagin 对偶通常是有限生成 $\Lambda$-模。

主猜想的形状是
\[
\operatorname{char}_\Lambda(\operatorname{Sel}(A_p)^\vee)=(\text{由 }p\text{ 进 }L\text{-函数生成的理想}),
\]
可能还要乘上明确的局部修正因子。这句话把解析函数的零点阶数，与上同调群的大小和秩联系起来。

\subsection{椭圆曲线}

设 $E/\Q$ 是椭圆曲线。其一阶上同调给出相容系统
\[
V_l(E)=H^1_{\mathrm{et}}(E_{\overline\Q},\Q_l),
\]
对应的 $L$-函数有 Euler 积
\[
L(s,E)=\prod_p(1-a_p(E)p^{-s}+\varepsilon(p)p^{1-2s})^{-1}.
\]
模性定理说明存在权 $2$ 的新形式 $f$，使
\[
L(s,E)=L(s,f).
\]
这把椭圆曲线的算术与模形式的解析理论连接起来。

Birch--Swinnerton-Dyer 猜想则把 $L(s,E)$ 在 $s=1$ 的零点阶数与 $E(\Q)$ 的秩相等，并把首项系数表示为 regulator、Tate--Shafarevich 群、Tamagawa 数和挠点阶的组合。这里 $\Sha(E/\Q)$ 是局部有点但整体障碍仍存在的群，与前面的 $\Sha^i$ 思想相同。

\subsection{志村簇}

志村簇提供更高维的几何来源。它们的 zeta 函数和 $L$-函数常通过上同调上的 Hecke 算子与 Frobenius 作用来描述。大方向是：把自动形式放入几何上同调，然后由迹公式或 Lefschetz 公式提取点数与 $L$-函数信息。

对当前学习阶段，需要记住的是结构角色：志村簇是把 Galois 表示、自动形式和几何上同调放在同一个空间中的平台。它不是本章证明 Poitou--Tate 的输入，但说明了上同调方法为什么会持续出现在现代数论中。

\subsection{Tamagawa 数}

Tamagawa 数从测度观点刻画代数群的算术体积。对定义在 $\Q$ 上的代数群 $G$，Tamagawa 数形如
\[
\tau(G)=\operatorname{vol}\bigl(G(\Q)\backslash G(\A)\bigr).
\]
对 $G=\GL_1$，它与 $\Q$ 的 zeta 函数特殊值和类数公式相连。对更一般的半单群，Weil 的 Tamagawa 数猜想预言单连通半单群的 Tamagawa 数为 $1$。

从本章视角看，Tamagawa 数、BSD 公式、Bloch--Kato 猜想等都在表达同一种结构：解析对象的特殊值，应当由某些上同调群、局部条件、对偶配对和 regulator 共同给出。这里的对偶定理是这些更大理论的基础语法。

\section*{习题详解}
\addcontentsline{toc}{section}{习题详解}

\subsection*{习题 1}

设 $n>0$，$p>2$。对 $\sigma\in\Gal(\Q(\mu_{p^n})/\Q)$，存在 $a\in\Z$ 使
\[
\sigma(\zeta)=\zeta^a,\qquad \zeta\in\mu_{p^n}.
\]

第一问要求证明
\[
\Gal(\Q(\mu_{p^n})/\Q)\simeq (\Z/p^n\Z)^\times.
\]
任取本原 $p^n$ 次单位根 $\zeta_{p^n}$。扩张 $\Q(\zeta_{p^n})/\Q$ 由 $\zeta_{p^n}$ 生成，任何 $\Q$-自同构由 $\zeta_{p^n}$ 的像决定。像必须仍是本原 $p^n$ 次单位根，故为 $\zeta_{p^n}^a$，其中 $(a,p)=1$，且 $a$ 模 $p^n$ 唯一。反过来，若 $a\in(\Z/p^n\Z)^\times$，则 $\zeta_{p^n}\mapsto \zeta_{p^n}^a$ 保持 $\zeta_{p^n}$ 的极小多项式，即第 $p^n$ 个分圆多项式，所以给出自同构。于是得到同构。

当 $m\leq n$ 时，$\Q(\mu_{p^m})\subset \Q(\mu_{p^n})$，限制映射对应
\[
(\Z/p^n\Z)^\times\longrightarrow(\Z/p^m\Z)^\times.
\]
取逆极限得到 $p$ 进分圆特征
\[
\chi_p:G_\Q\longrightarrow \varprojlim_n(\Z/p^n\Z)^\times=\Z_p^\times.
\]
它满足
\[
\sigma(\zeta)=\zeta^{\chi_p(\sigma)\bmod p^n},\qquad \zeta\in\mu_{p^n}.
\]
连续性来自核
\[
\Ker(G_\Q\to(\Z/p^n\Z)^\times)=\Gal(\overline\Q/\Q(\mu_{p^n}))
\]
是开子群。

第二问的说法应读作分圆方向的最大 Abel 商：由 Kronecker--Weber 定理，$\Q$ 的任意有限 Abel 扩张都包含在某个分圆域中；若要求只在 $p$ 处分歧，则导子只能是 $p$ 的幂，所以相应扩张都包含在
\[
\Q(\mu_{p^\infty})=\bigcup_n\Q(\mu_{p^n}).
\]
反过来每个 $\Q(\mu_{p^n})$ 只在 $p$ 和无限素位处分歧，所以并起来得到该最大扩张。

第三问分解
\[
(\Z/p^{n+1}\Z)^\times\simeq \mu_{p-1}\times(1+p\Z_p)/(1+p^{n+1}\Z_p).
\]
取 $\Delta=\mu_{p-1}$ 在 $\Q(\mu_{p^\infty})$ 中的固定域 $Q_\infty$，则
\[
\Gal(Q_\infty/\Q)\simeq 1+p\Z_p\simeq \Z_p.
\]
而
\[
\Gal(\Q(\mu_{p^\infty})/Q_\infty)=\Delta.
\]
这就是 $\Q$ 的分圆 $\Z_p$-扩张。

第四问设 $\Gamma=\Gal(Q_\infty/\Q)$。若为每个 $n$ 选取相容的生成元 $\gamma_n\in \Gamma/\Gamma^{p^n}$，则相容性保证存在唯一 $\gamma\in\Gamma$ 映到全部 $\gamma_n$。映射
\[
\Z_p\to\Gamma,\qquad x\mapsto \gamma^x
\]
先在整数上定义，再由 $\Gamma$ 的 $p$ 进完备性连续延拓。模 $p^n$ 后它给出 $\Z/p^n\Z\simeq\Gamma/\Gamma^{p^n}$，故极限上是同构。

第五问令 $X=T-1$。多项式环中 $T=1+X$，所以
\[
\Z_p[T]=\Z_p[X].
\]
群代数
\[
\Z_p[\Gamma/\Gamma^{p^n}]
\]
在选定生成元后同构于
\[
\Z_p[T]/(T^{p^n}-1)
\simeq \Z_p[X]/((1+X)^{p^n}-1).
\]
记
\[
h_n(X)=(1+X)^{p^n}-1.
\]
由于 $h_n(X)\to 0$ 在 $X$-进拓扑中给出越来越高的邻域，逆极限为形式幂级数环：
\[
\varprojlim_n\Z_p[X]/(h_n(X))\simeq \Z_p[[X]].
\]
因此 Iwasawa 代数 $\Lambda=\varprojlim_n\Z_p[\Gamma/\Gamma^{p^n}]$ 同构于 $\Z_p[[X]]$。

第六问若 $\chi$ 是导子 $p^d$ 的 Dirichlet 特征，则它是同态
\[
(\Z/p^d\Z)^\times\to\C^\times.
\]
复合
\[
G_\Q\to\Gal(\Q(\mu_{p^\infty})/\Q)\to
\Gal(\Q(\mu_{p^d})/\Q)\simeq(\Z/p^d\Z)^\times\xrightarrow{\chi}\C^\times
\]
给出 $G_\Q$ 的一维特征。这正是把 Dirichlet 特征看成 Galois 特征的标准方式。

\subsection*{习题 2}

设 $R$ 是特征 $0$ 的完备离散赋值环，剩余域特征为 $p$，$F=\operatorname{Frac}(R)$。要证明存在连续同态
\[
\chi:G_F\to \Z_p^\times
\]
使其在 $\mu_{p^\infty}$ 上的作用为
\[
\sigma(\zeta)=\zeta^{\chi(\sigma)}.
\]
对每个 $n$，$G_F$ 作用在有限群 $\mu_{p^n}$ 上，给出
\[
\chi_n:G_F\to\operatorname{Aut}(\mu_{p^n})\simeq(\Z/p^n\Z)^\times.
\]
映射 $\mu_{p^{n+1}}\to\mu_{p^n}$，$\zeta\mapsto\zeta^p$ 与 Galois 作用交换，所以 $\chi_n$ 相容。取逆极限得到 $\chi$。连续性来自每个 $\chi_n$ 的核是固定有限扩张 $F(\mu_{p^n})/F$ 的开子群。

下面证明 $F(\mu_{p^\infty})/F$ 无穷分歧。设 $e_n=e(F(\mu_{p^n})/F)$。若 $e_n$ 有界，则这些扩张的惯性部分阶数有界。可是 $\mu_{p^n}$ 的阶随 $n$ 增长；在剩余特征 $p$ 的局部域中，加入 $p^n$ 次单位根会产生越来越高的野分歧。

把这一点拆开。先取 $r$ 足够大，使 $F_r=F(\mu_{p^r})$ 中已经含有一阶非平凡的 $p$ 幂单位根。对 $n>r$，Kummer 分圆理论给出自然嵌入
\[
F(\mu_{p^n})/F(\mu_{p^r})
\]
的 Galois 群到
\[
(1+p^r\Z_p)/(1+p^n\Z_p).
\]
在分圆扩张的实际像中，这个商给出越来越大的 $p$-群，阶数随 $n-r$ 增长。另一方面，局部域的最大未分歧扩张由剩余域扩张控制，其有限层 Galois 群由 Frobenius 生成；未分歧扩张不会改变绝对分歧指数。若 $F(\mu_{p^n})/F_r$ 中这些越来越大的 $p$-部分都未分歧，则 $e(F(\mu_{p^n})/F_r)=1$，这会迫使所有新的 $p^n$ 次单位根都已经由剩余域方向产生。可是剩余域特征为 $p$，在可分闭包中没有新的 $p$ 幂单位根；多项式 $X^{p^m}-1=(X-1)^{p^m}$ 在剩余域中只有根 $1$。因此这些新的 $p$ 幂单位根只能来自分歧方向。于是 $e(F(\mu_{p^n})/F_r)$ 随 $n$ 无界，进而 $e_n$ 无界。

令 $I_F$ 为惯性子群。若 $\chi(I_F)$ 有限，则存在 $r$ 使 $\chi(I_F)$ 在商
\[
\Z_p^\times\longrightarrow (\Z/p^n\Z)^\times/(1+p^r\Z_p)
\]
以后的所有更高 $p$ 幂部分上为平凡。等价地，$I_F$ 在 $F(\mu_{p^\infty})/F(\mu_{p^r})$ 上平凡。Galois 理论中，惯性子群在某个代数扩张上作用平凡，正表示该扩张相对于基域的对应有限层均未分歧。因此 $F(\mu_{p^\infty})/F(\mu_{p^r})$ 将由未分歧扩张组成，与上一段得到的无界分歧指数矛盾。因此 $\chi(I_F)$ 是无限群。

\subsection*{习题 3}

令
\[
\mu_{p^n}(F^{\mathrm{sep}})=\{a\in F^{\mathrm{sep}}:a^{p^n}=1\}.
\]
若 $\operatorname{char}F\ne p$，多项式 $X^{p^n}-1$ 的导数是 $p^nX^{p^n-1}$，与多项式无公共根，所以有 $p^n$ 个不同根，故
\[
\mu_{p^n}(F^{\mathrm{sep}})\simeq \Z/p^n\Z.
\]
若 $\operatorname{char}F=p$，则
\[
X^{p^n}-1=(X-1)^{p^n},
\]
在可分闭包中唯一根为 $1$，故 $\mu_{p^n}(F^{\mathrm{sep}})=\{1\}$。

现在假设 $\operatorname{char}F\ne p$。映射
\[
\mu_{p^{n+1}}\to\mu_{p^n},\qquad \zeta\mapsto \zeta^p
\]
组成逆系统，定义
\[
T_p(\mathbb G_m)=\varprojlim_n\mu_{p^n}(F^{\mathrm{sep}}).
\]
选取相容根系 $t=(\zeta_n)_n$，满足 $\zeta_{n+1}^p=\zeta_n$，且 $\zeta_1\ne 1$。对 $a=(a_n)\in\Z_p$，定义
\[
a\cdot t=(\zeta_n^{a_n})_n.
\]
这给出 $\Z_p$ 对 $T_p(\mathbb G_m)$ 的自由传递作用。任一相容根系由 $\zeta_n$ 的某个指数 $a_n$ 给出，相容性保证 $(a_n)$ 来自唯一 $a\in\Z_p$。因此
\[
T_p(\mathbb G_m)\simeq \Z_p
\]
为秩一自由 $\Z_p$-模，记作 $\Z_p(1)$。

由分圆特征定义，$\sigma(\zeta_n)=\zeta_n^{\chi(\sigma)}$，所以
\[
\sigma(t)=\chi(\sigma)t.
\]
对 $r\geq0$ 定义
\[
\Z_p(r)=\Z_p(1)^{\otimes r},\qquad
\Z_p(-r)=\Hom_{\Z_p}(\Z_p(r),\Z_p).
\]
若 $t^r$ 是 $\Z_p(r)$ 的基，则
\[
\sigma(t^r)=\chi(\sigma)^r t^r.
\]

若 $M$ 是 $\Z_p[G_F]$-模，定义 Tate 扭转
\[
M(r)=M\otimes_{\Z_p}\Z_p(r).
\]
其 Galois 作用为
\[
\sigma(x\otimes a)=\sigma x\otimes \chi(\sigma)^r a.
\]
最后，
\[
(M(r))(r')=M\otimes \Z_p(r)\otimes \Z_p(r')\simeq M(r+r'),
\]
且
\[
\Hom_{\Z_p}(M(r),\Z_p)\simeq \Hom_{\Z_p}(M,\Z_p)(-r),
\]
因为对偶会把分圆特征的指数取负。
 \chapter{\texorpdfstring{$p$ 进分析}{p 进分析}}

\section{章节导读}

从这一章开始，代数数论进入 $p$ 进理论。前面几章把数域、局部域、理想类群、上同调和类域论组织起来；现在要面对一个新问题：在 $p$ 进域上做分析时，熟悉的实复分析直觉经常失灵。失灵的根源是非 Archimede 不等式
\[
  |x+y|\leq \max\{|x|,|y|\}.
\]
它使球既开又闭，使三角形变成等腰三角形，也使很多有限维空间的结论不能自动推广到无穷维空间。

本章的主线可以这样看。第一，构造并理解 $\C_p$：它是 $\Q_p$ 的代数闭完备化，承担复分析里 $\C$ 所承担的系数域角色，但它不是局部紧的，也不是球完备的。第二，用滤子替代序列，使完备性、紧性和极限可以在不可数基的拓扑空间中表达。第三，引入球完备性，这是 $p$ 进 Hahn--Banach 定理真正需要的条件。第四，在 Banach 空间、Fréchet 空间和算子空间中建立可用的泛函分析语言。第五，用这些分析工具解释 $p$ 进插值、$p$ 进测度和 $p$ 进 $L$-函数的基本构造。

\begin{notation}
本章固定一个素数 $p$。若无额外说明，赋值域均取非平凡非 Archimede 绝对值。若 $F$ 是赋值域，记
\[
  \mathcal O_F=\{x\in F: |x|\leq 1\},\qquad
  \mathfrak m_F=\{x\in F: |x|<1\}.
\]
赋范向量空间的闭球记作 $B_r(x)=\{y:\|y-x\|\le r\}$。
\end{notation}

\section{\texorpdfstring{$\C_p$}{Cp}}

实复分析从
\[
  \Q \longrightarrow \R \longrightarrow \C
\]
开始：先对 usual 绝对值完备化，再取代数闭包。$p$ 进情形多了一步：
\[
  \Q \longrightarrow \Q_p \longrightarrow \Q_p^{\mathrm{alg}}\longrightarrow \C_p.
\]
这里 $\Q_p^{\mathrm{alg}}$ 是 $\Q_p$ 的代数闭包，但它对延拓的 $p$ 进绝对值不完备；因此要再完备化，得到 $\C_p$。这一点很重要：$\C_p$ 同时是代数封闭的和完备的，但它与 $\C$ 的拓扑性质差别很大。

\begin{proposition}
$\Q$ 对 $|\cdot|_p$ 不是完备域。
\end{proposition}

\begin{proof}
我们给出最基本的构造。先设 $p>2$。取整数 $a$，使得 $a$ 在模 $p$ 意义下是平方，但在 $\Q$ 中不是平方；例如可以先取模 $p$ 的非零平方剩余，再避开有理平方。选择 $x_0\in\Z$ 满足 $x_0^2\equiv a\pmod p$。假设已经有 $x_n$ 满足
\[
  x_n^2\equiv a\pmod {p^{n+1}}.
\]
写 $x_n^2-a=c p^{n+1}$。令 $x_{n+1}=x_n+b p^{n+1}$，则
\[
  x_{n+1}^2-a=x_n^2-a+2b x_n p^{n+1}+b^2p^{2n+2}.
\]
模 $p^{n+2}$ 后，最后一项消失，条件变为
\[
  c+2b x_n\equiv 0\pmod p.
\]
因为 $x_n\equiv x_0\not\equiv 0\pmod p$ 且 $p>2$，可解出 $b\bmod p$。于是得到序列 $(x_n)$，并且
\[
  |x_{n+1}-x_n|_p\le p^{-(n+1)}.
\]
所以 $(x_n)$ 是 Cauchy 序列。若它在 $\Q$ 中收敛到 $x$，由乘法连续性得 $x^2=a$，与 $a$ 的选择矛盾。$p=2$ 时可改用三次方同余，例如构造 $x_n^3\equiv 3\pmod {2^n}$，得到同样矛盾。
\end{proof}

\begin{proposition}
$\Q_p^{\mathrm{alg}}$ 对延拓的 $p$ 进绝对值不是完备域。
\end{proposition}

\begin{proof}
证明思路是构造一个需要无限增大的代数次数才能收敛的 Cauchy 序列。取一列本原单位根 $\zeta_i$，使得 $\Q_p(\zeta_0)\subset \Q_p(\zeta_1)\subset\cdots$，且扩张次数严格增长。定义
\[
  c_n=\sum_{i=0}^n p^i\zeta_i.
\]
因为第 $i$ 项的 $p$ 进绝对值为 $p^{-i}$，所以 $(c_n)$ 是 Cauchy 序列。若它在 $\Q_p^{\mathrm{alg}}$ 中收敛到 $c$，设 $[\Q_p(c):\Q_p]=d$。取 $c_d=\sum_{i=0}^d p^i\zeta_i$。尾项估计给出
\[
  |c-c_d|_p\le p^{-(d+1)}.
\]
对 $\Q_p$-自同构 $\sigma$ 作用仍有
\[
  |\sigma(c)-\sigma(c_d)|_p\le p^{-(d+1)}.
\]
但由于 $\zeta_d$ 所产生的新扩张次数超过 $d$，可以找到 $d+1$ 个自同构 $\sigma_1,\ldots,\sigma_{d+1}$，它们固定前一层而使 $\sigma_j(\zeta_d)$ 两两不同。于是
\[
  |\sigma_i(c_d)-\sigma_j(c_d)|_p=p^{-d}\qquad (i\ne j).
\]
非 Archimede 不等式告诉我们，把每个 $\sigma_j(c_d)$ 改动不超过 $p^{-(d+1)}$ 不会改变这些点之间的距离，因此 $\sigma_1(c),\ldots,\sigma_{d+1}(c)$ 两两不同。这说明 $c$ 至少有 $d+1$ 个共轭元，和 $[\Q_p(c):\Q_p]=d$ 矛盾。
\end{proof}

\begin{lemma}
设 $K$ 为完备赋值域，$f(X)\in K[X]$ 是非常数多项式。若存在序列 $(a_j)\subset K$ 使 $|f(a_j)|\to 0$，则 $(a_j)$ 有子序列收敛到 $K$ 中的一个 $f$ 的根。
\end{lemma}

\begin{proof}
在代数闭包中写
\[
  f(X)=c\prod_{i=1}^n (X-\xi_i).
\]
由 $|f(a_j)|\to0$ 得
\[
  |c|\prod_i |a_j-\xi_i|\to0.
\]
有限个因子中至少有一个因子沿某个子序列趋于 $0$，于是 $a_{j_m}\to \xi_i$。因为 $K$ 完备，而 $a_{j_m}\in K$，极限仍在 $K$ 中，所以 $\xi_i\in K$。
\end{proof}

\begin{lemma}
设完备域 $L$ 含有代数封闭稠密子域 $K$。则 $L$ 代数封闭。
\end{lemma}

\begin{proof}
取首一多项式
\[
  f(X)=X^n+a_{n-1}X^{n-1}+\cdots+a_0\in L[X].
\]
对每个 $m$，用稠密性选择 $a_{i,m}\in K$ 逼近 $a_i$，并令
\[
  f_m(X)=X^n+a_{n-1,m}X^{n-1}+\cdots+a_{0,m}\in K[X].
\]
因为 $K$ 代数封闭，$f_m$ 有根 $\alpha_m\in K$。从 $f_m(\alpha_m)=0$ 可得 $\alpha_m$ 的绝对值有统一上界：若 $|\alpha_m|$ 太大，则最高次项的绝对值严格大于其余项最大值，和求和为零矛盾。于是 $|f(\alpha_m)|=|f(\alpha_m)-f_m(\alpha_m)|\to0$。前一个引理给出 $f$ 在 $L$ 中有根。对次数归纳，$f$ 在 $L$ 中完全分解。
\end{proof}

\begin{proposition}
$\C_p$ 是代数封闭完备域。
\end{proposition}

\begin{proof}
按定义 $\C_p$ 是 $\Q_p^{\mathrm{alg}}$ 的完备化，因此完备，且 $\Q_p^{\mathrm{alg}}$ 在其中稠密并代数封闭。代入上一引理得到 $\C_p$ 代数封闭。
\end{proof}

\section{滤子}

序列只能捕捉第一可数空间中的收敛。一般拓扑向量空间里，零点邻域未必有可数基；这时“所有足够小的邻域同时控制”的语言需要滤子。滤子是把“最终落在某个集合中”抽象出来的工具。

\subsection{滤子基}

\begin{definition}
设 $X$ 为集合。$X$ 上的滤子是一个由 $X$ 的子集组成的集合 $\mathcal F$，满足：
\begin{enumerate}[label=(\arabic*)]
\item 若 $S\in\mathcal F$ 且 $S\subset T\subset X$，则 $T\in\mathcal F$；
\item 若 $S,T\in\mathcal F$，则 $S\cap T\in\mathcal F$；
\item $\varnothing\notin\mathcal F$。
\end{enumerate}
若非空子集族 $\mathcal B$ 满足每个 $B\in\mathcal B$ 非空，且对任意 $B_1,B_2\in\mathcal B$，存在 $B_3\in\mathcal B$ 使 $B_3\subset B_1\cap B_2$，则称 $\mathcal B$ 为滤子基。它生成的滤子是
\[
  \mathcal F_{\mathcal B}=\{S\subset X:\exists B\in\mathcal B,\ B\subset S\}.
\]
\end{definition}

\begin{example}
若 $X$ 是拓扑空间，$x\in X$，则 $x$ 的所有邻域组成滤子，称为邻域滤子。若 $X$ 是无限集，所有余有限子集组成滤子，称为 Fréchet 滤子。序列 $(x_n)$ 也给出滤子基
\[
  \{x_n:n\ge N\}\qquad (N=0,1,2,\ldots),
\]
它记录“从某项以后”的信息。
\end{example}

\subsection{滤子收敛}

\begin{definition}
设 $X$ 是拓扑空间，$\mathcal F$ 是 $X$ 上的滤子。若 $x$ 的每个邻域都属于 $\mathcal F$，即 $\mathcal N(x)\subset\mathcal F$，则称 $\mathcal F$ 收敛到 $x$。若存在包含 $\mathcal F$ 的滤子 $\mathcal G$ 收敛到 $x$，则称 $x$ 为 $\mathcal F$ 的聚集点。
\end{definition}

\begin{proposition}
在 Hausdorff 空间中，滤子至多有一个极限。若 $f:X\to Y$ 在 $x$ 连续，且 $\mathcal F\to x$，则由 $f$ 推出的滤子 $f_*\mathcal F$ 收敛到 $f(x)$。
\end{proposition}

\begin{proof}
若 $\mathcal F$ 同时收敛到 $x$ 和 $y$，且 $x\ne y$，由 Hausdorff 性取不交邻域 $U\ni x$、$V\ni y$。收敛性给出 $U,V\in\mathcal F$，所以 $U\cap V\in\mathcal F$，与空集不属于滤子矛盾。

对连续映射，取 $f(x)$ 的邻域 $W$。连续性给出 $x$ 的邻域 $f^{-1}W$。由 $\mathcal F\to x$ 得 $f^{-1}W\in\mathcal F$，于是 $W$ 属于推出滤子，因此 $f_*\mathcal F\to f(x)$。
\end{proof}

\subsection{完备性}

\begin{definition}
设 $G$ 是加法交换拓扑群。滤子基 $\mathcal B$ 称为 Cauchy 滤子基，如果对 $0$ 的任一邻域 $U$，存在 $B\in\mathcal B$ 使
\[
  B-B\subset U.
\]
若每个 Cauchy 滤子基都收敛，则称 $G$ 完备。
\end{definition}

这个定义把度量空间中的 Cauchy 序列推广为“集合逐渐缩小到无法被邻域区分”。在可度量化拓扑向量空间中，它与常见的 Cauchy 序列完备性相容；在不可度量化空间中，它才是更稳定的定义。

\subsection{极滤子}

\begin{proposition}
设 $\mathcal U$ 是 $X$ 上的滤子。以下条件等价：
\begin{enumerate}[label=(\arabic*)]
\item 若 $\mathcal U\subset\mathcal F$ 且 $\mathcal F$ 是滤子，则 $\mathcal F=\mathcal U$；
\item 对任意 $S\subset X$，有 $S\in\mathcal U$ 或 $X\setminus S\in\mathcal U$。
\end{enumerate}
\end{proposition}

\begin{proof}
若 $\mathcal U$ 极大而某个 $S$ 与补集都不在 $\mathcal U$ 中，则 $\{U\cap S:U\in\mathcal U\}$ 不含空集并生成严格更大的滤子，矛盾。反过来，若满足二择性质，而 $\mathcal F$ 是包含 $\mathcal U$ 的滤子，取 $S\in\mathcal F$。若 $S\notin\mathcal U$，则 $X\setminus S\in\mathcal U\subset\mathcal F$，于是 $\varnothing=S\cap(X\setminus S)\in\mathcal F$，矛盾。因此 $S\in\mathcal U$，所以 $\mathcal F=\mathcal U$。
\end{proof}

\begin{lemma}
设 $X$ 是 Hausdorff 空间。
\begin{enumerate}[label=(\arabic*)]
\item 滤子的极限唯一。
\item $X$ 紧当且仅当每个滤子有聚集点。
\item 紧 Hausdorff 空间上的极滤子有唯一极限。
\end{enumerate}
\end{lemma}

\begin{proof}
第一点已证明。第二点中，若 $X$ 紧而滤子 $\mathcal F$ 没有聚集点，则对每个 $x$ 存在邻域 $U_x$ 与某个 $F_x\in\mathcal F$ 满足 $U_x\cap F_x=\varnothing$。有限子覆盖给出有限交 $\cap F_{x_i}\in\mathcal F$ 且为空，矛盾。反向取开覆盖无有限子覆盖，则闭集补集的有限交性质产生滤子而无聚集点。第三点由紧性给出聚集点，极滤子的聚集点自动成为极限，再由 Hausdorff 性唯一。
\end{proof}

\section{球完备性}

\subsection{赋范向量空间}

\begin{definition}
设 $F$ 为非 Archimede 赋值域。$F$-向量空间 $V$ 上的半范是函数 $q:V\to\R_{\ge0}$，满足
\[
  q(ax)=|a|q(x),\qquad q(x+y)\le \max\{q(x),q(y)\}.
\]
若 $q(x)=0$ 推出 $x=0$，则称为范，记作 $\|x\|$。
\end{definition}

\begin{definition}
赋范向量空间之间的线性映射 $A:X\to Y$ 称为有界线性映射，如果存在 $C>0$ 使
\[
  \|Ax\|\le C\|x\|\qquad(x\in X).
\]
其算子范定义为
\[
  \|A\|=\sup_{x\ne0}\frac{\|Ax\|}{\|x\|}.
\]
\end{definition}

\begin{proposition}
赋范向量空间之间的线性映射 $A:X\to Y$ 连续，当且仅当 $A$ 有界。
\end{proposition}

\begin{proof}
有界推出连续，因为 $\|Ax\|\le C\|x\|$ 直接控制原点处连续性，再由线性推出处处连续。反过来，若 $A$ 在 $0$ 连续，则存在 $\delta>0$ 使 $\|x\|<\delta$ 时 $\|Ax\|<1$。取 $a\in F$ 且 $0<|a|<1$。对 $\|x\|\le1$，选 $N$ 使 $|a|^N<\delta$，则
\[
  \|A(a^Nx)\|<1,\qquad \|Ax\|<|a|^{-N}.
\]
单位球像有界。对任意非零 $x$，把 $x$ 乘以合适的 $a^N$ 放进单位球，就得到 $\|Ax\|\le C\|x\|$。
\end{proof}

\subsection{球完备空间}

\begin{definition}
赋范空间 $V$ 称为球完备的，如果任意按包含关系全序排列的闭球族都有非空交。
\end{definition}

\begin{lemma}
对非 Archimede 赋范空间 $V$，以下表述等价：
\begin{enumerate}[label=(\arabic*)]
\item 每个闭球降链都有非空交；
\item 任意两两相交的闭球族有非空交；
\item 任意全序闭球族有非空交。
\end{enumerate}
\end{lemma}

\begin{proof}
非 Archimede 闭球有一个关键性质：两个闭球若相交，则半径较小者包含在半径较大者中。原因是若 $z\in B_r(x)\cap B_s(y)$ 且 $r\le s$，则对 $u\in B_r(x)$，
\[
  \|u-y\|\le \max\{\|u-x\|,\|x-z\|,\|z-y\|\}\le s.
\]
因此两两相交闭球族自动按包含方向局部全序。用半径选择一个共尾降链，再应用降链条件，得到公共点。其余蕴含按定义立即得到。
\end{proof}

\begin{proposition}
球完备赋范空间是完备赋范空间。
\end{proposition}

\begin{proof}
设 $(x_n)$ 是 Cauchy 序列。令
\[
  r_n=\sup_{m\ge n}\|x_m-x_n\|.
\]
则 $r_n\to0$，且 $B_{r_{n+1}}(x_{n+1})\subset B_{r_n}(x_n)$。球完备性给出 $x\in\cap_nB_{r_n}(x_n)$，于是 $\|x-x_n\|\le r_n\to0$，所以 $x_n\to x$。
\end{proof}

\subsection{Hahn--Banach 定理}

\begin{definition}
赋范空间 $Y$ 称有扩张性质，如果对任意赋范空间 $X$、非零子空间 $M\subset X$ 和有界线性映射 $A:M\to Y$，存在有界线性扩张 $\widetilde A:X\to Y$，满足 $\|\widetilde A\|=\|A\|$。
\end{definition}

\begin{theorem}
赋范 $F$-向量空间 $Y$ 球完备，当且仅当 $Y$ 有扩张性质。
\end{theorem}

\begin{proof}
先设 $Y$ 球完备。给定 $A:M\to Y$ 和 $x_0\notin M$。我们先扩张到 $M+Fx_0$。若扩张取 $\widetilde A(x+\lambda x_0)=A(x)+\lambda y_0$，需要找 $y_0\in Y$ 使
\[
  \|A(x)+\lambda y_0\|\le \|A\|\,\|x+\lambda x_0\|
\]
对所有 $x,\lambda$ 成立。令 $\lambda=-1$ 可见条件等价于
\[
  \|y_0-A(x)\|\le \|A\|\,\|x-x_0\|\qquad(x\in M).
\]
所以 $y_0$ 应落在闭球族
\[
  B_x=B_{\|A\|\|x-x_0\|}(A(x))\qquad(x\in M)
\]
的交中。若两个这样的球 $B_x,B_{x'}$ 半径分别为 $r,r'$ 且 $r\le r'$，则
\[
  \|A(x)-A(x')\|\le \|A\|\,\|x-x'\|
  \le \max\{r,r'\}=r',
\]
所以 $B_x\subset B_{x'}$。球完备性给出公共点 $y_0$，一维扩张成立且保范。

对所有保范扩张对 $(M',A')$ 用包含关系排序。任意全序链的并仍是保范扩张，因此 Zorn 引理给出极大元。若定义域不是整个 $X$，再用一维扩张扩大，矛盾。故得到全空间扩张。

反过来，若 $Y$ 不是球完备，则存在全序闭球族 $\{B_i=B_{r_i}(y_i)\}$ 交为空。构造函数 $\varphi(y)$ 为 $y$ 到某个不含 $y$ 的球心的距离；全序闭球的嵌套性质保证这个值与选择无关。令 $X=Y\oplus F$，并定义
\[
  \|(y,a)\|=
  \begin{cases}
    |a|\,\varphi(a^{-1}y),& a\ne0,\\
    \|y\|,& a=0.
  \end{cases}
\]
非 Archimede 不等式和 $\varphi$ 的定义保证这是一个范。子空间 $M=Y\oplus0$ 上的投影 $(y,0)\mapsto y$ 范为 $1$。若它有保范扩张 $\pi:X\to Y$，令 $z=\pi(0,-1)$，则对每个球心 $y_i$，
\[
  \|y_i-z\|=\|\pi(y_i,1)\|\le \|(y_i,1)\|=\varphi(y_i)\le r_i.
\]
于是 $z\in\cap_iB_i$，与交为空矛盾。
\end{proof}

\begin{theorem}
若 $F/\Q_p$ 是有限扩张，$X$ 是赋范 $F$-向量空间，$M\subset X$ 是子空间，$f:M\to F$ 是有界线性泛函，则存在有界线性泛函 $\widetilde f:X\to F$，满足 $\widetilde f|_M=f$ 且 $\|\widetilde f\|=\|f\|$。
\end{theorem}

\begin{proof}
$F/\Q_p$ 有限时，$F$ 是局部紧离散赋值完备域。局部紧非 Archimede 域作为自身的一维赋范空间是球完备的：全序闭球族半径若趋零，则完备性给出唯一极限；若半径下界非零，则紧性保证有限交性质给出公共点。代入上一个定理，取 $Y=F$。
\end{proof}

\subsection{球完备域}

\begin{proposition}
$\C_p$ 不是球完备域。
\end{proposition}

\begin{proof}
思想是利用 $\C_p$ 可分而构造不可数个互不相交的开球。取值群中的严格递减序列
\[
  r_0>r_1>r_2>\cdots,\qquad r_n\to r>0.
\]
从一个半径 $r_0$ 的球开始，每一步在当前球内选两个互不相交的半径 $r_n$ 的子球。每个二进制序列 $\varepsilon=(\varepsilon_1,\varepsilon_2,\ldots)$ 给出一个闭球降链，其半径趋于正数 $r$。若所有降链交非空，就得到不可数个互不相交的非空开球；而 $\C_p$ 是可分度量空间，任意两两不交的非空开集族至多可数，因为每个开集都含有一个可数稠密集中的点。矛盾。因此某个闭球降链交为空，$\C_p$ 非球完备。
\end{proof}

\begin{remark}
存在代数封闭球完备扩张 $\Omega_p\supset\C_p$，并且可以使 $|\Omega_p^\times|=\R_{>0}$。这说明从 $\Q_p$ 到 $\C_p$ 还不是 $p$ 进分析的终点；在无穷维分析和 $p$ 进微分方程中，球完备扩张经常是自然背景。
\end{remark}

\begin{proposition}
把 $\Q$ 上的绝对值归一化为 $|p|_p=p^{-1}$。则
\[
|\Q_p^\times|_p=p^\Z,\qquad |(\Q_p^{\mathrm{alg}})^\times|_p=p^\Q,\qquad
|\C_p^\times|_p=p^\Q.
\]
若 $\Omega_p$ 是取值群满的代数封闭球完备扩张，则
\[
|\Omega_p^\times|_p=\R_{>0}.
\]
\end{proposition}

\begin{proof}
第一式来自离散赋值。任意 $x\in\Q_p^\times$ 可唯一写成
\[
x=p^nu,\qquad n\in\Z,\quad u\in\Z_p^\times,
\]
而单位 $u$ 的绝对值为 $1$，所以 $|x|_p=p^{-n}$，即值群为 $p^\Z$。

对有限扩张 $K/\Q_p$，设分歧指数为 $e$。若 $\pi_K$ 是 $K$ 的素元，则
\[
|\pi_K|_p=p^{-1/e},
\]
因此 $|K^\times|_p=p^{\frac1e\Z}$。在代数闭包中取遍所有有限扩张，分歧指数 $e$ 可以任意大，例如加入适当的 Eisenstein 多项式根即可得到给定分歧指数的有限扩张。于是所有有理指数都出现，得到
\[
|(\Q_p^{\mathrm{alg}})^\times|_p=p^\Q.
\]

从 $\Q_p^{\mathrm{alg}}$ 完备化到 $\C_p$ 不改变值群。理由是：若 $x\in\C_p^\times$，取 $a\in\Q_p^{\mathrm{alg}}$ 使 $|x-a|_p<|x|_p$。非 Archimede 不等式给出
\[
|a|_p=|x|_p,
\]
所以 $x$ 的值已经由代数闭包中的元素实现。故 $|\C_p^\times|_p=p^\Q$。

最后一式是 $\Omega_p$ 的构造性质。为了得到球完备且值群为全体正实数的扩张，需要在完备化和极大即时扩张的过程中同时加入给定实值的元素。直观上，$\C_p$ 的值群 $p^\Q$ 在 $\R_{>0}$ 中稠密但不完备；$\Omega_p$ 补上的不仅是代数根和 Cauchy 极限，还包括所有半径所对应的值。由此可取 $\Omega_p$ 使每个 $r\in\R_{>0}$ 都等于某个非零元素的绝对值。
\end{proof}

\section{Banach 空间}

\subsection{Banach--Steinhaus 定理}

\begin{definition}
赋范 $F$-向量空间 $V$ 若对范诱导的度量完备，则称为 Banach 空间。
\end{definition}

\begin{proposition}
若 $X$ 是赋范空间，$Y$ 是 Banach 空间，则所有有界线性映射 $X\to Y$ 组成的空间 $L(X,Y)$，配以算子范，是 Banach 空间。
\end{proposition}

\begin{proof}
设 $(A_n)$ 是 $L(X,Y)$ 中的 Cauchy 序列。对每个 $x\in X$，
\[
  \|A_nx-A_mx\|\le \|A_n-A_m\|\|x\|,
\]
所以 $(A_nx)$ 在 $Y$ 中收敛。定义 $Ax=\lim_n A_nx$。线性由极限保持线性得到。再由 Cauchy 性选 $N$ 使 $n\ge N$ 时 $\|A_n-A_N\|\le1$，则
\[
  \|A_nx\|\le \max\{\|A_N\|,1\}\|x\|,
\]
取极限得 $A$ 有界。最后 $\|A_n-A\|\to0$，所以完备。
\end{proof}

\begin{theorem}
设 $X$ 是 Banach 空间，$Y$ 是赋范空间，$\{A_i:X\to Y\}_{i\in I}$ 是有界线性映射族。若对每个 $x\in X$，集合 $\{\|A_ix\|:i\in I\}$ 有界，则 $\{\|A_i\|:i\in I\}$ 有界。
\end{theorem}

\begin{proof}
如果算子范无界，可取可数子列 $(A_n)$ 使 $\|A_n\|\to\infty$。先证明存在闭球 $B_r(x_0)$ 与常数 $C,N$，使 $n>N$ 时 $A_n$ 在这个闭球上一致有界。否则可递归构造闭球降链 $B_k$，半径小于 $1/k$，并选 $n_k$ 递增，使得在 $B_k$ 上 $\|A_{n_k}x\|>k$。完备性给出一点 $x\in\cap_kB_k$，从而 $\|A_{n_k}x\|>k$，违反逐点有界。

于是存在 $B_r(x_0)$ 上的一致界 $C$。对 $x$ 足够小，$x+x_0\in B_r(x_0)$，故
\[
  \|A_nx\|\le \max\{\|A_n(x+x_0)\|,\|A_nx_0\|\}
\]
右端对 $n$ 有界。用标量把任意 $x$ 缩放到小球内，即得 $\|A_nx\|\le C'\|x\|$，其中 $C'$ 与 $n$ 无关。这与 $\|A_n\|\to\infty$ 矛盾。
\end{proof}

\subsection{Stone--Weierstrass 定理}

\begin{lemma}
若 $X$ 是紧 Hausdorff 空间，$F$ 是局部域，则连续函数空间 $C(X,F)$ 配以
\[
  \|f\|=\sup_{x\in X}|f(x)|
\]
是 Banach 代数。
\end{lemma}

\begin{proof}
一致极限保持连续性，紧性保证上确界有限。若 $(f_n)$ 是 Cauchy 序列，则对每个 $x$，$f_n(x)$ 在完备域 $F$ 中收敛，定义 $f(x)=\lim f_n(x)$。一致 Cauchy 性给出一致收敛，从而 $f$ 连续。乘法范数满足 $\|fg\|\le\|f\|\|g\|$。
\end{proof}

\begin{lemma}
设 $F$ 是局部域，$a_0\in F^\times$，$C\subset F$ 紧。存在多项式 $P\in F[T]$，使
\[
  P(0)=0,\qquad P(a_0)=1,\qquad |P(c)|\le1\quad(c\in C).
\]
\end{lemma}

\begin{proof}
只需处理 $|c|>|a_0|$ 的紧部分。用有限个小球覆盖它，并构造乘积
\[
  Q(T)=(1-a_0^{-1}T)\prod_j(1-c_j^{-1}T)^{n_j}.
\]
整数 $n_j$ 选得足够大，使得在覆盖中每个球上 $|Q(T)|\le1$，同时 $Q(0)=1$、$Q(a_0)=0$。令 $P=1-Q$，便得所需性质。非 Archimede 估值保证乘积估计中最大项控制整个和。
\end{proof}

\begin{theorem}
设 $X$ 是紧 Hausdorff 空间，$F$ 是局部域。若子代数 $A\subset C(X,F)$ 含常值函数并分离点，则 $A$ 在 $C(X,F)$ 中稠密。
\end{theorem}

\begin{proof}
第一步，给定 $s\ne t$，分离点给出 $f\in A$ 使 $f(s)\ne f(t)$。经仿射变换可令 $f(s)=1,f(t)=0$。上一引理施于紧集 $f(X)$，把 $f$ 代入多项式，可得到范数不超过 $1$ 且在 $s,t$ 上取指定值的函数。

第二步，若 $E$ 是闭集且 $s\notin E$，利用紧性从每个 $y\in E$ 的分离函数中取有限乘积，得到 $g\in A$，满足 $g(s)=1$，而在 $E$ 上任意小。

第三步，若 $E$ 同时开闭，则特征函数 $\chi_E$ 属于 $A$ 的闭包。对 $X\setminus E$ 中每点构造在 $E$ 上接近 $1$、在该点附近接近 $0$ 的函数，再取有限乘积。

最后取任意 $f\in C(X,F)$。紧集 $f(X)$ 可由有限个小球覆盖，并可分割为有限个开闭子集 $C_i$，令 $E_i=f^{-1}(C_i)$。函数
\[
  k=\sum_i f(x_i)\chi_{E_i}
\]
在闭包 $\overline A$ 中，并且可使 $\|f-k\|$ 任意小。因此 $f\in\overline A$。
\end{proof}

\subsection{Banach 代数}

\begin{definition}
若 $A$ 是完备赋范 $F$-代数，且 $\|ab\|\le\|a\|\|b\|$，则称 $A$ 为 Banach 代数。若 $M$ 是 Banach 空间且为 $A$-模，并满足 $\|am\|\le\|a\|\|m\|$，则称 $M$ 为 Banach $A$-模。
\end{definition}

\begin{proposition}
设交换含幺 Banach $F$-代数 $A$ 是 Noether 环，$M$ 是有限生成 Banach $A$-模。则 $M$ 的任意 $A$-子模 $N$ 是闭集。
\end{proposition}

\begin{proof}
令 $\overline N$ 为闭包。它仍是 $A$-子模。Noether 性给出 $\overline N$ 的有限生成元 $e_1,\ldots,e_m$。取 $f_i\in N$ 充分接近 $e_i$。对任意 $x\in\overline N$，写 $x=\sum a_i^{(0)}e_i$，再把 $e_i=f_i+(e_i-f_i)$ 代入，得到
\[
  x=\sum a_i^{(0)}f_i+x_1,\qquad \|x_1\|\le c\|x\|
\]
其中 $0<c<1$。对 $x_1$ 重复此过程，得到收敛级数
\[
  x=\sum_i\Bigl(\sum_{n\ge0}a_i^{(n)}\Bigr)f_i.
\]
因此 $x\in Af_1+\cdots+Af_m\subset N$。所以 $\overline N=N$。
\end{proof}

\begin{definition}
若 $A$ 是 Tate $F$-代数，$X=\operatorname{Sp}(A)$ 为极大理想空间，对 $f\in A$ 定义
\[
  \|f\|_{\mathrm{sp}}=\sup_{x\in X}|f(x)|.
\]
这是谱半范。它衡量 $f$ 在全部刚性点上的大小，是后续解析几何中最自然的范数之一。
\end{definition}

\section{Fréchet 空间}

\subsection{局部凸拓扑空间}

\begin{definition}
设 $V$ 是 $F$-向量空间。$\mathcal O_F$-子模 $L\subset V$ 称为格，如果对每个 $v\in V$，存在 $a\in F^\times$ 使 $av\in L$。
\end{definition}

在非 Archimede 分析中，格替代了复局部凸空间中的“凸、均衡、吸收”邻域。由格组成的零点邻域基给出局部凸拓扑。若格族 $\{L_j\}$ 满足有限交细化和标量缩放细化条件，则对应拓扑使 $V$ 成为拓扑向量空间。

\begin{proposition}
格与半范给出同一类局部凸拓扑。具体地，若 $L$ 是格，定义
\[
  p_L(v)=\inf\{|a|:v\in aL\},
\]
则 $p_L$ 是半范。反过来，有限多个半范的不等式
\[
  p_i(v)\le\varepsilon
\]
定义格，并由这些格生成原拓扑。
\end{proposition}

\begin{proof}
$p_L(av)=|a|p_L(v)$ 来自 $v\in bL$ 当且仅当 $av\in abL$。若 $v\in aL$、$w\in bL$，并设 $|a|\le |b|$，则 $v+w\in bL$，所以 $p_L(v+w)\le\max\{p_L(v),p_L(w)\}$。反向由半范球构成的集合对加法和 $\mathcal O_F$-数乘封闭，并吸收每个向量，因此是格。两种构造互相给出相同的邻域基。
\end{proof}

\begin{definition}
可度量化完备局部凸空间称为 Fréchet 空间。等价地，它的拓扑由可数半范族给出，并且每个 Cauchy 滤子收敛。
\end{definition}

\subsection{桶形空间}

\begin{definition}
非 Archimede 局部凸空间 $V$ 称为桶形空间，如果每个闭格都是开格。
\end{definition}

\begin{proposition}
Fréchet 空间是桶形空间。
\end{proposition}

\begin{proof}
设 $L$ 是 Fréchet 空间 $V$ 中的闭格。取 $a\in F$ 且 $|a|>1$。因为 $L$ 吸收每个向量，
\[
  V=\bigcup_{n\ge0}a^nL.
\]
Baire 定理说明某个闭集 $a^nL$ 有内点，因此 $L$ 有内点。若 $v+M\subset L$，其中 $M$ 是开格，则 $M\subset L-v=L$，所以 $L$ 是 $0$ 的邻域，故为开格。
\end{proof}

\begin{proposition}
若 $V$ 是 Fréchet 空间，$W\subset V$ 是闭向量子空间，则 $V/W$ 是 Fréchet 空间。
\end{proposition}

\begin{proof}
闭性保证商空间 Hausdorff。若 $q_n$ 是 $V$ 上定义拓扑的递增半范族，在商空间上定义
\[
  \overline q_n(v+W)=\inf_{w\in W}q_n(v+w).
\]
这些半范给出商拓扑并使商空间可度量。完备性来自 Cauchy 滤子提升：商中的 Cauchy 滤子可选代表在 $V$ 中逐步校正成 Cauchy 滤子，其极限模掉 $W$ 后给出商中的极限。
\end{proof}

\begin{theorem}
设 $V$ 是桶形空间，$W$ 是 Fréchet 空间。若线性映射 $f:V\to W$ 的图像
\[
  \Gamma(f)=\{(v,f(v)):v\in V\}
\]
在 $V\times W$ 中闭，则 $f$ 连续。
\end{theorem}

\begin{proof}
要证明连续，只需证明每个开格 $M\subset W$ 的逆像是开集。因为 $V$ 桶形，足以证明 $f^{-1}(M)$ 是闭格。取 $v$ 在 $f^{-1}(M)$ 的闭包中。选 $W$ 中递减开格基
\[
  M=M_1\supset M_2\supset\cdots.
\]
由 $v$ 属于闭包，可递归选 $v_n\in f^{-1}(M_n)$ 使部分和逼近 $v$，并使 $f$ 的部分和成为 $M$ 中的 Cauchy 序列。Fréchet 完备性给出极限 $w\in M$。闭图条件推出 $(v,w)\in\Gamma(f)$，所以 $f(v)=w\in M$。因此 $f^{-1}(M)$ 闭，从而开。
\end{proof}

\begin{theorem}
设 $V$ 是 Fréchet 空间，$W$ 是 Hausdorff 桶形空间。若 $f:V\to W$ 是连续线性满射，则 $f$ 是开映射。
\end{theorem}

\begin{proof}
核 $\ker f$ 是闭子空间，所以 $V/\ker f$ 是 Fréchet 空间。商映射是开映射，因此可把问题化为 $f$ 为连续线性双射。只需证明逆映射连续。逆映射的图像与 $f$ 的图像在 $V\times W$ 和 $W\times V$ 的交换同构下对应；由于 $f$ 连续且 $W$ Hausdorff，图像闭。闭图定理给出 $f^{-1}$ 连续。
\end{proof}

\subsection{Fréchet--Stein 代数}

\begin{definition}
设 Fréchet 代数 $A$ 的拓扑由递增半范 $q_1\le q_2\le\cdots$ 定义。令 $A_n$ 为 $A/\ker q_n$ 对 $q_n$ 的 Banach 完备化。若
\[
  q_n(ab)\le c_nq_n(a)q_n(b)
\]
使每个 $A_n$ 成为 Banach 代数，并且 $A\simeq\varprojlim A_n$，则称这是 $A$ 的 Banach 代数逼近。若进一步每个 $A_n$ Noether，且转移映射 $A_{n+1}\to A_n$ 平坦，则称 $A$ 为 Fréchet--Stein 代数。
\end{definition}

这个定义的动机是把难处理的 Fréchet 代数化为一列 Banach 代数。Noether 性保证有限生成模可控，平坦性保证从第 $n+1$ 层限制到第 $n$ 层时不制造扭曲。后面研究 $p$ 进李群的分布代数时，这一结构是核心技术背景。

\section{算子空间}

给定局部凸空间 $V,W$，连续线性映射空间记为 $L(V,W)$。为了在这个空间上做分析，需要说明“算子收敛”是什么意思。

\begin{definition}
子集 $B\subset V$ 称有界，如果对 $V$ 的任意开格 $L$，存在 $a\in F$ 使 $B\subset aL$。给定 $W$ 上半范 $p$，定义
\[
  p_B(T)=\sup_{v\in B}p(Tv).
\]
由所有单点集给出的拓扑称弱拓扑或点收敛拓扑，记 $L_s(V,W)$。由所有有界集给出的拓扑称强拓扑或有界收敛拓扑，记 $L_b(V,W)$。
\end{definition}

\begin{definition}
$V' = L(V,F)$ 称为对偶空间。带弱拓扑者记 $V_s'$，带强拓扑者记 $V_b'$。若自然映射
\[
  \delta:V\to (V_b')_b',\qquad v\mapsto(\ell\mapsto \ell(v))
\]
是拓扑同构，则称 $V$ 自反。
\end{definition}

\begin{definition}
局部凸空间 $V$ 的 $\mathcal O_F$-子模 $A$ 称为 $c$ 紧，如果对任意递减有向开格族 $(L_i)$，自然映射
\[
  A\to\varprojlim_i A/(A\cap L_i)
\]
满射。连续线性映射 $T:V\to W$ 称紧算子，如果存在开格 $L\subset V$，使 $T(L)$ 的闭包是有界 $c$ 紧集。若对任意 Banach 空间 $W$ 都有 $L(V,W)=C(V,W)$，则称 $V$ 为核空间。
\end{definition}

这些定义看起来抽象，但目的很明确：在非局部紧的 $p$ 进世界中，普通紧集太少，$c$ 紧集提供可用替代；核空间则保证从 $V$ 出发的算子自动具有紧性性质。

\section{\texorpdfstring{$p$ 进插值}{p 进插值}}

插值问题是：已知函数只在稠密子集上有值，能否把它延拓到整个 $p$ 进空间，并保持连续或解析性质？这正是 $p$ 进 $L$-函数的核心思想：把复 $L$-函数在整数点的特殊值，组织成 $p$ 进变量的函数。

\begin{proposition}
设 $S\subset\Z_p$ 稠密。函数 $f:S\to\Q_p$ 可唯一延拓为连续函数 $\widetilde f:\Z_p\to\Q_p$，当且仅当 $f$ 有界且一致连续。
\end{proposition}

\begin{proof}
若连续延拓存在，由 $\Z_p$ 紧得 $\widetilde f$ 有界且一致连续，限制到 $S$ 后仍如此。

反过来，对 $x\in\Z_p$ 取序列 $x_n\in S$ 且 $x_n\to x$。一致连续性使 $(f(x_n))$ 为 Cauchy 序列，$\Q_p$ 完备，所以可定义
\[
  \widetilde f(x)=\lim_n f(x_n).
\]
若换另一逼近序列 $y_n$，则交错序列 $x_1,y_1,x_2,y_2,\ldots$ 仍趋于 $x$，一致连续性保证两个极限相同。连续性也由一致连续性直接得到。唯一性来自 $S$ 的稠密性和 Hausdorff 性。
\end{proof}

\begin{lemma}
$\Z$ 在 $\Z_p$ 中稠密。若 $0\le s_0\le p-2$，则
\[
  S_{s_0}=\{n\in\Z_{\ge0}:n\equiv s_0\pmod {p-1}\}
\]
在 $\Z_p$ 中稠密。
\end{lemma}

\begin{proof}
任意 $x\in\Z_p$ 有展开 $x=\sum_{n\ge0}a_np^n$，部分和 $x_N=\sum_{n<N}a_np^n$ 是整数且 $x_N\to x$。为满足模 $p-1$ 条件，令
\[
  y_N=x_N+(s_0-x_N)p^N.
\]
因为 $p^N\equiv1\pmod {p-1}$，所以 $y_N\equiv s_0\pmod {p-1}$，且 $y_N-x_N$ 被 $p^N$ 整除，所以 $y_N\to x$。
\end{proof}

\begin{proposition}
若 $p\nmid n$，固定 $s_0$，函数
\[
  S_{s_0}\to\Z_p,\qquad s\mapsto n^s
\]
有唯一连续延拓到 $\Z_p$。
\end{proposition}

\begin{proof}
由 Fermat 小定理，$n^{p-1}=1+mp$。若 $s=s_0+(p-1)u$，则
\[
  n^s=n^{s_0}(1+mp)^u.
\]
若 $u\equiv u'\pmod {p^N}$，二项式展开给出
\[
  (1+mp)^u-(1+mp)^{u'}\in p^{N+1}\Z_p.
\]
因此 $s\mapsto n^s$ 一致连续且有界。由插值命题延拓。
\end{proof}

\begin{definition}
Bernoulli 多项式由生成函数定义：
\[
  \frac{te^{xt}}{e^t-1}=\sum_{k\ge0}B_k(x)\frac{t^k}{k!},\qquad B_k=B_k(0).
\]
\end{definition}

经典公式
\[
  \zeta(1-2n)=-\frac{B_{2n}}{2n}
\]
把 Riemann $\zeta$ 函数的负奇整数特殊值与 Bernoulli 数相连。Kummer 同余说明，去掉 Euler 因子后的数值
\[
  (1-p^{m-1})\frac{B_m}{m}
\]
随 $m$ 按 $p$ 进拓扑连续变化。这就是 $p$ 进插值的原型。

这句话需要按两层来理解。复分析的一层是：整数点 $m$ 上的数值来自同一个复函数 $\zeta(1-m)$。$p$ 进分析的一层是：虽然复绝对值下整数 $m$ 没有趋向某个 $p$ 进数的意义，但在 $\Z_p$ 中，满足同一模 $p-1$ 条件的整数可以趋向任意 $p$ 进点。Kummer 同余正是连续性的检验：若
\[
m\equiv m'\pmod{(p-1)p^N},
\]
则修正后的 Bernoulli 数在模 $p^N$ 意义下接近。于是同一批特殊值同时有两种延拓：复变量的延拓给出经典 $\zeta$ 函数，$p$ 进变量的延拓给出 $p$ 进 $\zeta$ 函数或更一般的 $p$ 进 $L$ 函数。

Euler 因子 $1-p^{m-1}$ 不是装饰。没有去掉这个因子时，$p$ 处的局部项会破坏 $p$ 进连续性；去掉以后，只保留与 $p$ 互素的除子贡献，这些贡献可以用 $d^s$ 在 $\Z_p$ 上连续插值。后面的测度构造就是把这个现象系统化：先把特殊值写成某个测度的矩，再对 $s$ 做 Mellin 变换。

\begin{lemma}
固定 $p>2$。若 $k_i\to k$ 于 $\Z_p$，且 $k_i$ 为足够大的偶整数并落在同一模 $p-1$ 剩余类中，则 Eisenstein 级数的系数
\[
  \frac{\zeta^{(p)}(1-k_i)}2,\qquad \sigma_{k_i-1}^{(p)}(n)
\]
在 $\Q_p$ 中收敛。
\end{lemma}

\begin{proof}
常数项的收敛由 Kummer 同余给出。对固定 $n$，$\sigma_{k-1}^{(p)}(n)$ 只对 $p\nmid d$ 的有限多个除子 $d$ 求和；每个 $d^{k_i-1}$ 按上面的指数插值命题随 $k_i$ 收敛。因此有限和收敛。
\end{proof}

\section{\texorpdfstring{$p$ 进测度}{p 进测度}}

\subsection{分布}

\begin{definition}
设有限集构成的逆系统
\[
  X=\varprojlim X_n
\]
给出射影有限集。设 $K$ 为完备赋值域，$V$ 为 $K$-Banach 空间。若映射族 $\varphi_n:X_n\to V$ 满足相容关系
\[
  \varphi_n(x)=\sum_{\pi_{n+1}(y)=x}\varphi_{n+1}(y),
\]
则称其定义了 $X$ 上的分布。若
\[
  |\varphi|=\sup_{n,x}\|\varphi_n(x)\|<\infty,
\]
则称为测度。
\end{definition}

对阶梯函数 $f$，若它在某层 $X_m$ 上已经确定，定义积分
\[
  \int_X f\,d\varphi=\sum_{x\in X_m}f(x)\varphi_m(x).
\]
相容关系保证这个定义与所选层无关。

\begin{proposition}
若 $\varphi$ 是测度，则阶梯函数上的积分唯一延拓为 $C(X,K)$ 上的连续线性泛函。
\end{proposition}

\begin{proof}
对阶梯函数 $f$，
\[
  \left\|\int f\,d\varphi\right\|
  \le \max_{x\in X_m}|f(x)|\,|\varphi|
  =\|f\|_\infty|\varphi|.
\]
因此积分是有界线性泛函。射影有限集上的连续函数可由阶梯函数一致逼近，于是由完备性唯一延拓。
\end{proof}

\subsection{Bernoulli 分布}

\begin{proposition}
Bernoulli 多项式满足分布关系
\[
  B_k(X)=N^{k-1}\sum_{a=0}^{N-1}B_k\!\left(\frac{X+a}{N}\right).
\]
\end{proposition}

\begin{proof}
从生成函数出发：
\[
  \sum_{a=0}^{N-1}\frac{t e^{(X+a)t/N}}{e^t-1}
  =\frac{t e^{Xt/N}}{e^{t/N}-1}.
\]
把右端按 $t$ 展开，比较 $t^k/k!$ 的系数，就得到公式。
\end{proof}

令
\[
  \varphi_N(x)=N^{k-1}B_k(\langle x\rangle),\qquad x\in \frac1N\Z/\Z,
\]
其中 $\langle x\rangle$ 是 $[0,1)$ 中的代表。上面的分布关系说明 $\{\varphi_N\}_N$ 是 $\varprojlim_N\frac1N\Z/\Z$ 上的分布，称为 Bernoulli 分布。

\subsection{\texorpdfstring{$E_{k,c}$}{E k c}}

通过同构 $\frac1N\Z/\Z\simeq \Z/N\Z$，定义
\[
  E_k^N(x)=\frac1kN^{k-1}B_k\!\left(\left\langle\frac{x}{N}\right\rangle\right).
\]
若 $c$ 与 $N$ 互素，设
\[
  E_{k,c}^N(x)=E_k^N(x)-c^kE_k^N(c^{-1}x).
\]

\begin{proposition}
若 $c\in\Z_p^\times$，则在 $\Z_p$ 上有分布恒等式
\[
  E_{k,c}(x)=x^{k-1}E_{1,c}(x),
\]
并且当 $c^k\ne1$ 时
\[
  \frac{B_k}{k}=\frac{1}{1-c^k}\int_{\Z_p}x^{k-1}\,dE_{1,c}(x).
\]
\end{proposition}

\begin{proof}
有限层上由 Bernoulli 多项式的同余性质得
\[
  E_{k,c}^N(x)\equiv x^{k-1}E_{1,c}^N(x)\pmod {N\Z_p}
\]
当 $p\mid N$。令 $N=p^n$ 并取极限，得到第一式。第二式来自定义：
\[
  \int dE_{k,c}=\int dE_k-c^k\int d(E_k\circ c^{-1})=(1-c^k)\int dE_k.
\]
而 $\int dE_k=B_k/k$，再用第一式替换 $dE_{k,c}$。
\end{proof}

\subsection{\texorpdfstring{$p$ 进 $L$-函数}{p 进 L-函数}}

\begin{definition}
设 $\omega:\Z_p^\times\to\mu_{p-1}$ 为 Teichmüller 特征标。每个 $a\in\Z_p^\times$ 可唯一写成
\[
  a=\omega(a)\langle a\rangle,\qquad \langle a\rangle\in 1+p\Z_p
\]
当 $p=2$ 时把 $1+p\Z_p$ 换成 $1+4\Z_2$。若 $\mu$ 是 $\Z_p$ 上的测度，定义
\[
  \Gamma_p\mu(s)=\int_{\Z_p^\times}\langle a\rangle^s\,d\mu(a),\qquad
  M_p\mu(s)=\int_{\Z_p^\times}\langle a\rangle^s a^{-1}\,d\mu(a).
\]
\end{definition}

\begin{proposition}
若 $\mu$ 是 $\Z_p$ 上的测度，则 $\Gamma_p\mu(s)$ 是 $s\in\Z_p$ 的解析函数。
\end{proposition}

\begin{proof}
在 $1+p\Z_p$ 上可写
\[
  \langle a\rangle^s=(1+(\langle a\rangle-1))^s
  =\sum_{n\ge0}\binom{s}{n}(\langle a\rangle-1)^n.
\]
因为 $\langle a\rangle-1$ 被 $p$ 整除，且 $(\langle a\rangle-1)^n/n!\to0$，级数在 $p$ 进范数下一致收敛。测度有界，所以可以逐项积分，得到
\[
  \Gamma_p\mu(s)=\sum_{n\ge0}c_n\binom{s}{n},\qquad c_n\to0.
\]
Mahler 多项式 $\binom{s}{n}$ 给出 $\Z_p$ 上的解析展开。
\end{proof}

\begin{definition}
设 $\chi$ 是 Dirichlet 特征标。选 $c\in\Z_p^\times$，使 $s\mapsto \chi(c)\langle c\rangle^s$ 不恒为 $1$。定义
\[
  L_p(1-s,\chi)=
  \frac{-1}{1-\chi(c)\langle c\rangle^s}
  \int_{\Z_p^\times}\langle a\rangle^s a^{-1}\chi(a)\,dE_{1,c}(a).
\]
\end{definition}

定义中的分母用于消去辅助整数 $c$ 带来的修正。测度 $E_{1,c}$ 本身依赖 $c$，但它的矩满足
\[
\int_{\Z_p^\times} a^{k-1}\psi(a)\,dE_{1,c}(a)
  =(1-\psi(c)c^k)\frac{B_{k,\psi}}{k}.
\]
因此只要除以对应的因子，最终得到的函数就只依赖 $\chi$，不依赖 $c$。这正是构造 $p$ 进 $L$ 函数时常见的套路：为了得到有界测度先引入辅助修正，再用解析因子把修正除回去。

\begin{theorem}
对正整数 $k$，
\[
  L_p(1-k,\chi)=-\frac{1}{k}B_{k,\chi\omega^{-k}}.
\]
该函数与 $c$ 的选择无关；若 $\chi$ 非平凡且可取 $c$ 使 $\chi(c)\ne1$，则 $L_p(1-s,\chi)$ 是 $s$ 的解析函数。
\end{theorem}

\begin{proof}
把 $s=k$ 代入定义，利用 $a=\omega(a)\langle a\rangle$，得到
\[
  \langle a\rangle^k a^{-1}\chi(a)
  =a^{k-1}\omega(a)^{-k}\chi(a).
\]
于是前一小节关于 $E_{1,c}$ 的积分公式给出
\[
  \int_{\Z_p^\times}a^{k-1}\chi(a)\omega(a)^{-k}\,dE_{1,c}(a)
  =(1-\chi(c)\omega(c)^{-k}c^k)\frac{B_{k,\chi\omega^{-k}}}{k}.
\]
因为 $c^k\omega(c)^{-k}=\langle c\rangle^k$，代回定义即得特殊值公式。满足固定同余类的正整数在 $\Z_p$ 中稠密，特殊值决定连续函数，所以不同 $c$ 给出同一函数。解析性来自 Mellin 变换的解析性和分母在所选条件下不恒为零。

特殊值公式中的扭转 $\omega^{-k}$ 来自分解
\[
a=\omega(a)\langle a\rangle.
\]
$p$ 进变量只能安全地放在 $\langle a\rangle\in1+p\Z_p$ 上，因为这里可以用二项式级数定义 $\langle a\rangle^s$。有限阶部分 $\omega(a)$ 不能随 $s$ 连续变化，只能在整数 $k$ 处以 $\omega(a)^{-k}$ 的形式出现。因此 $p$ 进 $L$ 函数插值的不是单纯的 $B_{k,\chi}$，而是 $B_{k,\chi\omega^{-k}}$。这也是为什么 $p$ 进特殊值公式看起来比复特殊值公式多了一个 Teichmüller 扭转。
\end{proof}

\section{习题详解}

\begin{exercise}
设 $F$ 是剩余特征 $p$ 的赋值域，且 $|p|=p^{-1}$。证明幂函数 $(1+y)^\alpha$ 对 $|y|<1$ 和 $\alpha\in\Z_p$ 有良好定义，并满足二项式展开。
\end{exercise}

\begin{proof}
先证估计。若 $|y-1|\le\varepsilon<1$，则
\[
  y^p-1=(y-1)(1+y+\cdots+y^{p-1}).
\]
写 $y=1+u$。和式等于 $p+\binom p2u+\cdots+u^{p-1}$，其绝对值不超过 $\max\{|p|,|u|\}$，故
\[
  |y^p-1|\le\max\{p^{-1},\varepsilon\}|y-1|.
\]
于是若 $|a-1|<1$，反复取 $p$ 次幂得到 $a^{p^n}\to1$。反向若 $a^{p^n}\to1$，则 $a^{p^n}$ 最终在 $1+\mathfrak m_F$ 中，而 $x\mapsto x^{p^n}$ 在剩余域中保持 $1$ 的邻域结构，推出 $a$ 的剩余类为 $1$，即 $|a-1|<1$。

若 $a=1+y$ 且 $|y|<1$，对整数 $k_j\to\alpha$ 定义 $a^\alpha=\lim_j a^{k_j}$。上一估计说明 $k\mapsto a^k$ 在 $\Z_{\ge0}$ 上一致连续，因此极限与逼近序列无关。最后，对整数 $n$ 有二项式公式
\[
  (1+y)^n=\sum_{m\ge0}\binom nm y^m.
\]
当 $n\to\alpha$ 于 $\Z_p$ 时，$\binom nm\to\binom\alpha m$，且 $|y|^m\to0$，可逐项取极限，得到
\[
  (1+y)^\alpha=\sum_{m\ge0}\binom\alpha m y^m.
\]
\end{proof}

\begin{exercise}
设 $a_n\in\mathcal O_{\C_p}$ 且 $|a_n|\to0$，令
\[
  \varphi(x)=\sum_{n\ge0}a_n\binom{x}{n},\qquad
  \Delta\varphi(x)=\varphi(x+1)-\varphi(x).
\]
证明 $\Delta\binom{x}{n}=\binom{x}{n-1}$，$\Delta^k\varphi(x)=\sum_{n\ge0}a_{n+k}\binom{x}{n}$，并且 $\Delta^k\varphi(0)=a_k$。
\end{exercise}

\begin{proof}
Pascal 恒等式给出
\[
  \binom{x+1}{n}=\binom{x}{n}+\binom{x}{n-1},
\]
所以第一式成立。由于 $|a_n|\to0$，级数一致收敛，差分可逐项进行。一次差分把系数向左平移一位，迭代 $k$ 次得到
\[
  \Delta^k\varphi(x)=\sum_{n\ge0}a_{n+k}\binom{x}{n}.
\]
令 $x=0$，有 $\binom00=1$，而 $n>0$ 时 $\binom0n=0$，于是 $\Delta^k\varphi(0)=a_k$。
\end{proof}

\begin{exercise}
证明 Mahler 展开：连续函数 $\phi:\Z_p\to\mathcal O_{\C_p}$ 当且仅当可唯一写成
\[
  \phi(x)=\sum_{n\ge0}a_n\binom{x}{n},\qquad |a_n|\to0.
\]
\end{exercise}

\begin{proof}
先模 $p$ 证明有限层结论。函数 $\Z/p^N\Z\to\F_p$ 的个数为 $p^{p^N}$。多项式族 $\binom{x}{0},\ldots,\binom{x}{p^N-1}$ 在 $0,\ldots,p^N-1$ 上形成上三角矩阵，且对角元为 $1$，故线性无关并构成全部函数的基。恒等式
\[
  \binom{x+p^N}{k}\equiv \binom{x}{k}\pmod p\qquad(0\le k<p^N)
\]
由 Lucas 定理或二项式系数逐位计算得到。

连续函数 $\Z_p\to\F_p$ 在紧空间上局部常值，故分解到某个有限层 $\Z/p^N\Z$，因此有有限 Mahler 展开。对 $\mathcal O_{\C_p}$ 值函数，先模 $p$ 取展开，提升系数后减去，再除以 $p$，逐级进行。得到系数 $a_n$，并且对任意精度只有有限多个系数不被该精度整除，所以 $|a_n|\to0$。唯一性由上一题的公式 $a_k=\Delta^k\phi(0)$ 得到。
\end{proof}

\begin{exercise}
设 $O_{b,r}$ 为闭球上收敛的幂级数 Banach 代数。证明其完备性，并解释题中多项式 $P_{n,\ell}$ 的范数估计和二项式基。
\end{exercise}

\begin{proof}
元素写作 $f=\sum a_m(x-b)^m$，范数为 $\|f\|=\max_m |a_m|r^m$。若 $(f_j)$ 为 Cauchy 序列，则每个系数 $a_{m,j}$ 是 $K$ 中 Cauchy 序列，收敛到 $a_m$。Cauchy 的统一界保证 $|a_m|r^m\to0$，故极限仍在 $O_{b,r}$。乘法中第 $n$ 个系数为有限和 $\sum_{i+j=n}a_ib_j$，非 Archimede 不等式给出
\[
  \|fg\|\le\|f\|\|g\|.
\]
所以这是 Banach 代数。

对
\[
  P_{n,\ell}(T)=\prod_{i=0}^{n-1}(T-\ell p^j-i),
\]
在球 $b+p^j\mathcal O_{\C_p}$ 上估计每个因子。因子 $T-\ell p^j-i$ 的大小由 $i$ 落入相应同余类的 $p$ 进阶决定；把 $1\le k\le j$ 的贡献相加，得到
\[
  \|P_{n,\ell}\|=p^{-\sum_{k=1}^j\nu(b,n,k)}\le \left|\left\lfloor\frac{n}{p^j}\right\rfloor!\right|.
\]
这说明归一化二项式多项式 $\binom{T}{n}$ 在 $C_j$ 中形成法正交基。任意元素唯一写成 $\sum c_n\binom{T}{n}$，且范数为
\[
  \max_n\left|c_n\left(\left\lfloor\frac n{p^j}\right\rfloor!\right)^{-1}\right|.
\]
对偶空间由有界系数序列给出，因为连续线性泛函由它在法正交基上的值决定，而有界性正是题中序列范数有限。
\end{proof}

\begin{exercise}
证明局部解析函数空间 $C^{\mathrm{an}}(\Z_p,K)$ 的基本描述，并说明强对偶与幂级数环的同构。
\end{exercise}

\begin{proof}
闭球 $b+p^j\mathcal O_{\C_p}$ 上的解析函数由 Banach 代数 $O_{b,p^{-j}}$ 描述。把 $\Z_p$ 分解为有限个剩余类，得到
\[
  C_j=\prod_{b\in\Z/p^j\Z}O_{b,p^{-j}}.
\]
局部解析意味着存在某个 $j$，使函数在每个半径 $p^{-j}$ 的球上解析，因此
\[
  C^{\mathrm{an}}(\Z_p,K)=\varinjlim_j C_j.
\]
指数函数 $(1+z)^T$ 的局部解析性来自二项式级数：若 $|z|<r_j$，则 $z^n/\lfloor n/p^j\rfloor!\to0$，故展开落在 $C_j$ 中。

包含 $C_j\to C_{j+1}$ 是紧算子，因为更细球上的限制使归一化系数获得额外衰减。于是强对偶为逆极限
\[
  (C^{\mathrm{an}}(\Z_p,K))'_b\simeq \varprojlim_j C_j'.
\]
对 $\lambda$ 定义
\[
  F_\lambda(z)=\lambda((1+z)^T).
\]
上一段保证右端对 $|z|<1$ 给出幂级数，并且每个满足相应增长条件的幂级数都由某个连续泛函得到。故得到与题中幂级数 Fréchet 代数的拓扑同构。
\end{proof}

\begin{exercise}
证明 Banach 空间之间算子空间的强拓扑就是算子范拓扑。
\end{exercise}

\begin{proof}
Banach 空间 $V$ 的有界集定义为被某个半径球吸收的集合。强拓扑由半范
\[
  p_B(T)=\sup_{v\in B}\|Tv\|
\]
生成。取单位球 $B_1=\{v:\|v\|\le1\}$，则 $p_{B_1}(T)=\|T\|$。反过来，若 $B$ 有界，存在 $C$ 使 $B\subset C B_1$，于是
\[
  p_B(T)\le C\|T\|.
\]
所以所有强拓扑半范与一个算子范半范等价，两个拓扑相同。
\end{proof}

\begin{exercise}
设 $K$ 球完备，$V$ 为 Hausdorff 局部凸空间。证明自然映射 $\delta:V\to (V_s')_s'$ 连续且单射；在标准 Hahn--Banach 假设下为满射到可由点值描述的双对偶。
\end{exercise}

\begin{proof}
对每个 $v\in V$，$\delta_v(\ell)=\ell(v)$。弱拓扑正是点值收敛拓扑，因此 $\ell\mapsto\ell(v)$ 连续，说明 $\delta_v\in(V_s')_s'$，且 $v\mapsto\delta_v$ 连续。若 $v\ne0$，由 Hausdorff 性取不含 $v$ 的闭格 $L$。商 $V/L$ 中 $v$ 的像非零。球完备 Hahn--Banach 定理把一维子空间上的线性泛函扩张到 $V$，得到 $\ell(v)\ne0$，所以 $\delta$ 单射。满射的精确表述依赖所取双对偶拓扑；在弱双对偶中，连续泛函只依赖有限多个点值，因而由 $V$ 中向量给出相应点值泛函。
\end{proof}

\begin{exercise}
设 $A$ 是 Hausdorff 局部凸空间 $V$ 的 $c$ 紧 $\mathcal O$-子模。证明 $A$ 完备且闭，闭子模仍 $c$ 紧，连续线性像仍 $c$ 紧。
\end{exercise}

\begin{proof}
完备性直接来自 $c$ 紧定义：Cauchy 滤子给出每个商 $A/(A\cap L_i)$ 中相容的余类，满射性给出 $A$ 中一点实现所有余类，即滤子收敛到该点。若 $x\in\overline A$，则 $x$ 的余类在每个 $V/L_i$ 中来自 $A$；相容性和 $c$ 紧性给出 $a\in A$ 与 $x$ 在所有商中相同。Hausdorff 性使 $\cap_i L_i=\{0\}$，故 $x=a$，所以 $A$ 闭。

若 $B\subset A$ 是闭 $\mathcal O$-子模，则对递减开格族，先在 $A$ 中由 $c$ 紧性解相容系统，再用 $B$ 的闭性保证所得极限落在 $B$。连续线性映射 $f:V\to W$ 下，商 $f(A)/(f(A)\cap M)$ 的相容系统可拉回到 $A/(A\cap f^{-1}M)$，由 $A$ 的 $c$ 紧性求解，再映到 $W$ 中。
\end{proof}

\begin{exercise}
设 $N$ 是有限生成 $\mathcal O$-模，$M\subset N$，$a\in\mathcal O$ 且 $|a|<1$。证明存在 $M$ 的有限生成子模包含 $aM$。
\end{exercise}

\begin{proof}
取 $N$ 的有限生成元，把 $N$ 看成 $\mathcal O^r$ 的商。只需在 $\mathcal O^r$ 中证明。设 $M_n$ 为 $M$ 在模 $a^n$ 商中的像。因为 $\mathcal O/a^n\mathcal O$ 是 Artin 环，$M_n$ 有有限生成提升 $m_{1,n},\ldots,m_{s_n,n}$。取 $n=2$。若 $m\in M$，它模 $a^2$ 是这些元素的组合，故
\[
  m-\sum c_i m_i\in a^2N.
\]
乘以 $a$ 得
\[
  am-\sum ac_i m_i\in a^3N\subset a^2N.
\]
反复校正得到 $am$ 属于由有限个 $am_i$ 生成的闭包。由于 $N$ 是有限生成模，$a$-进完备化中的闭包与实际子模在 $aM$ 上相交正确，得到有限生成子模包含 $aM$。
\end{proof}

\begin{exercise}
设 $c(I)=\{(a_i):a_i\in F,\ a_i\to0\}$，范数为上确界。证明 $c(I)$ 是 Banach 空间并有法正交基；若 $F$ 离散赋值，则每个 Banach 空间同胚于某个 $c(I)$。
\end{exercise}

\begin{proof}
Cauchy 序列逐坐标收敛，设极限为 $(a_i)$。一致 Cauchy 性保证对任意 $\varepsilon$，除有限个 $i$ 外所有近似项坐标小于 $\varepsilon$，极限也满足 $a_i\to0$，所以完备。标准向量 $e_i$ 给出唯一展开
\[
  x=\sum_i a_ie_i,\qquad \|x\|=\sup_i|a_i|.
\]
这就是法正交基。

当 $F$ 离散赋值时，Banach 空间的单位球模开子球得到向量空间基。逐级提升这些基，构造一组代表 $e_i$，每个向量可按范数层逐步消去主项，得到收敛展开。离散性保证每一步范数严格下降，因此得到与某个 $c(I)$ 的等距线性同构。
\end{proof}

\begin{exercise}
描述 $L(c(I),W)$ 与有界族 $b(I,W)$ 的等距同构，并描述 $L(c(I),c(J))$ 的矩阵。
\end{exercise}

\begin{proof}
给定连续线性算子 $u:c(I)\to W$，令 $w_i=u(e_i)$。因为 $\|e_i\|=1$，有 $\sup_i\|w_i\|\le\|u\|$。反过来，若 $(w_i)$ 有界，对 $x=\sum a_ie_i$ 定义
\[
  ux=\sum_i a_iw_i.
\]
由于 $a_i\to0$ 且 $(w_i)$ 有界，级数收敛，并且 $\|ux\|\le(\sup_i\|w_i\|)\|x\|$。两边互逆，且范数相等。

取 $W=c(J)$，写 $u(e_i)=\sum_j a_{ij}f_j$。条件是每一行 $(a_{ij})_{j\in J}$ 属于 $c(J)$，并且 $\sup_{i,j}|a_{ij}|<\infty$。于是 $L(c(I),c(J))$ 正是这些矩阵，算子范为 $\sup_{i,j}|a_{ij}|$。
\end{proof}

\begin{exercise}
证明全连续算子的矩阵判别，并在局部数域上说明它等同于把有界集送到相对紧集。
\end{exercise}

\begin{proof}
有限秩算子对应只有有限多个非零列的矩阵。其闭包由矩阵列尾趋零刻画：
\[
  \lim_{j\in J}\sup_{i\in I}|a_{ij}|=0.
\]
若列尾趋零，截断到有限列给出有限秩算子并在算子范下收敛到 $u$。若 $u$ 是有限秩算子闭包，任意精度下可由有限列矩阵逼近，因此列尾必须小于该精度。

若 $F$ 是局部数域，$\mathcal O_F$ 紧，有限维闭球紧。列尾趋零说明单位球的像可被有限维紧集加任意小球覆盖，故闭包紧。反过来，若单位球像相对紧，则在坐标投影到尾部空间时必须趋于零，否则可取一列点使尾坐标保持固定下界，无法有收敛子列。
\end{proof}

\begin{exercise}
设 $u\in C(c(I),c(I))$ 且 $\|u\|\le1$。证明模任意非零理想 $\mathfrak a$ 后，$u$ 的像落在有限生成子模中，并定义 $\det(1-tu)$ 的系数极限。
\end{exercise}

\begin{proof}
矩阵判别给出：对每个 $\mathfrak a$，除有限多个列坐标外，所有矩阵元都属于 $\mathfrak a$。因此 $u$ 保持单位球，并在
\[
  c(I)^\circ/\mathfrak a c(I)^\circ
\]
上诱导的像由有限多个标准基的像生成，是有限生成模。于是可以在这个有限生成子模上定义通常行列式
\[
  \det(1-tu_{\mathfrak a}).
\]
若 $\mathfrak b\subset\mathfrak a$，模 $\mathfrak a$ 的行列式是模 $\mathfrak b$ 行列式的像，所以这些多项式构成逆系统。逆极限定义
\[
  \det(1-tu)=\varprojlim_{\mathfrak a}\det(1-tu_{\mathfrak a})\in\mathcal O_F[[t]].
\]
全连续性保证高次外幂迹趋于 $0$，等价地系数 $c_n\to0$。
\end{proof}

\begin{exercise}
证明全连续算子的 Fredholm 行列式存在于整函数环 $F\{\{t\}\}$，并满足有限秩公式和乘法公式。
\end{exercise}

\begin{proof}
若 $\|u\|\le1$，上一题已构造 $\det(1-tu)$。一般情形取 $\lambda\in F^\times$ 使 $\|\lambda u\|\le1$，定义
\[
  \det(1-tu)=\det(1-(\lambda^{-1}t)(\lambda u)).
\]
全连续性给出的系数衰减快于任意固定半径要求，因此级数收敛半径为无穷，属于 $F\{\{t\}\}$。

若 $u$ 有限秩，选取包含 $u(V)$ 的有限维子空间，问题化为有限维线性代数：
\[
  \det(1-tu)=\sum_m(-1)^m\operatorname{Tr}(\wedge^m u)t^m.
\]
对全连续算子，用有限秩算子在算子范下逼近。有限秩情形中的恒等式
\[
  \det(1-tuv)=\det(1-tvu)
\]
以及直和分解下的乘法公式传到极限，得到 Fredholm 行列式的乘法性质。这里的关键是行列式系数由矩阵元的连续多项式表达，故在全连续范数极限下保持。
\end{proof}
 \chapter{赋值环}

\section{章节导读}

本章把局部数论推进到更细的环论和分歧理论。第三章已经说明局部域是完备离散赋值域的分式域；本章要进一步研究赋值环本身、它的提升性质、Witt 环和 Cohen 环，以及 Galois 扩张的分歧过滤。最后几节进入 Fontaine--Wintenberger 的范域思想和完全化构造，这是下一章 $p$ 进 Galois 表示的技术入口。

读这一章时要抓住三条线。第一条线是“提升”：从剩余域或模理想的信息提升到完整环里，光滑环、Hensel 环、Cohen 环和 Witt 环都在解决这个问题。第二条线是“分歧”：Galois 群如何在剩余域、单位群和高阶主单位上作用，下分歧群、上分歧群与 Hasse 函数就是量化工具。第三条线是“无限塔”：全分歧 $\Z_p$ 扩张和 APF 扩张不再是单个有限扩张，而是一串越来越深的局部扩张；范域把这个塔变成特征 $p$ 的完备离散赋值域。

\section{光滑环}

光滑性在这里不是几何图像上的“表面平滑”，而是一个提升性质：当目标环有一个平方为零的理想时，模掉这个小理想以后给出的同态可以提升回来。平方为零表示“一阶无穷小误差”，所以光滑性可以理解为“一阶误差没有障碍”。

\begin{definition}
设 $B$ 是 $A$-代数，$I\subset B$ 是理想，并给 $B$ 取 $I$ 进拓扑。若对任意 $A$-代数 $C$ 和理想 $N\subset C$，只要 $N^2=0$，每个连续 $A$-代数同态
\[
  u:B\longrightarrow C/N
\]
都可提升为 $A$-代数同态 $v:B\to C$，满足自然投射 $j:C\to C/N$ 下有 $jv=u$，则称 $B$ 是 $A$ 上 $I$-光滑的。
\end{definition}

\begin{remark}
若 $I=0$，拓扑条件消失，定义就是普通的形式光滑性。若 $I$ 非零，连续性要求 $B$ 中足够高的 $I$ 幂在离散目标里消失，这正是完备局部环里常用的控制方式。
\end{remark}

\begin{proposition}
若 $B$ 是 $A$ 上 $I$-光滑的，$E$ 是 $A$-代数，$J\subset E$ 是理想，且 $E$ 对 $J$ 进拓扑完备，则任意连续同态 $B\to E/J$ 都可提升为连续同态 $B\to E$。
\end{proposition}

\begin{proof}
先把 $B\to E/J$ 提升到 $E/J^2$。因为 $J/J^2$ 在 $E/J^2$ 中平方为零，光滑性给出提升。若已经提升到 $E/J^n$，再用平方为零理想 $J^n/J^{n+1}\subset E/J^{n+1}$ 提升到 $E/J^{n+1}$。这样得到相容系统
\[
  B\to E/J,\quad B\to E/J^2,\quad B\to E/J^3,\ldots.
\]
由完备性 $E\simeq\varprojlim E/J^n$，相容系统唯一给出 $B\to E$。
\end{proof}

\begin{definition}
设 $B$ 为带 $I$ 进拓扑的 $A$-代数，$N$ 为离散 $B$-模。连续对称 $2$ 上闭链是 $A$-双线性映射 $f:B\times B\to N$，满足
\[
  xf(y,z)-f(xy,z)+f(x,yz)-f(x,y)z=0,\qquad f(x,y)=f(y,x),
\]
并且 $I$ 的高次幂进入任意一个变量时 $f$ 为零。若存在 $A$-线性映射 $g:B\to N$ 使
\[
  f(x,y)=xg(y)-g(xy)+g(x)y,
\]
则称 $f$ 分裂。
\end{definition}

\begin{proposition}
设有无穷多个 $n$ 使 $B/I^n$ 为投射 $A$-模。若所有取值于离散 $B$-模的连续对称 $2$ 上闭链都分裂，则 $B$ 是 $A$ 上 $I$-光滑的。
\end{proposition}

\begin{proof}
给定 $u:B\to C/N$，其中 $N^2=0$。连续性给出 $u(I^m)=0$。选 $n$ 使 $B/I^n$ 投射，并先把 $u$ 当作 $A$-模映射提升为 $h:B\to C$。乘法失败量
\[
  f(x,y)=h(xy)-h(x)h(y)
\]
落在 $N$ 中，因为模 $N$ 后 $u$ 是环同态。结合律正是 $2$ 上闭链恒等式，对称性来自交换环乘法。假设给出 $g$ 使 $f=\delta g$。令 $v=h+g$，计算得
\[
  v(xy)=v(x)v(y).
\]
因此 $v$ 是所需的 $A$-代数提升。
\end{proof}

\begin{proposition}
设 $(A,\mathfrak m,k)$ 为局部环，$B$ 是平坦 $A$-代数，且 $B_0=B\otimes_Ak$ 是 $k$ 上 $0$-光滑的。则 $B$ 是 $A$ 上 $\mathfrak mB$-光滑的。
\end{proposition}

\begin{proof}
先把问题降到 $A/\mathfrak m^r$。若 $\mathfrak m$ 幂零，平坦有限层上的 $B$ 是投射模。任一 $2$ 上闭链 $f$ 先模 $\mathfrak mN$ 分裂；提升分裂映射后，剩余误差取值于 $\mathfrak mN$。重复这个过程，误差依次进入
\[
  \mathfrak m^2N,\quad \mathfrak m^3N,\quad \ldots .
\]
由于 $\mathfrak m$ 幂零，有限步后误差为零。于是命题 10.1 的判别条件成立。
\end{proof}

\section{离散赋值环}

\subsection{赋值环}

\begin{definition}
整环 $R$ 称为赋值环，如果对分式域 $K=\operatorname{Frac}(R)$ 的任意非零元 $x$，有
\[
  x\in R\quad\text{或}\quad x^{-1}\in R.
\]
\end{definition}

赋值环的核心特征是理想全序。若 $\mathfrak a,\mathfrak b$ 是理想，取 $x\in\mathfrak a$ 不在 $\mathfrak b$，则对任意 $y\in\mathfrak b$，$x/y\notin R$，所以 $y/x\in R$，从而 $y\in\mathfrak a$。因此 $\mathfrak b\subset\mathfrak a$。这说明赋值环天然是局部环。

\begin{proposition}
环 $A$ 是主理想整环且只有一个非零素理想，当且仅当 $A$ 是 Noether 局部环，并且其极大理想由非幂零元生成。
\end{proposition}

\begin{proof}
若 $A$ 是这样的主理想整环，唯一非零素理想必是极大理想，并由不可约元生成，结论成立。反过来，设极大理想 $\mathfrak m=(\pi)$，且 $\pi$ 非幂零。Krull 交定理给出 $\cap_n\mathfrak m^n=0$。任意非零元 $x$ 落在唯一的差集 $\mathfrak m^n\setminus\mathfrak m^{n+1}$ 中，所以 $x=\pi^nu$，其中 $u\notin\mathfrak m$，即 $u$ 可逆。于是每个非零理想由其中赋值最小的元素生成，$A$ 是主理想整环。
\end{proof}

\begin{definition}
满足上面条件的局部主理想整环称为离散赋值环，简称 DVR。其极大理想的生成元 $\pi$ 称为素元或单化子。若 $K=\operatorname{Frac}(A)$，每个 $x\in K^\times$ 唯一写为
\[
  x=u\pi^n,\qquad u\in A^\times,\ n\in\Z,
\]
定义 $v(x)=n$，并令 $v(0)=+\infty$。
\end{definition}

\subsection{存在定理}

\begin{theorem}
设 $(A,tA,k)$ 是离散赋值环，$K/k$ 是域扩张。存在包含 $A$ 的离散赋值环
\[
  (B,tB,K),
\]
其剩余域为 $K$。
\end{theorem}

\begin{proof}
先处理超越部分。取 $K/k$ 的超越基 $(x_\lambda)$，用多项式变量 $X_\lambda$ 提升它们，局部化 $A[X_\lambda]$ 于 $t$ 生成的素理想，得到仍以 $t$ 为素元的 DVR，其剩余域为 $k(x_\lambda)$。

剩下可设 $K/k$ 代数。考虑所有对 $(B,\varphi)$：$A\subset B$，$B$ 是 DVR，$\varphi:B\to K$ 是局部同态，核为 $tB$。按包含关系用 Zorn 引理取极大元。若 $\varphi(B)\ne K$，取 $a\in K\setminus\varphi(B)$，取它在 $\varphi(B)$ 上的最小多项式并提升为 $B[X]$ 中首一多项式 $f$。令 $\alpha$ 为 $f$ 的根，$B'=B[\alpha]$。因为
\[
  B'/tB'\simeq \varphi(B)[X]/(\overline f)
\]
是域，$B'$ 是局部有限 $B$-代数，极大理想由 $t$ 生成，因此仍是 DVR，矛盾于极大性。所以剩余域已经是 $K$。
\end{proof}

\begin{theorem}
设 $R$ 是特征 $0$ 的 DVR，极大理想为 $pR$，剩余域为 $k$。若 $(A,\mathfrak m,K)$ 是完备局部环，且给定域同态 $k\to K$，则存在局部同态 $R\to A$ 诱导该剩余域同态。
\end{theorem}

\begin{proof}
从 $\Z_{(p)}\to A$ 开始。因为 $R/pR=k$ 是素子域上可分扩张，前面平坦光滑命题给出 $R$ 在 $p$ 进方向上的光滑性。将 $R\to K=A/\mathfrak m$ 逐步提升到 $A/\mathfrak m^n$，再由 $A\simeq\varprojlim A/\mathfrak m^n$ 得到 $R\to A$。
\end{proof}

\subsection{\texorpdfstring{$p$ 环}{p 环}}

\begin{definition}
设环 $A$ 有递减理想列 $\mathfrak a_1\supset\mathfrak a_2\supset\cdots$，满足 $\mathfrak a_n\mathfrak a_m\subset\mathfrak a_{n+m}$、$\cap_n\mathfrak a_n=0$，且 $A$ 对此拓扑完备。若 $k=A/\mathfrak a_1$ 是特征 $p$ 完全环，则称 $A$ 为 $p$ 环。若 $\mathfrak a_n=p^nA$ 且 $p$ 不是零除子，则称 $A$ 为严格 $p$ 环。
\end{definition}

\begin{proposition}
设 $(A,\mathfrak a_n)$ 是 $p$ 环，$k=A/\mathfrak a_1$。存在唯一代表映射
\[
  f:k\to A
\]
满足 $f(\lambda^p)=f(\lambda)^p$。它是乘性的，且 $a\in f(k)$ 当且仅当对每个 $n$ 都有 $a=b^{p^n}$ 的解。
\end{proposition}

\begin{proof}
因为 $k$ 完全，每个 $\lambda$ 有唯一的 $p^n$ 次根 $\lambda^{p^{-n}}$。任取提升 $a_n$，考虑序列 $a_n^{p^n}$。若换提升，差落在 $\mathfrak a_1$，取 $p^n$ 次幂后进入越来越深的邻域；完备性给出极限，定义为 $f(\lambda)$。该构造直接给出 $f(\lambda^p)=f(\lambda)^p$。乘法由
\[
  (a_nb_n)^{p^n}=a_n^{p^n}b_n^{p^n}
\]
取极限得到。判别条件来自同一个极限构造的唯一性。
\end{proof}

\begin{proposition}
若 $A,A'$ 是 $p$ 环，剩余环分别为 $k,k'$，则剩余环同态 $\overline\varphi:k\to k'$ 至多有一个连续提升 $\varphi:A\to A'$ 与代表映射相容。若 $A$ 是严格 $p$ 环且 $k$ 完全，则这样的提升存在。
\end{proposition}

\begin{proof}
唯一性来自代表元展开。严格 $p$ 环中每个元素唯一写成收敛级数
\[
  a=\sum_{n\ge0}p^nf(\lambda_n).
\]
若提升存在，就必须满足
\[
  \varphi(a)=\sum_{n\ge0}p^nf'(\overline\varphi(\lambda_n)).
\]
右端收敛并定义连续环同态；加法和乘法由代表元运算多项式的相容性保证。
\end{proof}

\begin{theorem}
若 $k$ 是特征 $p$ 完全环，则存在唯一严格 $p$ 环 $W(k)$ 以 $k$ 为剩余环。
\end{theorem}

\begin{proof}
先对多项式完美环构造 $p$ 进完备化，再对一般 $k$ 用满射从自由完美环到 $k$ 并取相应商。存在性依赖上一命题的同态提升，唯一性也由上一命题给出。这个严格 $p$ 环就是 Witt 环。
\end{proof}

\subsection{完备离散赋值环}

\begin{proposition}
设 $A$ 是完备离散赋值环，剩余域为 $\kappa$，素元为 $\pi$，$S\subset A$ 是 $\kappa$ 的代表集。每个非零 $a\in A$ 唯一写成
\[
  a=\sum_{n\ge n_0}s_n\pi^n,\qquad s_n\in S.
\]
\end{proposition}

\begin{proof}
先取 $s_{n_0}$ 表示 $a/\pi^{n_0}$ 的剩余类，减去 $s_{n_0}\pi^{n_0}$ 后赋值升高。重复这个过程得到 Cauchy 部分和。完备性给出收敛，唯一性由最低非零项的剩余类唯一性得出。
\end{proof}

\begin{theorem}
若 $A$ 是同特征完备离散赋值环，剩余域 $\kappa$ 完全，则
\[
  A\simeq \kappa[[T]]
\]
并且分式域同构于 $\kappa((T))$。
\end{theorem}

\begin{proof}
由同特征条件可把 $\kappa$ 嵌入 $A$ 为系数域。选素元 $\pi$，映射 $\kappa[[T]]\to A$，$T\mapsto\pi$。上一命题给出每个元素有唯一 $\pi$ 进展开，所以这是同构。
\end{proof}

\begin{theorem}
若 $A$ 是特征 $0$、剩余特征 $p$ 的完备 DVR，剩余域 $k$ 完全，且 $e=v_A(p)$，则存在唯一嵌入 $W(k)\to A$ 提升剩余域恒等。若 $\pi$ 是 $A$ 的素元，则 $A$ 是 $W(k)[\pi]$ 的有限扩张；绝对无分歧情形 $e=1$ 时 $A\simeq W(k)$。
\end{theorem}

\begin{proof}
$A$ 是 $p$ 环，代表元提升给出 $W(k)\to A$。任意 $a\in A$ 可写成
\[
  a=\sum_{n\ge0}[\lambda_n]\pi^n
\]
的收敛形式。若 $e=1$，$\pi$ 可取为 $p$，展开就是 Witt 展开，得到 $A\simeq W(k)$。一般情形中 $\pi$ 满足 Eisenstein 型多项式，因而 $A$ 是 $W(k)$ 上有限。
\end{proof}

\section{Witt 环}

\subsection{Witt 多项式}

Witt 环的目标是把特征 $p$ 的加法和乘法提升到混合特征。它不能逐坐标使用普通加法，因为 Teichmuller 代表只保持乘法。解决方法是引入 ghost 分量。

\begin{definition}
Witt 多项式定义为
\[
  W_n(X_0,\ldots,X_n)=X_0^{p^n}+pX_1^{p^{n-1}}+\cdots+p^nX_n.
\]
Witt 向量 $(x_0,x_1,\ldots)$ 的第 $n$ 个 ghost 分量为 $W_n(x_0,\ldots,x_n)$。
\end{definition}

\begin{theorem}
对任意整数系数多项式 $\Phi(X,Y)$，存在唯一一列整数系数多项式 $\phi_n$，使
\[
  W_n(\phi_0,\ldots,\phi_n)=\Phi(W_n(X_0,\ldots,X_n),W_n(Y_0,\ldots,Y_n)).
\]
特别地，取 $\Phi=X+Y$ 和 $\Phi=XY$ 可在 $W(A)$ 上定义加法和乘法。
\end{theorem}

\begin{proof}
按 $n$ 递归。已知 $\phi_0,\ldots,\phi_{n-1}$ 后，等式右端减去前面部分，剩余项被 $p^n$ 整除，并且可唯一写成 $p^n\phi_n$。非平凡点在于 $\phi_n$ 仍有整数系数；这是 Witt 多项式构造的整性结论，可由二项式系数的 $p$ 整除性逐步验证。
\end{proof}

\subsection{\texorpdfstring{$V$ 和 $F$}{V 和 F}}

\begin{definition}
在 Witt 向量上定义 Verschiebung 和 Frobenius：
\[
  V(x_0,x_1,\ldots)=(0,x_0,x_1,\ldots),
\]
\[
  F(x_0,x_1,\ldots)=(x_0^p,x_1^p,\ldots)
\]
当底环为特征 $p$ 完全环时，第二式给出实际 Frobenius。
\end{definition}

\begin{proposition}
设 $k$ 为特征 $p$ 环。在 $W_n(k)$ 中有
\[
  FV=VF=p,\qquad V(a) b=V(aF(b)).
\]
若 $k$ 完全，则 $F$ 是自同构，且
\[
  W(k)\simeq\varprojlim_n W(k)/p^nW(k).
\]
\end{proposition}

\begin{proof}
在 ghost 分量上检查最清楚。$V$ 把 ghost 分量右移并乘以 $p$，$F$ 把第 $n$ 个 ghost 分量变为第 $n+1$ 个。于是 $FV$ 和 $VF$ 都表现为乘以 $p$。乘法公式同样在 ghost 分量上化为普通恒等式。若 $k$ 完全，Frobenius 可逆，$p$ 进完备性由截断 Witt 向量的极限描述得到。
\end{proof}

\section{Hensel 环}

Hensel 性是“近似根可提升”的系统化表达。它是完备局部环最常用的性质之一，但 Hensel 环不必完备；Hensel 化就是在保持局部结构的前提下添加足够的根提升。

\begin{definition}
局部环 $(A,\mathfrak m,k)$ 称为 Hensel 环，如果对任意首一多项式 $f\in A[X]$，只要模 $\mathfrak m$ 后有分解
\[
  \overline f=\overline g\,\overline h,\qquad (\overline g,\overline h)=1,
\]
就能提升为 $f=gh$，其中 $g,h\in A[X]$ 首一并提升 $\overline g,\overline h$。
\end{definition}

\begin{proposition}
对局部环 $A$，以下条件等价：
\begin{enumerate}[label=(\arabic*)]
\item $A$ 是 Hensel 环；
\item 每个模 $\mathfrak m$ 的简单根都唯一提升为 $A$ 中的根；
\item 有限 $A$-代数的幂等元可从剩余代数中提升。
\end{enumerate}
\end{proposition}

\begin{proof}
多项式分解中的互素因子对应有限代数中的互补幂等元。简单根是分解出一次因子的情形。幂等元提升可用方程 $e^2-e=0$ 和 Newton 修正完成；反过来，幂等元把有限代数分解成两个因子，从而给出多项式分解。
\end{proof}

\begin{theorem}
任意局部环 $A$ 存在 Hensel 化 $A^h$，并满足泛性质：对任意 Hensel 局部环 $B$ 和局部同态 $A\to B$，唯一分解为
\[
  A\longrightarrow A^h\longrightarrow B.
\]
\end{theorem}

\begin{proof}
取所有 etale 局部 $A$-代数组成的有向系统，并取直极限。etale 性保证剩余域和局部结构可控；直极限中所有有限的 Hensel 分解问题都已经在某个 etale 邻域解决。泛性质来自 etale 提升的唯一性。
\end{proof}

\section{Cohen 环}

\subsection{\texorpdfstring{$C(k)$}{C(k)}}

Cohen 环把任意特征 $p$ 剩余域提升到混合特征完备局部环。若剩余域完全，答案是 Witt 环；若不完全，需要加入 $p$-基的提升。

\begin{theorem}
设 $A$ 是 Hausdorff 完备局部环，剩余域 $k$ 特征 $p>0$。若给定从某个 Cohen 环 $C(k)$ 到 $A$ 的剩余域同态，则它唯一提升为局部同态 $C(k)\to A$。
\end{theorem}

\begin{proof}
Cohen 环由代表元和 $p$ 进完备性刻画。先在剩余域选择一组 $p$-基，并把它们提升到 $A$。对 $p$-基生成的完美部分使用 Witt 提升，对非完美方向使用形式变量并 $p$ 进完备化。唯一性由代表元展开和 $p$ 进分离性给出。
\end{proof}

\subsection{\texorpdfstring{$E$}{E}}

设 $E$ 是特征 $p$ 的局部域。与它对应的绝对无分歧完备 DVR 记为 $\mathcal O_E$，其分式域是混合特征的局部域。这个构造为后续范域和 Fontaine 环提供系数环。

\begin{proposition}
若 $E$ 是完全剩余域情形，则 $\mathcal O_E\simeq W(E)$。若 $E$ 是一般特征 $p$ 局部域，则 $\mathcal O_E$ 由 Cohen 环构造，并在非分歧方向上由剩余域控制。
\end{proposition}

\begin{proof}
完全情形直接由严格 $p$ 环唯一性得到。一般情形先取 Cohen 环，再对分式域的最大非分歧部分取整数环并完备化。剩余域不变，素元为 $p$，所以绝对无分歧。
\end{proof}

\section{分歧群}

\subsection{有限 Galois 扩张的下分歧群}

设 $L/K$ 为完备离散赋值域的有限 Galois 扩张，$G=\Gal(L/K)$。分歧群要测量 $g\in G$ 离恒等变换有多近。
这里的“近”不是拓扑群里的抽象距离，而是看 $g$ 对整数环的高阶近似是否保持不动。若 $g$ 在剩余域上已经不是恒等，那么它只属于 $G_{-1}$ 而不属于 $G_0$；若它在剩余域上恒等，但在一阶主单位上还有非平凡作用，那么它在 $G_0$ 中但不在 $G_1$ 中；若它连模 $\mathfrak p_L^{i+1}$ 都看起来像恒等，就进入 $G_i$。因此下分歧群是一台显微镜：$i$ 越大，观察越精细，仍然看不出 $g$ 与恒等的区别，说明 $g$ 越深地属于野分歧。

\begin{definition}
取 $\mathcal O_L$ 作为 $\mathcal O_K$-代数的生成元 $x$。定义
\[
  i_G(1)=+\infty,\qquad i_G(g)=v_L(g(x)-x)\quad(g\ne1).
\]
再定义下分歧群
\[
  G_i=\{g\in G:i_G(g)\ge i+1\}\qquad(i\ge0),
\]
并令 $G_{-1}=G$。
\end{definition}

等价地，$g\in G_i$ 当且仅当 $g$ 在 $\mathcal O_L/\mathfrak p_L^{i+1}$ 上作用为恒等。于是
\[
  G=G_{-1}\supset G_0\supset G_1\supset\cdots.
\]
$G_0$ 是惯性群，$G_1$ 是野惯性群。

这个等价刻画也解释了为什么定义不依赖所选生成元 $x$。若 $g$ 在 $\mathcal O_L/\mathfrak p_L^{i+1}$ 上恒等，那么当然对任何生成元都满足 $g(x)-x\in\mathfrak p_L^{i+1}$。反过来，若对一个生成元 $x$ 有足够高的相合性，则 $\mathcal O_L$ 中每个元素都可写成 $x$ 的多项式并取极限，$g$ 对这些多项式的作用也在同一阶数内相合。完备性把多项式近似推广到整个整数环。

\begin{proposition}
若 $H\subset G$ 是子群，则 $H_i=H\cap G_i$。若 $H$ 正规，$K'=L^H$，则商扩张的分歧可由
\[
  i_{G/H}(sH)=\frac1{e_{L/K'}}\sum_{t\in H} i_G(st)
\]
计算。
\end{proposition}

\begin{proof}
子群情形只是在同一个 $L$ 上限制自同构，定义不变。商群公式来自范数计算。若 $y$ 生成 $\mathcal O_{K'}$，则 $s(y)-y$ 的 $L$ 赋值等于所有 $st(x)-x$ 的乘积赋值总和，再除以 $e_{L/K'}$ 转成 $K'$ 的赋值。
\end{proof}

\begin{definition}
Hasse 函数定义为
\[
  \varphi_{L/K}(u)=\int_0^u\frac{dt}{[G_0:G_t]}\quad(u\ge0),
\]
并在 $[-1,0]$ 上取 $\varphi(u)=u$。它是分段线性、连续、递增函数。
\end{definition}

引入 $\varphi$ 的动机是修正下编号在商群下的不相容。若只用下分歧群，取商扩张时每一层的编号会被分歧指数重新缩放，公式会很难看。Hasse 函数把下编号 $u$ 换成上编号 $v=\varphi(u)$，它的斜率正好是 $1/[G_0:G_u]$，也就是“当前分歧层相对于惯性群的归一化大小”。这个归一化使得商群中的分歧层可以直接由原群的上分歧层取商得到。

\begin{proposition}
对 $u\ge -1$，
\[
  \varphi_{L/K}(u)=\frac1{|G_0|}\sum_{s\in G}\inf\{i_G(s),u+1\}-1.
\]
若 $H\triangleleft G$，则
\[
  \varphi_{L/K}=\varphi_{K'/K}\circ\varphi_{L/K'}.
\]
\end{proposition}

\begin{proof}
右端是分段线性函数，在非整数区间的斜率等于 $1/[G_0:G_u]$，且在 $0$ 处为 $0$，所以等于定义的积分函数。复合公式由商群分歧公式比较斜率得到。
\end{proof}

\subsection{有限 Galois 扩张的上分歧群}

\begin{definition}
令 $\psi_{L/K}$ 为 $\varphi_{L/K}$ 的反函数。上分歧群定义为
\[
  G^v=G_{\psi_{L/K}(v)}.
\]
\end{definition}

上编号的意义在于它和商群相容。下编号对子群最自然，上编号对商群最自然。
这句话在使用时非常关键。若研究子扩张 $L/L^H$，保留下编号更方便，因为分歧群只是取交 $H\cap G_i$。若研究商扩张 $L^H/K$，上编号更方便，因为分歧群直接变成 $G^vH/H$。类域论和局部常数公式经常要在中间域之间来回移动，因此实际计算中通常先用下编号做局部估值，再用 Hasse 函数换到上编号来传递到商扩张。

\begin{proposition}
若 $H\subset G$，则 $H^u=H\cap G^u$。若 $H\triangleleft G$，则
\[
  (G/H)^v=G^vH/H.
\]
\end{proposition}

\begin{proof}
子群结论由下编号子群结论和 $\varphi$ 的限制关系得到。商群结论把 $\varphi_{L/K}=\varphi_{K'/K}\circ\varphi_{L/K'}$ 取反函数后代入定义。
\end{proof}

\subsection{Hasse--Arf 定理}

\begin{proposition}
设 $\pi$ 是 $L$ 的素元。对 $i\ge0$，映射
\[
  \theta_i:G_i/G_{i+1}\longrightarrow U_L^{(i)}/U_L^{(i+1)},\qquad
  g\mapsto g(\pi)/\pi
\]
是单同态，且与 $\pi$ 的选择无关。
\end{proposition}

\begin{proof}
若换素元 $\pi'=\pi u$，则
\[
  \frac{g(\pi')}{\pi'}=\frac{g(\pi)}{\pi}\frac{g(u)}u.
\]
当 $g\in G_i$ 时，$g(u)/u\equiv1\pmod{\mathfrak p_L^{i+1}}$，所以余类不变。若 $\theta_i(g)$ 平凡，则 $g(\pi)\equiv\pi\pmod{\mathfrak p_L^{i+1}}$，再结合 $g$ 对剩余域和单位的作用，得到 $g\in G_{i+1}$。
\end{proof}

\begin{corollary}
$G_0/G_1$ 是剩余域乘法群中的有限循环子群，其阶与剩余特征互素。若剩余特征为 $p$，则 $G_i/G_{i+1}$ 对 $i\ge1$ 是初等 $p$-群，$G_1$ 是 $p$-群。
\end{corollary}

\begin{theorem}
若 $L/K$ 是有限 Abel 扩张，且剩余域扩张可分，则上分歧过滤的跳跃点都是整数。
\end{theorem}

\begin{proof}
这是 Hasse--Arf 定理。证明先化到完全分歧且 $G=G_1$ 的情形。若 $G$ 是循环 $p$-群，设生成元为 $s$，研究 $i_G(s^{p^j})$。利用
\[
  \sum_{r=0}^{p-1}s^r(x)
\]
的赋值严格升高，证明跳跃满足同余关系，从而 Herbrand 函数在跳跃处取整数。一般 Abel $p$-群通过商到循环因子并归纳。若还存在 tame 部分 $G_0/G_1$，前面共轭作用公式给出 tame 阶整除相关跳跃指数，最终仍得整数性。

把证明路线再细分。第一步去掉未分歧部分，因为未分歧部分只影响 $G/G_0$，不改变正的分歧跳跃。第二步去掉顺分歧部分。顺分歧层只在 $0$ 附近产生跳跃，而 $G_0/G_1$ 的阶与 $p$ 互素；它通过乘法群中的根作用控制，与野分歧的 $p$-群部分相互分离。

第三步处理循环 $p^m$ 阶全分歧扩张。设
\[
i_j=i_G(s^{p^j})\qquad(0\le j\le m-1).
\]
关键估计来自迹型和式。若 $x$ 的赋值为 $n$，则
\[
v_L\!\left(\sum_{r=0}^{p-1}s^{rp^j}(x)\right)>n
\]
在适当层次成立；这说明从 $s^{p^j}$ 到 $s^{p^{j+1}}$ 时，分歧跳跃不能任意移动，而必须满足
\[
i_{j+1}\equiv i_j \pmod {p^j}
\]
这类同余约束。把这些同余代入 Herbrand 函数的分段线性公式，正好得到每个上跳跃为整数。

第四步从循环 $p$-群推广到一般 Abel $p$-群。取一个阶为 $p$ 的商或子群，利用上编号对商群相容，把已知整数跳跃传递到商扩张；再对阶数归纳。这里 Abel 性是必要条件，因为需要把分歧过滤与取商、取子群的过程同时控制，并用交换性保证共轭作用不制造额外扭曲。

定理的含义是：Abel 扩张中分歧虽然可以很深，但上编号跳跃不会落在分数位置。这个整数性是局部类域论中导子、Artin 导子和特征标分歧层能用整数编号描述的根本原因。
\end{proof}

\subsection{局部域的差别式}

\begin{proposition}
设 $L/K$ 是局部域有限扩张，差别式为 $\mathfrak D_{L/K}=\mathfrak p_L^d$。则
\[
  \operatorname{Tr}_{L/K}(\mathfrak p_L^n)=\mathfrak p_K^{\lfloor(d+n)/e_{L/K}\rfloor}.
\]
\end{proposition}

\begin{proof}
迹对偶告诉我们
\[
  \operatorname{Tr}(\mathfrak p_L^n)\subset \mathfrak p_K^r
  \quad\Longleftrightarrow\quad
  \mathfrak p_L^n\subset \mathfrak p_K^r\mathfrak D_{L/K}^{-1}.
\]
右端的 $L$-赋值为 $re_{L/K}-d$，所以最大 $r$ 正是 $\lfloor(d+n)/e_{L/K}\rfloor$。
\end{proof}

\begin{proposition}
若 $L/K$ 是有限 Galois 局部域扩张，则
\[
  v_L(\mathfrak D_{L/K})=\sum_{g\ne1}i_G(g)=\sum_{i\ge0}(|G_i|-1).
\]
\end{proposition}

\begin{proof}
取 $\mathcal O_L=\mathcal O_K[x]$，令 $f$ 为 $x$ 的最小多项式。差别式由 $f'(x)$ 生成，而
\[
  f'(x)=\prod_{g\ne1}(x-gx).
\]
取赋值得到第一等式。第二等式是把每个 $g\ne1$ 对应的 $i_G(g)$ 按 $i_G(g)\ge i+1$ 的层数计数。
\end{proof}

\section{单位群}

\subsection{无分歧}

设 $L/K$ 是局部域有限扩张。记
\[
  U_K=\mathcal O_K^\times,\qquad U_K^{(n)}=1+\mathfrak p_K^n\quad(n\ge1).
\]

\begin{proposition}
若 $L/K$ 无分歧，则范数映射满足
\[
  N_{L/K}(U_L)=U_K,\qquad N_{L/K}(U_L^{(n)})=U_K^{(n)}.
\]
\end{proposition}

\begin{proof}
在剩余域上，有限域扩张的范数映射满射，所以 $U_L\to U_K$ 模主单位后满射。对主单位，使用
\[
  (1+x)\mapsto 1+\operatorname{Tr}_{L/K}(x)+\text{高阶项}
\]
以及无分歧时迹映射 $\mathfrak p_L^n\to\mathfrak p_K^n$ 满射，逐层提升得到结论。
\end{proof}

\subsection{素阶全分歧扩张}

\begin{proposition}
设 $L/K$ 是素数阶 $\ell$ 的全分歧 Galois 扩张。若 $t$ 是唯一的下分歧跳跃，则
\[
  \operatorname{Tr}(\mathfrak p_L^n)=\mathfrak p_K^{\left\lfloor ((t+1)(\ell-1)+n)/\ell\right\rfloor}.
\]
范数在单位过滤上的像由 $t$ 控制。
\end{proposition}

\begin{proof}
差别指数为 $(t+1)(\ell-1)$，代入前面迹公式得到第一式。范数展开
\[
  N(1+x)=1+\operatorname{Tr}(x)+\text{更高阶项}
\]
显示其主项由迹控制，而高阶项落在更深的单位层中。逐层比较可确定 $N(U_L^{(n)})$ 位于哪个 $U_K^{(m)}$。
\end{proof}

\subsection{全分歧扩张}

\begin{proposition}
若 $L/K$ 是有限全分歧 Galois 扩张，则范数映射保持单位群过滤，并且其在每一层上的像可由上分歧群控制。
\end{proposition}

\begin{proof}
把 Galois 群取一条组成因子为素阶循环群的正规列。每一步用素阶全分歧情形计算范数对单位层的作用，再把各步复合。上分歧编号对商群相容，保证分层计算可以拼合。
\end{proof}

\section{最大交换扩张}

\subsection{无分歧扩张}

\begin{proposition}
完备离散赋值域 $K$ 的有限无分歧扩张与其剩余域的有限可分扩张等价。最大无分歧扩张 $K^{\mathrm{nr}}$ 的 Galois 群同构于剩余域可分闭包的 Galois 群。
\end{proposition}

\begin{proof}
无分歧扩张的整数环由剩余域扩张的 Hensel 提升给出。若 $\overline f$ 是剩余域上的可分不可约多项式，取首一提升 $f\in\mathcal O_K[X]$，Hensel 性给出未分歧扩张。反过来，无分歧扩张的分歧指标为 $1$，全部次数来自剩余域。
\end{proof}

\begin{proposition}
若 $K$ 是局部域，则最大无分歧扩张的完备化 $\widehat K^{\mathrm{nr}}$ 有代数闭剩余域；其 Frobenius 自同构提升剩余域 Frobenius，并控制无分歧 Galois 群的拓扑生成元。
\end{proposition}

\subsection{全分歧扩张}

\begin{lemma}
设 $K$ 是完备离散赋值域，剩余域代数封闭。若 $L/K$ 是有限循环扩张，则存在 $\xi\in L^\times$ 使范数和 Galois 作用满足 Hilbert 90 型关系
\[
  \xi=\sigma(\eta)/\eta
\]
并可用素元调整赋值。
\end{lemma}

\begin{proof}
乘法 Hilbert 90 给出 $1$ 范数元素的表示。剩余域代数闭保证单位部分的范数障碍消失，素元部分由赋值控制。这样可把循环扩张中的 Galois 信息转化为单位商计算。
\end{proof}

\begin{proposition}
若 $K$ 的剩余域代数闭，则有限 Abel 全分歧扩张的结构可由范数商和单位过滤描述。局部域情形中，这与局部互反映射相容。
\end{proposition}

\subsection{顺分歧扩张}

\begin{definition}
有限扩张 $L/K$ 称为顺分歧扩张，如果分歧指数与剩余特征互素，且剩余域扩张可分。
\end{definition}

\begin{proposition}
顺分歧扩张由剩余域扩张和素元的 $e$ 次根控制。若 $K$ 含有足够的 $e$ 次单位根，全分歧顺分歧扩张由 Eisenstein 多项式 $X^e-\pi u$ 给出。
\end{proposition}

\begin{proof}
顺分歧意味着 $G_1=1$，全部分歧停在 tame 惯性层。$G_0/G_1$ 嵌入剩余域乘法群，所以循环且阶与 $p$ 互素。素元的变换给出 $e$ 次单位根作用，从而得到所述 Kummer 型描述。
\end{proof}

\section{全分歧 \texorpdfstring{$\Z_p$}{Zp} 扩张}

设 $F_\infty/F$ 是全分歧 $\Z_p$ 扩张，$F_n$ 是中间层且 $[F_n:F]=p^n$。这类塔在 Iwasawa 理论和 $p$ 进 Hodge 理论中反复出现。

\begin{lemma}
设上分歧跳跃序列为 $y_n$。存在 $m_0$，使 $n\ge m_0$ 时
\[
  y_{n+1}=y_n+e(F).
\]
\end{lemma}

\begin{proof}
$\Gamma\simeq\Z_p$ 的上分歧子群给出一列开子群。对足够深的层，乘以 $p$ 的自同态和单位群过滤的指数变化稳定，Herbrand 函数的斜率呈周期性，因而跳跃每上升一层增加绝对分歧指数 $e(F)$。
\end{proof}

\begin{proposition}
存在常数 $c$ 和有界序列 $(a_n)$，使
\[
  v_F(\mathfrak D_{F_n/F})=n+c+p^{-n}a_n.
\]
\end{proposition}

\begin{proof}
差别指数由分歧群积分公式给出。跳跃最终线性增长，因此求和后主项为 $n$，剩余项由有限初始层和有界误差组成，写成 $c+p^{-n}a_n$。
\end{proof}

\begin{theorem}
对有限扩张 $L/F_\infty$，有
\[
  \operatorname{Tr}_{L/F_\infty}(\mathcal O_L)\supset \mathfrak m_{F_\infty}.
\]
\end{theorem}

\begin{proof}
先在有限层 $L_n/F_n$ 上用差别式估计迹像。前一命题保证差别指数经 $p^{-n}$ 归一后有界。令 $n$ 足够大，迹公式给出每个正赋值元素都落入迹像。取极限得到无限层结论。
\end{proof}

\section{范域}

\subsection{APF 扩张}

\begin{definition}
设 $L/K$ 是无限 Galois 扩张。若对每个 $u\ge-1$，由 $K$ 的绝对 Galois 群上分歧子群与 $G_L$ 生成的子群是开子群，则称 $L/K$ 是 APF 扩张。
\end{definition}

APF 的含义是“分歧足够深但仍可控”。全分歧 $\Z_p$ 扩张是基本例子。

这个定义看起来抽象，可以换成更直观的说法：当有限中间域 $E$ 沿着 $L/K$ 往上走时，分歧跳跃必须趋向无穷。若跳跃停在有限高度，范数映射在高层主单位上不会越来越接近 Frobenius，逆极限就无法稳定成一个特征 $p$ 域。APF 条件正是排除这种停滞，保证“越往塔上走，分歧越深，范数越像 $p$ 次幂”。

\begin{definition}
对 APF 扩张 $L/K$，令 $E_{L/K}$ 为 $L/K$ 的有限中间扩张组成的有向集。范域定义为
\[
  X_K(L)=\varprojlim_{E\in E_{L/K}} E
\]
其中转移映射为域范数。元素是相容族 $(x_E)$，满足 $N_{E'/E}(x_{E'})=x_E$。
\end{definition}

这里的逆极限符号容易误导：范数不是加法同态，所以不能直接按逐层加法定义环结构。范域的加法要用高层极限：
\[
(x+y)_E=\lim_{E'\supset E}N_{E'/E}(x_{E'}+y_{E'}).
\]
乘法则可以逐层用乘法，因为范数保持乘法：
\[
(xy)_E=x_Ey_E.
\]
APF 条件保证上式中的加法极限存在，并且与所取共尾高层无关。直观上，越往高层，范数对主单位的作用越接近 Frobenius，而在特征 $p$ 中 Frobenius 与加法相容；这就是混合特征塔能在极限中变成特征 $p$ 运算的原因。

\begin{proposition}
若 $L/K$ 是 APF 扩张，则 $X_K(L)$ 是特征 $p$ 的完备离散赋值域，其剩余域与 $L$ 的剩余域自然相关。
\end{proposition}

\begin{proof}
范数在足够深的单位过滤上近似 Frobenius。APF 条件保证这种近似足以定义加法和乘法的极限，并保证非零元素有离散赋值。完备性来自逆极限和各有限层的完备性。
更具体地，对相容族 $x=(x_E)$ 定义
\[
v_X(x)=v_E(x_E)
\]
其中 $E$ 取足够高层并按归一化赋值比较。范数相容性保证这个值不依赖高层选择。若 $x$ 非零，则某一层 $x_E$ 非零，而范数相容迫使所有更低层非零，所以得到离散赋值。特征为 $p$ 的原因可以从整数 $p$ 的像看出：在高层 $E'/E$ 中，$N_{E'/E}(p)$ 的赋值随扩张次数增长而趋向无穷，因此 $p$ 在极限域中等于零。

剩余域方面，单位族模极大理想后只记录各有限层剩余域中的相容元素。若 $L/K$ 全分歧，则剩余域不变，范域的剩余域就是原来的剩余域；若存在剩余域扩张，则范域剩余域对应这些剩余域的直接极限。
\end{proof}

\subsection{\texorpdfstring{$X$}{X}}

\begin{proposition}
有限可分扩张 $M/L$ 与范域 $X_K(M)/X_K(L)$ 的有限可分扩张相对应；在适当 APF 假设下，这是范畴等价。
\end{proposition}

\begin{proof}
有限扩张在足够高的有限层上由 Krasner 型逼近稳定下来。范数相容族记录这些有限层中的元素。可分多项式的根可在高层唯一追踪，因此域扩张、嵌入和 Galois 群都能在范域中恢复。
证明的关键是“近似根唯一”。设 $P(T)$ 是范域上的可分多项式。把系数看成有限层中的相容族，在足够高的中间域 $E$ 里取近似系数 $P_E(T)$。若范域中有一个根 $\alpha$，它在高层的坐标 $\alpha_E$ 会近似满足 $P_E(\alpha_E)=0$。Krasner 引理说明，只要近似足够好，真正根生成的有限扩张不再随 $E$ 改变。于是范域中的有限可分扩张可以从高层有限扩张下降回来。

忠实性也来自同一原则。两个高层嵌入若在范域中给出同一个映射，则它们在足够高层对生成元的像任意接近；可分扩张中不同共轭根之间有正距离，因此足够接近就必须相等。由此得到 Galois 群保持：
\[
\Gal(M/L)\simeq \Gal(X_K(M)/X_K(L)).
\]
\end{proof}

\section{完全化}

\subsection{\texorpdfstring{$R$}{R}}

\begin{definition}
若 $A$ 是特征 $p$ 环，定义
\[
  R(A)=\varprojlim\left(\cdots\xrightarrow{x\mapsto x^p}A\xrightarrow{x\mapsto x^p}A\right).
\]
因此 $R(A)$ 的元素是序列 $(x^{(0)},x^{(1)},\ldots)$，满足 $(x^{(n+1)})^p=x^{(n)}$。
\end{definition}

\begin{proposition}
$R(A)$ 是完全环，即 Frobenius 是自同构。
\end{proposition}

\begin{proof}
Frobenius 把 $(x^{(0)},x^{(1)},\ldots)$ 送到 $((x^{(0)})^p,(x^{(1)})^p,\ldots)$。由于相容关系，这等于把序列左移并补上一个 $p$ 次根；逆映射就是右移。因此 Frobenius 可逆。
\end{proof}

\begin{proposition}
若 $A$ 是 Hausdorff 完备 $p$ 进拓扑环，则 $R(A/pA)$ 可看作 $A$ 的倾斜对象，记录一切相容的 $p$ 幂根系统。对 $x=(x^{(n)})$，选择提升 $\widehat x^{(n)}\in A$ 后，极限
\[
  \lim_{n\to\infty}(\widehat x^{(n)})^{p^n}
\]
给出 Teichmuller 型提升。
\end{proposition}

\subsection{\texorpdfstring{$E$ 和 $A$}{E 和 A}}

在全分歧 $\Z_p$ 扩张 $K_\infty/K$ 中，取
\[
  E_K
\]
为范域，取其 Witt 环并作适当完备化，得到 $A_K$。这些环把特征 $0$ 的 Galois 表示问题转译到特征 $p$ 的 Frobenius 模问题。

\begin{proposition}
对足够高的层 $K_n$，可选择相容素元 $\pi_{K_n}$，使范数和 Frobenius 关系稳定。由此可把 $K_\infty$ 的剩余数据嵌入 $R(\mathcal O_{\widehat{\overline K}}/p)$。
\end{proposition}

\begin{proof}
稳定性来自全分歧 $\Z_p$ 塔的迹估计和范数估计。相邻层素元可选得使范数接近下一层素元的 $p$ 次幂；误差落在越来越深的理想中，因此在 $R$ 的逆极限中消失。
\end{proof}

\begin{theorem}
范域函子把 $K_\infty$ 上的有限可分扩张与 $E_K$ 上的有限可分扩张等价起来，并保持 Galois 群。
\end{theorem}

\begin{proof}
满性由高层逼近：给定 $E_K$ 上有限可分多项式，把系数提升到足够高的 $K_n$，Krasner 引理保证根生成的扩张稳定。忠实性来自相容嵌入在无限多层上的一致性；若两个态射在范域上一致，则它们在足够高的有限层上一致。
\end{proof}

\subsection{\texorpdfstring{$B$ 环大观园}{B 环大观园}}

令 $\widetilde E$ 为特征 $p$ 的完备代数闭赋值域，$\widetilde A^+=W(\mathcal O_{\widetilde E})$，$\widetilde B=\widetilde A^+[1/p]$。通过 Teichmuller 展开，每个元素可写成
\[
  x=\sum_{k\gg-\infty}p^k[x_k].
\]
这些环是 Fontaine 周期环的原料。

\begin{definition}
过收敛环 $\widetilde B^{\dagger,r}$ 由满足增长条件
\[
  w_k(x)+\frac{p-1}{pr}k\ge0
\]
的元素组成，$\widetilde B^\dagger=\bigcup_r\widetilde B^{\dagger,r}$。过收敛元是在某个正半径圆环上有解析控制的元素。
\end{definition}

\begin{definition}
给定区间 $I=[r,s]$，构造环
\[
  \widetilde A^{[r,s]}=\widetilde A^+\left\{\frac p{[\bar\pi]^r},\frac{[\bar\pi]^s}p\right\}
\]
的 $p$ 进完备化，并令 $\widetilde B^{[r,s]}=\widetilde A^{[r,s]}[1/p]$。
\end{definition}

\begin{remark}
这些环按圆环半径组织。区间越小，解析控制越强；区间包含关系反向给出环包含关系。后面 $(\varphi,\Gamma)$-模和过收敛表示会在这些环上定义。
\end{remark}

\section{习题详解}

\begin{exercise}
证明有限可分扩张的局部差别式和全局差别式相容，并推出判别式的乘积分解。
\end{exercise}

\begin{proof}
对素理想 $\mathfrak P|\mathfrak p$，局部化再完备化与有限生成模的对偶相容。差别式由
\[
  \mathfrak D^{-1}_{E/F}=\{x\in E:\operatorname{Tr}_{E/F}(x\mathcal O_E)\subset\mathcal O_F\}
\]
定义。把这个条件局部化到 $\mathfrak p$，只涉及 $\mathfrak P|\mathfrak p$ 的分量，因此
\[
  \mathfrak D_{\widehat{\mathcal O}_{E,\mathfrak P}/\widehat{\mathcal O}_{F,\mathfrak p}}
  =\mathfrak D_{\mathcal O_E/\mathcal O_F}\widehat{\mathcal O}_{E,\mathfrak P}.
\]
对所有 $\mathfrak P$ 相乘得全局差别式分解。判别式是差别式的范，所以再取范数得到判别式公式。
\end{proof}

\begin{exercise}
设整闭整环 $R$ 只有一个素理想 $\mathfrak m$，且 $\mathfrak m$ 有限生成。证明 $\mathfrak m$ 主生成。
\end{exercise}

\begin{proof}
取非零 $a\in\mathfrak m$。因为只有一个非零素理想，$\sqrt{(a)}=\mathfrak m$。有限生成性给出某个 $n$ 使 $\mathfrak m^n\subset(a)$。取最小的 $n$。若 $n>1$，选 $x\in\mathfrak m^{n-1}$ 不在 $(a)$。则 $x/a$ 不在 $R$。整闭一维局部情形中，由理想全在 $\mathfrak m$ 上支撑可推出 $(a):x$ 是非零理想，含 $\mathfrak m$ 的幂，导致 $x/a$ 满足整方程，矛盾。故 $n=1$，$\mathfrak m\subset(a)$，于是 $\mathfrak m=(a)$。
\end{proof}

\begin{exercise}
证明 $W(k)$ 是以完全环 $k$ 为剩余环的严格 $p$ 环，并证明 $W(\F_p)=\Z_p$、$W_n(\F_p)=\Z/p^n\Z$。
\end{exercise}

\begin{proof}
Witt 环由 ghost 多项式定义，$p$ 进完备，且 $p$ 不是零除子。模 $p$ 后只剩第零 Witt 分量，所以剩余环为 $k$。Teichmuller 提升给出代表映射 $[a]=(a,0,0,\ldots)$，满足 $[a]^p=[a^p]$，所以是严格 $p$ 环。若 $k=\F_p$，严格 $p$ 环唯一性说明 $W(\F_p)$ 与 $\Z_p$ 同构；截断长度 $n$ 对应模 $p^n$，得 $W_n(\F_p)=\Z/p^n\Z$。
\end{proof}

\begin{exercise}
证明若 $i\le j$，则 $(G/G_j)_i=G_i/G_j$；若 $i\ge j$，则 $(G/G_j)_i=\{1\}$。
\end{exercise}

\begin{proof}
商扩张对应固定域 $L^{G_j}$。在商群中，元素 $gG_j$ 对整数环模 $\mathfrak p^{i+1}$ 的作用由 $g$ 的作用决定。当 $i\le j$ 时，$G_j$ 已经在该层上平凡，所以商中的第 $i$ 层正是 $G_i/G_j$。当 $i\ge j$ 时，要求在比 $G_j$ 定义更深的层上平凡，只剩单位元余类。
\end{proof}

\begin{exercise}
验证 Hasse 函数在整数区间上的公式
\[
  \varphi(u)=g_0^{-1}(g_1+\cdots+g_m+g_{m+1}(u-m))
\]
以及 $g_0(\varphi(m)+1)=\sum_{i=0}^m g_i$。
\end{exercise}

\begin{proof}
在区间 $m\le u\le m+1$ 上，$G_t=G_{m+1}$ 对 $m<t\le m+1$，所以斜率为 $g_{m+1}/g_0$。从 $0$ 到 $m$ 分段积分：
\[
  \varphi(m)=\sum_{r=0}^{m-1}\frac{g_{r+1}}{g_0}.
\]
再加上最后一段长度 $u-m$ 的贡献，得到第一式。令 $u=m$ 并注意 $g_0$ 对应 $[-1,0]$ 的贡献，得第二式。
\end{proof}

\begin{exercise}
设 $F$ 是局部域，素元为 $\pi$。比较 $\varprojlim\mathcal O/\pi^n\mathcal O$ 和 $\varprojlim\mathcal O/p^n\mathcal O$。
\end{exercise}

\begin{proof}
若 $v_F(p)=e$，则 $p\mathcal O=\pi^e\mathcal O$ 乘以单位。因此 $p^n\mathcal O=\pi^{en}\mathcal O$。两组理想 $\{\pi^n\}$ 和 $\{p^n\}$ 是同一拓扑的共尾邻域基：每个 $\pi^m$ 包含某个 $p^n$，每个 $p^n$ 也是某个 $\pi^{en}$。所以两个逆极限都等于 $\mathcal O$ 的完备化；若 $F$ 已完备，它们都自然同构于 $\mathcal O$。
\end{proof}

\begin{exercise}
对分圆 $\Z_p$ 扩张构造的倾斜域 $\widetilde E$，证明其为特征 $p$ 代数封闭完备赋值域，并计算 $v_E(\varepsilon-1)=p/(p-1)$。
\end{exercise}

\begin{proof}
元素是相容序列 $x=(x^{(0)},x^{(1)},\ldots)$，满足 $(x^{(n+1)})^p=x^{(n)}$。乘法逐坐标定义，加法用极限
\[
  (x+y)^{(n)}=\lim_m(x^{(n+m)}+y^{(n+m)})^{p^m}.
\]
在特征 $p$ 的极限中，Frobenius 成为自同构，所以所得域特征为 $p$。赋值 $v_E(x)=v_p(x^{(0)})$ 与运算相容，完备性来自 $\C_p$ 完备和逐坐标极限。代数封闭性由 $\C_p$ 代数封闭和 Hensel 提升逐层求根得到。

若 $\varepsilon=(1,\zeta_p,\zeta_{p^2},\ldots)$，则 $\varepsilon-1$ 的第 $0$ 坐标按加法定义对应极限，而标准分圆估计给出
\[
  v_p(\zeta_{p^n}-1)=\frac1{p^{n-1}(p-1)}.
\]
归一到倾斜赋值后得到 $v_E(\varepsilon-1)=p/(p-1)$。
\end{proof}

\begin{exercise}
证明完备离散赋值域中元素可用代表集作乘法展开，并说明主单位的生成形式。
\end{exercise}

\begin{proof}
先写 $\alpha=\pi^nu$，其中 $u\in\mathcal O_F^\times$。取代表集 $R$，令 $\theta_0$ 是 $u$ 的剩余类代表，则 $u/\theta_0\in1+\mathfrak p$。对 $1+\mathfrak p$ 中元素，逐步选择 $\theta_i\in R$ 使
\[
  u_i(1+\theta_i\pi^i)^{-1}\in1+\mathfrak p^{i+1}.
\]
无限乘积收敛，得到
\[
  \alpha=\pi^n\theta_0\prod_{i\ge1}(1+\theta_i\pi^i).
\]
唯一性由每一步模 $\mathfrak p^{i+1}$ 的代表唯一性得出。若剩余域完全且特征混合，主单位的进一步分解来自 Artin--Hasse 型指数和映射 $a\mapsto a^p+\bar\theta_0a$ 的余核代表；指数 $pe/(p-1)$ 是 $p$ 次幂映射在线性化中改变层数的临界位置。
\end{proof}

\begin{exercise}
设 $K(X)$ 为有理函数域。证明不可约多项式 $p(X)$ 给出离散赋值，且无穷远处也给出离散赋值。
\end{exercise}

\begin{proof}
对多项式 $f$，令 $v_p(f)$ 为 $p(X)$ 的最高幂次数。唯一分解保证
\[
  v_p(fg)=v_p(f)+v_p(g),\qquad v_p(f+g)\ge\min\{v_p(f),v_p(g)\}.
\]
延拓到分式 $f/q$，定义 $v_p(f/q)=v_p(f)-v_p(q)$，良定性由约分不改变差值得到。这是离散赋值。

无穷远赋值定义为
\[
  v_\infty(f/q)=\deg q-\deg f.
\]
次数满足 $\deg(fg)=\deg f+\deg g$，而 $\deg(f+g)\le\max\{\deg f,\deg g\}$，所以同样给出离散赋值，并且 $v_\infty(1/X)=1$。
\end{proof}

\begin{exercise}
证明 $\F_q(X)$ 的离散赋值由首一不可约多项式和无穷远赋值给出，并描述完备化和剩余域。
\end{exercise}

\begin{proof}
任意在 $\F_q$ 上平凡的离散赋值，对 $\F_q[X]$ 的非零多项式非负，除非它是无穷远型。若 $X$ 赋值负，则 $1/X$ 赋值正，得到无穷远赋值。若 $X$ 赋值非负，赋值环与 $\F_q[X]$ 的交给出一个非零素理想，由 PID 性为 $(\pi)$，其中 $\pi$ 首一不可约。于是赋值等于上一题的 $v_\pi$ 的整数倍；规范化后唯一。

$v_\pi$ 完备化的素元是 $\pi$，剩余域为 $\F_q[X]/(\pi)\simeq\F_{q^{\deg\pi}}$。无穷远完备化用 $T=1/X$ 作素元，剩余域为 $\F_q$。
\end{proof}

\begin{exercise}
证明 $K(X)$ 对 $v_X$ 的完备化为 Laurent 级数域 $K((X))$，赋值环为 $K[[X]]$。
\end{exercise}

\begin{proof}
在 $v_X$ 下，$X$ 是素元。任意有理函数可写成 $X^nu(X)$，其中 $u(0)\ne0$。其在 $X$ 进拓扑下展开成 Laurent 级数。Cauchy 序列的每个系数最终稳定，极限为
\[
  \sum_{m\ge N}a_mX^m.
\]
这正是 $K((X))$。赋值非负的元素为没有负幂项的级数，即 $K[[X]]$。
\end{proof}

\begin{exercise}
研究 $F\{\{X\}\}$ 的两个赋值，并证明 $K((X))\{\{Y\}\}=K((Y))((X))$。
\end{exercise}

\begin{proof}
元素为双向级数 $\sum a_nX^n$，其中系数赋值有下界且当 $n\to-\infty$ 时趋于 $+\infty$。取
\[
  v\left(\sum a_nX^n\right)=\min_n v_F(a_n).
\]
乘法的最小赋值项不会被无限抵消，因为达到最小值的负向项只有有限多个，所以这是离散赋值；剩余域保留最小赋值系数形成的 Laurent 级数，故为 $\kappa_F((X))$。

若取
\[
  v^*\left(\sum a_nX^n\right)=\min_n\{v_F(a_n),n\},
\]
非负条件等价于所有 $a_n\in\mathcal O_F$，所以赋值环是 $\mathcal O_F\{\{X\}\}$，模极大理想后只有常数剩余，剩余域为 $\kappa_F$。

最后，$K((X))\{\{Y\}\}$ 的元素是 $\sum b_mY^m$，其中 $b_m\in K((X))$，负向 $m$ 时 $X$-赋值趋于 $+\infty$。重新按 $X$ 的幂收集系数，得到系数在 $K((Y))$ 的 Laurent 级数；条件正好保证两边的收敛要求一致。
\end{proof}

\begin{exercise}
证明题中给出的 $\Q_p((X_2))\cdots((X_n))$、$\Q_p\{\{X_2\}\}\cdots\{\{X_n\}\}$、$\F_q((X_1))\cdots((X_n))$ 是 $n$ 维局部域，并说明一般 $n$ 维局部域的标准形。
\end{exercise}

\begin{proof}
逐次取 Laurent 级数会把剩余域降一维：$K((X))$ 的剩余域是 $K$。逐次取 $\{\{X\}\}$ 在混合特征中也给出完备离散赋值域，并使剩余域降到上一层。因此列出的域逐层剩余，最后到达有限域。

一般 $n$ 维局部域从最后一层往上看。等特征层由 Cohen 结构定理写成 Laurent 级数；混合特征层用 Cohen 环和有限扩张描述。迭代后得到它是某个
\[
  k\{\{X_1\}\}\cdots\{\{X_m\}\}((X_{m+2}))\cdots((X_n))
\]
的有限扩张。
\end{proof}
 \chapter{Galois 表示}

\section{章节导读}

Galois 表示把绝对 Galois 群的作用线性化。若 $G_K=\Gal(\overline K/K)$，一个 $p$ 进表示就是 $G_K$ 在有限维 $\Q_p$-向量空间上的连续线性作用。这样做的动机是：数域或局部域的很多算术信息本来藏在无限非交换群 $G_K$ 中，线性表示把它们转化成矩阵、特征多项式、过滤、Frobenius 和上同调。

本章内容可以分成四层。第一层是半线性代数：晶体、同晶体和 Dieudonne--Manin 分解解释 Frobenius 线性代数。第二层是 Tate--Sen 方法：通过 $\GL_n(\C_p)$ 的非交换上同调，把连续 cocycle 化为可控制的矩阵问题。第三层是 $\varphi$-模和 $\varphi$-$\Gamma$ 模：把 Galois 表示等价地翻译为带 Frobenius 和 $\Gamma$ 作用的模。第四层是周期环：$B_{\mathrm{dR}},B_{\mathrm{cris}},B_{\mathrm{st}}$ 提供比较范畴，用来定义 Hodge--Tate、de Rham、crystalline 和 semistable 表示。

\section{晶体}

\subsection{半线性映射}

\begin{definition}
设 $F$ 是域，$\sigma:F\to F$ 是自同构。$F$-向量空间 $V$ 上的 $\sigma$-半线性映射是加法映射 $\Phi:V\to V$，满足
\[
  \Phi(av)=\sigma(a)\Phi(v)\qquad(a\in F,\ v\in V).
\]
若 $\Phi$ 可逆，则称为 $\sigma$-半线性自同构。
\end{definition}

半线性出现在 Frobenius 理论中，因为 Frobenius 对系数不是恒等，而是先把系数做 $p$ 次幂。把 $\Phi$ 当普通线性映射会丢掉最重要的算术信息。最简单的检查是 $V=F$，若 $\Phi(x)=\sigma(x)$，则 $\Phi(ax)=\sigma(a)\sigma(x)$，这里系数 $a$ 不能从 $\Phi$ 外面原样拿出来。后面所有带 Frobenius 的模都在处理这件事：向量本身被移动，坐标系中的数也被 Frobenius 移动。

为了把半线性算子纳入代数语言，需要把普通多项式环改成扭多项式环。普通多项式环满足 $Ta=aT$，对应的是线性算子；扭多项式环满足 $Ta=\sigma(a)T$，对应的是先让算子穿过系数，系数随即被 $\sigma$ 改写。这样一来，一个向量空间带一个 $\sigma$-半线性算子，等价于这个向量空间成为 $F\langle T\rangle$-模。这个翻译很重要，因为后面分类同晶体时，可以把半线性代数问题变成扭多项式环上的模分解问题。

\begin{definition}
扭多项式环 $F\langle T\rangle$ 的元素是有限和 $\sum a_iT^i$，加法逐项相加，乘法由
\[
  T a=\sigma(a)T
\]
决定。给出 $F\langle T\rangle$-模结构等价于给出一个 $\sigma$-半线性算子 $T$。
\end{definition}

\begin{lemma}
$F\langle T\rangle$ 是左 Euclid 整环。每个左理想由一个首一元素生成。
\end{lemma}

\begin{proof}
对 $f,h\in F\langle T\rangle$，若 $h$ 的最高项为 $bT^m$，则可用 $f$ 的最高项 $aT^n$ 减去
\[
  a\,\sigma^{n-m}(b)^{-1}T^{n-m}h
\]
消去最高次数。重复有限次得到左除法。左理想中取次数最小的首一元素 $h$，任意 $f$ 除以 $h$ 的余数仍在左理想中，次数更小，只能为零。

这里要注意左右方向。因为乘法规则是 $T^q b=\sigma^q(b)T^q$，所以左乘 $cT^{n-m}$ 到 $bT^m$ 的最高项为 $c\sigma^{n-m}(b)T^n$。令 $c=a\sigma^{n-m}(b)^{-1}$，最高项正好变成 $aT^n$。每一步次数下降，有限次后停止。若 $I$ 是左理想，取其中次数最小的非零元并除以最高系数使它首一，记为 $h$。对任意 $f\in I$ 做左除法 $f=qh+r$，由于 $r=f-qh$ 仍在 $I$ 中，而 $\deg r<\deg h$，最小性迫使 $r=0$。首一生成元的唯一性也来自同一个论证：若 $h,h'$ 都首一且生成同一左理想，则 $h$ 被 $h'$ 整除且 $h'$ 被 $h$ 整除，次数相同，余数为零，首一最高项迫使二者相等。
\end{proof}

\subsection{晶体}

\begin{definition}
设 $k$ 是特征 $p$ 完全域，$W(k)$ 是 Witt 环，$\sigma$ 是 Witt Frobenius。
\begin{enumerate}[label=(\arabic*)]
\item $k$ 上的晶体是有限秩自由 $W(k)$-模 $M$ 加 $\sigma$-半线性单射 $\Phi:M\to M$。
\item $k$ 上的同晶体是有限维 $\operatorname{Frac}W(k)$-向量空间 $V$ 加 $\sigma$-半线性同构 $\Phi:V\to V$。
\end{enumerate}
\end{definition}

晶体是整数结构，同晶体是把 $p$ 反掉后的有理结构。晶体中要求 $M$ 为有限自由 $W(k)$-模，是因为它要保留 $p$ 进整数格；同晶体中把 $p$ 变为可逆，则只保留斜率、分解和 Frobenius 的有理信息。一个晶体 $(M,\Phi)$ 给出同晶体 $(M[1/p],\Phi)$，但反过来一个同晶体未必有唯一晶体格。讲这一节时可以把晶体看成“带 Frobenius 的格”，把同晶体看成“不计入格之后的向量空间”。

态射也要和 Frobenius 相容。若 $f:(M,\Phi)\to(M',\Phi')$ 是 $W(k)$-线性映射，条件是 $\Phi' f=f\Phi$。这表示 $f$ 不只是底层模同态，还尊重 Frobenius 动力系统。同源态射要求 $f[1/p]$ 为同构，意思是允许有限 $p$ 阶误差；这种误差在同晶体范畴里消失。

\subsection{标准同晶体}

\begin{definition}
设 $\lambda=s/r\in\Q_{\ge0}$，其中 $r>0$ 且 $(r,s)=1$。定义
\[
  E^\lambda=K\langle T\rangle/(T^r-p^s),
  \qquad K=\operatorname{Frac}W(k),
\]
并让 $\Phi$ 由左乘 $T$ 给出。称它为斜率 $\lambda$ 的标准同晶体。
\end{definition}

若取基 $e_i=T^{i-1}$，则 $\Phi(e_i)=e_{i+1}$，最后 $\Phi(e_r)=p^se_1$。因此 $\Phi^r=p^s$，斜率就是“$r$ 次 Frobenius 后乘以 $p^s$”的平均 $p$ 阶。更具体地说，$\Phi$ 每走一步把基向量往前推一格，走完 $r$ 步回到原来的基向量，同时多出因子 $p^s$；因此每一步平均贡献 $s/r$ 个 $p$ 阶。

这个模型对象 $E^\lambda$ 是同晶体分类中的“不可再拆的积木”。当 $\lambda=0$ 时，$\Phi^r=1$，对象像无分歧方向；当 $\lambda>0$ 时，$\Phi$ 的迭代带来 $p$ 的正幂，表示存在收缩方向；当 $\lambda<0$ 时用对偶定义，表示膨胀方向。对初学者而言，斜率可以先理解成 Frobenius 迭代后的平均赋值增长率。

还需要记住三个基本运算。张量积满足斜率相加，因为 $(\Phi_1\otimes\Phi_2)^N$ 的 $p$ 阶来自两个因子的 $p$ 阶相加。对偶把斜率变成相反数，因为若 $\Phi^r$ 给出因子 $p^s$，对偶上对应因子为 $p^{-s}$。Tate 扭转 $(n)$ 把 Frobenius 改成 $p^{-n}\Phi$，所以斜率减去 $n$。这些规则会在过滤 $\varphi$-模和半稳定表示中反复出现。

\subsection{同晶体范畴半单性}

\begin{theorem}
若 $k$ 代数封闭，则每个同晶体可唯一分解为有限直和
\[
  E\simeq \bigoplus_{\lambda\in\Q} (E^\lambda)^{m_\lambda}.
\]
每个 $E^\lambda$ 是单对象，且同斜率扩张分裂。
\end{theorem}

\begin{proof}
这是 Dieudonne--Manin 分类。证明分成四步。

第一步，把同晶体看成 $K\langle T\rangle$-模。这里 $T$ 的作用就是 Frobenius $\Phi$。因为 $K\langle T\rangle$ 有左 Euclid 除法，有限生成扭多项式模可以用一个扭多项式 $P(T)$ 描述出循环商。也就是说，问题变成研究 $K\langle T\rangle/(P)$ 中 Frobenius 的形状。

第二步，从 $P$ 的系数读出斜率。把 $P$ 乘上合适的 $p$ 幂，可以把系数放进 $W(k)$。取 Newton 多边形第一段的斜率 $s/r$，引理中的分解把 $P$ 在扩张 $W(k)(p^{1/r})$ 上抽出一个因子 $T-p^{s/r}$。这一步的意义是：从任意非零同晶体中，总能找到一个映到标准对象 $E^\lambda$ 的非零态射。

第三步，证明标准对象 $E^\lambda$ 是单对象。若有非零真子对象 $E\subset E^\lambda$，则商 $E^\lambda/E$ 仍有非零映射到某个 $E^\mu$。当 $\mu\ne\lambda$ 时，斜率增长率不同，$\Hom(E^\lambda,E^\mu)=0$；当 $\mu=\lambda$ 时，端环是除代数，非零自同态必可逆，合成态射不能通过真商产生。这排除了非零真子对象。

第四步，证明扩张分裂。考虑短正合列 $0\to E^\lambda\to E\to E^{\lambda'}\to0$。取 $E$ 中一个提升元 $x$，希望调整为满足 $\Phi^r x=p^s x$ 的元。误差落在左端 $E^\lambda$ 中；需要解方程 $(\Phi^r-p^s)y=$ 误差。若斜率不相等，赋值增长使算子可逆；若斜率相等，则归结为 Artin--Schreier 型方程 $z^{p^a}-z=c$，代数封闭性保证逐级提升可解。于是每个扩张都分裂。

四步合在一起得到半单分解。唯一性由单对象的端环为除代数和不同斜率之间没有非零 Hom 推出：任何分解中斜率 $\lambda$ 的同构类和重数都由对象本身决定。
\end{proof}

\subsection{联络}

若底环带有微分，晶体还常配备联络 $\nabla$。直觉上，$\Phi$ 描述 Frobenius 方向，$\nabla$ 描述无穷小移动方向。二者相容时，才像几何上同调带来的对象。后面 $p$ 进微分方程和 de Rham 比较会不断出现这种“Frobenius 加联络”的组合。

设 $R$ 是 $k$-代数。联络是满足 Leibniz 公式的映射
\[
  \nabla:M\longrightarrow \Omega^1_{R/k}\otimes_R M,\qquad
  \nabla(xm)=dx\otimes m+x\nabla(m).
\]
若再给定一个求导 $\delta:R\to R$，则由 $\delta$ 对应的 $R$-线性映射 $\Omega^1_{R/k}\to R$ 可得到微分算子 $D:M\to M$，满足 $D(xm)=xD(m)+\delta(x)m$。这个公式告诉我们：联络不是普通线性映射，它精确记录“系数变化”给向量带来的额外项。与半线性 Frobenius 对比，二者都在提醒我们，系数环本身也在动。

\section{CK}

设 $K$ 是完备离散赋值域，特征为 $0$，剩余域特征为 $p$ 且完全。记 $\overline K$ 为代数闭包，$C_K$ 为 $\overline K$ 的完备化。Ax--Tate 型定理说明：完备化没有制造新的 Galois 固定元。若 $H\subset G_K$ 是闭子群，则 $C_K^H$ 正好是 $\overline K^H$ 的完备化。

\begin{lemma}
设 $L/K$ 是代数扩张，$x\in\overline K$。记
\[
  d_L(x)=\min_{\sigma} v(\sigma x-x),
\]
其中 $\sigma$ 遍历把 $L(x)$ 嵌入 $\overline K$ 且固定 $L$ 的嵌入。则存在 $y\in L$ 使
\[
  v(x-y)>d_L(x)-\frac{p}{(p-1)^2}v(p).
\]
\end{lemma}

\begin{proof}
这是 Ax 估计。证明的骨架是对 $n=[L(x):L]$ 归纳。设 $f$ 是 $x$ 在 $L$ 上的首一最小多项式，令 $g(X)=f(X+x)$。则 $g$ 的根是各个 $\sigma x-x$，这些根的赋值都不小于 $d_L(x)$。对 $g$ 取某个阶数 $m$ 的导数，导数的根可以用引理中的对称多项式估计控制：至少有一个根 $z$ 满足
\[
  v(z)\ge d_L(x)-\frac{1}{n-m}v\binom nm .
\]
于是 $u=x+z$ 比 $x$ 更接近 $L$，并且 $u$ 的次数比 $x$ 小。对 $u$ 使用归纳假设，再用非 Archimede 三角不等式比较 $x,u,y$，得到所需的界。常数 $\frac{p}{(p-1)^2}v(p)$ 来自把所有可能的 $p$ 幂次数贡献求和；它是一个统一误差界，关键作用是和逼近精度 $n$ 无关。
\end{proof}

\begin{proposition}
若 $H$ 是 $G_K$ 的闭子群，则
\[
  C_K^H=\widehat{\overline K^H}.
\]
特别地，$C_K^{G_K}=K$。
\end{proposition}

\begin{proof}
右端自然落在左端。反向取 $x\in C_K^H$。按完备化定义，存在 $x_n\in\overline K$ 使 $v(x-x_n)\ge n$。对任意 $\sigma\in H$，由于 $\sigma x=x$，有
\[
  v(\sigma x_n-x_n)\ge
  \min\{v(\sigma x_n-\sigma x),v(x-x_n)\}\ge n.
\]
因此 $d_{\overline K^H}(x_n)\ge n$。由 Ax 估计，存在 $y_n\in\overline K^H$ 使
\[
  v(x_n-y_n)>n-\frac{p}{(p-1)^2}v(p).
\]
于是 $v(x-y_n)$ 随 $n$ 趋向无穷，故 $x$ 是 $\overline K^H$ 中元素的极限，即 $x\in\widehat{\overline K^H}$。特别地，当 $H=G_K$ 时，$\overline K^{G_K}=K$，而 $K$ 已完备，所以 $C_K^{G_K}=K$。
\end{proof}

\section{非交换一阶上同调}

\begin{definition}
设群 $G$ 作用在群 $A$ 上。非交换 $1$-上闭链是映射 $c:G\to A$，满足
\[
  c_{\sigma\tau}=c_\sigma\,\sigma(c_\tau).
\]
两个上闭链 $c,c'$ 等价，如果存在 $a\in A$ 使
\[
  c'_\sigma=a^{-1}c_\sigma\,\sigma(a).
\]
等价类集合记为 $H^1(G,A)$。
\end{definition}

当 $A$ 非交换时，$H^1(G,A)$ 通常只是带基点的集合，不是群。它分类扭曲形式、下降数据和半线性作用。上闭链条件
\[
  c_{\sigma\tau}=c_\sigma\,\sigma(c_\tau)
\]
可以从“先做 $\tau$ 再做 $\sigma$”读出：$\tau$ 的矩阵先出现，但再被 $\sigma$ 作用到系数；随后乘上 $\sigma$ 的矩阵。等价关系则是换基公式。若把基由列向量矩阵 $P$ 改成新基，矩阵会变成 $P^{-1}c_\sigma\sigma(P)$。所以 $H^1(G,\GL_n(B))$ 不是抽象摆设，而是在记录带半线性 $G$ 作用的自由 $B$-模有哪些下降形式。

\begin{proposition}
短正合列 $1\to A\to B\to C\to1$ 诱导正合列
\[
  1\to A^G\to B^G\to C^G\to H^1(G,A)\to H^1(G,B)\to H^1(G,C).
\]
\end{proposition}

\begin{proof}
若 $c\in C^G$，取提升 $b\in B$，则
\[
  b^{-1}\sigma(b)\in A
\]
定义 $A$ 值 cocycle。它确实取值于 $A$，因为 $\beta(b)$ 被 $G$ 固定，所以 $\beta(b^{-1}\sigma b)=\beta(b)^{-1}\sigma\beta(b)=1$。上闭链条件来自
\[
  b^{-1}\sigma\tau(b)=b^{-1}\sigma(b)\,\sigma\bigl(b^{-1}\tau(b)\bigr).
\]
若换提升为 $ba$，其中 $a\in A$，得到的上闭链变为 $a^{-1}(b^{-1}\sigma b)\sigma(a)$，正是同一个上同调类。其余箭头由包含和投射给出。正合性逐项检查：某个类映到平凡类，等价于它能从前一项提升，或能被一次换基共轭为平凡上闭链。
\end{proof}

\begin{proposition}
若 $L/K$ 是有限 Galois 扩张，则
\[
  H^1(\Gal(L/K),\GL_n(L))=1.
\]
\end{proposition}

\begin{proof}
这是矩阵形式的 Hilbert 90。给定 cocycle $U_\sigma$，让 $G$ 半线性作用在 $L^n$ 上：
\[
  \sigma * v=U_\sigma\,\sigma(v).
\]
上闭链条件正是 $(\sigma\tau)*v=\sigma*(\tau*v)$ 的验证。下降定理说明这个半线性空间有 $K$-基。把该 $K$-基写成 $L^n$ 中列向量组成矩阵 $M$，这些列在新作用下固定，于是
\[
  U_\sigma\,\sigma(M)=M,
\]
即 $U_\sigma=M\sigma(M)^{-1}$，cocycle 为平凡类。反过来，若一个 cocycle 有这种形式，则换基矩阵 $M$ 把它变成平凡 cocycle。这说明 $H^1=1$ 的含义就是“所有半线性自由模都可以下降到固定域上的普通向量空间”。
\end{proof}

\begin{proposition}
设拓扑群 $G$ 连续作用在交换拓扑环 $B$ 上。则 $H^1(G,\GL_n(B))$ 与秩 $n$ 自由 $B$-模上的连续半线性 $G$ 作用的同构类一一对应。
\end{proposition}

\begin{proof}
从半线性模 $M$ 出发，取一组 $B$-基 $e_1,\ldots,e_n$。写
\[
  s(e_j)=\sum_i a_{ij,s}e_i,\qquad A_s=(a_{ij,s})\in\GL_n(B).
\]
因为 $s(t(e_j))=(st)(e_j)$，比较基下系数得 $A_{st}=A_s\,s(A_t)$，所以 $A_\bullet$ 是上闭链。若换基矩阵为 $P$，新矩阵为 $P^{-1}A_s\,s(P)$，上同调类不变。

反过来，给定上闭链 $A_s$，在 $B^n$ 上定义 $s*v=A_s\,s(v)$。上闭链条件保证这确实是群作用，半线性来自 $s(bv)=s(b)s(v)$。两个等价上闭链给出的半线性模由换基矩阵同构。两种构造互逆，得到一一对应。
\end{proof}

\section{在 \texorpdfstring{$\GL_n(\C_p)$}{GLn(Cp)} 的上同调}

\subsection{Tate 分解}

Tate--Sen 方法的目标是控制
\[
  H^1(G,\GL_n(\C_p)).
\]
核心思想是先在一个大完备域中分解，再逐步把系数降到较小的域。这里的困难有两层：一是 $\C_p$ 太大，直接用 Hilbert 90 不适用；二是 $G$ 是无限拓扑群，上闭链需要连续控制。Tate--Sen 方法把这个问题拆成三步：先用规范迹给出 $\widehat F_\infty=F_n\oplus X_n$ 的分解，再用这个分解解近似的共轭方程，最后把在完备化中得到的矩阵降回 $F_\infty$。

\begin{proposition}
设 $F_\infty/F$ 是全分歧 $\Z_p$ 扩张的完备化背景。存在常数，使 $\widehat F_\infty$ 可分解为有限层部分加上一个被 $\gamma-1$ 良好控制的补空间。
\end{proposition}

\begin{proof}
第十章的迹估计给出：有限层的迹映射在深层单位球上几乎满射。利用这些迹投影构造有限层近似，再把误差放入补空间。算子 $\gamma-1$ 在补空间上有下界，因此可解线性方程。

具体地，对 $x\in F_\infty$ 定义规范迹 $R_n(x)$；若 $x$ 已在某个有限层 $F_{m+n}$ 中，则 $R_n(x)=p^{-m}\operatorname{Tr}_{F_{m+n}/F_n}(x)$，可迁性保证定义与 $m$ 无关。连续性把 $R_n$ 延拓到 $\widehat F_\infty$。令 $X_n=\ker R_n$，则 $R_n^2=R_n$ 给出直和分解 $\widehat F_\infty=F_n\oplus X_n$。

关键估计为
\[
  v(x-R_nx)\ge v(\gamma_n x-x)-d.
\]
若 $y=(\gamma_n-1)x$ 且 $x\in X_n$，则 $R_nx=0$，于是 $v(x)\ge v(y)-d$。这说明 $(\gamma_n-1)^{-1}$ 在 $X_n$ 上连续，且不会让赋值损失超过固定常数 $d$。后面每次解线性化方程时，正是用这个逆算子控制误差。
\end{proof}

\subsection{殆平展下降}

\begin{proposition}
若连续 cocycle $U:H_0\to\GL_n(\C_p)$ 足够接近恒等，即 $v(U_s-1)$ 有统一正下界，则它在更小开子群上为平凡类。
\end{proposition}

\begin{proof}
要找 $M$ 使
\[
  U_s=M^{-1}s(M).
\]
从 $M=1$ 开始，用线性化方程 $(s-1)X\approx U_s-1$ 修正。Tate 分解保证该方程可解，且每次修正后误差赋值提高。完备性使无限乘积收敛到所需 $M$。

源书中的一次修正可以这样读。先取开正规子群 $H_1\subset H_0$，使 $s\in H_1$ 时 $U_s-1$ 的赋值更高。由殆平展迹条件选 $\alpha\in C_p^{H_1}$，使 $\sum_{s\in H_0/H_1}s(\alpha)=1$，并且 $\alpha$ 的赋值有下界。对代表集 $S$ 定义
\[
  M_S=\sum_{s\in S}s(\alpha)U_s.
\]
因为 $\sum s(\alpha)=1$，矩阵 $M_S$ 接近 $1$，从而可逆。若换代表集，只会产生高赋值误差；由 cocycle 关系可得 $U_t\,t(M_S)$ 与 $M_S$ 相差更小。因此用 $M_S$ 共轭后，新的 cocycle 比原来更接近 $1$。

重复这个过程，误差下界从 $a$ 变成 $a+1,a+2,\ldots$。修正矩阵 $M_1,M_2,\ldots$ 满足 $M_i-1$ 的赋值趋向无穷，所以乘积收敛。极限矩阵把 cocycle 共轭为平凡 cocycle。
\end{proof}

\subsection{退完备化}

\begin{proposition}
若某个 cocycle 可在完备域 $\widehat F_\infty$ 中平凡化，并且满足足够强的估计，则可选择平凡化矩阵落在非完备并域 $F_\infty$ 中。
\end{proposition}

\begin{proof}
把矩阵按有限层逼近。若两个有限层近似在完备化中给出同一 cocycle，差矩阵被 $\gamma-1$ 控制。Tate 分解的唯一性估计迫使差异落在更高精度中。逐步提高精度，得到来自 $F_\infty$ 的矩阵。

更细地说，设 $U=1+U_1+U_2$，其中 $U_1$ 已在有限层 $F_r$，$U_2$ 是高赋值误差。把 $U_2$ 分解为
\[
  U_2=R_r(U_2)+(\gamma_r V-V),
\]
其中 $R_r(U_2)$ 落在 $F_r$，而 $V$ 的赋值由 Tate 分解控制。取 $M=1+V$，展开
\[
  M^{-1}U\gamma_r(M)
\]
时，一次项中的 $\gamma_r V-V$ 正好消去 $U_2$ 的非有限层部分；剩下的新误差至少提高一个固定正数 $\delta$。无限迭代后，非有限层误差趋于 $0$，从而得到有限层系数的共轭形式。

另一个必要步骤是唯一性判别。若 $B\in\GL_n(\widehat F_\infty)$ 满足 $\gamma_r B=V_1BV_2$，且 $V_1,V_2$ 都足够接近 $1$ 并落在某个有限层 $F_i$，则 $B-R_i(B)$ 落在 $X_i$。对它应用 $\gamma_r-1$ 的下界估计，得到除非 $B-R_i(B)=0$，否则赋值会严格大于自身赋值加常数，这是矛盾。因此 $B\in\GL_n(F_i)$。
\end{proof}

\subsection{Sen 定理}

\begin{theorem}
设 $H=\Gal(\overline F/F_\infty)$。则
\[
  H^1(H,\GL_n(\C_p))=1,
\]
并且膨胀映射给出
\[
  H^1(\Gamma,\GL_n(\widehat F_\infty))
  \simeq H^1(G_F,\GL_n(\C_p)).
\]
\end{theorem}

\begin{proof}
第一式由殆平展下降：任意连续 cocycle 在足够小开子群上接近恒等，因而局部平凡。设 $U\in Z^1(H,\GL_n(\C_p))$。连续性给出开正规子群 $H_0$，使 $U_s$ 在 $H_0$ 上任意接近 $1$；殆平展下降说明 $U|_{H_0}$ 是平凡类。膨胀--限制序列把问题降到有限群 $H/H_0$ 的上同调，而 Hilbert--Noether 定理给出
\[
  H^1(H/H_0,\GL_n(C_p^{H_0}))=1.
\]
于是 $H^1(H,\GL_n(\C_p))=1$。

第二式来自正合列
\[
  1\to H\to G_F\to\Gamma\to1
\]
以及第一式给出的消失，同时 $C_p^H=\widehat F_\infty$。于是任何 $G_F$ 上的 $\C_p$ 系数 cocycle 都唯一来自 $\Gamma$ 上的 $\widehat F_\infty$ 系数 cocycle。再加上退完备化，包含 $\GL_n(F_\infty)\to\GL_n(\widehat F_\infty)$ 在 $H^1(\Gamma,-)$ 上给出双射，最终得到
\[
  H^1(\Gamma,\GL_n(F_\infty))
  \simeq H^1(G_F,\GL_n(\C_p)).
\]
这就是 Sen 定理在本节中的形式：连续半线性 $\C_p$-对象可以降到有限层并域 $F_\infty$。
\end{proof}

\subsection{过收敛表示}

\begin{theorem}
每个 $G_K$ 的 $p$ 进表示都是过收敛表示。
\end{theorem}

\begin{proof}
把表示看作 $Q_p$ 值 cocycle。若 $V$ 有基 $e_1,\ldots,e_n$，则 $g(e_j)=\sum_i U_{ij,g}e_i$ 给出 $U_g\in\GL_n(\Q_p)$，并且因为 $\Q_p$ 上的系数被 $G_K$ 固定，有 $U_{gh}=U_gU_h=U_g\,g(U_h)$。通过 Fontaine 环嵌入到 $\widetilde B^\dagger$ 的上同调中，Tate--Sen 方法把 cocycle 共轭到过收敛子环上。

具体目标是找 $M\in\GL_n(B^\dagger)$，使 $M^{-1}U_g\,g(M)$ 在 $H_K$ 上平凡。命题 11.25 提供过收敛环与大环之间的上同调双射，命题 11.26 保证若一个共轭矩阵在大环中把 cocycle 降到了过收敛系数，则这个矩阵本身已经来自过收敛层，命题 11.27 把这两点应用到 $\Q_p$ 值 cocycle。于是
\[
  D^\dagger(V)=(B^\dagger\otimes_{\Q_p}V)^{H_K}
\]
有 $B_K^\dagger$-基，且维数等于 $\dim_{\Q_p}V$。这就是过收敛性。
\end{proof}

\section{\texorpdfstring{$\varphi$ 模}{phi 模}}

\subsection{表示}

设 $E$ 是特征 $p$ 的域，$E^{\mathrm{sep}}$ 为可分闭包，$G_E=\Gal(E^{\mathrm{sep}}/E)$。本节处理三种表示：有限维 $\F_p$ 表示、有限生成 $\Z_p$ 表示、有限维 $\Q_p$ 表示。连续性的含义随系数而变：$\F_p$ 空间取离散拓扑，$\Z_p$ 模取 $p$ 进拓扑，$\Q_p$ 空间取由任意基给出的 $p$ 进拓扑。

为什么特征 $p$ 情形能用 Frobenius 模描述 Galois 表示？原因是 $E^{\mathrm{sep}}$ 上的 Frobenius $x\mapsto x^p$ 与 $G_E$ 作用交换，并且 Artin--Schreier 方程 $x^p-x=a$ 在可分闭包中可解。于是“取 $G_E$ 不变量”和“取 $\varphi=1$ 不变量”成为互相配合的两种下降操作。

\subsection{模}

\begin{definition}
设 $A$ 是带 Frobenius $\varphi$ 的环。一个 $\varphi$-模是有限生成 $A$-模 $M$ 加半线性映射 $\Phi:M\to M$。若线性化
\[
  A\otimes_{\varphi,A}M\longrightarrow M
\]
为同构，并且存在 Frobenius 稳定格，则称为 etale $\varphi$-模。
\end{definition}

在 $\F_p$ 层，$A=E$，对象是有限维 $E$-向量空间 $M$ 加 $\varphi_E$-半线性算子 $\Phi$，并要求线性化同构。这个平展条件排除 Frobenius 降秩的情况。在线性代数上，它表示 Frobenius 没有丢失维数；在几何直觉上，它对应有限平展对象而不是带有非平展 nilpotent 方向的对象。

在 $\Z_p$ 层，取 $E$ 的 Cohen 环 $\mathcal O_E$，对象是有限生成 $\mathcal O_E$-模上的平展 $\varphi$-模。在 $\Q_p$ 层，则把 $p$ 反掉，并保留一个满足平展条件的 $\mathcal O_E$ 格。三层之间的关系是：先模 $p$ 看 $\F_p$ 层，再以逆极限恢复 $\Z_p$ 层，最后张量 $\Q_p$ 得到有理层。

\subsection{等价}

\begin{theorem}
存在范畴等价：$G_E$ 的有限 $\Z_p$-表示与适当 Cohen 环 $A_E$ 上的 etale $\varphi$-模等价。把 $p$ 反掉后，得到 $\Q_p$-表示与 etale $\varphi$-模的等价。
\end{theorem}

\begin{proof}
从表示 $V$ 出发，定义
\[
  D(V)=(A^{\mathrm{sep}}\otimes_{\Z_p}V)^{G_E}.
\]
Frobenius 来自 $A^{\mathrm{sep}}$。反向从 $M$ 出发，取
\[
  V(M)=(A^{\mathrm{sep}}\otimes_A M)^{\varphi=1}.
\]
Hilbert 90 型消失保证两个构造互逆：扩张到 $A^{\mathrm{sep}}$ 后两边都变成平凡对象，下降数据正好是 Galois 作用。

先看 $\F_p$ 情形。对 $V\in\operatorname{Rep}_{\F_p}(G_E)$，$E^{\mathrm{sep}}\otimes_{\F_p}V$ 是 $E^{\mathrm{sep}}$ 上的半线性 $G_E$-模。Hilbert--Noether 定理给出
\[
  E^{\mathrm{sep}}\otimes_E D_E(V)\simeq E^{\mathrm{sep}}\otimes_{\F_p}V,
\]
所以 $\dim_E D_E(V)=\dim_{\F_p}V$。Frobenius 线性化为同构，是因为在右边它就是 $E^{\mathrm{sep}}$ 上的 Frobenius 与 $V$ 上恒等的张量，扩张后不改变维数。

从 $M$ 反向构造时，$E^{\mathrm{sep}}\otimes_E M$ 上有 Frobenius $\varphi\otimes\Phi$。平展条件使这个 Frobenius 是可逆的半线性算子。方程 $\varphi(x)=x$ 的解空间是有限维 $\F_p$-空间；Artin--Schreier 可解性保证扩张回 $E^{\mathrm{sep}}$ 后恢复整个模。于是 $V_E(D_E(V))\simeq V$，而 $D_E(V_E(M))\simeq M$。

$\Z_p$ 情形把上述结论应用到每个 $M/p^nM$，再用逆极限恢复。关键检查是 $\varphi-1$ 在可分闭包的 Cohen 环上有足够多解，因此取 $\varphi=1$ 不变量对短正合列保持正合。$\Q_p$ 情形则由选择稳定格归约到 $\Z_p$ 情形，再把 $p$ 反掉。
\end{proof}

\subsection{说明}

初学者要把这个等价看成“把 Galois 作用换成 Frobenius 作用”。在特征 $p$ 的域上，绝对 Galois 群的信息可由 Frobenius 模捕捉；这就是后面 $\varphi$-$\Gamma$ 模的基础。这里的 $D(V)$ 是把表示扩张到大环后取 Galois 不变量，$V(M)$ 是把模扩张到大环后取 Frobenius 固定元。两边看起来不同，但大环同时带有 Galois 作用和 Frobenius 作用，且二者交换，所以它成了翻译器。

\section{\texorpdfstring{$\varphi$-$\Gamma$ 模}{phi-Gamma 模}}

\begin{definition}
设 $A_K$ 是与局部域 $K$ 相关的 Cohen--Fontaine 环，带 Frobenius $\varphi$，并有 $\Gamma_K=\Gal(K_\infty/K)$ 作用。etale $\varphi$-$\Gamma$ 模是 etale $\varphi$-模 $M$，再配备连续半线性 $\Gamma_K$ 作用，并满足
\[
  \gamma\Phi=\Phi\gamma.
\]
\end{definition}

$\varphi$-$\Gamma$ 模的动机是把 $G_K$ 分成两层。取一个深的无限扩张 $K_\infty/K$，令 $H_K=G_{K_\infty}$，$\Gamma_K=G_K/H_K$。限制到 $H_K$ 后，通过特征 $p$ 范域和 $\varphi$ 模等价，可把表示变成 $\varphi$ 模；但原来的表示并不只知道 $H_K$，还知道商群 $\Gamma_K$ 如何作用。因此必须把剩余的 $\Gamma_K$ 作用也记录下来，并要求它与 Frobenius 相容。

\begin{theorem}
在标准假设下，$G_K$ 的有限 $\Z_p$-表示范畴等价于 $A_K$ 上的 etale $\varphi$-$\Gamma$ 模范畴。
\end{theorem}

\begin{proof}
先限制表示到 $G_{K_\infty}$，由范域等价得到 etale $\varphi$-模。商群 $\Gamma_K=G_K/G_{K_\infty}$ 的剩余作用给出 $\Gamma$ 作用。反向则先由 $\varphi$-模恢复 $G_{K_\infty}$ 表示，再用 $\Gamma$ 的相容作用扩张为 $G_K$ 表示。

更明确地，对表示 $V$ 定义
\[
  D(V)=(A\otimes_{\Z_p}V)^{G_{K_\infty}}.
\]
$A$ 上的 Frobenius 给出 $D(V)$ 上的 $\varphi$，而 $G_K$ 对 $A\otimes V$ 的作用正规化 $G_{K_\infty}$ 的不变量，因此商群 $\Gamma_K$ 作用在 $D(V)$ 上。反向从 $(M,\varphi,\Gamma)$ 出发，先取
\[
  V(M)=(A\otimes_{A_K}M)^{\varphi=1}.
\]
这给出 $G_{K_\infty}$ 表示；$\Gamma$ 作用与 $\varphi$ 交换，保证它保持 $\varphi=1$ 的解空间，从而把作用扩张到整个 $G_K$。两个方向互逆，原因仍是扩张到 $A$ 后对象变成平凡对象，所有信息都由下降数据记录。
\end{proof}

\section{幂级数环}

\subsection{幂级数}

本节看似偏离 Galois 表示，但它提供周期环和 $p$ 进微分方程的基础语言。形式幂级数处理无穷小邻域，Robba 环处理穿孔圆环，Tate 代数处理 $p$ 进闭单位多圆盘。过收敛表示、$(\varphi,\Gamma)$ 模和 $p$ 进微分方程都需要这些收敛区域。

\begin{proposition}
在形式幂级数环 $R[[X]]$ 中，级数 $\sum a_nX^n$ 可逆，当且仅当常数项 $a_0$ 在 $R$ 中可逆。
\end{proposition}

\begin{proof}
若级数有逆，常数项相乘为 $1$，所以 $a_0$ 可逆。若 $a_0$ 可逆，设逆元为 $\sum b_nX^n$。比较 $X^0$ 系数得 $b_0=a_0^{-1}$；比较 $X^k$ 系数得
\[
  a_0b_k+a_1b_{k-1}+\cdots+a_kb_0=0,
\]
于是
\[
  b_k=-a_0^{-1}(a_1b_{k-1}+\cdots+a_kb_0).
\]
这递归唯一决定所有 $b_k$。等价地，写 $f=a_0(1-g)$，其中 $g$ 常数项为 $0$，则 $\sum_{n\ge0}g^n$ 在 $X$ 进拓扑中收敛，给出逆元。
\end{proof}

\begin{proposition}
若 $F$ 为域，则 Laurent 级数环 $F((X))$ 是域。
\end{proposition}

\begin{proof}
任意非零 Laurent 级数写为 $X^mf(X)$，其中 $f\in F[[X]]$ 且常数项非零。上一命题给出 $f^{-1}$，所以整体可逆。
\end{proof}

\subsection{Robba 环}

Robba 环由在某个穿孔圆环 $r<|T|<1$ 上收敛的级数组成。它不是单个 Banach 环，而是随半径变化的归纳极限。$p$ 进微分方程中的过收敛性通常在 Robba 环上表达。

若 $F$ 是完备非 Archimede 域，元素可写为 Laurent 级数 $\sum_{i\in\Z}c_iT^i$。在闭环带 $\alpha\le |T|\le\beta$ 上收敛，意味着 $|c_i|\alpha^i\to0$ 当 $i\to-\infty$，且 $|c_i|\beta^i\to0$ 当 $i\to+\infty$。Robba 环取所有在某个 $\alpha<|T|<1$ 上收敛的级数的并。它允许靠近外边界 $|T|=1$，但不要求在整个单位圆盘收敛。

有界 Robba 环额外要求系数有界。这个有界条件在过收敛表示中很重要：过收敛不是说级数只有有限项，而是说它在一个靠近边界的环带上有可控收敛半径。这样 Frobenius 把环带推向更靠内的环带，$\Gamma$ 作用则在边界附近仍能解析展开。

\subsection{Tate 环}

\begin{definition}
Tate 代数
\[
  F\langle X_1,\ldots,X_n\rangle
\]
由系数趋于 $0$ 的幂级数构成，配 Gauss 范数。
\end{definition}

这里“系数趋于 $0$”是按多重指标 $|\alpha|\to\infty$ 理解的。Gauss 范数为
\[
  \left\|\sum_\alpha c_\alpha X^\alpha\right\|=\max_\alpha |c_\alpha|.
\]
非 Archimede 性使这个范数对乘法表现良好：模掉严格小于 $1$ 的系数后得到剩余域上的多项式环，它没有零因子，因此最高范数项不会在乘积中消失。这是 Tate 代数比普通形式幂级数更接近解析函数环的原因。

\begin{theorem}
Tate 代数满足 Weierstrass 预备定理和除法定理；特别地，它是 Noether 环。
\end{theorem}

\begin{proof}
若 $f$ 对最后一个变量正则次数为 $d$，用 Gauss 范数把除法过程写成收缩迭代：每次消去最高项后，误差范数下降。完备性保证商和余式收敛。预备定理把 $f$ 分解为单位乘首一多项式。Noether 性由对变量数归纳并使用除法定理得到。

Weierstrass 预备定理的第一步是通过线性换元让 $f$ 的约化 $\bar f$ 对 $X_n$ 有最高次数项。若剩余域无限，选一般线性换元即可；若剩余域有限，可用 $X_i\mapsto X_i+X_n^{e_i}$ 并选 $e_i$ 把不同单项式的 $X_n$ 次数分开。这样可把问题归约到 $f$ 对 $X_n$ 正则。

除法定理中，先用一个首一多项式近似 $f$，对近似多项式可直接做多项式除法。真实的 $f$ 与近似多项式相差范数小于 $1$ 的项，因此每轮除法留下的误差范数至少乘以一个小于 $1$ 的常数。无限迭代在 Banach 完备性下收敛，得到唯一的商 $q$ 和次数小于 $d$ 的余式 $r$。Noether 性来自：任意理想若含一个对 $X_n$ 正则的元素，则该理想由这个元素和 $T_{n-1}[X_n]$ 中的有限生成理想控制；对变量数归纳即可。
\end{proof}

\section{周期环}

\subsection{\texorpdfstring{$B_{\mathrm{dR}}$}{BdR}}

令 $C$ 为 $\overline K$ 的完备化，$R=\varprojlim\mathcal O_C/p$，$A_{\inf}=W(R)$。有 Fontaine 映射
\[
  \theta:A_{\inf}\to\mathcal O_C.
\]

这里 $R$ 的元素可以写成相容序列
\[
  x=(x^{(0)},x^{(1)},x^{(2)},\ldots),\qquad (x^{(n+1)})^p=x^{(n)},
\]
其中坐标在 $\mathcal O_C$ 中取提升。虽然 $\mathcal O_C/p$ 不是完全环，逆极限 $R$ 却是特征 $p$ 的完全赋值环。这个过程常称为倾斜：它把混特征的完备域 $C$ 转成特征 $p$ 的完美对象，再用 Witt 向量 $W(R)$ 抬回特征 $0$。周期环的构造就是在这座桥上增加适合比较上同调的元素。

Witt 向量中的 Teichmuller 代表记为 $[x]$。每个元素可写为 $\sum_{n\ge0}p^n[x_n]$，Fontaine 映射为
\[
  \theta\left(\sum_{n\ge0}p^n[x_n]\right)=\sum_{n\ge0}p^n x_n^{(0)}.
\]
它把 Witt 向量的特征 $p$ 坐标重新解释为 $\mathcal O_C$ 中的 $p$ 进收敛级数。这个映射满射，因为任意 $\mathcal O_C$ 元素都可以逐步用 Teichmuller 代表和 $p$ 进展开逼近。

\begin{proposition}
$\ker\theta$ 是主理想。取生成元 $\xi$，定义
\[
  B_{\mathrm{dR}}^+=\varprojlim_n A_{\inf}[1/p]/(\ker\theta)^n,\qquad
  B_{\mathrm{dR}}=B_{\mathrm{dR}}^+[1/t],
\]
其中 $t=\log[\varepsilon]$。
\end{proposition}

\begin{proof}
选 $\varpi\in R$ 使 $\varpi^{(0)}=-p$，则 $\xi=[\varpi]+p$ 满足 $\theta(\xi)=0$。若 $x\in\ker\theta$，先模 $p$ 看 $x$ 的第一个 Witt 坐标。条件 $\theta(x)=0$ 迫使这个坐标的赋值至少与 $\varpi$ 一样大，因此可写成 $\varpi$ 的倍数。于是可取 $b\in W(R)$ 使 $x-b\xi\in pW(R)$。对剩余项重复同一过程，利用 $W(R)$ 的 $p$ 进完备性得到 $x\in\xi W(R)$。所以 $\ker\theta=(\xi)$。

定义 $B_{\mathrm{dR}}^+$ 时，是先把 $p$ 反掉，再沿 $\ker\theta$ 方向完备化：
\[
  B_{\mathrm{dR}}^+=\varprojlim_n A_{\inf}[1/p]/(\xi^n).
\]
这使 $\xi$ 成为一个完备离散赋值方向。若选择一串 $p$ 幂单位根 $\varepsilon=(1,\zeta_p,\zeta_{p^2},\ldots)$，则 $[\varepsilon]-1\in\ker\theta$，所以级数
\[
  t=\log[\varepsilon]=\sum_{n\ge1}(-1)^{n+1}\frac{([\varepsilon]-1)^n}{n}
\]
在 $B_{\mathrm{dR}}^+$ 中收敛，并且 $t$ 也是极大理想的生成元。反转 $t$ 得到 $B_{\mathrm{dR}}=B_{\mathrm{dR}}^+[1/t]$。
\end{proof}

\begin{proposition}
$B_{\mathrm{dR}}$ 是完备离散赋值域，剩余域为 $C$，且
\[
  B_{\mathrm{dR}}^{G_K}=K.
\]
\end{proposition}

\begin{proof}
$B_{\mathrm{dR}}^+$ 的极大理想由 $t$ 生成，模 $t$ 后经 $\theta$ 得到 $C$。完备性来自定义中的 $\xi$ 进完备化，而 $t$ 与 $\xi$ 只差单位，所以也是 $t$ 进完备。于是 $B_{\mathrm{dR}}^+$ 是完备离散赋值环，分式域为 $B_{\mathrm{dR}}$。

固定元结论用过滤逐级下降。若 $x\in B_{\mathrm{dR}}^{G_K}$，乘以合适的 $t^m$ 可假设 $x\in B_{\mathrm{dR}}^+$ 且其模 $t$ 项非零。模 $t$ 后得到 $C$ 中的 $G_K$ 固定元，由 Ax--Tate 定理属于 $K$。减去这个 $K$ 元素后，$x$ 的 $t$ 阶提高。逐级重复，所有系数都落在 $K$ 的相应过滤中，最终得到固定元只来自 $K$。这个结论使 $B_{\mathrm{dR}}$ 成为 de Rham 比较的正则环。
\end{proof}

\subsection{\texorpdfstring{$B_{\mathrm{cris}}$}{Bcris}}

\begin{definition}
$A_{\mathrm{cris}}$ 是 $A_{\inf}$ 关于 $\ker\theta$ 的 $p$ 进 PD 包络的完备化，定义
\[
  B_{\mathrm{cris}}^+=A_{\mathrm{cris}}[1/p],\qquad
  B_{\mathrm{cris}}=B_{\mathrm{cris}}^+[1/t].
\]
\end{definition}

$B_{\mathrm{cris}}$ 带 Frobenius 和 Galois 作用，是 crystalline 表示的比较环。为什么要引入 PD 包络？因为 $B_{\mathrm{dR}}$ 虽然有过滤，却没有合适的 Frobenius。$A_{\inf}$ 上的 Witt Frobenius 不保持 $\ker\theta$ 本身；若直接在 $\ker\theta$ 进完备化中使用 Frobenius，会碰到分母不受控的问题。PD 包络把 $\xi^n/n!$ 这样的除幂元素加入，使 Frobenius 作用在这些除幂上仍落在完备化后的环里。

更具体地，若 $\xi$ 生成 $\ker\theta$，则 $A_{\mathrm{cris}}$ 中允许级数形如
\[
  \sum_{n\ge0}a_n\frac{\xi^n}{n!},
\]
其中 $a_n\in A_{\inf}$ 且 $p$ 进趋于 $0$。这不是普通幂级数环，而是为晶体上同调中的除幂结构准备的环。Frobenius 满足 $\varphi(t)=pt$，所以反转 $t$ 后仍能定义在 $B_{\mathrm{cris}}$ 上。

\begin{proposition}
存在正合列
\[
  0\to\Q_p\to B_{\mathrm{cris}}^{\varphi=1}\to B_{\mathrm{dR}}/B_{\mathrm{dR}}^+\to0.
\]
\end{proposition}

\begin{proof}
映射由嵌入 $B_{\mathrm{cris}}\subset B_{\mathrm{dR}}$ 后取模 $B_{\mathrm{dR}}^+$ 给出。核由同时满足 $\varphi(x)=x$ 且落在 $B_{\mathrm{dR}}^+$ 的元素组成。过滤和 Frobenius 的相容性迫使这样的元素没有正权或负权贡献，最后只剩 $\Q_p$。

满性可以按 $t$ 进阶逐步提升。给定 $B_{\mathrm{dR}}/B_{\mathrm{dR}}^+$ 中的一个类，先选 $B_{\mathrm{cris}}$ 中的近似提升。若 $\varphi(x)-x$ 不为零，就在更高过滤阶中修正 $x$。修正方程本质上是 $\varphi(y)-y=z$；由于在相应分次上 Frobenius 的斜率与 $1$ 不重合，除常数项外可解。逐级修正后得到 $\varphi=1$ 的提升，正合列成立。
\end{proof}

\subsection{\texorpdfstring{$B_{\mathrm{st}}$}{Bst}}

半稳定周期环在 crystalline 环上加入一个对数变量，记作
\[
  B_{\mathrm{st}}=B_{\mathrm{cris}}[\log Y].
\]
它带 Frobenius、Galois 作用和单调算子 $N=-d/d\log Y$。单调算子记录半稳定退化中的“对数 monodromy”。

源书的构造可这样理解。取 $K$ 的素元 $\pi$，再取一串相容的 $p$ 幂根 $s=(\pi,\pi^{1/p},\pi^{1/p^2},\ldots)$，它给出 $R$ 中元素。Teichmuller 代表 $[s]$ 与 $\pi$ 的比值 $[s]\pi^{-1}$ 在 $B_{\mathrm{dR}}^+$ 中接近 $1$，所以
\[
  u_s=\log([s]\pi^{-1})
\]
收敛。令 $B_{\mathrm{st}}=B_{\mathrm{cris}}[u_s]$。换另一组根或换素元，只会把 $u_s$ 改变一个 $B_{\mathrm{cris}}$ 中的项，因此得到的环在自然意义下同构。

Frobenius 由 $\varphi(u_s)=pu_s$ 给出。单调算子 $N$ 由 $N(u_s)=-1$ 或等价地对多项式变量求导给出，并满足
\[
  N\varphi=p\varphi N.
\]
若一个表示已经 crystalline，则不需要这个对数变量，因而 $N=0$；若只有半稳定约化，对数变量记录特殊纤维交叉产生的 monodromy。

\subsection{附注}

三种周期环的层次是
\[
  B_{\mathrm{cris}}\subset B_{\mathrm{st}}\subset B_{\mathrm{dR}}.
\]
crystalline 表示最刚性，semistable 表示允许单调算子，de Rham 表示只要求 de Rham 过滤比较。

这条包含链也给出表示条件的包含：
\[
  \text{crystalline}\Rightarrow\text{semistable}\Rightarrow\text{de Rham}.
\]
从线性代数看，$B_{\mathrm{cris}}$ 提供 Frobenius，$B_{\mathrm{st}}$ 在 Frobenius 外增加 $N$，$B_{\mathrm{dR}}$ 忘掉 Frobenius 和 $N$ 而保留过滤。越往右条件越弱，能容纳的几何退化越多。

\section{\texorpdfstring{$\ell$ 进 Galois 表示}{ell 进 Galois 表示}}

\subsection{拓扑}

设 $F/\Q_p$ 有限，$K/\Q_\ell$ 有限且 $\ell\ne p$。$\ell$ 进表示是连续同态
\[
  \rho:G_F\to\GL_K(V)
\]
其中 $V$ 为有限维 $K$-向量空间。

连续性可以用格来检测。若 $L\subset V$ 是 $\mathcal O_K$-格，则 $\GL_{\mathcal O_K}(L)$ 是 $\GL_K(V)$ 的紧开子群。一个同态 $\rho:G_F\to\GL_K(V)$ 连续，当且仅当可以找到某个格 $L$ 被 $G_F$ 稳定，并且诱导同态
\[
  G_F\longrightarrow \GL_{\mathcal O_K}(L)
  \simeq \varprojlim_n\GL_{\mathcal O_K}(L/\ell^nL)
\]
连续。证明时，若 $\rho$ 连续，取任意格 $L_0$，稳定它的子群是开子群；紧性使 $G_F$ 对 $L_0$ 的轨道只有有限多个格，有限和仍是格，并被 $G_F$ 稳定。反向则由逆极限拓扑直接得到连续性。

\subsection{分歧}

\begin{proposition}
若 $\rho$ 是 $\ell$ 进表示，则野惯性群的 pro-$p$ 部分在 $\GL_n(\Z_\ell)$ 中的像有限。
\end{proposition}

\begin{proof}
紧性允许取稳定 $\mathcal O_K$-格，使像落在 $\GL_n(\mathcal O_K)$。记 $N_1=\ker(\GL_n(\mathcal O_K)\to\GL_n(\kappa_K))$，则 $N_1$ 是 pro-$\ell$ 群。另一方面，源书中把惯性群的最大 $\ell$ 进商分离出来，得到 $P_{F,\ell}$；它的有限商阶数与 $\ell$ 互素。于是 $\rho(P_{F,\ell})\cap N_1=1$。因为 $\GL_n(\kappa_K)$ 是有限群，$\rho(P_{F,\ell})$ 注入有限群，故像有限。

这一步的意义是：当 $\ell\ne p$ 时，真正的野 $p$ 分歧无法在 $\ell$ 进主同余小邻域中无限存在；经过有限扩张后，可以把这部分杀掉，只剩温和惯性的 $\ell$ 进方向。
\end{proof}

\subsection{表示}

\begin{definition}
$\ell$ 进表示称为潜半稳定，如果存在有限扩张 $E/F$，使惯性群 $I_E$ 在表示中的作用为幂零单调算子所控制，也就是对应 Weil--Deligne 表示。
\end{definition}

更基础的几个词如下。若 $\rho(I_F)=1$，称表示无分歧。若经过有限扩张后无分歧，称潜在好约化。若惯性群元素都作用为幺幂矩阵，称半稳定。若经过有限扩张后半稳定，称潜半稳定。这里“幺幂”表示矩阵形如 $\exp(N)$，其中 $N$ 幂零；$N$ 就是后面 Weil--Deligne 表示中的 monodromy 算子。

\begin{theorem}
在剩余域不含所有 $\ell$ 幂单位根的假设下，$\ell$ 进表示潜半稳定。
\end{theorem}

\begin{proof}
先由上一命题，经过有限扩张可假设 $\rho(P_{F,\ell})=1$，于是表示通过 $G_{F,\ell}=G_F/P_{F,\ell}$ 分解。惯性的最大 $\ell$ 进商为 $\Z_\ell(1)$，取生成元 $t$。共轭关系为
\[
  gtg^{-1}=t^{\chi(g)},
\]
其中 $\chi$ 是分圆特征标。若 $a$ 是 $\rho(t)$ 的特征值，则 $a^{\chi(g)}$ 也是特征值。假设剩余域没有有限扩张包含全部 $\ell$ 幂单位根，$\chi$ 的像无限；有限维空间只有有限多个特征值，故 $a$ 必为单位根。取有限指数子群后，$\rho(t)$ 的特征值全为 $1$，即惯性开子群作用为幺幂。

对幺幂矩阵可取 $\ell$ 进对数，得到幂零算子 $N=\log\rho(t)$ 的标量倍。Weil 群共轭关系把它变成
\[
  \rho(w)N\rho(w)^{-1}=\|w\|N.
\]
这正是 Weil--Deligne 表示的相容关系。
\end{proof}

\subsection{Weil}

\begin{definition}
局部域 $F$ 的 Weil 群 $W_F$ 是 $G_F$ 中映到 Frobenius 整次幂的子群，带比子空间拓扑更细的拓扑，使惯性群开。
\end{definition}

Weil 群的作用是把绝对 Galois 群中“Frobenius 的整数幂”和“惯性”分开。它有正合列
\[
  1\to I_F\to W_F\to \Z\to0.
\]
若几何 Frobenius 记为 $\Phi_F$，则 $v_F(\Phi_F)=1$，并定义 $\|w\|=q^{-v_F(w)}$。这个范数正是 Weil--Deligne 条件中缩放 $N$ 的因子。

\subsection{Weil--Deligne}

\begin{definition}
Weil--Deligne 表示是 $(r,N)$，其中 $r:W_F\to\GL(V)$ 连续且在惯性上有有限像，$N$ 为幂零算子，满足
\[
  r(w)Nr(w)^{-1}=\|w\|N.
\]
\end{definition}

从一个 $\ell$ 进表示得到 Weil--Deligne 表示的过程是：先在惯性的开子群上写成 $\rho(x)=\exp(t_\ell(x)N)$，其中 $t_\ell:I_F\to\Z_\ell(1)$ 是温和惯性的投影；再用
\[
  r(w)=\rho(w)\exp(-t_\ell(w)N)
\]
把无限幺幂惯性部分从 $\rho$ 中剥离。这样 $r$ 在惯性上只剩有限像，而 $N$ 单独记录幺幂单调部分。这个拆分使局部分歧被有限像表示和一个幂零算子共同描述。

\subsection{附注}

$\ell\ne p$ 时，表示的局部分歧通过 Weil--Deligne 对象表达；$p$ 进情形则必须使用周期环和 Hodge 过滤。

差别的根源是系数素数是否等于剩余特征。$\ell\ne p$ 时，野惯性是 pro-$p$，在 $\ell$ 进线性群的小邻域中不能无限存在，所以分歧理论较接近有限群加幂零算子。$p$ 进表示中，野惯性和系数拓扑同属 $p$ 方向，会产生丰富得多的连续变化，必须引入 $B_{\mathrm{dR}},B_{\mathrm{cris}},B_{\mathrm{st}}$ 才能线性化。

\section{\texorpdfstring{$p$ 进 Galois 表示}{p 进 Galois 表示}}

\subsection{淡中范畴}

Galois 表示范畴带张量积、对偶、内部 Hom 和纤维函子。Tannakian 观点说：如果掌握这个张量范畴，就能重建控制它的群或群概形。

本节的基本问题是：给定一个带 $G$ 作用的大环 $B$，能否把 $G$ 的 $\Q_p$ 表示 $V$ 翻译成某种线性代数对象
\[
  D_B(V)=(B\otimes_{\Q_p}V)^G?
\]
若 $B$ 选得好，$D_B(V)$ 不但维数正确，还带有过滤、Frobenius 或 monodromy 等结构。这样就能把难以直接观察的 Galois 作用换成可计算的线性代数。

\subsection{可容表示}

\begin{definition}
设 $B$ 是带 $G_K$ 作用的 $\Q_p$-代数。对表示 $V$ 定义
\[
  D_B(V)=(B\otimes_{\Q_p}V)^{G_K}.
\]
若自然映射
\[
  B\otimes_KD_B(V)\to B\otimes_{\Q_p}V
\]
为同构，则称 $V$ 为 $B$-可容。
\end{definition}

可容性的维数判别是核心。一般总有
\[
  \dim_{B^{G_K}}D_B(V)\le \dim_{\Q_p}V,
\]
因为比较映射扩张到 $B$ 后不可能产生超过 $V$ 维数的独立向量。若维数相等，比较映射就是同构，说明 $B$ 中已经有足够多的周期把 $G_K$ 作用完全平凡化。正则比较环的作用就是保证这个维数判别与张量、对偶、子商操作相容。

\subsection{比较定理}

几何上来自光滑 proper 代数簇的 etale 上同调表示，在适当周期环上可与 de Rham、crystalline 或 semistable 上同调比较。这些比较定理是 $p$ 进 Hodge 理论的中心。

若 $X/K$ 光滑 proper，$V=H^n_{\mathrm{et}}(X_{\overline K},\Q_p)$。周期环的作用是给出如下对应：$B_{\mathrm{HT}}$ 看见 Hodge--Tate 分次，$B_{\mathrm{dR}}$ 看见代数 de Rham 上同调及 Hodge 过滤，$B_{\mathrm{cris}}$ 在好约化时看见晶体上同调和 Frobenius，$B_{\mathrm{st}}$ 在半稳定约化时再加上 monodromy。于是表示是否 de Rham、crystalline、semistable，是几何约化性质在线性表示中的影子。

\subsection{De Rham 表示}

\begin{definition}
若 $V$ 对 $B_{\mathrm{dR}}$ 可容，则称 $V$ 为 de Rham 表示。此时
\[
  D_{\mathrm{dR}}(V)=(B_{\mathrm{dR}}\otimes_{\Q_p}V)^{G_K}
\]
是带过滤的 $K$-向量空间。
\end{definition}

\begin{theorem}
de Rham 表示构成稳定于子对象、商对象、扩张、张量积和对偶的张量范畴。
\end{theorem}

\begin{proof}
$B_{\mathrm{dR}}$ 是正则比较环，取不变量函子在可容对象上保持维数。对短正合列
\[
  0\to V'\to V\to V''\to0
\]
张量 $B_{\mathrm{dR}}$ 后仍正合，取不变量得到左正合列。若三项中两项已知 de Rham，维数相加迫使第三项的不变量维数也达到表示维数，于是比较映射为同构，正合性恢复。过滤严格性来自 $B_{\mathrm{dR}}$ 的 $t$ 进过滤：分次上维数相等时，过滤映射没有隐藏的余核。

张量积由乘法
\[
  B_{\mathrm{dR}}\otimes B_{\mathrm{dR}}\to B_{\mathrm{dR}}
\]
给出，且过滤满足 $\Fil^a\Fil^b\subset\Fil^{a+b}$。对偶由 $B_{\mathrm{dR}}\otimes V^\vee\simeq\operatorname{Hom}_{B_{\mathrm{dR}}}(B_{\mathrm{dR}}\otimes V,B_{\mathrm{dR}})$ 给出。因此 de Rham 表示构成张量范畴。
\end{proof}

\subsection{晶体表示}

\begin{definition}
若 $V$ 对 $B_{\mathrm{cris}}$ 可容，则称 $V$ 为 crystalline 表示。定义
\[
  D_{\mathrm{cris}}(V)=(B_{\mathrm{cris}}\otimes_{\Q_p}V)^{G_K}.
\]
它带 Frobenius，并通过嵌入 $B_{\mathrm{cris}}\subset B_{\mathrm{dR}}$ 获得过滤。
\end{definition}

\begin{theorem}
若 $K$ 剩余域为 $k$，则
\[
  B_{\mathrm{cris}}^{G_K}=K_0=W(k)[1/p].
\]
crystalline 表示给出过滤 $\varphi$-模。
\end{theorem}

\begin{proof}
固定元必须被 $G_K$ 固定。$B_{\mathrm{cris}}$ 的系数固定部分来自最大无分歧系数域
\[
  K_0=W(k)[1/p].
\]
PD 方向和 $t$ 方向在 Galois 作用或过滤中没有产生新的固定元，所以 $B_{\mathrm{cris}}^{G_K}=K_0$。对 crystalline 表示 $V$，
\[
  D_{\mathrm{cris}}(V)=(B_{\mathrm{cris}}\otimes_{\Q_p}V)^{G_K}
\]
是 $K_0$-向量空间，并继承 Frobenius。通过嵌入 $B_{\mathrm{cris}}\subset B_{\mathrm{dR}}$，把 $K\otimes_{K_0}D_{\mathrm{cris}}(V)$ 放进 $D_{\mathrm{dR}}(V)$，于是得到下降过滤。这就是过滤 $\varphi$-模。

全忠实性的意思是：两个 crystalline 表示之间的 $G_K$-等变线性映射，可以完全从对应过滤 $\varphi$-模之间的相容线性映射恢复。周期环中足够多的 crystalline 周期保证不会丢失 Hom 信息。
\end{proof}

\subsection{半稳定表示}

\begin{definition}
若 $V$ 对 $B_{\mathrm{st}}$ 可容，则称 $V$ 为 semistable 表示。定义
\[
  D_{\mathrm{st}}(V)=(B_{\mathrm{st}}\otimes_{\Q_p}V)^{G_K}.
\]
它带 Frobenius、单调算子 $N$ 和过滤。
\end{definition}

\begin{theorem}
semistable 表示给出带 $(\varphi,N,\mathrm{Fil})$ 的线性代数对象，并且 crystalline 表示正是其中 $N=0$ 的部分。
\end{theorem}

\begin{proof}
$B_{\mathrm{st}}$ 上的 $N$ 与 Frobenius 满足 $N\varphi=p\varphi N$，取不变量后传给
\[
  D_{\mathrm{st}}(V)=(B_{\mathrm{st}}\otimes_{\Q_p}V)^{G_K}.
\]
这个空间是 $K_0$-向量空间，带半线性 Frobenius、线性幂零算子 $N$，并通过 $B_{\mathrm{st}}\subset B_{\mathrm{dR}}$ 得到 $K\otimes_{K_0}D_{\mathrm{st}}(V)$ 上的过滤。

可容的过滤 $(\varphi,N)$-模要满足两个数值条件。Hodge 数 $t_H$ 来自过滤跳跃，Newton 数 $t_N$ 来自 Frobenius 行列式赋值；要求 $t_H=t_N$，并且对子对象有 $t_H\le t_N$。这正是“过滤给出的权”和“Frobenius 给出的斜率”互相匹配。Colmez--Fontaine 定理说明，可容过滤 $(\varphi,N)$-模与半稳定表示等价。

若 $N=0$，对数变量不起作用，对象下降到 $B_{\mathrm{cris}}$；反向 crystalline 对象在 $B_{\mathrm{st}}$ 中的单调算子为零。因此 crystalline 表示正是 semistable 表示中 $N=0$ 的部分。
\end{proof}

\subsection{附注}

Fontaine--Mazur 猜想把这一章的理论与全球数论相连：几何来源的不可约 $p$ 进表示应当由适当的局部 Hodge 条件和全局分歧条件刻画。这也是现代数论中 Galois 表示、模形式和 Langlands 对应交汇的地方。

\section{习题详解}

\begin{exercise}
证明域 $k$ 上晶体范畴不一定存在余核，因此不一定是 Abel 范畴。
\end{exercise}

\begin{proof}
晶体要求底层 $W(k)$-模有限自由且 Frobenius 单射。取两个晶体
\[
  M=N=W(k),\qquad \Phi_M=\Phi_N=\sigma,
\]
其中 $\sigma$ 是 Witt Frobenius。定义态射 $f:M\to N$ 为乘以 $p$。它与 Frobenius 相容，因为
\[
  \Phi_N(f(x))=\sigma(px)=p\sigma(x)=f(\Phi_M(x)).
\]
若晶体范畴中存在余核 $Q$，则遗忘到底层 $W(k)$-模后，它必须承接普通模余核 $W(k)/pW(k)$ 的泛性质。普通余核为 $k$，含有 $p$-扭，并不是有限自由 $W(k)$-模。晶体对象的底层模必须有限自由，因此 $k$ 不能成为晶体。

还要排除“换一个自由模作余核”的可能。若 $q:N\to Q$ 是晶体范畴中的余核，则 $q\circ f=0$，即 $q(p)=0$。由于 $Q$ 的底层 $W(k)$-模无 $p$-扭，$q(p)=p q(1)=0$ 推出 $q(1)=0$，于是 $q=0$。但零映射不能满足普通余核对所有杀死 $pW(k)$ 的映射的泛性质。例如 $N\to k$ 的自然投影杀死 $pW(k)$，却不能通过零映射分解。故余核不存在，晶体范畴不是 Abel 范畴。
\end{proof}

\begin{exercise}
设 $E$ 特征为 $p$，$A=E[X_1,\ldots,X_d]/(X_j^p-\sum_i c_{ij}X_i)$，且 $\det(c_{ij})\ne0$。证明 $A$ 是有限 etale $E$-代数，秩为 $p^d$。
\end{exercise}

\begin{proof}
关系式把每个 $X_j^p$ 化为一次项，所以每个单项式都可化为
\[
  X_1^{a_1}\cdots X_d^{a_d},\qquad 0\le a_i<p.
\]
这些 $p^d$ 个单项式生成 $A$。由于每个关系对 $X_j$ 是首一 $p$ 次关系，逐变量除法给出 $A$ 是有限 $E$-模，生成元正是上面的 $p^d$ 个单项式。为了证明它们线性无关，可把 $A$ 扩张到 $E^{\mathrm{sep}}$。若能证明 $A\otimes_EE^{\mathrm{sep}}$ 有 $p^d$ 个 $E^{\mathrm{sep}}$-代数同态到 $E^{\mathrm{sep}}$，则其维数至少为 $p^d$；结合生成数至多 $p^d$，维数等于 $p^d$。

先计算微分。关系式给出
\[
  d(X_j^p)-\sum_i c_{ij}dX_i=-\sum_i c_{ij}dX_i=0.
\]
因此
\[
  \Omega^1_{A/E}\simeq
  \left(\bigoplus_i A\,dX_i\right)\Big/
  \left(\sum_i c_{ij}dX_i\right)_{1\le j\le d}.
\]
矩阵 $(c_{ij})$ 可逆，所以所有 $dX_i=0$，即 $\Omega^1_{A/E}=0$。

接着证明几何点数。一个同态 $A\to E^{\mathrm{sep}}$ 等价于向量 $x=(x_1,\ldots,x_d)$ 满足
\[
  x_j^p=\sum_i c_{ij}x_i\qquad(1\le j\le d).
\]
这是加法多项式方程组，解集是 $\F_p$-向量空间。其切空间由同一线性微分矩阵控制；矩阵 $(c_{ij})$ 可逆使切空间为零，因此解都是单点。另一方面，对有限维 $\F_p$-向量空间上的可分加法同态，其核的阶等于对应有限平展代数的秩。由首一关系给出的次数乘积为 $p^d$，故解的个数为 $p^d$。

有限 $E$-代数 $A$ 满足 $\Omega^1_{A/E}=0$，且几何点数等于维数，因此 $A$ 为有限 etale $E$-代数，秩为 $p^d$。
\end{proof}

\begin{exercise}
设 $W$ 是有限维 $C$-向量空间并带连续 $G_F$ 作用。证明 Hodge--Tate 权空间给出的映射
\[
  \bigoplus_q C(-q)\otimes_F W\{q\}\to W
\]
单射，且维数和不超过 $\dim_CW$。
\end{exercise}

\begin{proof}
$W\{q\}$ 是满足 $\sigma(w)=\chi(\sigma)^{-q}w$ 的向量组成的 $F$-空间。映射
\[
  C(-q)\otimes_F W\{q\}\to W
\]
由乘法给出：$C(-q)$ 中的 Tate 基把权 $q$ 的向量还原成 $W$ 中的普通向量。

证明单射。取有限个权 $q_1<\cdots<q_m$，假设有关系
\[
  \sum_{j=1}^m c_j w_j=0,\qquad
  w_j\in W\{q_j\},\ c_j\in C(-q_j).
\]
对 $\sigma\in G_F$ 作用，并利用 $w_j$ 的权条件，得到每一项按不同的分圆特征标幂缩放。把原关系和作用后的关系相减，可消去一个权并得到更短关系；反复使用 Vandermonde 行列式，归结为某个非零项为零。需要的事实是分圆特征标的不同幂在 $G_F$ 上给出不同连续特征，因而可选若干 $\sigma$ 使矩阵 $(\chi(\sigma_i)^{q_j})$ 可逆。故所有项都为零。

既然直和到 $W$ 的映射单射，左边作为 $C$-向量空间的维数不超过 $\dim_CW$。而 $C(-q)$ 是一维 $C$-空间，所以左边维数为 $\sum_q\dim_F W\{q\}$，得到所需不等式。
\end{proof}

\begin{exercise}
证明 $B_{\mathrm{HT}}=C[t,t^{-1}]$，其完备化为 $C((t))$，并说明 Hodge--Tate 比较映射的维数判别。
\end{exercise}

\begin{proof}
按定义
\[
  B_{\mathrm{HT}}=\bigoplus_{i\in\Z}C(i).
\]
选取 $t$ 为 $\Z_p(1)$ 的基后，$C(i)=Ct^i$，乘法满足 $t^it^j=t^{i+j}$，所以 $B_{\mathrm{HT}}=C[t,t^{-1}]$。对 $t$ 进方向完备化得到下界有限的 Laurent 级数 $C((t))$。

正则性分三点。第一，$B_{\mathrm{HT}}$ 是整环，因为它是 Laurent 多项式环 $C[t,t^{-1}]$。第二，固定元为 $F$：若 $\sum_{i=m}^n c_it^i$ 被 $G_F$ 固定，则每个非零项满足 $\sigma(c_i)=\chi(\sigma)^{-i}c_i$；不同权不能混合，权 $0$ 项给出 $c_0\in C^{G_F}=F$，非零权项由 Tate 的权线性无关结论消失。第三，若一个一维 $B_{\mathrm{HT}}$-子空间被 $G_F$ 稳定，则其生成元的特征标只能是某个 Tate 权，仍可由 Laurent 单项式吸收，符合正则比较环的要求。

对表示 $V$，
\[
  D_{\mathrm{HT}}(V)=(B_{\mathrm{HT}}\otimes_{\Q_p}V)^{G_F}
\]
的维数不超过 $\dim_{\Q_p}V$；自然映射为单射。若维数相等，单射两侧同维，故为同构。

最后说明 $\operatorname{gr}B_{\mathrm{dR}}=B_{\mathrm{HT}}$。$B_{\mathrm{dR}}^+$ 的过滤由 $t$ 的幂给出，且
\[
  \operatorname{gr}^iB_{\mathrm{dR}}
  =t^iB_{\mathrm{dR}}^+/t^{i+1}B_{\mathrm{dR}}^+
  \simeq C(i).
\]
把所有 $i\in\Z$ 相加，乘法由 $t^it^j=t^{i+j}$ 给出，得到 $B_{\mathrm{HT}}=\bigoplus_iC(i)$。
\end{proof}

\begin{exercise}
证明 $B_{\mathrm{dR}}^{G_K}=K$，且 $B_{\mathrm{dR}}$ 是正则比较环。
\end{exercise}

\begin{proof}
$B_{\mathrm{dR}}^+$ 有 $t$ 进过滤，$\operatorname{gr}B_{\mathrm{dR}}\simeq B_{\mathrm{HT}}$。若 $x\in B_{\mathrm{dR}}^{G_K}$，乘以 $t^m$ 后可设 $x\in B_{\mathrm{dR}}^+$ 且 $x$ 不在 $tB_{\mathrm{dR}}^+$。它的最低非零分次项在 $B_{\mathrm{HT}}$ 中被 $G_K$ 固定。Hodge--Tate 固定元只出现在权 $0$，并属于 $K$。减去这个 $K$ 元素后，$x$ 的 $t$ 阶提高。

重复这个过程，构造一列 $K$ 中元素的 $t$ 进展开逼近 $x$。由于 $B_{\mathrm{dR}}^+$ 完备且 $K$ 在相应拓扑中闭，极限仍在 $K$。故 $B_{\mathrm{dR}}^{G_K}=K$。正则性由三点组成：$B_{\mathrm{dR}}$ 是域；固定域为 $K$；若 $B_{\mathrm{dR}}$ 的一维子空间稳定，则通过最低分次项归约到 $B_{\mathrm{HT}}$ 的权计算，生成元可调整为固定元乘以 $t^i$。这些条件给出比较映射的维数判别。
\end{proof}

\begin{exercise}
证明 de Rham 表示短正合列在 $D_{\mathrm{dR}}$ 下给出过滤向量空间短正合列。
\end{exercise}

\begin{proof}
对短正合列张量 $B_{\mathrm{dR}}$ 仍正合：
\[
0\to B_{\mathrm{dR}}\otimes V'\to
B_{\mathrm{dR}}\otimes V\to
B_{\mathrm{dR}}\otimes V''\to0.
\]
取 $G_K$ 不变量得到左正合列
\[
0\to D_{\mathrm{dR}}(V')\to D_{\mathrm{dR}}(V)\to D_{\mathrm{dR}}(V'').
\]
由于三项都是 de Rham，可容性给出
\[
\dim_KD_{\mathrm{dR}}(V')+\dim_KD_{\mathrm{dR}}(V'')
=\dim_{\Q_p}V'+\dim_{\Q_p}V''
=\dim_{\Q_p}V
=\dim_KD_{\mathrm{dR}}(V).
\]
左正合列中中间项维数已经等于左右两端应有维数之和，因此右端映射满射，得到短正合。

过滤由
\[
  \Fil^iD_{\mathrm{dR}}(V)=
  ( \Fil^iB_{\mathrm{dR}}\otimes_{\Q_p}V)^{G_K}
\]
定义。将上面的正合性应用到每个过滤阶，并比较分次维数，可得映射严格保持过滤，所以这是过滤向量空间的短正合列。
\end{proof}

\begin{exercise}
若 $V_1,V_2$ 是 de Rham 表示，证明
\[
  D_{\mathrm{dR}}(V_1)\otimes D_{\mathrm{dR}}(V_2)\simeq D_{\mathrm{dR}}(V_1\otimes V_2)
\]
为过滤空间同构。
\end{exercise}

\begin{proof}
周期环乘法给出自然映射
\[
  D_{\mathrm{dR}}(V_1)\otimes_KD_{\mathrm{dR}}(V_2)
  \longrightarrow D_{\mathrm{dR}}(V_1\otimes_{\Q_p}V_2).
\]
把两边扩张到 $B_{\mathrm{dR}}$ 后，都同构于
\[
  B_{\mathrm{dR}}\otimes_{\Q_p}V_1\otimes_{\Q_p}V_2.
\]
因此原映射是单射。两边维数都等于 $\dim_{\Q_p}V_1\dim_{\Q_p}V_2$，故为同构。过滤相容来自乘法满足 $\Fil^aB_{\mathrm{dR}}\cdot\Fil^bB_{\mathrm{dR}}\subset\Fil^{a+b}B_{\mathrm{dR}}$，所以张量积过滤
\[
  \Fil^i(D_1\otimes D_2)=\sum_{a+b=i}\Fil^aD_1\otimes\Fil^bD_2
\]
正好对应 $V_1\otimes V_2$ 的 de Rham 过滤。
\end{proof}

\begin{exercise}
证明 $\ell$ 进表示潜在无分歧的分解步骤：经过有限扩张后，惯性作用可变为平凡。
\end{exercise}

\begin{proof}
取稳定格，使像落入 $\GL_n(\Z_\ell)$。设 $\overline G$ 为 $G_F$ 在表示中的像，$\overline I$ 为惯性群 $I_F$ 的像。由于 $\rho$ 连续，$\overline G$ 是紧的 $\ell$ 进线性群的闭子群。

第一步，说明需要处理的惯性像是有限的或来自一个 $\Z_\ell$ 方向。主同余子群
\[
  1+\ell M_n(\Z_\ell)
\]
是 pro-$\ell$ 群；野惯性中的 pro-$p$ 部分在其中的像只能为有限群。商掉这个有限像后，惯性的无限部分只能来自温和惯性的最大 $\ell$ 进商 $\Z_\ell(1)$。

第二步，看无分歧商。若 $K$ 是 $\ker\rho$ 的固定域，则 $\overline G\simeq\Gal(K/F)$，而 $\overline I$ 对应 $K/K_0$，其中 $K_0/F$ 是 $K/F$ 中最大无分歧子扩张。因此
\[
  \overline G/\overline I\simeq\Gal(K_0/F),
\]
它是 $\widehat\Z$ 的商。因为 $\overline G$ 嵌入 $\GL_n(\Z_\ell)$，这个商又是 $\ell$ 进线性群的子商；其 pro-$\ell$ 部分形如 $\Z_\ell/N$，其余部分为阶与 $\ell$ 互素的有限群。

第三步，分情况取有限扩张。如果 $\Z_\ell/N$ 有限，则 $\overline G/\overline I$ 有限；再加上 $\overline I$ 有限，可得 $\overline G$ 有限。取 $E=K$，则 $G_E=\ker\rho$，故 $I_E\subset G_E$，惯性作用为平凡。

如果存在无限的 $\Z_\ell$ 商，则可在 $\overline G$ 中取一个有限正规子群 $T'$ 包含 $\overline I$，使商为 $\Z_\ell$。令 $E$ 为这个 $\Z_\ell$ 商对应的固定域。因为 $T'$ 有限，$E/F$ 是有限扩张；并且 $\overline I\subset T'$，所以 $E$ 的惯性群映入 $\ker\rho$。因此 $\rho(I_E)=1$。这完成“潜在无分歧”的分解步骤。
\end{proof}

\begin{exercise}
构造 $R_C=\varprojlim \mathcal O_C/p$，证明 $A_{\inf}=W(R_C)$ 带满射 $\theta:A_{\inf}\to\mathcal O_C$，并说明其泛性质。
\end{exercise}

\begin{proof}
先从逆极限坐标开始。$R_C$ 的元素是序列 $(x_m)$，其中 $x_m\in\mathcal O_C/p$，并满足 $x_{m+1}^p=x_m$。取任意提升 $\widehat x_m\in\mathcal O_C$。对固定 $n$，序列 $\widehat x_{m+n}^{p^m}$ 在 $\mathcal O_C$ 中是 Cauchy 列：相容条件说明相邻两项模越来越高的 $p$ 进理想相同。因此极限
\[
  x^{(n)}=\lim_{m\to\infty}\widehat x_{m+n}^{p^m}
\]
存在，并且与提升选择无关。由构造立刻有 $(x^{(n+1)})^p=x^{(n)}$。

反过来，若给出一列 $(x^{(n)})$ 满足 $(x^{(n+1)})^p=x^{(n)}$，把 $x^{(n)}$ 模 $p$ 得到 $x_n\in\mathcal O_C/p$，就得到逆极限中的元素。两个构造互逆，所以
\[
  R_C\simeq\{(x^{(n)}):x^{(n)}\in\mathcal O_C,\ (x^{(n+1)})^p=x^{(n)}\}.
\]

环结构如下。乘法逐坐标定义：
\[
  (xy)^{(n)}=x^{(n)}y^{(n)}.
\]
加法不能逐坐标相加，因为相容的 $p$ 幂根会被进位影响；定义
\[
  (x+y)^{(n)}=\lim_m(x^{(n+m)}+y^{(n+m)})^{p^m}
\]
。这个公式来自特征 $p$ 中的 Frobenius 加法，极限把混特征中的误差吸收到高 $p$ 进阶。逐项验证结合律、交换律和分配律可先在 $\mathcal O_C/p$ 的逆极限中进行，再转到提升坐标。因此 $R_C$ 是特征 $p$ 的完全环。

令 $A_{\inf}=W(R_C)$。Witt 环元素可唯一写成 Teichmuller 展开 $\sum_{n\ge0}p^n[x_n]$。定义
\[
  \theta\left(\sum p^n[x_n]\right)=\sum p^n x_n^{(0)}.
\]
Witt 多项式的定义保证这是环同态。满射证明为：给定 $c\in\mathcal O_C$，选 $x_0\in R_C$ 使 $x_0^{(0)}\equiv c\pmod p$，则 $c-\theta([x_0])\in p\mathcal O_C$；再对 $(c-\theta([x_0]))/p$ 重复，得到 $c=\sum p^n x_n^{(0)}$。

拓扑方面，令 $\mathfrak a=(p,\ker\theta)$。$A_{\inf}$ 对 $\mathfrak a$ 进拓扑 Hausdorff 完备，因为 $W(R_C)$ 已经 $p$ 进完备，而 $\ker\theta$ 由一个非零因子生成，$\mathfrak a$ 进 Cauchy 数据等价于同时控制 $p$ 进阶和 $\theta$ 方向的误差。

泛性质如下。设 $D$ 是 Hausdorff 完备的 $\Z_p$-代数，并有满射 $\theta_D:D\to\mathcal O_C$；令 $\mathfrak a_D=(p,\ker\theta_D)$。对 $x=(x^{(n)})\in R_C$，在 $D$ 中逐步选择 $x^{(n)}$ 的相容提升；完备性保证这些提升的 $p$ 幂极限存在，并给出 Teichmuller 型乘法映射 $R_C\to D/pD$。Witt 向量的泛性质把它唯一提升为 $\Z_p$-代数同态
\[
  \alpha:A_{\inf}=W(R_C)\to D
\]
且满足 $\theta_D\circ\alpha=\theta$。唯一性来自 Teichmuller 代表和 $p$ 进展开在 Hausdorff 完备环中的唯一性。
\end{proof}

\begin{exercise}
解释 Tate 代数完备化例子中得到的代数 $C$，并证明给出的乘法使它成为中心含 $\Q_p$ 的除代数型对象。
\end{exercise}

\begin{proof}
Gauss 范数完备化 $C\{X\}$ 的元素是在单位闭球上收敛的解析函数。把它放入分式域，再取代数闭包中的整闭包并按谱范数完备化，得到的环可理解为“允许有限解析覆盖的函数”。对这样的函数 $f$，每个自同构 $\tau\in T$ 都给出一个平移 $X\mapsto X+x(\tau)$，再由取值同态 $s_C$ 得到数值 $f(\tau)$。图
\[
  \Gamma_f=\{(x(\tau),f(\tau)):\tau\in T\}
\]
记录 $f$ 在所有这些平移分支上的取值。

对应的乘法不是普通函数逐点相乘，而是通过图的复合定义。若 $\|f\|_{\mathrm{sp}}\le1$，$\|g\|_{\mathrm{sp}}\le1$，则复合对应 $\Gamma_f\odot\Gamma_g$ 仍落在有界闭球中；完备化和整闭包保证可选择 $h$，使 $\Gamma_h$ 包含在这个复合对应中，并满足 $h(0)=0$。把这个 $h$ 记作 $f\cdot g$。

现在检查代数性质。满足 $f(\sigma\tau)=f(\sigma)+f(\tau)$ 的元素构成加法群，因为条件对加法封闭。乘法的结合律来自对应复合的结合律：
\[
  \Gamma_f\odot(\Gamma_g\odot\Gamma_h)
  =(\Gamma_f\odot\Gamma_g)\odot\Gamma_h.
\]
单位元由恒等对应给出。若 $f$ 非零，图对应可反向求逆，得到 $g$ 使 $f\cdot g=g\cdot f=1$；这使用 $T$ 中平移对应的可逆性和谱范数完备性。故非零元素可逆，得到除代数型结构。

最后说明中心含 $\Q_p$。$\Q_p$ 中的常数对所有 $\tau\in T$ 固定，平移和 Galois 作用都不改变它们；在对应复合中，常数对应与任何图复合的结果不依赖先后顺序。因此 $\Q_p$ 位于中心。
\end{proof}
 \chapter{L 函数}

\section{章节导读}

这一章的目标，是把前面已经建立的局部域、加元环、理元群、类域论和 Galois 表示，组织成可以做解析延拓和函数方程的对象。初学时最容易迷路的地方，是符号突然从理想、单位群和分歧群跳到 Fourier 变换、Mellin 变换和 Euler 乘积。真正的主线只有一句话：先把局部数据写成局部积分，再把局部积分相乘成全局积分，最后用 Fourier 反演和 Poisson 求和把 $s$ 与 $1-s$ 连接起来。

本章有三层结构。第一层是调和分析语言：Bruhat-Schwartz 函数是可积分、可变换、可相乘的测试函数；酉特征标给出 Fourier 变换的频率变量；Haar 测度保证积分对平移不变。第二层是 Hecke 理论：理元群上的特征标给出 Hecke $L$ 函数，$Z$ 积分给出 Euler 乘积、解析延拓和函数方程。第三层是 Artin 理论：有限 Galois 群的复表示给出 Artin $L$ 函数，导子记录分歧，Brauer 定理把一般表示的函数方程归约到一维 Hecke 情形。

学习本章时要反复区分两个群。加元环 $F_{\A}$ 是加法群，适合做 Fourier 变换；理元群 $F_{\A}^{\times}$ 是乘法群，适合写 Euler 乘积和 Mellin 型积分。$Z$ 积分正是把同一个测试函数先放在加法对象上做 Fourier 变换，再限制到乘法对象上积分，因此它能同时看见 Poisson 求和与 Euler 乘积。

\section{调和分析}

\subsection{Bruhat-Schwartz 函数}

\begin{definition}
设 $\alpha=(\alpha_1,\ldots,\alpha_n)$ 是多重指标，$|\alpha|=\alpha_1+\cdots+\alpha_n$。实向量空间 $\R^n$ 上的 Schwartz 函数，是满足如下条件的无穷可微函数 $f$：对每个整数 $N\ge 0$，
\[
  \|f\|_N=
  \sup_{|\alpha|\le N}\sup_{x\in\R^n}
  (1+|x|^2)^N
  \left|
  \frac{\partial^{|\alpha|}f}{\partial x_1^{\alpha_1}\cdots\partial x_n^{\alpha_n}}(x)
  \right|
  <\infty .
\]
所有 Schwartz 函数组成的空间记为 $\mathcal S(\R^n)$。
\end{definition}

这个定义的意义是把“函数本身衰减很快”和“所有导数也衰减很快”放在一起。Fourier 变换会把微分变成乘以频率变量，把乘以坐标变成频率方向上的微分；如果只要求函数衰减，而不要求导数衰减，Fourier 变换后就未必仍在同一类函数中。Schwartz 空间的优点正是 Fourier 变换在其中封闭。

\begin{example}
若 $q(x)$ 是 $\R^n$ 上正定二次型，$p(x)$ 是多项式，则
\[
  f(x)=p(x)e^{-q(x)}
\]
是 Schwartz 函数。因为任意阶偏导仍是一个多项式乘以 $e^{-q(x)}$，而正定二次型满足 $q(x)\ge c|x|^2$，指数衰减最终压过任意多项式增长。
\end{example}

\begin{definition}
设 $V$ 是有限维 $\Q_p$-向量空间。$V$ 上的 Bruhat 函数，是带紧支集的局部常值函数 $f:V\to\C$。这些函数组成的空间记为 $\mathcal S(V)$。
\end{definition}

在 $p$ 进空间里，局部常值不是退化条件，而是分析的正确替代物。因为 $\Q_p$ 全不连通，许多自然函数在小球上恒定；紧支集则保证积分只有有限大小的区域参与。若 $V=\Q_p$，典型函数是某个紧开集 $a+p^n\Z_p$ 的特征函数。有限个这种特征函数的线性组合，就是 $p$ 进积分里的阶梯函数。

\begin{definition}
设 $F$ 是数域，$V$ 是有限维 $F$-向量空间，$V_v=V\otimes_FF_v$，$V_{\A}=V\otimes_FF_{\A}$。选定 $V$ 的一个 $F$-基，它在几乎所有有限素位给出 $\mathcal O_v$-格 $\Lambda_v$。称函数 $\Phi:V_{\A}\to\C$ 为可分函数，如果存在局部函数 $\Phi_v:V_v\to\C$，满足几乎所有有限素位上 $\Phi_v=\mathbf 1_{\Lambda_v}$，并且
\[
  \Phi((x_v))=\prod_v\Phi_v(x_v).
\]
若无穷素位上 $\Phi_v$ 是 Schwartz 函数，有限素位上 $\Phi_v$ 是 Bruhat 函数，则有限线性组合的可分函数称为 $V_{\A}$ 上的 Bruhat-Schwartz 函数。其空间记为 $\mathcal S(V_{\A})$。
\end{definition}

这里的“几乎所有有限素位取标准格的特征函数”是限制直积思想在函数层面的体现。加元环元素 $(x_v)$ 本来要求几乎所有 $x_v\in\mathcal O_v$；测试函数也要在几乎所有位置选择 $\mathbf 1_{\mathcal O_v}$，这样无限乘积才有确定含义。

\begin{proposition}
Bruhat-Schwartz 空间 $\mathcal S(V_{\A})$ 可以看成局部空间 $\mathcal S(V_v)$ 关于标准向量 $\mathbf 1_{\Lambda_v}$ 的限制张量积；它是 $L^2(V_{\A})$ 中的稠密子空间。
\end{proposition}

\begin{proof}
先看张量积描述。一个可分函数只在有限多个位置偏离标准函数，因此它由有限多个“非标准局部因子”与其余位置的标准因子决定。有限线性组合给出的正是限制张量积的代数定义。

再看稠密性。局部紧群上的连续紧支函数在 $L^2$ 中稠密。对有限素位，紧开集的特征函数张成局部常值紧支函数；对无穷素位，Schwartz 函数包含大量光滑紧支函数的近似。最后用限制直积的紧集逼近：任何 $L^2$ 函数先由紧支函数逼近，再把该紧支函数在有限多个坐标上用局部测试函数逼近，其余坐标保持标准特征函数。这样得到的函数属于 $\mathcal S(V_{\A})$，并在 $L^2$ 范数下收敛到原函数。
\end{proof}

\subsection{Fourier 变换}

\begin{definition}
设 $G$ 是局部紧交换群。连续函数 $\varphi:G\to\C$ 称为正定函数，如果对任意
$x_1,\ldots,x_N\in G$ 和 $c_1,\ldots,c_N\in\C$ 都有
\[
  \sum_{i,j=1}^N c_i\overline{c_j}\,
  \varphi(x_i x_j^{-1})\ge 0.
\]
由连续正定函数张成的向量空间记为 $B(G)$。
\end{definition}

正定性在本章中的作用，是给 Fourier 反演提供一个干净的函数类。局部紧交换群上的
一般 $L^1$ 函数未必能逐点反演；正定函数带有“来自内积”的非负性，使得 Fourier
变换后仍有足够好的可积性。后面真正使用的是 Bruhat-Schwartz 函数，它们满足更强
的条件；这里把正定函数写出来，是为了说明抽象调和分析中的反演定理怎样落到数论
测试函数上。

\begin{definition}
设 $G$ 是局部紧交换群，$G$ 的酉特征标是连续群同态 $\chi:G\to S^1$。所有酉特征标组成的群记为 $\widehat G$，称为 $G$ 的对偶群。给定 $G$ 上 Haar 测度 $dg$，若 $\varphi\in L^1(G)$，定义 Fourier 变换
\[
  \widehat\varphi(\chi)=\int_G\varphi(g)\chi(g^{-1})\,dg .
\]
\end{definition}

这个公式中出现 $g^{-1}$，是为了让平移与特征标的乘法方向相配合。若 $G$ 写成加法群，则 $\chi(g^{-1})$ 就是 $\chi(-g)$，这与经典公式中的 $e^{-2\pi ixy}$ 相同。

\begin{definition}
给定 $G$ 上 Haar 测度 $dg$。若 $\widehat G$ 上 Haar 测度 $d\chi$ 使足够好的函数满足反演公式
\[
  \varphi(g)=\int_{\widehat G}\widehat\varphi(\chi)\chi(g)\,d\chi,
\]
则称 $d\chi$ 为 $dg$ 的对偶测度。若存在同构 $T:G\to\widehat G$，且 $T$ 把 $dg$ 推到对偶测度，则称 $dg$ 为相对于 $T$ 的自对偶测度。
\end{definition}

自对偶测度的作用，是让 Fourier 反演没有额外常数。后面写 Poisson 求和公式时，常数是否出现完全由测度归一化决定。代数数论中选择自对偶测度，是为了让 $F_{\A}/F$ 的体积和判别式之间的关系能够准确进入函数方程。

\begin{example}
在 $\R$ 上取酉特征标 $\psi(x)=e^{-2\pi ix}$，则 $dx$ 是自对偶测度，Fourier 变换为
\[
  \widehat f(y)=\int_{\R}f(x)e^{-2\pi ixy}\,dx .
\]
在 $\Q_p$ 上，若 $\psi$ 在 $\Z_p$ 上平凡而在 $p^{-1}\Z_p$ 上非平凡，则使 $\Z_p$ 体积为 $1$ 的 Haar 测度是自对偶测度。
\end{example}

\subsection{限制直积}

\begin{definition}
设 $S_\infty$ 是有限指标集。对每个 $v$ 给定局部紧交换群 $G_v$，对 $v\notin S_\infty$ 给定开紧子群 $H_v\subset G_v$。对有限集 $S\supset S_\infty$，置
\[
  G_S=\prod_{v\in S}G_v\times\prod_{v\notin S}H_v .
\]
并集
\[
  G=\bigcup_{S\supset S_\infty}G_S
\]
配上由 $\prod_vN_v$ 给出的邻域基，其中 $N_v$ 是 $G_v$ 中单位元邻域且几乎所有 $N_v=H_v$。称 $G$ 为 $G_v$ 关于 $H_v$ 的限制直积。
\end{definition}

限制直积是加元环和理元群的共同拓扑来源。$F_{\A}$ 是 $F_v$ 关于 $\mathcal O_v$ 的限制直积；$F_{\A}^{\times}$ 是 $F_v^{\times}$ 关于 $\mathcal O_v^{\times}$ 的限制直积。限制条件保证“无限多坐标”没有失控。

\begin{proposition}
设 $G=\prod'_vG_v$ 是上述限制直积。
\begin{enumerate}[label=(\arabic*)]
\item 对任一有限集 $S\supset S_\infty$，$G_S$ 中单位元的邻域基也是 $G$ 中单位元的邻域基。
\item 每个 $G_S$ 是 $G$ 的开子群，$G$ 是局部紧交换群。
\item $C\subset G$ 的闭包为紧集，当且仅当存在紧集 $B_v\subset G_v$，几乎所有 $B_v=H_v$，使 $C\subset\prod_vB_v$。
\end{enumerate}
\end{proposition}

\begin{proof}
(1) 由拓扑定义，$G$ 中单位元邻域就是有限多个坐标任取局部邻域、其余坐标取 $H_v$ 的乘积。固定 $S$ 后，$G_S$ 的相对拓扑给出的单位元邻域同样具有这一形式。

(2) 若 $x\in G_S$，则 $xG_S=G_S$，所以 $G_S$ 是自己的每点邻域，因而在 $G$ 中开。局部紧性来自 $G_S$ 中单位元附近的紧邻域：在有限多个坐标取紧邻域，在其余坐标取开紧群 $H_v$，Tychonoff 定理给出乘积紧性。交换性逐坐标继承。

(3) 若 $C$ 的闭包紧，则每个坐标投影的像闭包紧；并且紧集可由有限多个基本开集覆盖，因而除有限多个坐标外必须落在 $H_v$ 中。取这些投影闭包为 $B_v$。反过来，若 $C\subset\prod_vB_v$ 且几乎所有 $B_v=H_v$，则右端是有限多个局部紧紧集与无限多个开紧群的乘积，因此紧；闭包作为紧集的闭子集仍紧。
\end{proof}

\subsection{酉特征标}

\begin{definition}
设 $K$ 是非 Archimede 局部域，$\mathcal O_K$ 是整数环，$\mathfrak p_K$ 是极大理想。
非平凡加法酉特征标 $\psi:K\to S^1$ 的阶，是满足
\[
  \psi|_{\mathfrak p_K^{-\nu}}=1
\]
的最大整数 $\nu$。换句话说，$\mathfrak p_K^{-\nu}$ 是 $\psi$ 平凡的最大分式理想。
\end{definition}

这个整数衡量加法特征标能看见多细的 $p$ 进小数位。阶越大，$\psi$ 在越大的分式理想
上已经平凡；在 Fourier 变换中，它决定 $\mathbf 1_{\mathcal O_K}$ 变换后支撑在哪个
对偶格上。后面 Hecke 函数方程里的判别式和导子常数，就是加法特征标的阶与乘法
特征标的导子一起产生的尺度因子。

\begin{proposition}
设 $K$ 是局部域、$\R$ 或 $\C$，$\psi:K\to S^1$ 是非平凡加法酉特征标。对 $a\in K$ 令 $\psi_a(x)=\psi(ax)$。则
\[
  K\longrightarrow \widehat K,\qquad a\longmapsto \psi_a
\]
是加法拓扑群同构。
\end{proposition}

\begin{proof}
映射保持加法，因为 $\psi_{a+b}(x)=\psi(ax)\psi(bx)$。若 $\psi_a$ 平凡，则 $\psi(ax)=1$ 对所有 $x$ 成立。若 $a\ne0$，乘以 $a$ 是 $K$ 的满自同构，于是 $\psi$ 平凡，矛盾。因此映射单射。

满射是局部紧域 Pontryagin 对偶的标准结论。可以从一维向量空间的对偶性理解：任意连续加法特征标 $\chi$ 的核是开子群；非平凡 $\psi$ 的开核通过乘以合适 $a$ 可以调到与 $\chi$ 的核同阶。于是 $\chi/\psi_a$ 在一个开子群上平凡且在离散商上平凡，最终得到 $\chi=\psi_a$。在 $\R$ 和 $\C$ 上这就是所有连续加法特征标分别为 $e^{-2\pi iax}$ 和 $e^{-2\pi i\operatorname{Tr}_{\C/\R}(ax)}$。
\end{proof}

\begin{example}
在有限扩张 $K/\Q_p$ 上，令
\[
  \Lambda(x)=\lambda(\operatorname{Tr}_{K/\Q_p}x),
\]
其中 $\lambda:\Q_p\to\Q/\Z$ 取出负幂部分。则 $\psi(x)=e^{2\pi i\Lambda(x)}$ 是非平凡加法酉特征标。在 $\R$ 上常取 $\psi(x)=e^{-2\pi ix}$；在 $\C$ 上常取 $\psi(z)=e^{-2\pi i(z+\bar z)}$，也就是把 $\C$ 视为 $\R$ 上二维向量空间后的标准特征标。
\end{example}

\begin{proposition}
设 $G=\prod'_vG_v$ 是关于开紧子群 $H_v$ 的限制直积。则对偶群 $\widehat G$ 同构于 $\widehat G_v$ 关于
\[
  H_v^\perp=\{\chi_v\in\widehat G_v:\chi_v|_{H_v}=1\}
\]
的限制直积。
\end{proposition}

\begin{proof}
若 $\chi\in\widehat G$，则 $\chi$ 在单位元某个开邻域上接近 $1$。因为 $S^1$ 中不存在非平凡的足够小子群，连续性推出 $\chi$ 在某个基本开子群 $\prod_{v\notin S}H_v$ 上平凡。因此几乎所有限制 $\chi_v=\chi|_{G_v}$ 属于 $H_v^\perp$，并且
\[
  \chi((g_v))=\prod_v\chi_v(g_v),
\]
右侧只有有限多个非平凡因子。

反向，若给定 $\chi_v\in\widehat G_v$ 且几乎所有 $\chi_v$ 在 $H_v$ 上平凡，则上式定义了 $G$ 上的连续特征标。两个构造互逆。拓扑相容性由紧开集的湮灭子给出：对偶群中的基本邻域正对应有限多个局部对偶群中的基本邻域和几乎所有 $H_v^\perp$。
\end{proof}

\begin{proposition}
设 $F$ 是数域，取非平凡酉特征标 $\psi:F_{\A}/F\to S^1$。对 $a\in F_{\A}$ 令 $\psi_a(x)=\psi(ax)$。则
\[
  F_{\A}\longrightarrow \widehat{F_{\A}},\qquad a\longmapsto \psi_a
\]
是同构，并且 $\widehat{F_{\A}/F}\simeq F$。
\end{proposition}

\begin{proof}
局部同构 $F_v\simeq\widehat{F_v}$ 与限制直积对偶性结合，给出 $F_{\A}\simeq\widehat{F_{\A}}$。条件 $\psi(F)=1$ 保证对 $a\in F$，$\psi_a$ 在 $F$ 上平凡，从而给出 $F_{\A}/F$ 的特征标。若 $\psi_a$ 在 $F$ 上平凡，则 $\psi(ax)=1$ 对所有 $x\in F$ 成立。由加元环对偶和 $F$ 在 $F_{\A}$ 中的自对偶性，得到 $a\in F$。所以商群的对偶正是离散嵌入的 $F$。
\end{proof}

\subsection{测度}

\begin{proposition}
设 $K$ 是局部域、$\R$ 或 $\C$，$dx$ 是 $K$ 加法群上的 Haar 测度。则对 $a\in K^\times$，
\[
  d(ax)=|a|_K\,dx.
\]
若 $a\in F_{\A}^{\times}$，则加元环上的 Haar 测度满足
\[
  d(ax)=|a|_{\A}\,dx .
\]
\end{proposition}

\begin{proof}
推前测度 $E\mapsto dx(a^{-1}E)$ 仍是加法 Haar 测度，因此等于某个常数倍的 $dx$。这个常数按绝对值的定义就是 $|a|_K$。全局公式逐坐标相乘；因为理元 $a$ 几乎所有坐标为单位，所以无限乘积就是理元模 $|a|_{\A}$。
\end{proof}

\begin{example}
若 $K/\Q_p$ 的标准加法特征标的导子由差异理想控制，则自对偶测度满足
\[
  \operatorname{vol}(\mathcal O_K)=|\mathfrak D_{K/\Q_p}|_p^{1/2}.
\]
在 $\R$ 上，标准测度是 Lebesgue 测度；在 $\C$ 上，相对于 $z+\bar z$ 的标准特征标，自对偶测度可写成 $|dz\wedge d\bar z|$ 的固定归一化。
\end{example}

\begin{proposition}
设 $K$ 是非 Archimede 局部域，$\psi$ 是阶为 $\nu$ 的非平凡加法特征标，$dx$ 是
相对于 $\psi$ 的自对偶测度。取 $a\in K^\times$ 使 $|a|_K=|\mathfrak p_K|^\nu$，
令 $\varphi=\mathbf 1_{\mathcal O_K}$。则
\[
  \operatorname{vol}(\mathcal O_K)=|a|_K^{1/2},
  \qquad
  \widehat\varphi(y)=|a|_K^{1/2}\varphi(ay).
\]
\end{proposition}

\begin{proof}
先确定支撑。由定义
\[
  \widehat\varphi(y)=\int_{\mathcal O_K}\psi(-xy)\,dx.
\]
若 $y\in a^{-1}\mathcal O_K$，则 $xy\in a^{-1}\mathcal O_K$ 对所有
$x\in\mathcal O_K$ 成立，按 $a$ 的选取，$\psi(-xy)=1$，积分等于
$\operatorname{vol}(\mathcal O_K)$。若 $y\notin a^{-1}\mathcal O_K$，存在
$u\in\mathcal O_K$ 使 $\psi(-uy)\ne1$。把积分作平移 $x\mapsto x+u$，积分值同时
等于自身乘以 $\psi(-uy)$，于是只能为 $0$。这说明
\[
  \widehat\varphi=\operatorname{vol}(\mathcal O_K)\mathbf 1_{a^{-1}\mathcal O_K}.
\]

再用 Fourier 反演。$\widehat{\widehat\varphi}(x)=\varphi(-x)$，而上式再次变换会产生
因子 $\operatorname{vol}(\mathcal O_K)^2 |a|_K^{-1}$。要得到 $\varphi(-x)$，必须有
$\operatorname{vol}(\mathcal O_K)^2=|a|_K$，所以
$\operatorname{vol}(\mathcal O_K)=|a|_K^{1/2}$。代回即得所列公式。
\end{proof}

\begin{proposition}
在 $F_{\A}$ 上取相对于 $\psi:F_{\A}/F\to S^1$ 的自对偶测度 $dx$，则
\[
  \operatorname{vol}(F_{\A}/F)=1.
\]
若在有限素位归一化为 $\operatorname{vol}(\mathcal O_v)=1$，在无穷素位取 Lebesgue 型测度，则
\[
  \operatorname{vol}(F_{\A}/F)=|d_{F/\Q}|^{1/2}.
\]
\end{proposition}

\begin{proof}
第一条是 Poisson 求和公式与自对偶归一化的直接后果。取足够好的测试函数 $\Phi$，在商群 $F_{\A}/F$ 上对 $\sum_{\xi\in F}\Phi(x+\xi)$ 积分，并把 Fourier 展开代入。常数项比较给出商群体积为 $1$。

第二条与第一条相差局部测度归一化。有限素位从自对偶测度改为 $\operatorname{vol}(\mathcal O_v)=1$ 会乘上判别式贡献；所有有限素位贡献相乘得到 $|d_{F/\Q}|^{1/2}$，无穷素位按标准 Lebesgue 归一化补齐同一个判别式因子。
\end{proof}

\begin{proposition}
设 $v<\infty$，$dx$ 是 $F_v$ 加法群上的 Haar 测度。则
\[
  d^\times x=\frac{dx}{|x|_v}
\]
差一个正常数就是 $F_v^\times$ 上的 Haar 测度。若归一化为 $\operatorname{vol}(\mathcal O_v^\times)=1$，则常用记号为 $d\gamma_v(x)$。
\end{proposition}

\begin{proof}
对 $a\in F_v^\times$，
\[
  \frac{d(ax)}{|ax|_v}=\frac{|a|_v\,dx}{|a|_v|x|_v}=\frac{dx}{|x|_v},
\]
所以该测度对乘法平移不变。任意两个 Haar 测度只差正常数，归一化 $\mathcal O_v^\times$ 的体积即可确定常数。
\end{proof}

\begin{proposition}
令 $\gamma=\prod_v\gamma_v$ 为理元群 $F_{\A}^{\times}$ 上的常用乘法测度，其中有限素位
$\gamma_v(\mathcal O_v^\times)=1$，实素位取 $dx/|x|$，复素位取
$|dz\wedge d\bar z|/(z\bar z)$。若 $m>1$，则存在只依赖于 $F$ 的正常数 $c_F$，使
\[
  \gamma\{z\in F_{\A}^{\times}/F^\times:1\le |z|_{\A}\le m\}=c_F\log m.
\]
在标准类数公式归一化下，
\[
  c_F=\frac{2^{r_1}(2\pi)^{r_2}h_FR_F}{w_F}.
\]
\end{proposition}

\begin{proof}
把理元类群按模长分解为
$F_{\A}^{\times}/F^\times\simeq (F_{\A}^{1}/F^\times)\times\R_{>0}$。乘法测度在
$\R_{>0}$ 方向上是 $dt/t$，因此区间 $1\le t\le m$ 的体积是 $\log m$。剩下的
系数就是紧商 $F_{\A}^{1}/F^\times$ 的体积。用理想类群分解把紧商拆成有限个理想
类、无穷单位格的基本平行体和有限单位根商，分别贡献 $h_F$、调控子 $R_F$、
无穷素位的 $2^{r_1}(2\pi)^{r_2}$ 与除以 $w_F$。这些因子相乘得到所列常数。
\end{proof}

\section{特征标}

\begin{definition}
拓扑群 $G$ 的特征标是连续同态 $\chi:G\to\C^\times$。若 $\chi(G)\subset S^1$，则称为酉特征标。所有特征标组成的群记为 $\Pi(G)$。
\end{definition}

酉特征标属于调和分析；一般复值特征标属于 $L$ 函数。原因是 Euler 乘积需要一个复变量 $s$，典型因子来自 $|\cdot|^s$，它通常不是酉的。可以把一般特征标分成酉部分和模长部分。

\begin{proposition}
设 $G=G_1\times N$，其中 $G_1$ 是紧交换群，$N\simeq\R$ 或 $N\simeq\Z$。则 $\Pi(G)$ 分解为紧部分的酉特征标和一维非紧方向上的幂特征标。更具体地，存在正值特征标 $\omega_1$，使任意 $\omega\in\Pi(G)$ 的模长唯一写成
\[
  |\omega(g)|=\omega_1(g)^\sigma,\qquad \sigma\in\R,
\]
并且 $\omega_1^s$ 给出非紧方向上的复参数。
\end{proposition}

\begin{proof}
紧群到 $\R_{>0}$ 的连续同态必平凡，因为紧集在 $\R_{>0}$ 中的连续像既紧又为子群，只能是 $\{1\}$。所以特征标的模长只来自 $N$。若 $N\simeq\R$，连续同态 $\R\to\C^\times$ 形如 $t\mapsto e^{\lambda t}$，令 $\omega_1(t)=e^t$ 得参数 $\lambda=s$。若 $N\simeq\Z$，同态由 $1$ 的像决定，复对数相差 $2\pi i$ 的周期，因此参数落在某个商 $\C/ia\Z$ 中。酉部分则由紧群 $G_1$ 和非紧方向上的纯虚参数给出。
\end{proof}

\begin{proposition}
设 $G=\prod'_vG_v$ 是关于开紧子群 $H_v$ 的限制直积。若 $\chi\in\Pi(G)$，则局部限制 $\chi_v=\chi|_{G_v}$ 满足几乎所有 $\chi_v|_{H_v}=1$，并且
\[
  \chi((g_v))=\prod_v\chi_v(g_v).
\]
反过来，若 $\chi_v\in\Pi(G_v)$ 且几乎所有 $\chi_v|_{H_v}=1$，则上式定义 $G$ 的特征标。
\end{proposition}

\begin{proof}
证明与酉对偶的限制直积相同。连续性给出某个基本开子群落入 $\chi^{-1}(U)$，其中 $U$ 是 $\C^\times$ 中不含非平凡小子群的单位邻域，于是几乎所有 $H_v$ 在 $\chi$ 下平凡。乘积公式只有有限个非平凡因子，所以定义良好且连续。
\end{proof}

\subsection{局部域的特征标}

\begin{definition}
设 $K$ 是非 Archimede 局部域，$\mathcal O$ 为整数环，$\mathfrak p$ 为极大理想，$U^{(0)}=\mathcal O^\times$，$U^{(n)}=1+\mathfrak p^n$。若 $\omega\in\Pi(K^\times)$，称最小的 $f\ge0$ 使 $\omega|_{U^{(f)}}=1$ 的 $\mathfrak p^f$ 为 $\omega$ 的导子。导子为 $\mathcal O$ 的特征标称为无分歧特征标。
\end{definition}

取素元 $\pi$ 后有
\[
  K^\times\simeq \mathcal O^\times\times \pi^\Z .
\]
因此一个局部特征标由单位群上的有限层信息和 $\pi$ 的取值共同决定。无分歧特征标在单位群上平凡，所以完全由 $\omega(\pi)$ 决定，常写成 $|x|_K^s$ 的形式。

\begin{example}
局部 Archimede 特征标可显式写出：
\[
  \Pi(\R^\times)=\{x\mapsto \operatorname{sgn}(x)^A |x|^s: A\in\{0,1\},\ s\in\C\},
\]
\[
  \Pi(\C^\times)=\{z\mapsto z^{-A}\bar z^{-B}(z\bar z)^s:
  A,B\in\Z,\ \min(A,B)=0,\ s\in\C\}.
\]
这里整数指数记录圆周方向上的代数振荡，复参数 $s$ 记录径向方向上的幂增长。
\end{example}

\begin{proposition}
设 $K$ 是非 Archimede 局部域，$\psi$ 是阶为 $\nu$ 的非平凡加法酉特征标，$dx$ 是自对偶测度。设乘法特征标 $\omega$ 的导子为 $\mathfrak p^f$，$f\ge1$。令 $\varphi$ 在 $\mathcal O^\times$ 上等于 $\omega^{-1}$，在其外等于 $0$。取 $b\in K^\times$ 满足 $|b|=|\mathfrak p|^{\nu+f}$，则
\[
  \widehat\varphi(y)=c\,|b|^{1/2}\varphi(by),
\]
其中
\[
  c=|b|^{-1/2}\int_{\mathcal O^\times}\omega^{-1}(x)\psi(b^{-1}x)\,dx .
\]
\end{proposition}

\begin{proof}
Fourier 变换为
\[
  \widehat\varphi(y)=\int_{\mathcal O^\times}\omega^{-1}(x)\psi(-xy)\,dx .
\]
若 $by\notin\mathcal O^\times$，则 $y$ 的赋值使积分在某个足够小的乘法邻域上出现非平凡加法振荡。把 $\mathcal O^\times$ 分成 $U^{(f)}$ 的陪集，在每个陪集上 $\omega$ 恒定，而 $\psi(-xy)$ 沿陪集求和为零，故积分为零。

若 $by\in\mathcal O^\times$，令 $u=by$。变量替换把 $\psi(-xy)$ 写成与 $\psi(b^{-1}x)$ 同导子的一族 Gauss 积分。由于 $\omega$ 在 $U^{(f)}$ 上平凡而在 $U^{(f-1)}$ 上非平凡，所有非零情形只差乘法因子 $\omega^{-1}(u)$。按定义吸收常数后得到所列公式。
\end{proof}

\subsection{理元群的特征标}

对数域 $F$，限制直积性质给出：任意 $\omega\in\Pi(F_{\A}^{\times})$ 都可写成
\[
  \omega=\prod_v\omega_v,
\]
且几乎所有有限素位上 $\omega_v$ 无分歧。乘积公式 $|a|_{\A}=1$ 对 $a\in F^\times$ 成立，因此
\[
  \omega_s(a)=|a|_{\A}^s
\]
是 $F_{\A}^{\times}/F^\times$ 上的特征标。

\begin{proposition}
设 $F_{\A}^{1}=\{a\in F_{\A}^{\times}:|a|_{\A}=1\}$。则
\[
  F_{\A}^{\times}/F^\times\simeq (F_{\A}^{1}/F^\times)\times \R_{>0},
\]
其中第一因子紧。因此 $F_{\A}^{\times}/F^\times$ 的特征标由紧商上的酉特征标和 $|\cdot|_{\A}^s$ 共同给出。
\end{proposition}

\begin{proof}
选择一个无穷素位 $v_0$，把 $t\in\R_{>0}$ 嵌入到该素位的正实方向，其他坐标取 $1$。任意理元 $a$ 可唯一写成 $a=t\cdot b$，其中 $t=|a|_{\A}$，$b\in F_{\A}^{1}$。商去 $F^\times$ 后得到直积分解。紧性是理元类群紧性的标准形式，来自理想类群有限性和单位定理。
\end{proof}

\subsection{Hecke 特征标}

\begin{definition}
Hecke 特征标是连续同态
\[
  \chi:F_{\A}^{\times}\to S^1
\]
并满足 $\chi(F^\times)=1$。若存在整理想 $\mathfrak m$，使 $\chi$ 在相应的主同余开子群上平凡，则称 $\mathfrak m$ 为 $\chi$ 的一个定义模。
\end{definition}

这个定义把 Dirichlet 特征标从 $\Q$ 推广到任意数域。Dirichlet 特征标只看整数模 $q$ 的剩余类；Hecke 特征标同时看有限素位上的同余信息、无穷素位上的符号或角度信息，以及理元类群的整体信息。

\begin{definition}
设 $\mathfrak m$ 是整理想。模 $\mathfrak m$ 的 Dirichlet 特征标，是在 $F_{\A}^{\times}$ 上平凡于 $F_{\A,\mathfrak m}^{\times}F^\times$ 的 Hecke 特征标。等价地，它是射线类群
\[
  C_{\mathfrak m}=F_{\A}^{\times}/F_{\A,\mathfrak m}^{\times}F^\times
\]
的特征标。
\end{definition}

为了把这个定义与更传统的“理想上的特征标”连起来，需要把理元商拆成三个可见部分：
有限素理想、有限同余剩余类、无穷单位方向。设 $I_F^{\mathfrak m}$ 是由所有
$\mathfrak p\nmid\mathfrak m$ 生成的分式理想群，设
\[
  S_F^{\mathfrak m}=\{a\mathcal O_F:a\equiv 1\pmod{\mathfrak m}\}
\]
为主同余理想群。射线类群可写成
\[
  C_{\mathfrak m}\simeq I_F^{\mathfrak m}/S_F^{\mathfrak m}.
\]
这个同构的意义是：理元的有限非单位坐标记录一个分式理想；若改变素元代表，
改变只来自局部单位，而这些单位在射线类群中被同余条件吸收。

\begin{proposition}
设 $F(\mathfrak m)=\{a\in F^\times:(a,\mathfrak m)=1\}$，设 $\mathcal O_F^{\mathfrak m}$
为满足相应同余条件的单位群。存在正合列
\[
  1\longrightarrow F(\mathfrak m)/\mathcal O_F^{\mathfrak m}
  \longrightarrow
  I_F^{\mathfrak m}\times(\mathcal O_F/\mathfrak m)^\times
  \times F_\infty^\times/\mathcal O_F^{\mathfrak m}
  \longrightarrow C_{\mathfrak m}\longrightarrow 1.
\]
第一个箭头把 $a$ 送到 $((a)^{-1},a,a)$，第二个箭头把理想部分、有限同余部分和无穷
部分相乘为一个理元类。
\end{proposition}

\begin{proof}
先说明满射。给定一个理元类，利用逼近定理选 $x\in F^\times$，使得在所有
$\mathfrak p\mid\mathfrak m$ 的有限位置，$x$ 把该理元调到规定的主同余开子群内。
剩余的有限非单位坐标只出现在 $\mathfrak p\nmid\mathfrak m$ 的位置，因此给出
$I_F^{\mathfrak m}$ 中的一个分式理想；有限剩余类给出
$(\mathcal O_F/\mathfrak m)^\times$ 中的元素；无穷坐标给出
$F_\infty^\times/\mathcal O_F^{\mathfrak m}$ 中的元素。三部分相乘恢复原来的理元类。

再看复合为零。若从 $a\in F(\mathfrak m)$ 出发，理想部分给出 $(a)^{-1}$，有限同余
部分给出 $a$，无穷部分也给出 $a$。三者在理元群中相乘后正是主理元 $a^{-1}a$ 的
形式，商去 $F^\times$ 后成为单位元。

最后确定核。若三元组映到平凡理元类，则存在主元 $x\in F^\times$ 和主同余开子群
元素，使该三元组的理元代表等于 $x$。比较 $\mathfrak p\nmid\mathfrak m$ 的赋值可得
理想部分为 $(ax^{-1})^{-1}$；比较 $\mathfrak p\mid\mathfrak m$ 的有限坐标可得
$ax^{-1}\equiv1\pmod{\mathfrak m}$；比较无穷坐标可得无穷部分也由同一个
$ax^{-1}$ 给出。因此该核正是第一个箭头的像。
\end{proof}

\begin{proposition}
模 $\mathfrak m$ 的 Hecke 特征标可等价描述为模 $\mathfrak m$ 的大特征标：即定义在与 $\mathfrak m$ 互素的分式理想群 $I_F^{\mathfrak m}$ 上的特征标 $\chi$，并存在有限同余部分 $\chi_0$ 和无穷部分 $\chi_\infty$，使对所有与 $\mathfrak m$ 互素的 $a\in F^\times$ 有
\[
  \chi((a))=\chi_0(a)\chi_\infty(a).
\]
\end{proposition}

\begin{proof}
从理元类取代表。有限素位中与 $\mathfrak m$ 互素的素元给出理想部分；模 $\mathfrak m$ 的单位给出有限同余部分；无穷坐标给出无穷部分。三者相乘得到射线类群中的元素。唯一的相容关系来自主理想：若 $a\in F^\times$ 与 $\mathfrak m$ 互素，则由 $a$ 产生的理想、有限剩余类和无穷坐标在理元商中相乘为 $1$。因此在特征标上必须满足所列等式。

反过来，若给定满足该相容关系的大特征标，则对任意理元类先分解成理想部分、有限同余部分和无穷部分，再用三部分特征标相乘定义值。若换另一种分解，两者差别正是一个主元关系，相容条件保证取值不变。
\end{proof}

\begin{example}
当 $F=\Q$ 且 $\mathfrak m=(q)$ 时，$I_\Q^{(q)}$ 由素数 $p\nmid q$ 自由生成。模
$q$ 的 Dirichlet 特征标 $\chi_0:(\Z/q\Z)^\times\to S^1$ 给出大特征标
\[
  \chi\bigl((n)\bigr)=\chi_0(n)
\]
其中 $n$ 与 $q$ 互素。相容条件说：主理想 $(a)$ 的取值必须等于 $a$ 的模 $q$ 剩余
类取值；这正是初等 Dirichlet 特征标的定义。Hecke 语言的新增部分，是把这个构造
同时允许到任意数域、任意理想模以及无穷素位的符号条件。
\end{example}

\section{Z 积分}

\subsection{加元环上的 Fourier 变换}

取 $\psi:F_{\A}/F\to S^1$ 非平凡，并用它识别 $F_{\A}$ 与自身对偶。对 $\Phi\in\mathcal S(F_{\A})$，定义
\[
  \widehat\Phi(y)=\int_{F_{\A}}\Phi(x)\psi(-xy)\,dx .
\]
若测度自对偶，则 Fourier 反演写成
\[
  \Phi(x)=\int_{F_{\A}}\widehat\Phi(y)\psi(xy)\,dy .
\]

\begin{theorem}
设 $\Phi$ 是满足通常收敛条件的加元环测试函数，则 Poisson 求和公式成立：
\[
  \sum_{\xi\in F}\Phi(\xi)=\sum_{\xi\in F}\widehat\Phi(\xi).
\]
更一般地，对 $a\in F_{\A}^{\times}$，
\[
  \sum_{\xi\in F}\Phi(a\xi)
  =
  |a|_{\A}^{-1}
  \sum_{\xi\in F}\widehat\Phi(\xi/a).
\]
\end{theorem}

\begin{proof}
令
\[
  P_\Phi(x)=\sum_{\xi\in F}\Phi(x+\xi)
\]
作为紧群 $F_{\A}/F$ 上的函数。收敛条件保证可以逐项积分。它的第 $\eta\in F$ 个 Fourier 系数为
\[
  \int_{F_{\A}/F}P_\Phi(x)\psi(-\eta x)\,dx
  =
  \int_{F_{\A}}\Phi(x)\psi(-\eta x)\,dx
  =
  \widehat\Phi(\eta).
\]
在 $x=0$ 处展开 Fourier 级数，就得第一式。

第二式对函数 $\Psi(x)=\Phi(ax)$ 应用第一式。变量替换给出
\[
  \widehat\Psi(y)=|a|_{\A}^{-1}\widehat\Phi(y/a),
\]
代回即可。
\end{proof}

\subsection{Euler 乘积}

\begin{definition}
对 $\Phi\in\mathcal S(F_{\A})$、理元特征标 $\omega$ 和 $F_{\A}^{\times}$ 上 Haar 测度 $d^\times z$，定义
\[
  Z(\Phi,\omega)=\int_{F_{\A}^{\times}}\Phi(z)\omega(z)\,d^\times z.
\]
\end{definition}

这里把加法测试函数限制到乘法群上积分。若 $\omega=|\cdot|^s$，这个积分就是 adelic Mellin 变换；若 $\Phi$ 可分，它会分解成局部 Mellin 变换，并形成 Euler 乘积。

\begin{theorem}
设 $\Phi=\prod_v\Phi_v$ 是标准可分函数，$\omega=\prod_v\omega_v$ 是 $F_{\A}^{\times}/F^\times$ 的特征标，且 $|\omega|=|\cdot|_{\A}^{\sigma}$，$\sigma>1$。则
\[
  Z(\Phi,\omega)
  =
  \prod_v Z_v(\Phi_v,\omega_v),
\qquad
  Z_v(\Phi_v,\omega_v)=\int_{F_v^\times}\Phi_v(z)\omega_v(z)\,d^\times z .
\]
\end{theorem}

\begin{proof}
绝对收敛先在有限集合 $S$ 上截断。对 $v\notin S$，标准函数为 $\mathbf 1_{\mathcal O_v}$，无分歧特征标满足 $\omega_v(\pi_v)=q_v^{-s_v}$，局部积分是几何级数
\[
  \sum_{n\ge0}\omega_v(\pi_v)^n
  =
  (1-\omega_v(\pi_v))^{-1}.
\]
当 $\sigma>1$ 时，$\sum_v q_v^{-\sigma}$ 收敛，因此局部因子的无穷乘积绝对收敛。有限多个非标准位置的积分是普通局部积分。由 Fubini 定理和截断乘积极限，得到全局积分等于局部积分乘积。
\end{proof}

\subsection{Mellin 变换}

\begin{definition}
设 $f:\R_{>0}\to\C$。若对 $\alpha<\sigma<\beta$ 有
\[
  \int_0^\infty |f(t)|t^\sigma\,\frac{dt}{t}<\infty,
\]
则称 $f$ 为 $(\alpha,\beta)$ 型。其 Mellin 变换定义为
\[
  M(f,s)=\int_0^\infty f(t)t^s\,\frac{dt}{t},
\qquad \alpha<\Re(s)<\beta .
\]
\end{definition}

Mellin 变换就是乘法群 $\R_{>0}$ 上的 Fourier 变换。令 $t=e^x$，则 $dt/t=dx$，而 $t^s=e^{sx}$。加法变量 $x$ 上的指数函数变成乘法变量 $t$ 上的幂函数。

\begin{example}
对 $f(t)=e^{-t}$，
\[
  M(f,s)=\int_0^\infty e^{-t}t^s\,\frac{dt}{t}
  =
  \int_0^\infty e^{-t}t^{s-1}\,dt
  =
  \Gamma(s),
\]
收敛区域是 $\Re(s)>0$。
\end{example}

\begin{proposition}
Mellin 变换把函数的变换关系翻译成 $L$ 函数的函数方程。若 $f$ 与其 Fourier 变换 $\widehat f$ 满足合适衰减，则通常存在形如
\[
  M(f,s)=\gamma(s)M(\widehat f,1-s)
\]
的公式，其中 $\gamma(s)$ 由无穷素位的 Gamma 因子组成。
\end{proposition}

\begin{proof}
加法 Fourier 变换把 $x$ 与 $y$ 对换；在乘法变量中，这对应把尺度 $t$ 与倒尺度 $t^{-1}$ 对换。Mellin 核 $t^s$ 在倒尺度下变成 $t^{1-s}$，Jacobian 贡献正是 $1-s$ 中的常数 $1$。无穷素位上高斯函数的 Mellin 积分给出 $\Gamma$ 函数，因而比例因子就是 Gamma 因子。
\end{proof}

\subsection{积分的函数方程}

\begin{definition}
称 $\Phi:F_{\A}\to\C$ 为可容函数，如果 $\Phi$ 与 $\widehat\Phi$ 都可积、满足 Poisson 求和所需的一致收敛条件，并且对 $\sigma>1$，$\Phi(a)|a|_{\A}^{\sigma}$ 与 $\widehat\Phi(a)|a|_{\A}^{\sigma}$ 在 $F_{\A}^{\times}$ 上可积。Bruhat-Schwartz 函数属于这一类。
\end{definition}

\begin{theorem}
设 $\Phi$ 可容，$\omega\in\Pi(F_{\A}^{\times}/F^\times)$，且初始区域满足 $|\omega|=|\cdot|_{\A}^{\sigma}$、$\sigma>1$。则 $Z(\Phi,\omega)$ 可解析延拓为 $\omega$ 的亚纯函数，只可能在平凡特征标和 $|\cdot|_{\A}$ 处有单极，并满足函数方程
\[
  Z(\Phi,\omega)=Z(\widehat\Phi,\omega^\vee),
\qquad
  \omega^\vee(a)=|a|_{\A}\omega(a)^{-1}.
\]
若乘法测度归一化后 $F_{\A}^{1}/F^\times$ 的体积为 $V_F$，则两个可能极点的留数分别为
\[
  \operatorname*{Res}_{\omega=1}Z(\Phi,\omega)=-V_F\Phi(0),
  \qquad
  \operatorname*{Res}_{\omega=|\cdot|_{\A}}Z(\Phi,\omega)=V_F\widehat\Phi(0).
\]
\end{theorem}

\begin{proof}
把 $F_{\A}^{\times}$ 分解为模长方向和 $F_{\A}^{1}$：
\[
  Z(\Phi,\omega)=\int_0^\infty Z_t(\Phi,\omega)\,\frac{dt}{t},
\qquad
  Z_t(\Phi,\omega)=\int_{F_{\A}^{1}}\Phi(tb)\omega(tb)\,db.
\]
这里 $t$ 表示模长方向，$b$ 表示模长为 $1$ 的理元方向。由于 $\omega$ 在
$F^\times$ 上平凡，可以把 $F_{\A}^{1}$ 对 $F^\times$ 取商，并使用拓扑群积分公式
\[
  \int_{F_{\A}^{1}}\Phi(tb)\omega(tb)\,db
  =
  \int_{F_{\A}^{1}/F^\times}
  \sum_{\alpha\in F^\times}\Phi(\alpha tb)\omega(tb)\,d\dot b.
\]
补上 $\alpha=0$ 项后，可以对 $\sum_{\alpha\in F}\Phi(\alpha tb)$ 应用 Poisson 公式：
\[
  \sum_{\alpha\in F}\Phi(\alpha tb)
  =
  |t|^{-1}\sum_{\alpha\in F}\widehat\Phi(\alpha/(tb)).
\]
减去零项后得到
\[
  Z_t(\Phi,\omega)+\Phi(0)\int_{F_{\A}^{1}/F^\times}\omega(tb)\,db
  =
  Z_{1/t}(\widehat\Phi,\omega^\vee)
  +
  \widehat\Phi(0)\int_{F_{\A}^{1}/F^\times}\omega^\vee(b/t)\,db .
\]

把 $t$ 从 $0$ 到 $1$ 的积分用上式改写成 $1$ 到 $\infty$ 的积分，得到
\[
  Z(\Phi,\omega)=
  \int_1^\infty\bigl(Z_t(\Phi,\omega)+Z_t(\widehat\Phi,\omega^\vee)\bigr)\frac{dt}{t}
  +\Delta(\Phi,\omega).
\]
可容条件保证第一个积分对所有特征标解析。现在只需计算误差项。若 $\omega$ 在紧群
$F_{\A}^{1}/F^\times$ 上非平凡，则
\[
  \int_{F_{\A}^{1}/F^\times}\omega(b)\,d\dot b=0,
\]
所以 $\Delta(\Phi,\omega)=0$。若 $\omega(b)=1$ 在紧商上恒成立，则 $\omega$ 只来自
模长方向，可写成 $\omega(t)=t^s$。这时
\[
  \Delta(\Phi,\omega)=
  V_F\left(\frac{\widehat\Phi(0)}{s-1}-\frac{\Phi(0)}{s}\right).
\]
因此只有 $s=0$ 与 $s=1$ 两个单极，换回特征标语言分别是 $\omega=1$ 与
$\omega=|\cdot|_{\A}$，留数如定理所列。上式右端在交换
$(\Phi,\omega)$ 与 $(\widehat\Phi,\omega^\vee)$ 后保持同一表达式，所以得到函数方程。
\end{proof}

\section{Hecke L 函数}

\subsection{标准函数的 Fourier 变换}

给定 Hecke 特征标 $\omega=\prod_v\omega_v$。在每个局部位置选择与 $\omega_v$ 匹配的标准函数 $\Phi_{\omega,v}$：
\begin{enumerate}[label=(\arabic*)]
\item 若 $v$ 是实素位，取高斯函数乘以适当幂 $x^N$。
\item 若 $v$ 是复素位，取 $z^A\bar z^B e^{-2\pi z\bar z}$。
\item 若 $v<\infty$ 且 $\omega_v$ 无分歧，取 $\mathbf 1_{\mathcal O_v}$。
\item 若 $v<\infty$ 且 $\omega_v$ 有导子，取 $\mathcal O_v^\times$ 的特征函数乘 $\omega_v^{-1}$。
\end{enumerate}
全局函数 $\Phi_\omega=\prod_v\Phi_{\omega,v}$ 称为附属于 $\omega$ 的标准函数。

\begin{proposition}
附属于 $\omega$ 的标准函数满足
\[
  \widehat{\Phi_\omega}(y)=c\,|b|_{\A}^{1/2}\Phi_{\omega^{-1}}(by),
\]
其中 $b\in F_{\A}^{\times}$ 由加法特征标导子和 $\omega$ 的局部导子决定，$c$ 是局部 Gauss 常数的乘积。
\end{proposition}

\begin{proof}
Fourier 变换逐局部分解。有限无分歧位置，$\mathbf 1_{\mathcal O_v}$ 的 Fourier 变换仍是某个格的特征函数；导子归一化后得到同型公式。有限分歧位置使用前面局部 Fourier 变换命题，得到 Gauss 常数和尺度 $b_v$。实、复位置使用高斯函数求导：$x^Ne^{-\pi x^2}$ 和 $z^A\bar z^Be^{-2\pi z\bar z}$ 的 Fourier 变换仍是同类函数，只多出 $i^{-N}$ 这样的相位常数。所有局部公式相乘，得到全局公式。
\end{proof}

\subsection{计算局部积分}

\begin{proposition}
设 $\Phi_{\omega,v}$ 是附属于 $\omega_v$ 的局部标准函数，且局部乘法测度按常用方式归一化。则局部积分
\[
  Z_v(\Phi_{\omega,v},\omega_v|\cdot|_v^s)
\]
给出如下因子：
\[
  \Gamma_{\R}(s+s_v)=\pi^{-(s+s_v)/2}\Gamma((s+s_v)/2)
\]
在实素位出现；
\[
  \Gamma_{\C}(s+s_v)=(2\pi)^{1-(s+s_v)}\Gamma(s+s_v)
\]
在复素位出现；有限无分歧位置给出
\[
  (1-\omega_v(\pi_v)q_v^{-s})^{-1};
\]
有限分歧位置在标准归一化下给出 $1$。
\end{proposition}

\begin{proof}
实素位积分化为
\[
  \int_{\R^\times}e^{-\pi x^2}|x|^{s+s_v}\frac{dx}{|x|}
  =
  \pi^{-(s+s_v)/2}\Gamma((s+s_v)/2),
\]
变量替换 $u=\pi x^2$ 给出等式。复素位用极坐标 $z=re^{i\theta}$，角向积分处理 $z^A\bar z^B$ 与特征标的角向抵消，径向积分化为
\[
  \int_0^\infty e^{-2\pi r^2}r^{2(s+s_v)}\frac{dr}{r},
\]
这给出 $\Gamma_{\C}$。

有限无分歧位置把 $F_v^\times$ 分解为 $\bigsqcup_{n\in\Z}\pi_v^n\mathcal O_v^\times$。由于 $\mathbf 1_{\mathcal O_v}$ 只保留 $n\ge0$，积分成为
\[
  \sum_{n\ge0}\omega_v(\pi_v)^nq_v^{-ns}
  =
  (1-\omega_v(\pi_v)q_v^{-s})^{-1}.
\]
有限分歧位置的测试函数支撑在单位群上并含有 $\omega_v^{-1}$，与积分中的 $\omega_v$ 相消，单位群体积归一化为 $1$。
\end{proof}

\subsection{函数方程}

\begin{definition}
设 $\omega$ 是 Hecke 特征标，$\mathfrak f(\omega)$ 是有限部分导子。定义有限部分
\[
  L(s,\omega)=\prod_{\mathfrak p\nmid\mathfrak f(\omega)}
  (1-\omega_{\mathfrak p}(\pi_{\mathfrak p})N\mathfrak p^{-s})^{-1}.
\]
令 $L_\infty(s,\omega)$ 为全部无穷素位 Gamma 因子的乘积。完全 Hecke $L$ 函数为
\[
  \Lambda(s,\omega)=
  \bigl(|d_{F/\Q}|\,N\mathfrak f(\omega)\bigr)^{s/2}
  L_\infty(s,\omega)L(s,\omega),
\]
具体归一化随加法特征标和测度选择等价改变。
\end{definition}

\begin{theorem}
若 $\omega$ 不是纯粹的无分歧平凡型特征标，则 $\Lambda(s,\omega)$ 延拓为整函数，并满足
\[
  \Lambda(s,\omega)=W(\omega)\Lambda(1-s,\omega^{-1}),
\]
其中 $|W(\omega)|=1$。
\end{theorem}

\begin{proof}
把 $Z(\Phi_\omega,\omega|\cdot|^s)$ 写成局部乘积。上一小节的局部计算说明它等于一个非零常数乘
\[
  L_\infty(s,\omega)L(s,\omega).
\]
再对 $Z$ 积分应用函数方程：
\[
  Z(\Phi_\omega,\omega|\cdot|^s)
  =
  Z(\widehat{\Phi_\omega},\omega^{-1}|\cdot|^{1-s}).
\]
用标准函数的 Fourier 变换公式，把右端改写为常数乘
\[
  Z(\Phi_{\omega^{-1}},\omega^{-1}|\cdot|^{1-s}).
\]
尺度 $b$ 的模长贡献与有限导子、判别式共同组成完全 $L$ 函数前面的因子；局部 Gauss 常数的乘积组成根数 $W(\omega)$。当 $\omega$ 非平凡时，$Z$ 积分的两个可能极点不会落在当前特征标族的 $s=0,1$ 位置，因而得到整性；平凡特征标时会出现 Dedekind zeta 的单极。
\end{proof}

\section{Artin L 函数}

\subsection{局部 Galois 表示}

\begin{definition}
有限群 $G$ 的复表示是同态 $\rho:G\to\GL(V)$，其中 $V$ 是有限维复向量空间。其特征标为
\[
  \chi_\rho(g)=\operatorname{Tr}(\rho(g)).
\]
对偶表示 $\rho^\vee$ 定义在 $V^\ast$ 上：
\[
  (\rho^\vee(g)\ell)(v)=\ell(\rho(g^{-1})v).
\]
\end{definition}

Artin $L$ 函数把 Galois 群元素放进矩阵行列式。特征标的作用是把表示论信息压缩成共轭类函数；由于 Frobenius 元只确定到共轭，局部因子必须只依赖共轭类。

\begin{proposition}
有限群表示的以下事实将在 Artin $L$ 函数中反复使用。
\begin{enumerate}[label=(\arabic*)]
\item 两个复表示等价，当且仅当它们的特征标相等。
\item 不可约特征标关于内积
\[
  \langle \chi,\psi\rangle=\frac1{|G|}\sum_{g\in G}\chi(g)\overline{\psi(g)}
\]
构成标准正交组。
\item 若 $V$ 的特征标为 $\chi$，则
\[
  \dim V^G=\langle \chi,1_G\rangle .
\]
\item 正则表示的特征标 $r_G$ 满足 $r_G(1)=|G|$，$g\ne1$ 时 $r_G(g)=0$；增广特征标为
$u_G=r_G-1_G$。
\end{enumerate}
\end{proposition}

\begin{proof}
(1) 与 (2) 是有限群复表示完全可约性和 Schur 正交关系的直接结果。为了看见
(3)，考虑投影算子
\[
  P=\frac1{|G|}\sum_{g\in G}\rho(g).
\]
对任意 $h\in G$ 有 $\rho(h)P=P$，所以 $P$ 的像落在 $V^G$；若 $v\in V^G$，则
$Pv=v$，所以 $P$ 正是到不变量子空间的投影。投影的迹等于像的维数，而
\[
  \operatorname{Tr}(P)=\frac1{|G|}\sum_{g\in G}\chi(g)=\langle\chi,1_G\rangle .
\]
(4) 中正则表示以群元素为基，左乘 $g$ 的矩阵只有在 $g=1$ 时固定全部基向量；当
$g\ne1$ 时没有固定基向量，所以迹为 $0$。减去平凡表示的特征标就是增广特征标。
\end{proof}

\begin{definition}
设 $L/K$ 是局部域的有限 Galois 扩张，$G=\Gal(L/K)$，$G_i$ 为下编号分歧群，$g_i=|G_i|$，$f=[\kappa_L:\kappa_K]$。定义 Artin 类函数 $a_G$：若 $s\ne1$，
\[
  a_G(s)=-f\,i_G(s),
\]
其中 $i_G(s)$ 表示 $s$ 所在的分歧层深度；再用 $\langle 1_G,a_G\rangle=0$ 决定 $a_G(1)$。
\end{definition}

\begin{proposition}
Artin 类函数可写成
\[
  a_G=\sum_{i\ge0}\frac{g_0}{g_i}\operatorname{Ind}_{G_i}^{G}(u_{G_i}),
\]
其中 $u_{G_i}$ 是 $G_i$ 的增广特征标。因此对任意表示特征标 $\chi$，
\[
  \langle \chi,a_G\rangle\ge0.
\]
\end{proposition}

\begin{proof}
若 $s\ne1$ 且 $s\in G_k\setminus G_{k+1}$，则对 $i\le k$，诱导到 $G$ 的增广特征标在 $s$ 上贡献 $-g_0/g_i$；对 $i>k$，贡献为 $0$。乘以系数后总和为 $-f(k+1)$，这正是 $a_G(s)$ 的定义。$s=1$ 的值由与平凡特征标正交确定。

对表示 $\rho$ 的空间 $V$，Frobenius 互反给出
\[
  \langle \chi,\operatorname{Ind}_{G_i}^{G}u_{G_i}\rangle
  =
  \dim V-\dim V^{G_i}\ge0.
\]
加权求和得到非负性。
\end{proof}

\begin{theorem}
对有限 Galois 扩张 $L/K$，$a_G$ 是某个真正表示的特征标。
\end{theorem}

\begin{proof}
前一命题给出 $\langle\chi,a_G\rangle$ 是非负有理数。要证明它是整数，可用 Brauer 定理把任意不可约特征标写成若干子群上一维特征标诱导的整数线性组合。对子群上一维特征标，配对值由 Hasse-Arf 定理给出的整数分歧跳跃控制，因此是整数。于是 $a_G$ 在不可约特征标基底中的系数都是非负整数，故它是某个表示的特征标。
\end{proof}

\subsection{导子}

\begin{definition}
设 $E/F$ 是数域的有限 Galois 扩张，$G=\Gal(E/F)$。对有限素理想 $\mathfrak p$，取 $\mathfrak P|\mathfrak p$，令 $D_{\mathfrak P}$ 为分解群。局部 Artin 类函数诱导到 $G$ 后记为 $a_{\mathfrak p}$. 对 $G$ 的表示特征标 $\chi$，定义局部导子指数
\[
  f_{\mathfrak p}(\chi)=\langle\chi,a_{\mathfrak p}\rangle .
\]
Artin 导子为
\[
  \mathfrak f(\chi)=\prod_{\mathfrak p}\mathfrak p^{f_{\mathfrak p}(\chi)}.
\]
\end{definition}

\begin{proposition}
Artin 导子满足：
\[
  \mathfrak f(\chi+\chi')=\mathfrak f(\chi)\mathfrak f(\chi'),\qquad
  \mathfrak f(1)=1.
\]
若 $H\subset G$，$K=E^H$，$\psi$ 是 $H$ 的表示特征标，则
\[
  \mathfrak f(\operatorname{Ind}_{H}^{G}\psi,E/F)
  =
  \mathfrak d_{K/F}^{\psi(1)}
  N_{K/F}(\mathfrak f(\psi,E/K)).
\]
\end{proposition}

\begin{proof}
第一式来自内积对特征标加法的线性性。平凡表示在每个分歧群不变量空间都是全空间，因此与增广部分配对为零。

诱导公式逐素位检查。局部分解群限制到中间域后，Artin 类函数满足
\[
  a_G|_H=v_K(\mathfrak d_{K/F})r_H+f_{K/F}a_H,
\]
其中 $r_H$ 是正则特征标。与 $\psi$ 配对并用 Frobenius 互反，第一项给出判别式指数乘 $\psi(1)$，第二项给出中间扩张的导子指数，再把局部指数汇总为理想公式。
\end{proof}

\begin{corollary}
设 $E/F$ 是有限 Galois 扩张，则
\[
  \mathfrak d_{E/F}=\prod_{\chi}\mathfrak f(\chi)^{\chi(1)},
\]
乘积对 $G$ 的所有不可约特征标取。
\end{corollary}

\begin{proof}
取平凡子群的平凡表示并诱导到 $G$，得到正则表示。正则特征标分解为
\[
  r_G=\sum_{\chi}\chi(1)\chi .
\]
把导子的乘法性应用于这个分解，同时诱导公式说明正则表示的导子就是判别式。
\end{proof}

\subsection{L 函数结构}

\begin{definition}
设 $E/F$ 是有限 Galois 扩张，$G=\Gal(E/F)$，$\rho:G\to\GL(V)$ 是复表示。对有限素理想 $\mathfrak p$，取 $\mathfrak P|\mathfrak p$，$D_{\mathfrak P}$ 为分解群，$I_{\mathfrak P}$ 为惯性群，$\phi_{\mathfrak P}$ 为几何或算术 Frobenius 的固定约定下的生成元。局部 Artin 因子定义为
\[
  L_{\mathfrak p}(s,\rho)
  =
  \det\left(1-N\mathfrak p^{-s}\rho(\phi_{\mathfrak P})
  \mid V^{I_{\mathfrak P}}\right)^{-1}.
\]
Artin $L$ 函数为
\[
  L(s,\rho)=\prod_{\mathfrak p}L_{\mathfrak p}(s,\rho).
\]
\end{definition}

若换另一个 $\mathfrak P|\mathfrak p$，分解群、惯性群和 Frobenius 同时共轭。行列式在共轭下不变，所以局部因子定义良好。若 $\mathfrak p$ 无分歧，则 $I_{\mathfrak P}=1$，局部因子就是 $\det(1-N\mathfrak p^{-s}\rho(\phi_{\mathfrak P}))^{-1}$。

\begin{definition}
无穷素位给出 Gamma 因子 $L_\infty(s,\rho)$。再令
\[
  c(\rho)=|d_{F/\Q}|^{\dim\rho}N_{F/\Q}\mathfrak f(\rho).
\]
完全 Artin $L$ 函数定义为
\[
  \Lambda(s,\rho)=c(\rho)^{s/2}L_\infty(s,\rho)L(s,\rho).
\]
\end{definition}

\begin{theorem}
Artin $L$ 函数具有三条形式结构：
\[
  \Lambda(s,\rho_1\oplus\rho_2)=
  \Lambda(s,\rho_1)\Lambda(s,\rho_2),
\]
若 $\rho$ 通过商群膨胀到更大的 Galois 群，则 $\Lambda$ 不变；若 $K=E^H$ 且 $\psi$ 是 $H$ 的表示，则
\[
  \Lambda(s,\operatorname{Ind}_{H}^{G}\psi)=\Lambda(s,\psi).
\]
\end{theorem}

\begin{proof}
直和性质来自行列式分块相乘，导子和 Gamma 因子也按维数与特征标加法相乘。膨胀性质来自局部分解群与惯性群在商映射下对应，Frobenius 元映到中间扩张的 Frobenius 元，不变量空间保持一致。

诱导性质是核心。固定 $\mathfrak p$，把它在 $K$ 中分解为 $\mathfrak q_i$。诱导表示的空间由陪集分量组成，Frobenius 在这些分量上循环移动；每个循环长度等于剩余次数 $f(\mathfrak q_i/\mathfrak p)$。循环矩阵的行列式满足
\[
  \det(1-\phi t\mid \operatorname{Ind}W)
  =
  \prod_i \det(1-\phi_i^{f_i}t^{f_i}\mid W_i).
\]
这正好等于 $K$ 上各 $\mathfrak q_i$ 的局部因子乘积。导子部分由上一小节诱导公式给出，无穷 Gamma 因子按实、复嵌入的诱导分解逐项相同。
\end{proof}

\begin{theorem}
设 $E/F$ 是有限交换扩张，$\rho:\Gal(E/F)\to\C^\times$ 是非平凡一维表示，
$\chi$ 是其特征标。通过类域论的范剩符号，$\chi$ 对应一个 Hecke 特征标
\[
  \widetilde\chi:F_{\A}^{\times}/F^\times\longrightarrow S^1.
\]
若 $\mathfrak f$ 是扩张和特征标共同确定的导子，则除去有限多个分歧素位的修正因子后，
Artin $L$ 函数与 Hecke $L$ 函数的 Euler 因子相同；在完全化归一化下有
\[
  \Lambda(s,\rho)=\Lambda(s,\widetilde\chi).
\]
\end{theorem}

\begin{proof}
一维表示的空间是 $\C$。对有限素理想 $\mathfrak p$，取 $\mathfrak P|\mathfrak p$。
若 $\mathfrak p$ 在 $E/F$ 中无分歧，则惯性群 $I_{\mathfrak P}$ 平凡，于是
\[
  L_{\mathfrak p}(s,\rho)=
  (1-\chi(\operatorname{Frob}_{\mathfrak P})N\mathfrak p^{-s})^{-1}.
\]
类域论把未分歧素理想 $\mathfrak p$ 的理想类送到 Frobenius 元；因此
\[
  \widetilde\chi(\mathfrak p)=\chi(\operatorname{Frob}_{\mathfrak P})
\]
在所有未分歧且不除导子的素位成立。于是这些局部 Euler 因子逐项相等。

若 $\mathfrak p$ 分歧，则 Artin 局部因子取决于 $V^{I_{\mathfrak P}}$。一维情形中，
若 $\chi$ 在惯性群上非平凡，则 $V^{I_{\mathfrak P}}=0$，局部因子为 $1$；若它在惯性群
上平凡，则留下一个有限修正因子。Hecke 侧的导子定义正是把非平凡惯性作用记录进
$\mathfrak f$，所以有限分歧处的差别由同一组导子修正控制。

最后比较完全化因子。无穷素位的 Gamma 因子由一维表示在复共轭或实复嵌入上的
符号决定，这与对应 Hecke 特征标的无穷部分相同；有限导子由 Artin 导子和 Hecke
导子的一维相容性给出。因此完全 Artin $L$ 函数等于对应完全 Hecke $L$ 函数。
\end{proof}

\begin{remark}
这条定理说明前面花力气证明 Hecke $L$ 函数的函数方程不是孤立技巧。对一维
Galois 表示，Artin $L$ 函数已经被 Hecke 理论完全控制；对高维表示，Brauer 归纳把
高维特征标写成一维特征标的诱导组合，于是一般 Artin 函数方程可以从一维情形推出。
\end{remark}

\begin{theorem}
完全 Artin $L$ 函数可亚纯延拓到全平面，并满足
\[
  \Lambda(s,\rho)=W(\rho)\Lambda(1-s,\rho^\vee),
\qquad |W(\rho)|=1.
\]
\end{theorem}

\begin{proof}
由 Brauer 定理，$\rho$ 的特征标可写成
\[
  \chi_\rho=\sum_i n_i\operatorname{Ind}_{H_i}^{G}\chi_i,
\]
其中 $\chi_i$ 是子群 $H_i$ 的一维特征标，$n_i\in\Z$。形式结构给出
\[
  \Lambda(s,\rho)=\prod_i\Lambda(s,\chi_i)^{n_i}.
\]
一维情形通过类域论对应到 Hecke 特征标，前面已经证明 Hecke $L$ 函数的解析延拓和函数方程。有限个这样的函数取整数幂乘积，得到亚纯延拓。把每个一维函数方程相乘，根数相乘为 $W(\rho)$，对偶表示对应特征标取逆，于是得到所列函数方程。
\end{proof}

\begin{remark}
Artin 猜想说：若 $\rho$ 不可约且非平凡，则 $L(s,\rho)$ 应为整函数。上面的定理只给出亚纯性，因为 Brauer 分解中可能出现负系数，乘积表达式可能带来表面上的极点。证明这些极点抵消，是 Artin 猜想的内容。
\end{remark}

\section{习题详解}

\begin{exercise}
设 $F$ 是数域，$S$ 是包含所有无穷素位的有限集合。设 $\omega$ 是 $F_{\A}^{\times}$ 上非平凡特征标，且 $\omega|_{F^\times}=1$；对 $v\notin S$，$\omega_v$ 无分歧。证明
\[
  \prod_{v\notin S}(1-\omega_v(\pi_v)q_v^{-s})^{-1}
\]
在 $\Re(s)>1$ 绝对收敛；当 $s\to1$ 时收敛到有限值；若 $\omega^2$ 非平凡，则该极限非零。
\end{exercise}

\begin{proof}
记 $\sigma=\Re(s)$。因为 $\omega_v$ 在无分歧位置由 $\omega_v(\pi_v)$ 决定，且酉特征标满足 $|\omega_v(\pi_v)|=1$，有
\[
  \sum_{v\notin S}|\omega_v(\pi_v)q_v^{-s}|
  \le
  \sum_v q_v^{-\sigma}.
\]
右侧在 $\sigma>1$ 收敛，这是 Dedekind zeta Euler 乘积绝对收敛的标准判据。因此无穷乘积绝对收敛，并定义无零解析函数。

令 $P_S(s)$ 表示该乘积。它是删去有限因子的 Hecke $L$ 函数。非平凡 Hecke 特征标的 $L(s,\omega)$ 在 $s=1$ 没有极点；有限删去因子在 $s=1$ 有有限值。因此 $P_S(s)$ 在 $s\to1$ 时有有限极限。

再证明非零性。对 $\sigma>1$，绝对收敛的 Euler 乘积每个因子非零，所以 $P_S(s)\ne0$。若 $\omega^2$ 非平凡而 $P_S(1)=0$，则 $\log P_S(s)$ 在 $s=1$ 附近必须趋向 $-\infty$。但
\[
  \log P_S(s)=
  \sum_{v\notin S}\sum_{m\ge1}
  \frac{\omega_v(\pi_v)^m}{m}q_v^{-ms}.
\]
$m\ge2$ 的部分在 $s=1$ 绝对收敛；可能发散的只有 $m=1$ 部分。若该部分造成零点，则平方特征标的实部会强迫 $\sum_v\omega_v(\pi_v)^2q_v^{-s}$ 出现与平凡特征标相同的对数发散。因为 $\omega^2$ 非平凡，对应 Hecke $L$ 函数在 $s=1$ 无极点，这种发散不存在。故极限非零。
\end{proof}

\begin{exercise}
设 $F$ 是数域，$\omega$ 是 $F_{\A}^{\times}$ 的非平凡特征标，并且 $\omega|_{F^\times}=1$。证明 $L(1,\omega)\ne0$。
\end{exercise}

\begin{proof}
先把 $\omega$ 写成酉部分乘以 $|\cdot|_{\A}^{it}$ 和有限阶的无穷方向因子。若非紧参数 $t\ne0$，则 $s=1$ 对应 Hecke $L$ 函数在直线 $\Re(s)=1$ 上的非平凡点。Hecke 的零点排除定理说明非平凡特征标在 $\Re(s)=1$ 上无零点，因此 $L(1,\omega)\ne0$。

若 $\omega$ 为有限阶特征标，则类域论把它对应到某个有限交换扩张 $E/F$ 的 Galois 群的一维特征标。Dedekind zeta 函数分解为
\[
  \zeta_E(s)=\prod_{\chi\in\widehat{\Gal(E/F)}}L(s,\chi).
\]
左侧在 $s=1$ 有单极，右侧中平凡特征标因子是 $\zeta_F(s)$，也有单极。商
\[
  \frac{\zeta_E(s)}{\zeta_F(s)}
  =
  \prod_{\chi\ne1}L(s,\chi)
\]
在 $s=1$ 有有限非零极限；这是解析类数公式给出的正数比值。有限个因子的乘积非零，故每个非平凡一维因子在 $s=1$ 非零。于是 $L(1,\omega)\ne0$。
\end{proof}

\begin{exercise}
设 $f(x)=(e^x-1)^{-1}$。证明当 $x\to0^+$ 时 $f(x)=O(x^{-1})$；对任意 $b>0$，当 $x\to\infty$ 时 $f(x)=O(x^{-b})$。并证明
\[
  M(f,s)=\Gamma(s)\zeta(s)
\]
在 $\Re(s)>1$ 解析。
\end{exercise}

\begin{proof}
当 $x\to0^+$ 时，$e^x-1=x+x^2/2+\cdots$，所以存在 $\delta>0$，使 $0<x<\delta$ 时 $e^x-1\ge x/2$。于是
\[
  0<f(x)\le 2x^{-1}.
\]
这就是 $O(x^{-1})$。

当 $x\to\infty$ 时，$e^x-1\ge e^x/2$，所以 $f(x)\le2e^{-x}$。指数衰减压过任意幂衰减：对给定 $b>0$，函数 $x^be^{-x}$ 在 $[1,\infty)$ 有界，故 $e^{-x}=O(x^{-b})$，从而 $f(x)=O(x^{-b})$。

对 $\Re(s)>1$，
\[
  \frac1{e^x-1}=\sum_{n=1}^{\infty}e^{-nx}
\]
绝对收敛。于是
\[
  M(f,s)=\int_0^\infty \frac{x^s}{e^x-1}\frac{dx}{x}
  =
  \sum_{n=1}^{\infty}\int_0^\infty e^{-nx}x^{s-1}\,dx.
\]
变量替换 $u=nx$ 给出
\[
  \int_0^\infty e^{-nx}x^{s-1}\,dx=n^{-s}\Gamma(s).
\]
因此
\[
  M(f,s)=\Gamma(s)\sum_{n=1}^{\infty}n^{-s}=\Gamma(s)\zeta(s).
\]
级数在任意半平面 $\Re(s)\ge1+\epsilon$ 一致绝对收敛，$\Gamma(s)$ 解析，所以乘积在 $\Re(s)>1$ 解析。
\end{proof}

\begin{exercise}
设 $D(s)=\sum_{n\ge1}f(n)n^{-s}$，且 $f(mn)=f(m)f(n)$。证明
\[
  D(s)^{-1}=\sum_{n\ge1}\mu(n)f(n)n^{-s},
\qquad
  -\frac{D'(s)}{D(s)}=\sum_{n\ge1}\Lambda(n)f(n)n^{-s}.
\]
\end{exercise}

\begin{proof}
在绝对收敛区域，乘法性给出 Euler 乘积
\[
  D(s)=\prod_p(1-f(p)p^{-s})^{-1}.
\]
取倒数：
\[
  D(s)^{-1}=\prod_p(1-f(p)p^{-s}).
\]
展开乘积时，每个素数 $p$ 只能选 $1$ 或 $-f(p)p^{-s}$。因此只有无平方因子整数出现，系数为 $(-1)^r f(p_1\cdots p_r)=\mu(n)f(n)$，得到第一式。

对 Euler 乘积取对数：
\[
  \log D(s)=
  \sum_p\sum_{k\ge1}\frac{f(p)^k}{k}p^{-ks}.
\]
逐项求导得
\[
  -\frac{D'(s)}{D(s)}
  =
  \sum_p\sum_{k\ge1}(\log p)f(p)^k p^{-ks}.
\]
由于 $f(p^k)=f(p)^k$，且 Mangoldt 函数满足 $\Lambda(p^k)=\log p$、其他整数上为 $0$，上式就是
\[
  \sum_{n\ge1}\Lambda(n)f(n)n^{-s}.
\]
\end{proof}

\begin{exercise}
设 $\chi$ 是模 $q$ 的 Dirichlet 特征标。证明对 $\sigma>1$，
\[
  -\frac{L'}{L}(\sigma+it,\chi)
  =
  \sum_{n\ge1}\Lambda(n)n^{-\sigma}\chi(n)e^{-it\log n}.
\]
再用 $3+4\cos\theta+\cos2\theta=2(1+\cos\theta)^2\ge0$ 推出标准零点排除不等式，并说明 Dirichlet $L$ 函数的零点区域结论。
\end{exercise}

\begin{proof}
第一式由上一题应用于 $f(n)=\chi(n)$ 得到，因为 $n^{-s}=n^{-\sigma}e^{-it\log n}$。

令
\[
  z_n=\chi(n)e^{-it\log n}.
\]
当 $(n,q)=1$ 时 $|z_n|=1$，写 $z_n=e^{i\theta_n}$；当 $(n,q)>1$ 时 $\chi(n)=0$，该项不贡献。于是
\[
  3+4\Re z_n+\Re z_n^2\ge0.
\]
乘以非负权 $\Lambda(n)n^{-\sigma}$ 并求和，得到
\[
  3\sum_n\Lambda(n)n^{-\sigma}
  +4\Re\sum_n\Lambda(n)\chi(n)n^{-\sigma-it}
  +\Re\sum_n\Lambda(n)\chi(n)^2n^{-\sigma-2it}
  \ge0.
\]
把三项分别写成对数导数，就是
\[
3\left(-\frac{L'}{L}(\sigma,\chi_0)\right)
+4\left(-\Re\frac{L'}{L}(\sigma+it,\chi)\right)
+\left(-\Re\frac{L'}{L}(\sigma+2it,\chi^2)\right)\ge0.
\]

零点区域的思想如下。若 $L(s,\chi)$ 在 $s=\beta+it$ 有零点且 $\beta$ 很接近 $1$，则 $-L'/L(\sigma+it,\chi)$ 在 $\sigma>1$ 附近有大正贡献，大小约为 $(\sigma-\beta)^{-1}$。主特征标项 $-L'/L(\sigma,\chi_0)$ 的主要奇性来自 $\zeta(s)$ 在 $1$ 的极点，大小约为 $(\sigma-1)^{-1}$。上面的非负不等式迫使零点贡献不能超过极点贡献和有界误差。取
\[
  \sigma=1+\frac{A}{\log(q(|t|+2))}
\]
并估计其余 Euler 因子与 Gamma 因子，得到存在常数 $c>0$，使复特征标在
\[
  \sigma\ge 1-\frac{c}{\log(q(|t|+2))}
\]
无零点。实特征标时 $\chi^2=\chi_0$，不等式允许至多一个实零点靠近 $1$；若存在，它必须为单零点，因为二重零点会给对数导数两倍的不可抵消贡献。
\end{proof}

\begin{exercise}
设 $E/F$ 是数域的有限 Galois 扩张。证明 $\zeta_E(s)/\zeta_F(s)$ 是整函数。
\end{exercise}

\begin{proof}
设 $G=\Gal(E/F)$。正则表示 $r_G$ 分解为
\[
  r_G=1_G\oplus u_G,
\]
其中 $u_G$ 是增广表示。Artin 形式结构给出
\[
  \zeta_E(s)=L(s,r_G,E/F),\qquad
  \zeta_F(s)=L(s,1_G,E/F).
\]
因此
\[
  \frac{\zeta_E(s)}{\zeta_F(s)}=L(s,u_G,E/F).
\]
由 Brauer 定理，$u_G$ 的 Artin $L$ 函数可写为 Hecke $L$ 函数的有限乘积和商。Artin 整性定理在这个特殊情形给出无极点性；也可以用 Dedekind zeta 的类数公式检查 $s=1$ 处极点已经完全由平凡表示贡献，其他点的可能极点由函数方程转移并相互抵消。故该商在全平面解析。
\end{proof}

\begin{exercise}
设 $E/F$ 是数域的有限 Galois 扩张，$s_0$ 是 $\zeta_E(s)$ 的单零点。证明对 $G=\Gal(E/F)$ 的任意不可约表示 $\rho$，Artin $L$ 函数 $L(s,\rho,E/F)$ 在 $s=s_0$ 解析。
\end{exercise}

\begin{proof}
正则表示分解为
\[
  r_G=\bigoplus_{\rho\in\widehat G}(\dim\rho)\rho.
\]
于是
\[
  \zeta_E(s)=\prod_{\rho\in\widehat G}L(s,\rho)^{\dim\rho}.
\]
Artin $L$ 函数已知亚纯。若某个 $L(s,\rho)$ 在 $s_0$ 有极点，设极点阶为 $m>0$，则它在上式中贡献极点阶 $m\dim\rho$。其他因子只能贡献整数阶零点或极点。由于左侧在 $s_0$ 是单零点，总阶为 $+1$。极点贡献为负阶，必须由其他零点抵消，抵消后总阶要保持 $1$。

使用 Artin 形式结构的诱导关系，正则表示对应的 $\zeta_E$ 给出各不可约因子零点阶的加权和。若某个不可约因子在 $s_0$ 有负阶，则把该不可约特征标与适当诱导自一维特征标的 Brauer 关系配对，会构造出一个 Hecke $L$ 函数组合，使该组合在 $s_0$ 出现由负阶强迫的极点。另一方面，一维 Hecke 因子的解析性和 $\zeta_E$ 的单零点假设把总阶限制为 $1$，没有足够的正阶项去抵消这个被强迫的极点。Heilbronn 引理正是把这段阶数比较形式化：当 $\zeta_E$ 在 $s_0$ 为单零点时，所有不可约 Artin 因子在 $s_0$ 的阶数都非负。于是每个 $L(s,\rho)$ 在 $s_0$ 解析。
\end{proof}

\begin{exercise}
设 $F$ 是数域。证明当 $x>1$ 时
\[
  \prod_{\mathfrak p}(1-N\mathfrak p^{-x})^{-1}
\]
收敛，并证明当 $x\to+\infty$ 时该乘积趋于 $1$。
\end{exercise}

\begin{proof}
对 $0<u\le1/2$，有 $-\log(1-u)\le2u$。除有限多个素理想外，$N\mathfrak p^{-x}\le1/2$。因此乘积收敛等价于
\[
  \sum_{\mathfrak p}N\mathfrak p^{-x}
\]
收敛。该级数由 Dedekind zeta 的 Dirichlet 级数 $\sum_{\mathfrak a}N\mathfrak a^{-x}$ 控制，后者在 $x>1$ 收敛。

当 $x\to+\infty$ 时，
\[
  0\le \log\prod_{\mathfrak p}(1-N\mathfrak p^{-x})^{-1}
  \le
  C\sum_{\mathfrak p}N\mathfrak p^{-x}.
\]
右侧趋于 $0$。这是因为对任意固定下界 $M$，有限多个 $N\mathfrak p\le M$ 的项趋于 $0$；尾部由某个 $x_0>1$ 的收敛级数一致控制。故乘积的对数趋于 $0$，乘积趋于 $1$。
\end{proof}

\begin{exercise}
设 $S$ 是包含所有无穷素位的有限集，$|\theta_{\mathfrak p}|\le1$。证明
\[
  E(s)=\prod_{\mathfrak p\notin S}(1-\theta_{\mathfrak p}N\mathfrak p^{-s})^{-1}
\]
在 $\Re(s)>1$ 绝对收敛为非零解析函数，并且当 $\Re(s)\to+\infty$ 时，一致趋于 $1$。
\end{exercise}

\begin{proof}
令 $\sigma=\Re(s)$。由于 $|\theta_{\mathfrak p}N\mathfrak p^{-s}|\le N\mathfrak p^{-\sigma}$，在 $\sigma>1$ 时
\[
  \sum_{\mathfrak p\notin S}|\theta_{\mathfrak p}N\mathfrak p^{-s}|
  \le
  \sum_{\mathfrak p}N\mathfrak p^{-\sigma}<\infty .
\]
无穷乘积绝对收敛。每个局部因子非零，且对数级数
\[
  \log E(s)=\sum_{\mathfrak p\notin S}\sum_{m\ge1}
  \frac{\theta_{\mathfrak p}^m}{m}N\mathfrak p^{-ms}
\]
在紧子域 $\sigma\ge1+\epsilon$ 上一致收敛，因此 $E(s)$ 解析且无零。

一致趋于 $1$：给定 $\epsilon>0$，选 $A>1$ 使 $\sum_{\mathfrak p}N\mathfrak p^{-A}<\epsilon$。当 $\sigma\ge A$ 时，对数绝对值由常数倍的该尾和控制，且与 $\Im(s)$ 无关。因此 $\log E(s)\to0$ 一致，故 $E(s)\to1$ 一致。
\end{proof}

\begin{exercise}
设 $F/K$ 是数域扩张，$M$ 是 $F$ 的有限素位集合的一个子集。假设除有限多个例外外，对应素位的剩余次数 $f(\mathfrak p)>1$。若 $|\theta_{\mathfrak p}|\le1$，证明
\[
  \prod_{\mathfrak p\in M}(1-\theta_{\mathfrak p}N\mathfrak p^{-s})^{-1}
\]
在 $\Re(s)>1/2$ 绝对收敛为非零解析函数。
\end{exercise}

\begin{proof}
去掉有限多个例外不会影响收敛半平面和解析非零性，所以设所有 $\mathfrak p\in M$ 都满足 $f(\mathfrak p)>1$。令 $\mathfrak q$ 是 $K$ 中位于 $\mathfrak p$ 下方的素理想，则
\[
  N_F\mathfrak p=(N_K\mathfrak q)^{f(\mathfrak p)}\ge (N_K\mathfrak q)^2.
\]
对固定 $\mathfrak q$，其上方素理想个数至多 $[F:K]$。若 $\sigma=\Re(s)>1/2$，则
\[
  \sum_{\mathfrak p\in M}N_F\mathfrak p^{-\sigma}
  \le
  [F:K]\sum_{\mathfrak q}(N_K\mathfrak q)^{-2\sigma}.
\]
右侧在 $2\sigma>1$ 时收敛。于是
\[
  \sum_{\mathfrak p\in M}|\theta_{\mathfrak p}N\mathfrak p^{-s}|
\]
收敛，Euler 乘积绝对收敛。对数级数在任意 $\sigma\ge1/2+\epsilon$ 上一致收敛，因此乘积为无零解析函数。
\end{proof}
 
\backmatter
\chapter*{参考文献}
\addcontentsline{toc}{chapter}{参考文献}
\begin{thebibliography}{9}
\bibitem{Li} 黎景辉，《代数数论》。
\bibitem{Neukirch} J. Neukirch, \emph{Algebraic Number Theory}.
\bibitem{CasselsFrohlich} J. W. S. Cassels and A. Fröhlich (eds.), \emph{Algebraic Number Theory}.
\bibitem{SerreLocal} J.-P. Serre, \emph{Local Fields}.
\bibitem{SerreCohom} J.-P. Serre, \emph{Galois Cohomology}.
\bibitem{MilneCFT} J. S. Milne, \emph{Class Field Theory}.
\bibitem{Washington} L. C. Washington, \emph{Introduction to Cyclotomic Fields}.
\end{thebibliography}
\end{document}
